\documentclass[10pt,oneside]{book}

\usepackage[T1]{fontenc}
\usepackage[utf8]{inputenc}
\usepackage[french]{babel}
\usepackage{amsmath,amssymb,mathrsfs}
\usepackage[all]{xy}
\usepackage[a4paper,margin=22mm]{geometry}
\usepackage[
  colorlinks=true,
  linkcolor=blue,
  hypertexnames=false,
  unicode=true,
  pdfencoding=auto
]{hyperref}
\hypersetup{
  pdftitle={Éléments de géométrie algébrique IV, première partie -- édition française},
  pdfauthor={Alexander Grothendieck et Jean Dieudonné},
  pdfsubject={Transcription diplomatique française depuis l'autorité NUMDAM}
}

\renewcommand{\thesection}{\arabic{section}}
\renewcommand{\thesubsection}{\thesection.\arabic{subsection}}
\newcommand{\gr}{\operatorname{gr}}
\newcommand{\sref}[1]{(\ref{#1})}
\makeatletter
\newcommand{\setmanualcurrentlabel}[1]{\def\@currentlabel{#1}}
\makeatother
\newcommand{\oldpage}[2][]{%
  \setcounter{footnote}{0}%
  \marginpar{\scriptsize\textbf{#1}~|~#2}\ignorespaces}
\newenvironment{env}[1][]{\setmanualcurrentlabel{#1}\par\medskip\noindent\textbf{(#1)}\ }{\par}
\newenvironment{definition}[1][]{%
  \setmanualcurrentlabel{#1}\par\medskip\noindent\itshape Définition (#1). ---\ }{\par}
\newenvironment{proposition}[1][]{%
  \setmanualcurrentlabel{#1}\par\medskip\noindent\itshape Proposition (#1). ---\ }{\par}
\newenvironment{theorem}[1][]{%
  \setmanualcurrentlabel{#1}\par\medskip\noindent\itshape Théorème (#1). ---\ }{\par}
\newenvironment{corollary}[1][]{%
  \setmanualcurrentlabel{#1}\par\medskip\noindent\itshape Corollaire (#1). ---\ }{\par}
\newenvironment{remark}[1][]{%
  \setmanualcurrentlabel{#1}\par\medskip\noindent\itshape Remarque (#1). ---\ }{\par}
\newenvironment{remarks}[1][]{%
  \setmanualcurrentlabel{#1}\par\medskip\noindent\itshape Remarques (#1). ---\ }{\par}
\newenvironment{lemma}[1][]{%
  \setmanualcurrentlabel{#1}\par\medskip\noindent\itshape Lemme (#1). ---\ }{\par}
\newenvironment{example}[1][]{%
  \setmanualcurrentlabel{#1}\par\medskip\noindent\textit{Exemple (#1). ---}\ }{\par}
\newenvironment{examples}[1][]{%
  \setmanualcurrentlabel{#1}\par\medskip\noindent\textit{Exemples (#1). ---}\ }{\par}
\newenvironment{proof}{%
  \par\medskip\noindent\textit{Démonstration. ---\ }}{\par}

\makeatletter
\newcommand{\EGAReaderTarget}[2]{\hypertarget{#1}{}\@bsphack\begingroup\def\@currentHref{#1}\label{#2}\endgroup\@esphack}
\makeatother
\begin{document}
% BEGIN RIGHTS AND AI TRANSCRIPTION NOTICE
\clearpage
\begingroup
\small
\noindent\textbf{Droits, disponibilité et provenance}

Les droits sur les œuvres originales et sur toute contribution protégée à ces traductions et transcriptions dérivées restent ceux de leurs titulaires respectifs. Ce projet ne place pas ces œuvres ou contributions dans le domaine public et n'accorde aucune licence sur elles. Les mentions de droits de tiers et les autorisations déjà accordées restent applicables.

Ces éditions sont mises à disposition gratuitement, à des fins non commerciales et exclusivement pédagogiques, dans la présomption que les titulaires de droits ne s'opposeraient pas à la libre disponibilité de ces traductions et transcriptions aux côtés des originaux français sur NUMDAM (www.numdam.org/journals/PMIHES/). Cette présomption ne signifie pas qu'une autorisation a été obtenue ; le libre accès aux originaux n'autorise pas, à lui seul, leur traduction, leur adaptation ou leur republication.

Si vous détenez des droits sur un élément présenté ici et vous opposez à sa disponibilité, veuillez ouvrir une demande sur GitHub (github.com/KokunoYumeto/ega-fr/issues). Le responsable de la publication retirera l'élément concerné à la demande du titulaire des droits.

La transcription moderne, les corrections, la composition et les notes mathématiques sont produites par IA. Les travaux documentés comprennent ceux d'OpenAI Codex, GPT-6 Astra, effort Ultra. Les indications de provenance des contributions antérieures restent applicables. Les textes originaux restent l'œuvre d'Alexandre Grothendieck et Jean Dieudonné ; ils ne sont pas attribués à l'IA. Cette édition n'est approuvée ni par les auteurs, ni par l'IHÉS, ni par NUMDAM.
\par
\endgroup
\clearpage
% END RIGHTS AND AI TRANSCRIPTION NOTICE

%
% DEBUT SOURCE ega4/ega0-14-fr.tex
\clearpage
\thispagestyle{plain}
\markboth{PRÉLIMINAIRES}{PRÉLIMINAIRES}
\oldpage[IV-1]{5}
\phantomsection
\label{chapter:0IV1-fr}

\begin{center}
{\large CHAPITRE 0 (suite)\footnotemark\par}
\vspace{1.5em}
{\Large\bfseries PRÉLIMINAIRES\par}
\vspace{2em}
{\large\bfseries Sommaire\par}
\end{center}

\footnotetext{On rappelle que le symbole $0_N$ (resp. $\operatorname{Err}_N$)
réfère à la partie du chapitre 0 (resp. la liste d'Errata) jointe au chapitre
$N$.}

\begin{center}
\begin{tabular}{@{}r l@{}}
\S~14. & Dimension combinatoire d'un espace topologique.\\
\S~15. & Suites M-régulières et suites $\mathscr F$-régulières.\\
\S~16. & Dimension et profondeur dans les anneaux locaux noethériens.\\
\S~17. & Anneaux réguliers.\\
\S~18. & Compléments sur les extensions d'algèbres.\\
\S~19. & Algèbres formellement lisses et anneaux de Cohen.\\
\S~20. & Dérivations et différentielles.\\
\S~21. & Différentielles dans les anneaux de caractéristique $p$.\\
\S~22. & Critères différentiels de lissité formelle et de régularité.\\
\S~23. & Anneaux japonais.
\end{tabular}
\end{center}

La presque totalité des paragraphes précédents est consacrée à l'exposé de
notions d'Algèbre commutative qui seront utilisées au cours du chapitre IV.
Bien qu'une bonne partie de ces notions figure déjà dans plusieurs ouvrages
([1], [12], [13], [17], [30]), on a pensé qu'il serait plus commode pour le
lecteur d'en avoir une exposition suivie et à peu près indépendante. Joints
aux \S\S~5, 6 et 7 du chapitre IV (où est utilisé le langage des schémas), ces
paragraphes constituent, à l'intérieur de notre Traité, un petit Traité
spécial, à peu près indépendant des chapitres I à III, et qui vise à présenter
sous une forme cohérente les propriétés des anneaux qui «~se comportent bien~»
relativement à des opérations telles que la complétion ou la fermeture
intégrale, en rattachant systématiquement ces propriétés à quelques
conceptions générales\footnote{La plupart des propriétés dont il s'agit ont
été découvertes par Chevalley, Zariski, Nagata et Serre. La méthode utilisée
ici a d'abord été développée en automne 1961, dans un cours professé à Harvard
University, par A.~Grothendieck.}.

\setcounter{section}{13}
\section{Dimension combinatoire d'un espace topologique}
\label{section:0.14-fr}

\setcounter{subsection}{0}
\subsection{Dimension combinatoire d'un espace topologique}
\label{subsection:0.14.1-fr}

\oldpage[0\textsubscript{IV-1}]{6}
\begin{env}[14.1.1]
\label{0.14.1.1-fr}
Soit $I$ un ensemble ordonné~; une \emph{chaîne} d'éléments de $I$ est par
définition une suite finie strictement croissante
$i_0<i_1<\cdots<i_n$ d'éléments de $I$ ($n\geq 0$)~; par définition la
\emph{longueur} de cette chaîne est $n$. Si $X$ est un espace topologique,
l'ensemble de ses parties \emph{fermées irréductibles} est ordonné par
inclusion, d'où la notion de \emph{chaîne} de parties fermées irréductibles de
$X$.
\end{env}

\begin{definition}[14.1.2]
\label{0.14.1.2-fr}
Soit $X$ un espace topologique. On appelle \emph{dimension combinatoire de
$X$} (ou simplement \emph{dimension de $X$} s'il ne peut y avoir de
confusion) et on note $\dim_c(X)$, ou simplement $\dim(X)$, la borne
supérieure des longueurs des chaînes de parties fermées irréductibles de $X$.
Pour tout $x\in X$, on appelle \emph{dimension combinatoire de $X$ en $x$}
(ou simplement \emph{dimension de $X$ en $x$}) et on note $\dim_x(X)$ le
nombre $\inf_U(\dim(U))$, où $U$ parcourt l'ensemble des voisinages ouverts de
$x$ dans $X$.
\end{definition}

Il résulte de cette définition que l'on a
\[
  \dim(\varnothing)=-\infty
\]
(la borne supérieure dans $\overline{\mathbb R}$ de la partie vide étant
$-\infty$). Si $(X_\alpha)$ est la famille des composantes irréductibles de
$X$, on a
\[
  \dim(X)=\sup_\alpha\dim(X_\alpha)
  \tag{14.1.2.1}\label{0.14.1.2.1-fr}
\]
car toute chaîne de parties fermées irréductibles de $X$ est par définition
contenue dans une composante irréductible de $X$, et inversement ces
composantes sont fermées dans $X$, donc toute partie fermée irréductible d'un
$X_\alpha$ est une partie fermée irréductible de $X$.

\begin{definition}[14.1.3]
\label{0.14.1.3-fr}
On dit qu'un espace topologique $X$ est \emph{équidimensionnel} si toutes ses
composantes irréductibles ont même dimension (égale alors à $\dim(X)$ par
(14.1.2.1)).
\end{definition}

\begin{proposition}[14.1.4]
\label{0.14.1.4-fr}
\begin{enumerate}
  \item[\rm(i)] Pour toute partie $Y$ d'un espace topologique $X$, on a
  $\dim(Y)\leq\dim(X)$.
  \item[\rm(ii)] Si un espace topologique $X$ est réunion finie de parties
  fermées $X_i$, on a $\dim(X)=\sup_i\dim(X_i)$.
\end{enumerate}
\end{proposition}

\begin{proof}
Pour toute partie fermée irréductible $Z$ de $Y$, l'adhérence $\overline Z$ de
$Z$ dans $X$ est irréductible ($0_{\mathrm I}$, 2.1.2) et
$\overline Z\cap Y=Z$, d'où (i). D'autre part, si
$X=\bigcup_{i=1}^n X_i$, où les $X_i$ sont fermés, toute partie fermée
irréductible de $X$ est contenue dans un des $X_i$ ($0_{\mathrm I}$, 2.1.1),
et par suite toute chaîne de parties fermées irréductibles de $X$ est contenue
dans un des $X_i$~; d'où (ii).
\end{proof}

On déduit de (14.1.4, (i)) que pour tout $x\in X$, on peut aussi écrire
\[
  \dim_x(X)=\lim_U\dim(U)
  \tag{14.1.4.1}\label{0.14.1.4.1-fr}
\]
la limite étant prise suivant l'ensemble filtrant décroissant des voisinages
ouverts de $x$ dans $X$.

\oldpage[0\textsubscript{IV-1}]{7}
\begin{corollary}[14.1.5]
\label{0.14.1.5-fr}
Soient $X$ un espace topologique, $x$ un point de $X$, $U$ un voisinage de
$x$, $Y_i$ ($1\leq i\leq n$) des parties fermées de $U$ telles que
$x\in Y_i$ pour tout $i$ et que $U$ soit réunion des $Y_i$. Alors on a
\[
  \dim_x(X)=\sup_i\bigl(\dim_x(Y_i)\bigr).
  \tag{14.1.5.1}\label{0.14.1.5.1-fr}
\]
\end{corollary}

\begin{proof}
Il résulte de (14.1.4, (ii)) que l'on a
$\dim_x(X)=\inf_V(\sup_i(\dim(Y_i\cap V)))$ lorsque $V$ parcourt l'ensemble
des voisinages ouverts de $x$ contenus dans $U$~; de même on a
$\dim_x(Y_i)=\inf_V(\dim(Y_i\cap V))$ pour tout $i$. Le corollaire est donc
évident si
\[
  \sup_i\bigl(\dim_x(Y_i)\bigr)=+\infty~;
\]
dans le cas contraire, il y a un voisinage ouvert $V_0\subset U$ de $x$ tel
que $\dim(Y_i\cap V)=\dim_x(Y_i)$ pour $1\leq i\leq n$ et pour tout
$V\subset V_0$, d'où la conclusion.
\end{proof}

\begin{proposition}[14.1.6]
\label{0.14.1.6-fr}
Pour tout espace topologique $X$, on a
$\dim(X)=\sup_{x\in X}\dim_x(X)$.
\end{proposition}

\begin{proof}
Il résulte de la définition (14.1.2) et de la prop. (14.1.4) que
$\dim_x(X)\leq\dim(X)$ pour tout $x\in X$. D'autre part, soit
$Z_0\subset Z_1\subset\cdots\subset Z_n$ une chaîne de parties fermées
irréductibles de $X$, et soit $x\in Z_0$~; pour tout ouvert $U\subset X$
contenant $x$, $U\cap Z_i$ est irréductible ($0_{\mathrm I}$, 2.1.6) et
fermé dans $U$, et comme on a $\overline{U\cap Z_i}=Z_i$ dans $X$, les
$U\cap Z_i$ sont distincts deux à deux~; donc on a $\dim(U)\geq n$, ce qui
achève la démonstration.
\end{proof}

\begin{corollary}[14.1.7]
\label{0.14.1.7-fr}
Si $(X_\alpha)$ est un recouvrement ouvert de $X$, ou un recouvrement fermé
localement fini de $X$, on a
$\dim(X)=\sup_\alpha(\dim(X_\alpha))$.
\end{corollary}

\begin{proof}
Si $X_\alpha$ est un voisinage de $x\in X$, on a
$\dim_x(X)\leq\dim(X_\alpha)$, d'où la première assertion. D'autre part, si
les $X_\alpha$ sont fermés et si $U$ est un voisinage de $x\in X$ qui ne
rencontre qu'un nombre fini d'ensembles $X_\alpha$, on a
\[
  \dim_x(X)\leq\dim(U)=\sup_\alpha\bigl(\dim(U\cap X_\alpha)\bigr)
  \leq\sup_\alpha\bigl(\dim(X_\alpha)\bigr)
\]
d'après (14.1.4), d'où la seconde assertion.
\end{proof}

\begin{corollary}[14.1.8]
\label{0.14.1.8-fr}
Soit $X$ un espace de Kolmogoroff ($0_{\mathrm I}$, 2.1.3) noethérien, et
soit $F$ l'ensemble des points fermés de $X$. Alors
$\dim(X)=\sup_{x\in F}\dim_x(X)$.
\end{corollary}

\begin{proof}
Avec les notations de la démonstration de (14.1.6), il suffit de remarquer
qu'il existe dans $Z_0$ un point fermé ($0_{\mathrm I}$, 2.1.3).
\end{proof}

\begin{proposition}[14.1.9]
\label{0.14.1.9-fr}
Soit $X$ un espace de Kolmogoroff noethérien non vide. Pour que
$\dim(X)=0$, il faut et il suffit que $X$ soit fini et discret.
\end{proposition}

\begin{proof}
Si un espace $X$ est séparé (et \emph{a fortiori} si $X$ est un espace
discret), les seuls ensembles fermés irréductibles de $X$ sont réduits à un
point, donc $\dim(X)=0$. Inversement, supposons que $X$ soit un espace de
Kolmogoroff noethérien tel que $\dim(X)=0$~; comme toute composante
irréductible de $X$ contient un point fermé ($0_{\mathrm I}$, 2.1.3), elle
doit être réduite à ce point. Comme $X$ n'a qu'un nombre fini de composantes
irréductibles, il est fini et discret.
\end{proof}

\begin{corollary}[14.1.10]
\label{0.14.1.10-fr}
Soit $X$ un espace de Kolmogoroff noethérien. Pour qu'un point $x\in X$ soit
isolé, il faut et il suffit que $\dim_x(X)=0$.
\end{corollary}

\begin{proof}
La condition est évidemment nécessaire (sans hypothèse sur $X$). Elle est
suffisante,
\oldpage[0\textsubscript{IV-1}]{8}
car il en résulte que $\dim(U)=0$ pour un voisinage ouvert $U$ de $x$, et
comme $U$ est un espace de Kolmogoroff noethérien, $U$ est fini et discret.
\end{proof}

\begin{proposition}[14.1.11]
\label{0.14.1.11-fr}
La fonction $x\mapsto\dim_x(X)$ est semi-continue supérieurement dans $X$.
\end{proposition}

\begin{proof}
Il est clair que cette fonction est semi-continue supérieurement en tout
point où elle vaut $+\infty$. Supposons donc $\dim_x(X)=n<+\infty$~; alors,
la formule (14.1.4.1) montre qu'il existe un voisinage ouvert $U_0$ de $x$ tel
que $\dim(U)=n$ pour tout voisinage ouvert $U\subset U_0$ de $x$. Cela étant,
pour tout $y\in U_0$ et tout voisinage ouvert $V\subset U_0$ de $y$, on a
$\dim(V)\leq\dim(U_0)=n$ (14.1.4)~; on déduit donc de (14.1.4.1) que
$\dim_y(X)\leq n$.
\end{proof}

\begin{remark}[14.1.12]
\label{0.14.1.12-fr}
Si $X$, $Y$ sont deux espaces topologiques, et $f:X\to Y$ une application
continue, on peut avoir $\dim(f(X))>\dim(X)$~; on en a un exemple en prenant
pour $X$ un espace discret à 2 éléments $a$, $b$, pour $Y$ l'ensemble
$\{a,b\}$ muni de la topologie pour laquelle les ensembles fermés sont
$\varnothing$, $\{a\}$ et $\{a,b\}$~; si $f:X\to Y$ est l'application
identique, on a $\dim(Y)=1$ et $\dim(X)=0$. On notera que $Y$ est le spectre
d'un anneau de valuation discrète $A$, dont $a$ est l'unique point fermé et
$b$ le point générique~; si $K$ et $k$ sont le corps des fractions et le
corps résiduel de $A$, $X$ est le spectre de l'anneau $k\times K$ et $f$
l'application continue correspondant à l'homomorphisme
$(\varphi,\psi):A\to k\times K$, où $\varphi:A\to k$ et $\psi:A\to K$ sont
les homomorphismes canoniques (cf. IV, 5.4.3).
\end{remark}

\subsection{Codimension d'une partie fermée}
\label{subsection:0.14.2-fr}

\begin{definition}[14.2.1]
\label{0.14.2.1-fr}
Étant donnée une partie fermée irréductible $Y$ d'un espace topologique $X$,
on appelle \emph{codimension combinatoire} (ou simplement \emph{codimension})
de $Y$ dans $X$, et on note $\operatorname{codim}(Y,X)$ la borne supérieure
des longueurs des chaînes de parties fermées irréductibles de $X$, dont $Y$
est le plus petit élément. Si $Y$ est une partie fermée quelconque de $X$, on
appelle \emph{codimension de $Y$ dans $X$}, et on note encore
$\operatorname{codim}(Y,X)$ la borne inférieure des codimensions dans $X$ des
composantes irréductibles de $Y$. On dit que $X$ est
\emph{équicodimensionnel} si toutes les parties fermées irréductibles minimales
de $X$ ont même codimension dans $X$.
\end{definition}

Il résulte de cette définition que
$\operatorname{codim}(\varnothing,X)=+\infty$, la borne inférieure de la
partie vide de $\overline{\mathbb R}$ étant $+\infty$. Si $Y$ est fermé dans
$X$ et si $(X_\alpha)$ (resp. $(Y_\beta)$) est la famille des composantes
irréductibles de $X$ (resp. $Y$), tout $Y_\beta$ est contenu dans un
$X_\alpha$, et plus généralement toute chaîne de parties fermées
irréductibles de $X$ dont $Y_\beta$ est le plus petit élément est formée de
parties d'un $X_\alpha$~; on a donc
\[
  \operatorname{codim}(Y,X)=
  \inf_\beta\left(\sup_\alpha
  \bigl(\operatorname{codim}(Y_\beta,X_\alpha)\bigr)\right)
  \tag{14.2.1.1}\label{0.14.2.1.1-fr}
\]
où pour chaque $\beta$, $\alpha$ parcourt l'ensemble des indices tels que
$Y_\beta\subset X_\alpha$.

\begin{proposition}[14.2.2]
\label{0.14.2.2-fr}
Soit $X$ un espace topologique.
\begin{enumerate}
  \item[\rm(i)] Si $\Phi$ est l'ensemble des parties fermées irréductibles de
  $X$, on a
  \[
    \dim(X)=\sup_{Y\in\Phi}\bigl(\operatorname{codim}(Y,X)\bigr).
    \tag{14.2.2.1}\label{0.14.2.2.1-fr}
  \]
  \oldpage[0\textsubscript{IV-1}]{9}
  \item[\rm(ii)] Pour toute partie fermée non vide $Y$ de $X$, on a
  \[
    \dim(Y)+\operatorname{codim}(Y,X)\leq\dim(X).
    \tag{14.2.2.2}\label{0.14.2.2.2-fr}
  \]
  \item[\rm(iii)] Si $Y$, $Z$, $T$ sont trois parties fermées de $X$ telles
  que $Y\subset Z\subset T$, on a
  \[
    \operatorname{codim}(Y,Z)+\operatorname{codim}(Z,T)
    \leq\operatorname{codim}(Y,T).
    \tag{14.2.2.3}\label{0.14.2.2.3-fr}
  \]
  \item[\rm(iv)] Pour qu'une partie fermée $Y$ de $X$ soit telle que
  $\operatorname{codim}(Y,X)=0$, il faut et il suffit que $Y$ contienne une
  composante irréductible de $X$.
\end{enumerate}
\end{proposition}

\begin{proof}
Les assertions (i) et (iv) sont des conséquences immédiates de la définition
(14.2.1). Pour démontrer (ii), on peut se borner au cas où $Y$ est
irréductible, et alors la formule résulte des définitions (14.1.1) et
(14.2.1). Enfin, pour démontrer (iii), on peut, en vertu de la définition
(14.2.1), se borner d'abord au cas où $Y$ est irréductible~; alors
$\operatorname{codim}(Y,Z)=
\sup_\alpha(\operatorname{codim}(Y,Z_\alpha))$ pour les composantes
irréductibles $Z_\alpha$ de $Z$ contenant $Y$~; il est clair que
$\operatorname{codim}(Y,T)\geq\operatorname{codim}(Y,Z)$, donc l'inégalité
est vraie si $\operatorname{codim}(Y,Z)=+\infty$~; sinon, il existe un
$\alpha$ tel que
$\operatorname{codim}(Y,Z)=\operatorname{codim}(Y,Z_\alpha)$ et en vertu de
(14.2.1), on peut se borner au cas où $Z$ est lui aussi irréductible~; mais
alors l'inégalité (14.2.2.3) est conséquence évidente de la définition
(14.2.1).
\end{proof}

\begin{proposition}[14.2.3]
\label{0.14.2.3-fr}
Soient $X$ un espace topologique, $Y$ une partie fermée de $X$. Pour tout
ouvert $U$ dans $X$, on a
\[
  \operatorname{codim}(Y\cap U,U)\geq\operatorname{codim}(Y,X).
  \tag{14.2.3.1}\label{0.14.2.3.1-fr}
\]
En outre, pour que les deux membres de (14.2.3.1) soient égaux, il faut et il
suffit que, si $(Y_\alpha)$ est la famille des composantes irréductibles de
$Y$ rencontrant $U$, on ait
$\operatorname{codim}(Y,X)=
\inf_\alpha(\operatorname{codim}(Y_\alpha,X))$.
\end{proposition}

\begin{proof}
On sait ($0_{\mathrm I}$, 2.1.6) que $Z\mapsto\overline Z$ est une bijection
de l'ensemble des parties fermées irréductibles de $U$ sur l'ensemble des
parties fermées irréductibles de $X$ rencontrant $U$, et en particulier fait
correspondre aux composantes irréductibles de $Y\cap U$ les composantes
irréductibles de $Y$ rencontrant $U$~; si $Y_\alpha$ est une de ces
dernières, on a donc
$\operatorname{codim}(Y_\alpha,X)=
\operatorname{codim}(Y_\alpha\cap U,U)$, et la proposition résulte alors de
la définition (14.2.1).
\end{proof}

\begin{definition}[14.2.4]
\label{0.14.2.4-fr}
Soient $X$ un espace topologique, $Y$ une partie fermée de $X$, $x$ un point
de $X$. On appelle \emph{codimension de $Y$ dans $X$ au point $x$} et on note
$\operatorname{codim}_x(Y,X)$ le nombre
$\sup_U(\operatorname{codim}(Y\cap U,U))$, où $U$ parcourt l'ensemble des
voisinages ouverts de $x$ dans $X$.
\end{definition}

En vertu de (14.2.3), on peut aussi écrire
\[
  \operatorname{codim}_x(Y,X)=
  \lim_U\bigl(\operatorname{codim}(Y\cap U,U)\bigr)
  \tag{14.2.4.1}\label{0.14.2.4.1-fr}
\]
la limite étant prise suivant l'ensemble filtrant décroissant des voisinages
ouverts de $x$ dans $X$. On notera que l'on a
\[
  \operatorname{codim}_x(Y,X)=+\infty\quad\text{si}\quad x\in X-Y.
\]

\oldpage[0\textsubscript{IV-1}]{10}
\begin{proposition}[14.2.5]
\label{0.14.2.5-fr}
Si $(Y_i)_{1\leq i\leq n}$ est une famille finie de parties fermées d'un
espace topologique $X$, et $Y$ la réunion de cette famille, on a
\[
  \operatorname{codim}(Y,X)=
  \inf_i\bigl(\operatorname{codim}(Y_i,X)\bigr).
  \tag{14.2.5.1}\label{0.14.2.5.1-fr}
\]
\end{proposition}

\begin{proof}
En effet, toute composante irréductible de l'un des $Y_i$ est contenue dans
une composante irréductible de $Y$, et inversement toute composante
irréductible de $Y$ est aussi composante irréductible de l'un des $Y_i$
($0_{\mathrm I}$, 2.1.1)~; la conclusion résulte donc de la définition
(14.2.1) et de l'inégalité (14.2.2.3).
\end{proof}

\begin{corollary}[14.2.6]
\label{0.14.2.6-fr}
Soient $X$ un espace topologique, $Y$ un sous-espace fermé localement
noethérien de $X$.
\begin{enumerate}
  \item[\rm(i)] Pour tout $x\in X$, il n'existe qu'un nombre fini de
  composantes irréductibles $Y_i$ ($1\leq i\leq n$) de $Y$ contenant $x$,
  et l'on a
  $\operatorname{codim}_x(Y,X)=
  \inf_i(\operatorname{codim}(Y_i,X))$.
  \item[\rm(ii)] La fonction
  $x\mapsto\operatorname{codim}_x(Y,X)$ est semi-continue inférieurement
  dans $X$.
\end{enumerate}
\end{corollary}

\begin{proof}
En effet, par hypothèse il y a un voisinage ouvert $U_0$ de $x$ dans $X$ tel
que $Y\cap U_0$ soit noethérien, donc n'ait qu'un nombre fini de composantes
irréductibles, qui sont traces sur $U_0$ de composantes irréductibles de $Y$~;
\emph{a fortiori} il n'y a qu'un nombre fini de composantes irréductibles
$Y_i$ ($1\leq i\leq n$) de $Y$ contenant $x$, et on peut, en remplaçant $U_0$
par un voisinage ouvert $U\subset U_0$ de $x$ ne rencontrant aucun des $Y_j$
qui ne contiennent pas $x$, supposer que les $Y_i\cap U$ sont les composantes
irréductibles de $Y\cap U$~; pour tout voisinage ouvert $V\subset U$ de $x$
dans $X$, les $Y_i\cap V$ sont alors les composantes irréductibles de
$Y\cap V$, et (14.2.3) montre alors que
$\operatorname{codim}(Y_i,X)=\operatorname{codim}(Y_i\cap V,V)$, ce qui
prouve (i). En outre, pour tout $x'\in U$, les composantes irréductibles de
$Y$ contenant $x'$ sont certaines des $Y_i$, donc
$\operatorname{codim}_{x'}(Y,X)\geq\operatorname{codim}_x(Y,X)$, ce qui
démontre (ii).
\end{proof}

\subsection{La condition des chaînes}
\label{subsection:0.14.3-fr}

\begin{env}[14.3.1]
\label{0.14.3.1-fr}
Dans un espace topologique $X$, nous dirons qu'une chaîne
$Z_0\subset Z_1\subset\cdots\subset Z_n$ de parties fermées irréductibles est
\emph{saturée} s'il n'existe pas de partie fermée irréductible $Z'$ distincte
des $Z_i$ et telle que $Z_k\subset Z'\subset Z_{k+1}$ pour un indice $k$.
\end{env}

\begin{proposition}[14.3.2]
\label{0.14.3.2-fr}
Soit $X$ un espace topologique tel que, pour deux parties fermées
irréductibles quelconques $Y$, $Z$ de $X$ telles que $Y\subset Z$, on ait
$\operatorname{codim}(Y,Z)<+\infty$. Les deux conditions suivantes sont
équivalentes~:
\begin{enumerate}
  \item[\rm(a)] Deux chaînes saturées de parties fermées irréductibles de $X$,
  ayant mêmes extrémités, ont même longueur.
  \item[\rm(b)] Si $Y$, $Z$, $T$ sont trois parties fermées irréductibles de
  $X$ telles que $Y\subset Z\subset T$, on a
  \[
    \operatorname{codim}(Y,T)=
    \operatorname{codim}(Y,Z)+\operatorname{codim}(Z,T).
    \tag{14.3.2.1}\label{0.14.3.2.1-fr}
  \]
\end{enumerate}
\end{proposition}

\begin{proof}
Il est immédiat que a) entraîne b). Réciproquement, supposons b) vérifiée, et
démontrons que si deux chaînes saturées de mêmes extrémités ont pour longueurs
$m$ et $n\leq m$, on a nécessairement $m=n$. Raisonnons par récurrence sur
$n$, la proposition étant évidente pour $n=1$. Supposons donc $n>1$, $n<m$,
et soit $Z_0\subset Z_1\subset\cdots\subset Z_n$ une
\oldpage[0\textsubscript{IV-1}]{11}
chaîne saturée telle qu'il existe une autre chaîne saturée d'extrémités
$Z_0$, $Z_n$ et de longueur $m$. Comme
$\operatorname{codim}(Z_0,Z_n)\geq m>n$ et que
$\operatorname{codim}(Z_0,Z_1)=1$, il résulte de b) que
$\operatorname{codim}(Z_1,Z_n)=\operatorname{codim}(Z_0,Z_n)-1>n-1$, ce qui
contredit l'hypothèse de récurrence.
\end{proof}

Lorsque les conditions de (14.3.2) sont remplies, on dit que $X$ satisfait à
la \emph{condition des chaînes}, ou encore est un \emph{espace caténaire}. Il
est clair que tout sous-espace fermé d'un espace caténaire est caténaire.

\begin{proposition}[14.3.3]
\label{0.14.3.3-fr}
Soit $X$ un espace de Kolmogoroff noethérien de dimension finie. Les
conditions suivantes sont équivalentes~:
\begin{enumerate}
  \item[\rm(a)] Deux chaînes maximales de parties fermées irréductibles de $X$
  ont même longueur.
  \item[\rm(b)] $X$ est équidimensionnel, équicodimensionnel, et caténaire.
  \item[\rm(c)] $X$ est équidimensionnel, et pour deux parties fermées
  irréductibles $Y$, $Z$ de $X$ telles que $Y\subset Z$, on a
  \[
    \dim(Z)=\dim(Y)+\operatorname{codim}(Y,Z).
    \tag{14.3.3.1}\label{0.14.3.3.1-fr}
  \]
  \item[\rm(d)] $X$ est équicodimensionnel et pour deux parties fermées
  irréductibles $Y$, $Z$ de $X$ telles que $Y\subset Z$, on a
  \[
    \operatorname{codim}(Y,X)=
    \operatorname{codim}(Y,Z)+\operatorname{codim}(Z,X).
    \tag{14.3.3.2}\label{0.14.3.3.2-fr}
  \]
\end{enumerate}
\end{proposition}

\begin{proof}
Les hypothèses sur $X$ entraînent que les extrémités d'une chaîne maximale de
parties fermées irréductibles de $X$ sont nécessairement un point fermé et
une composante irréductible de $X$ ($0_{\mathrm I}$, 2.1.3)~; en outre, toute
chaîne saturée d'extrémités $Y$, $Z$ (avec $Y\subset Z$) est contenue dans une
chaîne maximale dont les éléments autres que ceux de la chaîne donnée, ou
bien sont contenus dans $Y$ ou bien contiennent $Z$. Ces remarques
établissent aussitôt l'équivalence de a) et b), et prouvent aussi que si a)
est vérifiée, on a, pour toute partie fermée irréductible $Y$ de $X$
\[
  \dim(Y)+\operatorname{codim}(Y,X)=\dim(X)~;
  \tag{14.3.3.3}\label{0.14.3.3.3-fr}
\]
comme (14.3.2.1) a lieu, on déduit aussitôt (14.3.3.1) et (14.3.3.2) de
(14.3.3.3). Inversement, (14.3.3.1) entraîne (14.3.2.1), donc (14.3.3.1)
entraîne la condition des chaînes en vertu de (14.3.2)~; en outre, en
appliquant (14.3.3.1) au cas où $Y$ est réduit à un point fermé $x$ de $X$ et
$Z$ à une composante irréductible de $X$, on obtient
$\operatorname{codim}(\{x\},X)=\dim(Z)$~; on en conclut que c) entraîne b).
De même, (14.3.3.2) entraîne (14.3.2.1), donc la condition des chaînes~; en
outre, avec le même choix de $Y$ et $Z$ que ci-dessus, (14.3.3.2) entraîne
encore $\operatorname{codim}(\{x\},X)=\dim(Z)$, donc (toute composante
irréductible de $X$ contenant un point fermé en vertu de ($0_{\mathrm I}$,
2.1.3)), d) entraîne b).
\end{proof}

On dit qu'un espace de Kolmogoroff noethérien est \emph{biéquidimensionnel}
s'il est de \emph{dimension finie} et s'il vérifie les conditions équivalentes
de (14.3.3).

\begin{corollary}[14.3.4]
\label{0.14.3.4-fr}
Soit $X$ un espace de Kolmogoroff noethérien biéquidimensionnel~; alors, pour
tout point fermé $x$ de $X$ et toute composante irréductible $Z$ de $X$, on a
\[
  \dim(X)=\dim(Z)=\operatorname{codim}(\{x\},X)=\dim_x(X).
  \tag{14.3.4.1}\label{0.14.3.4.1-fr}
\]
\end{corollary}

\oldpage[0\textsubscript{IV-1}]{12}
\begin{proof}
La dernière égalité résulte de ce que si
$Y_0=\{x\}\subset Y_1\subset\cdots\subset Y_m$ est une chaîne maximale de
parties fermées irréductibles de $X$ et $U$ un voisinage ouvert de $x$, les
$U\cap Y_i$ sont des parties fermées irréductibles de $U$ deux à deux
distinctes (puisque $\overline{U\cap Y_i}=Y_i$), d'où
$\dim(U)=\dim(X)$ en vertu de (14.1.4).
\end{proof}

\begin{corollary}[14.3.5]
\label{0.14.3.5-fr}
Soit $X$ un espace de Kolmogoroff noethérien~; si $X$ est
biéquidimensionnel, il en est de même de toute réunion de composantes
irréductibles de $X$ et de toute partie fermée irréductible de $X$. En outre,
pour toute partie fermée $Y$ de $X$, on a alors
\[
  \dim(Y)+\operatorname{codim}(Y,X)=\dim(X).
  \tag{14.3.5.1}\label{0.14.3.5.1-fr}
\]
\end{corollary}

\begin{proof}
Toute chaîne de parties fermées irréductibles de $X$ étant contenue dans une
composante irréductible de $X$, la première assertion découle aussitôt de
(14.3.3). En outre, si $X'$ est une partie fermée irréductible de $X$, $X'$
vérifie trivialement les conditions (14.3.3, c)), d'où la seconde assertion.

Enfin, pour démontrer (14.3.5.1), remarquons que l'on a vu dans la
démonstration de (14.3.3) que cette relation est vérifiée lorsque $Y$ est
irréductible~; si $Y_i$ ($1\leq i\leq m$) sont les composantes irréductibles
de $Y$, celle des $Y_i$ pour laquelle $\dim(Y_i)$ est le plus grand est aussi
celle pour laquelle $\operatorname{codim}(Y_i,X)$ est le plus petit~; donc
(14.3.5.1) résulte des définitions de $\dim(Y)$ et de
$\operatorname{codim}(Y,X)$.
\end{proof}

\begin{remark}[14.3.6]
\label{0.14.3.6-fr}
Le lecteur remarquera que la démonstration de (14.3.2) s'applique à un
ensemble ordonné quelconque, le fait qu'il s'agit de parties fermées
irréductibles d'un espace topologique n'intervenant pas dans cette
démonstration. Il en est de même de la démonstration de (14.3.3) dans un
ensemble ordonné $E$ tel que pour tout $x\in E$, il existe un $z\leq x$ qui
soit \emph{minimal} dans $E$, et où la longueur des chaînes d'éléments de $E$
est bornée.
\end{remark}

% FIN SOURCE ega4/ega0-14-fr.tex

%
% DEBUT SOURCE ega4/ega0-15-fr.tex
\section{Suites \texorpdfstring{$M$}{M}-régulières et suites
\texorpdfstring{$\mathscr F$}{F}-régulières}
\label{section:0.15-fr}

\subsection{Suites \texorpdfstring{$M$}{M}-régulières et suites
\texorpdfstring{$M$}{M}-quasi-régulières}
\label{subsection:0.15.1-fr}
\label{0.15.1-fr}

\begin{env}[15.1.1]
\label{0.15.1.1-fr}
\oldpage[0\textsubscript{IV-1}]{12}
Soient $A$ un anneau, $N$ un $A$-module~; dans ce qui suit, nous noterons
$N[T_1,\ldots,T_n]$, pour abréger, le $A$-module gradué
$N\otimes_A A[T_1,\ldots,T_n]$ ($T_i$ indéterminées), où $N$ est gradué par la
graduation triviale et $A[T_1,\ldots,T_n]$ par le degré total.

Soient $M$ un $A$-module, $f_i$ ($1\leq i\leq n$) des éléments de $A$,
$\mathfrak{J}=\sum_{i=1}^n f_iA$ l'idéal de $A$ engendré par les $f_i$, et
munissons $M$ de la filtration $\mathfrak{J}$-préadique~; on définit alors un
\emph{homomorphisme gradué surjectif de $A$-modules gradués}, de degré $0$,
\[
  \varphi : (\gr_0(M))[T_1,\ldots,T_n]\longrightarrow\gr_\bullet(M)
  \tag{15.1.1.1}\label{0.15.1.1.1-fr}
\]
de la façon suivante~: si $\xi\in\gr_0(M)$ est la classe mod.
$\mathfrak{J}M$ d'un élément $x\in M$, au couple formé de $\xi$ et d'un
polynôme $P[T_1,\ldots,T_n]$ homogène de degré $k$, on fait correspondre la
classe mod. $\mathfrak{J}^{k+1}M$ de $P(f_1,\ldots,f_n)x$~; il est immédiat
que cette classe
\oldpage[0\textsubscript{IV-1}]{13}
ne dépend que de $\xi$ et de $P$, et si $S_k$ est l'ensemble des polynômes
homogènes de degré $k$ dans $A[T_1,\ldots,T_n]$, on a ainsi défini une
application $A$-linéaire
\[
  \varphi_k:\gr_0(M)\otimes_A S_k\longrightarrow\gr_k(M),
\]
qui est surjective par définition de $\mathfrak{J}^kM$.

Nous dirons dans ce qui suit qu'un élément $f\in A$ est \emph{non diviseur de
zéro dans un $A$-module $M$} si l'homothétie $f_M:x\mapsto fx$ de $M$ est
\emph{injective}.
\end{env}

\begin{lemma}[15.1.2]
\label{0.15.1.2-fr}
Soient $M$ un $A$-module, $f$ un élément de $A$, et munissons $M$ de la
filtration $(fA)$-préadique. Si $f$ est non diviseur de zéro dans $M$,
l'homomorphisme canonique défini dans (15.1.1)~:
\[
  \varphi:(M/fM)[T]\longrightarrow\gr_\bullet(M)
  \tag{15.1.2.1}\label{0.15.1.2.1-fr}
\]
est bijectif. La réciproque est vraie si $M$ est séparé pour la topologie
$(fA)$-préadique.
\end{lemma}

\begin{proof}
On a ici $\gr_k(M)=f^kM/f^{k+1}M$~; dire que $\varphi$ est injective signifie
que pour tout $x\in M$ et tout $k$, la relation $f^kx\in f^{k+1}M$ implique
$x\in fM$~; il est clair qu'il en est bien ainsi si $f$ est non diviseur de
zéro dans $M$. Inversement, remarquons que l'homothétie
$f_M:x\mapsto fx$ définit par passage aux quotients, des homomorphismes
$\gr_k(M)\to\gr_{k+1}(M)$, donc un homomorphisme gradué $g$ de degré $1$ de
$\gr_\bullet(M)$ dans lui-même~; et dire que $\varphi$ est injectif exprime
encore que $g$ est injectif. Si on suppose la filtration de $M$ séparée, il en
résulte que $f_M$ est injective (Bourbaki, \emph{Alg.~comm.}, chap.~III,
\S~2, no~8, cor.~1 du th.~1).
\end{proof}

\begin{corollary}[15.1.3]
\label{0.15.1.3-fr}
Soient $A$ un anneau noethérien, $M$ un $A$-module de type fini. Pour que $f$
soit non diviseur de zéro dans $M$, il faut et il suffit que $\varphi$ soit
bijectif.
\end{corollary}

\begin{proof}
On doit seulement démontrer que la condition est suffisante~; pour prouver
que $f_M$ est injective, il suffit de montrer que pour tout idéal premier
$\mathfrak p$ de $A$, l'homothétie de rapport
$f'=f/1\in A'=A_{\mathfrak p}$ dans le $A_{\mathfrak p}$-module
$M_{\mathfrak p}=M'$ est injective (Bourbaki, \emph{Alg.~comm.}, chap.~II,
\S~3, no~4, th.~1). Or on a
$f'^kM'/f'^{k+1}M'=(f^kM/f^{k+1}M)\otimes_A A'$, et il est immédiat que
l'homomorphisme canonique
$\varphi':(M'/f'M')[T]\to\gr_\bullet(M')$ est égal à
$\varphi\otimes 1_{A'}$~; si $\varphi$ est injectif, il en est donc de même de
$\varphi'$ ($0_{\mathrm I}$, 1.3.2). On est donc ramené au cas où $A$ est un
anneau noethérien \emph{local}, dont nous désignerons par $\mathfrak m$
l'idéal maximal. Si $f\notin\mathfrak m$, alors $f$ est inversible et il n'y a
rien à démontrer. Si $f\in\mathfrak m$, on sait que $M$ est séparé pour la
topologie $(fA)$-préadique ($0_{\mathrm I}$, 7.3.5), et il suffit donc
d'appliquer (15.1.2).
\end{proof}

\begin{definition}[15.1.4]
\label{0.15.1.4-fr}
Soient $A$ un anneau, $M$ un $A$-module. On dit qu'un élément $f\in A$ est
\emph{$M$-régulier} (resp. \emph{$M$-quasi-régulier}) si $f$ est non diviseur
de zéro dans $M$ (resp. si l'homomorphisme (15.1.2.1) est bijectif).
\end{definition}

Nous venons donc de démontrer que tout élément $M$-régulier est
$M$-quasi-régulier, et que la réciproque est vraie si $A$ est noethérien et
$M$ un $A$-module de type fini. Lorsqu'on ne suppose plus $M$ de type fini,
cette réciproque peut être en défaut~: par exemple, si $A$ est un anneau
principal qui n'est pas un corps, $K$ son corps des fractions et $M$ le
$A$-module $K/A$, on a $fM=M$ pour tout $f\neq0$ dans $A$, donc les deux
membres de (15.1.2.1) sont réduits à $0$, mais $f$ est alors diviseur de zéro
dans $M$.

\begin{env}[15.1.5]
\label{0.15.1.5-fr}
\oldpage[0\textsubscript{IV-1}]{14}
Supposons maintenant $M$ muni d'une \emph{filtration} $(M_k)$ telle que
$M_0=M$, et désignons par
$\gr_\bullet(M)=\bigoplus_{k\geq0}M_k/M_{k+1}$ le $A$-module gradué associé.
Les notations étant celles de (15.1.1), posons, pour tout $k\geq0$,
\[
  M'_k=M_k+\mathfrak{J}M_{k-1}+\cdots+
  \mathfrak{J}^{k-1}M_1+\mathfrak{J}^kM_0.
  \tag{15.1.5.1}\label{0.15.1.5.1-fr}
\]
Il est clair que $(M'_k)$ est une filtration sur $M$. Lorsque
$M_k=\mathfrak{Q}^kM$, où $\mathfrak Q$ est un idéal de $A$, la filtration
$(M'_k)$ n'est autre que la filtration
$(\mathfrak{J}+\mathfrak{Q})$-préadique. Nous noterons
$\gr'_\bullet(M)$ le $A$-module gradué associé à $(M'_k)$. Notons encore
\[
  (\gr_\bullet(M)\otimes_A(A/\mathfrak{J}))[T_1,\ldots,T_n]
\]
le produit tensoriel de $A$-modules gradués ($\mathrm{II}$, 2.1.2)
\[
  (\gr_\bullet(M)\otimes_A(A/\mathfrak{J}))
  \otimes_A A[T_1,\ldots,T_n].
\]
On définit un \emph{homomorphisme gradué surjectif de $A$-modules gradués}, de
degré $0$,
\[
  \psi:(\gr_\bullet(M)\otimes_A(A/\mathfrak{J}))[T_1,\ldots,T_n]
  \longrightarrow\gr'_\bullet(M)
  \tag{15.1.5.2}\label{0.15.1.5.2-fr}
\]
de la façon suivante~: si $\alpha\in A/\mathfrak{J}$ est la classe d'un
élément $a\in A$, $\xi\in M_k/M_{k+1}$ la classe d'un élément $x\in M_k$ et
$P[T_1,\ldots,T_n]$ un polynôme homogène de degré $h$, on fait correspondre au
triplet $(\alpha,\xi,P)$ la classe mod. $M'_{h+k+1}$ de
$P(f_1,\ldots,f_n)ax$~; on vérifie encore aussitôt que cet élément ne dépend
bien que de $\xi$, $\alpha$ et $P$, et cela définit l'homomorphisme gradué
(15.1.5.2), qui est surjectif en vertu de la définition (15.1.5.1). On
retrouve l'homomorphisme (15.1.1.1) en considérant la filtration sur $M$ telle
que $M_1=0$.
\end{env}

\begin{lemma}[15.1.6]
\label{0.15.1.6-fr}
On garde les notations de (15.1.5) et on suppose $n=1$, en posant $f=f_1$. Si
$f$ est $\gr_\bullet(M)$-régulier, l'homomorphisme (15.1.5.2) est bijectif.
\end{lemma}

\begin{proof}
Soit $Q_k$ (resp. $Q'_k$) le sous-$A$-module des termes de degré $k$ au
premier (resp. second) membre de (15.1.5.2). Munissons $Q_k$ de la filtration
\[
  (Q_k)_i=\sum_{j\leq k-i}
  (\gr_{k-j}(M)\otimes_A(A/fA))T^j
  \tag{15.1.6.1}\label{0.15.1.6.1-fr}
\]
de sorte que $(Q_k)_0=Q_k$ et $(Q_k)_{k+1}=0$, et pour cette filtration, on a
\[
  \gr_i(Q_k)=(\gr_i(M)\otimes_A(A/fA))T^{k-i}.
\]
Munissons $Q'_k$ de la filtration formée des
$\psi((Q_k)_i)=(Q'_k)_i$. Comme ces filtrations sont \emph{finies}, il suffit
de prouver que les homomorphismes
$\psi_{ki}:\gr_i(Q_k)\to\gr_i(Q'_k)$ déduits de $\psi_k$ sont injectifs
(Bourbaki, \emph{Alg.~comm.}, chap.~III, \S~3, no~8, cor.~1 du th.~1). Or on
a
\[
  \gr_i(Q_k)=((M_i/M_{i+1})\otimes_A(A/fA))T^{k-i}~;
\]
d'autre part, $(Q'_k)_{i+1}$ est l'image de
\[
  M_k+fM_{k-1}+\cdots+f^{k-i-1}M_{i+1}
\]
\oldpage[0\textsubscript{IV-1}]{15}
dans $M'_k/M'_{k+1}$. On est donc ramené à écrire que, pour $x\in M_i$, la
relation
\[
  f^{k-i}x\in M_k+fM_{k-1}+\cdots+
  f^{k-i-1}M_{i+1}+M'_{k+1}
  \tag{15.1.6.2}\label{0.15.1.6.2-fr}
\]
implique $x\in fM_i+M_{i+1}$. Or, le second membre de (15.1.6.2) est contenu
dans $M_{i+1}+M'_{k+1}$, et en vertu de (15.1.5.1), $M'_{k+1}$ est contenu
dans $M_{i+1}+f^{k+1-i}M$. L'hypothèse sur $f$ entraîne que $f$ est non
diviseur de zéro dans $M/M_h$ pour tout $h$ (Bourbaki, \emph{loc.~cit.}). La
relation $f^{k-i}x\in f^{k+1-i}M+M_{i+1}$ entraîne donc (en raisonnant dans
$M/M_{i+1}$) l'existence d'un $y\in M$ tel que $x-fy\in M_{i+1}$, et comme
$x\in M_i$, on a $fy\in M_i$, d'où (en raisonnant dans $M/M_i$), $y\in M_i$,
et finalement $x\in fM_i+M_{i+1}$. C.Q.F.D.
\end{proof}

\begin{definition}[15.1.7]
\label{0.15.1.7-fr}
Soit $(f_i)_{1\leq i\leq n}$ une suite d'éléments de $A$. On dit qu'elle est
\emph{$M$-régulière} si, pour $1\leq i\leq n$, $f_i$ est
$\left(M/(\sum_{1\leq j\leq i-1}f_jM)\right)$-régulier. On dit que la suite
$(f_i)$ est \emph{$M$-quasi-régulière} si
l'homomorphisme canonique $\varphi$ défini dans (15.1.1.1) est bijectif.

Lorsque $M=A$, on dit simplement «~suite régulière~» et
«~suite quasi-régulière~».
\end{definition}

\begin{proposition}[15.1.8]
\label{0.15.1.8-fr}
Les notations étant celles de (15.1.5), on suppose que la suite
$(f_i)_{1\leq i\leq n}$ est $\gr_\bullet(M)$-régulière. Alors
l'homomorphisme canonique (15.1.5.2) est bijectif.
\end{proposition}

\begin{proof}
La proposition n'est autre que (15.1.6) pour $n=1$~; raisonnons par récurrence
sur $n$. Posons
\[
  \mathfrak{J}''=\sum_{i=1}^{n-1}f_iA,
\]
de sorte que $\mathfrak{J}=\mathfrak{J}''+f_nA$, et soit
\[
  M''_k=M_k+\mathfrak{J}''M_{k-1}+\cdots+
  \mathfrak{J}''^{k-1}M_1+\mathfrak{J}''^kM_0~;
\]
on peut donc écrire
\[
  M'_k=M''_k+f_nM''_{k-1}+\cdots+f_n^{k-1}M''_1+f_n^kM''_0.
\]
Notons $\gr''_\bullet(M)$ le $A$-module gradué associé à la filtration
$(M''_k)$ de $M$. On peut écrire, à des isomorphismes canoniques près
\[
  (\gr_\bullet(M)\otimes(A/\mathfrak{J}))\otimes_A A[T_1,\ldots,T_n]
  =
  (\gr_\bullet(M)\otimes(A/\mathfrak{J}''))[T_1,\ldots,T_{n-1}]
  \otimes(A/f_nA)\otimes_A A[T_n]
\]
et par suite (15.1.5.2) se factorise en
\[
  (\gr''_\bullet(M)\otimes(A/f_nA))[T_n]\longrightarrow\gr'_\bullet(M)
  \tag{15.1.8.1}\label{0.15.1.8.1-fr}
\]
et
\[
  ((\gr_\bullet(M)\otimes(A/\mathfrak{J}''))[T_1,\ldots,T_{n-1}])
  \otimes(A/f_nA)\otimes_A A[T_n]
  \xrightarrow{\psi'\otimes1}
  \gr''_\bullet(M)\otimes(A/f_nA)\otimes_A A[T_n]
\]
$\psi'$ étant l'application (15.1.5.2) où l'on remplace $n$ par $n-1$,
$\mathfrak{J}$ par $\mathfrak{J}''$ et $\gr'_\bullet(M)$ par
$\gr''_\bullet(M)$. L'hypothèse de récurrence entraîne que $\psi'$ est
bijective et il reste donc à voir qu'il en est de même de (15.1.8.1). En
appliquant (15.1.6), il suffit de montrer que $f_n$ est
$\gr''_\bullet(M)$-régulier. Mais $\gr''_\bullet(M)$ s'identifie à
\[
  (\gr_\bullet(M)\otimes(A/\mathfrak{J}''))
  [T_1,\ldots,T_{n-1}]~;
\]
par hypothèse, l'homothétie de rapport $f_n$ dans
$\gr_\bullet(M)\otimes(A/\mathfrak{J}'')$ est bijective~; il en est par suite
de même de l'homothétie de rapport $f_n$ dans $\gr''_\bullet(M)$, puisque
$A[T_1,\ldots,T_{n-1}]$ est un $A$-module libre.
\end{proof}

\begin{proposition}[15.1.9]
\label{0.15.1.9-fr}
Soient $M$ un $A$-module, $(f_i)_{1\leq i\leq n}$ une suite d'éléments de $A$.
Si la suite $(f_i)$ est $M$-régulière, elle est $M$-quasi-régulière. La
réciproque est vraie lorsqu'on suppose en outre que les $A$-modules
\[
  M,\quad M/f_1M,\quad\ldots,\quad
  M/\left(\sum_{1\leq j\leq n-1}f_jM\right)
\]
sont séparés pour la topologie $\mathfrak{J}$-préadique (où
$\mathfrak{J}=\sum_{i=1}^n f_iA$).
\end{proposition}

\begin{proof}
La première assertion résulte de (15.1.8), où l'on prend $M_1=0$, et par suite
l'homomorphisme (15.1.5.2) se réduit à (15.1.1.1). Inversement, supposons que
l'homomorphisme $\varphi$ de (15.1.1.1) soit bijectif. Notons que si on pose
\[
  \mathfrak{J}''=\sum_{i=1}^{n-1}f_iA,
\]
\oldpage[0\textsubscript{IV-1}]{16}
et si on désigne par $\gr''_\bullet(M)$ le $A$-module gradué associé à $M$
filtré par la filtration $\mathfrak{J}''$-préadique, $\varphi$ se factorise en
\[
  (\gr''_\bullet(M)\otimes(A/f_nA))[T_n]\longrightarrow\gr_\bullet(M)
  \tag{15.1.9.1}\label{0.15.1.9.1-fr}
\]
et en $\varphi'\otimes1$, où $\varphi'$ est l'homomorphisme canonique
\[
  (\gr''_0(M))[T_1,\ldots,T_{n-1}]
  \longrightarrow\gr''_\bullet(M).
\]
Comme ces deux homomorphismes sont surjectifs, ils sont bijectifs si
$\varphi$ est supposé bijectif, donc $\varphi'$ est nécessairement injectif,
et par suite bijectif puisqu'il est surjectif. On peut alors raisonner par
récurrence sur $n$ (le cas $n=1$ résultant de (15.1.2)), car sur
\[
  M,\quad M/f_1M,\quad\ldots,\quad
  M/\left(\sum_{j=1}^{n-2}f_jM\right)
\]
la topologie $\mathfrak{J}''$-préadique est plus fine que la topologie
$\mathfrak{J}$-préadique, et par suite est séparée. On voit donc que la suite
$(f_i)_{1\leq i\leq n-1}$ est $M$-régulière. D'autre part, montrons que
l'hypothèse entraîne que si $y\in M/\mathfrak{J}''M$ est tel que
\[
  f_n^ky\in f_n^{k+1}(M/\mathfrak{J}''M),
\]
alors on a $y\in f_n(M/\mathfrak{J}''M)$. En effet, si $x$ est un représentant
de $y$, montrons que l'on a $f_n^kx\in\mathfrak{J}^{k+1}M$~; notre assertion
résultera alors du fait que (15.1.9.1) est injectif. L'hypothèse entraîne
\[
  f_n^kx\in f_n^{k+1}M+\mathfrak{J}''M\subset\mathfrak{J}M,
\]
d'où, puisque $\varphi$ est injectif,
$f_n^{k-1}x\in\mathfrak{J}^kM$~; par récurrence, on obtient ainsi
$x\in\mathfrak{J}M$, d'où finalement
$f_n^kx\in\mathfrak{J}^{k+1}M$. On termine le raisonnement comme dans
(15.1.2)~: sur $M/\mathfrak{J}''M$, la filtration
$\mathfrak{J}$-préadique est identique à la filtration
$(f_nA)$-préadique, donc est séparée par hypothèse, et on conclut que
l'homothétie de rapport $f_n$ dans $M/\mathfrak{J}''M$ est injective, ce qui
achève de montrer que la suite $(f_i)_{1\leq i\leq n}$ est $M$-régulière.
\end{proof}

\begin{corollary}[15.1.10]
\label{0.15.1.10-fr}
Avec les notations de (15.1.9), supposons que pour tout $A$-module $N$ de type
fini, $M\otimes_A N$ soit séparé pour la topologie
$\mathfrak{J}$-préadique. Alors, pour que la suite $(f_i)$ soit
$M$-régulière, il faut et il suffit qu'elle soit $M$-quasi-régulière.
\end{corollary}

On notera que l'hypothèse de (15.1.10) entraîne que, si la suite
$(f_i)_{1\leq i\leq n}$ est $M$-régulière, il en est de même de la suite
$(f_{\sigma(i)})$ pour toute permutation $\sigma$ de l'intervalle $[1,n]$.

\begin{corollary}[15.1.11]
\label{0.15.1.11-fr}
Soient $A$ un anneau noethérien, $M$ un $A$-module de type fini,
$(f_i)_{1\leq i\leq n}$ une suite finie d'éléments contenus dans le radical de
$A$. Alors, pour que la suite $(f_i)$ soit $M$-régulière, il faut et il suffit
qu'elle soit $M$-quasi-régulière.
\end{corollary}

\begin{proof}
C'est un cas particulier de (15.1.10), puisque tout $A$-module de type fini
est alors séparé pour la topologie $\mathfrak{J}$-préadique
($0_{\mathrm I}$, 7.3.5).
\end{proof}

\begin{remarks}[15.1.12]
\label{0.15.1.12-fr}
\mbox{}
\begin{enumerate}
  \item[(i)] Lorsque $A$ est un anneau local noethérien, $M$ un $A$-module de
  type fini, et lorsque les $f_i$ sont contenus dans l'idéal maximal de $A$,
  la notion de suite $M$-régulière a été introduite par J.-P.~Serre sous le
  nom de $M$-\emph{suite} [17].
  \item[(ii)] Même si $A$ est noethérien et $M=A$, une suite $(f,g)$ de deux
  éléments de $A$ qui est quasi-régulière n'est pas nécessairement régulière.
  Par exemple, la suite $(0,1)$ est quasi-régulière, puisqu'alors
  $\mathfrak{J}=fA+gA=A$, donc $\gr_\bullet(A)=0$~; mais si $A\neq0$ elle
  n'est pas régulière~; on notera toutefois que dans ce cas la suite $(1,0)$
  est régulière, car $1$ n'est pas diviseur de $0$ dans $A$, et comme
  $A/A=0$, $0$ n'est pas diviseur de zéro dans ce module (cf. (15.1.1)).
  \item[(iii)]
  \oldpage[0\textsubscript{IV-1}]{17}
  Dans (15.1.6), on ne peut en général remplacer l'hypothèse que $f$ est
  $\gr_\bullet(M)$-régulier par l'hypothèse qu'il est
  $\gr_\bullet(M)$-quasi-régulier. Pour le voir, reprenons l'exemple (15.1.4)
  où $A$ est un anneau principal de corps des fractions $K\neq A$ et prenons
  $M=K$, la filtration $(M_k)$ de $M$ étant définie par
  $M_0=M=K$, $M_1=A$, $M_k=0$ pour $k\geq2$, d'où
  $\gr_0(M)=K/A$, $\gr_1(M)=A$, $\gr_k(M)=0$ pour $k\geq2$. Si $f\neq0$ dans
  $A$, on a vu que $f$ n'est pas $\gr_\bullet(M)$-régulier, mais comme on a
  \[
    f^k\gr_\bullet(M)/f^{k+1}\gr_\bullet(M)=f^kA/f^{k+1}A,
  \]
  il est immédiat que l'homomorphisme (15.1.2.1) (où $M$ est remplacé par
  $\gr_\bullet(M)$) est injectif, donc $f$ est
  $\gr_\bullet(M)$-quasi-régulier. Mais avec les notations de (15.1.6), on a
  alors $M'_k=K$ pour tout $k\geq0$, donc $\gr'_\bullet(M)=0$, tandis que le
  premier membre de (15.1.5.2) n'est pas réduit à $0$.
\end{enumerate}
\end{remarks}

\begin{proposition}[15.1.13]
\label{0.15.1.13-fr}
Soient $A$ un anneau, $M$, $N$ deux $A$-modules, $(f_i)_{1\leq i\leq n}$ une
suite d'éléments de $A$. Si la suite $(f_i)$ est $M$-régulière et si $N$ est
un $A$-module plat, alors la suite $(f_i)$ est
$M\otimes_A N$-régulière. La réciproque est vraie si $N$ est un $A$-module
fidèlement plat.
\end{proposition}

\begin{proof}
En effet (sans hypothèse sur $N$),
\[
  (M\otimes_A N)/
  \left(\sum_{j=1}^{i-1}f_j(M\otimes_A N)\right)
\]
est isomorphe à
\[
  \left(M/\left(\sum_{j=1}^{i-1}f_jM\right)\right)\otimes_A N~;
\]
il suffit alors d'appliquer ($0_{\mathrm I}$, 6.1.1) et
($0_{\mathrm I}$, 6.4.1) à l'homothétie de rapport $f_i$ dans
$M/(\sum_{j=1}^{i-1}f_jM)$.
\end{proof}

\begin{corollary}[15.1.14]
\label{0.15.1.14-fr}
Soient $A$ un anneau, $M$ un $A$-module, $(f_i)_{1\leq i\leq n}$ une suite
finie d'éléments de $A$. Soient $\rho:A\to B$ un homomorphisme d'anneaux,
$M'=M_{(B)}=M\otimes_A B$, $f'_i=\rho(f_i)$ pour $1\leq i\leq n$. Si $B$ est
un $A$-module plat et si la suite $(f_i)$ est $M$-régulière, alors la suite
$(f'_i)$ est $M'$-régulière~; la réciproque est vraie si $B$ est un
$A$-module fidèlement plat.
\end{corollary}

\begin{proof}
En effet,
\[
  M'/\left(\sum_{j=1}^{i-1}f'_jM'\right)
\]
s'identifie à
\[
  \left(M/\left(\sum_{j=1}^{i-1}f_jM\right)\right)\otimes_A B,
\]
et l'homothétie de rapport $f'_i$ dans le $B$-module
$M'/(\sum_{j=1}^{i-1}f'_jM')$ à l'homothétie de rapport $f_i$ dans le
$A$-module
\[
  \left(M/\left(\sum_{j=1}^{i-1}f_jM\right)\right)\otimes_A B~;
\]
l'assertion résulte donc de (15.1.13).
\end{proof}

\begin{proposition}[15.1.15]
\label{0.15.1.15-fr}
Soient $A$ un anneau, $B$ une $A$-algèbre commutative,
$(f_i)_{1\leq i\leq n}$ une suite finie d'éléments de $B$, $M$ un
$B$-module. Supposons que la suite $(f_i)$ soit $M$-régulière et que chacun
des $A$-modules
\[
  M/\left(\sum_{j=1}^{i-1}f_jM\right)\quad(1\leq i\leq n)
\]
soit plat. Alors, pour tout homomorphisme d'anneaux $\rho:A\to A'$, si l'on
pose $B'=B\otimes_A A'$, $M'=M\otimes_A A'$ et $f'_i=f_i\otimes1$
($1\leq i\leq n$), la suite $(f'_i)$ est $M'$-régulière et les
$A'$-modules
\[
  M'/\left(\sum_{j=1}^{i-1}f'_jM'\right)\quad(1\leq i\leq n)
\]
sont plats.
\end{proposition}

\begin{proof}
Considérons la suite (exacte par hypothèse)
\[
  0\longrightarrow M\xrightarrow{f_1}M\longrightarrow M/f_1M
  \longrightarrow0
\]
(la flèche $M\xrightarrow{f_1}M$ étant l'homothétie de rapport $f_1$ dans
$M$). L'hypothèse que $M/f_1M$ est $A$-plat entraîne que la suite
\[
  0\longrightarrow M\otimes_A A'
  \xrightarrow{f_1\otimes1}M\otimes_A A'
  \longrightarrow(M/f_1M)\otimes_A A'\longrightarrow0
\]
\oldpage[0\textsubscript{IV-1}]{18}
est exacte ($0_{\mathrm I}$, 6.1.2). Cela montre que l'homothétie de rapport
$f'_1$ dans $M'$ est injective.

Posons ensuite, pour $1\leq i\leq n-1$,
\[
  M_i=M/\left(\sum_{j=1}^i f_jM\right),\qquad
  M'_i=M'/\left(\sum_{j=1}^i f'_jM'\right)~;
\]
on a $M'_i=M_i\otimes_A A'$, $M_i/f_{i+1}M_i=M_{i+1}$,
$M'_i/f'_{i+1}M'_i=M'_{i+1}$~; il suffit de remplacer dans le raisonnement
précédent $M$ et $f_1$ par $M_i$ et $f_{i+1}$ pour conclure que l'homothétie
de rapport $f'_{i+1}$ dans $M'_i$ est injective, en raison de la platitude de
$M_i/f_{i+1}M_i$. La dernière assertion résulte de ($0_{\mathrm I}$, 6.2.1).
\end{proof}

\begin{proposition}[15.1.16]
\label{0.15.1.16-fr}
Soient $A$, $B$ deux anneaux locaux noethériens, $k$ le corps résiduel de
$A$, $\varphi:A\to B$ un homomorphisme local, $M$ un $B$-module de type fini.
Soit $(f_i)_{1\leq i\leq n}$ une suite d'éléments de l'idéal maximal de $B$,
et pour tout $i$, soit $g_i$ l'image de $f_i$ dans $B\otimes_A k$. Les
conditions suivantes sont équivalentes~:
\begin{enumerate}
  \item[\rm(a)] La suite $(f_i)$ est $M$-régulière et les quotients
  \[
    M_i=M/\left(\sum_{j=1}^i f_jM\right)
  \]
  sont des $A$-modules plats pour $1\leq i\leq n$.
  \item[\rm(b)] La suite $(f_i)$ est $M$-régulière et le $A$-module $M_n$ est
  plat.
  \item[\rm(c)] Le $A$-module $M$ est plat et la suite
  $(g_i)_{1\leq i\leq n}$ est $(M\otimes_A k)$-régulière.
  \item[\rm(d)] Le $A$-module $M$ est plat, et pour tout homomorphisme
  d'anneaux $\rho:A\to A'$, si on pose $B'=B\otimes_A A'$,
  $M'=M\otimes_A A'$, $f'_i=f_i\otimes1$ ($1\leq i\leq n$), la suite
  $(f'_i)$ est $M'$-régulière.
\end{enumerate}
\end{proposition}

\begin{proof}
Il est clair que a) entraîne b)~; pour montrer que b) entraîne d), il suffit,
en vertu de (15.1.15), de montrer que b) entraîne que les $M_i$ sont des
$A$-modules plats pour $0\leq i\leq n-1$ (en posant $M_0=M$). Comme
$M_{i+1}=M_i/f_{i+1}M_i$, on est ramené, par récurrence descendante, à prouver
que si $h$ est un élément de l'idéal maximal de $B$ tel que l'homothétie de
rapport $h$ dans $M$ soit injective et que $M/hM$ soit un $A$-module plat,
alors $M$ est un $A$-module plat~; mais cela n'est autre que
($0_{\mathrm{III}}$, 10.2.7). Il est clair que c) est le cas particulier de
d) obtenu en prenant $A'=k$. Reste donc à prouver que c) entraîne a). Comme
$M$ est un $A$-module plat et un $B$-module de type fini, il résulte de c) et
de ($0_{\mathrm{III}}$, 10.2.4) que l'homothétie de rapport $f_1$ dans $M$
est injective et que son conoyau $M/f_1M=M_1$ est un $A$-module plat.
Raisonnons par récurrence sur $i$ et supposons démontré que $M_i$ est un
$A$-module plat~; comme l'homothétie de rapport $g_{i+1}$ dans
$M_i\otimes_A k$ est injective par hypothèse, il résulte encore de
($0_{\mathrm{III}}$, 10.2.4) que l'homothétie de rapport $f_{i+1}$ dans
$M_i$ est injective et que $M_i/f_{i+1}M_i=M_{i+1}$ est un $A$-module plat,
ce qui achève la démonstration. On notera que la démonstration du fait que c)
entraîne a) n'utilise pas le fait que les $f_i$ appartiennent à l'idéal
maximal de $B$.
\end{proof}

\begin{corollary}[15.1.17]
\label{0.15.1.17-fr}
Les hypothèses sur $A$ et $B$ étant celles de (15.1.16), soient $M$ un
$B$-module de type fini, plat sur $A$, $R$ un sous-$B$-module de $M$ tel que
$N=M/R$ soit un $A$-module plat. Soit $(f_i)_{1\leq i\leq n}$ une suite
d'éléments de l'idéal maximal de $B$ ayant les propriétés suivantes~:
\begin{enumerate}
  \item[(i)] $f_iM\subset R$ pour $1\leq i\leq n$.
  \item[(ii)] Si $g_i$ est l'image de $f_i$ dans $B\otimes_A k$, la suite
  $(g_i)_{1\leq i\leq n}$ est $(M\otimes_A k)$-régulière.
  \item[(iii)] La somme des sous-$(B\otimes_A k)$-modules
  $g_i(M\otimes_A k)$ de $M\otimes_A k$ est l'image canonique de
  $R\otimes_A k$ dans $M\otimes_A k$ (égale au noyau de l'homomorphisme
  canonique $M\otimes_A k\to N\otimes_A k$).
\end{enumerate}
Alors la suite $(f_i)$ est $M$-régulière et l'on a
$R=\sum_{i=1}^n f_iM$.
\end{corollary}

\begin{proof}
\oldpage[0\textsubscript{IV-1}]{19}
Comme $N$ est un $A$-module plat, il résulte de l'exactitude de la suite
$0\to R\to M\to N\to0$ que la suite
\[
  0\longrightarrow R\otimes_A k\longrightarrow M\otimes_A k
  \longrightarrow N\otimes_A k\longrightarrow0
\]
est exacte ($0_{\mathrm I}$, 6.1.2)~; on a donc
$R\otimes_A k=\sum_i g_i(M\otimes_A k)$. Comme $M$ est un $B$-module de type
fini et que $B$ est noethérien, $R$ est aussi un $B$-module de type fini,
d'où $R=\sum_i f_iM$ en vertu du lemme de Nakayama, les $f_i$ appartenant à
l'idéal maximal de $B$.
\end{proof}

\begin{lemma}[15.1.18]
\label{0.15.1.18-fr}
Soient $A$ un anneau, $M$ un $A$-module muni d'une filtration séparée et
exhaustive $(M_k)_{k\geq0}$, et soit $f\in A$ un élément
$\gr_\bullet(M)$-régulier. Alors~:
\begin{enumerate}
  \item[(i)] $f$ est $M$-régulier.
  \item[(ii)] Pour tout $k$, on a $M_k\cap fM=fM_k$, et l'image canonique
  $M'_k$ de $M_k$ dans $M'=M/fM$ est donc canoniquement isomorphe à
  $M_k/fM_k$.
  \item[(iii)] La suite
  \[
    0\longrightarrow M_{k+1}/fM_{k+1}\longrightarrow M_k/fM_k
    \longrightarrow\gr_k(M)/f\cdot\gr_k(M)\longrightarrow0
  \]
  est exacte, et si l'on munit $M'$ de la filtration quotient $(M'_k)$,
  $\gr_\bullet(M')$ est donc canoniquement isomorphe à
  $\gr_\bullet(M)/f\cdot\gr_\bullet(M)$.
\end{enumerate}
\end{lemma}

\begin{proof}
L'assertion (i) résulte de l'hypothèse sur $f$ et de Bourbaki,
\emph{Alg.~comm.}, chap.~III, \S~2, no~8, cor.~1 du th.~1. Comme
$\gr_\bullet(M/M_k)$ est facteur direct de $\gr_\bullet(M)$, $f$ est
$\gr_\bullet(M/M_k)$-régulier pour tout $k$, donc $(M/M_k)$-régulier par
(i), ce qui établit (ii). Pour prouver (iii), il suffit de voir que
l'application $M_{k+1}/fM_{k+1}\to M_k/fM_k$ est injective~: or, si
$x\in M_{k+1}$ est tel que $x\in fM_k$, il résulte de (ii) que
$x\in fM_{k+1}$, ce qui achève de prouver le lemme.
\end{proof}

\begin{proposition}[15.1.19]
\label{0.15.1.19-fr}
Soient $A$ un anneau, $M$ un $A$-module muni d'une filtration
$(M_k)_{k\geq0}$ telle que $M_0=M$, $(f_i)_{1\leq i\leq n}$ une suite
$\gr_\bullet(M)$-régulière d'éléments de $A$. Supposons en outre que pour
$0\leq i\leq n-1$, la filtration quotient de $(M_k)$ sur
\[
  M/\left(\sum_{j\leq i}f_jM\right)
\]
soit séparée. Alors la suite $(f_i)$ est $M$-régulière, et si l'on pose
\[
  M'=M/\left(\sum_{i\leq n}f_iM\right)
\]
et que l'on munit $M'$ de la filtration quotient de $(M_k)$,
$\gr_\bullet(M')$ est canoniquement isomorphe à
\[
  \gr_\bullet(M)/
  \left(\sum_{i\leq n}f_i\gr_\bullet(M)\right).
\]
\end{proposition}

\begin{proof}
Pour $n=1$, la proposition résulte du lemme (15.1.18). Raisonnons par
récurrence sur $n$, en posant
\[
  N=M/\left(\sum_{i\leq n-1}f_iM\right)
\]
et en désignant par $(N_k)$ la filtration quotient de $(M_k)$ sur $N$, qui
est séparée par hypothèse~; en outre $\gr_\bullet(N)$ est isomorphe à
\[
  \gr_\bullet(M)/
  \left(\sum_{i\leq n-1}f_i\gr_\bullet(M)\right).
\]
Par hypothèse $f_n$ est donc $\gr_\bullet(N)$-régulier, et il résulte donc du
lemme (15.1.18) appliqué à $N$ que $f_n$ est $N$-régulier et
$\gr_\bullet(M')$ isomorphe à
$\gr_\bullet(N)/f_n\gr_\bullet(N)$, ce qui démontre la proposition.
\end{proof}

La proposition (15.1.19) s'appliquera en particulier pour les filtrations
vérifiant l'une ou l'autre hypothèse suivante~:
\begin{enumerate}
  \item[$1^\circ$] La filtration $(M_k)$ est finie et séparée (puisque cela
  entraîne $M_k=0$ pour $k$ assez grand)~;
  \item[$2^\circ$] $A$ est un anneau noethérien, $\mathfrak{J}$ un idéal
  contenu dans le radical de $A$, $M$ un $A$-module de type fini et $(M_k)$
  la filtration $\mathfrak{J}$-préadique ($0_{\mathrm I}$, 7.3.5).
\end{enumerate}

\begin{corollary}[15.1.20]
\label{0.15.1.20-fr}
\oldpage[0\textsubscript{IV-1}]{20}
Soient $A$ un anneau, $M$ un $A$-module, $(f_i)_{1\leq i\leq n}$ une suite
$M$-régulière d'éléments de $A$. Alors, quels que soient les entiers
$\nu_i\geq1$, la suite $(f_i^{\nu_i})$ est $M$-régulière et
$M/(\sum_{i\leq n}f_i^{\nu_i}M)$ a une suite de composition dont les
quotients sont isomorphes à $M/(\sum_{i\leq n}f_iM)$.
\end{corollary}

\begin{proof}
Raisonnons par récurrence sur $n$~; pour $n=1$ les assertions sont immédiates
puisque $f^jM/f^{j+1}M$ est isomorphe à $M/fM$, l'homothétie de rapport $f$
dans $M$ étant injective par hypothèse. Posons alors
\[
  N=M/\left(\sum_{i\leq n-1}f_i^{\nu_i}M\right)~;
\]
par l'hypothèse de récurrence, $N$ admet une filtration finie $(N_k)$ pour
laquelle les $\gr_k(N)$ sont isomorphes à
$M/(\sum_{i\leq n-1}f_iM)$~; par hypothèse $f_n$ est donc
$\gr_\bullet(N)$-régulier~; il résulte par suite de (15.1.19) que $f_n$ est
$N$-régulier. Appliquant le cas $n=1$ démontré au début, on en conclut que
$f_n^{\nu_n}$ est $N$-régulier et que $N/f_n^{\nu_n}N$ admet une suite de
composition dont les quotients sont isomorphes à $N/f_nN$~; mais pour la
filtration de $N/f_nN$ quotient de la filtration $(N_k)$,
$\gr_k(N/f_nN)$ est isomorphe à $\gr_k(N)/f_n\gr_k(N)$ par (15.1.19) pour
tout $k$, donc à $M/(\sum_{i\leq n}f_iM)$, ce qui prouve le corollaire.
\end{proof}

\begin{proposition}[15.1.21]
\label{0.15.1.21-fr}
Soient $A$, $B$ deux anneaux locaux noethériens, $\varphi:A\to B$ un
homomorphisme local, $M$ un $B$-module de type fini, $(f_i)_{1\leq i\leq n}$
une suite $A$-régulière formée d'éléments de l'idéal maximal de $A$. Les
propriétés suivantes sont équivalentes~:
\begin{enumerate}
  \item[a)] $M$ est un $A$-module plat.
  \item[b)] Si $\mathfrak{J}$ est l'idéal $\sum_{i\leq n}f_iA$,
  $M/\mathfrak{J}M$ est un $A/\mathfrak{J}$-module plat et la suite $(f_i)$
  est $M$-régulière.
\end{enumerate}
\end{proposition}

\begin{proof}
Le fait que a) entraîne b) résulte de (15.1.13) et de
($0_{\mathrm{III}}$, 10.2.2). Inversement, si la suite $(f_i)$ est
$M$-régulière, elle est $M$-quasi-régulière (15.1.9), donc l'homomorphisme
(15.1.1.1) est bijectif. Mais comme la suite $(f_i)_{1\leq i\leq n}$ est aussi
supposée $A$-régulière, elle est $A$-quasi-régulière (15.1.9), et
l'homomorphisme canonique (15.1.1.1) correspondant
$(A/\mathfrak{J})[T_1,\ldots,T_n]\to\gr_\bullet(A)$ est bijectif. On en
conclut que l'homomorphisme canonique ($0_{\mathrm{III}}$, 10.1.1.2)
\[
  \gr_0(M)\otimes_{A/\mathfrak{J}}\gr_\bullet(A)
  \longrightarrow\gr_\bullet(M)
\]
est bijectif. La conclusion résulte alors de ($0_{\mathrm{III}}$, 10.2.2),
les $f_i$ étant dans l'idéal maximal de $A$ et l'homomorphisme $\varphi$ étant
local.
\end{proof}

\subsection{Suites \texorpdfstring{$\mathscr{F}$}{F}-régulières}
\label{subsection:0.15.2-fr}

\begin{env}[15.2.1]
\label{0.15.2.1-fr}
Soient $X$ un espace annelé, $\mathscr{F}$ un $\mathcal{O}_X$-Module,
$f_i$ ($1\leq i\leq n$) une suite de sections de $\mathcal{O}_X$ au-dessus
de $X$, $\mathscr{J}$ l'Idéal de $\mathcal{O}_X$ engendré par les $f_i$
(image de $\mathcal{O}_X^n$ par l'homomorphisme
$\mathcal{O}_X^n\to\mathcal{O}_X$ défini canoniquement par les $f_i$
($0_{\mathrm I}$, 5.1.1)). Les sous-$\mathcal{O}_X$-Modules
$\mathscr{J}^k\mathscr{F}$ définissent sur $\mathscr{F}$ une filtration, et
on pose encore
$\gr_k(\mathscr{F})=\mathscr{J}^k\mathscr{F}/\mathscr{J}^{k+1}\mathscr{F}$~;
si $S_k$ est l'ensemble des polynômes homogènes de degré total $k$ dans
l'anneau $\mathbb{Z}[T_1,\ldots,T_n]$, $S_k$ peut être considéré comme un
faisceau simple sur $X$, et on définit comme dans (15.1.1) un homomorphisme
canonique
\[
\label{0.15.2.1.1-fr}
  \varphi_k:\gr_0(\mathscr{F})\otimes_{\mathbb Z}S_k
  \longrightarrow\gr_k(\mathscr{F}).
\tag{15.2.1.1}
\]
\oldpage[0\textsubscript{IV-1}]{21}
de la façon suivante~: pour tout ouvert $U$ de $X$ (ou d'une base de la
topologie de $X$), on considère l'homomorphisme (15.1.1.1)
\[
  \varphi_{k,U}:\gr_0(\Gamma(U,\mathscr{F}))\otimes_{\mathbb Z}S_k
  \longrightarrow\gr_k(\Gamma(U,\mathscr{F}))
\]
défini par les $n$ sections $f_i|U\in\Gamma(U,\mathcal{O}_X)$, le
$\Gamma(U,\mathcal{O}_X)$-module $\Gamma(U,\mathscr{F})$ étant filtré par les
sous-modules $(\Gamma(U,\mathscr{J}))^k\Gamma(U,\mathscr{F})$. Le
$\mathcal{O}_X$-Module $\gr_k(\mathscr{F})$ étant associé au préfaisceau
$U\mapsto\gr_k(\Gamma(U,\mathscr{F}))$ pour tout $k$, les $\varphi_{k,U}$
définissent un homomorphisme de $\mathcal{O}_X$-Modules (15.2.1.1). En vertu
des propriétés d'exactitude des limites inductives de modules et de leur
commutation aux produits tensoriels, il résulte de ce qui précède que pour
tout $x\in X$, l'homomorphisme $(\varphi_k)_x$ des fibres des deux membres de
(15.2.1.1) au point $x$ s'identifie à l'homomorphisme de (15.1.1)
\[
\label{0.15.2.1.2-fr}
  \gr_0(\mathscr{F}_x)\otimes_{\mathbb Z}S_k
  \longrightarrow\gr_k(\mathscr{F}_x)
\tag{15.2.1.2}
\]
défini par la suite $((f_i)_x)_{1\leq i\leq n}$ d'éléments de
$\mathcal{O}_x$.
\end{env}

\begin{env}[15.2.2]
\label{0.15.2.2-fr}
Les notations étant celles de (15.2.1), nous dirons que la suite
$(f_i)_{1\leq i\leq n}$ est $\mathscr{F}$-régulière si, en désignant par
$f_j\mathscr{F}$ le sous-$\mathcal{O}_X$-Module de $\mathscr{F}$ image de
l'endomorphisme de $\mathscr{F}$ défini par $f_j$, l'endomorphisme de
$\mathscr{F}/(\sum_{j=1}^{i-1}f_j\mathscr{F})$ défini par $f_i$ est injectif
pour $1\leq i\leq n$~; nous dirons que la suite $(f_i)$ est
$\mathscr{F}$-quasi-régulière si les homomorphismes (15.2.1.1) sont injectifs.
Il résulte de ces définitions et des remarques de (15.2.1) que pour que la
suite $(f_i)$ soit $\mathscr{F}$-régulière (resp.
$\mathscr{F}$-quasi-régulière), il faut et il suffit que, pour tout $x\in X$,
la suite $((f_i)_x)_{1\leq i\leq n}$ soit $\mathscr{F}_x$-régulière (resp.
$\mathscr{F}_x$-quasi-régulière). Toute suite $\mathscr{F}$-régulière est donc
$\mathscr{F}$-quasi-régulière (15.1.9)~; la réciproque est vraie si
$\mathscr{F}$ est un $\mathcal{O}_X$-Module de type fini, si les
$\mathcal{O}_x$ sont des anneaux noethériens pour tout $x\in X$ et si en outre
les $(f_i)_x$ appartiennent au radical de $\mathcal{O}_x$ pour tout $x\in X$
(15.1.11).

Si $X=\operatorname{Spec}(A)$ est un schéma affine, et
$\mathscr{F}=\widetilde{M}$, il
résulte de (I, 1.3) que pour que la suite $(f_i)$ soit
$\mathscr{F}$-régulière (resp. $\mathscr{F}$-quasi-régulière), il faut et il
suffit qu'elle soit $M$-régulière (resp. $M$-quasi-régulière).

Lorsque $\mathscr{F}=\mathcal{O}_X$, on supprimera la mention de
$\mathscr{F}$ dans les définitions précédentes.
\end{env}

\begin{remark}[15.2.3]
\label{0.15.2.3-fr}
Lorsque $X=\operatorname{Spec}(A)$ et
$\mathfrak{J}=\sum_{i=1}^n f_iA$, la suite $(f_j)$ est
quasi-régulière pour tout $A$-module $M$ dans tout ouvert $U$ ne rencontrant
pas $V(\mathfrak{J})$, car en tout point $x\in U$, on a
$\mathfrak{J}_x=\mathcal{O}_x$ et les deux membres de (15.2.1.2) sont nuls. La
notion de suite $\widetilde{M}$-régulière n'a donc d'intérêt qu'au voisinage
des points de $V(\mathfrak{J})$.
\end{remark}

\begin{proposition}[15.2.4]
\label{0.15.2.4-fr}
Soient $X$ un espace annelé, $\mathscr{F}$ un $\mathcal{O}_X$-Module cohérent,
$x$ un point de $X$, $(f_i)_{1\leq i\leq n}$ une suite d'éléments de
$\Gamma(X,\mathcal{O}_X)$. Si la suite $((f_i)_x)_{1\leq i\leq n}$ est
$\mathscr{F}_x$-régulière, il existe un voisinage ouvert $U$ de $x$ tel que la
suite $(f_i|U)_{1\leq i\leq n}$ soit $(\mathscr{F}|U)$-régulière.
\end{proposition}

\begin{proof}
En effet, pour $1\leq i\leq n$, le $\mathcal{O}_X$-Module
$\mathscr{F}/(\sum_{j=1}^{i-1}f_j\mathscr{F})$ est cohérent
($0_{\mathrm I}$, 5.3.4), donc il en est de même du noyau de
l'endomorphisme de ce $\mathcal{O}_X$-Module défini par $f_i$
($0_{\mathrm I}$, 5.3.4)~; comme la fibre au point $x$ de ce noyau est nulle,
il en est de même de sa restriction à un voisinage de $x$
($0_{\mathrm I}$, 5.2.2).
\end{proof}

\begin{proposition}[15.2.5]
\label{0.15.2.5-fr}
\oldpage[0\textsubscript{IV-1}]{22}
Soit $u=(\psi,\theta):X\to Y$ un morphisme plat d'espaces annelés
($0_{\mathrm I}$, 6.7.1)~; soient $\mathscr{F}$ un
$\mathcal{O}_Y$-Module, $\mathscr{G}=u^*(\mathscr{F})$,
$(f_i)_{1\leq i\leq n}$ une suite d'éléments de
$\Gamma(Y,\mathcal{O}_Y)$, $(g_i)_{1\leq i\leq n}$ la suite de leurs images
$g_i=\theta(f_i)\in\Gamma(Y,u_*(\mathcal{O}_X))=\Gamma(X,\mathcal{O}_X)$. Si
la suite $(f_i)$ est $\mathscr{F}$-régulière, la suite $(g_i)$ est
$\mathscr{G}$-régulière. La réciproque est vraie si $u$ est un morphisme
fidèlement plat ($0_{\mathrm I}$, 6.7.8).
\end{proposition}

\begin{proof}
Soient $x$ un point de $X$, $y=u(x)=\psi(x)$~; on a
\[
  \mathscr{G}_x=\psi_x^*(\mathscr{F}_y)
  \otimes_{\psi_x^*(\mathcal{O}_y)}\mathcal{O}_x,
  \qquad\text{et}\qquad
  (g_i)_x=\theta_x^\#(\psi_x^*((f_i)_y))~;
\]
comme par hypothèse
$\theta_x^\#:\psi_x^*(\mathcal{O}_y)\to\mathcal{O}_x$ fait de
$\mathcal{O}_x$ un $\psi_x^*(\mathcal{O}_y)$-module plat, et que le foncteur
$\psi^*$ est exact, la première assertion résulte de l'assertion correspondante
de (15.1.14)~; il en est de même de la seconde, compte tenu de ce que
l'hypothèse que $u$ est fidèlement plat signifie que $\psi$ est surjectif et
que $\theta_x^\#$ fait de $\mathcal{O}_x$ un
$\psi_x^*(\mathcal{O}_y)$-module fidèlement plat.
\end{proof}

% FIN SOURCE ega4/ega0-15-fr.tex

%
% DEBUT SOURCE ega4/ega0-16-fr.tex
\section{Dimension et profondeur dans les anneaux locaux noethériens}
\label{section:0.16-fr}

\subsection{Dimension d'un anneau}
\label{subsection:0.16.1-fr}

\begin{env}[16.1.1]
\label{0.16.1.1-fr}
On appelle dimension (ou dimension de Krull) d'un anneau $A$ et on note
$\dim A$ la dimension (combinatoire) de son spectre
$\operatorname{Spec}(A)$ (14.2.1)~; comme les parties fermées irréductibles de
$\operatorname{Spec}(A)$ sont les
$V(\mathfrak{p})$, où $\mathfrak{p}$ est un idéal premier de $A$ (Bourbaki,
\emph{Alg.~comm.}, chap.~II, \S~4, no~3, cor.~1 de la prop.~14), il revient au
même de dire que $\dim(A)$ est la borne supérieure des longueurs des chaînes
(14.1.1)
\[
  \mathfrak{p}_0\subset\mathfrak{p}_1\subset\cdots\subset\mathfrak{p}_n
\]
d'idéaux premiers de $A$. Si $A\neq0$, on a $\dim A\geq0$. Un anneau artinien
$\neq0$ (en particulier un corps) est de dimension 0 (Bourbaki,
\emph{Alg.~comm.}, chap.~IV, \S~2, no~5, prop.~7)~; un anneau de Dedekind qui
n'est pas un corps (en particulier $\mathbb{Z}$ ou un anneau $k[T]$ de
polynômes à une indéterminée sur un corps $k$) est de dimension 1 (Bourbaki,
\emph{Alg.~comm.}, chap.~VII, \S~2). Un anneau noethérien n'est pas
nécessairement de dimension finie [30].

Si $(\mathfrak{p}_\alpha)$ est la famille des idéaux premiers minimaux de $A$,
on a (14.1.2.1)
\[
\label{0.16.1.1.1-fr}
  \dim(A)=\sup_\alpha(\dim(A/\mathfrak{p}_\alpha)).
\tag{16.1.1.1}
\]
\end{env}

\begin{env}[16.1.2]
\label{0.16.1.2-fr}
Pour tout idéal $\mathfrak{a}$ de $A$, les idéaux premiers de
$A/\mathfrak{a}$ correspondent biunivoquement aux idéaux premiers de $A$
contenant $\mathfrak{a}$, donc
\[
\label{0.16.1.2.1-fr}
  \dim(A/\mathfrak{a})\leq\dim(A).
\tag{16.1.2.1}
\]
Si en outre $\mathfrak{a}$ n'est contenu dans aucun idéal premier minimal
$\mathfrak{p}_\alpha$ de $A$, toute chaîne d'idéaux premiers de $A$ contenant
$\mathfrak{a}$ peut être complétée par un des $\mathfrak{p}_\alpha$, donc on a
\[
\label{0.16.1.2.2-fr}
  1+\dim(A/\mathfrak{a})\leq\dim(A).
\tag{16.1.2.2}
\]
\end{env}

\begin{env}[16.1.3]
\label{0.16.1.3-fr}
\oldpage[0\textsubscript{IV-1}]{23}
Si $S$ est une partie multiplicative de $A$, les idéaux premiers de $S^{-1}A$
correspondent biunivoquement aux idéaux premiers de $A$ ne rencontrant pas
$S$, donc on a
\[
\label{0.16.1.3.1-fr}
  \dim(S^{-1}A)\leq\dim(A).
\tag{16.1.3.1}
\]
Pour tout idéal premier $\mathfrak{p}$ de $A$, on a, par définition (14.2.1)
\[
\label{0.16.1.3.2-fr}
  \dim(A_{\mathfrak{p}})
  =\operatorname{codim}(V(\mathfrak{p}),\operatorname{Spec}(A)).
\tag{16.1.3.2}
\]
Ce nombre est encore appelé hauteur de l'idéal premier $\mathfrak{p}$ et noté
$\operatorname{ht}(\mathfrak{p})$. Plus généralement, pour tout idéal
$\mathfrak{a}$ de $A$, on appelle hauteur de $\mathfrak{a}$ le nombre
\[
\label{0.16.1.3.3-fr}
  \operatorname{ht}(\mathfrak{a})
  =\operatorname{codim}(V(\mathfrak{a}),\operatorname{Spec}(A)).
\tag{16.1.3.3}
\]
Si $(\mathfrak{p}_\lambda)$ est la famille des idéaux premiers minimaux parmi
ceux qui contiennent $\mathfrak{a}$, on a donc (14.2.1)
\[
\label{0.16.1.3.4-fr}
  \operatorname{ht}(\mathfrak{a})
  =\inf_\lambda(\operatorname{ht}(\mathfrak{p}_\lambda)).
\tag{16.1.3.4}
\]
On a évidemment
\[
\label{0.16.1.3.5-fr}
  \dim(A)=\sup_{\mathfrak{m}}(\dim(A_{\mathfrak{m}}))
\tag{16.1.3.5}
\]
où $\mathfrak{m}$ parcourt l'ensemble des idéaux maximaux de $A$.
\end{env}

\begin{env}[16.1.4]
\label{0.16.1.4-fr}
Pour tout idéal premier $\mathfrak{p}$ de $A$, on a (14.2.2.2)
\[
\label{0.16.1.4.1-fr}
  \dim(A_{\mathfrak{p}})+\dim(A/\mathfrak{p})\leq\dim A.
\tag{16.1.4.1}
\]
On dit que $A$ est caténaire (resp. équidimensionnel, équicodimensionnel,
biéquidimensionnel) lorsque $\operatorname{Spec}(A)$ est caténaire (resp.
équidimensionnel,
équicodimensionnel, biéquidimensionnel). Lorsque $A$ est noethérien et
biéquidimensionnel, les deux membres de (16.1.4.1) sont égaux pour tout idéal
premier de $A$ (14.3.5).

Pour tout idéal $\mathfrak{a}$ de $A$, les idéaux premiers de
$A/\mathfrak{a}$ correspondent biunivoquement aux idéaux premiers de $A$
contenant $\mathfrak{a}$, donc, si $A$ est caténaire, il en est de même de
$A/\mathfrak{a}$. De même, pour toute partie multiplicative $S$ de $A$, les
idéaux premiers de $S^{-1}A$ correspondent biunivoquement aux idéaux premiers
de $A$ ne rencontrant pas $S$, donc si $A$ est caténaire, il en est de même de
$S^{-1}A$.

En outre, comme les idéaux premiers de $A$ contenus dans un idéal premier
$\mathfrak{p}$ correspondent biunivoquement aux idéaux premiers de
$A_{\mathfrak{p}}$, pour que $A$ soit caténaire, il faut et il suffit que
$A_{\mathfrak{p}}$ le soit pour tout $\mathfrak{p}$. Dire que $A$ est
caténaire signifie donc ((14.3.2) et (16.1.3.2)) que pour tout couple d'idéaux
premiers $\mathfrak{p}$, $\mathfrak{q}$ de $A$ tels que
$\mathfrak{q}\subset\mathfrak{p}$, on a
\[
\label{0.16.1.4.2-fr}
  \dim(A_{\mathfrak{q}})
  +\dim(A_{\mathfrak{p}}/\mathfrak{q}A_{\mathfrak{p}})
  =\dim(A_{\mathfrak{p}}).
\tag{16.1.4.2}
\]
Si $A$ est biéquidimensionnel, il en est de même de $A_{\mathfrak{p}}$ pour
tout idéal premier $\mathfrak{p}$ de $A$.
\end{env}

\begin{proposition}[16.1.5]
\label{0.16.1.5-fr}
Soit $\rho:A\to B$ un homomorphisme d'anneaux faisant de $B$ une $A$-algèbre
entière. Alors on a $\dim(B)\leq\dim(A)$~; si de plus $\rho$ est injectif, on a
$\dim(B)=\dim(A)$.
\end{proposition}

\begin{proof}
\oldpage[0\textsubscript{IV-1}]{24}
Si $\mathfrak{n}$ est le noyau de $\rho$, $\rho(A)$ est isomorphe à
$A/\mathfrak{n}$, donc $\dim\rho(A)\leq\dim(A)$ par (16.1.2.1), et comme $B$
est entier sur $\rho(A)$, on peut se borner à démontrer la seconde assertion
lorsque $A\subset B$. Si
$\mathfrak{p}'_0\subset\mathfrak{p}'_1\subset\cdots\subset\mathfrak{p}'_n$
est une chaîne d'idéaux premiers de $B$, les $\mathfrak{p}'_i\cap A$ sont
alors des idéaux premiers distincts dans $A$ (Bourbaki, \emph{Alg.~comm.},
chap.~V, \S~2, no~1, cor.~1 de la prop.~1), donc forment une chaîne d'idéaux
premiers. Inversement, pour toute chaîne
$\mathfrak{p}_0\subset\mathfrak{p}_1\subset\cdots\subset\mathfrak{p}_m$
d'idéaux premiers de $A$, on sait qu'il existe pour chaque $i$ un idéal
premier $\mathfrak{p}'_i$ de $B$ tel que
$\mathfrak{p}'_i\cap A=\mathfrak{p}_i$ et que
$\mathfrak{p}'_i\subset\mathfrak{p}'_{i+1}$ pour $0\leq i\leq m-1$
(\emph{loc.~cit.}, cor.~2 du th.~1). D'où la conclusion.
\end{proof}

\begin{proposition}[16.1.6]
\label{0.16.1.6-fr}
Soient $B$ un anneau intègre, $A$ un sous-anneau local de $B$,
$\mathfrak{m}$ son idéal maximal. Supposons que $A$ soit intégralement clos et
que $B$ soit entier sur $A$~; alors, pour tout idéal maximal $\mathfrak{n}$ de
$B$, on a $\dim(B_{\mathfrak{n}})=\dim(A)$.
\end{proposition}

\begin{proof}
La première partie du raisonnement de (16.1.5) montre que
$\dim(B_{\mathfrak{n}})\leq\dim(A)$. Inversement, considérons une chaîne
d'idéaux premiers
$\mathfrak{p}_0\subset\mathfrak{p}_1\subset\cdots\subset
\mathfrak{p}_k=\mathfrak{m}$ de $A$~; on sait que
$\mathfrak{n}\cap A=\mathfrak{m}$ (Bourbaki, \emph{Alg.~comm.}, chap.~V,
\S~2, no~1, prop.~1) et en vertu des hypothèses faites, on peut appliquer le
second théorème de Cohen-Seidenberg (Bourbaki, \emph{Alg.~comm.}, chap.~V,
\S~2, no~4, th.~3), qui prouve l'existence d'une chaîne d'idéaux premiers
$\mathfrak{p}'_0\subset\mathfrak{p}'_1\subset\cdots\subset
\mathfrak{p}'_k=\mathfrak{n}$ de $B$ telle que
$\mathfrak{p}'_i\cap A=\mathfrak{p}_i$ pour tout $i$~; on en déduit aussitôt
$\dim(B_{\mathfrak{n}})\geq\dim(A)$, d'où la proposition.
\end{proof}

\begin{env}[16.1.7]
\label{0.16.1.7-fr}
Soit $M$ un $A$-module~; on appelle dimension de $M$ et on note
$\dim_A(M)$ ou $\dim(M)$ la dimension de l'anneau $A/\operatorname{Ann}(M)$
(isomorphe à l'anneau des homothéties $A_M$ de $M$). On a donc
$\dim(M^k)=\dim(M)$ pour tout $k>0$. Si $M$ est de type fini, les idéaux
premiers de $A$ contenant $\operatorname{Ann}(M)$ sont exactement ceux qui
appartiennent à $\operatorname{Supp}(M)$ ($0_{\mathrm I}$, 1.7.4) et ce
dernier est fermé dans $\operatorname{Spec}(A)$, donc on a alors
\[
\label{0.16.1.7.1-fr}
  \dim(M)=\dim(\operatorname{Supp}(M))\leq\dim(A).
\tag{16.1.7.1}
\]
Si $M$, $N$ sont deux $A$-modules tels que $N\subset M$, on a
\[
\label{0.16.1.7.2-fr}
  \dim(N)\leq\dim(M),\qquad \dim(M/N)\leq\dim(M).
\tag{16.1.7.2}
\]
Si $S$ est une partie multiplicative de $A$ et si $M$ est de type fini, on a
\[
\label{0.16.1.7.3-fr}
  \dim_{S^{-1}A}(S^{-1}M)\leq\dim_A(M).
\tag{16.1.7.3}
\]
En effet, on a
$\operatorname{Ann}(S^{-1}M)=S^{-1}\operatorname{Ann}(M)$ (Bourbaki,
\emph{Alg.~comm.}, chap.~II, \S~2, no~4) et
$S^{-1}A/S^{-1}\operatorname{Ann}(M)=S^{-1}(A/\operatorname{Ann}(M))$.

Par ailleurs, toute chaîne maximale d'idéaux premiers contenant
$\operatorname{Ann}(M)$ se termine par un idéal maximal $\mathfrak{m}$, donc
sa longueur est égale à celle de la chaîne des idéaux premiers correspondants
de $A_{\mathfrak{m}}$ (qui contiennent
$\operatorname{Ann}(M_{\mathfrak{m}})=(\operatorname{Ann}(M))_{\mathfrak{m}}$)~;
de cette remarque et de (16.1.7.3), on tire
\[
\label{0.16.1.7.4-fr}
  \dim_A(M)=\sup_{\mathfrak{m}}
  (\dim_{A_{\mathfrak{m}}}(M_{\mathfrak{m}}))
\tag{16.1.7.4}
\]
où $\mathfrak{m}$ parcourt l'ensemble des idéaux maximaux de $A$.

On dit que $M$ est caténaire (resp. équidimensionnel, équicodimensionnel,
biéquidimensionnel) lorsque $A/\operatorname{Ann}(M)$ l'est.
\end{env}

\begin{proposition}[16.1.8]
\label{0.16.1.8-fr}
\oldpage[0\textsubscript{IV-1}]{25}
Soient $A$ un anneau, $M$ un $A$-module de type fini, $\mathfrak{p}$ un idéal
premier de $A$. On a alors
\[
\label{0.16.1.8.1-fr}
  \dim_{A_{\mathfrak{p}}}(M_{\mathfrak{p}})
  +\dim_A(M/\mathfrak{p}M)\leq\dim_A(M).
\tag{16.1.8.1}
\]
\end{proposition}

\begin{proof}
On a
$\operatorname{Ann}(M_{\mathfrak{p}})=(\operatorname{Ann}(M))_{\mathfrak{p}}$,
donc les idéaux premiers de $A_{\mathfrak{p}}$ contenant
$\operatorname{Ann}(M_{\mathfrak{p}})$ correspondent biunivoquement aux
idéaux premiers de $A$ contenant $\operatorname{Ann}(M)$ et contenus dans
$\mathfrak{p}$. D'autre part, si $\operatorname{Ann}(M)=\mathfrak{n}$,
$V(\operatorname{Ann}(M/\mathfrak{p}M))=V(\mathfrak{p}+\mathfrak{n})$
($0_{\mathrm I}$, 1.7.5), donc les idéaux premiers de $A$ contenant
$\operatorname{Ann}(M/\mathfrak{p}M)$ contiennent
$\mathfrak{p}+\mathfrak{n}$~; d'où la conclusion.

Cette dernière remarque montre en outre que l'on a
\[
\label{0.16.1.8.2-fr}
  \dim_{A/\mathfrak{p}}(M/\mathfrak{p}M)=\dim_A(M/\mathfrak{p}M).
\tag{16.1.8.2}
\]
\end{proof}

\begin{proposition}[16.1.9]
\label{0.16.1.9-fr}
Soient $A$ un anneau noethérien, $\rho:A\to B$ un homomorphisme d'anneaux,
$M$ un $B$-module tel que $M_{[\rho]}$ soit un $A$-module de type fini. Alors
on a
\[
\label{0.16.1.9.1-fr}
  \dim_A(M_{[\rho]})=\dim_B(M).
\tag{16.1.9.1}
\]
\end{proposition}

\begin{proof}
En effet, l'anneau $B_M$ des homothéties du $B$-module $M$ s'identifie à un
sous-anneau de $C=\operatorname{End}_A(M_{[\rho]})$, et $C$ est un $A$-module
de type fini (Bourbaki, \emph{Alg.~comm.}, chap.~III, \S~3, no~1, lemme~2),
donc il en est de même de $B_M$~; d'ailleurs l'anneau $A_M$ des homothéties du
$A$-module $M_{[\rho]}$ est un sous-anneau de $B_M$, image canonique de $A$
dans $C$, donc $B_M$ est fini sur $A_M$~; la conclusion résulte donc de
(16.1.5).
\end{proof}

\begin{env}[16.1.10]
\label{0.16.1.10-fr}
Soit $A$ un anneau noethérien~; si $M$ est un $A$-module de longueur finie, on
sait (Bourbaki, \emph{Alg.~comm.}, chap.~IV, \S~2, no~5, prop.~7 et cor.~1 de
la prop.~7) que $\operatorname{Supp}(M)$ est un espace fini discret, donc
$\dim(M)=0$ si $M\neq0$~; réciproquement, si $M$ est un $A$-module de type
fini tel que $\dim(M)=0$, tout point de $\operatorname{Supp}(M)$ est fermé,
autrement dit est un idéal maximal de $A$, donc (\emph{loc.~cit.}, prop.~7)
$M$ est de longueur finie.
\end{env}

\subsection[Dimension d'un anneau semi-local noethérien]{Dimension d'un
anneau semi-local noethérien\protect\footnote{Ce numéro reproduit l'exposé
donné par Serre dans un cours au Collège de France en 1958.}}
\label{subsection:0.16.2-fr}

\begin{env}[16.2.1]
\label{0.16.2.1-fr}
Soient $A$ un anneau semi-local noethérien, $\mathfrak{r}$ son radical~;
rappelons qu'un idéal de définition $\mathfrak{q}$ de $A$ est un idéal tel que
$\mathfrak{q}\subset\mathfrak{r}$ et que $\mathfrak{q}$ contienne une
puissance de $\mathfrak{r}$~; il revient au même de dire que les seuls idéaux
premiers contenant $\mathfrak{q}$ sont maximaux, ou que $A/\mathfrak{q}$ est
un anneau artinien. Soit $M$ un $A$-module de type fini~; pour tout entier
$n>0$, $M/\mathfrak{q}^nM$ est un $A$-module de longueur finie, égale à
$\sum_{i=0}^{n-1}\operatorname{long}(\gr_i(M))$, où $\gr_\bullet(M)$ est le
$(A/\mathfrak{q})$-module gradué associé à $M$ muni de la filtration
$\mathfrak{q}$-préadique. On sait qu'il existe un polynôme $Q(n)$ à
coefficients rationnels tel que $\operatorname{long}(\gr_n(M))=Q(n)$ pour $n$
assez grand (III, 2.5.4)~; on en déduit aussitôt qu'il existe un polynôme
$P_{\mathfrak{q}}(M,n)$ en $n$, à coefficients rationnels, tel que
\[
  \operatorname{long}(M/\mathfrak{q}^nM)=P_{\mathfrak{q}}(M,n)
\]
pour $n$ assez grand~; il est clair que ce polynôme est unique et que son
coefficient dominant est $>0$ si $M\neq0$~; on dit que c'est le polynôme de
Hilbert-Samuel de $M$ pour la filtration $\mathfrak{q}$-préadique.
\end{env}

\begin{lemma}[16.2.2.1]
\label{0.16.2.2.1-fr}
\oldpage[0\textsubscript{IV-1}]{26}
Soit $(M_n)$ une filtration $\mathfrak{q}$-bonne de $M$ (Bourbaki,
\emph{Alg.~comm.}, chap.~III, \S~3, no~1)~; alors
$\operatorname{long}(M/M_n)$ est égal, pour $n$ grand, à un polynôme en $n$
dont le degré et le coefficient dominant sont les mêmes que ceux de
$P_{\mathfrak{q}}(M,n)$.
\end{lemma}

\begin{proof}
En effet, il existe $n_0$ tel que $M_{n+1}=\mathfrak{q}M_n$ pour $n\geq n_0$,
d'où les inclusions
$\mathfrak{q}^{n+n_0}M\subset M_{n+n_0}=\mathfrak{q}^nM_{n_0}
\subset\mathfrak{q}^nM\subset M_n$ pour $n$ grand, ce qui entraîne
\[
  P_{\mathfrak{q}}(M,n+n_0)\geq
  \operatorname{long}(M/M_{n+n_0})\geq
  P_{\mathfrak{q}}(M,n)\geq\operatorname{long}(M/M_n)
\]
et par suite la conclusion.
\end{proof}

Si $\mathfrak{q}'$ est un second idéal de définition, on a
$\mathfrak{q}'\supset\mathfrak{q}^m$ pour un entier $m$, donc
$P_{\mathfrak{q}'}(M,n)\leq P_{\mathfrak{q}}(M,mn)$~; par suite le degré de
$P_{\mathfrak{q}}(M,n)$ ne dépend pas de l'idéal de définition
$\mathfrak{q}$ considéré~; on le note $d(M)$.

\begin{lemma}[16.2.2.2]
\label{0.16.2.2.2-fr}
Soit $0\to M'\to M\to M''\to0$ une suite exacte de $A$-modules de type
fini. Alors le polynôme
\[
  P_{\mathfrak{q}}(M,n)-P_{\mathfrak{q}}(M',n)
  -P_{\mathfrak{q}}(M'',n)
\]
a un degré $\leq d(M')-1\leq d(M)-1$.
\end{lemma}

\begin{proof}
En effet, la filtration des
$M'_n=M'\cap\mathfrak{q}^nM'$ est $\mathfrak{q}$-bonne (Bourbaki,
\emph{Alg.~comm.}, chap.~III, \S~3, no~1, prop.~1)~; comme
\[
  \operatorname{long}(M/\mathfrak{q}^nM)
  =\operatorname{long}(M''/\mathfrak{q}^nM'')
  +\operatorname{long}(M'/M'_n),
\]
le lemme résulte aussitôt de (16.2.2.1) appliqué à $M'$.
\end{proof}

\begin{theorem}[16.2.3]
\label{0.16.2.3-fr}
\textup{(Krull-Chevalley-Samuel).}
Soient $A$ un anneau semi-local noethérien, $\mathfrak{r}$ son radical,
$M\neq0$ un $A$-module de type fini, $d(M)$ le degré du polynôme de
Hilbert-Samuel de $M$ (pour un idéal de définition quelconque de $A$)~; soit
d'autre part $s(M)$ la borne inférieure des entiers $n$ tels qu'il existe des
éléments $x_1,\ldots,x_n$ de $\mathfrak{r}$ pour lesquels
$M/(x_1M+\cdots+x_nM)$ soit de longueur finie. Alors on a
\[
  d(M)=s(M)=\dim(M).
\]
\end{theorem}

\begin{lemma}[16.2.3.1]
\label{0.16.2.3.1-fr}
Pour tout $x\in\mathfrak{r}$, soit ${}_xM$ le sous-module de $M$ formé des
éléments annulés par $x$. Alors~:
\begin{enumerate}
  \item[\rm{(i)}] On a $s(M)\leq s(M/xM)+1$.
  \item[\rm{(ii)}] Soient $\mathfrak{p}_i$ les idéaux premiers appartenant à
  $\operatorname{Supp}(M)$ et tels que
  \[
    \dim(A/\mathfrak{p}_i)=\dim(M)\quad(1\leq i\leq m).
  \]
\end{enumerate}
Si $x\notin\mathfrak{p}_i$ pour tout $i$, on a
$\dim(M/xM)\leq\dim(M)-1$.
\begin{enumerate}
  \item[\rm{(iii)}] Si $\mathfrak{q}$ est un idéal de définition de $A$, le
  polynôme
  \[
    P_{\mathfrak{q}}({}_xM,n)-P_{\mathfrak{q}}(M/xM,n)
  \]
  est de degré $\leq d(M)-1$.
\end{enumerate}
\end{lemma}

\begin{proof}
L'assertion (i) est triviale, car si $(x_i)_{1\leq i\leq m}$ est tel que pour
le module $N=M/xM$, $N/(x_1N+\cdots+x_mN)$ soit de longueur finie, il suffit
d'observer que ce module est isomorphe à
$M/(xM+x_1M+\cdots+x_mM)$.

Pour prouver (ii), on remarque que, si $\mathfrak{a}$ est l'annulateur de
$M$, les idéaux premiers qui contiennent l'annulateur de $M/xM$ sont ceux qui
contiennent $\mathfrak{a}+Ax$ ($0_{\mathrm I}$, 1.7.5)~; aucun des
$\mathfrak{p}_i$ ne peut donc contenir $\mathfrak{a}+Ax$ et par définition de
$\dim(M)$ (16.1.7) toute chaîne d'idéaux premiers contenant
$\mathfrak{a}+Ax$ a donc une longueur $\leq\dim(M)-1$.

\oldpage[0\textsubscript{IV-1}]{27}

Enfin, pour démontrer (iii), on note que l'on a deux suites exactes
\[
  0\to{}_xM\to M\xrightarrow{x}xM\to0,\qquad
  0\to xM\to M\to M/xM\to0,
\]
et l'assertion découle aussitôt du lemme (16.2.2.2).
\end{proof}

La démonstration de (16.2.3) se fait en trois étapes~:

\begin{env}[16.2.3.2]
\label{0.16.2.3.2-fr}
On a $\dim(M)\leq d(M)$.
\end{env}

Raisonnons par récurrence sur $d(M)$, le cas $d(M)=0$ étant trivial.
Supposons donc $d(M)\geq1$, et soit $\mathfrak{p}_0\in\operatorname{Supp}(M)$
tel que $\dim(A/\mathfrak{p}_0)=\dim(M)$. Cela entraîne que
$\mathfrak{p}_0$ est un élément \emph{minimal} de $\operatorname{Supp}(M)$,
donc il est associé à $M$ (Bourbaki, \emph{Alg.~comm.}, chap.~IV, \S~1,
no~4, th.~2) et $M$ contient par suite un sous-module $N$ isomorphe à
$A/\mathfrak{p}_0$ (\emph{loc.~cit.}, \S~1, no~1)~; comme $d(M)\geq d(N)$
par (16.2.2.2), on est ramené à prouver que $\dim(N)\leq d(N)$. Soit donc
$\mathfrak{p}_0\subset\mathfrak{p}_1\subset\cdots\subset\mathfrak{p}_n$ une
chaîne d'idéaux premiers distincts dans $A$, et montrons que $n\leq d(N)$.
C'est évident si $n=0$. Sinon, il existe
$x\in\mathfrak{p}_1\cap\mathfrak{r}$ tel que $x\notin\mathfrak{p}_0$, car la
relation $\mathfrak{p}_0\supset\mathfrak{p}_1\cap\mathfrak{r}$ entraînerait
$\mathfrak{p}_0\supset\mathfrak{r}$ puisque $\mathfrak{p}_0$ ne contient pas
$\mathfrak{p}_1$, donc $\mathfrak{p}_0$ serait maximal, ce qui est absurde.
On a $\mathfrak{p}_i\in\operatorname{Supp}(N/xN)$, pour $i\geq1$, donc
$n-1\leq\dim(N/xN)$ (16.1.7). Comme ${}_xN=0$ en vertu du choix de $x$, on a
$d(N/xN)\leq d(N)-1$ par (16.2.3.1, (iii)). L'hypothèse de récurrence implique
alors $d(N/xN)=\dim(N/xN)\geq n-1$, donc $n-1\leq d(N)-1$.

\begin{env}[16.2.3.3]
\label{0.16.2.3.3-fr}
On a $d(M)\leq s(M)$.
\end{env}

Soit en effet $(x_i)_{1\leq i\leq n}$ une famille d'éléments de
$\mathfrak{r}$ telle que, si l'on pose
$\mathfrak{a}=x_1A+\cdots+x_nA$, $M/\mathfrak{a}M$ soit de longueur finie.
Les idéaux associés à $M/\mathfrak{a}M$ sont alors maximaux (Bourbaki,
\emph{Alg.~comm.}, chap.~IV, \S~2, no~5, prop.~7) et comme ce sont les seuls
idéaux premiers contenant
$\mathfrak{q}=\mathfrak{a}+(\mathfrak{r}\cap\operatorname{Ann}(M))$,
$A/\mathfrak{q}$ est artinien (\emph{loc.~cit.}, prop.~9) et $\mathfrak{q}$
est un idéal de définition de $A$. Mais on a
$\mathfrak{q}^mM/\mathfrak{q}^{m+1}M=
\mathfrak{a}^mM/\mathfrak{a}^{m+1}M$, et si $M$ est engendré par $r$ éléments
$z_j$, $\mathfrak{a}^mM/\mathfrak{a}^{m+1}M$ est un $(A/\mathfrak{q})$-module
engendré par les images canoniques des
$x_1^{\nu_1}x_2^{\nu_2}\cdots x_n^{\nu_n}z_j$ où
$\sum_i\nu_i=m$, donc par $r\binom{m+n-1}{n-1}$ éléments~; sa longueur est
par suite $\leq ar\binom{m+n-1}{n-1}$, où $a$ est une constante, et on en
déduit aussitôt que le degré de $P_{\mathfrak{q}}(M,n)$ est $\leq n$~; d'où
$d(M)\leq s(M)$.

\begin{env}[16.2.3.4]
\label{0.16.2.3.4-fr}
On a $s(M)\leq\dim(M)$.
\end{env}

Notons d'abord que $n=\dim(M)$ est fini par (16.2.3.2). Raisonnons par
récurrence sur $n$~; la proposition est évidente pour $n=0$, puisque $M$ est
alors de longueur finie (16.1.10). Supposons $n\geq1$, et soient
$\mathfrak{p}_i$ $(1\leq i\leq m)$ les idéaux premiers de
$\operatorname{Supp}(M)$ tels que $\dim(A/\mathfrak{p}_i)=n$~; la définition
des $\mathfrak{p}_i$ montre que ce sont des éléments minimaux de
$\operatorname{Supp}(M)$, donc en nombre fini (Bourbaki, \emph{Alg.~comm.},
chap.~IV, \S~1, no~4, th.~2). Comme $n\geq1$, les $\mathfrak{p}_i$ ne sont pas
maximaux, et il existe donc $x\in\mathfrak{r}$ tel que
$x\notin\mathfrak{p}_i$ pour tout $i$ (Bourbaki, \emph{Alg.~comm.}, chap.~II,
\S~1, no~1, prop.~2). En vertu de (16.2.3.1), on a
$s(M)\leq s(M/xM)+1$ et $\dim(M)\geq\dim(M/xM)+1$~; l'hypothèse de récurrence
entraîne que $s(M/xM)\leq\dim(M/xM)$, d'où la conclusion.

\begin{corollary}[16.2.4]
\label{0.16.2.4-fr}
Sous les hypothèses de (16.2.3), on a
\begin{equation}
\label{0.16.2.4.1-fr}
\tag{16.2.4.1}
  \dim_{\widehat{A}}(\widehat{M})=\dim_A(M).
\end{equation}
\end{corollary}

\oldpage[0\textsubscript{IV-1}]{28}

\begin{proof}
En effet, $M/\mathfrak{r}^nM$ et
$\widehat{M}/\mathfrak{r}^n\widehat{M}$ sont isomorphes, donc
$d(\widehat{M})=d(M)$.
\end{proof}

\begin{corollary}[16.2.5]
\label{0.16.2.5-fr}
La dimension d'un anneau semi-local noethérien $A$ est finie et égale au
nombre minimum d'éléments du radical de $A$ engendrant un idéal de définition.
\end{corollary}

\begin{proof}
C'est l'égalité $s(M)=\dim(M)$ pour $M=A$.
\end{proof}

En particulier~:

\begin{corollary}[16.2.6]
\label{0.16.2.6-fr}
Soient $A$ un anneau local noethérien, $\mathfrak{m}$ son idéal maximal,
$k=A/\mathfrak{m}$ son corps résiduel~; on a
\begin{equation}
\label{0.16.2.6.1-fr}
\tag{16.2.6.1}
  \dim(A)\leq\operatorname{rg}_k(\mathfrak{m}/\mathfrak{m}^2).
\end{equation}
\end{corollary}

\begin{proof}
En effet, on sait que si les $x_i$ $(1\leq i\leq n)$ sont des éléments de
$\mathfrak{m}$ dont les classes mod.~$\mathfrak{m}^2$ forment une base du
$k$-espace vectoriel $\mathfrak{m}/\mathfrak{m}^2$, les $x_i$ engendrent
$\mathfrak{m}$ (Bourbaki, \emph{Alg.~comm.}, chap.~II, \S~3, no~2, cor.~2 de
la prop.~4).
\end{proof}

\begin{proposition}[16.2.7]
\label{0.16.2.7-fr}
Soient $k$ un corps, $B$ une $k$-algèbre graduée de type fini à degrés
positifs engendrée par ses éléments homogènes de degré 1 et telle que $B_0=k$~;
soit $E$ un $B$-module gradué de type fini, de sorte (III, 2.5.4) que, pour $n$
assez grand, $\operatorname{rg}_k(E_n)=r(n)$ est de la forme
$\Phi(n)-\Phi(n-1)$, où $\Phi$ est un polynôme en $n$ de degré $d$. Alors il
existe dans $B$ une famille $(t_i)_{1\leq i\leq d}$ de $d$ éléments homogènes
de degrés $\geq1$ tels que $E/(\sum_i t_iE)$ soit de rang fini sur $k$.
\end{proposition}

\begin{proof}
Soit $\mathfrak{q}$ l'idéal maximal $\bigoplus_{n\geq1}B_n$ de $B$, et, pour
tout $n$, soit $E'_n$ le sous-$B$-module $\bigoplus_{h\geq n+1}E_h$ de $E$.
Tout élément de $B-\mathfrak{q}$ détermine une homothétie injective dans le
$B$-module artinien $E/E'_n$, et cette homothétie est donc bijective
(Bourbaki, \emph{Alg.}, chap.~VIII, \S~2, no~2, lemme~3), ce qui entraîne que
$(E/E'_n)_{\mathfrak{q}}$ s'identifie à $E/E'_n$~; comme le corps résiduel de
$B_{\mathfrak{q}}$ est $k$, on déduit de l'hypothèse que
$\operatorname{long}_{B_{\mathfrak{q}}}((E/E'_n)_{\mathfrak{q}})$ est un
polynôme en $n$ de degré $d$. D'autre part, comme $E$ est un $B$-module de
type fini et que $B$ est engendré par ses éléments homogènes de degré 1, la
filtration $(E'_n)$ sur $E$ est $\mathfrak{q}$-bonne (Bourbaki,
\emph{Alg.~comm.}, chap.~III, \S~1, no~3, lemme~1), donc la filtration
$((E'_n)_{\mathfrak{q}})$ sur $E_{\mathfrak{q}}$ est
$\mathfrak{q}B_{\mathfrak{q}}$-bonne. Puisque $B_{\mathfrak{q}}$ est un anneau
local noethérien, on conclut de (16.2.3) et de (16.2.2.1) que l'on a
$\dim(E_{\mathfrak{q}})=d$. Raisonnons alors par récurrence sur $d$, le lemme
étant évident si $d=0$. Supposons $d\geq1$, et observons que les idéaux premiers
minimaux $\mathfrak{p}_i$ $(1\leq i\leq r)$ dans
$\operatorname{Ass}(E)$ sont \emph{gradués} (Bourbaki, \emph{Alg.~comm.},
chap.~IV, \S~3, no~1, prop.~1), et que $\mathfrak{q}$ est distinct de la
réunion des $\mathfrak{p}_i$ puisque $d\geq1$~; il y a donc un élément
\emph{homogène} $t_1\in\mathfrak{q}$ n'appartenant à aucun des
$\mathfrak{p}_i$ (Bourbaki, \emph{Alg.~comm.}, chap.~III, \S~1, no~4,
prop.~8). Comme les idéaux premiers $(\mathfrak{p}_i)_{\mathfrak{q}}$ sont les
éléments minimaux de $\operatorname{Ass}(E_{\mathfrak{q}})$ (Bourbaki,
\emph{Alg.~comm.}, chap.~IV, \S~1, no~2, prop.~5), on a
$\dim(E_{\mathfrak{q}}/t_1E_{\mathfrak{q}})=d-1$ (16.3.4). Il suffit alors
d'appliquer l'hypothèse de récurrence au $B$-module gradué $E/t_1E$ pour
obtenir la conclusion.
\end{proof}

\subsection{Systèmes de paramètres dans un anneau local noethérien}
\label{subsection:0.16.3-fr}

\begin{proposition}[16.3.1]
\label{0.16.3.1-fr}
Soient $A$ un anneau noethérien, $\mathfrak{p}$ un idéal premier de $A$, $n$ un
entier. Les conditions suivantes sont équivalentes~:
\begin{enumerate}
  \item[$a)$] $\operatorname{ht}(\mathfrak{p})=
  \dim(A_{\mathfrak{p}})\leq n$~;
  \item[$b)$] Il existe $n$ éléments $x_i$ de $A$ $(1\leq i\leq n)$ tels que
  $\mathfrak{p}$ soit minimal dans l'ensemble des idéaux premiers contenant
  l'idéal $\mathfrak{a}$ engendré par les $x_i$ (autrement dit
  $V(\mathfrak{p})$ est une composante irréductible de $V(\mathfrak{a})$).
\end{enumerate}
\end{proposition}

\begin{proof}
La condition $b)$ entraîne que $\mathfrak{a}A_{\mathfrak{p}}$ est un idéal de
définition de $A_{\mathfrak{p}}$, d'où $a)$ en vertu de (16.2.5).
Inversement, si $a)$ est vérifiée, il existe un idéal de définition
$\mathfrak{b}$ de $A_{\mathfrak{p}}$, engendré par $n$ éléments $x_i/s$, avec
$s\in A-\mathfrak{p}$. En vertu de la correspondance biunivoque entre idéaux
premiers de $A_{\mathfrak{p}}$ et idéaux premiers de $A$ contenus dans
$\mathfrak{p}$, l'idéal $\mathfrak{a}$ de $A$ engendré par les $x_i$ vérifie
$b)$.
\end{proof}

\oldpage[0\textsubscript{IV-1}]{29}

Pour $n=1$, on obtient le cas particulier~:

\begin{corollary}[16.3.2]
\label{0.16.3.2-fr}
\textup{(``Hauptidealsatz'').}
Pour qu'un idéal premier dans un anneau noethérien soit de hauteur $\leq1$, il
faut et il suffit qu'il soit minimal dans l'ensemble des idéaux premiers
contenant un idéal principal convenable.
\end{corollary}

\begin{corollary}[16.3.3]
\label{0.16.3.3-fr}
\textup{(Artin-Tate).}
Soit $A$ un anneau intègre noethérien. Les conditions suivantes sont
équivalentes~:
\begin{enumerate}
  \item[$a)$] $A$ est un anneau semi-local de dimension $\leq1$.
  \item[$b)$] Il existe $f\neq0$ dans $A$ tel que $A_f$ soit un corps.
\end{enumerate}
\end{corollary}

\begin{proof}
La condition $a)$ équivaut à dire que $A$ n'a qu'un nombre fini d'idéaux
premiers $\mathfrak{m}_i\neq0$, et que ces idéaux sont tous maximaux ($0$ étant
un idéal premier). Comme le produit des $\mathfrak{m}_i$ n'est pas nul, puisque
$A$ est intègre, un élément $f\neq0$ de ce produit appartient à tous les
$\mathfrak{m}_i$, donc $(0)$ est le seul idéal premier de l'anneau intègre
$A_f$, autrement dit $A_f$ est un corps. (On notera que cette partie de la
démonstration n'utilise pas le fait que $A$ est noethérien.) Prouvons
maintenant que $b)$ entraîne $a)$~; l'hypothèse $b)$ signifie que tout idéal
premier $\mathfrak{p}\neq0$ de $A$ contient $f$. Soient $\mathfrak{p}_i$
$(1\leq i\leq n)$ les idéaux premiers minimaux parmi ceux qui contiennent $f$
(qui sont en nombre fini puisque $A/fA$ est noethérien)~; il suffit de prouver
que les $\mathfrak{p}_i$ sont des idéaux \emph{maximaux}, car alors tout idéal
premier $\neq0$ contient nécessairement un des $\mathfrak{p}_i$, donc lui est
égal. Supposons donc qu'il existe un idéal maximal $\mathfrak{m}$ contenant un
des $\mathfrak{p}_i$ et distinct des $\mathfrak{p}_i$~; $\mathfrak{m}$ n'est
donc pas contenu dans la réunion des $\mathfrak{p}_i$ (Bourbaki,
\emph{Alg.~comm.}, chap.~II, \S~1, no~1, prop.~2), autrement dit il existe
$g\in\mathfrak{m}$ n'appartenant à aucun des $\mathfrak{p}_i$. Soit
$\mathfrak{q}$ un des idéaux premiers minimaux parmi ceux qui contiennent
$g$~; il résulte du ``Hauptidealsatz'' (16.3.2) que $\mathfrak{q}$ est de
hauteur 1, donc $\neq0$~; il contient par suite $f$, donc aussi un des
$\mathfrak{p}_i$, et comme $g\in\mathfrak{q}$ et que $g$ n'appartient à aucun
des $\mathfrak{p}_i$, $\mathfrak{q}$ est distinct des $\mathfrak{p}_i$~; il
serait donc de hauteur $\geq2$, ce qui est contradictoire.
\end{proof}

\begin{proposition}[16.3.4]
\label{0.16.3.4-fr}
Soient $A$ un anneau semi-local noethérien, $\mathfrak{r}$ son radical, $M$ un
$A$-module de type fini, $\mathfrak{p}_i$ les éléments de
$\operatorname{Supp}(M)$ tels que
$\dim(A/\mathfrak{p}_i)=\dim(M)$. Alors les $\mathfrak{p}_i$ sont des éléments
minimaux de $\operatorname{Ass}(M)$, et l'on a
$\dim(M/\mathfrak{p}_iM)=\dim(M)$. Pour tout $x\in\mathfrak{r}$, on a
\begin{equation}
\label{0.16.3.4.1-fr}
\tag{16.3.4.1}
  \dim(M/xM)\geq\dim(M)-1.
\end{equation}
En outre, pour que les deux membres de (16.3.4.1) soient égaux, il faut et il
suffit que $x$ n'appartienne à aucun des $\mathfrak{p}_i$.
\end{proposition}

\begin{proof}
La première assertion est immédiate, car les $\mathfrak{p}_i$ sont par
définition des éléments minimaux de $\operatorname{Supp}(M)$ (16.1.6), et ces
derniers sont aussi éléments minimaux de $\operatorname{Ass}(M)$ (Bourbaki,
\emph{Alg.~comm.}, chap.~IV, \S~1, no~4, th.~2). En outre,
$\operatorname{Supp}(M/\mathfrak{p}_iM)$ est l'ensemble des idéaux premiers
contenant $\mathfrak{p}_i$ ($0_{\mathrm I}$, 1.7.5), donc
$\dim(M/\mathfrak{p}_iM)=\dim(A/\mathfrak{p}_i)=\dim(M)$.

Posons $N=M/xM$. Si $x_1,\ldots,x_n$ sont des éléments de $\mathfrak{r}$ tels
que $N/(x_1N+\cdots+x_nN)$ soit de longueur finie, cela signifie que
$M/(xM+x_1M+\cdots+x_nM)$ est de longueur finie, d'où l'inégalité (16.3.4.1)
en vertu de (16.2.3).

Le fait que les deux membres de (16.3.4.1) soient égaux lorsque $x$
n'appartient à aucun des $\mathfrak{p}_i$ résulte de (16.3.4.1) et de
(16.2.3.1). Inversement, si
\oldpage[0\textsubscript{IV-1}]{30}
$\dim(M/xM)=\dim(M)-1$, aucun des $\mathfrak{p}_i$ ne peut appartenir à
$\operatorname{Supp}(M/xM)$, et les idéaux premiers de ce support sont ceux
qui contiennent $\operatorname{Ann}(M)+Ax$~; comme les $\mathfrak{p}_i$
contiennent $\operatorname{Ann}(M)$, ils ne peuvent contenir $x$.
\end{proof}

\begin{corollary}[16.3.5]
\label{0.16.3.5-fr}
Avec les notations de (16.3.4), pour tout idéal
$\mathfrak{a}\subset\mathfrak{p}_i$, on a
$\dim(M/\mathfrak{a}M)=\dim(M)$ et si l'on pose
$N=M/\mathfrak{a}M$, la relation $\dim(M/xM)=\dim(M)-1$ entraîne
$\dim(N/xN)=\dim(N)-1$.
\end{corollary}

\begin{proof}
En effet, $\operatorname{Supp}(N)$ est l'ensemble des idéaux premiers
contenant $\mathfrak{a}+\operatorname{Ann}(M)$, donc $\mathfrak{p}_i$
appartient à $\operatorname{Supp}(N)$ et comme
$\dim(N)\leq\dim(M)=\dim(A/\mathfrak{p}_i)$, on a
$\dim(A/\mathfrak{p}_i)=\dim(N)$~; en outre les idéaux premiers
$\mathfrak{q}\in\operatorname{Supp}(N)$ tels que
$\dim(A/\mathfrak{q})=\dim(N)$ sont certains des $\mathfrak{p}_j$~; si $x$
n'appartient à aucun des $\mathfrak{p}_j$, on a donc
$\dim(N/xN)=\dim(N)-1$ en vertu de (16.3.4).
\end{proof}

\begin{definition}[16.3.6]
\label{0.16.3.6-fr}
Soient $A$ un anneau semi-local noethérien, $\mathfrak{r}$ son radical, $M$ un
$A$-module de type fini, et posons $n=\dim(M)$. On appelle \emph{système de
paramètres pour $M$} tout système $(x_i)_{1\leq i\leq n}$ de $n$ éléments de
$\mathfrak{r}$ tel que $M/(x_1M+\cdots+x_nM)$ soit de longueur finie.
\end{definition}

On a vu (16.2.3) qu'il existe toujours de tels systèmes.

On a vu au cours de la démonstration de (16.2.3.3) que si
$\mathfrak{J}=\sum_{i=1}^n x_iA$ est tel que $M/\mathfrak{J}M$ soit de
longueur finie, $\mathfrak{q}=\mathfrak{J}+\operatorname{Ann}(M)$ est un
\emph{idéal de définition} de $A$, et réciproquement, s'il en est ainsi,
$M/\mathfrak{J}M$ est un module de type fini sur l'anneau artinien
$A/\mathfrak{q}$, donc est de longueur finie. Il revient donc au même de dire
que $(x_i)_{1\leq i\leq n}$ est un système de paramètres \emph{pour $M$} ou
\emph{pour $A/\operatorname{Ann}(M)$} (ou encore que leurs images dans
$A/\operatorname{Ann}(M)$ forment un système de paramètres pour cet anneau).
Si $(x_i)_{1\leq i\leq n}$ est un système de paramètres pour $M$, il en est de
même de $(x_i^k)$ pour tout entier $k\geq1$, car on peut se borner, en vertu de
ce qui précède, au cas où $M=A$, et si $\mathfrak{J}$ est un idéal de
définition de $A$, l'idéal $\sum_{i=1}^n x_i^kA$ contient
$\mathfrak{J}^{kn}$, donc est aussi un idéal de définition.

\begin{proposition}[16.3.7]
\label{0.16.3.7-fr}
Soient $A$ un anneau semi-local noethérien, $\mathfrak{r}$ son radical,
$x_1,\ldots,x_k$ des éléments de $\mathfrak{r}$, $M$ un $A$-module de type
fini. On a alors
\begin{equation}
\label{0.16.3.7.1-fr}
\tag{16.3.7.1}
  \dim(M/(x_1M+\cdots+x_kM))\geq\dim(M)-k.
\end{equation}
En outre, pour que les deux membres de (16.3.7.1) soient égaux, il faut et il
suffit qu'il existe des éléments $x_i$
$(k+1\leq i\leq n=\dim(M))$ de $\mathfrak{r}$ tels que
$(x_i)_{1\leq i\leq n}$ soit un système de paramètres pour $M$.
\end{proposition}

\begin{proof}
L'inégalité (16.3.7.1) résulte de (16.3.4.1) par récurrence sur $k$. Si les
deux membres de (16.3.7.1) sont égaux, et si $x_{k+1},\ldots,x_n$ est un
système de paramètres pour $M/(x_1M+\cdots+x_kM)$, il est clair que
$M/(x_1M+\cdots+x_kM+x_{k+1}M+\cdots+x_nM)$ est de longueur finie, donc
$(x_i)_{1\leq i\leq n}$ est un système de paramètres pour $M$~;
réciproquement, s'il existe des $x_i$ $(k+1\leq i\leq n)$ ayant cette
propriété et si l'on pose $N=M/(x_1M+\cdots+x_kM)$, il est clair que
$N/(x_{k+1}N+\cdots+x_nN)$ est de longueur finie, donc
$\dim(N)\leq n-k$, et les deux membres sont égaux en vertu de (16.3.7.1).
\end{proof}

\oldpage[0\textsubscript{IV-1}]{31}

\begin{proposition}[16.3.8]
\label{0.16.3.8-fr}
Soient $A$ un anneau semi-local noethérien, $X=\operatorname{Spec}(A)$ son
spectre, $\mathfrak{a}$ un idéal de $A$ distinct de $A$ tel que
$\operatorname{codim}(V(\mathfrak{a}),X)=r>0$. Il existe alors $r$ éléments
$x_1,\ldots,x_r$ de $\mathfrak{a}$, faisant partie d'un système de paramètres
de $A$, tels que
\[
  \operatorname{codim}(V(\mathfrak{a}),
  V(Ax_1+\cdots+Ax_r))=0.
\]
\end{proposition}

\begin{proof}
Soient $\mathfrak{p}_i$ $(1\leq i\leq m)$ les idéaux premiers minimaux de
l'anneau $A$, $\mathfrak{q}_j$ $(1\leq j\leq n)$ les idéaux premiers minimaux
parmi ceux qui contiennent $\mathfrak{a}$~; on a (14.2.1)
\[
  \operatorname{codim}(V(\mathfrak{a}),X)
  =\inf_j(\operatorname{codim}(V(\mathfrak{q}_j),X))
  =\inf_j(\dim(A_{\mathfrak{q}_j}))
\]
(16.1.3.2). L'hypothèse $r>0$ entraîne que $\mathfrak{a}$ n'est contenu dans
aucun des $\mathfrak{p}_i$~; par suite il existe $x\in\mathfrak{a}$ qui
n'appartient à aucun des $\mathfrak{p}_i$ (Bourbaki, \emph{Alg.~comm.},
chap.~II, \S~1, no~1, prop.~2). On a comme ci-dessus
\[
  \operatorname{codim}(V(\mathfrak{a}),V(Ax))
  =\inf_j(\dim((A/Ax)_{\mathfrak{q}_j/Ax}))
  =\inf_j(\dim(A_{\mathfrak{q}_j}/A_{\mathfrak{q}_j}x))~;
\]
mais comme $x/1$ n'est contenu dans aucun des idéaux premiers minimaux de
$A_{\mathfrak{q}_j}$, on a
$\dim(A_{\mathfrak{q}_j}/A_{\mathfrak{q}_j}x)
=\dim(A_{\mathfrak{q}_j})-1$ (16.3.4), d'où
$\operatorname{codim}(V(\mathfrak{a}),V(Ax))=r-1$. Pour la même raison, on a
$\dim(A/Ax)=\dim(A)-1$. Raisonnons maintenant par récurrence sur $r$~; si
$r>1$, il existe $r-1$ éléments $x'_i$ $(2\leq i\leq r)$ de $A'=A/Ax$,
faisant partie d'un système de paramètres de cet anneau, appartenant à
$\mathfrak{a}/Ax$ et tels que
\[
  \operatorname{codim}(V(\mathfrak{a}/Ax),
  V(A'x'_2+\cdots+A'x'_r))=0.
\]
Soit $(x'_i)_{r+1\leq i\leq s}$ un système d'éléments de $A'$ tel que
$(x'_i)_{2\leq i\leq s}$ soit un système de paramètres de $A'$~; pour tout
$i$ tel que $2\leq i\leq s$, soit $x_i$ un élément de la classe $x'_i$ dans
$A$, avec $x_i\in\mathfrak{a}$ pour $2\leq i\leq r$~; il est immédiat que
$A/(Ax+Ax_2+\cdots+Ax_s)$ est de longueur finie, et comme $\dim(A)=s$, on voit
que $x_1=x$ et les $x_i$ d'indices $i\geq2$ répondent aux conditions de
l'énoncé.
\end{proof}

\begin{proposition}[16.3.9]
\label{0.16.3.9-fr}
Soient $A$, $B$ deux anneaux locaux noethériens, $\mathfrak{m}$ l'idéal maximal
de $A$, $k$ son corps résiduel, $\varphi:A\to B$ un homomorphisme local.
\begin{enumerate}
  \item[\rm{(i)}] On a
  \begin{equation}
  \label{0.16.3.9.1-fr}
  \tag{16.3.9.1}
    \dim(B)\leq\dim(A)+\dim(B\otimes_A k).
  \end{equation}
  \item[\rm{(ii)}] Soient $M\neq0$ un $A$-module de type fini, $N\neq0$ un
  $B$-module de type fini. On a
  \begin{equation}
  \label{0.16.3.9.2-fr}
  \tag{16.3.9.2}
    \dim_B(M\otimes_A N)
    \leq\dim_A(M)+\dim_{B\otimes_A k}(N\otimes_A k).
  \end{equation}
\end{enumerate}
\end{proposition}

\begin{proof}
Soient $m$ la dimension de $A$, $(s_i)_{1\leq i\leq m}$ un système de
paramètres de $A$ (16.3.6), de sorte que si
$\mathfrak{a}=\sum_i s_iA$, $A/\mathfrak{a}$ est un $A$-module de longueur
finie. Alors l'anneau $A/\mathfrak{a}$ est artinien, donc l'idéal maximal
$\mathfrak{m}/\mathfrak{a}$ de cet anneau est nilpotent~; l'idéal
$\mathfrak{m}B/\mathfrak{a}B$ de $B/\mathfrak{a}B$, image de
$(\mathfrak{m}/\mathfrak{a})\otimes_A B$, est donc lui aussi nilpotent, et par
suite on a
$\dim(B/\mathfrak{a}B)=\dim(B/\mathfrak{m}B)$ (les spectres de ces deux
anneaux ayant même espace sous-jacent). D'autre part,
$B/\mathfrak{a}B=B/(s_1B+\cdots+s_mB)$, et les images des $s_i$ dans $B$
appartiennent à l'idéal maximal de $B$. Par suite (16.3.7), on a
\[
  \dim(B)\leq m+\dim(B\otimes_A k),
\]
ce qui n'est autre que (16.3.9.1).

Pour prouver (16.3.9.2), notons que si $\mathfrak{r}$ (resp.~$\mathfrak{s}$)
est l'annulateur de $M$ (resp.~$N$)
\oldpage[0\textsubscript{IV-1}]{32}
dans $A$ (resp.~$B$), on a
$\dim_A(M)=\dim(A/\mathfrak{r})$, $\dim_B(N)=\dim(B/\mathfrak{s})$ par
définition (16.1.7)~; d'autre part, $\operatorname{Supp}(M\otimes_A N)$ est un
sous-ensemble fermé de $\operatorname{Spec}(B)$ qui s'identifie (I, 9.1.13) à
\[
  \operatorname{Supp}(M)\times_{\operatorname{Spec}(A)}\operatorname{Supp}(N)
  =\operatorname{Spec}((B/\mathfrak{s})\otimes_A(A/\mathfrak{r}))
  =\operatorname{Spec}(B/(\mathfrak{s}+\mathfrak{r}B)).
\]
Comme $N\otimes_A k$ est un $B$-module de type fini, on a
$\dim_{B\otimes_A k}(N\otimes_A k)=\dim_B(N\otimes_A k)$ (16.1.9), et
$\operatorname{Supp}(N\otimes_A k)$ est égal à
$\operatorname{Spec}(B/(\mathfrak{s}+\mathfrak{m}B))$ par le même raisonnement
que ci-dessus. Appliquant (16.3.9.1) à $A/\mathfrak{r}$ et à
$B/(\mathfrak{s}+\mathfrak{r}B)$ il vient
\[
  \dim(B/(\mathfrak{s}+\mathfrak{r}B))
  \leq\dim(A/\mathfrak{r})+
  \dim(B/(\mathfrak{s}+\mathfrak{m}B)),
\]
ce qui n'est autre que (16.3.9.2).
\end{proof}

\begin{corollary}[16.3.10]
\label{0.16.3.10-fr}
Sous les hypothèses de (16.3.9), si $\mathfrak{m}B$ est un idéal de définition
de $B$, on a $\dim(B)\leq\dim(A)$~; si de plus $A$ est intègre et
$\dim(A)=\dim(B)$, $\varphi$ est injectif.
\end{corollary}

\begin{proof}
La première assertion résulte de (16.3.9) puisque alors
$\dim(B/\mathfrak{m}B)=0$. On peut d'ailleurs remplacer $A$ par
$\varphi(A)$, donc $\dim(B)\leq\dim(\varphi(A))$~; si
$\mathfrak{a}=\operatorname{Ker}(\varphi)\neq0$, et si $A$ est intègre, on a
$\dim(\varphi(A))=\dim(A/\mathfrak{a})<\dim(A)$ (16.1.2.2), donc on ne peut
alors avoir $\dim(A)=\dim(B)$ que si $\mathfrak{a}=0$.
\end{proof}

\subsection[Profondeur et coprofondeur]{Profondeur et
coprofondeur\protect\footnote{Dans la 3e partie du chap.~III, nous développons
la notion de profondeur du point de vue cohomologique et dans un cadre plus
général.}}
\label{subsection:0.16.4-fr}

\begin{proposition}[16.4.1]
\label{0.16.4.1-fr}
Soient $A$ un anneau local noethérien, $\mathfrak{m}$ son idéal maximal, $M$ un
$A$-module de type fini. Alors toute suite $M$-régulière
$(x_i)_{1\leq i\leq r}$ (15.1.7) formée d'éléments de $\mathfrak{m}$, fait
partie d'un système de paramètres pour $M$.
\end{proposition}

\begin{proof}
Raisonnons par récurrence sur $r$~; considérons le $A$-module
\[
  N=M/(x_1M+\cdots+x_{r-1}M)~;
\]
par hypothèse $x_1,\ldots,x_{r-1}$ font partie d'un système de paramètres pour
$M$, donc $\dim(N)=\dim(M)-(r-1)$ (16.3.7). Comme par hypothèse l'homothétie
de rapport $x_r$ dans $N$ est injective, $x_r$ n'appartient à aucun des idéaux
premiers associés à $N$, donc (16.3.4) on a
$\dim(N/x_rN)=\dim(N)-1$, ce qui s'écrit aussi
\[
  \dim(M/(x_1M+\cdots+x_rM))=\dim(M)-r~;
\]
la conclusion résulte donc de (16.3.7).
\end{proof}

Pour un $A$-module $M\neq0$, une suite $M$-régulière d'éléments de
$\mathfrak{m}$ a donc \emph{au plus} $\dim(M)$ éléments.

\begin{lemma}[16.4.2]
\label{0.16.4.2-fr}
Soient $A$ un anneau local noethérien, $\mathfrak{m}$ son idéal maximal,
$k=A/\mathfrak{m}$ son corps résiduel, $M$ un $A$-module. Pour qu'une suite
$M$-régulière $(x_i)_{1\leq i\leq n}$ d'éléments de $\mathfrak{m}$ soit
maximale, il faut et il suffit que l'on ait
\[
  \operatorname{Hom}_A(k,M/(x_1M+\cdots+x_nM))\neq0.
\]
\end{lemma}

\begin{proof}
Dire que $(x_i)$ est maximale signifie que, pour aucun $x\in\mathfrak{m}$,
l'homothétie de rapport $x$ dans $N=M/(x_1M+\cdots+x_nM)$ n'est injective, ou
encore que $\mathfrak{m}$ est contenu dans la réunion des idéaux premiers
associés à $N$ (Bourbaki, \emph{Alg.~comm.}, chap.~IV, \S~1, no~1, cor.~2 de
la prop.~2), donc égal à un de ces derniers, puisque c'est un idéal
\oldpage[0\textsubscript{IV-1}]{33}
maximal (Bourbaki, \emph{Alg.~comm.}, chap.~II, \S~1, no~1, prop.~2)~;
autrement dit, il y a un élément $z\in N$ dont l'annulateur est
$\mathfrak{m}$, et par suite le sous-module $Az$ de $N$ est isomorphe à
$A/\mathfrak{m}=k$~; d'où le lemme.
\end{proof}

\begin{lemma}[16.4.3]
\label{0.16.4.3-fr}
Soient $A$ un anneau local noethérien, $\mathfrak{m}$ son idéal maximal,
$k=A/\mathfrak{m}$ son corps résiduel, $M$ un $A$-module,
$(x_i)_{1\leq i\leq n}$ une suite $M$-régulière d'éléments de
$\mathfrak{m}$. Alors les $A$-modules
\[
  \operatorname{Hom}_A(k,M/(x_1M+\cdots+x_nM))
  \quad\text{et}\quad \operatorname{Ext}^n_A(k,M)
\]
sont isomorphes, et pour $n\geq1$, ils sont aussi isomorphes à
$\operatorname{Ext}^{n-1}_A(k,M/x_1M)$.
\end{lemma}

\begin{proof}
Raisonnons par récurrence sur $n$, la proposition étant évidente pour $n=0$.
Posons $N=M/x_1M$~; la suite $(x_i)_{2\leq i\leq n}$ est $N$-régulière, donc
\[
  \operatorname{Hom}_A(k,M/(x_1M+\cdots+x_nM))
  =\operatorname{Hom}_A(k,N/(x_2N+\cdots+x_nN))
\]
est isomorphe à $\operatorname{Ext}^{n-1}_A(k,N)$ en vertu de l'hypothèse de
récurrence. Considérons alors la suite exacte
$0\to M\xrightarrow{x_1}M\to M/x_1M=N\to0$~; on en déduit la suite exacte des
Ext
\begin{equation}
\label{0.16.4.3.1-fr}
\tag{16.4.3.1}
  \operatorname{Ext}^{n-1}_A(k,M)
  \to\operatorname{Ext}^{n-1}_A(k,N)
  \to\operatorname{Ext}^n_A(k,M)
  \xrightarrow{x_1}\operatorname{Ext}^n_A(k,M).
\end{equation}
En vertu de l'hypothèse de récurrence, $\operatorname{Ext}^{n-1}_A(k,M)$ est
isomorphe à
\[
  \operatorname{Hom}_A(k,M/(x_1M+\cdots+x_{n-1}M)),
\]
qui est nul, puisque $(x_i)_{1\leq i\leq n-1}$ n'est pas une suite
$M$-régulière maximale (16.4.2). D'autre part, comme
$x_1\in\mathfrak{m}$, l'homothétie de rapport $x_1$ dans le $A$-module
$\operatorname{Hom}_A(k,T)$ est nulle pour tout $A$-module $T$, donc il en est
de même de l'homothétie de rapport $x_1$ dans tout $A$-module
$\operatorname{Ext}^n_A(k,T)$~; les assertions du lemme découlent donc
aussitôt de la suite exacte (16.4.3.1).
\end{proof}

\begin{corollary}[16.4.4]
\label{0.16.4.4-fr}
Soient $A$ un anneau local noethérien, $\mathfrak{m}$ son idéal maximal,
$k=A/\mathfrak{m}$ son corps résiduel, $M$ un $A$-module de type fini
$\neq0$. Toutes les suites $M$-régulières maximales d'éléments de
$\mathfrak{m}$ ont le même nombre d'éléments, qui est le plus petit entier $n$
tel que $\operatorname{Ext}^n_A(k,M)\neq0$.
\end{corollary}

\begin{definition}[16.4.5]
\label{0.16.4.5-fr}
Soient $A$ un anneau local noethérien, $\mathfrak{m}$ son idéal maximal, $M$ un
$A$-module de type fini. On appelle \emph{profondeur} de $M$, et l'on note
$\operatorname{prof}_A M$ ou $\operatorname{prof}(M)$ la borne supérieure du
nombre d'éléments d'une suite $M$-régulière formée d'éléments de
$\mathfrak{m}$. On a donc $\operatorname{prof}(M)=+\infty$ si et seulement si
$M=0$, et, par (16.4.1)
\begin{equation}
\label{0.16.4.5.1-fr}
\tag{16.4.5.1}
  \operatorname{prof}(M)\leq\dim(M)\quad\text{si}\quad M\neq0.
\end{equation}
On a évidemment $\operatorname{prof}(M^k)=\operatorname{prof}(M)$ pour tout
$k>0$.
\end{definition}

\begin{proposition}[16.4.6]
\EGAReaderTarget{canon-ref.0.16.4.6-fr}{0.16.4.6-fr}
Soient $A$ un anneau local noethérien, $\mathfrak{m}$ son idéal maximal, $M$ un
$A$-module de type fini.
\begin{enumerate}
  \item[\rm{(i)}] Pour que $\operatorname{prof}(M)=0$, il faut et il suffit que
  $\mathfrak{m}\in\operatorname{Ass}(M)$ (ou encore qu'il existe un sous-module
  de $M$ isomorphe au corps résiduel $k=A/\mathfrak{m}$ de $A$).
  \item[\rm{(ii)}] Pour tout élément $M$-régulier $x$ appartenant à
  $\mathfrak{m}$, on a
  \begin{equation}
  \label{0.16.4.6.1-fr}
  \tag{16.4.6.1}
    \operatorname{prof}(M/xM)=\operatorname{prof}(M)-1.
  \end{equation}
  \item[\rm{(iii)}] Si $M\neq0$, on a
  \begin{equation}
  \label{0.16.4.6.2-fr}
  \tag{16.4.6.2}
    \operatorname{prof}(M)
    \leq\inf_{\mathfrak{p}\in\operatorname{Ass}(M)}\dim(A/\mathfrak{p})
    \leq\dim(M).
  \end{equation}
  \oldpage[0\textsubscript{IV-1}]{34}
  \item[\rm{(iv)}] On a
  $\operatorname{prof}_A(M)=
  \operatorname{prof}_{\widehat{A}}(\widehat{M})$.
\end{enumerate}
\end{proposition}

\begin{proof}
(i) Dire que $\operatorname{prof}(M)=0$ signifie qu'il n'existe aucun
$x\in\mathfrak{m}$ tel que l'homothétie de rapport $x$ dans $M$ soit
injective, autrement dit que tout $x\in\mathfrak{m}$ appartient à un idéal
premier associé à $M$ (Bourbaki, \emph{Alg.~comm.}, chap.~IV, \S~1, no~1,
cor.~2 de la prop.~2)~; mais $\mathfrak{m}$ ne peut être réunion des idéaux
premiers associés à $M$ que s'il est égal à l'un d'eux (Bourbaki,
\emph{Alg.~comm.}, chap.~II, \S~1, no~1, prop.~2).

(ii) Si $(x_i)_{1\leq i\leq m}$ est une suite $(M/xM)$-régulière formée
d'éléments de $\mathfrak{m}$, la suite formée de $x$ et des $x_i$ est
$M$-régulière, et il résulte de la définition qu'elle est maximale lorsque la
suite $(x_i)_{1\leq i\leq m}$ est une suite $(M/xM)$-régulière maximale, d'où
la conclusion.

(iii) Raisonnons par récurrence sur $n=\operatorname{prof}(M)$, l'assertion
étant évidente pour $n=0$. Nous utiliserons le lemme suivant~:
\end{proof}

\begin{lemma}[16.4.6.3]
\label{0.16.4.6.3-fr}
Soient $A$ un anneau noethérien, $M$ un $A$-module de type fini, $t\in A$ un
élément $M$-régulier. Alors, pour tout
$\mathfrak{p}\in\operatorname{Ass}(M)$, tout idéal premier minimal parmi ceux
qui contiennent $\mathfrak{p}+At$ est associé à $M/tM$.
\end{lemma}

\begin{proof}
On sait qu'il existe une suite exacte
\[
  0\to M'\to M\to M''\to0
\]
telle que $\operatorname{Ass}(M')=\{\mathfrak{p}\}$,
$\operatorname{Ass}(M'')=\operatorname{Ass}(M)-\{\mathfrak{p}\}$ (Bourbaki,
\emph{Alg.~comm.}, chap.~IV, \S~1, no~1, prop.~4)~; comme $t$ est
$M$-régulier, il est aussi $M'$-régulier et $M''$-régulier
(\emph{loc.~cit.}, \S~1, no~1, cor.~2 de la prop.~2). On en déduit que la
suite
\[
  0\to M'/tM'\to M/tM\to M''/tM''\to0
\]
est exacte (15.1.18). Par suite, on a
$\operatorname{Ass}(M'/tM')\subset\operatorname{Ass}(M/tM)$. Mais les idéaux
premiers de $A$ contenant $\mathfrak{p}$ sont les points de
$\operatorname{Supp}(M')$ (\emph{loc.~cit.}, \S~1, no~3, prop.~7), donc les
idéaux premiers contenant $\mathfrak{p}+At$ sont les points du support de
$M'/tM'=M'\otimes_A(A/tA)$ ($0_{\mathrm I}$, 1.7.5)~; un élément minimal de
l'ensemble de ces idéaux premiers appartient donc à
$\operatorname{Ass}(M'/tM')$ (Bourbaki, \emph{loc.~cit.}, \S~1, no~3,
cor.~1 de la prop.~7), et \emph{a fortiori} à
$\operatorname{Ass}(M/tM)$.
\end{proof}

\begin{proof}[Suite de la démonstration de (16.4.6)]
Ce lemme étant établi, revenons à la démonstration de (iii). Soit
$x\in\mathfrak{m}$ un élément $M$-régulier, de sorte que
$\operatorname{prof}(M/xM)=n-1$ par (ii)~; pour tout
$\mathfrak{p}\in\operatorname{Ass}(M)$, il résulte du lemme qu'il y a un idéal
$\mathfrak{p}'$ associé à $M/xM$ et contenant $\mathfrak{p}+xA$~; mais comme
$x$ est $M$-régulier, on a $x\notin\mathfrak{p}$, d'où
$\mathfrak{p}'\neq\mathfrak{p}$ et par suite
$\dim(A/\mathfrak{p}')\leq\dim(A/\mathfrak{p})-1$. Or l'hypothèse de
récurrence montre que $n-1\leq\dim(A/\mathfrak{p}')$~; on en conclut que
$n\leq\dim(A/\mathfrak{p})$ pour tout
$\mathfrak{p}\in\operatorname{Ass}(M)$.

(iv) On peut se borner au cas où $M\neq0$. Rappelons que
$\widehat{M}=M\otimes_A\widehat{A}$ à un isomorphisme canonique près
($0_{\mathrm I}$, 7.3.3)~; comme $\widehat{A}$ est un $A$-module plat, on a des
isomorphismes canoniques
\[
  \operatorname{Ext}^i_A(k,M)\otimes_A\widehat{A}
  \xrightarrow{\sim}
  \operatorname{Ext}^i_{\widehat{A}}
  (k,M\otimes_A\widehat{A})
\]
puisque $k=k\otimes_A\widehat{A}$ à un isomorphisme canonique près
($0_{\mathrm{III}}$, 12.3.5). L'assertion (iv) résulte donc de (16.4.4) et du
fait que $\widehat{A}$ est un $A$-module fidèlement plat
($0_{\mathrm I}$, 7.3.5).
\end{proof}

\begin{env}[16.4.6.4]
\label{0.16.4.6.4-fr}
On notera que si $\mathfrak{p}$ est un idéal premier de $A$, on peut avoir
aussi bien
$\operatorname{prof}_{A_{\mathfrak{p}}}(M_{\mathfrak{p}})
<\operatorname{prof}_A(M)$ que
$\operatorname{prof}_{A_{\mathfrak{p}}}(M_{\mathfrak{p}})
>\operatorname{prof}_A(M)$~; pour un exemple du premier cas,
\oldpage[0\textsubscript{IV-1}]{35}
il suffit de prendre pour $A$ un anneau local intègre de dimension $\geq1$~; on a
$\operatorname{prof}(A)\geq1$ mais $\operatorname{prof}(K)=0$ pour le corps
des fractions $K$ de $A$. Pour un exemple du second cas, considérons un
anneau local noethérien intègre $A_0$ de dimension $\geq2$, d'idéal maximal
$\mathfrak{m}_0$ et soit $k=A_0/\mathfrak{m}_0$ son corps résiduel~; soit
$A=A_0\oplus\mathfrak{J}$ l'extension triviale de $A_0$ par le $A_0$-module
$k$ (18.2.3), l'idéal $\mathfrak{J}$ étant $A_0$-isomorphe à $k$ (de sorte
que la multiplication dans $A$ est donnée par
$(x,z)(x',z')=(xx',xz'+x'z)$). Il est clair que $\mathfrak{J}$ est le
nilradical de $A$ et $A/\mathfrak{J}$ est isomorphe à $A_0$, de sorte que
tout idéal premier $\mathfrak{p}$ de $A$ est de la forme
$\mathfrak{p}_0\oplus\mathfrak{J}$, où $\mathfrak{p}_0$ est un idéal premier
de $A_0$~; en particulier $\mathfrak{m}=\mathfrak{m}_0\oplus\mathfrak{J}$
est l'unique idéal maximal de $A$. Comme tout élément de $\mathfrak{m}_0$
annule $\mathfrak{J}$, on a
$\operatorname{prof}(A)=\operatorname{prof}(A_{\mathfrak{m}})=0$~; au
contraire, si $\mathfrak{p}_0\neq\mathfrak{m}_0$, les éléments de
$\mathfrak{J}$ sont annulés par ceux de
$\mathfrak{m}_0-\mathfrak{p}_0$, donc $A_{\mathfrak{p}}$ est canoniquement
isomorphe à $(A_0)_{\mathfrak{p}_0}$, et comme $A_0$ est intègre,
$\operatorname{prof}(A_{\mathfrak{p}})
=\operatorname{prof}((A_0)_{\mathfrak{p}_0})\geq1$ si
$\mathfrak{p}_0\neq0$.
\end{env}

\begin{corollary}[16.4.7]
\label{0.16.4.7-fr}
Pour qu'un anneau local noethérien réduit $A$ soit de profondeur $0$, il faut
et il suffit que $A$ soit un corps.
\end{corollary}

\begin{proof}
En effet, comme l'intersection des idéaux premiers minimaux de $A$ est $0$,
$A$ n'a pas d'idéaux premiers associés immergés~; en vertu de (16.4.6, (i)),
si $\operatorname{prof}(A)=0$, l'idéal maximal $\mathfrak{m}$ de $A$ est
aussi idéal premier minimal, donc c'est le seul idéal premier de $A$, et
comme $A$ est réduit, $\mathfrak{m}=0$.
\end{proof}

\begin{proposition}[16.4.8]
\label{0.16.4.8-fr}
Soient $A$, $B$ deux anneaux locaux noethériens,
$\rho:A\to B$ un homomorphisme local, $M$ un $B$-module tel que
$M_{[\rho]}$ soit un $A$-module de type fini. Alors on a
\begin{equation}
\label{0.16.4.8.1-fr}
\tag{16.4.8.1}
  \operatorname{prof}_A(M_{[\rho]})=\operatorname{prof}_B(M).
\end{equation}
\end{proposition}

\begin{proof}
On peut se borner au cas $M\neq0$~; supposons que
$\operatorname{prof}_A(M_{[\rho]})=n$. Si $(x_i)_{1\leq i\leq n}$ est une
suite $(M_{[\rho]})$-régulière maximale formée d'éléments de l'idéal maximal
de $A$, les $\rho(x_i)$ constituent une suite $M$-régulière formée d'éléments
de l'idéal maximal de $B$, et on a donc, en posant
$N=M/(\rho(x_1)M+\cdots+\rho(x_n)M)$,
$\operatorname{prof}_B(N)=\operatorname{prof}_B(M)-n$ (16.4.6, (ii))~; de
même
$N_{[\rho]}=M_{[\rho]}/(x_1M_{[\rho]}+\cdots+x_nM_{[\rho]})$ et
$\operatorname{prof}_A(N_{[\rho]})=0$~; on est ainsi ramené à démontrer la
proposition lorsque $n=0$. Soit $k$ le corps résiduel de $A$~; comme $k$ est
un quotient de $A$, le $A$-module
$P=\operatorname{Hom}_A(k,M_{[\rho]})$ est un sous-module de
$\operatorname{Hom}_A(A,M_{[\rho]})=M_{[\rho]}$ et l'hypothèse entraîne que
$P\neq0$ (16.4.6, (i))~; en outre, $P$ est un $A$-module de type fini
(puisque $M_{[\rho]}$ est un $A$-module de type fini), donc un $k$-espace
vectoriel de type fini, étant annulé par l'idéal maximal de $A$, et finalement
un $A$-module de longueur finie. D'autre part, $P$ est aussi un sous-module du
$B$-module $M$, et puisqu'il est de longueur finie en tant que $A$-module, il
l'est \emph{a fortiori} en tant que $B$-module. Comme il est $\neq0$, il
contient un sous-$B$-module simple, c'est-à-dire isomorphe au corps résiduel de
$B$~; on conclut alors de (16.4.6, (i)) que l'on a bien
$\operatorname{prof}_B(M)=0$.
\end{proof}

\begin{definition}[16.4.9]
\label{0.16.4.9-fr}
Soient $A$ un anneau local noethérien, $M$ un $A$-module de type fini. Si
$M\neq0$, on appelle \emph{coprofondeur} de $M$ et on note
$\operatorname{coprof}_A(M)$ ou $\operatorname{coprof}(M)$, l'entier fini
$\dim_A(M)-\operatorname{prof}_A(M)\geq0$ (16.4.5.1). Si $M=0$, on pose
$\operatorname{coprof}(M)=0$.

On a $\operatorname{coprof}(M^k)=\operatorname{coprof}(M)$ pour tout $k>0$.
\end{definition}

\begin{proposition}[16.4.10]
\label{0.16.4.10-fr}
Soient $A$ un anneau local noethérien, $M$ un $A$-module de type fini.
\oldpage[0\textsubscript{IV-1}]{36}
\begin{enumerate}
  \item[(i)] Pour tout élément $M$-régulier $x$ appartenant à l'idéal maximal
  de $A$, on a
  $\operatorname{coprof}(M/xM)=\operatorname{coprof}(M)$.
  \item[(ii)] On a
  $\operatorname{coprof}_A(M)=\operatorname{coprof}_{\widehat A}(\widehat M)$.
\end{enumerate}
\end{proposition}

\begin{proof}
En effet, (i) résulte de (16.3.4) et (16.4.6, (ii)), et (ii) de (16.2.4) et
(16.4.6, (iv)).
\end{proof}

Nous montrerons plus tard (IV, 6.11.5) que pour tout idéal premier
$\mathfrak{p}$ de $A$, on a
$\operatorname{coprof}_{A_{\mathfrak{p}}}(M_{\mathfrak{p}})
\leq\operatorname{coprof}_A(M)$.

\begin{proposition}[16.4.11]
\label{0.16.4.11-fr}
Sous les hypothèses de (16.4.8), on a
\begin{equation}
\label{0.16.4.11.1-fr}
\tag{16.4.11.1}
  \operatorname{coprof}_A(M_{[\rho]})=\operatorname{coprof}_B(M).
\end{equation}
\end{proposition}

\begin{proof}
Cela résulte de (16.1.9) et (16.4.8).
\end{proof}

\subsection{Modules de Cohen-Macaulay}
\label{subsection:0.16.5-fr}

\begin{definition}[16.5.1]
\label{0.16.5.1-fr}
Soit $A$ un anneau local noethérien. On dit qu'un $A$-module de type fini $M$
est un \emph{module de Cohen-Macaulay} si
$\operatorname{coprof}(M)=0$, autrement dit si $M=0$ ou si $M\neq0$ et
$\dim(M)=\operatorname{prof}(M)$. On dit que $A$ est un
\emph{anneau de Cohen-Macaulay} si le $A$-module $A$ est un module de
Cohen-Macaulay.
\end{definition}

Un $A$-module de longueur finie est un module de Cohen-Macaulay puisqu'il est
de dimension $0$ (16.1.10). Un anneau local réduit $A$ de dimension $1$ est un
anneau de Cohen-Macaulay, car alors son idéal maximal $\mathfrak{m}$ ne peut
être associé à $A$ (Bourbaki, \emph{Alg.~comm.}, chap.~IV, \S~2, no~5,
prop.~10), donc $\operatorname{prof}(A)=1$ en vertu de (16.4.6, (i)). Si $A$
est un \emph{anneau local intègre et intégralement clos de dimension $\geq2$},
on a $\operatorname{prof}(A)\geq2$~: en effet, si $x\neq0$ est un élément de
$\mathfrak{m}$, on sait (Bourbaki, \emph{Alg.~comm.}, chap.~VII, \S~1, no~4,
prop.~8) que les idéaux premiers associés à $A/xA$ sont de hauteur $1$, donc
distincts de $\mathfrak{m}$, et leur réunion est par suite distincte de
$\mathfrak{m}$ (Bourbaki, \emph{Alg.~comm.}, chap.~II, \S~1, no~1, prop.~2)~;
il existe donc des suites $(x,y)$ qui sont $A$-régulières et formées
d'éléments de $\mathfrak{m}$, d'où notre assertion. On en conclut qu'un anneau
local intégralement clos de dimension $2$ est un anneau de Cohen-Macaulay.

\begin{proposition}[16.5.2]
\label{0.16.5.2-fr}
Soient $A$ un anneau local noethérien, $M$ un $A$-module de type fini. Pour que
$M$ soit un module de Cohen-Macaulay, il faut et il suffit que $\widehat M$
soit un $\widehat A$-module de Cohen-Macaulay.
\end{proposition}

\begin{proof}
Cela résulte aussitôt de (16.4.10, (ii)).
\end{proof}

\begin{proposition}[16.5.3]
\label{0.16.5.3-fr}
Sous les hypothèses de (16.4.8), pour que $M$ soit un $B$-module de
Cohen-Macaulay, il faut et il suffit que $M_{[\rho]}$ soit un $A$-module de
Cohen-Macaulay.
\end{proposition}

\begin{proof}
Cela résulte de (16.4.11).
\end{proof}

\begin{proposition}[16.5.4]
\label{0.16.5.4-fr}
Soient $A$ un anneau local noethérien, $M$ un $A$-module de type fini $\neq0$.
Si $M$ est un $A$-module de Cohen-Macaulay, on a
$\dim(A/\mathfrak{p})=\dim(M)$ pour tout idéal premier
$\mathfrak{p}\in\operatorname{Ass}(M)$~; en particulier aucun idéal premier
associé à $M$ n'est immergé.
\end{proposition}

\begin{proof}
Cela résulte aussitôt des inégalités (16.4.6.2).
\end{proof}

On peut encore dire que $M$ (ou $\operatorname{Supp}(M)$) est
\emph{équidimensionnel} (16.1.7).

\oldpage[0\textsubscript{IV-1}]{37}

\begin{proposition}[16.5.5]
\label{0.16.5.5-fr}
Soient $A$ un anneau local noethérien, $\mathfrak{m}$ son idéal maximal, $M$
un $A$-module de type fini $\neq0$, $x$ un élément de $\mathfrak{m}$.
Supposons que $M$ soit un module de Cohen-Macaulay~; alors les conditions
suivantes sont équivalentes~:
\begin{enumerate}
  \item[a)] $x$ est $M$-régulier~;
  \item[b)] $\dim(M/xM)=\dim(M)-1$~;
  \item[c)] $x$ n'appartient à aucun élément minimal de
  $\operatorname{Supp}(M)$.
\end{enumerate}
En outre, $M/xM$ est alors un module de Cohen-Macaulay.
\end{proposition}

\begin{proof}
On a vu (16.5.4) que tous les éléments $\mathfrak{p}$ de
$\operatorname{Ass}(M)$ sont minimaux dans $\operatorname{Supp}(M)$ et que
$\dim(A/\mathfrak{p})=\dim(M)$~; l'équivalence de a) et c) résulte donc de
Bourbaki, \emph{Alg.~comm.}, chap.~IV, \S~1, no~1, cor.~2 de la prop.~2~;
l'équivalence de b) et c) résulte de (16.3.4). Le fait que $M/xM$ soit un
module de Cohen-Macaulay résulte enfin de (16.4.10, (i)).
\end{proof}

\begin{corollary}[16.5.6]
\label{0.16.5.6-fr}
Sous les hypothèses de (16.5.5), soit $(x_i)_{1\leq i\leq r}$ une suite
d'éléments de $\mathfrak{m}$. Les conditions suivantes sont équivalentes~:
\begin{enumerate}
  \item[a)] La suite $(x_i)$ est $M$-régulière.
  \item[b)] On a
  $\dim(M/(x_1M+\cdots+x_rM))=\dim(M)-r$.
  \item[c)] Les $x_i$ font partie d'un système de paramètres de $M$.
\end{enumerate}
En outre, lorsque ces conditions sont satisfaites,
$N=M/(x_1M+\cdots+x_rM)$ est un module de Cohen-Macaulay, donc, pour tout
idéal $\mathfrak{p}\in\operatorname{Ass}(N)$, on a
$\dim(A/\mathfrak{p})=\dim(M)-r$.
\end{corollary}

\begin{proof}
On sait déjà que b) et c) sont équivalentes (16.3.6) et que a) entraîne c)
(16.4.1) sans hypothèse sur $M$. Pour voir que c) implique a) et prouver la
dernière assertion de l'énoncé, raisonnons par récurrence sur $r$, l'assertion
étant triviale pour $r=0$. Puisque $x_1$ fait partie d'un système de
paramètres, $\dim(M/x_1M)=\dim(M)-1$ (16.3.6), donc $x_1$ est $M$-régulier et
$M/x_1M$ est un module de Cohen-Macaulay par (16.5.5). Comme la suite
$(x_i)_{2\leq i\leq r}$ fait partie d'un système de paramètres pour
$P=M/x_1M$, l'hypothèse de récurrence montre que $(x_i)_{2\leq i\leq r}$ est
une suite $P$-régulière et que
$P/(x_2P+\cdots+x_rP)=N$ est un module de Cohen-Macaulay~; en outre
$(x_i)_{1\leq i\leq r}$ est $M$-régulière.
\end{proof}

\begin{remark}[16.5.7]
\label{0.16.5.7-fr}
Il résulte de (16.5.6) qu'il y a identité entre les \emph{systèmes de
paramètres pour $M$} et les \emph{suites $M$-régulières maximales d'éléments de
$\mathfrak{m}$} lorsque $M$ est un module de Cohen-Macaulay. Inversement, il
résulte de (16.4.1) que si un $A$-module de type fini non nul est tel qu'il y
ait une suite $M$-régulière qui soit un système de paramètres pour $M$, alors
$M$ est un module de Cohen-Macaulay. La dernière propriété énoncée dans
(16.5.6) \emph{caractérise aussi} les modules de Cohen-Macaulay~:
\end{remark}

\begin{proposition}[16.5.8]
\label{0.16.5.8-fr}
Soient $A$ un anneau local noethérien, $M$ un $A$-module de type fini. On
suppose que pour toute suite $(x_i)_{1\leq i\leq r}$ d'éléments de $A$
faisant partie d'un système de paramètres pour $M$, et tout idéal premier
$\mathfrak{p}$ associé à
$N=M/(x_1M+\cdots+x_rM)$, on ait
$\dim(A/\mathfrak{p})=\dim(N)$. Alors $M$ est un module de Cohen-Macaulay.
\end{proposition}

\begin{proof}
Soit $(y_i)_{1\leq i\leq n}$ un système de paramètres pour $M$~; il suffit de
montrer que la suite $(y_i)$ est $M$-régulière. Raisonnons par récurrence sur
$n=\dim(M)$, la proposition étant triviale pour $n=0$. Par hypothèse
(appliquée avec $r=0$) les idéaux premiers $\mathfrak{p}_j$ associés à $M$
sont tous tels que $\dim(A/\mathfrak{p}_j)=\dim(M)$~; comme $y_1$ fait partie
d'un système de paramètres
\oldpage[0\textsubscript{IV-1}]{38}
pour $M$, il résulte de (16.3.7) et (16.3.4) que $y_1$ n'appartient à aucun
des $\mathfrak{p}_j$, donc est $M$-régulier. Si on pose $P=M/y_1M$,
$(y_i)_{2\leq i\leq n}$ est un système de paramètres pour $P$, et il est
immédiat que $P$ vérifie l'hypothèse de l'énoncé~; appliquant l'hypothèse de
récurrence, on voit donc que $(y_i)_{2\leq i\leq n}$ est une suite
$P$-régulière, ce qui démontre la proposition.
\end{proof}

\begin{proposition}[16.5.9]
\label{0.16.5.9-fr}
Soient $A$ un anneau local noethérien, $M$ un $A$-module de type fini~; on
suppose que $M$ soit un module de Cohen-Macaulay. Soit
$\mathfrak{p}\in\operatorname{Supp}(M)$, et posons
$r=\dim(M)-\dim(M/\mathfrak{p}M)$. Alors il existe une suite $M$-régulière
$(x_i)_{1\leq i\leq r}$ formée d'éléments de $\mathfrak{p}$~; pour toute
suite ayant ces propriétés, on a
\[
  \dim(M/(x_1M+\cdots+x_rM))
  =\dim(M/\mathfrak{p}M)=\dim(A/\mathfrak{p}),
\]
et $\mathfrak{p}$ est un élément minimal de
$\operatorname{Ass}(M/(x_1M+\cdots+x_rM))$.
\end{proposition}

\begin{proof}
Prouvons d'abord l'existence des $x_i$~; il n'y a rien à démontrer pour
$r=0$. Pour $r>0$, procédons par récurrence sur $r$. Comme pour tout idéal
premier $\mathfrak{q}_i$ associé à $M$ on a
$\dim(A/\mathfrak{q}_i)=\dim(M)$ (16.5.4), l'hypothèse
$\dim(M/\mathfrak{p}M)<\dim(M)$ entraîne que $\mathfrak{p}$ ne peut être
contenu dans un des $\mathfrak{q}_i$, ces derniers étant les idéaux premiers
minimaux parmi ceux qui contiennent $\operatorname{Ann}(M)$ (16.1.2.2 et
16.1.7)~; on en conclut que $\mathfrak{p}$ n'est pas non plus contenu dans la
réunion des $\mathfrak{q}_i$ (Bourbaki, \emph{Alg.~comm.}, chap.~II, \S~1,
no~1, prop.~2), et par suite il existe un $x_1\in\mathfrak{p}$ qui est
$M$-régulier (Bourbaki, \emph{Alg.~comm.}, chap.~IV, \S~1, no~1, cor.~2 de
la prop.~2). On déduit donc de (16.5.5) que $N=M/x_1M$ est un module de
Cohen-Macaulay, de dimension $\dim(M)-1$~; comme $x_1\in\mathfrak{p}$, on a
d'ailleurs $N/\mathfrak{p}N=M/\mathfrak{p}M$, d'où
$\dim(N)-\dim(N/\mathfrak{p}N)=r-1$~; on peut par suite appliquer à $N$
l'hypothèse de récurrence, et si $(x_i)_{2\leq i\leq r}$ est une suite
$N$-régulière formée d'éléments de $\mathfrak{p}$, il est clair que
$(x_i)_{1\leq i\leq r}$ est une suite $M$-régulière formée d'éléments de
$\mathfrak{p}$.

Si on pose $P=M/(x_1M+\cdots+x_rM)$, on a
$P/\mathfrak{p}P=M/\mathfrak{p}M$ puisque les $x_i$ sont dans
$\mathfrak{p}$, et $\dim(P)=\dim(P/\mathfrak{p}P)$ par (16.5.6)~; mais comme
$P$ est un module de Cohen-Macaulay, on a aussi
$\dim(P)=\dim(A/\mathfrak{p}')$ pour tout
$\mathfrak{p}'\in\operatorname{Ass}(P)$ (16.5.4). Comme
$\mathfrak{p}\in\operatorname{Supp}(P)$ ($0_{\mathrm I}$, 1.7.5),
$\mathfrak{p}$ contient un des idéaux
$\mathfrak{p}'\in\operatorname{Ass}(P)$ (Bourbaki, \emph{Alg.~comm.},
chap.~IV, \S~1, no~4, th.~2), et comme
\[
  \dim(P)=\dim(P/\mathfrak{p}P)
  \leq\dim(A/\mathfrak{p})
  \leq\dim(A/\mathfrak{p}')=\dim(P),
\]
on a nécessairement $\mathfrak{p}=\mathfrak{p}'$ (16.1.2.2), ce qui achève
la démonstration.
\end{proof}

\begin{corollary}[16.5.10]
\label{0.16.5.10-fr}
Sous les mêmes hypothèses que dans (16.5.9)~:
\begin{enumerate}
  \item[(i)] $M_{\mathfrak{p}}$ est un $A_{\mathfrak{p}}$-module de
  Cohen-Macaulay.
  \item[(ii)] On a
  \begin{equation}
  \label{0.16.5.10.1-fr}
  \tag{16.5.10.1}
    \dim(M)=\dim_A(M/\mathfrak{p}M)
    +\dim_{A_{\mathfrak{p}}}(M_{\mathfrak{p}}).
  \end{equation}
\end{enumerate}
\end{corollary}

\begin{proof}
Avec les notations de (16.5.9), soient $x_i'$ les images canoniques des $x_i$
dans $A_{\mathfrak{p}}$ $(1\leq i\leq r)$~; comme les $x_i'$ appartiennent à
l'idéal maximal de $A_{\mathfrak{p}}$ et forment une suite
$M_{\mathfrak{p}}$-régulière par platitude (15.1.14), on a
\[
  \operatorname{prof}_{A_{\mathfrak{p}}}(M_{\mathfrak{p}})
  \geq r=\dim(M)-\dim(M/\mathfrak{p}M)
  \geq\dim_{A_{\mathfrak{p}}}(M_{\mathfrak{p}})
\]
(en vertu de (16.1.8.1)). Mais compte tenu de (16.4.5.1), les trois termes de
cette inégalité sont nécessairement égaux, d'où le corollaire.
\end{proof}

\oldpage[0\textsubscript{IV-1}]{39}

\begin{corollary}[16.5.11]
\label{0.16.5.11-fr}
Soit $A$ un anneau local noethérien~; supposons qu'il existe un $A$-module de
type fini $M$ de support $\operatorname{Spec}(A)$, qui soit un module de
Cohen-Macaulay. Alors, pour tout idéal premier $\mathfrak{p}$ de $A$ on a
\begin{equation}
\label{0.16.5.11.1-fr}
\tag{16.5.11.1}
  \dim(A)=\dim(A/\mathfrak{p})+\dim(A_{\mathfrak{p}}).
\end{equation}
\end{corollary}

\begin{proof}
En effet, $\operatorname{Supp}(M/\mathfrak{p}M)
=\operatorname{Spec}(A/\mathfrak{p})$ ($0_{\mathrm I}$, 1.7.5) et
$\operatorname{Supp}(M_{\mathfrak{p}})
=\operatorname{Spec}(A_{\mathfrak{p}})$~; la relation (16.5.11.1) est donc un
cas particulier de (16.5.10.1).
\end{proof}

\begin{corollary}[16.5.12]
\label{0.16.5.12-fr}
Tout anneau quotient d'un anneau local noethérien $A$ vérifiant les conditions
de (16.5.11) (en particulier tout anneau quotient d'un anneau local de
Cohen-Macaulay) est \emph{caténaire}.
\end{corollary}

\begin{proof}
Il suffit évidemment de prouver que $A$ lui-même est caténaire c'est-à-dire
que, pour deux idéaux premiers $\mathfrak{p}\subset\mathfrak{q}$ de $A$, on a
(16.1.4.2)
\[
  \dim(A_{\mathfrak{q}})
  =\dim(A_{\mathfrak{p}})
  +\dim(A_{\mathfrak{q}}/\mathfrak{p}A_{\mathfrak{q}}).
\]
Or, en vertu de (16.5.10, (i)) $A_{\mathfrak{q}}$ vérifie les mêmes
hypothèses que $A$, et la relation précédente n'est donc autre que
(16.5.11.1) appliquée à $A_{\mathfrak{q}}$ et à l'idéal premier
$\mathfrak{p}A_{\mathfrak{q}}$ de $A_{\mathfrak{q}}$.
\end{proof}

\begin{env}[16.5.13]
\label{0.16.5.13-fr}
Soient $A$ un anneau noethérien, $M$ un $A$-module de type fini. On dit que
$M$ est un \emph{$A$-module de Cohen-Macaulay} si, pour tout idéal premier
$\mathfrak{p}$ de $A$, $M_{\mathfrak{p}}$ est un $A_{\mathfrak{p}}$-module
de Cohen-Macaulay~; en vertu de (16.5.10) cette définition coïncide avec
(16.5.1) lorsque $A$ est \emph{local}. On dit que $A$ est un
\emph{anneau de Cohen-Macaulay} si tous les $A_{\mathfrak{p}}$ le sont. Il est
clair que si $M$ (resp. $A$) est un $A$-module de Cohen-Macaulay (resp. un
anneau de Cohen-Macaulay), $S^{-1}M$ est un $S^{-1}A$-module de Cohen-Macaulay
(resp. $S^{-1}A$ est un anneau de Cohen-Macaulay) pour toute partie
multiplicative $S$ de $A$.
\end{env}

% FIN SOURCE ega4/ega0-16-fr.tex

%
% DEBUT SOURCE ega4/ega0-17-fr.tex
\section{Anneaux réguliers}
\label{section:0.17-fr}

\subsection{Définition des anneaux réguliers}
\label{subsection:0.17.1-fr}

\begin{proposition}[17.1.1]
\label{0.17.1.1-fr}
Soient $A$ un anneau local noethérien de dimension $n$, $\mathfrak{m}$ son
idéal maximal, $k=A/\mathfrak{m}$ son corps résiduel. Les conditions suivantes
sont équivalentes~:
\begin{enumerate}
  \item[\rm a)] L'homomorphisme canonique surjectif de $k$-modules gradués
  (15.1.1.1)
  \[
    \varphi:\mathbf{S}^*_k(\mathfrak{m}/\mathfrak{m}^2)
    \longrightarrow\operatorname{gr}^*_{\mathfrak{m}}(A)
  \]
  (où le second membre est le $k$-module gradué associé à $A$ muni de la
  filtration $\mathfrak{m}$-préadique) est bijectif.
  \item[\rm b)] On a
  $\operatorname{rg}_k(\mathfrak{m}/\mathfrak{m}^2)=n$.
  \item[\rm c)] L'idéal $\mathfrak{m}$ admet un système de générateurs de $n$
  éléments.
  \item[\rm d)] L'idéal $\mathfrak{m}$ admet un système de générateurs qui est
  une suite $A$-régulière.
\end{enumerate}
\end{proposition}


\begin{proof}
En vertu du lemme de Nakayama,
$\operatorname{rg}_k(\mathfrak{m}/\mathfrak{m}^2)$ est le plus petit nombre
d'éléments d'un système de générateurs de $\mathfrak{m}$, donc b) et c) sont
équivalentes. D'autre part, si $(x_i)_{1\leq i\leq r}$ est un système de
générateurs de $\mathfrak{m}$ dont les classes $\overline{x}_i$ mod.
$\mathfrak{m}^2$ forment une base du $k$-espace vectoriel
$\mathfrak{m}/\mathfrak{m}^2$, $\mathbf{S}^*(\mathfrak{m}/\mathfrak{m}^2)$
est isomorphe à l'anneau de
\oldpage[0\textsubscript{IV-1}]{40}
polynômes $k[\mathrm{T}_1,\ldots,\mathrm{T}_r]$~; compte tenu de ce que tout
$A$-module de type fini est séparé pour la filtration $\mathfrak{m}$-préadique
($0_{\mathrm I}$, 7.3.5), il résulte de (15.1.9) que les conditions a) et d)
sont équivalentes~; en outre, comme toute suite $A$-régulière d'éléments de
$\mathfrak{m}$ a au plus $n$ éléments (16.4.1), on voit que d) entraîne c).
Reste à prouver que c) entraîne a).

Posons pour abréger
$\mathrm{S}=\mathbf{S}^*(\mathfrak{m}/\mathfrak{m}^2)
=k[\mathrm{T}_1,\ldots,\mathrm{T}_n]$,
$\mathrm{G}=\operatorname{gr}^*_{\mathfrak{m}}(A)$, et considérons la suite
exacte
$0\longrightarrow\mathfrak{J}\longrightarrow\mathrm{S}
\xrightarrow{\varphi}\mathrm{G}\longrightarrow0$, où le noyau
$\mathfrak{J}$ de $\varphi$ est donc un idéal gradué de $\mathrm{S}$. Pour
tout entier $s>0$, on a donc
\begin{equation}
\label{0.17.1.1.1-fr}
\tag{17.1.1.1}
  \binom{s+n-1}{n-1}
  =\operatorname{long}(\mathrm{S}_s)
  =\operatorname{long}(\mathrm{G}_s)
  +\operatorname{long}(\mathfrak{J}_s).
\end{equation}

Supposons $\mathfrak{J}\neq0$, de sorte qu'il existe un élément homogène
$u\in\mathfrak{J}$ de degré $h>0$~; $\mathfrak{J}_s$ contient
$u\mathrm{S}_{s-h}$ pour $s\geq h$, et comme $u\neq0$ et que $\mathrm{S}$ est
intègre, $u\mathrm{S}_{s-h}$ est isomorphe à $\mathrm{S}_{s-h}$ comme
$k$-espace vectoriel~; on aurait donc
$\operatorname{long}(\mathfrak{J}_s)\geq
\operatorname{long}(\mathrm{S}_{s-h})=\binom{s-h+n-1}{n-1}$, d'où pour
$s\geq h$,
\begin{equation}
\label{0.17.1.1.2-fr}
\tag{17.1.1.2}
  \operatorname{long}(\mathrm{G}_s)
  \leq\binom{s+n-1}{n-1}-\binom{s-h+n-1}{n-1}.
\end{equation}
Or, le second membre de (17.1.1.2) est un polynôme en $s$ de degré
$\leq n-2$~; mais on a
$\mathrm{G}_s=\mathfrak{m}^s/\mathfrak{m}^{s+1}$, et pour $s$ assez grand,
$\operatorname{long}(\mathrm{G}_s)$ est un polynôme en $s$ de degré exactement
égal à $n-1$ dont le coefficient dominant est $>0$ (16.2.1), ce qui contredit
l'inégalité (17.1.1.2). C.Q.F.D.
\end{proof}

\begin{definition}[17.1.2]
\label{0.17.1.2-fr}
On dit qu'un anneau local noethérien qui vérifie les conditions équivalentes de
(17.1.1) est \emph{régulier}.
\end{definition}

\begin{corollary}[17.1.3]
\label{0.17.1.3-fr}
Un anneau local régulier est intègre, intégralement clos, et un anneau de
Cohen-Macaulay.
\end{corollary}

\begin{proof}
La dernière assertion résulte de (17.1.1, d))~; d'autre part, si $A$ est
régulier, $\operatorname{gr}^*_{\mathfrak{m}}(A)$ est intègre et complètement
intégralement clos, étant isomorphe à $k[\mathrm{T}_1,\ldots,\mathrm{T}_n]$~;
on en conclut que $A$ possède aussi ces deux propriétés (Bourbaki,
\emph{Alg.~comm.}, chap.~III, \S~2, no~3, cor. de la prop.~1 et chap.~V,
\S~1, no~5, prop.~15).
\end{proof}

\begin{examples}[17.1.4]
\label{0.17.1.4-fr}
\begin{enumerate}
  \item[\rm (i)] Un anneau local régulier de \emph{dimension} $0$, étant
  intègre par (17.1.3), est nécessairement un \emph{corps}, et réciproquement.
  \item[\rm (ii)] Pour qu'un anneau local noethérien de \emph{dimension} $1$
  soit régulier, il faut et il suffit qu'il soit un \emph{anneau de valuation
  discrète}~: en effet, dire qu'un anneau local noethérien de dimension $1$ est
  régulier signifie que son idéal maximal est \emph{principal} et la conclusion
  résulte de Bourbaki, \emph{Alg.~comm.}, chap.~VI, \S~3, no~6, prop.~9.
  \item[\rm (iii)] Soit $k$ un corps~; l'anneau de séries formelles
  $A=k[[\mathrm{T}_1,\ldots,\mathrm{T}_n]]$ est un \emph{anneau régulier de
  dimension $n$}~; en effet, il est clair que les $\mathrm{T}_i$
  $(1\leq i\leq n)$ engendrent l'idéal maximal $\mathfrak{m}$ de $A$ et
  forment une suite $A$-régulière, car
  $A/(\mathrm{T}_1A+\cdots+\mathrm{T}_iA)$ est isomorphe à
  $k[[\mathrm{T}_{i+1},\ldots,\mathrm{T}_n]]$.
\end{enumerate}
\end{examples}

\begin{proposition}[17.1.5]
\label{0.17.1.5-fr}
Pour qu'un anneau local noethérien $A$ soit régulier, il faut et il suffit que
son complété $\widehat{A}$ le soit.
\end{proposition}

\begin{proof}
\oldpage[0\textsubscript{IV-1}]{41}
En effet, l'idéal maximal de $\widehat{A}$ est $\mathfrak{m}\widehat{A}$ et on
sait que $\mathfrak{m}^h/\mathfrak{m}^{h+1}$ et
$(\mathfrak{m}\widehat{A})^h/(\mathfrak{m}\widehat{A})^{h+1}$ sont des
$k$-espaces vectoriels isomorphes~; la condition a) de (17.1.1) est donc la
même pour $A$ et $\widehat{A}$.
\end{proof}

\begin{definition}[17.1.6]
\label{0.17.1.6-fr}
Étant donné un anneau local noethérien $A$, d'idéal maximal $\mathfrak{m}$, on
appelle \emph{système régulier de paramètres de $A$} un système de paramètres
pour $A$ (16.3.5) qui engendre $\mathfrak{m}$.
\end{definition}

L'existence d'un tel système équivaut au fait que $A$ est régulier (17.1.1)~;
si $(x_i)_{1\leq i\leq n}$ est un système régulier de paramètres, c'est un
système minimal de générateurs de $\mathfrak{m}$, dont les classes mod.
$\mathfrak{m}^2$ forment une base de $\mathfrak{m}/\mathfrak{m}^2$ sur $k$~;
en vertu de (17.1.1, a)) et de (15.1.9), $(x_i)$ est une suite $A$-régulière
maximale (d'où la terminologie).

On prendra garde toutefois que, dans un anneau régulier, un système de
paramètres pour $A$ qui est aussi une suite $A$-régulière \emph{n'est pas
nécessairement} un système régulier de paramètres, comme le montre l'exemple
des puissances $k$-ièmes d'un système régulier de paramètres (15.1.20).

\begin{proposition}[17.1.7]
\label{0.17.1.7-fr}
Soient $A$ un anneau local noethérien, $\mathfrak{m}$ son idéal maximal,
$k=A/\mathfrak{m}$ son corps résiduel, $(x_i)_{1\leq i\leq r}$ une suite
d'éléments de $\mathfrak{m}$, $\mathfrak{J}$ l'idéal
$x_1A+\cdots+x_rA$. Les conditions suivantes sont équivalentes~:
\begin{enumerate}
  \item[\rm a)] $A$ est régulier et les $x_i$ font partie d'un système
  régulier de paramètres de $A$.
  \item[\rm a')] $A$ est régulier et les classes $\overline{x}_i$ des $x_i$
  mod. $\mathfrak{m}^2$ sont linéairement indépendantes sur $k$.
  \item[\rm b)] Les $x_i$ font partie d'un système de paramètres pour $A$ et
  $A/\mathfrak{J}$ est un anneau régulier.
\end{enumerate}
En outre, lorsque ces conditions sont vérifiées, $\mathfrak{J}$ est un idéal
premier.
\end{proposition}

\begin{proof}
La dernière assertion résulte trivialement du fait que $A/\mathfrak{J}$,
étant régulier, est intègre (17.1.3). Posons $n=\dim(A)$. L'équivalence de a)
et a') est immédiate~; en effet, les classes mod. $\mathfrak{m}^2$ des éléments
d'un système régulier de paramètres forment une base de
$\mathfrak{m}/\mathfrak{m}^2$ par rapport à $k$, et réciproquement, si les
$\overline{x}_i$ $(1\leq i\leq r)$ sont linéairement indépendants sur $k$, on
peut trouver $n-r$ éléments $x_i\in\mathfrak{m}$ $(r+1\leq i\leq n)$ tels que
les $\overline{x}_i$ $(1\leq i\leq n)$ forment une base de
$\mathfrak{m}/\mathfrak{m}^2$, donc $(x_i)_{1\leq i\leq n}$ est un système
régulier de paramètres. Pour voir que a) entraîne b), remarquons que puisque
les $x_i$ font partie d'un système de paramètres de $A$, on a
$\dim(A/\mathfrak{J})=n-r$ (16.3.6)~; soit
$\mathfrak{n}=\mathfrak{m}/\mathfrak{J}$ l'idéal maximal de
$A/\mathfrak{J}$. On a une suite exacte
\begin{equation}
\label{0.17.1.7.1-fr}
\tag{17.1.7.1}
  0\longrightarrow(\mathfrak{m}^2+\mathfrak{J})/\mathfrak{m}^2
  \longrightarrow\mathfrak{m}/\mathfrak{m}^2
  \longrightarrow\mathfrak{n}/\mathfrak{n}^2\longrightarrow0
\end{equation}
puisque $\mathfrak{n}^2=(\mathfrak{m}^2+\mathfrak{J})/\mathfrak{J}$~;
l'hypothèse a) entraîne que
$\operatorname{rg}_k((\mathfrak{m}^2+\mathfrak{J})/\mathfrak{m}^2)=r$, donc
$\operatorname{rg}_k(\mathfrak{n}/\mathfrak{n}^2)=n-r
=\dim(A/\mathfrak{J})$, ce qui prouve b) en vertu de (17.1.1, b)).

Inversement, pour prouver que b) entraîne a), notons que le fait que les $x_i$
font partie d'un système de paramètres entraîne que
$\dim(A/\mathfrak{J})=n-r$ (16.3.6)~; l'hypothèse que $A/\mathfrak{J}$ est
régulier implique d'autre part
$\operatorname{rg}_k(\mathfrak{n}/\mathfrak{n}^2)=n-r$ (17.1.1)~; par ailleurs
on a évidemment
$\operatorname{rg}_k((\mathfrak{m}^2+\mathfrak{J})/\mathfrak{m}^2)\leq r$,
donc on tire de (17.1.7.1) que
$\operatorname{rg}_k(\mathfrak{m}/\mathfrak{m}^2)\leq r+(n-r)=n$~; donc $A$
est régulier en vertu de (16.2.6) et (17.1.1). En outre la relation
$\operatorname{rg}_k(\mathfrak{m}/\mathfrak{m}^2)=n$ entraîne alors
$\operatorname{rg}_k((\mathfrak{m}^2+\mathfrak{J})/\mathfrak{m}^2)=r$, et par
suite les $\overline{x}_i$ sont linéairement indépendants sur $k$. C.Q.F.D.
\end{proof}

\oldpage[0\textsubscript{IV-1}]{42}

\begin{corollary}[17.1.8]
\label{0.17.1.8-fr}
Soient $A$ un anneau local noethérien, $\mathfrak{m}$ son idéal maximal, $t$ un
élément de $\mathfrak{m}$. Les conditions suivantes sont équivalentes~:
\begin{enumerate}
  \item[\rm a)] $A/tA$ est régulier et $t$ est non diviseur de zéro dans $A$.
  \item[\rm b)] $A$ est régulier et $t\notin\mathfrak{m}^2$.
\end{enumerate}
\end{corollary}

\begin{proof}
C'est le cas particulier $r=1$ de (17.1.7) (compte tenu de ce qu'un élément de
$\mathfrak{m}$ non diviseur de zéro fait partie d'un système de paramètres
(16.4.1)).
\end{proof}

\begin{corollary}[17.1.9]
\label{0.17.1.9-fr}
Soient $A$ un anneau local régulier, $\mathfrak{m}$ son idéal maximal,
$\mathfrak{J}$ un idéal de $A$ contenu dans $\mathfrak{m}$. Les conditions
suivantes sont équivalentes~:
\begin{enumerate}
  \item[\rm a)] L'anneau $A/\mathfrak{J}$ est régulier.
  \item[\rm b)] $\mathfrak{J}$ est engendré par une suite
  $(x_i)_{1\leq i\leq r}$ faisant partie d'un système régulier de paramètres
  de $A$.
\end{enumerate}
\end{corollary}

\begin{proof}
On a déjà vu (17.1.7) que b) entraîne a). Inversement, supposons que
$A/\mathfrak{J}$ soit régulier (ce qui entraîne que $\mathfrak{J}$ est premier)
et soient $n=\dim(A)$, $n-r=\dim(A/\mathfrak{J})$~; avec les notations de
(17.1.7), la suite exacte (17.1.7.1) donne
$\operatorname{rg}_k((\mathfrak{m}^2+\mathfrak{J})/\mathfrak{m}^2)=r$, donc il
existe une suite $(x_i)_{1\leq i\leq r}$ d'éléments de $\mathfrak{J}$ faisant
partie d'un système régulier de paramètres de $A$. Posons
$\mathfrak{J}'=x_1A+\cdots+x_rA$~; il résulte de (17.1.7) que $\mathfrak{J}'$
est un idéal premier de $A$ et que $A/\mathfrak{J}'$ est régulier et de
dimension $n-r$~; mais comme
$A/\mathfrak{J}=(A/\mathfrak{J}')/(\mathfrak{J}/\mathfrak{J}')$ et que
$\mathfrak{J}/\mathfrak{J}'$ est premier dans l'anneau intègre
$A/\mathfrak{J}'$, les dimensions de $A/\mathfrak{J}$ et de
$A/\mathfrak{J}'$ ne peuvent être égales que si
$\mathfrak{J}=\mathfrak{J}'$ (16.1.2.2).
\end{proof}

\subsection{Rappels sur la dimension projective et la dimension injective des
modules}
\label{subsection:0.17.2-fr}

\begin{env}[17.2.1]
\label{0.17.2.1-fr}
Soient $A$ un anneau, $M$ un $A$-module. Rappelons (M, VI, 2) qu'on appelle
\emph{dimension projective} (resp. \emph{dimension injective}) de $M$ et que
l'on note $\operatorname{dim.proj}(M)$ ou
$\operatorname{dim.proj}_A(M)$ (resp. $\operatorname{dim.inj}(M)$ ou
$\operatorname{dim.inj}_A(M)$) le plus petit $n$ (entier ou égal à $+\infty$)
tel qu'il existe une résolution projective gauche (resp. une résolution
injective droite) de $M$ qui soit \emph{de longueur $n$}. Il revient au même
(\emph{loc. cit.}) de dire que $n$ est le plus petit nombre tel que pour tout
$A$-module $N$, on ait $\operatorname{Ext}^i_A(M,N)=0$ (resp.
$\operatorname{Ext}^i_A(N,M)=0$) pour tout $i>n$, ou seulement pour $i=n+1$.
Pour tout \emph{facteur direct} $M'$ de $M$, on a donc
$\operatorname{dim.proj}(M')\leq\operatorname{dim.proj}(M)$ puisque
$\operatorname{Ext}^i_A(M',N)$ est facteur direct de
$\operatorname{Ext}^i_A(M,N)$~; de même
$\operatorname{dim.inj}(M')\leq\operatorname{dim.inj}(M)$. Une condition
équivalente à $\operatorname{dim.proj}(M)=n$ (resp.
$\operatorname{dim.inj}(M)=n$) est que $n$ est le plus petit nombre tel que,
pour toute suite exacte
\[
  0\longrightarrow R\longrightarrow P_{n-1}\longrightarrow\cdots
  \longrightarrow P_0\longrightarrow M\longrightarrow0
\]
(resp.
$0\longrightarrow M\longrightarrow Q_0\longrightarrow\cdots
\longrightarrow Q_{n-1}\longrightarrow C\longrightarrow0$), où tous les
$P_i$ pour $0\leq i<n$ sont projectifs (resp. tous les $Q_i$ pour $0\leq i<n$
sont injectifs), alors $R$ est projectif (resp. $C$ est injectif).
\end{env}

\begin{remarks}[17.2.2]
\label{0.17.2.2-fr}
\begin{enumerate}
  \item[\rm (i)] Dire que $M$ est un $A$-module \emph{projectif} (resp.
  \emph{injectif}) équivaut donc à dire que
  $\operatorname{dim.proj}(M)=0$ (resp. $\operatorname{dim.inj}(M)=0$).
  \item[\rm (ii)] Soit $T$ un \emph{foncteur covariant} additif de la catégorie
  des $A$-modules dans une catégorie abélienne. Il résulte aussitôt des
  définitions de (17.2.1) et de celle des
  \oldpage[0\textsubscript{IV-1}]{43}
  foncteurs dérivés que si $\operatorname{dim.proj}(M)\leq n$ (resp.
  $\operatorname{dim.inj}(M)\leq n$), on a $L^iT(M)=0$ (resp.
  $R^iT(M)=0$) \emph{pour tout $i>n$}.
  \item[\rm (iii)] Si on suppose $A$ \emph{noethérien} et $M$ \emph{de type
  fini}, la dernière interprétation de la dimension projective donnée dans
  (17.2.1) montre que si $\operatorname{dim.proj}(M)=n$, $M$ admet une
  résolution gauche de longueur $n$ formée de modules projectifs \emph{de type
  fini}. Si de plus $A$ est un anneau \emph{local} noethérien, $M$ admet une
  résolution de longueur $n$ par des modules \emph{libres de type fini}, un
  $A$-module projectif de type fini étant alors libre (Bourbaki,
  \emph{Alg.~comm.}, chap.~II, \S~3, no~2, cor.~2 de la prop.~5).
\end{enumerate}
\end{remarks}

\begin{lemma}[17.2.3]
\label{0.17.2.3-fr}
Soient $A$ un anneau, $M$ un $A$-module. Pour que
$\operatorname{dim.inj}(M)\leq n$, il faut et il suffit que pour tout
$A$-module monogène $N$, on ait $\operatorname{Ext}^{n+1}_A(N,M)=0$.
\end{lemma}

\begin{proof}
Avec les notations de (17.2.1), il s'agit de prouver que $C$ est injectif~;
pour tout $A$-module $N$, $\operatorname{Ext}^{n+1}_A(N,M)$ est isomorphe à
$\operatorname{Ext}^1_A(N,C)$ (M, V, 7), donc on a
$\operatorname{Ext}^1_A(N,C)=0$ pour tout $A$-module $N$ de la forme
$A/\mathfrak{J}$, où $\mathfrak{J}$ est un idéal quelconque de $A$. La suite
exacte des Ext montre alors que l'homomorphisme canonique
$\operatorname{Hom}(A,C)\longrightarrow\operatorname{Hom}(\mathfrak{J},C)$ est
surjectif pour tout idéal $\mathfrak{J}$ de $A$, ce qui implique que $C$ est un
$A$-module injectif (M, I, 3.2).
\end{proof}

\begin{lemma}[17.2.4]
\label{0.17.2.4-fr}
Soient $A$ un anneau noethérien, $M$ un $A$-module de type fini. Pour que
$\operatorname{dim.proj}(M)\leq n$, il faut et il suffit que pour tout
$A$-module monogène $N$, on ait $\operatorname{Ext}^{n+1}_A(M,N)=0$.
\end{lemma}

\begin{proof}
On sait en effet (M, VI, 2.5) que la condition
$\operatorname{dim.proj}(M)\leq n$ est équivalente à
$\operatorname{Ext}^{n+1}_A(M,N)=0$ pour tout $A$-module $N$ \emph{de type
fini}~; pour voir que la condition de l'énoncé est aussi suffisante, on
raisonne par récurrence sur le nombre de générateurs $m$ de $N$~: il y a un
sous-module $N_1$ de $N$ engendré par $m-1$ éléments et tel que
$N_2=N/N_1$ soit monogène~; de la suite exacte
$0\longrightarrow N_1\longrightarrow N\longrightarrow N_2\longrightarrow0$,
on déduit alors la suite exacte
$\operatorname{Ext}^{n+1}_A(M,N_1)\longrightarrow
\operatorname{Ext}^{n+1}_A(M,N)\longrightarrow
\operatorname{Ext}^{n+1}_A(M,N_2)$ et l'hypothèse de récurrence montre que la
condition de l'énoncé entraîne bien $\operatorname{Ext}^{n+1}_A(M,N)=0$.
\end{proof}

\begin{corollary}[17.2.5]
\label{0.17.2.5-fr}
Soient $A$ un anneau noethérien, $M$ un $A$-module. On a
\begin{equation}
\label{0.17.2.5.1-fr}
\tag{17.2.5.1}
  \operatorname{dim.inj}_A(M)
  =\sup_{\mathfrak{m}}
  \bigl(\operatorname{dim.inj}_{A_{\mathfrak{m}}}(M_{\mathfrak{m}})\bigr)
\end{equation}
et si $M$ est \emph{de type fini}
\begin{equation}
\label{0.17.2.5.2-fr}
\tag{17.2.5.2}
  \operatorname{dim.proj}_A(M)
  =\sup_{\mathfrak{m}}
  \bigl(\operatorname{dim.proj}_{A_{\mathfrak{m}}}(M_{\mathfrak{m}})\bigr)
\end{equation}
où $\mathfrak{m}$ parcourt l'ensemble des idéaux premiers (ou l'ensemble des
idéaux maximaux) de $A$.
\end{corollary}

\begin{proof}
En effet, si $N$ est un $A$-module de type fini (donc de présentation finie),
on a, pour toute partie multiplicative $S$ de $A$ et pour tout $i\geq0$
\[
  S^{-1}\operatorname{Ext}^i_A(N,M)
  =\operatorname{Ext}^i_{S^{-1}A}(S^{-1}N,S^{-1}M)
\]
par platitude, en considérant une résolution libre de $M$ et en utilisant le
fait que la relation précédente est vraie pour $i=0$ (Bourbaki,
\emph{Alg.~comm.}, chap.~II, \S~2, no~7, prop.~19). En particulier
$\operatorname{Ext}^i_{A_{\mathfrak{m}}}
(A_{\mathfrak{m}}/\mathfrak{J}A_{\mathfrak{m}},M_{\mathfrak{m}})
=(\operatorname{Ext}^i_A(A/\mathfrak{J},M))_{\mathfrak{m}}$ pour tout idéal
premier $\mathfrak{m}$ et tout idéal $\mathfrak{J}$ de $A$~; si l'on tient
compte de (17.2.3) et de ce que tout
\oldpage[0\textsubscript{IV-1}]{44}
idéal de $A_{\mathfrak{m}}$ est de la forme $\mathfrak{J}A_{\mathfrak{m}}$
pour un idéal convenable $\mathfrak{J}$ de $A$, on déduit la formule
(17.2.5.1) de ce qui précède et de Bourbaki, \emph{Alg.~comm.}, chap.~II,
\S~3, no~3, th.~1. On procède de même pour (17.2.5.2) en utilisant cette fois
(17.2.4) et en échangeant les rôles de $M$ et de $N$.
\end{proof}

Pour les anneaux noethériens et les modules de type fini, l'étude de la
dimension projective ou de la dimension injective est donc ramenée par
(17.2.5) au cas des anneaux \emph{locaux}. On a alors le

\begin{lemma}[17.2.6]
\label{0.17.2.6-fr}
Soient $A$ un anneau local noethérien, $k$ son corps résiduel, $M$ un
$A$-module de type fini. Pour que $\operatorname{dim.proj}(M)\leq n$, il est
nécessaire que $\operatorname{Tor}^A_i(M,k)=0$ pour $i>n$, et il suffit que
$\operatorname{Tor}^A_{n+1}(M,k)=0$.
\end{lemma}

\begin{proof}
La nécessité est un cas particulier de la remarque (17.2.2 (ii)), appliquée au
foncteur covariant $M\mapsto k\otimes_A M$. Pour prouver que la condition est
suffisante, il s'agit, avec les notations de (17.2.1), d'établir que $R$ est
projectif lorsque les $P_i$ sont supposés de type fini~; or
$\operatorname{Tor}^A_{n+1}(M,k)$ est isomorphe à
$\operatorname{Tor}^A_1(R,k)$ (M, V, 7)~; et on sait que, puisque $R$ est de
type fini, la condition $\operatorname{Tor}^A_1(R,k)=0$ entraîne que $R$ est
libre (Bourbaki, \emph{Alg.~comm.}, chap.~II, \S~3, no~2, cor.~2 de la
prop.~5).
\end{proof}

\begin{corollary}[17.2.7]
\label{0.17.2.7-fr}
Sous les hypothèses de (17.2.6), soit $x$ un élément $M$-régulier de $A$
appartenant à l'idéal maximal de $A$~; alors on a
\begin{equation}
\label{0.17.2.7.1-fr}
\tag{17.2.7.1}
  \operatorname{dim.proj}(M/xM)=\operatorname{dim.proj}(M)+1.
\end{equation}
\end{corollary}

\begin{proof}
En effet, de la suite exacte
$0\longrightarrow M\xrightarrow{x}M\longrightarrow M/xM\longrightarrow0$
définie par la multiplication par $x$ dans $M$, on déduit la suite exacte des
Tor~:
\[
  \operatorname{Tor}^A_i(M,k)\xrightarrow{x}
  \operatorname{Tor}^A_i(M,k)\longrightarrow
  \operatorname{Tor}^A_i(M/xM,k)\longrightarrow
  \operatorname{Tor}^A_{i-1}(M,k)\xrightarrow{x}
  \operatorname{Tor}^A_{i-1}(M,k)
\]
pour tout $i\geq1$. Or, pour tout $A$-module $N$, l'homothétie de rapport $x$
dans $M\otimes_A N$ provient aussi bien de l'homothétie de rapport $x$ dans
$M$ que de l'homothétie de rapport $x$ dans $N$~; on en conclut aussitôt,
puisque l'homothétie de rapport $x$ dans $k$ est nulle par définition, que,
pour tout $i$, l'homomorphisme
$\operatorname{Tor}^A_i(M,k)\xrightarrow{x}\operatorname{Tor}^A_i(M,k)$ est
\emph{nul}~; autrement dit, on a la suite exacte
\[
  0\longrightarrow\operatorname{Tor}^A_i(M,k)\longrightarrow
  \operatorname{Tor}^A_i(M/xM,k)\longrightarrow
  \operatorname{Tor}^A_{i-1}(M,k)\longrightarrow0.
\]
Si $\operatorname{dim.proj}(M)=n$, on a
$\operatorname{Tor}^A_n(M,k)\neq0$ et
$\operatorname{Tor}^A_{n+1}(M,k)=\operatorname{Tor}^A_{n+2}(M,k)=0$ en vertu
de (17.2.6). Il résulte donc de ce qui précède que l'on a
$\operatorname{Tor}^A_{n+1}(M/xM,k)\neq0$ et
$\operatorname{Tor}^A_{n+2}(M/xM,k)=0$, donc
$\operatorname{dim.proj}(M/xM)=n+1$ par (17.2.6).
\end{proof}

\begin{proposition}[17.2.8]
\label{0.17.2.8-fr}
\textnormal{(M.~Auslander)} Soit $A$ un anneau. Les conditions suivantes sont
équivalentes~:
\begin{enumerate}
  \item[\rm a)] Tout $A$-module est de dimension projective $\leq n$.
  \item[\rm a')] Tout $A$-module de type fini est de dimension projective
  $\leq n$.
  \item[\rm b)] Tout $A$-module est de dimension injective $\leq n$.
  \item[\rm c)] Pour tout couple de $A$-modules $M$, $N$ on a
  $\operatorname{Ext}^{n+1}_A(M,N)=0$.
  \item[\rm c')] Pour tout couple de $A$-modules $M$, $N$ tel que $M$ soit de
  type fini (ou seulement monogène), on a
  $\operatorname{Ext}^{n+1}_A(M,N)=0$.
\end{enumerate}
\end{proposition}

\begin{proof}
\oldpage[0\textsubscript{IV-1}]{45}
Cela résulte aussitôt de (17.2.1) et (17.2.3).
\end{proof}

Le plus petit nombre $n$ (entier ou $+\infty$) pour lequel les conditions
équivalentes de (17.2.8) sont satisfaites est appelé \emph{dimension
cohomologique globale} (ou simplement \emph{dimension cohomologique}) de $A$
et noté $\operatorname{dim.coh}(A)$.

\begin{proposition}[17.2.9]
\label{0.17.2.9-fr}
Soit $A$ un anneau noethérien. Les conditions suivantes sont équivalentes~:
\begin{enumerate}
  \item[\rm a)] $\operatorname{dim.coh}(A)\leq n$.
  \item[\rm b)] Tout $A$-module de type fini est de dimension injective
  $\leq n$.
  \item[\rm c)] Pour tout couple de $A$-modules de type fini, $M$, $N$, on a
  $\operatorname{Ext}^{n+1}_A(M,N)=0$.
\end{enumerate}
\end{proposition}

\begin{proof}
Cela résulte aussitôt de la définition et de (17.2.4).
\end{proof}

\begin{corollary}[17.2.10]
\label{0.17.2.10-fr}
Si $A$ est un anneau noethérien, on a
\begin{equation}
\label{0.17.2.10.1-fr}
\tag{17.2.10.1}
  \operatorname{dim.coh}(A)
  =\sup_{\mathfrak{m}}
  \bigl(\operatorname{dim.coh}(A_{\mathfrak{m}})\bigr)
\end{equation}
où $\mathfrak{m}$ parcourt le spectre de $A$ (ou l'ensemble des idéaux
maximaux de $A$).
\end{corollary}

\begin{proof}
Cela résulte de (17.2.9) et (17.2.5).
\end{proof}

\begin{proposition}[17.2.11]
\label{0.17.2.11-fr}
Soient $A$ un anneau local noethérien, $k$ son corps résiduel. Pour que
$\operatorname{dim.coh}(A)\leq n$, il est nécessaire que
$\operatorname{Tor}^A_i(k,k)=0$, pour $i>n$, et il suffit que
$\operatorname{Tor}^A_{n+1}(k,k)=0$.
\end{proposition}

\begin{proof}
Compte tenu de (17.2.6), il suffit de prouver que les relations
$\operatorname{dim.coh}(A)\leq n$ et
$\operatorname{dim.proj}_A(k)\leq n$ sont équivalentes. Il est clair que la
première entraîne la seconde par définition. Inversement, si
$\operatorname{dim.proj}_A(k)\leq n$, on a
$\operatorname{Tor}^A_i(M,k)=0$ pour $i>n$ et pour tout $A$-module $M$ en
vertu de (17.2.2, (ii))~; donc $\operatorname{dim.proj}(M)\leq n$, ce qui
prouve la proposition en vertu de (17.2.8).
\end{proof}

\begin{corollary}[17.2.12]
\label{0.17.2.12-fr}
Soit $A$ un anneau noethérien. Pour que l'on ait
$\operatorname{dim.coh}(A)\leq n$, il est nécessaire que, pour tout idéal
maximal $\mathfrak{m}$ de $A$, on ait
$\operatorname{Tor}^{A_{\mathfrak{m}}}_i(A/\mathfrak{m},A/\mathfrak{m})=0$
pour $i>n$, et il suffit que ces relations soient vérifiées pour $i=n+1$.
\end{corollary}

\begin{proof}
Cela résulte aussitôt de (17.2.11) et (17.2.10).
\end{proof}

\begin{proposition}[17.2.13]
\label{0.17.2.13-fr}
Soient $A$, $B$ deux anneaux locaux noethériens,
$\varphi:A\longrightarrow B$ un homomorphisme local faisant de $B$ un
$A$-module plat. Alors on a
\begin{equation}
\label{0.17.2.13.1-fr}
\tag{17.2.13.1}
  \operatorname{dim.coh}(A)\leq\operatorname{dim.coh}(B).
\end{equation}
\end{proposition}

\begin{proof}
Supposons que $\operatorname{dim.coh}(B)=n$ soit finie~; il suffit de prouver
que pour tout couple $(M,N)$ de $A$-modules de type fini, on a
$\operatorname{Tor}^A_i(M,N)=0$ pour $i>n$ (17.2.11). Comme $B$ est un
$A$-module fidèlement plat ($0_{\rm I}$, 6.6.2), il revient au même
($0_{\rm I}$, 6.4.1) de montrer que l'on a
\begin{equation}
\label{0.17.2.13.2-fr}
\tag{17.2.13.2}
  \operatorname{Tor}^A_i(M,N)\otimes_A B=0.
\end{equation}
Or, si $L_\bullet=(L_j)$ est une résolution droite de $M$ par des $A$-modules
libres, il résulte de ce que $B$ est un $A$-module plat que
$L_\bullet\otimes_A B=(L_j\otimes_A B)$ est une résolution droite du
$B$-module $M\otimes_A B$ par des $B$-modules libres~; en outre, on a
$(L_j\otimes_A B)\otimes_B(N\otimes_A B)
=(L_j\otimes_A N)\otimes_A B$, d'où on conclut aussitôt que le premier membre
\oldpage[0\textsubscript{IV-1}]{46}
de (17.2.13.2) est égal à
$\operatorname{Tor}^B_i(M\otimes_A B,N\otimes_A B)$~; l'hypothèse sur $B$
entraîne que ce $B$-module est nul pour $i>n$ (17.2.2, (ii)), d'où la
conclusion.
\end{proof}

\begin{env}[17.2.14]
\label{0.17.2.14-fr}
Soient $(X,\mathscr{O}_X)$ un espace annelé, $\mathscr{F}$ un
$\mathscr{O}_X$-Module~; on appelle \emph{dimension projective} (resp.
\emph{injective}) \emph{ponctuelle} de $\mathscr{F}$ et on note
$\operatorname{dim.proj}(\mathscr{F})$ (resp.
$\operatorname{dim.inj}(\mathscr{F})$) le nombre
$\sup_{x\in X}(\operatorname{dim.proj}(\mathscr{F}_x))$ (resp.
$\sup_{x\in X}(\operatorname{dim.inj}(\mathscr{F}_x))$. On appelle
\emph{dimension cohomologique ponctuelle de $X$} et on note
$\operatorname{dim.coh}(X)$ le nombre
$\sup_{x\in X}(\operatorname{dim.coh}(\mathscr{O}_x))$. Il résulte de
(17.2.5) et (17.2.10) que lorsque $A$ est un anneau noethérien, $M$ un
$A$-module de type fini et $X=\operatorname{Spec}(A)$, on a
$\operatorname{dim.proj}(\widetilde{M})=\operatorname{dim.proj}(M)$,
$\operatorname{dim.inj}(\widetilde{M})=\operatorname{dim.inj}(M)$ et
$\operatorname{dim.coh}(X)=\operatorname{dim.coh}(A)$. On appelle
\emph{dimension projective} (resp. \emph{injective}) \emph{de $\mathscr{F}$ en
un point $x\in X$} la dimension projective (resp. injective) de
$\mathscr{F}_x$, \emph{dimension cohomologique de $X$ en un point $x$} la
dimension cohomologique de $\mathscr{O}_x$.
\end{env}

\begin{proposition}[17.2.15]
\label{0.17.2.15-fr}
Soient $X$, $Y$ deux espaces annelés en anneaux locaux noethériens,
$f:X\longrightarrow Y$ un morphisme plat. Si
$\operatorname{dim.coh}(X)\leq n$, alors $Y$ est de dimension cohomologique
$\leq n$ en tout point de $f(X)$.
\end{proposition}

\begin{proof}
Cela résulte aussitôt de (17.2.13).
\end{proof}

\subsection{Théorie cohomologique des anneaux réguliers}
\label{subsection:0.17.3-fr}

\begin{theorem}[17.3.1]
\label{0.17.3.1-fr}
\textnormal{(Hilbert-Serre)} Soit $A$ un anneau local noethérien. Pour que $A$
soit de dimension cohomologique finie, il faut et il suffit que $A$ soit
régulier~; on a alors
\begin{equation}
\label{0.17.3.1.1-fr}
\tag{17.3.1.1}
  \operatorname{dim.coh}(A)=\dim(A).
\end{equation}
\end{theorem}

\begin{proof}
Supposons $A$ régulier~; soient $\mathfrak{m}$ son idéal maximal,
$\mathbf{x}=(x_i)_{1\leq i\leq n}$ un système régulier de paramètres de $A$
(17.1.6)~; considérons le complexe de l'algèbre extérieure
$K_\bullet(\mathbf{x})$ (\textbf{III}, 1.1.1), qui est formé de $A$-modules
libres, avec $K_i(\mathbf{x})=0$ pour $i>n$~; comme $(x_i)$ est une suite
régulière, on a $H_i(K_\bullet(\mathbf{x}))=0$ pour $i>0$ (\textbf{III},
1.1.4 et 1.1.3.3) et
$H_0(K_\bullet(\mathbf{x}))=A/(x_1A+\cdots+x_nA)=A/\mathfrak{m}=k$~; les
$K_i(\mathbf{x})$ forment par suite une \emph{résolution libre} de $k$ de
longueur $n$. Or, le fait que les $x_i$ appartiennent à $\mathfrak{m}$
entraîne aussitôt que dans le complexe
$K_\bullet(\mathbf{x},k)=K_\bullet(\mathbf{x})\otimes_A k$, l'opérateur bord
est nul en toutes dimensions, si bien que l'on a, par définition,
$\operatorname{Tor}^A_i(k,k)=\bigwedge^i(k^n)$~; l'égalité (17.3.1.1) résulte
donc aussitôt de (17.2.11) (ce résultat est essentiellement le «~théorème des
syzygies~» de Hilbert).

Montrons maintenant que si $A$ est un anneau local noethérien,
$\mathfrak{m}$ son idéal maximal, et si
$\operatorname{dim.proj}_A(\mathfrak{m})$ est \emph{finie}, alors $A$ est
régulier, ce qui achèvera de démontrer (17.3.1). Nous procéderons par
récurrence sur $n=\operatorname{rg}_k(\mathfrak{m}/\mathfrak{m}^2)$. Pour
$n=0$, on a $\mathfrak{m}=0$ et l'assertion est triviale.
\end{proof}

\begin{lemma}[17.3.1.2]
\label{0.17.3.1.2-fr}
\textnormal{(Nagata)} Si tout élément de
$\mathfrak{m}-\mathfrak{m}^2$ est diviseur de zéro dans $A$, il existe
$c\neq0$ dans $A$ tel que $c\mathfrak{m}=0$ (autrement dit on a
$\mathfrak{m}\in\operatorname{Ass}(A)$).
\end{lemma}

\begin{proof}
On peut se borner au cas où $\mathfrak{m}\neq0$, donc
$\mathfrak{m}\neq\mathfrak{m}^2$. L'hypothèse entraîne que
$\mathfrak{m}-\mathfrak{m}^2$ est contenu dans la réunion des idéaux
$\mathfrak{p}_i$ de $\operatorname{Ass}(A)$ (Bourbaki,
\emph{Alg.~comm.}, chap.~IV, \S~1, no~1, cor.~3 de la prop.~2)~; donc
$\mathfrak{m}$ est contenu dans la réunion de $\mathfrak{m}^2$
\oldpage[0\textsubscript{IV-1}]{47}
et des $\mathfrak{p}_i$, donc dans l'un des $\mathfrak{p}_i$, puisque
$\mathfrak{m}\not\subset\mathfrak{m}^2$ (Bourbaki, \emph{Alg.~comm.},
chap.~II, \S~1, no~1, prop.~2)~; $\mathfrak{m}$ étant maximal, cela prouve le
lemme.
\end{proof}

\begin{lemma}[17.3.1.3]
\label{0.17.3.1.3-fr}
Si $a\in\mathfrak{m}-\mathfrak{m}^2$, $\mathfrak{m}/Aa$ est isomorphe à un
facteur direct de $\mathfrak{m}/a\mathfrak{m}$.
\end{lemma}

\begin{proof}
Il existe un système minimal de générateurs de $\mathfrak{m}$ contenant $a$
(dont les images dans $\mathfrak{m}/\mathfrak{m}^2$ forment une base de ce
$k$-espace vectoriel)~; soit $\mathfrak{b}$ l'idéal engendré par les éléments
de ce système autres que $a$. Comme la relation $xa\in\mathfrak{b}$ entraîne
$x\in\mathfrak{m}$ (en considérant l'image de $xa$ dans
$\mathfrak{m}/\mathfrak{m}^2$), on a
$\mathfrak{b}\cap Aa\subset a\mathfrak{m}$ d'où un homomorphisme
$\mathfrak{b}/(\mathfrak{b}\cap Aa)\longrightarrow
\mathfrak{m}/a\mathfrak{m}$ déduit de l'injection
$\mathfrak{b}\longrightarrow\mathfrak{m}$ par passage aux quotients, et qui
est \emph{injectif}. En outre, $\mathfrak{b}+Aa=\mathfrak{m}$, et
l'homomorphisme composé
\[
  \mathfrak{m}/Aa=(\mathfrak{b}+Aa)/Aa
  \xrightarrow{\sim}\mathfrak{b}/(\mathfrak{b}\cap Aa)
  \longrightarrow\mathfrak{m}/a\mathfrak{m}
  \longrightarrow\mathfrak{m}/Aa
\]
est l'identité~; d'où le lemme.
\end{proof}

\begin{lemma}[17.3.1.4]
\label{0.17.3.1.4-fr}
Soient $A$ un anneau local noethérien, $\mathfrak{m}$ son idéal maximal, $E$
un $A$-module de type fini et de dimension projective finie. Si
$a\in\mathfrak{m}$ est $A$-régulier et $E$-régulier, alors $E/aE$ est un
$(A/aA)$-module de dimension projective finie, au plus égale à
$\operatorname{dim.proj}_A(E)$.
\end{lemma}

\begin{proof}
Raisonnons par récurrence sur $h=\operatorname{dim.proj}_A(E)$, le cas $h=0$
étant trivial puisque alors $E$ est un $A$-module projectif, donc $E/aE$ un
$(A/aA)$-module projectif. Il existe une suite exacte
\[
  0\longrightarrow N\longrightarrow L\longrightarrow E\longrightarrow0
\]
où $L$ est libre et $\operatorname{dim.proj}_A(N)=h-1$ (17.2.2, (iii)), $N$
étant de type fini. En outre, la suite
\[
  0\longrightarrow N/aN\longrightarrow L/aL\longrightarrow E/aE
  \longrightarrow0
\]
est exacte (15.1.18). Comme $a$ est $A$-régulier, il est aussi $N$-régulier
puisque $L$ est libre~; l'hypothèse de récurrence entraîne que $N/aN$ est un
$(A/aA)$-module de dimension projective $\leq h-1$, et comme $L/aL$ est un
$(A/aA)$-module libre, $E/aE$ est un $(A/aA)$-module de dimension projective
$\leq h$.
\end{proof}

Examinons maintenant deux cas~:

\medskip\noindent\textbf{I.} --- Supposons d'abord que tout élément de
$\mathfrak{m}-\mathfrak{m}^2$ soit diviseur de zéro dans $A$, auquel cas
(17.3.1.2) il existe $c\neq0$ dans $A$ tel que $c\mathfrak{m}=0$. Montrons
alors que $\mathfrak{m}=0$. S'il n'en était pas ainsi, notons d'abord que
$\mathfrak{m}$ ne pourrait être un $A$-module projectif, car il serait libre
(Bourbaki, \emph{Alg.~comm.}, chap.~II, \S~3, no~2, cor.~2 de la prop.~5),
ce qui contredit la relation $c\mathfrak{m}=0$. On aurait donc
$n=\operatorname{dim.coh}(A)\geq1$. Comme
$\mathfrak{m}\in\operatorname{Ass}(A)$, il existerait une suite exacte de
$A$-homomorphismes
\[
  0\longrightarrow k\longrightarrow A\longrightarrow E\longrightarrow0.
\]
Mais cela est absurde, car en vertu de la relation $n\geq1$, la suite exacte
des Tor donnerait la suite exacte
$0\longrightarrow\operatorname{Tor}^A_{n+1}(E,k)\longrightarrow
\operatorname{Tor}^A_n(k,k)\longrightarrow0$~; or on a
$\operatorname{Tor}^A_{n+1}(E,k)=0$ (17.2.2, (ii)) et
$\operatorname{Tor}^A_n(k,k)\neq0$ (17.2.11) et nous avons abouti à une
contradiction.

\medskip\noindent\textbf{II.} --- On peut donc se borner au cas où il existe
$a\in\mathfrak{m}-\mathfrak{m}^2$ qui est un élément $A$-régulier, et par
suite aussi $\mathfrak{m}$-régulier. Considérons l'anneau $A'=A/aA$ et son
idéal
\oldpage[0\textsubscript{IV-1}]{48}
maximal $\mathfrak{m}'=\mathfrak{m}/aA$~; il est clair que
$\operatorname{rg}_k(\mathfrak{m}'/\mathfrak{m}'^2)=n-1$. En vertu de
(17.3.1.4), $\mathfrak{m}/a\mathfrak{m}$ est un $A'$-module de dimension
projective finie, donc il en est de même de
$\mathfrak{m}'=\mathfrak{m}/aA$, qui en est un facteur direct (17.3.1.3 et
17.2.1). L'hypothèse de récurrence entraîne par suite que $A'=A/aA$ est
régulier, ce qui prouve par (17.1.8) que $A$ est régulier.

\begin{corollary}[17.3.2]
\label{0.17.3.2-fr}
Si $A$ est un anneau local régulier, $A_{\mathfrak{p}}$ est régulier pour tout
idéal premier $\mathfrak{p}$ de $A$.
\end{corollary}

\begin{proof}
En effet, on a vu (17.2.10) que
$\operatorname{dim.coh}(A_{\mathfrak{p}})\leq
\operatorname{dim.coh}(A)$, donc la conclusion résulte aussitôt de (17.3.1).
\end{proof}

\begin{proposition}[17.3.3]
\label{0.17.3.3-fr}
Soient $A$, $B$ deux anneaux locaux noethériens,
$\varphi:A\longrightarrow B$ un homomorphisme local, $\mathfrak{m}$ (resp.
$\mathfrak{n}$) l'idéal maximal de $A$ (resp. $B$),
$k=A/\mathfrak{m}$ (resp. $k'=B/\mathfrak{n}$) le corps résiduel de $A$
(resp. $B$). On a donc un homomorphisme canonique de $k'$-espaces vectoriels
\begin{equation}
\label{0.17.3.3.1-fr}
\tag{17.3.3.1}
  \psi:(\mathfrak{m}/\mathfrak{m}^2)\otimes_k k'
  \longrightarrow\mathfrak{n}/\mathfrak{n}^2.
\end{equation}
\begin{enumerate}
  \item[\rm (i)] Si $B$ est régulier et si $B$ est un $A$-module plat, $A$
  est régulier.
  \item[\rm (ii)] Les conditions suivantes sont équivalentes~:
  \begin{enumerate}
    \item[\rm a)] $B$ est régulier et l'homomorphisme $\psi$ (17.3.3.1) est
    injectif.
    \item[\rm b)] $A$ et $B$ sont réguliers, et pour un système régulier de
    paramètres $(x_i)_{1\leq i\leq m}$ de $A$, les $\varphi(x_i)$ font partie
    d'un système régulier de paramètres de $B$ (auquel cas cette propriété a
    lieu pour \emph{tout} système régulier de paramètres de $A$).
    \item[\rm c)] $B$ et $B\otimes_A k$ sont réguliers, et $B$ est un
    $A$-module plat.
    \item[\rm d)] $A$ et $B\otimes_A k$ sont réguliers, et $B$ est un
    $A$-module plat.
    \item[\rm e)] $A$ et $B\otimes_A k$ sont réguliers, et l'on a
    \begin{equation}
    \label{0.17.3.3.2-fr}
    \tag{17.3.3.2}
      \dim(B)=\dim(A)+\dim(B\otimes_A k).
    \end{equation}
  \end{enumerate}
\end{enumerate}
\end{proposition}

\begin{proof}
\textnormal{(i)} On a
$\operatorname{dim.coh}(A)\leq\operatorname{dim.coh}(B)$ par (17.2.13), donc
il suffit d'appliquer (17.3.1).

\textnormal{(ii)} Lorsque $A$ est supposé régulier et que $(x_i)$ est un
système régulier de paramètres de $A$, dire que $B$ est un $A$-module
\emph{plat} équivaut, en vertu de (15.1.21) à dire que la suite des
$y_i=\varphi(x_i)$ est $B$-régulière (puisque $B\otimes_A k$ est un
$k$-module plat). D'autre part, comme $\dim(A)=m$, et que la suite $(y_i)$
engendre l'idéal $\mathfrak{m}B$, la relation (17.3.3.2), qui s'écrit aussi
$\dim(B/\mathfrak{m}B)=\dim(B)-m$, équivaut à dire que la suite
$(y_i)_{1\leq i\leq m}$ fait partie d'un système de paramètres de $B$
(16.3.7). On voit donc que les conditions b), d) et e) équivalent
respectivement aux suivantes~:
\begin{enumerate}
  \item[\rm b')] $A$ et $B$ sont réguliers et $(y_i)_{1\leq i\leq m}$ fait
  partie d'un système régulier de paramètres de $B$.
  \item[\rm d')] $A$ est régulier,
  $B/(\sum_{i=1}^m y_iB)$ est régulier et la suite $(y_i)$ est
  $B$-régulière.
  \item[\rm e')] $A$ est régulier,
  $B/(\sum_{i=1}^m y_iB)$ est régulier et la suite $(y_i)$ fait partie d'un
  système de paramètres de $B$.
\end{enumerate}
Or, b') et e') sont équivalentes en vertu de (17.1.7), et comme d') entraîne
e') (16.4.1) et est entraînée par b') (17.1.7), elle leur est équivalente. La
conjonction
\oldpage[0\textsubscript{IV-1}]{49}
de b) et d) implique trivialement c), et en vertu de (i), c) entraîne d)~; on
a donc prouvé l'équivalence de b), c), d) et e). En outre, il est clair que b)
entraîne a), les classes des $x_i$ dans
$\mathfrak{m}/\mathfrak{m}^2$ formant alors une base de ce $k$-espace
vectoriel et les classes des $y_i$ dans $\mathfrak{n}/\mathfrak{n}^2$ un
système libre de ce $k'$-espace vectoriel. Reste donc à prouver que a)
entraîne e). Posons $V=\mathfrak{m}/\mathfrak{m}^2$,
$W=\mathfrak{n}/\mathfrak{n}^2$~; il est immédiat que l'on a
$\psi(V\otimes_k k')=(\mathfrak{n}^2+\mathfrak{m}B)/\mathfrak{n}^2$. On a
d'une part, en vertu de (16.3.9)
\begin{equation}
\label{0.17.3.3.3-fr}
\tag{17.3.3.3}
  \dim(B)\leq\dim(A)+\dim(B\otimes_A k)~;
\end{equation}
en second lieu, d'après (16.2.6), on a
\begin{equation}
\label{0.17.3.3.4-fr}
\tag{17.3.3.4}
  \dim(A)\leq\operatorname{rg}_{k'}(V\otimes_k k').
\end{equation}
Enfin, dans l'anneau local $B\otimes_A k=B/\mathfrak{m}B$,
$\mathfrak{n}'=\mathfrak{n}/\mathfrak{m}B$ est l'idéal maximal, $k'$ le corps
résiduel et $\mathfrak{n}'/\mathfrak{n}'^2$ est isomorphe à
$\mathfrak{n}/(\mathfrak{n}^2+\mathfrak{m}B)
=W/\psi(V\otimes_k k')$~; on a donc, d'après (16.2.6)
\begin{equation}
\label{0.17.3.3.5-fr}
\tag{17.3.3.5}
  \dim(B\otimes_A k)
  \leq\operatorname{rg}_{k'}W
  -\operatorname{rg}_{k'}\psi(V\otimes_k k').
\end{equation}
Enfin, puisque $B$ est supposé régulier, on a
$\dim(B)=\operatorname{rg}_{k'}W$ (17.1.1)~; on conclut donc de (17.3.3.3),
(17.3.3.4) et (17.3.3.5) que l'on a
\[
  \operatorname{rg}_{k'}W
  \leq\dim(A)+\dim(B\otimes_A k)
  \leq\operatorname{rg}_{k'}W
  +\operatorname{rg}_{k'}(V\otimes_k k')
  -\operatorname{rg}_{k'}\psi(V\otimes_k k')
\]
et dire que les deux termes extrêmes de cette inégalité sont égaux signifie
que $\psi$ est \emph{injectif}. La condition a) entraîne donc nécessairement
que dans chacune des relations (17.3.3.3), (17.3.3.4) et (17.3.3.5), les deux
membres sont \emph{égaux}~; or, l'égalité dans (17.3.3.4) (resp. (17.3.3.5))
signifie que $A$ (resp. $B\otimes_A k$) est \emph{régulier} (17.1.1)~; on a
donc bien prouvé que a) entraîne e).
\end{proof}

\begin{proposition}[17.3.4]
\label{0.17.3.4-fr}
Soient $A$ un anneau local régulier de dimension $n$, $\mathfrak{m}$ son idéal
maximal. Pour tout $A$-module non nul de type fini $M$, on a
\begin{equation}
\label{0.17.3.4.1-fr}
\tag{17.3.4.1}
  \operatorname{prof}(M)+\operatorname{dim.proj}(M)=n.
\end{equation}
\end{proposition}

\begin{proof}
Raisonnons par récurrence sur $m=\operatorname{prof}(M)$. Si $m=0$, on sait
(16.4.6, (i)) qu'il existe un sous-$A$-module $N$ de $M$ isomorphe à
$k=A/\mathfrak{m}$~; appliquant la suite exacte des Tor à la suite exacte
$0\longrightarrow N\longrightarrow M\longrightarrow M/N\longrightarrow0$,
on obtient une suite exacte
\[
  \operatorname{Tor}^A_{n+1}(M/N,k)\longrightarrow
  \operatorname{Tor}^A_n(N,k)\longrightarrow
  \operatorname{Tor}^A_n(M,k),
\]
et l'hypothèse que $A$ est régulier entraîne
$\operatorname{Tor}^A_{n+1}(M/N,k)=0$ ((17.3.1.1) et (17.2.6)), donc
$\operatorname{Tor}^A_n(k,k)$ est isomorphe à un sous-$A$-module de
$\operatorname{Tor}^A_n(M,k)$~; appliquant de nouveau (17.3.1.1) et (17.2.6),
on voit que $\operatorname{Tor}^A_n(M,k)\neq0$, donc
$\operatorname{dim.proj}(M)\geq n$ (17.2.6)~; mais comme
$\operatorname{dim.proj}(M)\leq n$ par (17.3.1.1), on a bien
$\operatorname{dim.proj}(M)=n$.

Supposons maintenant $m>0$, et soit $x$ un élément $M$-régulier appartenant à
$\mathfrak{m}$~; on sait alors (16.4.6, (i)) que l'on a
$\operatorname{prof}(M/xM)=m-1$, et d'autre part (17.2.7)
$\operatorname{dim.proj}(M/xM)=\operatorname{dim.proj}(M)+1$~; l'hypothèse de
récurrence prouve aussitôt la relation (17.3.4.1).
\end{proof}

\oldpage[0\textsubscript{IV-1}]{50}

\begin{corollary}[17.3.5]
\label{0.17.3.5-fr}
\begin{enumerate}
  \item[\rm(i)] Soit $A$ un anneau local régulier de dimension $n$~; pour
  qu'un $A$-module $M\neq0$ de type fini soit libre, il faut et il suffit que
  $M$ soit un $A$-module de Cohen-Macaulay de dimension $n$.
  \item[\rm(ii)] Soient $A$ un anneau local régulier, $B$ un anneau local,
  $\rho:A\to B$ un homomorphisme local faisant de $B$ un $A$-module de type
  fini. Pour que $B$ soit un $A$-module libre, il faut et il suffit que $B$
  soit un anneau de Cohen-Macaulay et que $\rho$ soit injectif (ou, ce qui
  revient au même (16.1.5), que $\dim(B)=\dim(A)$).
\end{enumerate}
\end{corollary}

\begin{proof}
(i) Cela résulte de (17.3.4.1) et du fait que pour un $A$-module de type fini,
il revient au même de dire que ce module est projectif ou libre (Bourbaki,
\emph{Alg. comm.}, chap.~II, \S~3, n\textsuperscript{o}~2, cor.~2 de la
prop.~5)~; les $A$-modules libres de type fini $M$ sont donc caractérisés par
la relation $\operatorname{dim.proj}(M)=0$ (17.2.2, (i)).

(ii) Dire que $B$ est un anneau de Cohen-Macaulay équivaut à dire que $B$ est
un $A$-module de Cohen-Macaulay (16.5.3), donc il suffit d'appliquer (i),
puisque $\dim(B)=\dim(A)$ (16.1.5).
\end{proof}

\begin{env}[17.3.6]
\label{0.17.3.6-fr}
On dit qu'un anneau noethérien $A$ est \emph{régulier} si pour tout idéal
premier $\mathfrak{p}$ de $A$, l'anneau local $A_{\mathfrak{p}}$ est
régulier~; lorsque $A$ lui-même est un anneau local, il résulte de (17.3.2)
que cette définition est équivalente à celle de (17.1.2). Pour que $A$ soit
régulier, il suffit, en vertu de (17.3.2), que $A_{\mathfrak{m}}$ soit
régulier pour tout idéal maximal $\mathfrak{m}$ de $A$. En outre, il résulte
aussitôt de cette définition que pour toute partie multiplicative $S$ de
$A$, $S^{-1}A$ est régulier.
\end{env}

\begin{proposition}[17.3.7]
\label{0.17.3.7-fr}
Si $A$ est un anneau noethérien régulier, tout anneau de polynômes
$A[T_1,\ldots,T_n]$ est régulier.
\end{proposition}

\begin{proof}
Il suffit évidemment de prouver que l'anneau de polynômes $B=A[T]$ est
régulier~; comme $B$ est un $A$-module libre, pour tout idéal premier
$\mathfrak{q}$ de $B$, $B_{\mathfrak{q}}$ est un
$A_{\mathfrak{p}}$-module plat, où $\mathfrak{p}=\mathfrak{q}\cap A$
($0_{\mathrm I}$, 6.3.1)~; il suffit donc, en vertu de (17.1.10), de prouver
que $B_{\mathfrak{q}}/\mathfrak{p}B_{\mathfrak{q}}$ est régulier, et comme
cet anneau est un anneau local en un idéal premier de
$A_{\mathfrak{p}}[T]/\mathfrak{p}A_{\mathfrak{p}}[T]=k[T]$ (où $k$ est le
corps résiduel $A_{\mathfrak{p}}/\mathfrak{p}A_{\mathfrak{p}}$), il suffit
de prouver que $C=k[T]$ est régulier~; or, $C$ étant principal, les anneaux
locaux en les idéaux premiers de $C$ sont des anneaux de valuation discrète
ou un corps (pour l'idéal $(0)$), donc réguliers (17.1.4), ce qui achève la
démonstration.
\end{proof}

\begin{corollary}[17.3.8]
\label{0.17.3.8-fr}
Si $A$ est un anneau régulier, tout anneau de séries formelles
$A[[T_1,\ldots,T_n]]$ est régulier.
\end{corollary}

\begin{proof}
Soit $\mathfrak{J}$ l'idéal engendré par les $T_i$ dans l'anneau de polynômes
$A[T_1,\ldots,T_n]$~; comme ce dernier est régulier par (17.3.7), et que
$A[[T_1,\ldots,T_n]]$ est le complété de $A[T_1,\ldots,T_n]$ pour la
topologie $\mathfrak{J}$-préadique (Bourbaki, \emph{Alg. comm.}, chap.~III,
\S~2, n\textsuperscript{o}~12, \emph{Exemple 1}), la conclusion résulte du
\end{proof}

\begin{lemma}[17.3.8.1]
\label{0.17.3.8.1-fr}
Soient $A$ un anneau régulier, $\mathfrak{J}$ un idéal de $A$,
$\widehat{A}$ le séparé complété de $A$ pour la topologie
$\mathfrak{J}$-préadique. Alors $\widehat{A}$ est régulier.
\end{lemma}

\begin{proof}
Il suffit en effet (17.3.6) de voir que pour tout idéal maximal
$\mathfrak{n}$ de $\widehat{A}$, l'anneau local
$\widehat{A}_{\mathfrak{n}}$ est régulier~; on sait (Bourbaki,
\emph{Alg. comm.}, chap.~III, \S~3, n\textsuperscript{o}~4, prop.~8) que
$\mathfrak{n}$
\oldpage[0\textsubscript{IV-1}]{51}
est de la forme $\mathfrak{m}\widehat{A}$, où $\mathfrak{m}$ est un idéal
maximal de $A$ contenant $\mathfrak{J}$, et que l'homomorphisme canonique
$A\to\widehat{A}$ donne un homomorphisme \emph{injectif}
$A_{\mathfrak{m}}\to\widehat{A}_{\mathfrak{n}}$, tel que la topologie
$\mathfrak{n}\widehat{A}_{\mathfrak{n}}$-préadique de
$\widehat{A}_{\mathfrak{n}}$ induise sur $A_{\mathfrak{m}}$ la topologie
$\mathfrak{m}A_{\mathfrak{m}}$-préadique, et que $A_{\mathfrak{m}}$ soit
dense dans $\widehat{A}_{\mathfrak{n}}$ pour la topologie
$\mathfrak{n}\widehat{A}_{\mathfrak{n}}$-préadique~; par suite les
\emph{complétés} des anneaux locaux noethériens $A_{\mathfrak{m}}$ et
$\widehat{A}_{\mathfrak{n}}$ sont canoniquement \emph{isomorphes}. Or, par
hypothèse $A_{\mathfrak{m}}$ est régulier, donc il en est de même de son
complété (17.1.5) et comme le complété de $\widehat{A}_{\mathfrak{n}}$ est
régulier, il en est de même de $\widehat{A}_{\mathfrak{n}}$ par (17.1.5).
\end{proof}

\begin{corollary}[17.3.9]
\label{0.17.3.9-fr}
Soit $A$ un anneau noethérien, quotient d'un anneau noethérien régulier $B$.
Si $C$ est une $A$-algèbre de type fini, tout anneau de fractions $S^{-1}C$
de $C$ est quotient d'un anneau régulier.
\end{corollary}

\begin{proof}
En effet, $C$ est quotient d'un anneau de polynômes $A[T_1,\ldots,T_n]$~;
comme $A=B/\mathfrak{J}$, où $\mathfrak{J}$ est un idéal de $B$, on a
$A[T_1,\ldots,T_n]=B[T_1,\ldots,T_n]/\mathfrak{J}B[T_1,\ldots,T_n]$~; en
vertu de (17.3.7), on peut donc se borner au cas où $C$ est quotient d'un
anneau régulier $B'$~; mais si $S'$ est l'image réciproque de $S$ dans $B'$,
$S'$ est une partie multiplicative de $B'$ et $C$ est un anneau quotient de
$S'^{-1}B'$, si bien que finalement tout revient à voir que lorsque $C$ est
\emph{régulier}, il en est de même de $S^{-1}C$, ce qu'on a vu en (17.3.6).
\end{proof}

L'importance des quotients d'anneaux réguliers tient entre autres à la
propriété précédente et au fait que ce sont des anneaux \emph{caténaires}
(16.5.12).

Tous les anneaux importants dans les applications à la Géométrie algébrique
sont des quotients d'anneaux réguliers.

% FIN SOURCE ega4/ega0-17-fr.tex

%
% DEBUT SOURCE ega4/ega0-18-fr.tex
\section{Compléments sur les extensions d'algèbres}
\label{section:0.18-fr}

Ce paragraphe rassemble un certain nombre de constructions fonctorielles sur
les anneaux, qui seront utilisées de façon répétée aux \S\S~19 et 20~; il ne
contient aucun résultat non trivial.

\subsection{Images réciproques d'anneaux augmentés}
\label{subsection:0.18.1-fr}

\begin{env}[18.1.1]
\label{0.18.1.1-fr}
Étant donné un anneau $A$ (non nécessairement commutatif), la catégorie des
\emph{$A$-anneaux} a pour objets les couples $(B,\rho)$ formés d'un anneau
$B$ et d'un homomorphisme d'anneaux $\rho:A\to B$, pour morphismes (dits aussi
\emph{$A$-homomorphismes}) les homomorphismes d'anneaux $u:B\to C$ tels que,
si $\rho:A\to B$ et $\sigma:A\to C$ sont les homomorphismes (dits
\emph{structuraux}) définissant la structure de $A$-anneau sur $B$ et $C$
respectivement, le diagramme
\begin{equation}
\label{0.18.1.1.1-fr}
\tag{18.1.1.1}
  \xymatrix{
    B\ar[r]^u & C\\
    & A\ar[lu]^\rho\ar[u]_\sigma
  }
\end{equation}
soit commutatif. Le \emph{noyau} $\mathfrak{J}$ de $u$ est un idéal bilatère
qui est un $A$-\emph{bimodule} (pour $\rho$).

Si $f:A'\to A$ est un homomorphisme d'anneaux, pour tout $A$-anneau
$(B,\rho)$, $(B,\rho\circ f)$ est un $A'$-anneau, et si $u:B\to C$ est un
$A$-homomorphisme du $A$-anneau $(B,\rho)$
\oldpage[0\textsubscript{IV-1}]{52}
dans le $A$-anneau $(C,\sigma)$, c'est aussi un $A'$-homomorphisme du
$A'$-anneau $(B,\rho\circ f)$ dans le $A'$-anneau $(C,\sigma\circ f)$~; on
définit ainsi un foncteur canonique de la catégorie des $A$-anneaux dans celle
des $A'$-anneaux.
\end{env}

\begin{env}[18.1.2]
\label{0.18.1.2-fr}
Soient $B$, $E$, $F$ trois $A$-anneaux, $f:E\to B$, $g:F\to B$ deux
$A$-homomorphismes. Rappelons que l'on appelle \emph{produit fibré} de $E$ et
$F$ sur $B$ (pour les $A$-homomorphismes $f$ et $g$) et que l'on note
$E\times_BF$ le sous-anneau $G$ de l'anneau produit $E\times F$, formé des
couples $(x,y)$ tels que $f(x)=g(y)$~; les restrictions $p_1:G\to E$,
$p_2:G\to F$ des projections $\operatorname{pr}_1$ et
$\operatorname{pr}_2$ à $G$ sont encore appelées les \emph{projections
canoniques}~; la structure de $A$-anneau de $G$ est définie par
l'homomorphisme $\alpha\mapsto(\rho(\alpha),\sigma(\alpha))$ (où
$\rho:A\to E$ et $\sigma:A\to F$ sont les homomorphismes structuraux), qui
applique effectivement $A$ dans $G$ en vertu de (18.1.1). La propriété
caractéristique de $E\times_BF$ est que, pour tout couple de
$A$-homomorphismes $u:C\to E$, $v:C\to F$ tels que
$f\circ u=g\circ v$, il existe un $A$-homomorphisme et un seul $w:C\to G$
tel que $u=p_1\circ w$, $v=p_2\circ w$. On peut encore dire que
$E\times_BF$ est \emph{limite projective} du système projectif formé de $B$,
$E$, $F$ et des $A$-homomorphismes $f$, $g$, dans la catégorie des
$A$-anneaux ($0_{\mathrm{III}}$, 8.1.9).
\end{env}

\begin{env}[18.1.3]
\label{0.18.1.3-fr}
Soit $\mathfrak{J}$ l'idéal bilatère de $E$, noyau de $f$~; il est immédiat
que le noyau $\mathfrak{J}'$ de $p_2:G\to F$ est l'idéal formé des éléments
$(x,0)$ où $x\in\mathfrak{J}$~; la restriction
$i_1:\mathfrak{J}'\to\mathfrak{J}$ de $p_1$ est donc un
\emph{isomorphisme} d'anneau \emph{sans élément unité}, et aussi un
isomorphisme de $A$-bimodules. De même, si $\mathfrak{R}$ est le noyau de
$g$, le noyau $\mathfrak{R}'$ de $p_1:G\to E$ est l'idéal des éléments
$(0,y)$ où $y\in\mathfrak{R}$, et la restriction
$i_2:\mathfrak{R}'\to\mathfrak{R}$ de $p_2$ est un isomorphisme (au même
sens). Enfin, il est clair que le noyau du $A$-homomorphisme
$q=f\circ p_1=g\circ p_2$ de $E\times_BF$ dans $B$ est
$\mathfrak{J}\times\mathfrak{R}=\mathfrak{J}'\oplus\mathfrak{R}'$, de sorte
que l'on a le diagramme commutatif
\begin{equation}
\label{0.18.1.3.1-fr}
\tag{18.1.3.1}
  \xymatrix@C=2.2pc@R=1.8pc{
    0\ar[dr] & & 0\ar[d] & 0\ar[d]\\
    & \mathfrak{J}'\oplus\mathfrak{R}'\ar[r]\ar[d]\ar[dr]
      & \mathfrak{R}'\ar[r]^{i_2}_{\sim}\ar[d]^j
      & \mathfrak{R}\ar[d]^{p_2}\\
    0\ar[r] & \mathfrak{J}'\ar[r]\ar[d]^{i_1}_{\sim}
      & G\ar[r]\ar[d]^{p_1}\ar[dr]^q & F\ar[d]^g\\
    0\ar[r] & \mathfrak{J}\ar[r] & E\ar[r]_f & B
  }
\end{equation}

Les définitions et résultats de (18.1.2) et (18.1.3) s'étendent aussitôt au
produit fibré d'une \emph{famille} quelconque $(E_\lambda)_{\lambda\in I}$ de
$A$-anneaux défini par une famille de $A$-homomorphismes
$f_\lambda:E_\lambda\to B$. Nous laissons la formulation de ces résultats au
lecteur.
\end{env}

\begin{env}[18.1.4]
\label{0.18.1.4-fr}
Conformément à la terminologie de (M, VIII), nous appellerons
\emph{$A$-anneau augmenté sur $B$} un $A$-anneau $E$ muni d'un
$A$-homomorphisme \emph{surjectif} (dit \emph{augmentation} de $E$),
$f:E\to B$~; le noyau $\mathfrak{J}$ de $f$ est appelé
\emph{l'idéal d'augmentation}. On dit que le $A$-anneau augmenté $E$ est
\emph{trivial} s'il existe un $A$-homomorphisme de $A$-anneaux $s:B\to E$
qui est \emph{inverse à droite} de l'augmentation $f:E\to B$ (autrement dit
$f\circ s=1_B$). La suite exacte de $A$-bimodules
\[
  0\longrightarrow\mathfrak{J}\longrightarrow E
  \xrightarrow{f}B\longrightarrow0
\]
\oldpage[0\textsubscript{IV-1}]{53}
est alors \emph{scindée}~; autrement dit, on peut identifier le $A$-bimodule
$E$ à $B\times\mathfrak{J}$, et la multiplication dans $E$ est alors donnée
par
\[
  (b,z)(b',z')=(bb',bz'+zb'+zz'),
\]
$\mathfrak{J}$ étant considéré comme $B$-bimodule au moyen de $s:B\to E$.
\end{env}

\begin{env}[18.1.5]
\label{0.18.1.5-fr}
Avec les notations de (18.1.2), supposons que le $A$-homomorphisme $f$ soit
\emph{surjectif}, autrement dit fasse de $E$ un $A$-anneau \emph{augmenté sur
$B$}. Alors il est clair que $p_2:G\to F$ est aussi \emph{surjectif},
autrement dit définit sur $G$ une structure de $A$-anneau \emph{augmenté sur
$F$}, que l'on appelle \emph{l'image réciproque} par $g:F\to B$ de l'anneau
augmenté $E$.
\end{env}

\begin{proposition}[18.1.6]
\label{0.18.1.6-fr}
Pour que l'image réciproque $G=E\times_BF$ par $g:F\to B$ du $A$-anneau
augmenté $E$ soit un $A$-anneau augmenté trivial, il faut et il suffit qu'il
existe un $A$-homomorphisme $u:F\to E$ rendant commutatif le diagramme
\[
  \xymatrix{
    & F\ar[ld]_u\ar[d]^g\\
    E\ar[r]_f & B
  }
\]
\end{proposition}

\begin{proof}
La condition est évidemment nécessaire, en prenant $u=p_1\circ s$, où
$s:F\to G$ est un $A$-homomorphisme inverse à droite de $p_2$. Inversement,
s'il existe un $A$-homomorphisme $u$ vérifiant la condition de l'énoncé,
l'existence de l'inverse à droite $s$ de $p_2$ résulte de la propriété
universelle du produit fibré (18.1.2) appliquée aux $A$-homomorphismes
$u:F\to E$ et $1_F:F\to F$.

Ce résultat entraîne en particulier que si $E$ est un anneau augmenté
\emph{trivial}, il en est de même de toutes ses images réciproques.
\end{proof}

\begin{env}[18.1.7]
\label{0.18.1.7-fr}
Reprenons la situation décrite en (18.1.2) et (18.1.3)~; on a évidemment sur
$\mathfrak{R}$ une structure de $G$-bimodule, provenant de sa structure de
$F$-bimodule et de l'homomorphisme d'anneaux $p_2:G\to F$. Comme $i_2$ est
bijectif on a en outre une injection
$\theta=j\circ i_2^{-1}:\mathfrak{R}\to G$, qui est un homomorphisme de
$G$-bimodules. Nous allons voir inversement que, lorsqu'on suppose en outre
que $f$ et $g$ sont \emph{surjectifs}, autrement dit que $E$ et $F$ sont des
$A$-anneaux \emph{augmentés} sur $B$, la donnée d'un tel homomorphisme
$\theta$ permet de reconstituer l'anneau augmenté $E$ (sur $B$) à partir de
l'anneau augmenté $G$ (sur $F$).

De façon précise, donnons-nous un $A$-anneau $F$ augmenté sur $B$ et un
$A$-anneau $G$ augmenté sur $F$, les idéaux d'augmentation étant notés
$\mathfrak{R}$ et $\mathfrak{J}'$ respectivement~:
\begin{equation}
\label{0.18.1.7.1-fr}
\tag{18.1.7.1}
  \xymatrix@C=2pc@R=1.5pc{
    & & 0\ar[d] & & \\
    & & \mathfrak{R}\ar[ld]^\theta\ar[d] & & \\
    0\ar[r] & \mathfrak{J}'\ar[r] & G\ar[r]_h & F\ar[r]\ar[d]^g & 0\\
    & & & B\ar[d] & \\
    & & & 0 &
  }
\end{equation}

\oldpage[0\textsubscript{IV-1}]{54}
L'homomorphisme $h$ définit sur tous les idéaux de $F$ (et en particulier sur
$\mathfrak{R}$) une structure de $G$-bimodule. Soit
$\theta:\mathfrak{R}\to G$ un homomorphisme de $G$-bimodules rendant
commutatif le diagramme (18.1.7.1)~; cela implique que $\theta$ est
\emph{injectif} et que $\mathfrak{J}'\cap\theta(\mathfrak{R})=0$~; comme
$\theta(\mathfrak{R})$ est un idéal bilatère de $G$, et que
$h(\theta(\mathfrak{R}))=\mathfrak{R}$, l'homomorphisme composé
$g\circ h:G\to B$ se factorise en
$G\xrightarrow{q}G/\theta(\mathfrak{R})\xrightarrow{f}B$, où $f$ est
surjectif~; en outre l'image $\mathfrak{J}$ de $\mathfrak{J}'$ dans
$E=G/\theta(\mathfrak{R})$ est le noyau de $f$ (celui de $g\circ h$ étant
$\mathfrak{J}'\oplus\theta(\mathfrak{R})$), et la restriction à
$\mathfrak{J}'$ de l'homomorphisme canonique $q:G\to E$ est injective. On
peut donc former le \emph{produit fibré} $G'=E\times_BF$, et comme les deux
$A$-homomorphismes $q:G\to E$ et $h:G\to F$ sont tels que
$f\circ q=g\circ h$, ils définissent un unique $A$-homomorphisme $u:G'\to G$
par la propriété universelle de $G'$ (18.1.2)~; nous allons voir que $u$ est
\emph{bijectif}. Il suffit de le prouver lorsque $u$ est considéré comme
homomorphisme de $A$-bimodules~; mais on notera alors que $u$ est compatible
avec les filtrations finies sur $G'$ et $G$ formées respectivement par $G'$
et $\mathfrak{J}''=\operatorname{Ker}(G'\to F)$, et par $G$ et
$\mathfrak{J}'$~; en outre, on a vu (18.1.3) que
$\operatorname{gr}_0u:F\to F$ est l'identité et que
$\operatorname{gr}_1u:\mathfrak{J}''\to\mathfrak{J}'$ est bijectif, donc $u$
lui-même est bijectif (Bourbaki, \emph{Alg. comm.}, chap.~III, \S~2,
n\textsuperscript{o}~8, cor.~3 du th.~1).
\end{env}

\subsection{Extensions d'un anneau par un bimodule}
\label{subsection:0.18.2-fr}

\begin{env}[18.2.1]
\label{0.18.2.1-fr}
Soient $E$ un $A$-anneau augmenté sur $B$, $f:E\to B$ l'augmentation,
$\mathfrak{J}=\operatorname{Ker}(f)$ l'idéal d'augmentation. Si l'on a
$\mathfrak{J}^2=0$, $\mathfrak{J}$ est non seulement un $E$-bimodule mais
aussi un $B$-\emph{bimodule} puisque $B$ est isomorphe à
$E/\mathfrak{J}$~; de façon précise, tout $b\in B$ est de la forme $f(x)$
avec $x\in E$, et si $z,z'$ sont dans $\mathfrak{J}$ on a
$(x+z)z'=xz'$ (resp. $z'(x+z)=z'x$) de sorte que la valeur de $xz'$ (resp.
$z'x$) ne dépend pas de l'élément $x$ de $f^{-1}(b)$, et peut s'écrire $bz'$
(resp. $z'b$) ce qui définit la structure de $B$-bimodule considérée.
Inversement, si $\mathfrak{J}$ est muni d'une structure de $B$-bimodule telle
que $xz'=f(x)z'$ et $z'x=z'f(x)$ pour $x\in E$ et
$z'\in\mathfrak{J}$, il est clair que $\mathfrak{J}^2=0$.
\end{env}

\begin{env}[18.2.2]
\label{0.18.2.2-fr}
On appelle \emph{$A$-extension d'un $A$-anneau $B$ par un $B$-bimodule $L$}
une \emph{suite exacte} d'homomorphismes de $A$-bimodules
\[
  0\longrightarrow L\xrightarrow{j}E\xrightarrow{f}B\longrightarrow0
\]
où $E$ est un $A$-anneau, $f$ un $A$-homomorphisme d'anneaux et l'on a, pour
$x\in E$ et $z\in L$,
\[
  j(f(x)z)=xj(z),\qquad j(zf(x))=j(z)x
\]
d'où résulte (18.2.1) que $j(L)$ est un idéal bilatère de carré nul de $E$.
Par abus de langage on dira aussi que $E$ est une extension de $B$ par $L$.
On dit que deux $A$-extensions $E$, $E'$ de $B$ par $L$ sont
\emph{$A$-équivalentes} s'il existe un isomorphisme de $A$-anneaux
$u:E\xrightarrow{\sim}E'$ (dit aussi \emph{$A$-équivalence de
$A$-extensions}) rendant commutatif le diagramme
\begin{equation}
\label{0.18.2.2.1-fr}
\tag{18.2.2.1}
  \xymatrix@C=2.1pc@R=1.2pc{
    & & E\ar[dr]\ar[dd]_u\ar@{}[dd]|{\wr} & & \\
    0\ar[r] & L\ar[ur]\ar[dr] & & B\ar[r] & 0\\
    & & E'\ar[ur] & &
  }
\end{equation}
\end{env}

\oldpage[0\textsubscript{IV-1}]{55}

\begin{env}[18.2.3]
\label{0.18.2.3-fr}
On dit qu'une $A$-extension $E$ de $B$ par un $B$-bimodule $L$ est
\emph{$A$-triviale} si $E$ est un $A$-anneau augmenté (sur $B$) trivial
(18.1.4). On définit sur le $A$-bimodule produit $B\times L$ une structure de
$A$-anneau en posant $(x,s)(y,t)=(xy,xt+sy)$, comme on le vérifie aussitôt,
et il est immédiat que les applications canoniques $j:L\to B\times L$ et
$f:B\times L\to B$ définissent une $A$-extension de $B$ par $L$ qui est
$A$-triviale, l'application canonique $g:B\to B\times L$ étant un
$A$-homomorphisme inverse à droite de $f$. On dit que cette extension est
\emph{l'extension triviale type} de $B$ par $L$ et on la note $D_B(L)$~; il
est immédiat que toute $A$-extension $A$-triviale de $B$ par $L$ est
$A$-équivalente à $D_B(L)$.

On notera que toute extension du $A$-anneau $A$ lui-même par un $A$-bimodule
est nécessairement $A$-triviale.
\end{env}

\begin{env}[18.2.4]
\label{0.18.2.4-fr}
Étant données deux $A$-extensions
$L\xrightarrow{j}E\xrightarrow{f}B$,
$L'\xrightarrow{j'}E'\xrightarrow{f'}B'$, un \emph{morphisme} de la première
dans la seconde est par définition un triplet d'homomorphismes de
$A$-bimodules $(u,v,w)$ tel que le diagramme
\begin{equation}
\label{0.18.2.4.1-fr}
\tag{18.2.4.1}
  \xymatrix{
    0\ar[r] & L\ar[r]^j\ar[d]_w & E\ar[r]^f\ar[d]_u
      & B\ar[r]\ar[d]_v & 0\\
    0\ar[r] & L'\ar[r]_{j'} & E'\ar[r]_{f'} & B'\ar[r] & 0
  }
\end{equation}
soit commutatif, $u$ et $v$ étant des $A$-homomorphismes \emph{d'anneaux} et
$w$ étant tel que $w(bz)=v(b)w(z)$ et $w(zb)=w(z)v(b)$ pour $z\in L$ et
$b\in B$ (en d'autres termes, le couple $(v,w)$ constitue un
\emph{di-homomorphisme} du $B$-bimodule $L$ dans le $B'$-bimodule $L'$)~; il
est clair que si $(u',v',w')$ est un morphisme de
$L'\to E'\to B'$ dans une $A$-extension $L''\to E''\to B''$,
$(u'\circ u,v'\circ v,w'\circ w)$ est encore un morphisme, ce qui justifie la
terminologie.

La considération des deux carrés commutatifs du diagramme (18.2.4.1) va nous
conduire à deux opérations sur les extensions de $A$-anneaux.
\end{env}

\begin{env}[18.2.5]
\label{0.18.2.5-fr}
En premier lieu, considérons une $A$-extension $E'$ de $B'$ par $L'$
\[
  0\longrightarrow L'\xrightarrow{j'}E'\xrightarrow{f'}B'\longrightarrow0
\]
et un $A$-homomorphisme d'anneaux $v:B\to B'$, et soit
$F=E'\times_{B'}B$ \emph{l'image réciproque} par $v$ du $A$-anneau augmenté
$E'$ (18.1.5), de sorte que l'on a un diagramme commutatif
\[
  \xymatrix{
    0\ar[r] & L_0\ar[r]\ar[d]^i\ar@{}[d]|{\wr} & F\ar[r]^{p_2}\ar[d]^{p_1}
      & B\ar[r]\ar[d]^v & 0\\
    0\ar[r] & L'\ar[r]_{j'} & E'\ar[r]_{f'} & B'\ar[r] & 0
  }
\]
dont les lignes sont exactes, $p_1$ et $p_2$ étant les homomorphismes
canoniques~; on a vu (18.1.3) que $i$ est bijectif et il résulte aussi de la
définition (18.1.2) que $L_0^2=0$, de sorte que l'on peut considérer que $F$
est une $A$-extension de $B$ par $L'$, que l'on appelle \emph{image
réciproque par $v$} de l'extension $E'$ de $B'$ par $L'$ ($L'$ étant
naturellement considéré comme $B$-bimodule au moyen de l'homomorphisme
d'anneaux $v$). Le caractère fonctoriel
\oldpage[0\textsubscript{IV-1}]{56}
du produit fibré vis-à-vis de chacun des facteurs montre en outre que si l'on
a un morphisme de deux extensions de $B'$
\[
  \xymatrix{
    0\ar[r] & L'_1\ar[r]\ar[d]^h & E'_1\ar[r]\ar[d]^g
      & B'\ar[r]\ar[d]^{1_{B'}} & 0\\
    0\ar[r] & L'_2\ar[r] & E'_2\ar[r] & B'\ar[r] & 0
  }
\]
on en déduit un morphisme \emph{image réciproque par $v$}
\[
  \xymatrix{
    0\ar[r] & L'_1\ar[r]\ar[d]^h & E'_1\times_{B'}B\ar[r]
      \ar[d]^{g\times_{B'}1_B} & B\ar[r]\ar[d]^{1_B} & 0\\
    0\ar[r] & L'_2\ar[r] & E'_2\times_{B'}B\ar[r] & B\ar[r] & 0.
  }
\]
En particulier, si $E'_1$ et $E'_2$ sont des $A$-extensions
$A$-équivalentes de $B'$ par $L'$, leurs images réciproques par $v$ sont des
$A$-extensions $A$-équivalentes de $B$ par $L'$.

La définition du produit fibré montre que lorsque l'on a un morphisme
(18.2.4.1) de $A$-extensions, il se \emph{factorise à travers} l'image
réciproque de $E'$ par $v$~; de façon précise, il existe un
$A$-homomorphisme unique $u_0:E\to F=E'\times_{B'}B$ rendant commutatif le
diagramme
\[
  \xymatrix{
    0\ar[r] & L\ar[r]^j\ar[d]^{w_0} & E\ar[r]^f\ar[d]^{u_0}
      & B\ar[r]\ar[d]^{1_B} & 0\\
    0\ar[r] & L_0\ar[r]\ar[d]^i\ar@{}[d]|{\wr} & F\ar[r]^{p_2}\ar[d]^{p_1}
      & B\ar[r]\ar[d]^v & 0\\
    0\ar[r] & L'\ar[r]_{j'} & E'\ar[r]_{f'} & B'\ar[r] & 0
  }
\]
où $w_0$ est la restriction de $u_0$ à $L$ et
$p_1\circ u_0=u$, $i\circ w_0=w$.
\end{env}

\begin{env}[18.2.6]
\label{0.18.2.6-fr}
Étudions en particulier les images réciproques d'extensions par des
$A$-homomorphismes \emph{surjectifs}. Considérons un $A$-homomorphisme
surjectif $v:B\to B'$, une $A$-extension $F$ de $B$ par un $B$-bimodule $L$,
et l'idéal bilatère $\mathfrak{R}$ de $B$, noyau de $v$ (qui peut être
considéré comme $F$-bimodule au moyen de l'augmentation $F\to B$)~; on a vu
(18.1.7) que tout homomorphisme de $F$-bimodules
$\theta:\mathfrak{R}\to F$ rendant commutatif le diagramme
\[
  \xymatrix{
    & \mathfrak{R}\ar[ld]^\theta\ar[d]\\
    F\ar[r] & B
  }
\]
détermine une extension $E'$ de $B'=B/\mathfrak{R}$ par $L$ dont l'image
réciproque par $v:B\to B'$ est équivalente à $F$, et que toute $A$-extension
$E'$ de $B'$ par $L$ ayant cette dernière propriété s'obtient ainsi (à une
$A$-équivalence près). En outre, pour que deux homomorphismes
$\theta_1,\theta_2$ de $F$-bimodules de $\mathfrak{R}$ dans $F$ donnent deux
  $A$-extensions $A$-équivalentes de $B$ par $L$, il faut et il suffit qu'il
existe une $A$-équivalence $u$ de la $A$-extension $F$ sur elle-même telle que
\oldpage[0\textsubscript{IV-1}]{57}
$\theta_2=u\circ\theta_1$~; cela résulte aussitôt de ce qu'on a vu dans
(18.2.5) et de la définition de la bijection canonique de $F$ sur le produit
fibré $E'\times_{B'}B$ (18.1.7).
\end{env}

\begin{env}[18.2.7]
\label{0.18.2.7-fr}
Considérons maintenant le carré de gauche de (18.2.4.1), et rappelons d'abord
la notion de \emph{somme amalgamée} dans la catégorie des $A$-bimodules~:
étant donnés trois $A$-bimodules $X$, $Y$, $Z$ et deux $A$-homomorphismes
$f:X\to Y$, $g:X\to Z$, la somme amalgamée $Y\oplus_XZ$ est \emph{limite
inductive} du système inductif formé par $X$, $Y$, $Z$ et les
$A$-homomorphismes $f$, $g$ dans la catégorie des $A$-bimodules
($0_{\mathrm{III}}$, 8.1.11). On définit ce $A$-bimodule comme quotient du
produit $Y\times Z$ par le sous-$A$-bimodule $M$ image de $X$ par
l'homomorphisme $x\mapsto(f(x),-g(x))$. Sa propriété caractéristique est que,
pour tout couple d'homomorphismes de $A$-bimodules $u:Y\to T$, $v:Z\to T$
tels que $u\circ f=v\circ g$, il existe un homomorphisme et un seul
$w:Y\oplus_XZ\to T$ tel que $u=w\circ j_1$ et $v=w\circ j_2$, où
$j_1:Y\to Y\oplus_XZ$ et $j_2:Z\to Y\oplus_XZ$ sont les applications
canoniques.
\end{env}

\begin{env}[18.2.8]
\label{0.18.2.8-fr}
Considérons maintenant une $A$-extension $E$ de $B$ par un $B$-bimodule
$L$~:
\[
  0\longrightarrow L\xrightarrow{j}E\xrightarrow{f}B\longrightarrow0
\]
et soit d'autre part $w:L\to L'$ un homomorphisme de $B$-bimodules. Soit $H$
le $A$-bimodule somme amalgamée $E\oplus_LL'$~; montrons comment on peut munir
ce $A$-bimodule d'une structure de $A$-anneau et définir une $A$-extension
\[
  0\longrightarrow L'\longrightarrow H\longrightarrow B\longrightarrow0.
\]

Notons pour cela que $L'$ est muni d'une structure de $E$-\emph{bimodule} au
moyen de l'homomorphisme d'augmentation $E\to B$~; on peut donc former la
$A$-extension triviale type $G=D_E(L')$ (18.2.3). Considérons alors
l'application $\theta:z\mapsto(j(z),-w(z))$ de $L$ dans $G$~; c'est un
homomorphisme de $G$-bimodules ($L$ étant considéré comme $G$-bimodule au
moyen de l'homomorphisme canonique $\rho:G\to E$). En effet, pour
$(x,z')\in G$ et $z\in L$, on a
$(j(z),-w(z))(x,z')=(j(z)x,j(z)z'-w(z)f(x))$~; or
$j(z)z'=f(j(z))z'=0$ par définition de la structure de $E$-bimodule sur
$L'$, et $w(z)f(x)=w(zf(x))$ par définition de la structure de $B$-bimodule
sur $L'$~; on vérifie de même que $\theta$ est un homomorphisme de
$G$-module à gauche. On peut alors appliquer au diagramme commutatif
\[
  \xymatrix@C=2pc@R=1.4pc{
    & & 0\ar[d] & & \\
    & & L\ar[ld]^\theta\ar[d]^j & & \\
    0\ar[r] & L'\ar[r] & G\ar[r]_p & E\ar[r]\ar[d]^f & 0\\
    & & & B\ar[d] & \\
    & & & 0 &
  }
\]
le résultat de (18.1.7). Comme $H=G/\theta(L)$ par définition, notre
assertion est une conséquence immédiate de (18.1.7).

\oldpage[0\textsubscript{IV-1}]{58}

On dit que la $A$-extension $H$ de $B$ par $L'$ est \emph{déduite de $E$ au
moyen de l'homomorphisme} $w:L\to L'$. Le caractère fonctoriel de la somme
amalgamée en chacun des sommandes montre en outre que si l'on a un morphisme
d'extensions
\[
  \xymatrix{
    0\ar[r] & L\ar[r]\ar[d]^{1_L}\ar@{}[d]|{\wr}
      & E_1\ar[r]\ar[d]^g & B_1\ar[r]\ar[d]^h & 0\\
    0\ar[r] & L\ar[r] & E_2\ar[r] & B_2\ar[r] & 0
  }
\]
on en déduit canoniquement un morphisme d'extensions
\[
  \xymatrix{
    0\ar[r] & L'\ar[r]\ar[d]^{1_{L'}}\ar@{}[d]|{\wr}
      & E_1\oplus_LL'\ar[r]\ar[d]^{g\oplus_L1_{L'}}
      & B_1\ar[r]\ar[d]^h & 0\\
    0\ar[r] & L'\ar[r] & E_2\oplus_LL'\ar[r] & B_2\ar[r] & 0.
  }
\]
En particulier, si $E_1$ et $E_2$ sont des $A$-extensions $A$-équivalentes de
$B$ par $L$, les extensions de $B$ par $L'$ qu'on en déduit au moyen de $w$
sont $A$-équivalentes.

Lorsque l'on a un morphisme (18.2.4.1) de $A$-extensions, il se \emph{factorise
à travers} la $A$-extension $H$ de $B$ par $L'$ déduite de $E$ au moyen de
l'homomorphisme $w:L\to L'$ ($L'$ étant considéré comme $B$-bimodule au moyen
de l'homomorphisme $v:B\to B'$)~: en effet, la définition de la somme amalgamée
montre qu'il existe un $A$-homomorphisme unique $u_0:H\to E'$ de
$A$-bimodules, rendant commutatif le diagramme
\[
  \xymatrix{
    0\ar[r] & L\ar[r]^j\ar[d]^{w_0} & E\ar[r]^f\ar[d]^{j_1}
      & B\ar[r]\ar[d]^{1_B}\ar@{}[d]|{\wr} & 0\\
    0\ar[r] & L'_0\ar[r]^{j_2}\ar[d]^{j_0}\ar@{}[d]|{\wr}
      & H\ar[r]^{f_0}\ar[d]^{u_0} & B\ar[r]\ar[d]^v & 0\\
    0\ar[r] & L'\ar[r]_{j'} & E'\ar[r]_{f'} & B'\ar[r] & 0
  }
\]
avec $u=u_0\circ j_1$, $w=j_0\circ w_0$, $j_1$ et $j_2$ étant les
homomorphismes canoniques~; on vérifie immédiatement que $u_0$ est aussi un
homomorphisme d'anneaux.

Notons enfin les propriétés fonctorielles relatives aux extensions triviales~:
\end{env}

\begin{proposition}[18.2.9]
\label{0.18.2.9-fr}
Soient $B$, $B'$ deux $A$-anneaux, $L$ un $B$-bimodule, $L'$ un
$B'$-bimodule, $v:B\to B'$ un $A$-homomorphisme d'anneaux, $w:L\to L'$ un
homomorphisme de $A$-bimodules tel que $(v,w)$ soit un di-homomorphisme de
bimodules. Alors, il existe un $A$-homomorphisme unique d'anneaux
$u:D_B(L)\to D_{B'}(L')$ rendant commutatif les diagrammes
\[
  \xymatrix@C=2.2pc@R=1.6pc{
    D_B(L)\ar[r]^u & D_{B'}(L') & & D_B(L)\ar[r]^u & D_{B'}(L')\\
    B\ar[u]\ar[r]_v & B'\ar[u] & & L\ar[u]\ar[r]_w & L'\ar[u]
  }
\]
où les flèches verticales sont les injections canoniques.
\end{proposition}

\oldpage[0\textsubscript{IV-1}]{59}

\begin{proof}
En effet, $u$ ne peut être que l'application
$(x,s)\mapsto(v(x),w(s))$ et il reste à vérifier que c'est un
$A$-homomorphisme d'anneaux, ce qui résulte trivialement de la définition
(18.2.3). On notera que $u$ rend aussi commutatif le diagramme
\[
  \xymatrix@C=2.2pc@R=1.5pc{
    D_B(L)\ar[r]^u\ar[d] & D_{B'}(L')\ar[d]\\
    B\ar[r]_v & B'
  }
\]
où cette fois les flèches verticales sont les augmentations.
\end{proof}

\begin{proposition}[18.2.10]
\label{0.18.2.10-fr}
Soient $B$ un $A$-anneau, $L$ un $B$-bimodule, $E$ une $A$-extension de $B$
par $L$. On définit une application bijective de l'ensemble $G$ des
homomorphismes de $B$-anneaux de $B$ dans $E$ (autrement dit, l'ensemble des
$A$-homomorphismes inverses à droite de l'augmentation $E\to B$) sur
l'ensemble $G'$ des $A$-équivalences de $D_B(L)$ sur $E$ en faisant
correspondre à tout $g\in G$ la $A$-équivalence
$g':(x,s)\mapsto g(x)+s$~; l'application réciproque fait correspondre à tout
$g'\in G'$ le $B$-homomorphisme $x\mapsto g'(x,0)$.
\end{proposition}

\begin{proof}
Ceci résulte aussitôt des définitions.
\end{proof}

\subsection{Le groupe des classes de $A$-extensions}
\label{subsection:0.18.3-fr}

\begin{env}[18.3.1]
\label{0.18.3.1-fr}
Considérons un $A$-anneau $B$ et un $B$-bimodule $L$ fixés~; alors la relation
«~$E$ et $E'$ sont $A$-équivalentes~» entre $A$-extensions $E$, $E'$ de $B$
par $L$ est une relation d'équivalence, et pour cette relation on peut parler
de l'\emph{ensemble} des classes de $A$-extensions $A$-équivalentes de $B$ par
$L$. Il suffit pour le voir de remarquer que si, pour tout $x\in B$, $c_x$ est
un élément de $E$ dont l'image dans $B$ est $x$, tout $z\in E$ s'écrit d'une
manière et d'une seule sous la forme $c_x+t$, où $t\in L$, et l'on peut écrire
\[
  c_{x+y}=c_x+c_y+\varphi(x,y),\qquad
  c_{xy}=c_xc_y+\psi(x,y),
\]
où $\varphi(x,y)$ et $\psi(x,y)$ sont des éléments de $L$, les applications
$\varphi$ et $\psi$ de $B\times B$ dans $L$ devant vérifier des conditions
exprimant que $E$ est un $A$-anneau, qu'il est inutile d'écrire ici. Toute
$A$-extension de $B$ par $L$ est donc $A$-équivalente à une $A$-extension dont
$B\times L$ est l'ensemble sous-jacent, d'où on tire aussitôt notre conclusion.
\end{env}

\begin{env}[18.3.2]
\label{0.18.3.2-fr}
Le $A$-anneau $B$ étant fixé, désignons provisoirement par $T(L)$ l'ensemble
des classes de $A$-extensions de $B$ par $L$. Pour tout $B$-homomorphisme
$w:L\to L'$ de $B$-bimodules, on définit canoniquement une application
$T(w):T(L)\to T(L')$ en faisant correspondre à la classe d'une $A$-extension
$E$ de $B$ par $L$ la classe de la $A$-extension $E\oplus_LL'$ déduite de $E$
par $w$, en vertu de (18.2.8). Si $w':L'\to L''$ est un second homomorphisme
de $B$-bimodules, on a en outre
\begin{equation}
\label{0.18.3.2.1-fr}
\tag{18.3.2.1}
  T(w'\circ w)=T(w')\circ T(w).
\end{equation}
En effet, on sait qu'il existe un isomorphisme canonique de $A$-bimodules
\[
  E\oplus_LL''\xrightarrow{\sim}(E\oplus_LL')\oplus_{L'}L''
\]
en vertu des propriétés générales de limites inductives (cf. par exemple
($\mathrm{I}$, 3.3.9)), et il est immédiat de vérifier que c'est bien une
$A$-équivalence de $A$-extensions, d'où

\oldpage[0\textsubscript{IV-1}]{60}

la relation (18.3.2.1). Si $\mathbf C_B$ désigne la catégorie des
$B$-bimodules, on voit qu'on a ainsi défini un \emph{foncteur covariant}
$T:\mathbf C_B\to\mathbf{Ens}$.
\end{env}

\begin{env}[18.3.3]
\label{0.18.3.3-fr}
Considérons maintenant une famille $(L_\alpha)_{\alpha\in I}$ de
$B$-bimodules et leur produit $L=\prod_{\alpha\in I}L_\alpha$~; les
projections $\operatorname{pr}_\alpha:L\to L_\alpha$ définissent des
applications
\[
  T(\operatorname{pr}_\alpha):T(L)\to T(L_\alpha),
\]
d'où une application canonique
\begin{equation}
\label{0.18.3.3.1-fr}
\tag{18.3.3.1}
  \prod_\alpha T(\operatorname{pr}_\alpha):T(L)
  \longrightarrow\prod_{\alpha\in I}T(L_\alpha).
\end{equation}
Nous allons voir que cette application est \emph{bijective}. En effet, pour
tout $\alpha\in I$, soit $E_\alpha$ une $A$-extension de $B$ par $L_\alpha$,
et soit $F$ le \emph{produit fibré} des $E_\alpha$ sur $B$ (18.1.3)~; il est
immédiat que $F$ est une $A$-extension de $B$ par
$L=\prod_{\alpha\in I}L_\alpha$ et que si l'on remplace chaque $E_\alpha$ par
une $A$-extension $A$-équivalente $E'_\alpha$, le produit fibré $F'$ des
$E'_\alpha$ sur $B$ est $A$-équivalente à $F$. On a ainsi défini une
application $\prod_{\alpha\in I}T(L_\alpha)\to T(L)$ et il est clair que cette
application est réciproque de (18.3.3.1), d'où notre assertion. On identifiera
canoniquement $T(L)$ à $\prod_{\alpha\in I}T(L_\alpha)$ par l'application
(18.3.3.1). On vérifie en outre immédiatement que si
$(L'_\alpha)_{\alpha\in I}$ est une seconde famille de $B$-bimodules et, pour
chaque $\alpha$, $w_\alpha:L_\alpha\to L'_\alpha$ un homomorphisme de
$B$-bimodules, alors, en posant
$w=\prod_\alpha w_\alpha:L\to L'=\prod_{\alpha\in I}L'_\alpha$, $T(w)$
s'identifie à $\prod_{\alpha\in I}T(w_\alpha)$ quand on fait l'identification
précédente.
\end{env}

\begin{env}[18.3.4]
\label{0.18.3.4-fr}
Cela étant, pour un $B$-bimodule $L$, l'addition est un homomorphisme
$s:L\times L\to L$ de $B$-bimodules, et il en est de même de la symétrie
$t:L\to L$ de la loi additive de $L$. On en déduit une loi de composition
\[
  T(s):T(L)\times T(L)\longrightarrow T(L)
\]
dans $T(L)$ en vertu de (18.3.3), et cette loi est une loi de \emph{groupe
commutatif} dont $T(t)$ est la symétrie, comme il résulte de la définition d'un
objet en groupes au moyen de diagrammes commutatifs
($\mathbf 0_{\mathrm{III}}$, 8.2.5 et 8.2.6). Nous désignerons par
$\operatorname{Exan}_A(B,L)$ le groupe commutatif ainsi défini et nous dirons
que c'est le \emph{groupe des classes de $A$-extensions de $B$ par $L$}.
\end{env}

\begin{env}[18.3.5]
\label{0.18.3.5-fr}
Désignons par $\mathbf K$ la catégorie dont les objets sont les triplets
$(A,B,L)$ où $A$ est un anneau, $B$ un $A$-anneau et $L$ un $B$-bimodule~;
les \emph{morphismes} de cette catégorie sont les triplets $(u,v,w)$ où
$u:A'\to A$ et $v:B'\to B$ sont deux homomorphismes d'anneaux rendant
commutatif le carré de gauche du diagramme
\[
  \xymatrix@C=2pc@R=1.5pc{
    A\ar[r] & B & L\ar[d]^w\\
    A'\ar[u]^u\ar[r] & B'\ar[u]^v & L'
  }
\]
où les flèches horizontales sont les homomorphismes structuraux~; enfin
$w:L\to L'$ est un homomorphisme de groupes commutatifs tel que l'on ait
$w(v(b')z)=b'w(z)$ et

\oldpage[0\textsubscript{IV-1}]{61}

$w(zv(b'))=w(z)b'$ quels que soient $z\in L$ et $b'\in B'$ (en d'autres
termes, $w$ est un homomorphisme de $B'$-bimodules lorsqu'on munit $L$ de la
structure de $B'$-bimodule définie par $v$). La composition des morphismes est
définie par
$(u',v',w')\circ(u,v,w)=(u\circ u',v\circ v',w'\circ w)$, qui se justifie
aussitôt. Nous nous proposons de montrer que
\begin{equation}
\label{0.18.3.5.1-fr}
\tag{18.3.5.1}
  (A,B,L)\longmapsto\operatorname{Exan}_A(B,L)
\end{equation}
est un \emph{foncteur covariant} de la catégorie $\mathbf K$ dans la catégorie
$\mathbf{Ab}$ des groupes commutatifs. Il s'agit donc pour tout triplet
$(u,v,w)$ comme ci-dessus, de définir un homomorphisme de groupes commutatifs
\[
  (u,v,w)_*:\operatorname{Exan}_A(B,L)
  \longrightarrow\operatorname{Exan}_{A'}(B',L').
\]
En vertu de la définition des morphismes dans $\mathbf K$, on peut écrire
\[
  (u,v,w)=(1_{A'},1_{B'},w)\circ(1_{A'},v,1_L)
  \circ(u,1_B,1_L)
\]
où, dans le premier facteur, $L$ est muni de sa structure de $B'$-bimodule
définie par $v$~; nous définirons donc d'abord $(u,v,w)_*$ lorsque deux des
homomorphismes $u$, $v$, $w$ sont réduits à l'identité.
\end{env}

\begin{env}[18.3.6]
\label{0.18.3.6-fr}
Nous prendrons d'abord pour $(1_A,1_B,w)_*$ l'application
\begin{equation}
\label{0.18.3.6.1-fr}
\tag{18.3.6.1}
  w_*:\operatorname{Exan}_A(B,L)\longrightarrow
  \operatorname{Exan}_A(B,L')
\end{equation}
notée $T(w)$ dans (18.3.2)~; il est immédiat de vérifier que c'est un
homomorphisme de groupes, cette propriété s'exprimant par la commutativité de
diagrammes, transformés par $T$ de diagrammes analogues pour $L$ et $L'$.

L'application $(1_A,v,1_L)_*$ est l'application
\begin{equation}
\label{0.18.3.6.2-fr}
\tag{18.3.6.2}
  v^*:\operatorname{Exan}_A(B,L)\longrightarrow
  \operatorname{Exan}_A(B',L)
\end{equation}
définie de la façon suivante~: si $E$ est une $A$-extension de $B$ par $L$,
on a vu que $E\times_BB'$ est une $A$-extension de $B'$ par $L$ (18.2.5) et
que si l'on remplace $E$ par une $A$-extension $A$-équivalente $E'$,
$E'\times_BB'$ est $A$-équivalente à $E\times_BB'$~; l'image par $v^*$ de la
classe de $E$ est la classe de $E\times_BB'$. On vérifie aussitôt que si
$w:L\to L'$ est un homomorphisme de $B$-bimodules, le diagramme
\begin{equation}
\label{0.18.3.6.3-fr}
\tag{18.3.6.3}
  \xymatrix@C=3pc@R=2.4pc{
    \operatorname{Exan}_A(B,L)\ar[r]^{v^*}\ar[d]_{w_*}
      & \operatorname{Exan}_A(B',L)\ar[d]^{w_*}\\
    \operatorname{Exan}_A(B,L')\ar[r]_{v^*}
      & \operatorname{Exan}_A(B',L')
  }
\end{equation}
est commutatif, $L$ et $L'$ étant considérés comme $B'$-bimodules au moyen
de $v$ dans la colonne de droite. Remplaçant $L$ et $L'$ respectivement par
$L\times L$ et $L$, et $w$ par l'addition $s$ dans $L$, on en conclut que
$v^*$ est bien un \emph{homomorphisme de groupes}.
\end{env}

\oldpage[0\textsubscript{IV-1}]{62}

\begin{env}[18.3.6.4]
\label{0.18.3.6.4-fr}
Enfin, l'application $(u,1_B,1_L)_*$ est l'application
\[
  u^*:\operatorname{Exan}_A(B,L)\longrightarrow
  \operatorname{Exan}_{A'}(B,L)
\]
qui s'obtient en faisant correspondre à une $A$-extension $E$ de $B$ par $L$
l'anneau $E$ considéré comme $A'$-\emph{anneau} au moyen de $u$ (18.1.1), qui
est évidemment une $A'$-extension de $B$ par $L$, $B$ étant aussi considéré
comme $A'$-anneau au moyen de $u$~; il est clair qu'une $A$-équivalence est
aussi une $A'$-équivalence, d'où l'application (18.3.6.4), qui, pour tout
homomorphisme $w:L\to L'$ de $B$-bimodules, rend encore commutatif le diagramme
\begin{equation}
\label{0.18.3.6.5-fr}
\tag{18.3.6.5}
  \xymatrix@C=3pc@R=2.4pc{
    \operatorname{Exan}_A(B,L)\ar[r]^{u^*}\ar[d]_{w_*}
      & \operatorname{Exan}_{A'}(B,L)\ar[d]^{w_*}\\
    \operatorname{Exan}_A(B,L')\ar[r]_{u^*}
      & \operatorname{Exan}_{A'}(B,L')
  }
\end{equation}
d'où l'on conclut comme ci-dessus que $u^*$ est un homomorphisme de groupes.

Cela étant, on pose
$(u,v,w)_*=(1_{A'},1_{B'},w)_*\circ(1_{A'},v,1_L)_*
\circ(u,1_B,1_L)_*$ et l'on vérifie facilement, en raison de la commutativité
des diagrammes (18.3.6.3) et (18.3.6.5), que l'on a
\[
  (u\circ u',v\circ v',w'\circ w)_*
  =(u',v',w')_*\circ(u,v,w)_*,
\]
ce qui achève de prouver que (18.3.4.1) est un foncteur.

L'existence de l'homomorphisme de groupes (18.3.6.1) montre en particulier
que si $E$ est une $A$-extension \emph{triviale} de $B$ par $L$, l'extension
$E\oplus_LL'$ de $B$ par $L'$ définie dans (18.2.8) est aussi triviale, ce que
l'on vérifie d'ailleurs sans peine de façon directe.

D'autre part, pour tout élément $z$ du \emph{centre} $Z$ de $B$, l'homothétie
$h_z:y\mapsto yz=zy$ est un endomorphisme du $B$-bimodule $L$, donc $(h_z)_*$
est un endomorphisme du groupe commutatif $\operatorname{Exan}_A(B,L)$, et par
fonctorialité, ces endomorphismes définissent sur
$\operatorname{Exan}_A(B,L)$ une structure canonique de $Z$-\emph{module}.
\end{env}

\begin{env}[18.3.7]
\label{0.18.3.7-fr}
Soient $A$, $A'$ deux anneaux, $u:A'\to A$ un homomorphisme, $B$ un
$A$-anneau et $L$ un $B$-bimodule. Le \emph{noyau} de l'homomorphisme de
groupes
\[
  u^*:\operatorname{Exan}_A(B,L)\longrightarrow
  \operatorname{Exan}_{A'}(B,L)
\]
est formé par définition des classes de $A$-extensions de $B$ par $L$ qui
sont $A'$-\emph{triviales} quand on les considère comme $A'$-extensions au
moyen de $u$. On désigne ce noyau par la notation
$\operatorname{Exan}_{A/A'}(B,L)$ quand cela ne prête pas à confusion.

Si $\Lambda$ est un anneau, et si on a un diagramme commutatif de
$\Lambda$-\emph{homomorphismes} de $\Lambda$-anneaux
\begin{equation}
\label{0.18.3.7.1-fr}
\tag{18.3.7.1}
  \xymatrix@C=2pc@R=1.5pc{
    B'\ar[r] & B\\
    A'\ar[u]\ar[r] & A\ar[u]
  }
\end{equation}

\oldpage[0\textsubscript{IV-1}]{63}

on en déduit canoniquement des homomorphismes
\begin{equation}
\label{0.18.3.7.2-fr}
\tag{18.3.7.2}
  \operatorname{Exan}_{A/\Lambda}(B,L)\longrightarrow
  \operatorname{Exan}_{A'/\Lambda}(B,L)\longrightarrow
  \operatorname{Exan}_{A'/\Lambda}(B',L)
\end{equation}
qui proviennent de la commutativité du diagramme
\[
  \xymatrix@C=2.8pc@R=2.3pc{
    \operatorname{Exan}_A(B,L)\ar[r]\ar[d]
      & \operatorname{Exan}_{A'}(B,L)\ar[r]\ar[d]
      & \operatorname{Exan}_{A'}(B',L)\ar[d]\\
    \operatorname{Exan}_{\Lambda}(B,L)\ar[r]
      & \operatorname{Exan}_{\Lambda}(B,L)\ar[r]
      & \operatorname{Exan}_{\Lambda}(B',L)
  }
\]
où les flèches sont déduites de celles de (18.3.7.1) par fonctorialité.
\end{env}

\begin{proposition}[18.3.8]
\label{0.18.3.8-fr}
Soient $B$ un anneau, $\mathfrak J$ un idéal bilatère de $B$,
$C=B/\mathfrak J$ l'anneau quotient~; $\mathfrak J/\mathfrak J^2$ est alors
canoniquement muni d'une structure de $C$-bimodule. Pour tout $C$-bimodule
$L$, soit $\operatorname{Hom}_C(\mathfrak J/\mathfrak J^2,L)$ le groupe
additif des homomorphismes de $C$-bimodules de $\mathfrak J/\mathfrak J^2$
dans $L$. On définit alors un isomorphisme canonique de groupes commutatifs
\begin{equation}
\label{0.18.3.8.1-fr}
\tag{18.3.8.1}
  \eta_L:\operatorname{Hom}_C(\mathfrak J/\mathfrak J^2,L)
  \xrightarrow{\sim}\operatorname{Exan}_B(C,L)
\end{equation}
en faisant correspondre à tout $C$-homomorphisme
$w:\mathfrak J/\mathfrak J^2\to L$ (qui est \emph{a fortiori} un
$B$-homomorphisme) la classe de l'extension
$(B/\mathfrak J^2)\oplus_{(\mathfrak J/\mathfrak J^2)}L$ déduite de
l'extension $B/\mathfrak J^2$ de $C$ par $\mathfrak J/\mathfrak J^2$ au moyen
de $w$~; l'isomorphisme réciproque fait correspondre à la classe d'une
$B$-extension $E$ de $C$ par $L$ l'homomorphisme $w$ tel que le composé
$\mathfrak J\to\mathfrak J/\mathfrak J^2\xrightarrow{w}L$ soit la restriction
à $\mathfrak J$ de l'homomorphisme structural $B\to E$.
\end{proposition}

\begin{proof}
Soit $F=B/\mathfrak J^2$, qui est une $B$-extension de $C$ par
$\mathfrak J/\mathfrak J^2$. Pour toute $B$-extension
\[
  0\longrightarrow L\xrightarrow{j}E\xrightarrow{p}C\longrightarrow0
\]
de $C$ par $L$, l'homomorphisme structural $f:B\to E$ est tel que le composé
$B\xrightarrow{f}E\xrightarrow{p}C$ soit l'homomorphisme canonique
$B\to B/\mathfrak J=C$. Comme l'image de $\mathfrak J$ par $p\circ f$ est
nulle, $f(\mathfrak J)$ est contenu dans le noyau de $p$, c'est-à-dire
$j(L)$, et comme $j(L)$ est de carré nul, on a $f(\mathfrak J^2)=0$~; donc
$f$ se factorise en
$B\to B/\mathfrak J^2=F\xrightarrow{u}E$, et si $w$ est la restriction de
$u$ à $\mathfrak J/\mathfrak J^2$, on a un diagramme commutatif
\[
  \xymatrix@C=2pc@R=1.5pc{
    0\ar[r] & \mathfrak J/\mathfrak J^2\ar[r]\ar[d]^w
      & F\ar[r]^g\ar[d]^u
      & C\ar[r]\ar[d]^{1_C}\ar@{}[d]|{\wr} & 0\\
    0\ar[r] & L\ar[r]_j & E\ar[r]_p & C\ar[r] & 0
  }
\]

autrement dit $(u,1_C,w)$ est un morphisme d'extensions (18.2.4). On en
déduit un morphisme d'extensions
\[
  \xymatrix@C=2pc@R=1.5pc{
    0\ar[r] & L\ar[r]\ar[d]^{1_L}\ar@{}[d]|{\wr}
      & E'\ar[r]\ar[d]^{u'}
      & C\ar[r]\ar[d]^{1_C}\ar@{}[d]|{\wr} & 0\\
    0\ar[r] & L\ar[r]_j & E\ar[r]_p & C\ar[r] & 0
  }
\]

\oldpage[0\textsubscript{IV-1}]{64}

où $E'=F\oplus_{(\mathfrak J/\mathfrak J^2)}L$ (18.2.8)~; il s'agit de
prouver que $u'$ est une $B$-\emph{équivalence}, ce qui établira que $\eta_L$
est bijectif. En vertu de la construction faite dans (18.2.8), tout revient à
prouver (compte tenu de (18.1.7)) que l'image réciproque $G$ de l'extension
$E$ de $C$ par l'homomorphisme canonique $g:F\to C$ est $F$-triviale, ce qui
est évident puisque $g$ se factorise en
$F\xrightarrow{u}E\xrightarrow{p}C$ (18.1.6).

Il reste à voir que $\eta_L$ est un homomorphisme de \emph{groupes}~; or, pour
tout $B$-homomorphisme $h:L\to L'$, il est immédiat que le diagramme
\[
  \xymatrix@C=3pc@R=2.4pc{
    \operatorname{Hom}_C(\mathfrak J/\mathfrak J^2,L)
      \ar[r]^{\eta_L}\ar[d]_{\operatorname{Hom}(1,h)}
      & \operatorname{Exan}_B(C,L)\ar[d]^{h_*}\\
    \operatorname{Hom}_C(\mathfrak J/\mathfrak J^2,L')
      \ar[r]_{\eta_{L'}}^{\sim}
      & \operatorname{Exan}_B(C,L')
  }
\]
est commutatif. Il suffit d'appliquer cette remarque à l'homomorphisme
$L\times L\to L$ définissant l'addition pour conclure.
\end{proof}

\subsection{Extensions d'algèbres}
\label{subsection:0.18.4-fr}

\begin{env}[18.4.1]
\label{0.18.4.1-fr}
Soit $A$ un anneau \emph{commutatif}. La catégorie des $A$-\emph{algèbres} est
alors une sous-catégorie pleine de celle des $A$-anneaux, caractérisée par le
fait que les homomorphismes structuraux $\rho:A\to B$ sont \emph{centraux}.

Si $B$ est une $A$-algèbre, $L$ un $B$-bimodule, il est immédiat que la
$A$-extension triviale type $D_B(L)$ est une $A$-algèbre. De même, dans la
construction de (18.2.5) (resp. de (18.2.8)), si $B$, $B'$ et $E'$ (resp.
$B$ et $E$) sont des $A$-algèbres, il en est de même de
$E'\times_{B'}B$ (resp. de $E\oplus_LL'$). Enfin, il est clair que si $B$ est
une $A$-algèbre et $E$ une $A$-extension de $B$ par un $B$-bimodule $L$ qui
est une $A$-algèbre, toute $A$-extension de $B$ par $L$, $A$-\emph{équivalente}
à $E$, est aussi une $A$-algèbre. On déduit alors aussitôt de la définition
(18.3.4) que les classes de $A$-extensions équivalentes de $B$ par $L$ qui
sont des $A$-\emph{algèbres} forment un \emph{sous-groupe}, noté
$\operatorname{Exal}_A(B,L)$, de $\operatorname{Exan}_A(B,L)$. Soit
$\mathbf K'$ la sous-catégorie pleine de la catégorie $\mathbf K$ définie dans
(18.3.5), dont les objets $(A,B,L)$ sont tels que $A$ soit commutatif et $B$
une $A$-algèbre. Alors ce qui précède montre que
\[
  (A,B,L)\longmapsto\operatorname{Exal}_A(B,L)
\]
est un \emph{foncteur covariant} de $\mathbf K'$ dans $\mathbf{Ab}$. Les
résultats de (18.3.7) sont inchangés lorsqu'on remplace partout
$\operatorname{Exan}$ par $\operatorname{Exal}$.
\end{env}

\begin{env}[18.4.2]
\label{0.18.4.2-fr}
Supposons toujours l'anneau $A$ commutatif. Les remarques de (18.4.1) restent
valables lorsque l'on remplace «~algèbre~» par «~algèbre commutative~» et
«~bimodule~» par «~module~». Si $B$ est une $A$-algèbre commutative et $L$ un
$B$-module,

\oldpage[0\textsubscript{IV-1}]{65}

les classes de $A$-extensions équivalentes de $B$ par $L$ qui sont des
$A$-\emph{algèbres commutatives} (ou, ce qui revient au même, des
$A$-\emph{anneaux commutatifs}) forment un \emph{sous-groupe} de
$\operatorname{Exal}_A(B,L)$, noté $\operatorname{Exalcom}_A(B,L)$. Si
$\mathbf K''$ est la sous-catégorie pleine de $\mathbf K'$ formée des triplets
$(A,B,L)$ où $A$ et $B$ sont commutatifs et $L$ un $B$-module,
\[
  (A,B,L)\longmapsto\operatorname{Exalcom}_A(B,L)
\]
est encore un \emph{foncteur covariant} de $\mathbf K''$ dans $\mathbf{Ab}$. On
peut aussi dans (18.3.7) remplacer partout $\operatorname{Exan}$ par
$\operatorname{Exalcom}$. Enfin, si dans (18.3.8), on suppose que $B$ soit un
anneau \emph{commutatif} et que $L$ soit un $C$-\emph{module}, le même
raisonnement donne un isomorphisme canonique
\begin{equation}
\label{0.18.4.2.1-fr}
\tag{18.4.2.1}
  \operatorname{Hom}_C(\mathfrak J/\mathfrak J^2,L)
  \xrightarrow{\sim}\operatorname{Exalcom}_B(C,L)
\end{equation}
où le premier membre est le groupe des homomorphismes de $C$-\emph{module}.
\end{env}

\begin{env}[18.4.3]
\label{0.18.4.3-fr}
Soit $A$ un anneau commutatif. Un cas important d'extensions de
$A$-algèbres est formé des $A$-algèbres $E$, extensions d'une $A$-algèbre $B$
par un $B$-bimodule $L$ telles que $E$, en tant que $A$-\emph{module}, soit une
extension \emph{triviale} du $A$-module $B$ par le $A$-module $L$~; autrement
dit, la suite exacte de $A$-\emph{modules}
$0\to L\to E\to B\to0$ est \emph{scindée}. On dit alors que $E$ est une
\emph{extension de Hochschild} de $B$ par $L$. Ce sera toujours le cas lorsque
$B$ est un $A$-module \emph{projectif}, et en particulier lorsque $A$ est un
\emph{corps} commutatif. En tant que $A$-module, on peut identifier $E$ à
$B\times L$, la multiplication dans $E$ étant donnée par
$(x,s)(y,t)=(xy,xt+sy+f(x,y))$ avec $f(x,y)\in L$. Si l'on écrit que cette
multiplication définit sur $B\times L$ une structure de $A$-algèbre, on trouve
(M, XIV, 2) que $f$ doit être une application $A$-\emph{bilinéaire} de
$B\times B$ dans $L$, telle que
\begin{equation}
\label{0.18.4.3.1-fr}
\tag{18.4.3.1}
  f(xy,z)+f(x,y)z=xf(y,z)+f(x,yz),
\end{equation}
autrement dit, $f$ est un $2$-\emph{cocycle} sur $B$ à valeurs dans $L$, au
sens de Hochschild~; pour que l'extension $E$ soit $A$-triviale, il faut et il
suffit que l'on ait
\begin{equation}
\label{0.18.4.3.2-fr}
\tag{18.4.3.2}
  f(x,y)=xg(y)-g(xy)+g(x)y,
\end{equation}
où $g$ est une application $A$-\emph{linéaire} de $B$ dans $L$, autrement dit
$f$ doit être un $2$-\emph{cobord} au sens de Hochschild. On en déduit aussitôt
que les classes d'extensions de Hochschild de $B$ par $L$ forment un
sous-groupe de $\operatorname{Exal}_A(B,L)$, isomorphe au \emph{groupe de
cohomologie} de Hochschild $\mathrm H_A^2(B,L)$.

Si $B$ est une $A$-algèbre \emph{commutative}, et $L$ un $B$-module, les
extensions de Hochschild \emph{commutatives} de $B$ par $L$ correspondent aux
$2$-cocycles \emph{symétriques}, c'est-à-dire tels que $f(y,x)=f(x,y)$. Les
classes d'extensions de Hochschild commutatives de $B$ par $L$ forment donc un
sous-groupe de $\operatorname{Exalcom}_A(B,L)$, isomorphe au sous-groupe de
$\mathrm H_A^2(B,L)$ image du groupe des $2$-cocycles symétriques, que nous
noterons $\mathrm H_A^2(B,L)^s$.
\end{env}

\begin{env}[18.4.4]
\label{0.18.4.4-fr}
Bornons-nous au cas où $A$ et $B$ sont \emph{commutatifs} et $L$ un
$B$-module, et rappelons dans ce cas la définition équivalente des groupes de
Hochschild $\mathrm H_A^i(B,L)$ (M, IX, 6). On considère un complexe
$\mathbf P_\bullet=(\mathbf P_n)_{n\geqslant0}$ de $B$-modules, où
$\mathbf P_n=B^{\otimes(n+1)}=B\otimes_A B\otimes_A\cdots\otimes_A B$
($n+1$ fois), la structure de $B$-module étant définie

\oldpage[0\textsubscript{IV-1}]{66}

par
\[
  x(y_1\otimes y_2\otimes\cdots\otimes y_{n+1})
  =(xy_1)\otimes y_2\otimes\cdots\otimes y_{n+1}\,;
\]
la dérivation, de degré $-1$, $d_n:\mathbf P_n\to\mathbf P_{n-1}$, est
définie par
\begin{align*}
  d_n(x_1\otimes x_2\otimes\cdots\otimes x_{n+1})
    &=(x_1x_2)\otimes x_3\otimes\cdots\otimes x_{n+1}
      -x_1\otimes(x_2x_3)\otimes\cdots\otimes x_{n+1}+\cdots\\
    &\quad+(-1)^{n-1}x_1\otimes x_2\otimes\cdots\otimes x_{n-1}
      \otimes(x_nx_{n+1})
      +(-1)^n(x_{n+1}x_1)\otimes x_2\otimes\cdots\otimes x_n,
\end{align*}
qui est bien $B$-linéaire puisque $B$ est commutatif.

Un $2$-cocycle de ce complexe à valeurs dans $L$ est une application
$B$-linéaire $h$ de $B\otimes B\otimes B$ dans $L$ telle que
$h(d_3(x\otimes y\otimes z\otimes t))=0$~; mais comme
$h(x\otimes y\otimes z)=xh(1\otimes y\otimes z)$, le cocycle $h$ est déterminé
par l'application $A$-bilinéaire
$(y,z)\longmapsto f(y,z)=h(1\otimes y\otimes z)$ de $B\times B$ dans $L$, et,
en écrivant la condition précédente pour $h$, on retrouve pour $f$ la
condition (18.4.3.1). De même, un $2$-cobord sera une application de la forme
$x\otimes y\otimes z\longmapsto h'(d_2(x\otimes y\otimes z))$, où
$h':B\otimes_A B\to L$ est linéaire, et ici encore $h'$ est déterminé par
l'application $B$-linéaire $g:y\longmapsto h'(1\otimes y)$ de $B$ dans $L$~;
on obtient alors
$h'(d_2(1\otimes x\otimes y))=xg(y)-g(xy)+yg(x)$, ce qui redonne (18.4.3.2).
On procède de même pour tout $i$, et l'on voit ainsi que l'on a
\begin{equation}
\label{0.18.4.4.1-fr}
\tag{18.4.4.1}
  \mathrm H_A^i(B,L)=\mathrm H^i(\operatorname{Hom}_B(\mathbf P_\bullet,L)).
\end{equation}
\end{env}

\begin{env}[18.4.5]
\label{0.18.4.5-fr}
Sous les conditions de (18.4.4), on peut interpréter de la même façon le
groupe $\mathrm H_A^2(B,L)^s$ (18.4.3). Pour cela, modifions le complexe
$\mathbf P_\bullet$ en degré $3$, en considérant un nouveau complexe
\[
  \mathbf P'_\bullet:\mathbf P'_3\xrightarrow{d'_3}\mathbf P_2
  \xrightarrow{d_2}\mathbf P_1\,;
\]
nous prendrons
$\mathbf P'_3=\mathbf P_3\oplus(B\otimes_A B\otimes_A B)$, $d'_3$ coïncidant
avec $d_3$ sur $\mathbf P_3$, et étant donné dans $B\otimes_A B\otimes_A B$
par
\[
  d'_3(x\otimes y\otimes z)=x\otimes y\otimes z-x\otimes z\otimes y.
\]

La relation $d_2d'_3=0$ résulte de la commutativité de $B$. Avec les notations
introduites ci-dessus, un $2$-cocycle de $\mathbf P'_\bullet$ correspond
maintenant à une application $A$-bilinéaire $f:B\times B\to L$ qui est
\emph{symétrique} et vérifie (18.4.3.1)~; par suite, on a
\[
  \mathrm H_A^2(B,L)^s
  =\mathrm H^2(\operatorname{Hom}_B(\mathbf P'_\bullet,L)).
\]
\end{env}

\begin{env}[18.4.6]
\label{0.18.4.6-fr}
Dans le cas particulier où l'on considère un \emph{corps} commutatif $k$, une
\emph{extension} $K$ de $k$, on a $\operatorname{Ext}_K^1(M,L)=0$ quels que
soient les $K$-espaces vectoriels $L$, $M$, et par suite (M, VI, 3.3 $a$)), on
a un isomorphisme canonique
\begin{equation}
\label{0.18.4.6.1-fr}
\tag{18.4.6.1}
  \mathrm H_k^i(K,L)\xrightarrow{\sim}
  \operatorname{Hom}_K(\mathrm H_i(\mathbf P_\bullet),L)
\end{equation}
et de même
\begin{equation}
\label{0.18.4.6.2-fr}
\tag{18.4.6.2}
  \mathrm H_k^2(K,L)^s\xrightarrow{\sim}
  \operatorname{Hom}_K(\mathrm H_2(\mathbf P'_\bullet),L).
\end{equation}
\end{env}

\subsection{Cas des anneaux topologiques}
\label{subsection:0.18.5-fr}

\begin{env}[18.5.1]
\label{0.18.5.1-fr}
Soient $A$, $B$ deux anneaux \emph{topologiques} dont la topologie est
\emph{linéaire}, $\rho:A\to B$ un homomorphisme continu, $L$ un $B$-bimodule
topologique, et supposons qu'il existe un idéal bilatère \emph{ouvert}
$\mathfrak R_0$ de $B$ tel que
$\mathfrak R_0.L=L.\mathfrak R_0=0$, de sorte que $L$ peut être

\oldpage[0\textsubscript{IV-1}]{67}

considéré comme un $(B/\mathfrak R)$-bimodule pour tout idéal ouvert
$\mathfrak R\subset\mathfrak R_0$. Soient alors
$\mathfrak R\subset\mathfrak R_0$ un idéal bilatère ouvert de $B$,
$\mathfrak J$ un idéal bilatère ouvert de $A$ tel que
$\rho(\mathfrak J)\subset\mathfrak R$, de sorte que $B/\mathfrak R$ peut être
considéré comme un $(A/\mathfrak J)$-anneau discret. La remarque précédente
prouve que le groupe
$\operatorname{Exan}_{A/\mathfrak J}(B/\mathfrak R,L)$ est défini~; en outre,
si $\mathfrak R'\subset\mathfrak R$, $\mathfrak J'\subset\mathfrak J$ sont
deux idéaux bilatères ouverts de $A$ et $B$ respectivement tels que
$\rho(\mathfrak J')\subset\mathfrak R'$, on a par (18.3.5.1) un homomorphisme
canonique
\begin{equation}
\label{0.18.5.1.1-fr}
\tag{18.5.1.1}
  \operatorname{Exan}_{A/\mathfrak J}(B/\mathfrak R,L)
  \longrightarrow
  \operatorname{Exan}_{A/\mathfrak J'}(B/\mathfrak R',L).
\end{equation}

L'ensemble des couples d'idéaux ouverts $(\mathfrak J,\mathfrak R)$ tels que
$\rho(\mathfrak J)\subset\mathfrak R\subset\mathfrak R_0$ est évidemment
ordonné filtrant à droite pour la relation «~$\mathfrak J\supset\mathfrak J'$
et $\mathfrak R\supset\mathfrak R'$~» et les applications (18.5.1.1)
définissent un \emph{système inductif} de groupes additifs ayant cet ensemble
pour ensemble d'indices. On pose, par abus de notation (car il ne s'agit plus
d'un groupe en correspondance biunivoque naturelle avec un ensemble
d'extensions)
\begin{equation}
\label{0.18.5.1.2-fr}
\tag{18.5.1.2}
  \operatorname{Exantop}_A(B,L)
  =\varinjlim\operatorname{Exan}_{A/\mathfrak J}(B/\mathfrak R,L).
\end{equation}

Dire que le second membre de (18.5.1.2) est \emph{nul} signifie donc que, pour
tout couple d'idéaux ouverts $\mathfrak J\subset A$, $\mathfrak R\subset B$
tels que $\rho(\mathfrak J)\subset\mathfrak R\subset\mathfrak R_0$ et toute
$(A/\mathfrak J)$-extension $E$ de $B/\mathfrak R$ par $L$, il existe deux
idéaux ouverts $\mathfrak J'\subset\mathfrak J$,
$\mathfrak R'\subset\mathfrak R$ tels que
$\rho(\mathfrak J')\subset\mathfrak R'$ et que l'image réciproque par
l'homomorphisme $B/\mathfrak R'\to B/\mathfrak R$ de $E$ soit \emph{triviale}.

Nous laissons au lecteur la définition analogue des limites inductives
$\operatorname{Exaltop}_A(B,L)$, $\operatorname{Exalcotop}_A(B,L)$ à partir de
$\operatorname{Exal}_{A/\mathfrak J}(B/\mathfrak R,L)$ et de
$\operatorname{Exalcom}_{A/\mathfrak J}(B/\mathfrak R,L)$ pour le cas où $A$
est commutatif et $B$ une $A$-algèbre (resp. une $A$-algèbre commutative)
topologique.
\end{env}

\begin{env}[18.5.2]
\label{0.18.5.2-fr}
Si l'on a un diagramme commutatif d'homomorphismes continus d'anneaux
\[
  \xymatrix@C=3pc@R=2pc{
    B'\ar[r] & B\\
    A'\ar[u]\ar[r] & A\ar[u]
  }
\]
on en déduit canoniquement deux homomorphismes de groupes additifs
\[
  \operatorname{Exantop}_A(B,L)
  \longrightarrow\operatorname{Exantop}_{A'}(B,L)
  \longrightarrow\operatorname{Exantop}_{A'}(B',L)
\]
par passage à la limite inductive à partir de (18.3.5.1).

En vertu de l'exactitude du foncteur $\varinjlim$ dans la catégorie des groupes
commutatifs, le noyau de l'homomorphisme
\[
  \operatorname{Exantop}_A(B,L)
  \longrightarrow\operatorname{Exantop}_{A'}(B,L)
\]
est la limite inductive des noyaux des homomorphismes
\[
  \operatorname{Exan}_{A/\mathfrak J}(B/\mathfrak R,L)
  \longrightarrow
  \operatorname{Exan}_{A'/\mathfrak J'}(B/\mathfrak R,L)
\]
où on a pris pour $\mathfrak J'$ l'image réciproque de $\mathfrak J$~; on note
ce noyau $\operatorname{Exantop}_{A/A'}(B,L)$. On définit de même
$\operatorname{Exaltop}_{A/A'}(B,L)$ et
$\operatorname{Exalcotop}_{A/A'}(B,L)$. Enfin, si on a un homomorphisme

\oldpage[0\textsubscript{IV-1}]{68}

continu de $B$-bimodules $L\to L'$, on en déduit canoniquement un
homomorphisme de groupes additifs
\[
  \operatorname{Exantop}_A(B,L)\longrightarrow
  \operatorname{Exantop}_A(B,L')
\]
par passage à la limite inductive à partir de (18.3.6.1).
\end{env}

\begin{env}[18.5.3]
\label{0.18.5.3-fr}
Étant donnés un anneau topologique $C$ et deux $C$-bimodules topologiques
$M$, $N$,\linebreak on note $\operatorname{Hom.\ cont}_C(M,N)$ le groupe additif des
$C$-homomorphismes \emph{continus} de $M$ dans $N$.
\end{env}

\begin{lemma}[18.5.3.1]
\label{0.18.5.3.1-fr}
Soient $C$ un anneau topologique, $E$, $L$ deux $C$-bimodules topologiques~;
on suppose que les topologies sont linéaires, que $L$ est discret et annulé
par un idéal bilatère ouvert de $C$. Alors on a un isomorphisme canonique
\begin{equation}
\label{0.18.5.3.2-fr}
\tag{18.5.3.2}
  \varinjlim\operatorname{Hom}_{C/\mathfrak R}(E/V,L)
  \xrightarrow{\sim}\operatorname{Hom.\ cont}_C(E,L)
\end{equation}
où dans le premier membre la limite inductive est prise suivant l'ensemble
ordonné filtrant des couples $(\mathfrak R,V)$ tels que $\mathfrak R$ soit un
idéal bilatère ouvert de $C$, $V$ un sous-$C$-bimodule ouvert de $E$, tels que
$\mathfrak R.L=L.\mathfrak R=0$, $\mathfrak R.E\subset V$,
$E.\mathfrak R\subset V$.
\end{lemma}

\begin{proof}
Comme $C/\mathfrak R$ et $E/V$ sont discrets, on a des homomorphismes
canoniques
\[
  w_{\mathfrak R,V}:\operatorname{Hom}_{C/\mathfrak R}(E/V,L)
  \longrightarrow\operatorname{Hom.\ cont}_C(E,L)
\]
formant un système inductif, d'où un homomorphisme (18.5.3.2) par passage à la
limite inductive. Comme l'homomorphisme $E/V'\to E/V$ est surjectif pour
$V'\subset V$, il résulte aussitôt de la définition que l'homomorphisme
\[
  \operatorname{Hom}_{C/\mathfrak R}(E/V,L)
  \longrightarrow\operatorname{Hom}_{C/\mathfrak R}(E/V',L)
\]
(avec $\mathfrak R.L=L.\mathfrak R=0$, $\mathfrak R.E\subset V'$,
$E.\mathfrak R\subset V'$) est injectif, et il en est évidemment de même de
l'homomorphisme
\[
  \operatorname{Hom}_{C/\mathfrak R}(E/V,L)
  \longrightarrow\operatorname{Hom}_{C/\mathfrak R'}(E/V,L)
\]
pour $\mathfrak R'\subset\mathfrak R$~; on en conclut que l'homomorphisme
(18.5.3.2) est injectif. D'autre part, si $u$ est un $C$-homomorphisme continu
de $E$ dans $L$, son noyau est un sous-bimodule ouvert $V_0$ de $E$, et si
$\mathfrak R_0$ est un idéal bilatère ouvert de $C$ tel que
$\mathfrak R_0.L=L.\mathfrak R_0=0$ et $\mathfrak R_0.E\subset V_0$,
$E.\mathfrak R_0\subset V_0$, il est clair que $u$ est l'image canonique d'un
$(C/\mathfrak R_0)$-homomorphisme de $E/V_0$ dans $L$, donc (18.5.3.2) est
surjectif.
\end{proof}

Cela étant, la proposition (18.3.8) se généralise comme suit aux anneaux
topologiques~:

\begin{proposition}[18.5.4]
\label{0.18.5.4-fr}
Soient $B$ un anneau topologique linéairement topologisé, $\mathfrak J$ un
idéal bilatère de $B$, $C=B/\mathfrak J$ l'anneau topologique quotient~;
$\mathfrak J/\mathfrak J^2$ (où $\mathfrak J$ est muni de la topologie induite
par celle de $B$ et $\mathfrak J/\mathfrak J^2$ de la topologie quotient de
celle de $\mathfrak J$) est alors canoniquement muni d'une structure de
$C$-bimodule topologique. Pour tout $C$-bimodule discret $L$ annulé par un
idéal ouvert de $C$, il existe alors un isomorphisme canonique
\begin{equation}
\label{0.18.5.4.1-fr}
\tag{18.5.4.1}
  \operatorname{Hom.\ cont}_C(\mathfrak J/\mathfrak J^2,L)
  \xrightarrow{\sim}\operatorname{Exantop}_B(C,L).
\end{equation}
\end{proposition}

\begin{proof}
En effet, pour tout idéal ouvert $\mathfrak R$ de $B$ tel que
$(\mathfrak J+\mathfrak R)/\mathfrak J$ annule $L$, on a, d'après (18.3.9),
un isomorphisme canonique
\[
  \operatorname{Hom}_{B/(\mathfrak J+\mathfrak R)}
    ((\mathfrak J+\mathfrak R)/(\mathfrak J^2+\mathfrak R),L)
  \xrightarrow{\sim}
  \operatorname{Exan}_{B/\mathfrak R}(B/(\mathfrak J+\mathfrak R),L)
\]
et il suffit de passer à la limite inductive, en tenant compte de (18.5.3.1).
\end{proof}

% FIN SOURCE ega4/ega0-18-fr.tex

%
% DEBUT SOURCE ega4/ega0-19-fr.tex
\section{Algèbres formellement lisses et anneaux de Cohen}
\label{section:0.19-fr}
\label{0.19-fr}

\setcounter{subsection}{-1}
\subsection{Introduction}
\label{subsection:0.19.0-fr}
\label{0.19.0-fr}

\oldpage[0\textsubscript{IV-1}]{69}

\begin{env}[19.0.1]
\label{0.19.0.1-fr}
Au chapitre IV, nous introduirons et étudierons entre autres une classe
importante de morphismes de préschémas, les morphismes \emph{lisses}\footnote{Aussi
appelés morphismes \emph{simples} dans certains travaux récents (en s'inspirant
de la terminologie classique de «~point simple~»)~; cette terminologie prête
cependant à confusion, notamment en théorie des groupes algébriques.}. Une de
leurs propriétés fondamentales (et qui, jointe à une condition de finitude, les
caractérise) est une propriété de \emph{relèvement des morphismes}~: si
$f:X\to Y$ est un morphisme lisse, $g:Y'\to Y$ un morphisme, alors pour tout
morphisme $h:Y''\to Y'$ faisant de $Y''$ un préschéma «~peu différent~» de
$Y'$, \emph{tout $Y$-morphisme $Y''\to X$ se factorise} en
$Y''\xrightarrow{h}Y'\to X$. De façon précise, bornons-nous au cas où
$Y=\operatorname{Spec}(A)$, $X=\operatorname{Spec}(B)$ sont \emph{affines}~;
$B$ est alors appelée une $A$-algèbre \emph{formellement lisse} si, pour toute
$A$-algèbre $C$ et tout idéal \emph{nilpotent} $\mathfrak J$ de $C$, tout
$A$-homomorphisme $B\to C/\mathfrak J$ \emph{se relève} en
$B\to C\to C/\mathfrak J$. En d'autres termes, l'application
\[
  \operatorname{Hom}_A(B,C)\longrightarrow
  \operatorname{Hom}_A(B,C/\mathfrak J)
\]
est \emph{surjective}. Dans beaucoup d'applications, $X$ apparaîtra comme objet
représentant un \emph{foncteur contravariant représentable} $Y'\mapsto F(Y')$
de la catégorie des $Y$-préschémas dans celle des ensembles, de sorte
($\mathbf{0}_{\mathrm{III}}$, 8.1.8) que l'on aura
$F(Y')=\operatorname{Hom}_Y(Y',X)$. Dans le cas affine, si l'on pose
$F^0(C)=F(\operatorname{Spec}(C))$, la vérification du fait que, sous les
conditions ci-dessus, $F^0(C)\to F^0(C/\mathfrak J)$ est \emph{surjectif}
(qui peut se faire même \emph{sans savoir} que $F$ est représentable) établira
que $f$ est lisse, pourvu que la condition additionnelle de finitude soit
remplie.
\end{env}

\begin{env}[19.0.2]
\label{0.19.0.2-fr}
Pour les \emph{algèbres topologiques} sur un anneau topologique $A$, il y a une
notion analogue d'algèbre \emph{formellement lisse} que nous ne préciserons pas
ici (cf.~19.3.1). L'étude de ces notions est faite d'abord d'un point de vue
élémentaire au §~19, puis, à l'aide des propriétés des différentielles qui
seront développées aux §§~20 et 21, des théorèmes plus délicats seront prouvés
au §~22. Résumons ici les principaux résultats sur les algèbres formellement
lisses~:

\medskip\noindent\textbf{I.} --- Soient $A$, $B$ deux anneaux locaux
noethériens, $\varphi:A\to B$ un homomorphisme local, $k$ le corps résiduel de
$A$, et soit $B_0=B\otimes_A k$~; on munit $A$, $B$ et $B_0$ de leurs
topologies préadiques et $k$ de la topologie discrète. Alors, pour que $B$ soit
une $A$-algèbre formellement lisse, il faut et il suffit que $B$ soit un
$A$-module \emph{plat} et que $B_0$ soit une $k$-algèbre \emph{formellement
lisse} (19.7.1). Ce théorème ramène donc la lissité formelle, pour les anneaux
locaux noethériens, à la même question pour les anneaux locaux noethériens qui
sont des algèbres sur un \emph{corps}.

\medskip\noindent\textbf{II.} --- Soient $k$ un corps, $A$ un anneau local
noethérien qui est une $k$-algèbre. Pour que $A$ soit formellement lisse, il
faut et il suffit que $A$ soit \emph{géométriquement régulier} sur $k$,
c'est-à-dire que pour toute extension \emph{finie} $k'$ de $k$, l'anneau
semi-local $A\otimes_k k'$ soit

\oldpage[0\textsubscript{IV-1}]{70}

\noindent régulier (22.5.8)~; le \emph{complété} $\widehat A$ de $A$ est alors
un anneau isomorphe à un anneau de séries formelles
$K[[T_1,\ldots,T_n]]$ (19.6.5). En outre, la structure de $k$-algèbre de $A$,
lorsque $A$ est supposé être \emph{complet} et formellement lisse sur $k$, est
entièrement déterminée par le corps résiduel $K$ de $A$ et la \emph{dimension}
de $A$~; celle-ci peut d'ailleurs être arbitraire à condition de satisfaire à
l'inégalité
\[
  \dim(A)\geq \operatorname{rg}_k(\Gamma_{K/k}),
\]
où $\Gamma_{K/k}$ est le «~module d'imperfection~» de $K$ (22.2.6).

En particulier, pour qu'une \emph{extension} $K$ de $k$ soit formellement
lisse, il faut et il suffit que $K$ soit une extension \emph{séparable} de $k$
(19.6.1).

\medskip\noindent\textbf{III.} --- Soient $A$ un anneau local noethérien,
$\mathfrak J$ un idéal de $A$ distinct de $A$, $A_0=A/\mathfrak J$, $B_0$ un
anneau local noethérien \emph{complet}, $A_0\to B_0$ un homomorphisme local
faisant de $B_0$ une $A_0$-algèbre formellement lisse. Alors il existe un
anneau local noethérien \emph{complet} $B$, un homomorphisme local $A\to B$
faisant de $B$ un $A$-module plat, et un $A_0$-\emph{isomorphisme}
\[
  B\otimes_A A_0\xrightarrow{\sim}B_0
\]
(de sorte que, par I), $B$ est une $A$-algèbre \emph{formellement lisse})~; en
outre $B$ est déterminé par ces propriétés à isomorphisme près (19.7.2). Ce
théorème contient en particulier les \emph{théorèmes de Cohen} sur la structure
des anneaux locaux noethériens \emph{complets} (19.8), qui joueront un rôle
important dans les §§~6 et 7 du chapitre IV.

\medskip\noindent\textbf{IV.} --- L'intérêt de l'étude des anneaux locaux
noethériens formellement lisses sur un autre provient de la caractérisation
«~ponctuelle~» suivante des morphismes lisses~: si $X$ et $Y$ sont des
préschémas localement noethériens, $f:X\to Y$ un morphisme localement de type
fini, alors, pour que $f$ soit lisse, il faut et il suffit que pour tout
$x\in X$, l'anneau $\mathcal O_x$ soit formellement lisse sur
$\mathcal O_{f(x)}$. En particulier, si $Y=\operatorname{Spec}(k)$, où $k$ est
un corps parfait, dire que $f$ est lisse équivaut à dire que $X$ est un
préschéma \emph{régulier}.

\medskip\noindent\textbf{V.} --- Enfin, nous verrons aux §§~20, 21 et 22 que
la notion d'algèbre formellement lisse s'introduit de façon naturelle dans la
théorie des \emph{différentielles de Kähler}, les deux théories s'éclairant
mutuellement.
\end{env}

\begin{env}[19.0.3]
\label{0.19.0.3-fr}
Dans tout ce paragraphe et les suivants, les anneaux et modules topologiques
seront supposés \emph{linéairement topologisés} ($\mathbf{0}_{\mathrm I}$,
7.1.1)~; les anneaux topologiques considérés seront supposés \emph{commutatifs},
sauf mention expresse du contraire. Rappelons que si $A$ et $B$ sont deux
anneaux topologiques, $\rho:A\to B$ un homomorphisme d'anneaux définissant sur
$B$ une structure de $A$-algèbre, on dit que $B$ est une \emph{$A$-algèbre
topologique} si $\rho$ est \emph{continu} pour les topologies envisagées.

Pour abréger, dans un anneau topologique $A$ (resp. un $A$-module topologique
$M$), nous dirons «~\emph{système fondamental d'idéaux} (resp.
\emph{sous-modules}) \emph{ouverts}~» au lieu de «~système fondamental de
voisinages de 0 formé d'idéaux (resp. sous-modules)~».

Étant donnés un anneau topologique $A$ et un $A$-module $M$, les ensembles
$\mathfrak J M$, où $\mathfrak J$ parcourt un système fondamental d'idéaux
ouverts, forment un système fondamental de sous-modules ouverts d'une topologie
sur $M$ faisant de $M$ un $A$-module topologique, et que l'on dit \emph{déduite}
de la topologie de $A$.

Soit $M$ un $A$-module topologique dont la topologie est \emph{moins fine} que
la topologie déduite de celle de $A$~; alors, si $N$ est un sous-module
\emph{ouvert} de $M$, le $A$-module \emph{discret} $M/N$ est \emph{annulé par
un idéal ouvert} de $A$, car par hypothèse il existe un tel idéal $\mathfrak R$
tel que $\mathfrak R.M\subset N$.

\oldpage[0\textsubscript{IV-1}]{71}

Si $M$ et $N$ sont deux $A$-modules topologiques dont les topologies sont toutes
deux déduites de celle de $A$, alors \emph{tout} $A$-homomorphisme $u:M\to N$
est \emph{continu}, car pour tout voisinage $V$ de 0 dans $N$, il existe par
hypothèse un idéal ouvert $\mathfrak J$ de $A$ tel que $\mathfrak JN\subset V$,
et l'on a donc $u(\mathfrak JM)\subset\mathfrak JN\subset V$.

Lorsque $B$ est une $A$-algèbre topologique, la topologie sur $B$ \emph{déduite}
de celle de $A$ est \emph{plus fine} que la topologie donnée, car pour tout
idéal ouvert $\mathfrak R$ de $B$, il y a par hypothèse un idéal ouvert
$\mathfrak J$ de $A$ tel que $\mathfrak JB\subset\mathfrak R$.
\end{env}

\subsection{Épimorphismes et monomorphismes formels}
\label{subsection:0.19.1-fr}

\begin{proposition}[19.1.1]
\label{0.19.1.1-fr}
Soient $A$ un anneau topologique, $M$, $N$ deux $A$-modules topologiques,
$(W_\lambda)$ un système fondamental de sous-modules ouverts dans $N$,
$u:M\to N$ un $A$-homomorphisme continu, $\widehat M$, $\widehat N$ les
séparés complétés de $M$ et $N$, $\widehat u:\widehat M\to\widehat N$ le
prolongement continu de $u$ aux séparés complétés.
\begin{enumerate}
  \item[(i)] Les conditions suivantes sont équivalentes~:
  \begin{enumerate}
    \item[a)] $u(M)$ est \emph{dense} dans $N$.
    \item[a$'$)] $\widehat u(\widehat M)$ est \emph{dense} dans $\widehat N$.
    \item[a$''$)] Pour tout $\lambda$, l'homomorphisme composé
    $M\xrightarrow{u}N\to N/W_\lambda$ est \emph{surjectif}.
  \end{enumerate}
  \item[(ii)] Les conditions suivantes sont équivalentes~:
  \begin{enumerate}
    \item[b)] L'image réciproque par $u$ de la topologie de $N$ est égale à la
    topologie de $M$.
    \item[b$'$)] $\widehat u$ est un \emph{isomorphisme} du $\widehat A$-module
    topologique $\widehat M$ sur le sous-$\widehat A$-module topologique
    $\widehat u(\widehat M)$ de $\widehat N$ (qui est nécessairement
    \emph{fermé}).
    \item[b$''$)] Les $u^{-1}(W_\lambda)$ forment un système fondamental de
    voisinages de 0 dans $M$.
  \end{enumerate}
\end{enumerate}
Cela résulte immédiatement de la définition des séparés complétés, et du fait
que $N/W_\lambda$ est discret.
\end{proposition}

\begin{env}[19.1.2]
\label{0.19.1.2-fr}
Lorsque les conditions équivalentes de (i) (resp. (ii)) dans (19.1.1) sont
remplies, on dit que $u$ est un \emph{épimorphisme formel} (resp. un
\emph{monomorphisme formel}). On dit que $u$ est un \emph{bimorphisme formel}
s'il est à la fois un monomorphisme formel et un épimorphisme formel~; il
revient au même, en vertu de (19.1.1), de dire que $\widehat u$ est un
\emph{isomorphisme} du $\widehat A$-module topologique $\widehat M$ sur le
$\widehat A$-module topologique $\widehat N$.
\end{env}

\begin{proposition}[19.1.3]
\label{0.19.1.3-fr}
Soient $A$ un anneau topologique, $M$, $N$ deux $A$-modules topologiques,
$u:M\to N$ un $A$-homomorphisme continu. On suppose qu'il existe deux systèmes
fondamentaux de sous-modules ouverts, $(V_\lambda)$, $(W_\lambda)$ dans $M$ et
$N$ respectivement, ayant même ensemble d'indices et tels que
$u(V_\lambda)\subset W_\lambda$ pour tout $\lambda$~; soit
$u_\lambda:M/V_\lambda\to N/W_\lambda$ l'homomorphisme déduit de $u$ par
passage aux quotients. Alors~:
\begin{enumerate}
  \item[(i)] Pour que $u$ soit un épimorphisme formel, il faut et il suffit que,
  pour tout $\lambda$, $u_\lambda$ soit surjectif.
  \item[(ii)] Pour que $u$ soit un monomorphisme formel, il faut et il suffit
  que, pour tout $\lambda$, il existe un $\mu$ tel que
  $V_\mu\subset V_\lambda$ et que $\operatorname{Ker}(u_\mu)$ soit contenu
  dans le noyau $V_\lambda/V_\mu$ de l'application canonique
  $M/V_\mu\to M/V_\lambda$.
  \item[(iii)] Pour que $u$ soit un bimorphisme formel, il faut et il suffit
  que, pour tout $\lambda$, $u_\lambda$ soit surjectif et qu'il existe un
  $\mu$ tel que $V_\mu\subset V_\lambda$ et que l'homomorphisme canonique
  $M/V_\mu\to M/V_\lambda$

  \oldpage[0\textsubscript{IV-1}]{72}

  \noindent se factorise en
  $M/V_\mu\xrightarrow{u_\mu}N/W_\mu\to M/V_\lambda$ (où
  l'homomorphisme $N/W_\mu\to M/V_\lambda$ est nécessairement unique).
\end{enumerate}
\end{proposition}

\begin{proof}
L'assertion (i) est immédiate~; d'autre part,
$\operatorname{Ker}(u_\mu)=u^{-1}(W_\mu)/V_\mu$, et le noyau de l'application
canonique $M/V_\mu\to M/V_\lambda$ est $V_\lambda/V_\mu$~; la condition de
(ii) revient donc à dire que $u^{-1}(W_\mu)\subset V_\lambda$, et l'assertion
(ii) en découle aussitôt~; enfin, lorsque $u_\mu$ est surjectif, il revient au
même de dire que $\operatorname{Ker}(u_\mu)$ est contenu dans
$V_\lambda/V_\mu$, ou de dire que $M/V_\mu\to M/V_\lambda$ se factorise en
$v\circ u_\mu$, où $v$ est un homomorphisme
$N/W_\mu\to M/V_\lambda$, puisque $N/W_\mu$ s'identifie alors à
$(M/V_\mu)/\operatorname{Ker}(u_\mu)$.
\end{proof}

\begin{corollary}[19.1.4]
\label{0.19.1.4-fr}
Soient $A$ un anneau topologique, $M$, $N$ deux $A$-modules topologiques dont
les topologies sont déduites de celle de $A$, $u:M\to N$ un épimorphisme
formel. Soit $(\mathfrak J_\lambda)$ un système fondamental d'idéaux ouverts
dans $A$. Pour que $u$ soit un bimorphisme formel, il faut et il suffit que pour
tout $\lambda$, l'homomorphisme
$u_\lambda:M/\mathfrak J_\lambda M\to N/\mathfrak J_\lambda N$ déduit de $u$
par passage aux quotients soit bijectif.
\end{corollary}

\begin{proof}
On a en effet $u(\mathfrak J_\lambda M)\subset\mathfrak J_\lambda N$ et on
peut appliquer le critère (19.1.3, (iii))~; mais si l'on a une factorisation
\[
  M/\mathfrak J_\mu M\xrightarrow{u_\mu}
  N/\mathfrak J_\mu N\xrightarrow{v}M/\mathfrak J_\lambda M,
\]
l'image de $\mathfrak J_\lambda M/\mathfrak J_\mu M$ par $u_\mu$ est
$\mathfrak J_\lambda N/\mathfrak J_\mu N$, donc l'image de
$\mathfrak J_\lambda N/\mathfrak J_\mu N$ par $v$ est 0 et $v$ se factorise en
\[
  N/\mathfrak J_\mu N\longrightarrow N/\mathfrak J_\lambda N
  \xrightarrow{v'}M/\mathfrak J_\lambda M~;
\]
mais alors l'automorphisme identique de $M/\mathfrak J_\lambda M$ se factorise
en
\[
  M/\mathfrak J_\lambda M\xrightarrow{u_\lambda}
  N/\mathfrak J_\lambda N\xrightarrow{v'}M/\mathfrak J_\lambda M,
\]
ce qui montre que $u_\lambda$ est injectif, et comme on sait déjà qu'il est
surjectif, cela prouve le corollaire.
\end{proof}

\begin{proposition}[19.1.5]
\label{0.19.1.5-fr}
Soient $A$ un anneau topologique, $M$, $N$ deux $A$-modules topologiques,
$u:M\to N$ un $A$-homomorphisme continu. Les conditions suivantes sont
équivalentes~:
\begin{enumerate}
  \item[a)] Pour tout $A$-module topologique discret $E$ et tout
  $A$-homomorphisme continu $v:M\to E$, il existe un $A$-homomorphisme continu
  $w:N\to E$ tel que $v=w\circ u$.
  \item[b)] Pour tout sous-module ouvert $V'$ de $M$, il existe un sous-module
  ouvert $W''$ de $N$, un sous-module ouvert
  $V''\subset V'\cap u^{-1}(W'')$ et un homomorphisme
  $h:N/W''\to M/V'$ tels que l'homomorphisme canonique
  $M/V''\to M/V'$ se factorise en
  \[
    M/V''\xrightarrow{u''}N/W''\xrightarrow{h}M/V',
  \]
  où $u''$ est déduit de $u$ par passage aux quotients.
\end{enumerate}
\end{proposition}

\begin{proof}
Pour voir que a) entraîne b), il suffit d'appliquer a) à $E=M/V'$ et à
l'homomorphisme canonique $v:M\to M/V'$~; alors $W''=w^{-1}(0)$ est un
sous-module ouvert de $N$ et $V''=u^{-1}(W'')$ un sous-module ouvert de $M$
contenu dans $V'$~; ces sous-modules et l'homomorphisme
$h:N/W''\to M/V'$ déduit de $w$ par passage au quotient répondent à la
question. Inversement, pour montrer que b) entraîne a), on peut se limiter au
cas où $v$ est surjectif, en remplaçant $E$ par $v(M)$~; alors
$V'=\operatorname{Ker}(v)$ est un sous-module ouvert de $M$, donc $E$ est
isomorphe au $A$-module discret $M/V'$, et en appliquant b), on obtient

\oldpage[0\textsubscript{IV-1}]{73}

le $A$-homomorphisme continu $w$ cherché en prenant le composé
$N\to N/W''\xrightarrow{h}M/V'$, le diagramme
\[
  \xymatrix{
    M\ar[r]^u\ar[d] & N\ar[d]\\
    M/V''\ar[r]_{u''} & N/W''
  }
\]
étant commutatif.
\end{proof}

Lorsque les conditions équivalentes de (19.1.5) sont remplies, nous dirons que
$u$ est \emph{formellement inversible à gauche}~; comme la condition b) de
(19.1.5) implique que $u^{-1}(W'')\subset V'$, $u$ est alors un monomorphisme
formel. La terminologie est en outre justifiée par les corollaires suivants~:

\begin{corollary}[19.1.6]
\label{0.19.1.6-fr}
S'il existe un $\widehat A$-homomorphisme continu
$s:\widehat N\to\widehat M$ tel que
$s\circ\widehat u=1_{\widehat M}$, alors $u$ est formellement inversible à
gauche.
\end{corollary}

On notera que $\widehat u$ est alors un isomorphisme topologique de
$\widehat M$ sur un \emph{facteur direct topologique} de $\widehat N$. Pour
prouver (19.1.6), il suffit de noter qu'avec les notations de (19.1.5, a)),
$v:M\to E$ se prolonge en un $\widehat A$-homomorphisme continu
$\widehat v:\widehat M\to E$ puisque $E$ est discret~; soient
$j:M\to\widehat M$ et $j':N\to\widehat N$ les homomorphismes canoniques~;
alors, en posant $w=\widehat v\circ s\circ j'$, on a
$w\circ u=\widehat v\circ s\circ\widehat u\circ j=\widehat v\circ j=v$.

\begin{corollary}[19.1.7]
\label{0.19.1.7-fr}
Supposons que les topologies de $M$ et de $N$ soient déduites de celle de $A$,
et soit $(\mathfrak J_\lambda)$ un système fondamental d'idéaux ouverts dans
$A$. Pour que $u$ soit formellement inversible à gauche il faut et il suffit
que, pour tout $\lambda$, l'homomorphisme
$u_\lambda:M/\mathfrak J_\lambda M\to N/\mathfrak J_\lambda N$, déduit de
$u$ par passage aux quotients, soit inversible à gauche (autrement dit, soit
un isomorphisme de $M/\mathfrak J_\lambda M$ sur un facteur direct de
$N/\mathfrak J_\lambda N$).
\end{corollary}

En effet, la condition est suffisante, car, pour
$V'=\mathfrak J_\lambda M$ dans la condition b) de (19.1.5), on répond à la
question en prenant $W''=\mathfrak J_\lambda N$, $V''=V'$ et $h$ tel que
$h\circ u_\lambda$ soit l'identité dans $M/\mathfrak J_\lambda M$.
Inversement, si $u$ est formellement inversible à gauche, alors, pour tout
$\lambda$, il y a, en vertu de (19.1.5, b)), un $\mu$ tel que
$\mathfrak J_\mu\subset\mathfrak J_\lambda$ et un homomorphisme
$h:N/\mathfrak J_\mu N\to M/\mathfrak J_\lambda M$ tels que l'homomorphisme
canonique $M/\mathfrak J_\mu M\to M/\mathfrak J_\lambda M$ se factorise en
\[
  M/\mathfrak J_\mu M\xrightarrow{u_\mu}N/\mathfrak J_\mu N
  \xrightarrow{h}M/\mathfrak J_\lambda M~;
\]
mais comme
$h(\mathfrak J_\lambda(N/\mathfrak J_\mu N))
\subset\mathfrak J_\lambda(M/\mathfrak J_\lambda M)=0$, $h$ se factorise
canoniquement en
\[
  N/\mathfrak J_\mu N\longrightarrow N/\mathfrak J_\lambda N
  \xrightarrow{s}M/\mathfrak J_\lambda M,
\]
et il est immédiat que $s$ est inverse à gauche de $u_\lambda$.

\begin{proposition}[19.1.8]
\label{0.19.1.8-fr}
Soit $A$ un anneau topologique admettant un système fondamental dénombrable
décroissant $(\mathfrak J_n)$ d'idéaux ouverts. Soient $M$, $N$ deux
$A$-modules topologiques dont les topologies sont déduites de celle de $A$,
$u:M\to N$ un $A$-homomorphisme. Posons
$M_n=M/\mathfrak J_nM$, $N_n=N/\mathfrak J_nN$ et soit
$u_n:M_n\to N_n$ le $(A/\mathfrak J_n)$-homomorphisme déduit de $u$ par
passage aux quotients. Supposons que, pour tout $n$,
$P_n=\operatorname{Coker}(u_n)$ soit un $(A/\mathfrak J_n)$-module projectif
et que $M$ soit séparé et complet. Alors, pour que $u$ soit formellement
inversible à gauche, il faut et il suffit que $u$ soit inversible à gauche (et
$u$ est alors un isomorphisme topologique de $M$ sur un facteur direct
topologique de $N$).
\end{proposition}

\oldpage[0\textsubscript{IV-1}]{74}

\begin{proof}
En vertu de (19.1.7), on a des diagrammes commutatifs
\[
  \xymatrix{
    0\ar[r] & M_n\ar[r]^{u_n} & N_n\ar[r]^{p_n} & P_n\ar[r] & 0\\
    0\ar[r] & M_{n+1}\ar[u]^f\ar[r]_{u_{n+1}} &
      N_{n+1}\ar[u]^g\ar[r]_{p_{n+1}} & P_{n+1}\ar[u]^h\ar[r] & 0
  }
\]
où les lignes sont des suites exactes \emph{scindées}, autrement dit, il existe
pour chaque $n$ un homomorphisme $s_n:P_n\to N_n$ tel que
$p_n\circ s_n=1_{P_n}$. Nous allons montrer que l'on peut, par récurrence sur
$n$, définir un homomorphisme $s'_n:P_n\to N_n$ tel que
$p_n\circ s'_n=1_{P_n}$ et que les diagrammes
\[
  \xymatrix{
    P_n\ar[r]^{s'_n} & N_n\\
    P_{n+1}\ar[u]^h\ar[r]_{s'_{n+1}} & N_{n+1}\ar[u]^g
  }
\]
soient commutatifs. En effet, $g\circ s_{n+1}-s'_n\circ h$ est un
homomorphisme de $P_{n+1}$ dans
\[
  u_n(M_n)=u_{n+1}(M_{n+1})/
  \mathfrak J_nu_{n+1}(M_{n+1})~;
\]
comme $P_{n+1}$ est un $(A/\mathfrak J_{n+1})$-module \emph{projectif}, cet
homomorphisme se factorise en
\[
  P_{n+1}\xrightarrow{t_{n+1}}u_{n+1}(M_{n+1})
  \longrightarrow
  u_{n+1}(M_{n+1})/\mathfrak J_nu_{n+1}(M_{n+1}),
\]
d'où l'on conclut aussitôt que $s'_{n+1}=s_{n+1}-t_{n+1}$ répond à la
question. Cela étant, de la décomposition de $N_n$ en somme directe de
$u_n(M_n)$ et de $s'_n(P_n)$, on déduit aussitôt un homomorphisme
$w_n:N_n\to M_n$ inverse à gauche de $u_n$ et tel que les diagrammes
\[
  \xymatrix{
    N_n\ar[r]^{w_n} & M_n\\
    N_{n+1}\ar[u]^g\ar[r]_{w_{n+1}} & M_{n+1}\ar[u]^f
  }
\]
soient commutatifs. Le système projectif $(w_n)$ admet alors une limite
projective $\widehat w:\widehat N\to\widehat M=M$, d'où par composition avec
l'homomorphisme canonique $N\to\widehat N$, un homomorphisme $w:N\to M$ tel
que, pour tout $n$, l'endomorphisme $(w\circ u)_n$ de
$M_n=M/\mathfrak J_nM$ déduit de $w\circ u$ par passage aux quotients soit
l'identité~; cela entraîne que $w\circ u=1_M$ puisque $M$ est séparé et
complet.
\end{proof}

\begin{proposition}[19.1.9]
\label{0.19.1.9-fr}
Soient $A$ un anneau topologique préadmissible ($\mathbf{0}_{\mathrm I}$,
7.1.2), $\mathfrak Q$ un idéal de définition de $A$,
$(\mathfrak J_\lambda)$ un système fondamental d'idéaux ouverts de $A$.
Soient $M$, $N$ deux $A$-modules topologiques dont les topologies sont déduites
de celles de $A$, et tels que pour tout $\lambda$,
$N/\mathfrak J_\lambda N$ soit un $(A/\mathfrak J_\lambda)$-module projectif
(cf. (19.2.3)). Soit $u:M\to N$ un $A$-homomorphisme. Alors les conditions
suivantes sont équivalentes~:
\begin{enumerate}
  \item[a)] $u$ est formellement inversible à gauche.

  \oldpage[0\textsubscript{IV-1}]{75}

  \item[b)] L'homomorphisme
  $u_0:M/\mathfrak QM\to N/\mathfrak QN$ déduit de $u$ par passage aux
  quotients est inversible à gauche.
\end{enumerate}
\end{proposition}

\begin{proof}
On a vu (19.1.7) que la condition a) équivaut à dire que $u_\lambda$ est
inversible à gauche pour tout $\lambda$~; comme $\mathfrak Q$ est un idéal
ouvert, donc contient un $\mathfrak J_\lambda$, on en déduit aussitôt que
$u_0$ est inversible à gauche. Pour montrer inversement que b) entraîne a),
notons que pour tout $\lambda$, $\mathfrak Q/\mathfrak J_\lambda$ est par
hypothèse un idéal nilpotent de $A/\mathfrak J_\lambda$. Notre assertion
résultera de la proposition suivante.
\end{proof}

\begin{proposition}[19.1.10]
\label{0.19.1.10-fr}
Soient $A$ un anneau, $M$, $N$ deux $A$-modules, $N$ étant projectif,
$u:M\to N$ un $A$-homomorphisme. Soit $\mathfrak J$ un idéal de $A$~; on
suppose vérifiée l'une des conditions suivantes~:
\begin{enumerate}
  \item[(i)] $\mathfrak J$ est nilpotent.
  \item[(ii)] $\mathfrak J$ est contenu dans le radical de $A$ et $M$ est de
  type fini.
\end{enumerate}
Alors, pour que $u$ soit inversible à gauche, il faut et il suffit que
l'homomorphisme
$u_0:M/\mathfrak JM\to N/\mathfrak JN$ de $(A/\mathfrak J)$-modules, déduit
de $u$ par passage aux quotients, soit inversible à gauche.
\end{proposition}

\begin{proof}
La condition étant évidemment nécessaire, prouvons qu'elle est suffisante.
Soit $v_0$ un inverse à gauche de $u_0$~; l'homomorphisme composé
$N\to N/\mathfrak JN\xrightarrow{v_0}M/\mathfrak JM$ se factorise en
$N\xrightarrow{v}M\to M/\mathfrak JM$ puisque $N$ est projectif~; alors
$w=v\circ u$ est un endomorphisme de $M$ tel que l'endomorphisme $w_0$ de
$M/\mathfrak JM$ déduit par passage aux quotients soit l'identité, et il suffit
de prouver que $w$ est lui-même bijectif (car alors $w^{-1}\circ v$ sera un
inverse à gauche de $u$). Distinguons maintenant les deux cas.

(i) Pour tout $n$ on a le diagramme commutatif
\[
  \xymatrix{
    (\mathfrak J^n/\mathfrak J^{n+1})\otimes_{A/\mathfrak J}
      (M/\mathfrak JM)\ar[r]\ar[d]_{1\otimes\operatorname{gr}^0(w)} &
      \mathfrak J^nM/\mathfrak J^{n+1}M\ar[d]^{\operatorname{gr}^n(w)}\\
    (\mathfrak J^n/\mathfrak J^{n+1})\otimes_{A/\mathfrak J}
      (M/\mathfrak JM)\ar[r] & \mathfrak J^nM/\mathfrak J^{n+1}M
  }
\]
où les flèches horizontales sont \emph{surjectives}, et comme
$\operatorname{gr}^0(w)=w_0$ est l'identité, il en est de même de
$\operatorname{gr}^n(w)$ qui, \emph{a fortiori}, est bijectif. La filtration
$\mathfrak J$-préadique sur $M$ étant finie puisque $\mathfrak J$ est
nilpotent, on en conclut que $w$ est bijectif (Bourbaki, \emph{Alg. comm.},
chap.~III, §~2, n\textsuperscript{o}~8, cor.~3 du th.~1).

(ii) Il suffit de montrer que pour tout idéal maximal $\mathfrak m$ de $A$,
l'endomorphisme $w_{\mathfrak m}$ de $M_{\mathfrak m}$ est bijectif
(Bourbaki, \emph{Alg. comm.}, chap.~II, §~3, n\textsuperscript{o}~3, th.~1)
et comme $\mathfrak JA_{\mathfrak m}\subset\mathfrak mA_{\mathfrak m}$, et
$A_{\mathfrak m}/\mathfrak JA_{\mathfrak m}=(A/\mathfrak J)_{\mathfrak m}$,
on est ramené à prouver la proposition lorsque $A$ est un anneau local. En
outre, on peut supposer que $\mathfrak J$ est l'idéal maximal
$\mathfrak J_0$ de $A$, car si $u_0$ est inversible

\oldpage[0\textsubscript{IV-1}]{76}

à gauche, il en est de même de
$u_{00}:M/\mathfrak J_0M\to N/\mathfrak J_0M$ obtenu par tensorisation de
$u_0$ avec $1_{A/\mathfrak J_0}$, puisque l'on a
\[
  M/\mathfrak J_0M=(M/\mathfrak JM)\otimes_{A/\mathfrak J}(A/\mathfrak J_0).
\]
Supposons donc $\mathfrak J$ maximal, de sorte que $A/\mathfrak J$ est un
corps. Il suffit évidemment de montrer que $M$ est un $A$-module \emph{libre}
sous les conditions de l'énoncé~: en effet, $\det(w_0)$ est alors l'image
canonique de $\det(w)$ dans $A/\mathfrak J$, donc $\det(w)$ n'appartient pas à
l'idéal $\mathfrak J$ et est par suite inversible. Or le
$(A/\mathfrak J)$-espace vectoriel $M/\mathfrak JM$ étant libre de type fini,
il y a un $A$-module de type fini $L$ et un $A$-homomorphisme $f:L\to M$ tel
que l'homomorphisme $f_0:L/\mathfrak JL\to M/\mathfrak JM$ déduit de $f$ par
passage aux quotients soit bijectif. Comme $M$ est de type fini, on en conclut
tout d'abord que $f$ est surjectif par le lemme de Nakayama (Bourbaki,
\emph{Alg. comm.}, chap.~II, §~3, n\textsuperscript{o}~2, cor.~1 de la
prop.~4)~; en outre, si $g=u\circ f$, l'homomorphisme
$g_0:L/\mathfrak JL\to N/\mathfrak JN$ déduit par passage aux quotients est
inversible à gauche, et puisqu'ici $L$ est libre, la remarque du début prouve
que $g$ est lui-même inversible à gauche~; or cela entraîne évidemment que
$f$ est \emph{injectif}, ce qui achève la démonstration.
\end{proof}

Mentionnons en passant les corollaires suivants de (19.1.10)~:

\begin{corollary}[19.1.11]
\label{0.19.1.11-fr}
Soient $A$ un anneau local, $k$ son corps résiduel, $M$ un $A$-module de type
fini, $N$ un $A$-module projectif, $u:M\to N$ un homomorphisme. Pour que $u$
soit inversible à gauche, il faut et il suffit qu'il existe un système de
générateurs $(x_i)_{1\leq i\leq m}$ de $M$ tel que les images par
$u\otimes 1:M\otimes_Ak\to N\otimes_Ak$ des $x_i\otimes1$ soient linéairement
indépendantes dans $N\otimes_Ak$~; les $x_i$ forment alors une base de $M$.
\end{corollary}

La condition est évidemment nécessaire, car si $u$ est inversible à gauche,
$M$ est un $A$-module projectif de type fini, donc libre. Inversement, si la
condition est satisfaite, il est clair que les $x_i\otimes1$ forment une base
de $M\otimes_Ak$ et que $u\otimes1$ est inversible à gauche~; il suffit donc
d'appliquer (19.1.10) à l'idéal maximal $\mathfrak J$ de $A$.

\begin{corollary}[19.1.12]
\label{0.19.1.12-fr}
Soient $A$ un anneau, $M$ un $A$-module de type fini, $N$ un $A$-module
projectif, $u:M\to N$ un homomorphisme. Pour tout idéal premier
$\mathfrak p\in\operatorname{Spec}(A)$, les conditions suivantes sont
équivalentes~:
\begin{enumerate}
  \item[a)] L'homomorphisme
  $u_{\mathfrak p}:M_{\mathfrak p}\to N_{\mathfrak p}$ est inversible à
  gauche.
  \item[b)] L'homomorphisme
  $u\otimes1:M\otimes_Ak(\mathfrak p)\to N\otimes_Ak(\mathfrak p)$ est
  injectif (on rappelle que $k(\mathfrak p)$ désigne le corps résiduel de $A$
  en $\mathfrak p$).
  \item[c)] Il existe un système fini d'éléments $x_i\in M$
  $(1\leq i\leq m)$ tel que les images des $x_i$ dans $M_{\mathfrak p}$
  engendrent $M_{\mathfrak p}$, et un système de $m$ formes linéaires
  $y_i^*$ sur $N$ $(1\leq i\leq m)$ tel que
  $\det(\langle y_i^*,u(x_j)\rangle)\notin\mathfrak p$.
  \item[d)] Il existe $f\in A-\mathfrak p$ tel que l'homomorphisme
  $u_f:M_f\to N_f$ soit inversible à gauche.
\end{enumerate}
En outre, l'ensemble des $\mathfrak p\in\operatorname{Spec}(A)$ vérifiant ces
conditions est ouvert dans $\operatorname{Spec}(A)$.
\end{corollary}

La dernière condition est conséquence triviale de d). Comme $N$ est
projectif, il est facteur direct d'un $A$-module libre $A^{(I)}$~; en outre,
comme $M$ est de type fini, $u(M)$ est contenu dans un sous-module de $A^{(I)}$
de la forme $A^n$ pour $n$ fini~; comme chacun des énoncés a), b), c), d) est
équivalent à l'énoncé correspondant où l'on remplace $N$ par $N\oplus P$ ($u$
étant toujours supposé appliquer $M$ dans $N$), on est ramené au cas où $N$
est \emph{libre de type fini}. Il est trivial que d) entraîne a), et a) et b)
sont équivalentes en vertu de (19.1.11)~; d'ailleurs a) entraîne que
$M_{\mathfrak p}$ est libre (19.1.11), donc a) entraîne c), en prenant pour les
$x_i$ une famille dont les images dans $M_{\mathfrak p}$ forment une base de ce

\oldpage[0\textsubscript{IV-1}]{77}

$A_{\mathfrak p}$-module et en notant que (puisque $N$ est libre de type fini)
toute forme linéaire sur le $A_{\mathfrak p}$-module $N_{\mathfrak p}$ s'écrit
$u=v/f$, où $f\in A-\mathfrak p$, et où $v$ est une forme linéaire sur le
$A$-module $N$. Il est clair que c) entraîne b), et il reste donc à voir que a)
entraîne d). Or, comme $N$ est de présentation finie,
$(\operatorname{Hom}_A(N,M))_{\mathfrak p}$ s'identifie canoniquement à
$\operatorname{Hom}_{A_{\mathfrak p}}(N_{\mathfrak p},M_{\mathfrak p})$
(Bourbaki, \emph{Alg. comm.}, chap.~II, §~2, n\textsuperscript{o}~7,
prop.~19). Si $w_{\mathfrak p}$ est inverse à gauche de $u_{\mathfrak p}$, il
existe donc un homomorphisme $w:N\to M$ et un élément $f\in A-\mathfrak p$
tels que $w=w_{\mathfrak p}\otimes(1/f)1_{A_{\mathfrak p}}$. La relation
$w_{\mathfrak p}\circ u_{\mathfrak p}=1_{M_{\mathfrak p}}$ s'écrit donc aussi
$(w\circ u)\otimes1_{A_{\mathfrak p}}=f.1_{M_{\mathfrak p}}$. Mais comme $M$
est un $A$-module de type fini, il existe $g\in A-\mathfrak p$ tel que
l'endomorphisme $g((w\circ u)-f.1_M)$ s'annule en tous les générateurs de $M$,
donc soit nul. Si on pose $h=fg$, et
$u_h=u\otimes1_{A_h}$, $w_h=w\otimes1_{A_h}$, on a donc
$w_h(u_h(z))=(f/1)z$ pour tout $z\in M_h$~; mais comme $h/1$ est inversible
dans $A_h$, il en est de même de $f/1=f'$, et $f'^{-1}w_h$ est par suite
inverse à gauche de $u_h$, ce qui achève la démonstration.

\begin{proposition}[19.1.13]
\label{0.19.1.13-fr}
Supposons vérifiées les hypothèses de (19.1.9) et supposons en outre que
$(\mathfrak J_\lambda)$ soit une suite décroissante $(\mathfrak J_n)$ et que
$M$ soit séparé et complet. Alors les conditions a) et b) de (19.1.9) sont
aussi équivalentes à~:
\begin{enumerate}
  \item[c)] $u$ est inversible à gauche.
\end{enumerate}
\end{proposition}

On sait déjà que c) entraîne a) (19.1.6). Inversement, si a) est vérifiée, on
sait (avec les notations de (19.1.8)) que $M_n$ est facteur direct du
$(A/\mathfrak J_n)$-module projectif $N_n$, donc $P_n$ est aussi isomorphe à
un facteur direct de $N_n$, et par suite est projectif~; il suffit donc
d'appliquer (19.1.8).

Notons enfin la proposition suivante~:

\begin{proposition}[19.1.14]
\label{0.19.1.14-fr}
Soient $A$ un anneau, $M$ un $A$-module de type fini, $N$ un $A$-module
projectif, $u:M\to N$ un $A$-homomorphisme.
\begin{enumerate}
  \item[(i)] Pour que $u$ soit inversible à gauche, il faut et il suffit que,
  pour tout idéal maximal $\mathfrak m$ de $A$,
  $u_{\mathfrak m}:M_{\mathfrak m}\to N_{\mathfrak m}$ soit inversible à
  gauche.
  \item[(ii)] Soit $A'$ une $A$-algèbre qui soit un $A$-module fidèlement
  plat. Pour que $u$ soit inversible à gauche, il faut et il suffit que
  $u\otimes1:M\otimes_AA'\to N\otimes_AA'$ soit inversible à gauche.
\end{enumerate}
\end{proposition}

Comme dans (19.1.12), on peut se borner au cas où $N$ est libre de type fini~;
dire que $u$ est inversible à gauche signifie alors que $u$ est injectif et
que le module quotient $P=N/u(M)$ est \emph{projectif}, car $u(M)$ sera alors
facteur direct de $N$. Notons en outre que puisque $M$ est de type fini, $P$
est de \emph{présentation finie}. Cela étant~:

(i) La condition est évidemment nécessaire. Inversement, si elle est
satisfaite, on sait que $u$ est injectif (Bourbaki, \emph{Alg. comm.},
chap.~II, §~3, n\textsuperscript{o}~3, th.~1) et comme
$P_{\mathfrak m}=N_{\mathfrak m}/u_{\mathfrak m}(M_{\mathfrak m})$ est
projectif pour tout $\mathfrak m$, on sait que cela entraîne que $P$ est
projectif (\emph{loc. cit.}, §~5, n\textsuperscript{o}~2, th.~1).

(ii) Ici encore, la condition est trivialement nécessaire. Inversement, si
elle est remplie, on sait que $u$ est injectif ($\mathbf{0}_{\mathrm I}$,
6.4.1) et comme $P\otimes_AA'=\operatorname{Coker}(u\otimes1)$ est projectif,
donc plat, on en déduit que $P$ est un $A$-module plat
($\mathbf{0}_{\mathrm I}$, 6.6.3), donc projectif puisqu'il est de présentation
finie (Bourbaki, \emph{Alg. comm.}, chap.~II, §~5,
n\textsuperscript{o}~2, cor.~2 du th.~1).

\oldpage[0\textsubscript{IV-1}]{78}

\begin{remark}[19.1.15]
\label{0.19.1.15-fr}
La notion «~duale~» de celle d'homomorphisme formellement inversible à gauche
est celle d'homomorphisme \emph{formellement inversible à droite}~; un tel
$A$-homomorphisme continu $u:M\to N$ vérifie, par définition, la condition
suivante~: pour tout sous-module ouvert $W'$ de $N$, il existe un sous-module
ouvert $V'\subset u^{-1}(W')$ de $M$, un sous-module ouvert $W''\subset W'$ de
$N$ et un homomorphisme $h:N/W''\to M/V'$ tels que l'homomorphisme canonique
$N/W''\to N/W'$ se factorise en
\[
  N/W''\xrightarrow{h}M/V'\xrightarrow{u'}N/W',
\]
où $u'$ est déduit de $u$ par passage aux quotients. Cela entraîne que $u$ est
un épimorphisme formel~; s'il existe un $\widehat A$-homomorphisme continu
$r:\widehat N\to\widehat M$ tel que $\widehat u\circ r=1_{\widehat N}$, on
vérifie aussitôt que $u$ est formellement inversible à droite. Si les
conditions de (19.1.7) sont remplies, pour que $u$ soit formellement inversible
à droite, il faut et il suffit que les $u_\lambda$ soient inversibles à droite
(c'est-à-dire que le noyau de $u_\lambda$ est facteur direct de
$M/\mathfrak J_\lambda M$ et que $u_\lambda$ est un isomorphisme d'un
supplémentaire de $\operatorname{Ker}(u_\lambda)$ sur
$N/\mathfrak J_\lambda N$). Nous laissons au lecteur le soin d'énoncer et de
démontrer les propositions correspondant à (19.1.8) et (19.1.9) (dans
l'analogue de (19.1.8), il faut supposer $M$ séparé et complet et que $M_n$ est
un $(A/\mathfrak J_n)$-module projectif~; dans l'analogue de (19.1.9), il faut
supprimer l'hypothèse sur les $N/\mathfrak J_\lambda N$, mais supposer par
contre que, pour tout $\lambda$, $M/\mathfrak J_\lambda M$ est un
$(A/\mathfrak J_\lambda)$-module projectif).
\end{remark}

\subsection{Modules formellement projectifs}
\label{subsection:0.19.2-fr}

\begin{definition}[19.2.1]
\label{0.19.2.1-fr}
Soient $A$ un anneau topologique. On dit qu'un $A$-module topologique $M$ est
\emph{formellement projectif} s'il vérifie la condition suivante~: pour tout
idéal ouvert $\mathfrak J$ de $A$, tout couple de $(A/\mathfrak J)$-modules
(discrets) $P$, $Q$, tout $A$-homomorphisme surjectif $u:P\to Q$ et tout
$A$-homomorphisme continu $v:M\to Q$, il existe un $A$-homomorphisme continu
$w:M\to P$ tel que $v=u\circ w$.
\end{definition}

\begin{env}[19.2.2]
\label{0.19.2.2-fr}
Pour vérifier la condition de (19.2.1), on peut évidemment se borner (en
remplaçant $Q$ par $v(M)$ et $P$ par $u^{-1}(v(M))$) au cas où $v$ est
lui-même surjectif~; alors $Q$ est isomorphe à $M/V$, où $V$ est un
sous-module ouvert de $M$ tel que $\mathfrak JM\subset V$~; la condition de
(19.2.1) équivaut alors à dire que pour tout $(A/\mathfrak J)$-module discret
$P$ et tout $A$-homomorphisme surjectif $u:P\to M/V$, il existe un sous-module
ouvert $V'\subset V$ dans $M$ et un $A$-homomorphisme $w:M/V'\to P$, tels que
l'homomorphisme canonique $M/V'\to M/V$ se factorise en
\[
  M/V'\xrightarrow{w}P\xrightarrow{u}M/V.
\]
On notera qu'il suffit de vérifier cette condition lorsque $\mathfrak J$
parcourt un système fondamental de voisinages de 0 dans $A$ formé d'idéaux.
\end{env}

\begin{env}[19.2.3]
\label{0.19.2.3-fr}
Supposons qu'il existe un système fondamental $(\mathfrak J_\lambda)$ d'idéaux
ouverts dans $A$ et un système fondamental $(V_\lambda)$ de sous-modules
ouverts de $M$, ayant \emph{même} ensemble d'indices que
$(\mathfrak J_\lambda)$ et tel que, pour tout $\lambda$, $M/V_\lambda$ soit un
$(A/\mathfrak J_\lambda)$-module \emph{projectif}. Alors $M$ est formellement
projectif~: il suffit en effet, avec les notations de (19.2.2), de prendre
$\lambda$ tel que $\mathfrak J_\lambda\subset\mathfrak J$ et
$V_\lambda\subset V$~; comme $P$ est aussi un
$(A/\mathfrak J_\lambda)$-module, la factorisation de l'homomorphisme canonique
$M/V_\lambda\to M/V$ en $M/V_\lambda\to P\to M/V$ résulte de ce qu'il s'agit
d'un $(A/\mathfrak J_\lambda)$-homomorphisme et de ce que $M/V_\lambda$ est
supposé être un $(A/\mathfrak J_\lambda)$-module projectif.

Lorsque la condition plus stricte de ce numéro est satisfaite, on dit que $M$
est \emph{strictement formellement projectif}.
\end{env}

\begin{proposition}[19.2.4]
\label{0.19.2.4-fr}
Soient $A$ un anneau topologique, $M$ un $A$-module topologique dont la
topologie est déduite de celle de $A$. Pour que $M$ soit formellement
projectif, il faut et il suffit qu'il soit strictement formellement projectif.
\end{proposition}

\oldpage[0\textsubscript{IV-1}]{79}

\begin{proof}
On vient de voir que la condition est suffisante. Inversement, supposons $M$
formellement projectif et soit $(\mathfrak J_\lambda)$ un système fondamental
d'idéaux ouverts dans $A$. Pour tout $\lambda$, soient $P$ un
$(A/\mathfrak J_\lambda)$-module libre et
$u:P\to M/\mathfrak J_\lambda M$ un
$(A/\mathfrak J_\lambda)$-homomorphisme surjectif. Il existe donc un
$\mathfrak J_\mu\subset\mathfrak J_\lambda$ tel que l'homomorphisme canonique
$M/\mathfrak J_\mu M\to M/\mathfrak J_\lambda M$ se factorise en
\[
  M/\mathfrak J_\mu M\xrightarrow{w}P\xrightarrow{u}M/\mathfrak J_\lambda M;
\]
mais comme
$w(\mathfrak J_\lambda(M/\mathfrak J_\mu M))\subset\mathfrak J_\lambda P=0$,
$w$ se factorise en $M/\mathfrak J_\lambda M\xrightarrow{v}P$, et il est
clair que $u\circ v$ est l'identité dans $M/\mathfrak J_\lambda M$, ce qui
prouve que ce $(A/\mathfrak J_\lambda)$-module est projectif.
\end{proof}

\begin{proposition}[19.2.5]
\label{0.19.2.5-fr}
Soient $A$ un anneau topologique, $M$ un $A$-module topologique.
\begin{enumerate}
  \item[(i)] Pour que $M$ soit formellement projectif (resp. strictement
  formellement projectif), il faut et il suffit que le $\hat A$-module
  topologique $\hat M$ le soit.
  \item[(ii)] Soit $A'$ une $A$-algèbre topologique. Si $M$ est formellement
  projectif (resp. strictement formellement projectif), alors
  $M\otimes_A A'$ (muni de la topologie produit tensoriel) est un $A'$-module
  topologique formellement projectif (resp. strictement formellement
  projectif).
\end{enumerate}
\end{proposition}

\begin{proof}
\begin{enumerate}
  \item[(i)] Il suffit de remarquer que lorsque $\mathfrak J$ (resp. $V$)
  parcourt l'ensemble des idéaux ouverts de $A$ (resp. l'ensemble des
  sous-modules ouverts de $M$), le séparé complété $\hat{\mathfrak J}$
  (resp. $\hat V$) parcourt l'ensemble des idéaux ouverts de $\hat A$
  (resp. l'ensemble des sous-modules ouverts de $\hat M$), et
  $\hat A/\hat{\mathfrak J}=A/\mathfrak J$ (resp.
  $\hat M/\hat V=M/V$) à un isomorphisme canonique près~; comme la notion
  de module formellement projectif (resp. strictement formellement projectif)
  ne fait intervenir que les $A/\mathfrak J$ et les $M/V$, on en déduit
  aussitôt (i).
  \item[(ii)] Supposons d'abord $M$ formellement projectif et posons
  $M'=M\otimes_A A'$~; soient $\mathfrak J'$ un idéal ouvert de $A'$, $P'$,
  $Q'$ deux $(A'/\mathfrak J')$-modules discrets, $u':P'\to Q'$ un
  $A'$-homomorphisme surjectif, $v':M'\to Q'$ un $A'$-homomorphisme continu.
  Il y a un idéal ouvert $\mathfrak J$ de $A$ tel que
  $\mathfrak JA'\subset\mathfrak J'$, donc $P'$ et $Q'$ sont aussi des
  $(A/\mathfrak J)$-modules discrets. Si l'on considère le
  $A$-homomorphisme composé $v:M\xrightarrow{j}M'\xrightarrow{v'}Q'$ qui est
  continu, l'hypothèse entraîne qu'il existe un $A$-homomorphisme continu
  $w:M\to P'$ tel que $v=u'\circ w$~; mais comme $P'$ est un $A'$-module
  topologique, $w$ se factorise en
  $M\xrightarrow{j}M'\xrightarrow{w'}P'$, où $w'$ est un $A'$-homomorphisme
  continu, et comme $v'\circ j=(u'\circ w')\circ j$, on en conclut que
  $v'=u'\circ w'$.

  Supposons ensuite que $M$ soit strictement formellement projectif~; soit
  $(\mathfrak J_\lambda)$ (resp. $(V_\lambda)$) un système fondamental
  d'idéaux ouverts (resp. de sous-modules ouverts) dans $A$ (resp. $M$) tel
  que $(M/V_\lambda)$ soit un $(A/\mathfrak J_\lambda)$-module projectif,
  et soit $(\mathfrak J'_\mu)$ un système fondamental d'idéaux ouverts dans
  $A'$. On a un système fondamental de voisinages de $0$ dans $M'$ en prenant
  les sous-modules
  \[
    \operatorname{Im}(M\otimes_A\mathfrak J'_\mu)
    +\operatorname{Im}(V_\lambda\otimes_A A')=W_{\lambda\mu},
  \]
  où $\lambda$ et $\mu$ sont tels que
  $\mathfrak J_\lambda A'\subset\mathfrak J'_\mu$. Comme
  \[
    (M\otimes_A A')/W_{\lambda\mu}
    =(M/V_\lambda)\otimes_{A/\mathfrak J_\lambda}
      (A'/\mathfrak J'_\mu)
  \]
  et que $(M/V_\lambda)$ est un
  $(A/\mathfrak J_\lambda)$-module projectif, $M'/W_{\lambda\mu}$
  est un $(A'/\mathfrak J'_\mu)$-module projectif.
\end{enumerate}
\end{proof}

\subsection{Algèbres formellement lisses}
\label{subsection:0.19.3-fr}

\begin{definition}[19.3.1]
\label{0.19.3.1-fr}
Soient $A$ un anneau topologique, $B$ une $A$-algèbre topologique. On dit que
$B$ est une $A$-algèbre \emph{formellement lisse} si, pour toute
$A$-algèbre topologique discrète $C$ et tout idéal nilpotent $\mathfrak J$ de
$C$, tout $A$-homomorphisme continu $u:B\to C/\mathfrak J$ se factorise en
\[
  B\xrightarrow{v}C\xrightarrow{\varphi}C/\mathfrak J,
\]
où $v$ est un homomorphisme continu et $\varphi$ l'homomorphisme canonique.
\end{definition}

\oldpage[0\textsubscript{IV-1}]{80}

La définition (19.3.1) revient à dire que la propriété suivante a lieu~:
\begin{enumerate}
  \item[(P)] Pour tout idéal ouvert $\mathfrak R$ de la $A$-algèbre $B$ et
  tout $A$-homomorphisme $u':B/\mathfrak R\to C/\mathfrak J$, il y a un
  idéal ouvert $\mathfrak R'\subset\mathfrak R$ de $B$ tel que
  l'homomorphisme
  $B\to B/\mathfrak R\xrightarrow{u'}C/\mathfrak J$ se factorise en
  \[
    B\to B/\mathfrak R'\xrightarrow{v'}C\to C/\mathfrak J,
  \]
  où $v'$ est un $A$-homomorphisme.
\end{enumerate}
En effet, si $u:B\to C/\mathfrak J$ est un $A$-homomorphisme continu, il a
pour noyau $\mathfrak R$ un idéal ouvert de $B$, donc $u$ se factorise en
$B\to B/\mathfrak R\xrightarrow{u'}C/\mathfrak J$, et si (P) est
vérifiée, il suffit de l'appliquer à $u'$ et de prendre pour $v$ la composée
$B\to B/\mathfrak R'\xrightarrow{v'}C$ pour satisfaire aux conditions de
(19.3.1). Inversement, supposons que $B$ soit une $A$-algèbre formellement
lisse~; donnons-nous un idéal ouvert $\mathfrak R$ de $B$ et un
$A$-homomorphisme $u':B/\mathfrak R\to C/\mathfrak J$, et appliquons la
définition (19.3.1) à
$u:B\to B/\mathfrak R\xrightarrow{u'}C/\mathfrak J$~; si
$v:B\to C$ est un $A$-homomorphisme continu tel que $u$ se factorise en
$B\xrightarrow{v}C\to C/\mathfrak J$, l'idéal
$\mathfrak R'=\operatorname{Ker}(v)\cap\mathfrak R$ de $B$ est ouvert et
contenu dans $\mathfrak R$~; par suite $v$ se factorise en
$B\to B/\mathfrak R'\xrightarrow{v'}C$, et $v'$ vérifie bien la condition
(P).

\begin{proposition}[19.3.2]
\label{0.19.3.2-fr}
Soient $A$ un anneau discret, $V$ un $A$-module projectif~; l'algèbre
symétrique $B=\mathbf S_A^\bullet(V)$, munie de la topologie discrète, est une
$A$-algèbre formellement lisse.
\end{proposition}

\begin{proof}
En effet, les notations $C$ et $\mathfrak J$ ayant le même sens que dans
(19.3.1), soit $u:B\to C/\mathfrak J$ un homomorphisme de $A$-algèbres, qui
par restriction à $V=\mathbf S_A^1(V)$, donne un homomorphisme de $A$-modules
$u_1:V\to C/\mathfrak J$~; comme $V$ est projectif, $u_1$ se factorise en
$V\xrightarrow{v_1}C\xrightarrow{\varphi}C/\mathfrak J$, et $v_1$ se
prolonge en un homomorphisme de $A$-algèbres
$v:\mathbf S_A^\bullet(V)\to C$, tel que le composé
$\varphi\circ v$ coïncide avec $u$ dans $V$, donc
$\varphi\circ v=u$, ce qui démontre la proposition.
\end{proof}

\begin{corollary}[19.3.3]
\label{0.19.3.3-fr}
Si $A$ est un anneau discret, toute algèbre de polynômes
$B=A[T_\alpha]_{\alpha\in I}$, munie de la topologie discrète, est une
$A$-algèbre formellement lisse.
\end{corollary}

\begin{proposition}[19.3.4]
\label{0.19.3.4-fr}
Soit $A$ un anneau topologique, et soit
$B=A[[T_\alpha]]_{\alpha\in I}$ un anneau de séries formelles
$\sum_{\lambda\in\mathbb{N}^{(I)}}c_\lambda T^\lambda$ (algèbre large sur
$A$ du monoïde $\mathbb{N}^{(I)}$, identifié en tant que $A$-module au produit
$A^{\mathbb{N}^{(I)}}$). Si l'on munit $B$ de la topologie produit, $B$ est une
$A$-algèbre formellement lisse.
\end{proposition}

Soit $(\mathfrak L_\mu)_{\mu\in M}$ un système fondamental d'idéaux ouverts
dans $A$. Pour toute partie finie $H$ de $I$, tout $\mu\in M$ et tout entier
$n$, désignons par $\mathfrak K_{H,\mu,n}$ le voisinage de $0$ dans $B$ formé
des $(c_\lambda)$ tels que pour tout
$\lambda=(\lambda_\alpha)_{\alpha\in I}$ tel que
$\lambda_\alpha=0$ pour $\alpha\notin H$ et
$\sum_{\alpha\in H}\lambda_\alpha\leq n$, on ait
$c_\lambda\in\mathfrak L_\mu$. On vérifie immédiatement que les
$\mathfrak K_{H,\mu,n}$ forment un système fondamental de voisinages de $0$
dans $B$ et sont des idéaux de $B$, donc la topologie produit est compatible
avec la structure de $A$-algèbre de $B$.

Notons d'abord, avec les mêmes notations, le
\begin{lemma}[19.3.4.1]
\label{0.19.3.4.1-fr}
Soit $E$ une $A$-algèbre discrète.
\begin{enumerate}
  \item[(i)] Si $f:B\to E$ est un $A$-homomorphisme continu, il existe une
  partie finie $H$ de $I$ telle que $f(T_\alpha)=0$ pour $\alpha\notin H$,
  et $f(T_\alpha)$ est nilpotent dans $E$ pour tout $\alpha\in H$.
  \item[(ii)] Inversement, soient $H$ une partie finie de $I$,
  $(z_\alpha)_{\alpha\in H}$ une famille d'éléments nilpotents de $E$.
  Il existe un $A$-homomorphisme continu $g:B\to E$ et un seul tel que
  $g(T_\alpha)=z_\alpha$ pour $\alpha\in H$ et
  $g(T_\alpha)=0$ pour $\alpha\notin H$.
\end{enumerate}
\end{lemma}

\oldpage[0\textsubscript{IV-1}]{81}

\begin{proof}
(i) résulte aussitôt de ce que $f^{-1}(0)$ est un voisinage de $0$ dans le
$A$-module produit $A^{\mathbb{N}^{(I)}}$, d'où $f(T^\lambda)=0$ sauf pour un
nombre fini de valeurs de $\lambda\in\mathbb{N}^{(I)}$. Pour prouver (ii), il
suffit de remarquer que l'anneau de polynômes
$B'=A[T_\alpha]_{\alpha\in I}$ est dense dans $B$~; l'existence et
l'unicité de la restriction $g|B'$ sont triviales et sa continuité résulte de
l'hypothèse que les $f(T_\alpha)$ pour $\alpha\in H$ sont nilpotents,
car si $(f(T_\alpha))^n=0$ pour tout $\alpha\in H$, on a
$g(T^\lambda)=0$ pour tout $\lambda=(\lambda_\alpha)_{\alpha\in I}$ de
support fini sauf ceux tels que $\lambda_\alpha=0$ pour
$\alpha\notin H$ et $\lambda_\alpha\leq n$ pour
$\alpha\in H$, c'est-à-dire sauf pour un nombre fini de valeurs de
$\lambda$.
\end{proof}

\begin{proof}[Démonstration de la proposition~{\rm(19.3.4)}]
Ce lemme étant établi, et les notations $C$ et $\mathfrak J$ ayant le même
sens que dans (19.3.1), on a $u(T_\alpha)=0$ sauf pour
$\alpha\in H$, $H$ étant une partie finie de $I$, et les
$z_\alpha=u(T_\alpha)$ pour $\alpha\in H$ sont nilpotents dans
$C/\mathfrak J$~; comme $\mathfrak J$ est nilpotent, il existe une famille
$(x_\alpha)_{\alpha\in H}$ d'éléments nilpotents de $C$ dont les images
canoniques dans $C/\mathfrak J$ sont les $z_\alpha$~; si $v$ est le
$A$-homomorphisme continu de $B$ dans $C$ tel que
$v(T_\alpha)=0$ pour $\alpha\notin H$, $v(T_\alpha)=x_\alpha$ pour
$\alpha\in H$, il est clair que $u$ se factorise en
$B\xrightarrow{v}C\to C/\mathfrak J$.
\end{proof}

\begin{proposition}[19.3.5]
\label{0.19.3.5-fr}
Soit $A$ un anneau topologique.
\begin{enumerate}
  \item[(i)] $A$ est une $A$-algèbre formellement lisse.
  \item[(ii)] Si $B$ est une $A$-algèbre formellement lisse et $C$ une
  $B$-algèbre formellement lisse, alors $C$ est une $A$-algèbre formellement
  lisse.
  \item[(iii)] Soient $B$ une $A$-algèbre formellement lisse, $A'$ une
  $A$-algèbre topologique~; alors la $A'$-algèbre topologique
  $B\otimes_A A'$ ($\mathbf 0_{\rm I}$, 7.7.5 et 7.7.6) est formellement
  lisse.
  \item[(iv)] Soient $B$ une $A$-algèbre topologique, $S$ (resp. $T$) une
  partie multiplicative de $A$ (resp. $B$) telle que l'image canonique de $S$
  dans $B$ soit contenue dans $T$. Si $B$ est une $A$-algèbre formellement
  lisse, alors $T^{-1}B$ est une $S^{-1}A$-algèbre formellement lisse.
  \item[(v)] Soient $B_i$ ($1\leq i\leq n$) des $A$-algèbres topologiques.
  Pour que $\prod_{i=1}^n B_i$ soit une $A$-algèbre formellement lisse, il
  faut et il suffit que chacune des $B_i$ le soit.
\end{enumerate}
\end{proposition}

\oldpage[0\textsubscript{IV-1}]{82}

\begin{proof}
\begin{enumerate}
  \item[(i)] Si $C$ est une $A$-algèbre discrète,
  $\varphi:C\to C/\mathfrak J$ l'homomorphisme canonique de $C$ sur une
  $A$-algèbre quotient quelconque de $C$, le \emph{seul} $A$-homomorphisme
  de $A$ dans $C/\mathfrak J$ est l'homomorphisme composé
  $A\xrightarrow{\psi}C\xrightarrow{\varphi}C/\mathfrak J$, où
  $\psi$ est l'homomorphisme définissant la structure de $A$-algèbre de
  $C$~; comme $\psi$ est continu, la condition de (19.3.1) est trivialement
  vérifiée.
  \item[(ii)] Soient $\alpha:A\to B$, $\beta:B\to C$ les homomorphismes
  continus définissant respectivement la structure de $A$-algèbre sur $B$ et
  celle de $B$-algèbre sur $C$, de sorte que $\beta\circ\alpha$ définit la
  structure de $A$-algèbre sur $C$. Soient $E$ une $A$-algèbre discrète,
  $\mathfrak L$ un idéal nilpotent de $E$, $u:C\to E/\mathfrak L$ un
  $A$-homomorphisme continu, de sorte que $u\circ\beta\circ\alpha$ est
  l'homomorphisme définissant la structure de $A$-algèbre de
  $E/\mathfrak L$. Comme $B$ est une $A$-algèbre formellement lisse, le
  $A$-homomorphisme continu $u\circ\beta:B\to E/\mathfrak L$ se
  factorise en $B\xrightarrow{v}E\to E/\mathfrak L$, où $v$ est un
  $A$-homomorphisme continu~; $v$ et $u\circ\beta$ définissent alors sur
  $E$ et $E/\mathfrak L$ respectivement des structures de $B$-algèbre
  topologique, pour lesquelles $E/\mathfrak L$ est encore l'algèbre quotient
  (discrète) de la $B$-algèbre $E$. En outre, $u$ est un $B$-homomorphisme
  continu, donc se factorise en $C\xrightarrow{w}E\to E/\mathfrak L$,
  où $w$ est un $B$-homomorphisme continu~; comme
  $v\circ\alpha$ est le $A$-homomorphisme définissant la structure de
  $A$-algèbre sur $E$, $w$ est bien un $A$-homomorphisme continu, d'où notre
  assertion.
  \item[(iii)] Soient $C$ une $A'$-algèbre topologique discrète,
  $\mathfrak J$ un idéal nilpotent de $C$,
  $u:B\otimes_A A'\to C/\mathfrak J$ un $A'$-homomorphisme continu. Si on
  compose $u$ et l'homomorphisme canonique
  $\rho:B\to B\otimes_A A'$, on obtient (puisque
  $A\to B\xrightarrow{\rho}B\otimes_A A'$ est égal au composé
  $A\to A'\xrightarrow{\sigma}B\otimes_A A'$) un $A$-homomorphisme
  continu qui, par hypothèse, se factorise donc en
  $B\xrightarrow{v}C\to C/\mathfrak J$, où $v$ est un $A$-homomorphisme
  continu. L'égalité des composés $A\to B\xrightarrow{v}C$ et
  $A\to A'\xrightarrow{w}C$ entraîne donc l'existence d'un homomorphisme
  continu d'anneaux $f:B\otimes_A A'\to C$ tel que
  $v=f\circ\rho$ et $w=f\circ\sigma$ ($\mathbf 0_{\rm I}$, 7.7.6)~;
  comme les homomorphismes composés
  $B\xrightarrow{\rho}B\otimes_A A'\xrightarrow{f}C\to C/\mathfrak J$
  et
  $B\xrightarrow{\rho}B\otimes_A A'\xrightarrow{u}C/\mathfrak J$
  (resp.
  $A'\xrightarrow{\sigma}B\otimes_A A'\xrightarrow{f}C\to C/\mathfrak J$
  et
  $A'\xrightarrow{\sigma}B\otimes_A A'\xrightarrow{u}C/\mathfrak J$)
  sont égaux, on a bien la factorisation
  $u:B\otimes_A A'\xrightarrow{f}C\to C/\mathfrak J$, ce qui établit
  (iii).
  \item[(iv)] La topologie considérée sur $S^{-1}A$ (resp. $T^{-1}B$) est
  naturellement celle pour laquelle un système fondamental de voisinages de
  $0$ est formé des $S^{-1}\mathfrak J$ (resp. $T^{-1}\mathfrak R$), où
  $\mathfrak J$ (resp. $\mathfrak R$) parcourt un système fondamental
  d'idéaux ouverts dans $A$ (resp. $B$) ($\mathbf 0_{\rm I}$, 7.6.1). Si
  $\alpha:A\to B$ est l'homomorphisme canonique, il est clair que
  l'homomorphisme canonique $\alpha':S^{-1}A\to T^{-1}B$ déduit de
  $\alpha$ est continu. Soient alors $C$ une $S^{-1}A$-algèbre topologique
  discrète, $\mathfrak J$ un idéal nilpotent de cette algèbre,
  $u:T^{-1}B\to C/\mathfrak J$ un $S^{-1}A$-homomorphisme continu~; alors
  le composé $B\to T^{-1}B\xrightarrow{u}C/\mathfrak J$ est un
  $A$-homomorphisme continu qui, par hypothèse, se factorise en
  $B\xrightarrow{v}C\to C/\mathfrak J$, où $v$ est un $A$-homomorphisme
  continu. Comme pour tout $t\in T$, $t/1$ est inversible dans $T^{-1}B$,
  $u(t/1)$ est inversible dans $C/\mathfrak J$. Comme $\mathfrak J$ est
  nilpotent, tout élément de la classe $u(t/1)$ dans $C$, et en particulier
  $v(t)$, est inversible dans $C$ ($\mathbf 0_{\rm I}$, 7.1.12), et par suite
  $v$ se factorise en $B\to T^{-1}B\xrightarrow{w}C$~; comme $v$ est
  continu, il en est de même de $w$ ($\mathbf 0_{\rm I}$, 7.6.6), et c'est
  un $S^{-1}A$-homomorphisme car le composé
  $A\to B\to T^{-1}B\xrightarrow{w}C$ est égal à
  \[
    A\to S^{-1}A\xrightarrow{\alpha'}T^{-1}B\xrightarrow{w}C,
  \]
  donc $S^{-1}A\xrightarrow{\alpha'}T^{-1}B\xrightarrow{w}C$ est
  l'homomorphisme canonique définissant sur $C$ la structure de
  $S^{-1}A$-algèbre. Enfin, les homomorphismes composés
  $B\to T^{-1}B\xrightarrow{w}C\to C/\mathfrak J$ et
  $B\to T^{-1}B\xrightarrow{u}C/\mathfrak J$ étant égaux, le même
  raisonnement montre que $u$ est bien égal au composé
  $T^{-1}B\xrightarrow{w}C\to C/\mathfrak J$, d'où (iv).
  \item[(v)] Enfin, (v) est immédiat, la donnée d'un $A$-homomorphisme
  continu de $B=\prod_{i=1}^n B_i$ dans $C$ (resp. $C/\mathfrak J$)
  équivalant à celle de $n$ $A$-homomorphismes continus $B_i\to C$
  (resp. $B_i\to C/\mathfrak J$) et tout $A$-homomorphisme continu
  $B_j\to C$ (resp. $B_j\to C/\mathfrak J$) donnant par composition un
  $A$-homomorphisme continu $B\to B_j\to C$
  (resp. $B\to B_j\to C/\mathfrak J$).
\end{enumerate}
\end{proof}

\oldpage[0\textsubscript{IV-1}]{83}

\begin{proposition}[19.3.6]
\label{0.19.3.6-fr}
Soient $A$ un anneau topologique, $B$ une $A$-algèbre topologique, $\hat A$
et $\hat B$ les séparés complétés respectifs de $A$ et $B$, de sorte que
$\hat B$ est une $\hat A$-algèbre topologique. Les conditions suivantes
sont équivalentes~:
\begin{enumerate}
  \item[a)] $B$ est une $A$-algèbre formellement lisse.
  \item[b)] $\hat B$ est une $A$-algèbre formellement lisse.
  \item[c)] $\hat B$ est une $\hat A$-algèbre formellement lisse.
\end{enumerate}
\end{proposition}

\begin{proof}
Bien entendu, la structure de $\hat A$-algèbre sur $\hat B$ est définie
par l'homomorphisme $\hat\varphi$, si $\varphi:A\to B$ est
l'homomorphisme définissant la structure de $A$-algèbre sur $B$. Comme toute
$A$-algèbre discrète $C$ est séparée et complète, elle est aussi une
$\hat A$-algèbre (par prolongement canonique de l'homomorphisme de $A$ dans
$C$), et un $A$-homomorphisme continu de $B$ dans $C$ donne par prolongement
canonique un $\hat A$-homomorphisme (et \emph{a fortiori} un
$A$-homomorphisme) de $\hat B$ dans $C$ (autrement dit, tout
$A$-homomorphisme $B\to C$ se factorise en $B\to\hat B\to C$ de façon
unique). Ces remarques et la définition (19.3.1) entraînent aussitôt la
proposition.
\end{proof}

\begin{corollary}[19.3.7]
\label{0.19.3.7-fr}
Sous les hypothèses de (19.3.5), (iv), la $A\{S^{-1}\}$-algèbre topologique
$B\{T^{-1}\}$ est formellement lisse.
\end{corollary}

Cela résulte des définitions ($\mathbf 0_{\rm I}$, 7.6.1 et 7.6.7) et de
(19.3.5), (iv) et (19.3.6).

\begin{proposition}[19.3.8]
\label{0.19.3.8-fr}
Soient $A$ un anneau topologique, $B$ une $A$-algèbre topologique, et supposons
que pour tout idéal ouvert $\mathfrak R$ de $B$, $\mathfrak R^2$ soit aussi
ouvert. Soient $A'$, $B'$ des anneaux topologiques obtenus en munissant $A$ et
$B$ de topologies plus fines que les topologies initiales, et supposons que
l'homomorphisme canonique $\varphi:A\to B$ soit encore un homomorphisme
continu de $A'$ dans $B'$. Alors, si $B'$ est une $A'$-algèbre formellement
lisse, $B$ est une $A$-algèbre formellement lisse.
\end{proposition}

Il suffit d'appliquer le lemme suivant~:

\begin{lemma}[19.3.8.1]
\label{0.19.3.8.1-fr}
Soient $C$ une $A$-algèbre discrète, $\mathfrak J$ un idéal nilpotent de $C$.
Supposons que le carré de tout idéal ouvert de $B$ soit ouvert. Alors, si un
homomorphisme composé
$v:B\xrightarrow{u}C\to C/\mathfrak J$ est continu, l'homomorphisme $u$
est continu.
\end{lemma}

\begin{proof}
En effet, $\mathfrak R=u^{-1}(\mathfrak J)$ est le noyau de $v$ et il est
ouvert par hypothèse~; comme il existe un entier $n$ tel que
$\mathfrak J^n=0$, $\mathfrak R^n$ est contenu dans
$\operatorname{Ker}(u)$~; mais par récurrence sur $h$, il résulte de
l'hypothèse que $\mathfrak R^{2^h}$ est ouvert, donc aussi
$\mathfrak R^n$ pour tout $n$, et par suite $\operatorname{Ker}(u)$ est lui
aussi ouvert, ce qui prouve notre assertion.
\end{proof}

On notera que l'hypothèse sur $B$ signifie que la topologie de $B$ est borne
supérieure des topologies $\mathfrak R$-préadiques, où $\mathfrak R$ parcourt
l'ensemble des idéaux ouverts de $B$. Si $B$ est un anneau préadique
($\mathbf 0_{\rm I}$, 7.1.9), cette condition est toujours vérifiée.

\begin{env}[19.3.9]
\label{0.19.3.9-fr}
Nous allons maintenant voir que la propriété d'être formellement lisse
implique des propriétés de «~relèvement~» d'homomorphismes sous des conditions
plus générales que celles de la définition (19.3.1).
\end{env}

\oldpage[0\textsubscript{IV-1}]{84}

\begin{proposition}[19.3.10]
\label{0.19.3.10-fr}
Soient $A$ un anneau topologique, $B$ une $A$-algèbre formellement lisse.
Soient $C$ une $A$-algèbre topologique, $\mathfrak J$ un idéal de $C$,
vérifiant les conditions suivantes~:
\begin{enumerate}
  \item[$1^\circ$] $C$ est métrisable et complet.
  \item[$2^\circ$] $\mathfrak J$ est fermé et la suite
  $(\mathfrak J^n)$ tend vers $0$.
\end{enumerate}
Alors, tout $A$-homomorphisme continu $u:B\to C/\mathfrak J$ se factorise en
$B\xrightarrow{v}C\to C/\mathfrak J$, où $v$ est un
$A$-homomorphisme continu.
\end{proposition}

\begin{proof}
Soit $(\mathfrak Q_n)$ une suite décroissante d'idéaux de $C$ formant un
système fondamental de voisinages de $0$. Pour tout $n$, considérons la
$A$-algèbre discrète $C/\mathfrak Q_n$ et l'idéal
$(\mathfrak J+\mathfrak Q_n)/\mathfrak Q_n$ de cette algèbre~; comme il
existe $k$ tel que $\mathfrak J^k\subset\mathfrak Q_n$,
$(\mathfrak J+\mathfrak Q_n)/\mathfrak Q_n$ est nilpotent. Pour tout
$n$, soit $u_n$ l'homomorphisme continu
\[
  B\xrightarrow{u}C/\mathfrak J\to C/(\mathfrak J+\mathfrak Q_n)
  =(C/\mathfrak Q_n)/((\mathfrak J+\mathfrak Q_n)/\mathfrak Q_n)~;
\]
par hypothèse $u_n$ se factorise en
\[
  B\xrightarrow{v_n}C/\mathfrak Q_n
  \xrightarrow{\varphi_n}C/(\mathfrak J+\mathfrak Q_n),
\]
où $v_n$ est un $A$-homomorphisme continu. Montrons qu'on peut choisir de
proche en proche les $v_n$ de sorte que $v_n$ se factorise en
\[
  B\xrightarrow{v_{n+1}}C/\mathfrak Q_{n+1}\to C/\mathfrak Q_n
\]
pour tout $n$ (autrement dit, que les $v_n$ forment un \emph{système
projectif} d'homomorphismes). Cela résultera du lemme suivant.
\end{proof}

\begin{lemma}[19.3.10.1]
\label{0.19.3.10.1-fr}
Soient $B$ une $A$-algèbre formellement lisse, $E$, $E'$ deux $A$-algèbres
topologiques discrètes, $\mathfrak R$ (resp. $\mathfrak R'$) un idéal
nilpotent de $E$ (resp. $E'$), $f:E\to E'$ un $A$-homomorphisme surjectif tel
que $f(\mathfrak R)\subset\mathfrak R'$, $f':E/\mathfrak R\to
E'/\mathfrak R'$ l'homomorphisme déduit de $f$ par passage aux quotients.
Soient $u:B\to E/\mathfrak R$ un $A$-homomorphisme continu, $u'$
l'homomorphisme composé
$B\xrightarrow{u}E/\mathfrak R\xrightarrow{f'}E'/\mathfrak R'$, et soit
$v':B\to E'$ un $A$-homomorphisme continu tel que $u'$ se factorise en
$B\xrightarrow{v'}E'\to E'/\mathfrak R'$. Alors il existe un
$A$-homomorphisme continu $v:B\to E$ tel que $v'$ se factorise en
$B\xrightarrow{v}E\xrightarrow{f}E'$.
\end{lemma}

Dans le diagramme commutatif
\[
\xymatrix{
  E \ar[r]^f \ar[d] & E' \ar[d] \\
  E/\mathfrak R \ar[r]_{f'} & E'/\mathfrak R'
}
\]
tous les homomorphismes sont surjectifs~; si $\mathfrak L$ est le noyau de
$f$, $E'$ s'identifie à $E/\mathfrak L$ et $\mathfrak R'$ à
$(\mathfrak R+\mathfrak L)/\mathfrak L$. Utilisons maintenant le lemme
élémentaire suivant.

\begin{lemma}[19.3.10.2]
\label{0.19.3.10.2-fr}
Soient $F$ un anneau (non nécessairement commutatif), $\mathfrak a$,
$\mathfrak b$ deux idéaux bilatères de $F$~; alors le produit fibré
\[
  (F/\mathfrak a)\times_{F/(\mathfrak a+\mathfrak b)}(F/\mathfrak b)
\]
s'identifie canoniquement à $F/(\mathfrak a\cap\mathfrak b)$.
\end{lemma}

\begin{proof}
Ce n'est autre qu'une application particulière de (18.1.7), où, dans le
diagramme (18.1.7.1), on remplace $G$ par
$F/(\mathfrak a\cap\mathfrak b)$, $\mathfrak J'$ par
$\mathfrak b/(\mathfrak a\cap\mathfrak b)$, $F$ par
$F/\mathfrak b$, $\mathfrak R$ par
$(\mathfrak a+\mathfrak b)/\mathfrak b$, $B$ par
$F/(\mathfrak a+\mathfrak b)$ et $\theta$ par le $F$-homomorphisme
bijectif canonique
$(\mathfrak a+\mathfrak b)/\mathfrak b\to
\mathfrak a/(\mathfrak a\cap\mathfrak b)$.
\end{proof}

\oldpage[0\textsubscript{IV-1}]{85}

\begin{proof}[Suite de la démonstration du lemme~{\rm(19.3.10.1)}]
Appliquant ce lemme à la situation de (19.3.10.1), la commutativité du
diagramme
\[
\xymatrix{
  B \ar[r]^{v'} \ar[d]_u & E'=E/\mathfrak L \ar[d] \\
  E/\mathfrak R \ar[r]_{f'} & E'/\mathfrak R'
}
\]
montre l'existence d'un homomorphisme
$w:B\to E/(\mathfrak R\cap\mathfrak L)$ tel que $v'$ et $u$ se
factorisent en
$B\xrightarrow{w}E/(\mathfrak R\cap\mathfrak L)\to E/\mathfrak L$
et
$B\xrightarrow{w}E/(\mathfrak R\cap\mathfrak L)\to E/\mathfrak R$
respectivement~; comme le noyau de $w$ contient l'intersection de ceux de
$v'$ et de $u$, il est ouvert dans $B$ et $w$ est continu. Enfin, il est clair
que $\mathfrak R\cap\mathfrak L$ est nilpotent, donc $w$ se factorise en
$B\xrightarrow{v}E\to E/(\mathfrak R\cap\mathfrak L)$, où $v$ est
un $A$-homomorphisme continu répondant à la question.
\end{proof}

\begin{proof}[Suite de la démonstration de la proposition~{\rm(19.3.10)}]
L'existence de $v_{n+1}$ résulte du lemme (19.3.10.1) appliqué en remplaçant
$E$, $E'$, $\mathfrak R$, $\mathfrak R'$ par
$C/\mathfrak Q_{n+1}$, $C/\mathfrak Q_n$,
$(\mathfrak J+\mathfrak Q_{n+1})/\mathfrak Q_{n+1}$ et
$(\mathfrak J+\mathfrak Q_n)/\mathfrak Q_n$ respectivement, $u$, $u'$
et $v'$ par $u_{n+1}$, $u_n$ et $v_n$ respectivement. Notons que, comme $C$
est séparé et complet, il est égal à $\varprojlim_n(C/\mathfrak Q_n)$, et
$v=\varprojlim_n v_n$ est un $A$-homomorphisme continu de $B$ dans $C$.
Comme $\mathfrak J$ est fermé dans $C$, on a
$\mathfrak J=\bigcap_n(\mathfrak J+\mathfrak Q_n)$~; comme
$C/\mathfrak J$ est métrisable et complet et que les
$(\mathfrak J+\mathfrak Q_n)/\mathfrak J$ forment un système fondamental
de voisinages de $0$ dans $C/\mathfrak J$, on a
\[
  C/\mathfrak J=\varprojlim_n C/(\mathfrak J+\mathfrak Q_n)
\]
et $\varprojlim_n\varphi_n$ est l'application canonique
$\varphi:C\to C/\mathfrak J$. Comme $\varprojlim_n u_n=u$, on a bien
$u=v\circ\varphi$, C.Q.F.D.
\end{proof}

\begin{corollary}[19.3.11]
\label{0.19.3.11-fr}
Soient $A$ un anneau topologique, $B$ une $A$-algèbre formellement lisse,
$C$ un anneau local noethérien complet, $\mathfrak J$ un idéal distinct de
$C$, $\varphi:A\to C$ un homomorphisme continu, faisant de $C$ une
$A$-algèbre topologique. Alors tout $A$-homomorphisme continu
$u_0:B\to C_0=C/\mathfrak J$ se factorise en
$B\xrightarrow{u}C\to C/\mathfrak J$, où $u$ est un $A$-homomorphisme
continu.
\end{corollary}

Toutes les conditions de (19.3.10) sont en effet remplies
($\mathbf 0_{\rm I}$, 7.3.5).

\begin{proposition}[19.3.12]
\label{0.19.3.12-fr}
Soient $A$ un anneau topologique, $B$ une $A$-algèbre formellement lisse,
$C$ une $A$-algèbre topologique, $\mathfrak J$ un idéal de $C$, vérifiant les
conditions suivantes~:
\begin{enumerate}
  \item[$1^\circ$] Il existe un système fondamental d'idéaux ouverts
  $\mathfrak Q_\lambda$ de $C$, tels que les
  $C_\lambda=C/\mathfrak Q_\lambda$ soient des anneaux artiniens et que
  l'homomorphisme canonique $C\to\varprojlim_\lambda C_\lambda$ soit un
  isomorphisme d'anneaux topologiques.
  \item[$2^\circ$] L'idéal $\mathfrak J$ est fermé dans $C$ et
  topologiquement nilpotent.
  \item[$3^\circ$] Le carré de tout idéal ouvert de $B$ est ouvert.
\end{enumerate}
Dans ces conditions, tout $A$-homomorphisme continu
$u:B\to C/\mathfrak J$ se factorise en
$B\xrightarrow{v}C\to C/\mathfrak J$, où $v$ est un $A$-homomorphisme
continu.
\end{proposition}

\begin{proof}
Soit $\mathfrak J_\lambda=(\mathfrak J+\mathfrak Q_\lambda)/\mathfrak Q_\lambda$
l'image canonique de $\mathfrak J$ dans $C_\lambda$, qui est un idéal
nilpotent, et soit $u_\lambda$ l'homomorphisme composé
$B\xrightarrow{u}C/\mathfrak J\to C_\lambda/\mathfrak J_\lambda
=C/(\mathfrak J+\mathfrak Q_\lambda)$. Montrons que tout
$A$-homomorphisme $w:B\to C$ tel que $u$ se factorise en
$B\xrightarrow{w}C\to C/\mathfrak J$ est nécessairement continu~; en
effet, pour tout $\lambda$, il y a un idéal ouvert $\mathfrak R$ de $B$
tel que $u(\mathfrak R)\subset
(\mathfrak J+\mathfrak Q_\lambda)/\mathfrak J$, d'où
$w(\mathfrak R)\subset\mathfrak J+\mathfrak Q_\lambda$, donc il y a
une puissance $\mathfrak R^{2^m}$ de $\mathfrak R$ telle que
$w(\mathfrak R^{2^m})\subset\mathfrak Q_\lambda$, ce qui établit
notre assertion puisque $\mathfrak R^{2^m}$ est ouvert dans $B$. On peut
donc se borner à trouver un $A$-homomorphisme $w$ ayant la propriété précédente
sans s'inquiéter de ses propriétés de continuité. Or, l'ensemble
$\operatorname{Hom}(B,C)$ de tous les $A$-homomorphismes d'algèbres de $B$
dans $C$ est égal à
$\varprojlim_\lambda\operatorname{Hom}(B,C_\lambda)$, et ce dernier est
\emph{fermé} dans le $C_\lambda$-module $C_\lambda^B$, muni de la topologie
produit, pour laquelle il est \emph{linéairement compact} puisque
  $C_\lambda$ est artinien. Pour tout $\lambda$, soit $W_\lambda$

\oldpage[0\textsubscript{IV-1}]{86}

l'ensemble des $w_\lambda\in\operatorname{Hom}(B,C_\lambda)$ tels que
$u_\lambda$ se factorise en
$B\xrightarrow{w_\lambda}C_\lambda\to
C_\lambda/\mathfrak J_\lambda$~; $W_\lambda$ est une variété
linéaire fermée dans le $C_\lambda$-module
  $\operatorname{Hom}(B,C_\lambda)$, non vide puisque $B$ est formellement
  lisse. Dans le produit
$\prod_{\mu\leq\lambda}\operatorname{Hom}(B,C_\mu)=E_\lambda$,
considérons la sous-variété linéaire $U_\lambda$ formée des
$(w_\mu)_{\mu\leq\lambda}$ tels que
$w_\mu=\varphi_{\mu\lambda}\circ w_\lambda$ pour
$\mu\leq\lambda$ et $w_\lambda\in W_\lambda$, où
$\varphi_{\mu\lambda}:C_\lambda\to C_\mu$ est
l'homomorphisme canonique, qui est aussi fermée. Enfin, soit $V_\lambda$
la variété linéaire dans le produit
$\prod_\mu\operatorname{Hom}(B,C_\mu)$, produit de $U_\lambda$
et des $\operatorname{Hom}(B,C_\mu)$ pour
$\mu\not\leq\lambda$, qui est encore fermée. Tout revient à voir que
l'intersection des $V_\lambda$ est \emph{non vide}, car un élément $w$ de
cette intersection appartiendra à
$\varprojlim_\lambda\operatorname{Hom}(B,C_\lambda)$ par définition.
D'ailleurs, comme $\mathfrak J$ est fermé, et
$C=\varprojlim_\lambda C_\lambda$ linéairement compact,
$C/\mathfrak J$ s'identifie à
$\varprojlim_\lambda C_\lambda/\mathfrak J_\lambda$ et
l'application canonique $\psi:C\to C/\mathfrak J$ à
$\varprojlim_\lambda\psi_\lambda$, où $\psi_\lambda$ est
l'application canonique
$C_\lambda\to C_\lambda/\mathfrak J_\lambda$. On conclut donc
comme dans (19.3.10) que $\psi\circ w=u$.
\end{proof}

\subsection{Premiers critères de lissité formelle}
\label{subsection:0.19.4-fr}

\begin{proposition}[19.4.1]
\label{0.19.4.1-fr}
Soient $A$ un anneau topologique, $B$ une $A$-algèbre topologique~; supposons
qu'il existe deux familles filtrantes décroissantes
$(\mathfrak J_\alpha)_{\alpha\in I}$,
$(\mathfrak R_\alpha)_{\alpha\in I}$ d'idéaux de $A$ et $B$
respectivement, telles que~:
\begin{enumerate}
  \item[$1^\circ$] $(\mathfrak J_\alpha)$ tend vers $0$ dans $A$ et
  $(\mathfrak R_\alpha)$ tend vers $0$ dans $B$~;
  \item[$2^\circ$] pour tout $\alpha\in I$ on a
  $\mathfrak J_\alpha B\subset\mathfrak R_\alpha$ (de sorte que
  $B/\mathfrak R_\alpha$ est une
  $(A/\mathfrak J_\alpha)$-algèbre topologique)~;
  \item[$3^\circ$] pour tout $\alpha\in I$,
  $B/\mathfrak R_\alpha$ est une
  $(A/\mathfrak J_\alpha)$-algèbre formellement lisse.
\end{enumerate}
Alors $B$ est une $A$-algèbre formellement lisse.
\end{proposition}

\begin{proof}
En effet, soient $C$ une $A$-algèbre discrète, $\mathfrak L$ un idéal
nilpotent de $C$, $\mathfrak R$ un idéal ouvert de $B$~; par hypothèse il y a
un $\alpha\in I$ tel que
$\mathfrak R_\alpha\subset\mathfrak R$, donc $B/\mathfrak R$ est
quotient de $B/\mathfrak R_\alpha$ par un idéal ouvert. Tout
$A$-homomorphisme $u':B/\mathfrak R\to C/\mathfrak L$ est aussi un
$(A/\mathfrak J_\alpha)$-homomorphisme, donc il existe un idéal ouvert
$\mathfrak R'$ de $B$ tel que
$\mathfrak R_\alpha\subset\mathfrak R'\subset\mathfrak R$ et que
$B/\mathfrak R_\alpha\to B/\mathfrak R\xrightarrow{u'}C/\mathfrak L$
se factorise en
$B/\mathfrak R_\alpha\to B/\mathfrak R'\to C\to C/\mathfrak L$,
où le second homomorphisme est un
$(A/\mathfrak J_\alpha)$-homomorphisme~; d'où la conclusion, en vertu de
la remarque suivant (19.3.1).
\end{proof}

\begin{corollary}[19.4.2]
\label{0.19.4.2-fr}
Soient $A$ un anneau topologique, $B$ une $A$-algèbre topologique,
$(\mathfrak J_\alpha)_{\alpha\in I}$ une famille filtrante décroissante
d'idéaux de $A$ tendant vers $0$. Pour que $B$ soit une $A$-algèbre
formellement lisse, il faut et il suffit que pour tout $\alpha\in I$, si
l'on pose $A_\alpha=A/\mathfrak J_\alpha$,
$B_\alpha=B/\mathfrak J_\alpha B$, $B_\alpha$ soit une
$A_\alpha$-algèbre formellement lisse.
\end{corollary}

La condition est suffisante par (19.4.1), et elle est aussi nécessaire par
(19.3.5), (iii).

\begin{proposition}[19.4.3]
\label{0.19.4.3-fr}
Soient $A$ un anneau topologique, $B$ une $A$-algèbre topologique. Supposons
que pour toute $A$-algèbre discrète $C$ et tout idéal $\mathfrak J$ de $C$
tel que $\mathfrak J^2=0$, tout $A$-homomorphisme continu
$u:B\to C/\mathfrak J$ se factorise en
$B\xrightarrow{v}C\to C/\mathfrak J$, où $v$ est un $A$-homomorphisme
continu. Alors $B$ est une $A$-algèbre formellement lisse.
\end{proposition}

\begin{proof}
En effet, avec les mêmes notations, soit $\mathfrak R$ un idéal nilpotent
quelconque de $C$, et considérons un $A$-homomorphisme continu
$u':B\to C/\mathfrak R$. Supposons que $\mathfrak R^m=0$, et posons
$C_i=C/\mathfrak R^i$ pour $1\leq i\leq m$, de sorte que $C_m=C$, que
$\mathfrak R^{i-1}/\mathfrak R^i$ est un idéal de carré nul
$\mathfrak J_i$ dans $C_i$, et $C_i=C_{i+1}/\mathfrak J_{i+1}$~;
l'hypothèse entraîne alors par récurrence sur $i$ l'existence de
$A$-homomorphismes continus $u_i:B\to C_i$ tels que $u_1=u$ et que $u_i$
se factorise en
$B\xrightarrow{u_{i+1}}C_{i+1}\to C_{i+1}/\mathfrak J_{i+1}=C_i$~;
d'où la conclusion.
\end{proof}

\begin{proposition}[19.4.4]
\label{0.19.4.4-fr}
Soient $A$ un anneau topologique, $B$ une $A$-algèbre topologique
(commutative). Pour que $B$ soit une $A$-algèbre formellement lisse, il faut et
il suffit que pour

\oldpage[0\textsubscript{IV-1}]{87}

tout $B$-module topologique discret $L$, annulé par un idéal ouvert de $B$, on
ait (cf. (18.5.1))
\[
  \operatorname{Exalcotop}_A(B,L)=0.
\]
\end{proposition}

\begin{proof}
Soit $(\mathfrak R_\lambda)$ un système fondamental décroissant d'idéaux
ouverts de $B$ et posons $B_\lambda=B/\mathfrak R_\lambda$.
Considérons une $A$-algèbre topologique discrète $C$ et un idéal
$\mathfrak J$ de $C$ tel que $\mathfrak J^2=0$, de sorte que $C$ est une
$A$-extension de $C/\mathfrak J$ par $\mathfrak J$~; supposons donné un
$A$-homomorphisme $u:B_\lambda\to C/\mathfrak J$ et formons la
$A$-extension $E_\lambda$ de $B_\lambda$ par $\mathfrak J$, image
réciproque de $C$ par l'homomorphisme $u$ (18.2.5)~; c'est une $A$-algèbre
topologique pour la topologie discrète. Si
$\operatorname{Exalcotop}_A(B,L)=0$, il existe par définition (18.5.1) un
$\mu$ tel que $\mathfrak R_\mu\subset\mathfrak R_\lambda$ et que
l'extension image réciproque $E_\mu$ soit $A$-triviale~; mais cela signifie
(18.1.6) qu'il existe un $A$-homomorphisme continu
$v:B_\mu\to E_\lambda$ tel que l'homomorphisme canonique
$B_\mu\to B_\lambda$ se factorise en
$B_\mu\xrightarrow{v}E_\lambda\to B_\lambda$~; on en conclut
aussitôt que $B_\mu\to B_\lambda\xrightarrow{u}C/\mathfrak J$ se
factorise en
$B_\mu\xrightarrow{v}E_\lambda\to C\to C/\mathfrak J$, et cela
prouve, en vertu de (19.3.1) et (19.4.3), que $B$ est une $A$-algèbre
formellement lisse. La réciproque est immédiate, en appliquant (19.3.1) au cas
où $C$ est une $A$-algèbre topologique qui est une $A$-extension de
$B_\lambda$ par $L$, et à l'homomorphisme identique
$B_\lambda\to B_\lambda=C/L$.
\end{proof}

Lorsque $A$ et $B$ sont des anneaux discrets, le critère de lissité formelle
(19.4.4) se réduit à
\begin{equation}
\tag{19.4.4.1}
\label{0.19.4.4.1-fr}
\operatorname{Exalcom}_A(B,L)=0\qquad\text{pour tout $B$-module $L$};
\end{equation}
autrement dit, \emph{toute $A$-extension commutative de $B$ par un
$B$-module doit être $A$-triviale}.

\begin{corollary}[19.4.5]
\label{0.19.4.5-fr}
Soient $A$ un anneau topologique, $B$ une $A$-algèbre topologique
(commutative).
\begin{enumerate}
  \item[(i)] Supposons que $B$ soit une $A$-algèbre formellement lisse.
  Alors, pour tout idéal ouvert $\mathfrak R$ de $B$, tout
  $(B/\mathfrak R)$-module $L$ et tout $2$-cocycle de Hochschild
  $A$-bilinéaire et symétrique $f$ de
  $(B/\mathfrak R)\times(B/\mathfrak R)$ dans $L$ (18.4.3), il existe un
  idéal ouvert $\mathfrak R'\subset\mathfrak R$ tel que, si
  $\varphi:B/\mathfrak R'\to B/\mathfrak R$ est l'homomorphisme
  canonique, $f\circ(\varphi\times\varphi)$ soit un $2$-cobord de
  Hochschild $A$-bilinéaire de
  $(B/\mathfrak R')\times(B/\mathfrak R')$ dans $L$.
  \item[(ii)] Si $B$ est un $A$-module formellement projectif (19.2.1) et si
  la condition de (i) est vérifiée, $B$ est une $A$-algèbre formellement
  lisse.
\end{enumerate}
\end{corollary}

\begin{proof}
\begin{enumerate}
  \item[(i)] Le $2$-cocycle $f$ définit une $A$-extension de Hochschild $C$
  de $B/\mathfrak R$ par $L$ (18.4.3). Appliquant (19.3.1) à $C$, à l'idéal
  de carré nul $L$ de $C$ et à l'homomorphisme identique
  $B/\mathfrak R\to B/\mathfrak R=C/L$, on en déduit la condition (i) en
  vertu de (18.4.3).
  \item[(ii)] Appliquons le critère (19.4.3), en considérant un idéal ouvert
  $\mathfrak R$ de $B$, un idéal ouvert $\mathfrak J$ de $A$ tel que
  $\mathfrak JB\subset\mathfrak R$, et une $(A/\mathfrak J)$-extension $E$
  de $B/\mathfrak R$ par $L$. Comme $B$ est un $A$-module formellement
  projectif, l'application $A$-linéaire continue canonique
  $\rho:B\to B/\mathfrak R$ se factorise en
  $B\xrightarrow{w}E\to B/\mathfrak R$, où $w$ est une application
  $A$-linéaire continue (19.2.1), qui se factorise donc elle-même en
  $B\to B/\mathfrak R'\xrightarrow{w'}E$ où $\mathfrak R'$ est un idéal
  ouvert de $B$ contenu dans $\mathfrak R$~; remplaçant $\mathfrak J$ par
  un idéal plus petit, on peut supposer que
  $\mathfrak JB\subset\mathfrak R'$. Alors l'image réciproque par
  l'homomorphisme canonique $B/\mathfrak R'\to B/\mathfrak R$ de
  l'extension $E$ de $B/\mathfrak R$ par $L$, est une
  $(A/\mathfrak J)$-extension \emph{de Hochschild} $E'$ de
  $B/\mathfrak R'$ par $L$~;

  \oldpage[0\textsubscript{IV-1}]{88}

  appliquant (i) à un cocycle définissant cette extension (18.4.3), on en
  conclut qu'il y a un idéal ouvert
  $\mathfrak R''\subset\mathfrak R'$ de $B$ tel que l'image réciproque
  $E''$ de $E$ par $B/\mathfrak R''\to B/\mathfrak R$ soit
  $A$-triviale, C.Q.F.D.
\end{enumerate}
\end{proof}

\begin{corollary}[19.4.6]
\label{0.19.4.6-fr}
Soient $A$ un anneau topologique, $B$ une $A$-algèbre topologique qui soit un
$A$-module formellement projectif. Soit $A'$ une $A$-algèbre \emph{munie de
la topologie déduite de celle de $A$}. Supposons en outre que $A'$ soit un
$A$-module fidèlement plat, et que l'une des conditions suivantes soit
remplie~:
\begin{enumerate}
  \item[$1^\circ$] Il existe un système fondamental
  $(\mathfrak J_\lambda)$ d'idéaux ouverts de $A$ et un système fondamental
  $(\mathfrak M_\lambda)$ d'idéaux ouverts de $B$, ayant même ensemble
  d'indices et tels que, pour tout $\lambda$, on ait
  $\mathfrak J_\lambda B\subset\mathfrak M_\lambda$ et que
  $B/\mathfrak M_\lambda$ soit un
  $(A/\mathfrak J_\lambda)$-module projectif de type fini.
  \item[$2^\circ$] $A'$ est un $A$-module projectif de type fini.
\end{enumerate}
Alors, pour que $B'=B\otimes_A A'$ (muni de la topologie produit tensoriel)
soit une $A'$-algèbre formellement lisse, il faut et il suffit que $B$ soit une
$A$-algèbre formellement lisse.
\end{corollary}

\begin{proof}
La suffisance de la condition est contenue dans (19.3.5), (iii), sans aucune
hypothèse supplémentaire sur $B$ ni $A'$. Pour démontrer la réciproque, nous
allons appliquer le critère (19.4.5)~; sous l'hypothèse $2^\circ$, nous
désignerons encore par $(\mathfrak M_\lambda)$ un système fondamental
d'idéaux ouverts de $B$, et, pour tout $\lambda$, par
$\mathfrak J_\lambda$ un idéal ouvert de $A$ tel que
$\mathfrak J_\lambda B\subset\mathfrak M_\lambda$~; dans les deux
cas, nous poserons
$A_\lambda=A/\mathfrak J_\lambda$,
$B_\lambda=B/\mathfrak M_\lambda$,
$A'_\lambda=A_\lambda\otimes_A A'$,
$B'_\lambda=B_\lambda\otimes_A A'$. Soit $f$ un $2$-cocycle de
Hochschild $A_\lambda$-bilinéaire symétrique de
$B_\lambda\times B_\lambda$ dans un $B_\lambda$-module $L$~; par
extension des scalaires, on en déduit un $2$-cocycle de Hochschild
$f'=f\otimes1$, $A'_\lambda$-bilinéaire symétrique, de
$B'_\lambda\times B'_\lambda$ dans $L'=L\otimes_A A'$. Comme par
hypothèse $B'$ est une $A'$-algèbre formellement lisse, il existe, par (19.4.5),
(i), un indice $\mu$ tel que
$\mathfrak M_\mu\subset\mathfrak M_\lambda$, et tel que, si
$\varphi':B'_\mu\to B'_\lambda$ est l'application canonique,
$f'\circ(\varphi'\times\varphi')$ soit un $2$-cobord de Hochschild de
$B'_\mu\times B'_\mu$ dans $L'$~; autrement dit, son image $c'$ dans le
groupe de Hochschild
$\mathrm H^2_{A'_\mu}(B'_\mu,L')^s$ est nulle~; il est clair que si
$\varphi:B_\mu\to B_\lambda$ est l'homomorphisme canonique, et $c$ la
classe de $f\circ(\varphi\times\varphi)$ dans le groupe de Hochschild
$\mathrm H^2_{A_\mu}(B_\mu,L)^s$, $c'$ est l'image canonique de $c$.
Or, si $P_\bullet$ est le complexe relatif aux anneaux $A_\mu$ et
$B_\mu$ défini dans (18.4.5), servant au calcul de
$\mathrm H^2_{A_\mu}(B_\mu,L)^s$, le complexe analogue relatif aux
anneaux $A'_\mu$ et $B'_\mu$ est évidemment
$P_\bullet\otimes_A A'$~; dans l'hypothèse $1^\circ$, la construction de
$P_\bullet$ montre que c'est un $A_\mu$-module projectif de type fini.
On conclut donc de Bourbaki, \emph{Alg.}, chap.~II, 3\textsuperscript{e} éd., §~5,
n\textsuperscript{o}~3, prop.~7 que, sous les deux hypothèses, on a
\[
  \operatorname{Hom}_{A'_\mu}(P_\bullet\otimes_A A',L\otimes_A A')
  =\operatorname{Hom}_{A_\mu}(P_\bullet,L)\otimes_A A'
\]
à un isomorphisme canonique près~; comme $A'$ est un $A$-module plat, on a par
suite (18.4.5)
\[
  \mathrm H^2_{A'_\mu}(B'_\mu,L')^s
  =\mathrm H^2_{A_\mu}(B_\mu,L)^s\otimes_A A'
\]
et l'on peut donc écrire $c'=c\otimes1$~; mais comme $A'$ est un $A$-module
fidèlement plat, l'hypothèse $c'=0$ entraîne $c=0$, ce qui achève la
démonstration.
\end{proof}

\begin{proposition}[19.4.7]
\label{0.19.4.7-fr}
Soient $A$ un anneau topologique préadmissible, $\mathfrak J$ un idéal de
définition de $A$ ($\mathbf 0_{\rm I}$, 7.1.2), $B$ une $A$-algèbre
topologique qui soit un $A$-module formellement projectif. Considérons les
anneaux topologiques quotients $A_0=A/\mathfrak J$,
$B_0=B/\mathfrak JB$~; alors, pour que $B$ soit une $A$-algèbre formellement
lisse, il faut et il suffit que $B_0$ soit une $A_0$-algèbre formellement
lisse.
\end{proposition}

\oldpage[0\textsubscript{IV-1}]{89}

\begin{proof}
La nécessité de la condition résulte de (19.4.2). Pour voir qu'elle est
suffisante, notons d'abord qu'en considérant un système fondamental de
voisinages ouverts de $0$ dans $A$ formés d'idéaux
$\mathfrak J_\alpha$ contenus dans $\mathfrak J$, on peut, en vertu de
(19.4.2), se ramener au cas où $A$ est discret puisque
$B/\mathfrak J_\alpha B$ est un
$(A/\mathfrak J_\alpha)$-module formellement projectif (19.2.5)~; par
définition d'un anneau préadmissible, $\mathfrak J$ est alors
\emph{nilpotent}. Il suffit en outre de démontrer la proposition lorsque
$\mathfrak J^2=0$, car si $\mathfrak J^m=0$, on l'appliquera successivement
aux anneaux $A_k=A/\mathfrak J^k$ et $B_k=B/\mathfrak J^kB$ et à l'idéal
$\mathfrak J^{k-1}/\mathfrak J^k$ de $A_k$ pour $k=2,3,\ldots,m$, en
notant (19.2.5) que $B_k=B\otimes_A A_k$ est un $A_k$-module formellement
projectif. Appliquons le critère (19.4.5), (ii) en considérant un idéal ouvert
$\mathfrak R$ de $B$ et un $2$-cocycle de Hochschild $A$-bilinéaire
symétrique $f$ de $(B/\mathfrak R)\times(B/\mathfrak R)$ dans un
$(B/\mathfrak R)$-module $L$. Considérons d'abord le cas particulier où
$\mathfrak JL=0$, de sorte que $L$ peut aussi être considéré comme un
$(B_0/\mathfrak RB_0)$-module, et $f$ se factorise en
\[
  (B/\mathfrak R)\times(B/\mathfrak R)\to
  (B_0/\mathfrak RB_0)\times(B_0/\mathfrak RB_0)
  \xrightarrow{f_0}L,
\]
où $f_0$ est un $2$-cocycle de Hochschild bilinéaire symétrique. En vertu de
l'hypothèse, il y a donc un idéal ouvert
$\mathfrak R'\subset\mathfrak R$ dans $B$ et une application
$A_0$-linéaire $g_0:B_0/\mathfrak R'B_0\to L$ telle que
\[
  f_0(\varphi_0(x),\varphi_0(y))
  =xg_0(y)-g_0(xy)+g_0(x)y
\]
pour $x,y$ dans $B_0/\mathfrak R'B_0$, $\varphi_0$ étant
l'homomorphisme canonique
$B_0/\mathfrak R'B_0\to B_0/\mathfrak RB_0$. On en conclut aussitôt que
l'application $A$-linéaire composée
$g:B/\mathfrak R'\to B_0/\mathfrak R'B_0\xrightarrow{g_0}L$ est telle
que $dg=f\circ(\varphi\times\varphi)$, où
$\varphi:B/\mathfrak R'\to B/\mathfrak R$ est l'homomorphisme
canonique.

Passons maintenant au cas général, et considérons d'abord le
$(B/\mathfrak R)$-module $L/\mathfrak JL=L'$, pour lequel on a
$\mathfrak JL'=0$~; si $f'$ est l'application $A$-bilinéaire composée
$(B/\mathfrak R)\times(B/\mathfrak R)\xrightarrow{f}L\to L'$, $f'$ est
encore un $2$-cocycle de Hochschild symétrique, et en vertu de ce qui précède,
il existe un idéal ouvert $\mathfrak R'\subset\mathfrak R$ dans $B$ et
une application $A$-linéaire $g':B/\mathfrak R'\to L'$ vérifiant
$dg'=f'\circ(\varphi'\times\varphi')$ pour l'application canonique
$\varphi':B/\mathfrak R'\to B/\mathfrak R$. Comme $B$ est un
$A$-module formellement projectif, il existe un idéal ouvert
$\mathfrak R'_1\subset\mathfrak R'$ et une application $A$-linéaire
$g_1:B/\mathfrak R'_1\to L$ telle que l'homomorphisme
$B/\mathfrak R'_1\to B/\mathfrak R'\xrightarrow{g'}L'$ se factorise en
$B/\mathfrak R'_1\xrightarrow{g_1}L\to L'$. Considérons alors le
$2$-cocycle de Hochschild
\[
  f_1(x,y)=f(\varphi_1(x),\varphi_1(y))-xg_1(y)+g_1(xy)-g_1(x)y,
\]
application $A$-bilinéaire symétrique de
$(B/\mathfrak R'_1)\times(B/\mathfrak R'_1)$ dans $L$. Le choix de
$g_1$ entraîne que $f_1$ prend ses valeurs dans $\mathfrak JL$. Comme
$\mathfrak J(\mathfrak JL)=0$, on peut de nouveau appliquer le premier cas,
et il y a donc un idéal ouvert
$\mathfrak R''\subset\mathfrak R'_1$ et une application $A$-linéaire
$g_2:B/\mathfrak R''\to\mathfrak JL$ telle que
\[
  f_1(\varphi_2(x),\varphi_2(y))=xg_2(y)-g_2(xy)+g_2(x)y
\]
dans $(B/\mathfrak R'')\times(B/\mathfrak R'')$,
$\varphi_2:B/\mathfrak R''\to B/\mathfrak R'_1$ étant l'application
canonique~; l'application $A$-linéaire
$g=g_2+g_1\circ\varphi_2:B/\mathfrak R''\to L$ vérifie donc
$dg=f\circ(\varphi\times\varphi)$ pour l'application canonique
$\varphi:B/\mathfrak R''\to B/\mathfrak R$. C.Q.F.D.
\end{proof}

\subsection{Lissité formelle et anneaux gradués associés}
\label{subsection:0.19.5-fr}

\begin{env}[19.5.1]
\label{0.19.5.1-fr}
\oldpage[0\textsubscript{IV-1}]{90}
Soient $C$ un anneau topologique (commutatif), $V$ un $C$-module topologique,
et considérons l'algèbre symétrique
$\mathbf S_C^*(V)=\bigoplus_{n\geq0}\mathbf S_C^n(V)$, que nous allons munir
canoniquement d'une \emph{topologie linéaire}, compatible avec sa structure de
$C$-algèbre. Pour cela, considérons un sous-module ouvert $U$ de $V$, et
l'\emph{idéal gradué} $U.\mathbf S_C^*(V)$ qu'il engendre dans
$\mathbf S_C^*(V)$~; nous prendrons comme système fondamental de voisinages de
$0$ dans $\mathbf S_C^*(V)$ l'ensemble des sommes
$\mathfrak a.\mathbf S_C^*(V)+U.\mathbf S_C^*(V)$, où $\mathfrak a$ (resp.
$U$) parcourt un système fondamental d'idéaux ouverts (resp. de sous-modules
ouverts) de $C$ (resp. $V$). On notera que si la topologie de $V$ est moins
fine que la topologie déduite de celle de $C$, on peut se borner aux couples
$(\mathfrak a,U)$ tels que $\mathfrak aV\subset U$, donc
$\mathfrak a.\mathbf S_C^*(V)+U.\mathbf S_C^*(V)
=\mathfrak a.1_C+U.\mathbf S_C^*(V)$~; dans ce cas la topologie induite sur
chacun des $\mathbf S_C^n(V)$ pour $n\geq1$ a pour système fondamental de
voisinages de $0$ les $U.\mathbf S_C^{n-1}(V)$, où $U$ parcourt les
sous-modules ouverts de $V$~; en particulier, sur $V=\mathbf S_C^1(V)$, elle
coïncide avec la topologie donnée (en général elle est moins fine que cette
dernière). En outre, dans tous les cas, l'\emph{algèbre topologique}
$\mathbf S_C^*(V)$ ainsi définie est, pour les catégories des $C$-modules
topologiques et des $C$-algèbres topologiques, la solution du même problème
universel que pour les catégories de $C$-modules et de $C$-algèbres. En effet,
soit $E$ une $C$-algèbre topologique, $u:V\to E$ un homomorphisme de
$C$-modules, $\mathbf S(u)=\widetilde u$ son extension canonique à
$\mathbf S_C^*(V)$. Supposons $u$ continu~; alors, si $\mathfrak b$ est un
idéal ouvert de $E$, son image réciproque $U=u^{-1}(\mathfrak b)$ est un
sous-$C$-module ouvert de $V$, et l'image par $\widetilde u$ de
$U.\mathbf S_C^*(V)$ est donc contenue dans $\mathfrak b$~; comme par ailleurs
il existe un idéal ouvert $\mathfrak a$ de $C$ tel que
$\mathfrak a.E\subset\mathfrak b$, donc
$\widetilde u(\mathfrak a.\mathbf S_C^*(V))\subset\mathfrak b$, cela prouve
que $\widetilde u$ est continu. Réciproquement, si $\widetilde u$ est continu,
et si $\mathfrak b$ est un idéal ouvert de $E$, il existe un sous-module ouvert
$U$ de $V$ tel que $\widetilde u(U.\mathbf S_C^*(V))\subset\mathfrak b$, et en
particulier $\widetilde u(U.1_C)\subset\mathfrak b$, c'est-à-dire
$u(U)\subset\mathfrak b$, donc $u$ est continu. Rappelons en outre que l'on a
un isomorphisme canonique de $C$-modules topologiques (discrets)
\begin{equation}
\label{0.19.5.1.1-fr}
  \mathbf S_C^n(V)/U.\mathbf S_C^{n-1}(V)
  \xrightarrow{\sim}\mathbf S_C^n(V/U).
  \tag{19.5.1.1}
\end{equation}
\end{env}

\begin{env}[19.5.2]
\label{0.19.5.2-fr}
Soient $A$ un anneau topologique, $B$ une $A$-algèbre topologique
(commutative), $\mathfrak J$ un idéal de $B$ (non nécessairement ouvert ni
fermé)~; dans tout ce qui suit, on munit $\mathfrak J$ et les
$\mathfrak J^n$ ($n\geq2$) de la topologie \emph{induite} par celle de $B$,
les quotients $C=B/\mathfrak J$,
$\mathfrak J^n/\mathfrak J^{n+1}=\gr_{\mathfrak J}^n(B)$ de la topologie
\emph{quotient}, de sorte que les $\mathfrak J^n/\mathfrak J^{n+1}$ sont des
$C$-modules topologiques~; l'injection canonique
$\mathfrak J/\mathfrak J^2\to\gr_{\mathfrak J}^*(B)$ se prolonge en un
homomorphisme (d'abord non topologique) de $C$-algèbres
\[
  \varphi:\mathbf S_C^*(\mathfrak J/\mathfrak J^2)
  \longrightarrow\gr_{\mathfrak J}^*(B)
\]
qui pour chaque $n$ donne un homomorphisme \emph{surjectif} de $C$-modules
\begin{equation}
\label{0.19.5.2.1-fr}
  \varphi_n:\mathbf S_C^n(\mathfrak J/\mathfrak J^2)
  \longrightarrow\mathfrak J^n/\mathfrak J^{n+1}
  =\gr_{\mathfrak J}^n(B).
  \tag{19.5.2.1}
\end{equation}

\oldpage[0\textsubscript{IV-1}]{91}

Lorsque $\mathbf S_C^n(\mathfrak J/\mathfrak J^2)$ est muni de la topologie
définie dans (19.5.1), les homomorphismes $\varphi_n$ sont continus, en vertu
de la propriété universelle (19.5.1) de
$\mathbf S_C^*(\mathfrak J/\mathfrak J^2)$ appliquée à la $C$-algèbre
topologique $E=\gr_{\mathfrak J/\mathfrak J^n}(B/\mathfrak J^n)$ munie de la
topologie produit de celles des $\mathfrak J^k/\mathfrak J^{k+1}$, et à
l'injection canonique $\mathfrak J/\mathfrak J^2\to E$. Notons qu'ici la
topologie de $\mathfrak J/\mathfrak J^2$ est moins fine que la topologie
déduite de celle de $C$ (cette dernière topologie sur
$\mathfrak J/\mathfrak J^2$ étant aussi la topologie déduite de celle de $B$).
\end{env}

\begin{theorem}[19.5.3]
\label{0.19.5.3-fr}
Soient $A$ un anneau topologique, $B$ une $A$-algèbre topologique,
$\mathfrak J$ un idéal de $B$, $C=B/\mathfrak J$ la $A$-algèbre topologique
quotient. On suppose que la $A$-algèbre discrète $C$ est formellement lisse.
Alors~:
\begin{enumerate}
  \item[\rm(i)] Si $B$ est une $A$-algèbre formellement lisse,
  $\mathfrak J/\mathfrak J^2$ est un $C$-module topologique formellement
  projectif (19.2.1).
  \item[\rm(ii)] Si $B$ est une $A$-algèbre formellement lisse et un anneau
  préadmissible ($\mathbf 0_{\rm I}$, 7.1.2), les homomorphismes $\varphi_n$
  (19.5.2) sont des bimorphismes formels (19.1.2).
  \item[\rm(iii)] Réciproquement, supposons que $B$ soit préadmissible, que la
  suite $(\mathfrak J^n)$ tende vers $0$ dans $B$, que
  $\mathfrak J/\mathfrak J^2$ soit un $C$-module formellement projectif et que
  les $\varphi_n$ soient des bimorphismes formels. Alors $B$ est une
  $A$-algèbre formellement lisse.
\end{enumerate}
\end{theorem}

La démonstration de ce théorème est longue et encombrée de détails techniques~;
aussi commencerons-nous par en démontrer un corollaire plus simple, où
apparaissent plus nettement les idées directrices~; ce corollaire est d'ailleurs
le cas particulier du théorème (19.5.3) qui sera le plus fréquemment employé
par la suite.

\begin{corollary}[19.5.4]
\label{0.19.5.4-fr}
Soient $A$ un anneau topologique, $B$ une $A$-algèbre topologique,
$\mathfrak J$ un idéal de $B$ tel que la topologie de $B$ soit la topologie
$\mathfrak J$-préadique. On suppose que la $A$-algèbre discrète
$C=B/\mathfrak J$ est formellement lisse. Alors les trois conditions suivantes
sont équivalentes~:
\begin{enumerate}
  \item[a)] $B$ est une $A$-algèbre formellement lisse.
  \item[b)] $\mathfrak J/\mathfrak J^2$ est un $C$-module projectif et
  l'homomorphisme canonique
  \[
    \varphi:\mathbf S_C^*(\mathfrak J/\mathfrak J^2)
    \longrightarrow\gr_{\mathfrak J}^*(B)
  \]
  est bijectif.
  \item[c)] Le séparé complété $\widehat B$ de $B$ est une
  $\widehat A$-algèbre topologique isomorphe à une $\widehat A$-algèbre de la
  forme $\widehat{B'}$, où $B'=\mathbf S_C^*(V)$, $V$ étant un $C$-module
  projectif et $B'$ étant muni de la topologie ${B'}^+$-préadique, où
  ${B'}^+$ est l'idéal d'augmentation.
\end{enumerate}
\end{corollary}

\begin{env}[19.5.4.1]
\label{0.19.5.4.1-fr}
Prouvons d'abord que $a)$ entraîne que $\mathfrak J/\mathfrak J^2$ est un
$C$-module projectif. Soient donc $P$ et $Q$ deux $C$-modules,
$u:P\to Q$ un $C$-homomorphisme surjectif, et
$v:\mathfrak J/\mathfrak J^2\to Q$ un $C$-homomorphisme. L'anneau
$B/\mathfrak J^2$ étant considéré comme une $B$-extension de $C$ par
$\mathfrak J/\mathfrak J^2$, on en déduit par $v$ une $B$-extension
$E=(B/\mathfrak J^2)\oplus_{(\mathfrak J/\mathfrak J^2)}Q$ de $C$ par $Q$
(18.2.8). Comme $C$ est une $A$-algèbre discrète formellement lisse,
l'extension $E$ est $A$-\emph{triviale} (19.4.4) et peut donc s'identifier à
$\mathrm D_C(Q)$ (18.2.3). L'homomorphisme surjectif $u:P\to Q$ définit alors
canoniquement un $A$-homomorphisme surjectif
$\psi:\mathrm D_C(P)\to\mathrm D_C(Q)$ (18.2.9) dont le noyau est un idéal
$\mathfrak L$ de l'extension $E'=\mathrm D_C(P)$, contenu dans $P$, et
\emph{a fortiori} de carré nul. Soit $f:B\to E=E'/\mathfrak L$
l'homomorphisme définissant la structure de $B$-algèbre de $E$~: comme
$\mathfrak L$ est de carré nul, et que $B$ est une $A$-algèbre formellement
lisse, $f$ se factorise en $B\xrightarrow{g}E'\to E'/\mathfrak L$, où $g$ est
un $A$-homomorphisme continu.

\oldpage[0\textsubscript{IV-1}]{92}

Le diagramme
\[
  \xymatrix{
    E' \ar[r]^{\psi} & E \\
    B \ar[u]^{g} \ar[r]_{\theta} & B/\mathfrak J^2 \ar[u]_{v'}
  }
\]
où $v'$ est déduit de $v$ (18.2.8), est commutatif. D'ailleurs l'image de
$\mathfrak J$ par $\psi\circ g$ est contenue dans $Q$, donc l'image de
$\mathfrak J$ par $g$ est contenue dans $P+\mathfrak L=P$, et l'image de
$\mathfrak J^2$ par $g$ est nulle. Restreignant $g$ et $\theta$ à
$\mathfrak J$, on obtient un diagramme commutatif
\[
  \xymatrix{
    P \ar[r]^{u} & Q \\
    \mathfrak J/\mathfrak J^2 \ar[u]^{w} \ar[ur]_{v}
  }
\]
ce qui prouve que le $C$-module $\mathfrak J/\mathfrak J^2$ est projectif.
\end{env}

\begin{env}[19.5.4.2]
\label{0.19.5.4.2-fr}
Prouvons en second lieu que $a)$ entraîne que $\varphi$ est bijectif, ce qui
achèvera de montrer que $a)$ entraîne $b)$. Posons
$E_n=B/\mathfrak J^{n+1}$, et désignons par $F_n$ le quotient de
$\mathbf S_C^*(\mathfrak J/\mathfrak J^2)$ par la puissance $(n+1)$-ième de
son idéal d'augmentation. Comme $\mathfrak J/\mathfrak J^{n+1}$ est nilpotent
dans $E_n$ et que $C$ est une $A$-algèbre discrète formellement lisse,
l'automorphisme identique de $C$ se factorise en
$C\xrightarrow{f}E_n\to C=E_n/(\mathfrak J/\mathfrak J^{n+1})$ où $f$ est un
$A$-homomorphisme~; d'autre part, $\mathfrak J/\mathfrak J^2$ étant un
$C$-module projectif d'après (19.5.4.1), l'automorphisme identique de
$\mathfrak J/\mathfrak J^2$ se factorise en
$\mathfrak J/\mathfrak J^2\xrightarrow{g}\mathfrak J/\mathfrak J^{n+1}
\to\mathfrak J/\mathfrak J^2$, où $g$ est une application $C$-linéaire~; à
partir de $f$ et $g$ on obtient canoniquement un homomorphisme de $C$-algèbres
$\mathbf S_C^*(\mathfrak J/\mathfrak J^2)\to E_n$, qui d'ailleurs (par
définition de $g$) s'annule sur la puissance $(n+1)$-ième de l'idéal
d'augmentation de $\mathbf S_C^*(\mathfrak J/\mathfrak J^2)$, d'où, en
passant au quotient, un $A$-homomorphisme surjectif d'algèbres
$v:F_n\to E_n$ tel que $\gr^0(v)$ et $\gr^1(v)$ soient les automorphismes
identiques de $C$ et de $\mathfrak J/\mathfrak J^2$. Par définition de
l'homomorphisme canonique $\varphi$, on voit que $\gr^j(v)=\varphi_j$ pour tout
$j\leq n$. Notons maintenant que le noyau $\mathfrak R$ de $v$ est un idéal
nilpotent de $F_n$, de sorte que $E_n$ s'identifie à $F_n/\mathfrak R$. Comme
$B$ est une $A$-algèbre formellement lisse, le $A$-homomorphisme canonique
$p_n:B\to E_n=B/\mathfrak J^{n+1}$, qui est continu, se factorise en
$B\xrightarrow{w}F_n\xrightarrow{v}E_n$, où $w$ est un $A$-homomorphisme
continu~; comme $\gr^0(v)$ est l'identité, $w(\mathfrak J)$ est contenu dans
l'idéal d'augmentation de $F_n$, d'où $w(\mathfrak J^{n+1})=0$, si bien que
$w$ se factorise en
$B\xrightarrow{p_n}B/\mathfrak J^n=E_n\xrightarrow{w'}F_n$, et le composé
$E_n\xrightarrow{w'}F_n\xrightarrow{v}E_n$ est l'identité. En outre, comme
$\gr^0(v)$ et $\gr^1(v)$ sont les automorphismes identiques de $C$ et de
$\mathfrak J/\mathfrak J^2$, il en est de même de $\gr^0(w')$ et de
$\gr^1(w')$. Comme $\gr^*(E_n)$ est engendré par $\gr^1(E_n)$,
l'homomorphisme composé
\[
  \gr^j(F_n)\xrightarrow{\varphi_j}\gr^j(E_n)
  \xrightarrow{\gr^j(w')}\gr^j(F_n)
\]
est l'identité pour tout $j\leq n$, puisque c'est vrai pour $j=0$ et $j=1$~;
faisant $j=n$, on prouve ainsi que $\varphi_n$ est injectif, ce qui achève de
montrer que $a)$ entraîne $b)$.
\end{env}

\begin{env}[19.5.4.3]
\label{0.19.5.4.3-fr}
\oldpage[0\textsubscript{IV-1}]{93}
Prouvons ensuite que $b)$ entraîne $a)$. Le même raisonnement qu'au début de
(19.5.4.2) prouve l'existence d'un $A$-homomorphisme surjectif d'algèbres
$v:F_n\to E_n$ tel que $\gr^j(v)=\varphi_j$ pour tout $j$~; comme $\varphi_j$
est bijectif pour tout $j$ et que les filtrations de $F_n$ et $E_n$ sont
finies, on en conclut que $v$ est \emph{bijectif} (Bourbaki,
\emph{Alg. comm.}, chap.~III, §~2, n\textsuperscript{o}~8, cor.~3 du th.~1).
Soient maintenant $G$ une $A$-algèbre topologique discrète, $\mathfrak R$ un
idéal de carré nul dans $G$, $f:B\to G/\mathfrak R$ un $A$-homomorphisme
continu d'algèbres. Comme $G$ est discrète, il existe un entier $m$ tel que
$f$ s'annule dans $\mathfrak J^m$, donc $f$ se factorise en
$B\to E_n\xrightarrow{f_n}G/\mathfrak R$, où l'on prend $n=2m$. On a donc par
composition un $A$-homomorphisme continu d'algèbres
$r:C\to F_n\xrightarrow{v}E_n\xrightarrow{f_n}G/\mathfrak R$, et comme $C$
est une $A$-algèbre formellement lisse discrète, $r$ se factorise en
$C\xrightarrow{r'}G\to G/\mathfrak R$, de sorte que $G$ est muni par $r'$ d'une
structure de $C$-algèbre. D'autre part, la restriction à
$\mathfrak J/\mathfrak J^2$ de l'homomorphisme
$f'_n:F_n\xrightarrow{v}E_n\xrightarrow{f_n}G/\mathfrak R$ est une application
$C$-linéaire $g:\mathfrak J/\mathfrak J^2\to G/\mathfrak R$. Puisque
$\mathfrak J/\mathfrak J^2$ est un $C$-module projectif, $g$ se factorise en
$\mathfrak J/\mathfrak J^2\xrightarrow{h}G\to G/\mathfrak R$, où $h$ est une
application $C$-linéaire~; par prolongement, on en déduit un homomorphisme de
$C$-algèbres $w:\mathbf S_C^*(\mathfrak J/\mathfrak J^2)\to G$, et par
construction tout élément de degré $m$ a pour image par $w$ un élément de
$\mathfrak R$, donc tout élément de degré $n=2m$ a une image nulle puisque
$\mathfrak R^2=0$~; autrement dit, $w$ se factorise en
$\mathbf S_C^*(\mathfrak J/\mathfrak J^2)\to F_n\xrightarrow{w'}G$. Par
construction, le composé $F_n\xrightarrow{w'}G\xrightarrow{\theta}G/\mathfrak R$
coïncide avec $f'_n$ dans $C$ et dans $\mathfrak J/\mathfrak J^2$, donc est
égal à $f'_n$. Finalement, le composé
$B\to E_n\xrightarrow{v^{-1}}F_n\xrightarrow{w'}G\to G/\mathfrak R$ étant
égal à $f$, on voit que $B$ est une $A$-algèbre formellement lisse.
\end{env}

\begin{env}[19.5.4.4]
\label{0.19.5.4.4-fr}
Il reste à prouver l'équivalence de $a)$ et $c)$. En premier lieu, $c)$ entraîne
$a)$~: en effet, $B'$ est une $C$-algèbre formellement lisse pour la topologie
discrète (19.3.2), donc aussi pour la topologie ${B'}^+$-préadique (19.3.8)~;
comme $C$ est une $A$-algèbre formellement lisse, $B'$ est aussi une
$A$-algèbre formellement lisse (19.3.5, (ii)) et finalement $\widehat{B'}$ est
une $\widehat A$-algèbre formellement lisse (19.3.6), donc $B$ est une
$A$-algèbre formellement lisse (19.3.6). Reste à voir que $a)$ entraîne $c)$.
Notons d'abord que puisque $C$ est une $A$-algèbre formellement lisse,
l'homomorphisme identique $C\to B/\mathfrak J$ se factorise en
$C\to\widehat B\to B/\mathfrak J$, où $C\to\widehat B$ est un
$A$-homomorphisme (19.3.10)~; $\widehat B$, et par suite tous les
$B/\mathfrak J^{n+1}$, sont donc munis de structures de $C$-algèbres. D'autre
part, comme $\mathfrak J/\mathfrak J^2$ est un $C$-module projectif par $b)$,
l'injection canonique $\mathfrak J/\mathfrak J^2\to B/\mathfrak J^2$ permet de
former un système projectif de $C$-homomorphismes
$u_n:\mathfrak J/\mathfrak J^2\to B/\mathfrak J^{n+1}$ pour $n\geq1$, donc
aussi (par la propriété universelle de l'algèbre symétrique) un système
projectif d'homomorphismes de $C$-algèbres
$v_n:\mathbf S_C^*(\mathfrak J/\mathfrak J^2)=B'\to B/\mathfrak J^{n+1}$~;
d'ailleurs il est clair que $v_n$ s'annule dans $({B'}^+)^{n+1}$ et comme
$\gr^0(v_n)=\varphi_0$ et $\gr^1(v_n)=\varphi_1$, on a
$\gr^j(v_n)=\varphi_j$ pour $0\leq j\leq n$. Comme les $\varphi_j$ sont des
isomorphismes par $b)$, et que les filtrations de
$B'/({B'}^+)^{n+1}$ et de $B/\mathfrak J^{n+1}$ sont finies, on en conclut que
$v_n$ est \emph{bijectif} pour

\oldpage[0\textsubscript{IV-1}]{94}

tout $n$ (Bourbaki, \emph{Alg. comm.}, chap.~III, §~2,
n\textsuperscript{o}~8, cor.~3 du th.~1)~; d'où $c)$ par passage à la limite
projective.
\end{env}

\begin{remark}[19.5.5]
\label{0.19.5.5-fr}
Sous les hypothèses générales de (19.5.4), supposons en outre que
$\mathfrak J$ soit un idéal \emph{maximal} de $B$, de sorte que
$C=B/\mathfrak J$ soit un corps, et que $\mathfrak J/\mathfrak J^2$ soit un
$C$-espace vectoriel de dimension \emph{finie}. Alors les conditions $a)$,
$b)$ et $c)$ de (19.5.4) sont encore équivalentes à la suivante~:
\begin{enumerate}
  \item[d)] \emph{Étant donnés une algèbre de polynômes}
  $E=C[T_1,\ldots,T_n]$ \emph{sur le corps} $C$, \emph{et en désignant par}
  $\mathfrak n$ \emph{l'idéal maximal de} $E$ \emph{engendré par les} $T_i$
  $(1\leq i\leq n)$, \emph{alors, pour tout idéal} $\mathfrak b$ \emph{de}
  $E$ \emph{contenant une puissance} $\mathfrak n^k$, \emph{tout
  homomorphisme} $v:B\to E/\mathfrak b$ \emph{de} $A$-\emph{algèbres}
  $C$-\emph{augmentées se factorise en}
  $B\xrightarrow{w}E/\mathfrak n^k\to E/\mathfrak b$, \emph{où} $w$ \emph{est
  un} $A$-\emph{homomorphisme}.
\end{enumerate}
En effet, $\mathfrak J/\mathfrak J^2$ étant ici un $C$-module libre, il suffit
de vérifier que la condition $d)$ entraîne que l'homomorphisme canonique
$\varphi:\mathbf S_C^*(\mathfrak J/\mathfrak J^2)\to\gr_{\mathfrak J}^*(B)$
est bijectif. Or ici $\mathbf S_C^*(\mathfrak J/\mathfrak J^2)$ s'identifie à
l'algèbre de polynômes $E=C[T_1,\ldots,T_n]$, où
$n=\operatorname{rg}_C(\mathfrak J/\mathfrak J^2)$, et l'idéal d'augmentation
de $E$ est l'idéal $\mathfrak n$ engendré par les $T_i$. Pour tout entier
$k\geq0$, l'hypothèse que $C$ est une $A$-algèbre discrète formellement lisse
entraîne, comme dans (19.5.4.2), l'existence d'un $A$-homomorphisme surjectif
$f:E/\mathfrak n^{k+1}\to B/\mathfrak J^{k+1}$, tel que
$\gr^j(f)=\varphi_j$ pour tout $j\leq k$. Si
$\mathfrak b/\mathfrak n^{k+1}=\operatorname{Ker}(f)$,
$B/\mathfrak J^{k+1}$ s'identifie donc à $E/\mathfrak b$, et l'hypothèse $d)$
permet de factoriser l'homomorphisme canonique $B\to B/\mathfrak J^{k+1}$ en
$B\xrightarrow{w}E/\mathfrak n^{k+1}\xrightarrow{f}B/\mathfrak J^{k+1}$~;
comme $\gr^0(f)$ est l'identité, on a
$w(\mathfrak J)\subset\mathfrak n/\mathfrak n^{k+1}$, donc
$w(\mathfrak J^{k+1})=0$~; on conclut comme dans (19.5.4.2) que $\varphi_k$
est injectif pour tout $k$, ce qui achève de prouver notre assertion.
\end{remark}

\begin{env}[19.5.6]
\label{0.19.5.6-fr}
\emph{Démonstration du théorème} (19.5.3). --- Soit
$(\mathfrak b_\alpha)$ un système fondamental décroissant d'idéaux ouverts de
$B$. Nous poserons
\[
  B_\alpha=B/\mathfrak b_\alpha,\qquad
  C_\alpha=B/(\mathfrak b_\alpha+\mathfrak J),\qquad
  \mathfrak J_\alpha=(\mathfrak b_\alpha+\mathfrak J)/\mathfrak b_\alpha,
\]
de sorte que
\[
  C_\alpha=B_\alpha/\mathfrak J_\alpha\qquad\text{et}\qquad
  \mathfrak J_\alpha^n/\mathfrak J_\alpha^{n+1}
  =(\mathfrak b_\alpha+\mathfrak J^n)/(\mathfrak b_\alpha+\mathfrak J^{n+1}).
\]
Les $C$-modules
$((\mathfrak b_\alpha\cap\mathfrak J^n)+\mathfrak J^{n+1})/
\mathfrak J^{n+1}$ de $\mathfrak J^n/\mathfrak J^{n+1}$ forment un système
fondamental de sous-modules ouverts du $C$-module topologique
$\mathfrak J^n/\mathfrak J^{n+1}$~; comme
$(\mathfrak b_\alpha\cap\mathfrak J^n)+\mathfrak J^{n+1}
=\mathfrak J^n\cap(\mathfrak b_\alpha+\mathfrak J^{n+1})$, on a
\[
  \mathfrak J^n/((\mathfrak b_\alpha\cap\mathfrak J^n)+\mathfrak J^{n+1})
  =\mathfrak J^n/(\mathfrak J^n\cap(\mathfrak b_\alpha+\mathfrak J^{n+1}))
  =(\mathfrak b_\alpha+\mathfrak J^n)/(\mathfrak b_\alpha+\mathfrak J^{n+1})
  =\mathfrak J_\alpha^n/\mathfrak J_\alpha^{n+1}.
\]
\end{env}

\begin{env}[19.5.6.1]
\label{0.19.5.6.1-fr}
Soient $P$, $Q$ deux $C_\alpha$-modules discrets,
$u:P\to Q$ un $C_\alpha$-homomorphisme surjectif. L'anneau discret
$B/(\mathfrak b_\alpha+\mathfrak J^2)$ est une $B$-extension de $C_\alpha$ par
l'idéal de carré nul $\mathfrak J_\alpha/\mathfrak J_\alpha^2$. Soit
$v:\mathfrak J/\mathfrak J^2\to Q$ un $C$-homomorphisme continu~; remplaçant
au besoin $\mathfrak b_\alpha$ par un idéal ouvert plus petit de $B$, on peut
supposer que le noyau de $v$ contient le sous-$C$-module ouvert
$((\mathfrak b_\alpha\cap\mathfrak J)+\mathfrak J^2)/\mathfrak J^2$~;
passant au quotient, on en déduit un $C_\alpha$-homomorphisme de modules
discrets $v':\mathfrak J_\alpha/\mathfrak J_\alpha^2\to Q$~; soit
$E_\alpha$ la $B_\alpha$-extension de $C_\alpha$ par $Q$ déduite de
$B/(\mathfrak b_\alpha+\mathfrak J^2)$ au moyen de $v'$, et soit
$v'_*:B/(\mathfrak b_\alpha+\mathfrak J^2)\to E_\alpha$ le
$B_\alpha$-homomorphisme correspondant (18.2.8)~; $E_\alpha$ est une
$A$-algèbre topologique discrète, et l'isomorphisme canonique
$C_\alpha\xrightarrow{\sim}E_\alpha/Q$ donne par composition un
$A$-homomorphisme continu $f:C\to C_\alpha\to E_\alpha/Q$. Mais comme $Q$ est
de carré nul dans $E_\alpha$ et que $C$ est une $A$-algèbre formellement
lisse, $f$ se factorise en $C\xrightarrow{g}E_\alpha\to E_\alpha/Q$ où $g$
est un $A$-homomorphisme continu. Comme $g$ est continu et $E_\alpha$ discret,
il existe un $\beta\geq\alpha$ tel que $g$ se factorise en
$C\to C_\beta\xrightarrow{g'}E_\alpha$. D'autre part, soit $E_\beta$ la
$A$-extension de $C_\beta$ par $Q$, image réciproque de $E_\alpha$ par
l'homomorphisme canonique $C_\beta\to C_\alpha$~; l'existence de $g'$ (tel que
le composé $C_\beta\to E_\alpha\to C_\alpha$ soit l'homomorphisme canonique)
équivaut au fait que $E_\beta$ est une extension $A$-\emph{triviale}, et peut
donc s'identifier à $\mathrm D_{C_\beta}(Q)$. Cela étant, l'homomorphisme
surjectif $u:P\to Q$ définit canoniquement un homomorphisme surjectif
$\mathrm D_{C_\beta}(P)\to\mathrm D_{C_\beta}(Q)$ (18.2.9), dont le noyau est
un idéal $\mathfrak L$ de l'extension $E'_\beta=\mathrm D_{C_\beta}(P)$
contenu dans $P$, et \emph{a fortiori} de carré nul. Soit
$f':B\to E_\beta=E'_\beta/\mathfrak L$ le $A$-homomorphisme continu définissant
la structure de $B$-algèbre topologique de $E_\beta$~; comme $\mathfrak L$ est
de carré nul et que $B$ est une $A$-algèbre formellement lisse, $f'$ se
factorise en $B\xrightarrow{h'}E'_\beta\to E'_\beta/\mathfrak L$, où $h'$ est
un $A$-homomorphisme continu. La construction de $E'_\beta$ donne par
composition un $A$-homomorphisme $\psi:E'_\beta\to E_\beta\to E_\alpha$, et
il est clair que le diagramme
\[
  \xymatrix{
    E'_\beta \ar[r]^{\psi} & E_\alpha \\
    B \ar[u]^{h'} \ar[r]_{\theta} &
    B/(\mathfrak b_\alpha+\mathfrak J^2) \ar[u]_{v'_*}
  }
\]
est commutatif. D'ailleurs, l'image par $\psi\circ h'$ de $\mathfrak J$ est
contenue dans $Q$, donc l'image par $h'$ de $\mathfrak J$ est contenue dans
$P$, et l'image de $\mathfrak J^2$ par $h'$ est nulle. Restreignant $h'$ et
$\theta$ à $\mathfrak J$, on obtient un diagramme commutatif
\[
  \xymatrix{
    P \ar[r]^{u} & Q \\
    \mathfrak J/\mathfrak J^2 \ar[u]^{w} \ar[ur]_{v}
  }
\]
où $w$ est continu, ce qui prouve que
$\mathfrak J/\mathfrak J^2$ est un $C$-module formellement projectif.
\end{env}

\oldpage[0\textsubscript{IV-1}]{95}

\begin{env}[19.5.6.2]
\label{0.19.5.6.2-fr}
Pour tout entier $n\geq0$, nous poserons
\[
  E_n=B/\mathfrak J^{n+1},
  \qquad\text{de sorte que}\qquad E_0=C~;
\]
les idéaux $(\mathfrak b_\alpha+\mathfrak J^{n+1})/\mathfrak J^{n+1}$ forment
un système fondamental d'idéaux ouverts dans $E_n$, et nous poserons
\[
  E_{\alpha,n}=B/(\mathfrak b_\alpha+\mathfrak J^{n+1})
  =B_\alpha/\mathfrak J_\alpha^{n+1},
\]
anneau quotient discret. De même, dans
$\gr_{\mathfrak J}^n(B)=\mathfrak J^n/\mathfrak J^{n+1}$, nous avons vu que les
$((\mathfrak b_\alpha\cap\mathfrak J^n)+\mathfrak J^{n+1})/
\mathfrak J^{n+1}$ forment un système fondamental de voisinages de $0$, le
quotient de $\mathfrak J^n/\mathfrak J^{n+1}$ par ce sous-module s'identifiant
canoniquement à $\mathfrak J_\alpha^n/\mathfrak J_\alpha^{n+1}$. Considérons
l'algèbre symétrique
$\mathbf S_{C_\alpha}^*(\mathfrak J_\alpha/\mathfrak J_\alpha^2)$~; nous
noterons $F_{\alpha,n}$ le quotient de
$\mathbf S_{C_\alpha}^*(\mathfrak J_\alpha/\mathfrak J_\alpha^2)$ par la
puissance $(n+1)$-ième de son idéal d'augmentation. Pour un $n\geq1$ fixé, il
résulte de (19.5.1.1) que les
$\mathbf S_{C_\alpha}^n(\mathfrak J_\alpha/\mathfrak J_\alpha^2)$ sont les
quotients de $\mathbf S_C^n(\mathfrak J/\mathfrak J^2)$ par un système
fondamental de sous-modules ouverts dans ce $C$-module topologique.

Pour abréger le langage, nous dirons que pour $\alpha\leq\beta$, les
homomorphismes canoniques $B_\beta\to B_\alpha$, $C_\beta\to C_\alpha$,
$\mathfrak J_\beta/\mathfrak J_\beta^2\to
\mathfrak J_\alpha/\mathfrak J_\alpha^2$, $E_{\beta,n}\to E_{\alpha,n}$,
$F_{\beta,n}\to F_{\alpha,n}$, etc., sont les \emph{homomorphismes de
transition}.
\end{env}

\begin{lemma}[19.5.6.3]
\label{0.19.5.6.3-fr}
Supposons que la $A$-algèbre $C$ soit formellement lisse, l'anneau $B$
préadmissible et le $C$-module $\mathfrak J/\mathfrak J^2$ formellement
projectif. Alors~:
\begin{enumerate}
  \item[\rm(i)] Pour tout $\alpha$, il existe $\beta\geq\alpha$ et un
  $A$-homomorphisme surjectif d'algèbres
  \[
    v_{\alpha\beta}:F_{\beta,n}\longrightarrow E_{\alpha,n}
  \]
  tel que $\gr^0(v_{\alpha\beta}):C_\beta\to C_\alpha$ et
  $\gr^1(v_{\alpha\beta}):\mathfrak J_\beta/\mathfrak J_\beta^2
  \to\mathfrak J_\alpha/\mathfrak J_\alpha^2$ soient les homomorphismes de
  transition.
  \item[\rm(ii)] Si $\beta$ et $v_{\alpha\beta}$ vérifient les conditions de
  \emph{(i)}, alors, pour tout $\gamma\geq\beta$, il existe
  $\delta\geq\gamma$ et un $A$-homomorphisme surjectif d'algèbres
  $v_{\gamma\delta}:F_{\delta,n}\to E_{\gamma,n}$ vérifiant les conditions de
  \emph{(i)} (pour $\gamma$ et $\delta$) et rendant commutatif le diagramme
  \begin{equation}
  \label{0.19.5.6.4-fr}
  \xymatrix@C=3em{
    F_{\beta,n} \ar[r]^{v_{\alpha\beta}} & E_{\alpha,n} \\
    F_{\delta,n} \ar[u] \ar[r]_{v_{\gamma\delta}} & E_{\gamma,n} \ar[u]
  }
  \tag{19.5.6.4}
  \end{equation}
  où les flèches verticales sont les homomorphismes de transition.
\end{enumerate}
\end{lemma}

\begin{env}[19.5.6.4]
\label{0.19.5.6.4-proof-fr}
\begin{enumerate}
  \item[\rm(i)] Dans la $A$-algèbre topologique discrète $E_{\alpha,n}$,
  l'idéal $\mathfrak J_\alpha/\mathfrak J_\alpha^{n+1}$ est nilpotent, et
  l'isomorphisme identique
  $C_\alpha\xrightarrow{\sim}
  E_{\alpha,n}/(\mathfrak J_\alpha/\mathfrak J_\alpha^{n+1})$ donne par
  composition un $A$-homomorphisme continu
  $C\to C_\alpha\to
  E_{\alpha,n}/(\mathfrak J_\alpha/\mathfrak J_\alpha^{n+1})$~;

\oldpage[0\textsubscript{IV-1}]{96}

  comme $C$ est une $A$-algèbre formellement lisse, cet homomorphisme se
  factorise en
  $C\xrightarrow{f_\alpha}E_{\alpha,n}\to
  E_{\alpha,n}/(\mathfrak J_\alpha/\mathfrak J_\alpha^{n+1})$, où $f_\alpha$
  est un $A$-homomorphisme continu~; par suite,
  $\mathfrak J_\alpha/\mathfrak J_\alpha^{n+1}$ devient, au moyen de
  $f_\alpha$, un $C$-module topologique discret annulé par un idéal ouvert de
  $C$. L'hypothèse que $\mathfrak J/\mathfrak J^2$ est un $C$-module
  formellement projectif entraîne alors l'existence d'une application
  $C$-linéaire continue
  $g_\alpha:\mathfrak J/\mathfrak J^2\to
  \mathfrak J_\alpha/\mathfrak J_\alpha^{n+1}$ rendant commutatif le diagramme
  \[
    \xymatrix{
      \mathfrak J_\alpha/\mathfrak J_\alpha^{n+1} \ar[r] &
      \mathfrak J_\alpha/\mathfrak J_\alpha^2 \\
      & \mathfrak J/\mathfrak J^2 \ar[ul]^{g_\alpha} \ar[u]
    }
  \]
  Comme $f_\alpha$ et $g_\alpha$ sont continus, il y a un
  $\beta\geq\alpha$ tel que ces homomorphismes se factorisent respectivement en
  \[
    C\to C_\beta\xrightarrow{f'_{\alpha\beta}}E_{\alpha,n},
  \]
  \[
    \mathfrak J/\mathfrak J^2\to
    \mathfrak J_\beta/\mathfrak J_\beta^2
    \xrightarrow{g'_{\alpha\beta}}
    \mathfrak J_\alpha/\mathfrak J_\alpha^{n+1}
  \]
  et à partir de $f'_{\alpha\beta}$ et $g'_{\alpha\beta}$, on obtient donc
  canoniquement un homomorphisme de $C_\beta$-algèbres
  \[
    \mathbf S_{C_\beta}^*(\mathfrak J_\beta/\mathfrak J_\beta^2)
    \longrightarrow E_{\alpha,n}
  \]
  qui d'ailleurs (par définition de $g'_{\alpha\beta}$) s'annule dans la
  puissance $(n+1)$-ième de l'idéal d'augmentation de
  $\mathbf S_{C_\beta}^*(\mathfrak J_\beta/\mathfrak J_\beta^2)$~; passant au
  quotient par cette puissance $(n+1)$-ième, on obtient l'homomorphisme
  $v_{\alpha\beta}$ cherché, compte tenu des définitions de $f_\alpha$ et
  $g_\alpha$~; la surjectivité de $v_{\alpha\beta}$ résulte en effet de celle
  des deux homomorphismes $\gr^0(v_{\alpha\beta})$ et
  $\gr^1(v_{\alpha\beta})$, puisque cela entraîne que
  $\gr(v_{\alpha\beta})$ est surjectif (l'algèbre
  $\gr^*(E_{\alpha,n})$ étant engendrée par $\gr^0(E_{\alpha,n})$ et
  $\gr^1(E_{\alpha,n})$), et comme les filtrations considérées sont finies, on
  peut appliquer Bourbaki, \emph{Alg. comm.}, chap.~III, §~2,
  n\textsuperscript{o}~8, cor.~3 du th.~1.

  \item[\rm(ii)] L'hypothèse que $B$ est \emph{préadmissible} signifie que
  l'on peut supposer tous les $\mathfrak b_\alpha$ contenus dans un même
  $\mathfrak b_{\alpha_0}$ dont les puissances \emph{tendent vers} $0$. Cela
  entraîne en particulier que le noyau de l'homomorphisme de transition
  $E_{\gamma,n}\to E_{\alpha,n}$, égal à
  $(\mathfrak b_\alpha+\mathfrak J^{n+1})/
  (\mathfrak b_\gamma+\mathfrak J^{n+1})$, est \emph{nilpotent}~; appliquant le
  lemme (19.3.10.1), on voit qu'on peut supposer $f_\gamma$ choisi de sorte que
  le diagramme
  \[
    \xymatrix{
      & E_{\alpha,n} \\
      C \ar[ur]^{f_\alpha} \ar[r]_{f_\gamma} & E_{\gamma,n} \ar[u]
    }
  \]
  (où la flèche verticale est l'homomorphisme de transition) soit commutatif.
  L'hypothèse que $\mathfrak J/\mathfrak J^2$ est un $C$-module formellement
  projectif permet d'autre part de choisir
  $g_\gamma:\mathfrak J/\mathfrak J^2\to
  \mathfrak J_\gamma/\mathfrak J_\gamma^{n+1}$ de sorte que le diagramme
  \[
    \xymatrix@C=5em@R=3em{
      \mathfrak J_\alpha/\mathfrak J_\alpha^{n+1} \ar[r] &
      \mathfrak J_\alpha/\mathfrak J_\alpha^2 \\
      \mathfrak J_\gamma/\mathfrak J_\gamma^{n+1} \ar[u] \ar[r] &
      \mathfrak J_\gamma/\mathfrak J_\gamma^2 \ar[u] \\
      & \mathfrak J/\mathfrak J^2
        \ar[ul]_{g_\gamma} \ar[u]_{p_\gamma} \ar[uul]^{g_\alpha}
    }
  \]
  soit commutatif (il suffit de remarquer que l'image par $(g_\alpha,p_\gamma)$
  de $\mathfrak J/\mathfrak J^2$ dans le module produit
  $(\mathfrak J_\alpha/\mathfrak J_\alpha^{n+1})
  \times(\mathfrak J_\gamma/\mathfrak J_\gamma^2)$ est contenue (vu la
  définition de $g_\alpha$ et de la relation $\beta\leq\gamma$) dans l'image
  canonique $Q$ dans ce module produit du module
  $P=\mathfrak J_\gamma/\mathfrak J_\gamma^{n+1}$, et d'appliquer la définition
  (19.2.1) à l'homomorphisme \emph{surjectif} $P\to Q$ et à l'homomorphisme
  $(g_\alpha,p_\gamma)$ de $\mathfrak J/\mathfrak J^2$ dans $Q$). Ce choix de
  $f_\gamma$ et de $g_\gamma$ permet alors, en construisant
  $v_{\gamma\delta}$ comme dans \emph{(i)}, d'obtenir en outre la commutativité
  du diagramme (19.5.6.4).
\end{enumerate}
\end{env}

\oldpage[0\textsubscript{IV-1}]{97}

\begin{lemma}[19.5.6.5]
\label{0.19.5.6.5-fr}
Supposons que les $A$-algèbres $B$ et $C$ soient formellement lisses et que
$B$ soit préadmissible (de sorte qu'en vertu de (19.5.6.1), les conditions de
(19.5.6.3) sont satisfaites). Pour tout système de deux indices
$\alpha\leq\beta$ et d'un homomorphisme $v_{\alpha\beta}$ vérifiant les
conditions de (19.5.6.3, (i)), il existe un indice $\lambda\geq\beta$ et un
$A$-homomorphisme d'algèbres
\[
  w_{\beta\lambda}:E_{\lambda,n}\longrightarrow F_{\beta,n}
\]
tel que~: $1^\circ$ $\gr^0(w_{\beta\lambda}):C_\lambda\to C_\beta$ et
$\gr^1(w_{\beta\lambda}):
\mathfrak J_\lambda/\mathfrak J_\lambda^2\to
\mathfrak J_\beta/\mathfrak J_\beta^2$ soient les homomorphismes de
transition~; $2^\circ$ le composé
$E_{\lambda,n}\xrightarrow{w_{\beta\lambda}}F_{\beta,n}
\xrightarrow{v_{\alpha\beta}}E_{\alpha,n}$ soit l'homomorphisme de transition.
\end{lemma}

\begin{proof}
Appliquons le lemme (19.5.6.3, (ii)) avec $\gamma=\beta$, ce qui donne un
$\delta\geq\beta$ et un
$v_{\beta\delta}:F_{\delta,n}\to E_{\beta,n}$. Rappelons que
$v_{\beta\delta}$ est surjectif~; d'autre part, son noyau $\mathfrak R$ est
\emph{nilpotent}~: en effet, l'idéal d'augmentation $F_{\delta,n}^+$ de
$F_{\delta,n}$ est nilpotent par définition, et le noyau de
$\gr^0(v_{\beta\delta}):C_\delta\to C_\beta$ est aussi nilpotent en vertu de
l'hypothèse que $B$ est préadmissible. On voit donc que $E_{\beta,n}$
s'identifie à $F_{\delta,n}/\mathfrak R$. Comme $B$ est une $A$-algèbre
formellement lisse, le $A$-homomorphisme canonique
$p_{\beta,n}:B\to E_{\beta,n}$, qui est continu, se factorise en
$B\xrightarrow{w}F_{\delta,n}\xrightarrow{v_{\beta\delta}}E_{\beta,n}$, où
$w$ est un $A$-homomorphisme continu. Comme $F_{\delta,n}$ est discret, il
existe $\lambda\geq\delta$ tel que $w$ soit nul dans $\mathfrak b_\lambda$,
donc $w$ se factorise en
$B\to B/\mathfrak b_\lambda\xrightarrow{w'_\lambda}F_{\delta,n}$. Considérons
l'homomorphisme composé
$B/\mathfrak b_\lambda\xrightarrow{w'_\lambda}F_{\delta,n}
\xrightarrow{q_{\beta\delta}}F_{\beta,n}$, où $q_{\beta\delta}$ est
l'homomorphisme de transition~; notons que
$\gr^0(q_{\beta\delta})=\gr^0(v_{\beta\delta}):C_\delta\to C_\beta$ est
l'homomorphisme de transition~; comme le composé
$v_{\alpha\beta}\circ q_{\beta\delta}\circ w'_\lambda$ est l'homomorphisme
canonique $B/\mathfrak b_\lambda\to
B/(\mathfrak b_\alpha+\mathfrak J^{n+1})$, cela montre que l'image de
$\mathfrak J_\lambda=(\mathfrak b_\lambda+\mathfrak J)/\mathfrak b_\lambda$
par $q_{\beta\delta}\circ w'_\lambda$ est contenue dans l'idéal d'augmentation
$F_{\beta,n}^+$, et par suite l'image de
$\mathfrak J_\lambda^{n+1}=(\mathfrak b_\lambda+\mathfrak J^{n+1})/
\mathfrak b_\lambda$ est nulle. Autrement dit,
$q_{\beta\delta}\circ w'_\lambda$ se factorise en
\[
  B/\mathfrak b_\lambda\to
  B/(\mathfrak b_\lambda+\mathfrak J^{n+1})=E_{\lambda,n}
  \xrightarrow{w_{\beta\lambda}}F_{\beta,n}
\]
tel que le composé
$E_{\lambda,n}\xrightarrow{w_{\beta\lambda}}F_{\beta,n}
\xrightarrow{v_{\alpha\beta}}E_{\alpha,n}$ soit l'homomorphisme de transition.
Le raisonnement précédent montre aussi que $\gr^0(w_{\beta\lambda})$, qui est
le composé
$B/(\mathfrak b_\lambda+\mathfrak J)\to\gr^0(F_{\delta,n})
\xrightarrow{\gr^0(q_{\beta\delta})}\gr^0(F_{\beta,n})$, est l'homomorphisme
de transition $C_\lambda\to C_\beta$, puisque
$v_{\beta\delta}\circ w$ est l'homomorphisme canonique. En outre, on a aussi
$\gr^1(q_{\beta\delta})=\gr^1(v_{\beta\delta})$ (19.5.6.3), donc le même
raisonnement prouve que $\gr^1(w_{\beta\lambda})$ est l'homomorphisme de
transition.
\end{proof}

\begin{env}[19.5.6.6]
\label{0.19.5.6.6-fr}
\emph{Démonstration de} (19.5.3, (ii)). --- Pour montrer que $\varphi_n$ est un
bimorphisme formel, nous allons appliquer le critère de (19.1.3, (iii)). Les
conditions du lemme (19.5.6.5) étant satisfaites par hypothèse, déterminons,
pour tout indice $\alpha$, $v_{\alpha\beta}$ et $w_{\beta\lambda}$ vérifiant
les conclusions de ce lemme. L'homomorphisme
\[
  \gr^n(v_{\alpha\beta}):\gr^n(F_{\beta,n})
  \longrightarrow\gr^n(E_{\alpha,n})
\]
n'est autre que l'homomorphisme
\[
  \varphi_{\alpha\beta,n}:
  \mathbf S_{C_\beta}^n(\mathfrak J_\beta/\mathfrak J_\beta^2)
  \longrightarrow\mathfrak J_\alpha^n/\mathfrak J_\alpha^{n+1}
\]
déduit de l'homomorphisme canonique $\varphi_n$ (19.5.2.1) par passage aux
quotients~; en effet, il résulte de (19.5.6.3) que $\gr^0(v_{\alpha\beta})$ et
$\gr^1(v_{\alpha\beta})$ coïncident respectivement avec
$\varphi_{\alpha\beta,0}$ et $\varphi_{\alpha\beta,1}$, et la définition de
l'homomorphisme canonique $\varphi$ montre alors, par récurrence sur $j\leq n$,
que $\gr^j(v_{\alpha\beta})$ et $\varphi_{\alpha\beta,j}$ sont égaux pour tout
$j$. Cela étant, comme $\varphi_n$ est surjectif, c'est \emph{a fortiori} un
épimorphisme formel~; en outre, pour $j\leq n$, l'homomorphisme composé
\[
  \gr^j(F_{\lambda,n})\xrightarrow{\varphi_{\lambda\lambda,j}}
  \gr^j(E_{\lambda,n})\xrightarrow{\gr^j(w_{\beta\lambda})}
  \gr^j(F_{\beta,n})
\]
est l'homomorphisme de transition, car c'est vrai pour $j=0$ et $j=1$ en
vertu de (19.5.6.5), et comme $\gr^*(E_{\lambda,n})$ est engendré par
$\gr^1(E_{\lambda,n})$ cela démontre l'assertion par récurrence sur $j$.
Composant avec l'homomorphisme de transition
$\gr^j(F_{\beta,n})\to\gr^j(F_{\alpha,n})$, on voit donc, pour $j=n$, que
l'on a factorisé l'homomorphisme de transition
$\gr^n(F_{\lambda,n})\to\gr^n(F_{\alpha,n})$ en
\[
  \gr^n(F_{\lambda,n})\xrightarrow{\varphi_{\lambda\lambda,n}}
  \gr^n(E_{\lambda,n})\longrightarrow\gr^n(F_{\alpha,n}),
\]
ce qui est la condition du critère (19.1.3) pour que $\varphi_n$ soit un
bimorphisme formel.
\end{env}

\oldpage[0\textsubscript{IV-1}]{98}

\begin{lemma}[19.5.6.7]
\label{0.19.5.6.7-fr}
Supposons que $B$ soit préadmissible, que $\mathfrak J/\mathfrak J^2$ soit un
$C$-module formellement projectif, que $C$ soit une $A$-algèbre formellement
lisse et que les $\varphi_n$ soient des bimorphismes formels. Alors, pour tout
couple d'indices $\alpha$, $\beta$ tel que $\alpha\leq\beta$ et tout
homomorphisme $v_{\alpha\beta}:F_{\beta,n}\to E_{\alpha,n}$ vérifiant les
conditions de (19.5.6.3, (i)), il existe un indice $\gamma\geq\beta$ tel que,
pour tout indice $\mu\geq\gamma$, il y ait un diagramme commutatif de
$A$-homomorphismes
\[
\xymatrix@C=4em@R=3em{
  F_{\beta,n}\ar[r]^{v_{\alpha\beta}} & E_{\alpha,n} \\
  & E_{\mu,n}\ar[ul]_{u_{\beta\mu}}\ar[u]_{p_{\alpha\mu}}
}
\]
où $p_{\alpha\mu}$ est l'homomorphisme de transition.
\end{lemma}

\begin{proof}
Appliquant le critère (19.1.3) à chacun des $\varphi_j$ pour $0\leq j\leq n$,
on voit qu'il existe un indice $\gamma\geq\beta$ et des homomorphismes
uniquement déterminés (et surjectifs) de $C_\gamma$-modules
\[
  w_{\beta\gamma}^j:\gr^j(E_{\gamma,n})\longrightarrow
  \gr^j(F_{\beta,n})
\]
tels que les composés
\[
  \gr^j(F_{\gamma,n})
  \xrightarrow{\varphi_{\gamma\gamma,j}}\gr^j(E_{\gamma,n})
  \xrightarrow{w_{\beta\gamma}^j}\gr^j(F_{\beta,n})
\]
soient les homomorphismes de transition (le fait qu'on puisse choisir le même
indice $\gamma$ pour tous les $\varphi_j$ résulte de ce qu'ils sont en nombre
fini). En outre, l'unicité des $w_{\beta\gamma}^j$ prouve (puisque $\varphi$ est
un homomorphisme d'algèbres graduées) que
$w_{\beta\gamma}=(w_{\beta\gamma}^j)_{0\leq j\leq n}$ est un homomorphisme de
$C_\gamma$-algèbres graduées
\[
  \gr^*(E_{\gamma,n})\longrightarrow
  \gr^*(F_{\beta,n})=F_{\beta,n}.
\]
En outre, comme $\varphi_{\gamma\gamma,0}$ et $\varphi_{\gamma\gamma,1}$ sont
les homomorphismes identiques,
$w_{\beta\gamma}^0:C_\gamma\to C_\beta$ et
$w_{\beta\gamma}^1:\mathfrak J_\gamma/\mathfrak J_\gamma^2
\to\mathfrak J_\beta/\mathfrak J_\beta^2$ sont les homomorphismes de
transition, et il en est donc de même de
$\gr^0(v_{\alpha\beta}\circ w_{\beta\gamma})$ et de
$\gr^1(v_{\alpha\beta}\circ w_{\beta\gamma})$~; comme
$\gr^*(E_{\gamma,n})$ est engendré par $\gr^1(E_{\gamma,n})$, on en conclut que
$\gr^j(v_{\alpha\beta}\circ w_{\beta\gamma})$ est aussi l'homomorphisme de
transition pour $0\leq j\leq n$.

Appliquons maintenant à $\alpha$, $\beta$ et $\gamma$ le lemme (19.5.6.3,
(ii)), ce qui donne le diagramme (19.5.6.4) avec $\delta\geq\gamma$~; puis
répétons le raisonnement du début en remplaçant $\alpha$ et $\beta$ par
$\gamma$ et $\delta$. On obtient ainsi un indice $\lambda\geq\delta$ et un
diagramme commutatif d'homomorphismes
\[
\xymatrix@C=4em@R=4em{
  \gr^*(E_{\gamma,n})\ar[r]^{w_{\beta\gamma}} &
  F_{\beta,n}\ar[r]^{v_{\alpha\beta}} & E_{\alpha,n} \\
  \gr^*(E_{\lambda,n})\ar[u]_{\gr^*(p_{\gamma\lambda})}
    \ar[r]_{w_{\delta\lambda}} &
  F_{\delta,n}\ar[u]_{q_{\beta\delta}}\ar[r]_{v_{\gamma\delta}} &
  E_{\gamma,n}\ar[u]_{p_{\alpha\gamma}}
    \ar@{-->}[ul]_{u_{\beta\gamma}}
}
\]
où les flèches verticales sont les homomorphismes de transition. Tout revient
à prouver l'existence de l'homomorphisme $u_{\beta\gamma}$ laissant le
diagramme commutatif, et pour cela il suffit évidemment de montrer que l'on a
\[
  \operatorname{Ker}(v_{\gamma\delta})
  \subset\operatorname{Ker}(q_{\beta\delta}),
\]
$v_{\gamma\delta}$ et $q_{\beta\delta}$ étant surjectifs. Comme
$w_{\delta\lambda}$ est surjectif, cette dernière relation équivaut à
\[
\begin{split}
  \operatorname{Ker}(v_{\gamma\delta}\circ w_{\delta\lambda})
  &\subset\operatorname{Ker}(q_{\beta\delta}\circ w_{\delta\lambda})\\
  &=\operatorname{Ker}\bigl(w_{\beta\gamma}\circ
  (\gr(p_{\gamma\lambda}))\bigr).
\end{split}
\]
Mais on a vu ci-dessus que
$\gr(p_{\gamma\lambda})=\gr(v_{\gamma\delta}\circ w_{\delta\lambda})$ et il
est clair que
\[
\begin{split}
  \operatorname{Ker}(v_{\gamma\delta}\circ w_{\delta\lambda})
  &\subset\operatorname{Ker}\bigl(\gr(v_{\gamma\delta}\circ
  w_{\delta\lambda})\bigr)\\
  &=\operatorname{Ker}(\gr(p_{\gamma\lambda}))
  \subset\operatorname{Ker}\bigl(w_{\beta\gamma}\circ
  (\gr(p_{\gamma\lambda}))\bigr),
\end{split}
\]
ce qui achève la démonstration du lemme, car pour tout $\mu\geq\gamma$, il
suffira de définir $u_{\beta\mu}$ comme le composé
\[
  E_{\mu,n}\xrightarrow{p_{\gamma\mu}}E_{\gamma,n}
  \xrightarrow{u_{\beta\gamma}}F_{\beta,n}.
\]
\end{proof}

\begin{env}[19.5.6.8]
\label{0.19.5.6.8-fr}
\emph{Démonstration de} (19.5.3, (iii)). --- Soient $G$ une $A$-algèbre
topologique discrète, $\mathfrak R$ un idéal de carré nul dans $G$,
$f:B\to G/\mathfrak R$ un $A$-homomorphisme continu d'algèbres. Comme
$G/\mathfrak R$ est discrète, il y a un indice $\alpha$ tel que $f$ s'annule
dans $\mathfrak b_\alpha$~; par hypothèse, il existe un entier $m$ tel que
$\mathfrak J^m\subset\mathfrak b_\alpha$, donc $f$ se factorise en
\[
  B\xrightarrow{p_\alpha}E_{\alpha,n}
  \xrightarrow{f_{\alpha,n}}G/\mathfrak R,
\]
où l'on prend $n=2m$. Les hypothèses du lemme (19.5.6.3) étant satisfaites, on
a tout d'abord un $\beta\geq\alpha$ et un $A$-homomorphisme composé
\begin{equation}
\label{0.19.5.6.9-fr}
  f_{\alpha,n}\circ v_{\alpha\beta}:F_{\beta,n}
  \xrightarrow{v_{\alpha\beta}}E_{\alpha,n}
  \xrightarrow{f_{\alpha,n}}G/\mathfrak R.
  \tag{19.5.6.9}
\end{equation}
Comme $F_{\beta,n}$ est une $C_\beta$-algèbre, donc \emph{a fortiori} une
$C$-algèbre topologique (discrète), cela donne par composition un
$A$-homomorphisme continu
$r:C\to F_{\beta,n}\to G/\mathfrak R$. Comme d'autre part $C$ est une
$A$-algèbre formellement lisse,

\oldpage[0\textsubscript{IV-1}]{99}

cet homomorphisme se factorise en
$r:C\xrightarrow{r'}G\to G/\mathfrak R$, où $r'$ est un
$A$-homomorphisme continu, si bien que $G$ est ainsi muni d'une structure de
$C$-algèbre topologique (discrète). D'autre part, par composition avec
l'homomorphisme canonique
\[
  \mathfrak J/\mathfrak J^2\longrightarrow
  \mathfrak J_\beta/\mathfrak J_\beta^2\longrightarrow F_{\beta,n},
\]
on déduit de (19.5.6.9) un $C$-homomorphisme continu
$g:\mathfrak J/\mathfrak J^2\to G/\mathfrak R$. Comme
$\mathfrak J/\mathfrak J^2$ est un $C$-module formellement projectif, $g$ se
factorise en
\[
  \mathfrak J/\mathfrak J^2\xrightarrow{h}G\longrightarrow G/\mathfrak R,
\]
où $h$ est un $C$-homomorphisme continu. Comme $G$ est discret, il existe un
indice $\gamma\geq\beta$ tel que l'image par $h$ de
$((\mathfrak b_\gamma\cap\mathfrak J)+\mathfrak J^2)/\mathfrak J^2$ soit
nulle, ainsi que l'image par $r'$ de
$(\mathfrak b_\gamma+\mathfrak J)/\mathfrak J$, de sorte que $h$ se factorise
en
\[
  \mathfrak J/\mathfrak J^2\longrightarrow
  \mathfrak J_\gamma/\mathfrak J_\gamma^2
  \xrightarrow{h'}G,
\]
où $h'$ est un $C_\gamma$-homomorphisme. Par prolongement, on déduit donc de
$h'$ un homomorphisme de $C_\gamma$-algèbres
\[
  w:\mathbf S_{C_\gamma}^*(\mathfrak J_\gamma/\mathfrak J_\gamma^2)
  \longrightarrow G
\]
et par construction, tout élément de degré $m$ a pour image par $w$ un élément
de $\mathfrak R$, donc tout élément de degré $n=2m$ a une image nulle, puisque
$\mathfrak R^2=0$~; autrement dit, $w$ se factorise en
\[
  \mathbf S_{C_\gamma}^*(\mathfrak J_\gamma/\mathfrak J_\gamma^2)
  \longrightarrow F_{\gamma,n}\xrightarrow{w'_\gamma}G.
\]
Appliquons maintenant à
\[
  v_{\alpha\gamma}:F_{\gamma,n}
  \xrightarrow{q_{\beta\gamma}}F_{\beta,n}
  \xrightarrow{v_{\alpha\beta}}E_{\alpha,n}
\]
le lemme (19.5.6.7), dont les hypothèses sont vérifiées~; il existe donc un
$\delta\geq\gamma$ et un homomorphisme
$u_{\gamma\delta}:E_{\delta,n}\to F_{\gamma,n}$ tel que le composé
$v_{\alpha\gamma}\circ u_{\gamma\delta}$ soit l'homomorphisme de transition
$p_{\alpha\delta}:E_{\delta,n}\to E_{\alpha,n}$. On obtient donc finalement un
diagramme commutatif de $A$-homomorphismes continus
\[
\xymatrix@C=3em@R=2.5em{
  && E_{\alpha,n}\ar[r]^{f_{\alpha,n}} & G/\mathfrak R \\
  && F_{\beta,n}\ar[u]_{v_{\alpha\beta}} & \\
  B\ar[r]_{p_\delta} & E_{\delta,n}\ar[uur]^{p_{\alpha\delta}}
    \ar[r]_{u_{\gamma\delta}} &
  F_{\gamma,n}\ar[u]_{q_{\beta\gamma}}\ar[r]^{w'_\gamma} & G\ar[uu]
}
\]
et le composé $f':B\to G$ des homomorphismes de la ligne inférieure est donc
tel que $f$ se factorise en $B\xrightarrow{f'}G\to G/\mathfrak R$, ce qui
prouve que $B$ est une $A$-algèbre formellement lisse.

La démonstration du th.~(19.5.3) est ainsi achevée.
\end{env}

\begin{corollary}[19.5.7]
\label{0.19.5.7-fr}
Soient $A$ un anneau topologique, $B$ une $A$-algèbre topologique,
$(\mathfrak b_\lambda)$ un système fondamental d'idéaux ouverts dans $B$,
$\mathfrak J$ un idéal de $B$, $C=B/\mathfrak J$ la $A$-algèbre topologique
quotient. On pose $C_\lambda=B/(\mathfrak b_\lambda+\mathfrak J)$. On suppose
que~: $1^\circ$ pour tout $n$, la topologie induite sur $\mathfrak J^n$ par
celle de $B$ est aussi la topologie du $B$-module $\mathfrak J^n$ déduite de
la topologie de $B$ (19.0.2) (cette condition sera remplie en particulier si
$B$ est noethérien et sa topologie préadique ($\mathbf 0_{\mathrm I}$,
7.3.2))~; $2^\circ$ $C$ est une $A$-algèbre formellement lisse. Sous ces
conditions~:
\begin{enumerate}
  \item[\rm(i)] Si $B$ est une $A$-algèbre formellement lisse, alors, pour tout
  $\lambda$, $(\mathfrak J/\mathfrak J^2)\otimes_C C_\lambda$ est un
  $C_\lambda$-module projectif.
  \item[\rm(ii)] Si $B$ est une $A$-algèbre formellement lisse et un anneau
  préadmissible, alors pour tout $\lambda$, l'homomorphisme canonique
  \begin{equation}
  \label{0.19.5.7.1-fr}
    \varphi_\lambda=\varphi\otimes 1_{C_\lambda}:
    \mathbf S_{C_\lambda}^*((\mathfrak J/\mathfrak J^2)\otimes_C C_\lambda)
    \longrightarrow\gr_{\mathfrak J}^*(B)\otimes_C C_\lambda
    \tag{19.5.7.1}
  \end{equation}
  est bijectif.

  \oldpage[0\textsubscript{IV-1}]{100}

  \item[\rm(iii)] Réciproquement, supposons que $B$ soit préadmissible, que la
  suite $(\mathfrak J^n)$ tende vers $0$, et que, pour tout $\lambda$,
  $(\mathfrak J/\mathfrak J^2)\otimes_C C_\lambda$ soit un $C_\lambda$-module
  projectif et l'homomorphisme (19.5.4.1) soit bijectif. Alors $B$ est une
  $A$-algèbre formellement lisse.
\end{enumerate}
\end{corollary}

\begin{proof}
En effet, la topologie de $\mathfrak J/\mathfrak J^2$ et celle des
$\mathbf S_C^n(\mathfrak J/\mathfrak J^2)$ sont alors aussi déduites de celle
de $B$ (19.5.1)~; les conclusions sont alors des conséquences immédiates de
(19.5.3) et des critères de (19.1.4) et (19.2.4) spéciaux à ce type de
topologies.
\end{proof}

\begin{remark}[19.5.8]
\label{0.19.5.8-fr}
Supposons que $\mathfrak J/\mathfrak J^2$ soit un $C$-module de type fini et
que, pour la topologie quotient de celle de $B$, $C$ soit un anneau de Zariski~;
soit $\mathfrak r$ un idéal de définition de $C$, de sorte que les
$\mathfrak r^n$ forment un système fondamental de voisinages de $0$ dans $C$.
Alors les conclusions de (i) et (ii) dans (19.5.7) peuvent être remplacées par
les suivantes~:
\begin{enumerate}
  \item[\rm(i')] $\mathfrak J/\mathfrak J^2$ est un $C$-module projectif.
  \item[\rm(ii')] L'homomorphisme canonique
  $\varphi:\mathbf S_C^*(\mathfrak J/\mathfrak J^2)
  \to\gr_{\mathfrak J}^*(B)$ est bijectif.
\end{enumerate}
En effet, il est clair que (i') entraîne la conclusion de (i) dans (19.5.7).
Inversement, si $(\mathfrak J/\mathfrak J^2)\otimes_C C_\lambda$ est un
$C_\lambda$-module projectif pour tout $\lambda$, alors
$(\mathfrak J/\mathfrak J^2)\otimes_C(C/\mathfrak r^n)$ est un
$(C/\mathfrak r^n)$-module projectif (donc plat) pour tout $n$~; on en conclut
que $\mathfrak J/\mathfrak J^2$ est un $C$-module plat
($\mathbf 0_{\mathrm{III}}$, 10.2.2), donc projectif puisqu'il est de
présentation finie (Bourbaki, \emph{Alg. comm.}, chap.~II, \textsection5,
n\textsuperscript{o}~2, cor.~2 du th.~1). D'autre part, les $C$-modules
$\mathbf S_C^n(\mathfrak J/\mathfrak J^2)$ et $\gr_{\mathfrak J}^n(B)$ sont
de type fini, et on sait que lorsque $C$ est un anneau de Zariski, il revient
au même de dire alors que $\varphi_n$ est bijectif ou que
$\widehat\varphi_n$ est bijectif (Bourbaki, \emph{Alg. comm.}, chap.~III,
\textsection3, n\textsuperscript{o}~5, prop.~9), donc (ii) est équivalent à
(ii').
\end{remark}

\subsection{Cas des algèbres sur un corps}
\label{subsection:0.19.6-fr}

\begin{theorem}[19.6.1) (Cohen]
\label{0.19.6.1-fr}
Soient $k$ un corps, $K$ une extension de $k$, $k$ et $K$
étant munis des topologies discrètes. Pour que $K$ soit une $k$-algèbre
formellement lisse, il faut et il suffit que $K$ soit une extension séparable
de $k$.
\end{theorem}

\begin{proof}
La nécessité de la condition sera établie en (19.6.5.1) (et naturellement ne
sera pas utilisée jusque-là)~; nous nous bornerons ici à prouver que la
condition est \emph{suffisante}. Distinguons deux cas~:

\emph{I. --- $K$ est une extension séparable de type fini de $k$.} On sait
alors (Bourbaki, \emph{Alg.}, chap.~V, \textsection9,
n\textsuperscript{o}~3, th.~2) qu'il existe une sous-extension pure
$K'=k(T_1,\ldots,T_n)$ de $K$ telle que $K$ soit une extension algébrique
finie séparable de $K'$. Compte tenu de (19.3.5, (iii)), on peut donc se
borner, soit au cas où $K=K'$, soit au cas où $K$ est algébrique fini sur $k$.
Dans le premier cas, on sait que $A=k[T_1,\ldots,T_n]$ est une $k$-algèbre
formellement lisse (19.3.3), donc il en est de même de
$k(T_1,\ldots,T_n)=S^{-1}A$, où $S=A-\{0\}$ (19.3.5, (iv)). Dans le second
cas, \emph{tous} les groupes de cohomologie de Hochschild
$\mathrm H_k^n(K,L)$, pour un $(K,K)$-bimodule quelconque $L$, sont nuls~: en
effet, si l'on considère la $k$-algèbre produit tensoriel $C=K\otimes_k K$,
$L$ est un $C$-module à gauche et le groupe de cohomologie
$\mathrm H_k^n(K,L)$ est égal à $\operatorname{Ext}_C^n(K,L)$, où $K$ est
aussi considéré comme $(K,K)$-bimodule,

\oldpage[0\textsubscript{IV-1}]{101}

donc comme $C$-module à gauche (M, IX, 4). Or, comme $K$ est une extension
finie séparable de $k$, on sait que $K\otimes_k K$ est composée directe
d'extensions de $k$, dont l'une est $K$ lui-même (Bourbaki, \emph{Alg.},
chap.~VIII, \textsection8, prop.~3)~; cela entraîne donc que $K$ est un
$C$-module \emph{projectif}, d'où notre assertion. Toute $k$-extension de $K$
de noyau $L$ est donc $k$-triviale, et \emph{a fortiori} les $k$-extensions
commutatives le sont, d'où le théorème dans ce cas (19.4.4).

\emph{II. --- Cas général.} --- Avec les notations de (18.4.5), il s'agit de
prouver que $\operatorname{Hom}_K(\mathrm H_2(P'_\bullet),L)=0$ pour
\emph{tout} $K$-espace vectoriel $L$, ce qui signifie évidemment que
$\mathrm H_2(P'_\bullet)=0$. Si $K$ est réunion d'une famille filtrante de
sous-extensions $K^{(\alpha)}$ de $k$, $P'_\bullet$ est \emph{limite
inductive} des complexes correspondants $P'^{(\alpha)}_\bullet$, puisque le
foncteur $\varinjlim$ commute au produit tensoriel dans la catégorie des
$k$-modules~; par l'exactitude du foncteur $\varinjlim$ dans cette catégorie,
on a donc
$\mathrm H_2(P'_\bullet)=\varinjlim\mathrm H_2(P'^{(\alpha)}_\bullet)$.
Comme l'hypothèse que $K$ est séparable entraîne qu'il en est de même de toute
sous-extension de $K$ et que $K$ est réunion des sous-extensions de type fini,
la première partie de la démonstration entraîne bien que $K$ est une
$k$-algèbre formellement lisse. C.Q.F.D.
\end{proof}

\begin{corollary}[19.6.2]
\label{0.19.6.2-fr}
Soient $A$ un anneau local séparé et complet, $K$ son corps résiduel, $k$ un
sous-corps de $K$ tel que $K$ soit une extension séparable de $k$. Alors il
existe un sous-corps $K'$ de $A$ contenant $k$, tel que la restriction à $K'$
de l'homomorphisme canonique $A\to K$ soit un isomorphisme de $K'$ sur $K$.
\end{corollary}

\begin{proof}
Soit $\mathfrak m$ l'idéal maximal de $A$. En vertu de l'hypothèse et de
(19.6.1), $K$ est une $k$-algèbre formellement lisse~; appliquant (19.3.10) en
remplaçant $C$ par $A$ et $\mathfrak J$ par $\mathfrak m$, on voit que
l'automorphisme identique de $K=A/\mathfrak m$ se factorise en
$K\xrightarrow{u}A\to A/\mathfrak m$, où $u$ est un $k$-homomorphisme, donc
nécessairement un $k$-isomorphisme de $K$ sur un sous-corps $K'$ contenant
$k$.
\end{proof}

\begin{corollary}[19.6.3]
\label{0.19.6.3-fr}
Soient $A$ un anneau local noethérien complet, $\mathfrak m$ son idéal maximal,
$k$ son corps résiduel. Les conditions suivantes sont équivalentes~:
\begin{enumerate}
  \item[a)] $A$ contient un sous-corps.
  \item[b)] Si $p$ est la caractéristique de $k$, on a $pA=0$.
  \item[c)] Il existe un corps $K$ tel que $A$ soit isomorphe à un anneau
  quotient d'un anneau de séries formelles $B=K[[T_1,\ldots,T_n]]$.
\end{enumerate}
Lorsqu'il en est ainsi, on peut supposer que $A$ est isomorphe à $B/\mathfrak b$,
où $\mathfrak b$ est contenu dans le carré de l'idéal maximal de $B$.
\end{corollary}

\begin{proof}
Il est immédiat que $c)$ implique $a)$, car $A$ est $\neq0$ puisque c'est un
anneau local~; comme $K$ est un corps et que l'homomorphisme composé
$K\xrightarrow{\varphi}B\to A$ (où $\varphi$ est l'injection canonique) n'est
pas nul, il est injectif. Pour voir que $a)$ entraîne $b)$, il suffit de
remarquer que si $K$ est un sous-corps de $A$, $K$ est isomorphe à un
sous-corps de $k$ et a par suite même caractéristique $p$~; comme $p.1=0$ dans
$K$, donc dans $A$, on a $pA=0$. Réciproquement, $b)$ entraîne $a)$, car le
composé des homomorphismes canoniques
$\mathbf Z\xrightarrow{j}A\to k$ a pour noyau $p\mathbf Z$, et l'on a
$j(p\mathbf Z)=0$, donc $\operatorname{Ker}(j)=p\mathbf Z$ et $A$ contient un
anneau $A'$ isomorphe à $\mathbf Z/p\mathbf Z$ et ne rencontrant pas
$\mathfrak m$~; si $p>0$, $A'$ est un corps~; sinon, tout

\oldpage[0\textsubscript{IV-1}]{102}

élément de $A'$ étant inversible dans $A$, le corps des fractions de $A'$
(isomorphe à $\mathbf Q$) est aussi contenu dans $A$.

Prouvons enfin que $a)$ entraîne $c)$~; la condition $a)$ entraîne que $A$
contient un sous-corps premier $k_0$, isomorphe au sous-corps premier de $k$,
et dont $k$ est par suite extension séparable. Appliquant (19.6.2), on voit
d'abord qu'il existe un isomorphisme $f$ de $k$ sur un sous-corps $K$ de $A$.
D'autre part, soit $(x_i)_{1\leq i\leq n}$ un système d'éléments de
$\mathfrak m$ tel que les classes mod.~$\mathfrak m^2$ des $x_i$ forment une
base (sur $k$) de $\mathfrak m/\mathfrak m^2$. Puisque $A$ est complet, il y a
alors un homomorphisme continu d'anneaux
$u:k[[T_1,\ldots,T_n]]\to A$ tel que $u$ soit égal à $f$ dans $k$ et que
$u(T_i)=x_i$ pour tout $i$, et cet homomorphisme est surjectif en vertu du
choix des $x_i$ (Bourbaki, \emph{Alg. comm.}, chap.~III, \textsection2,
n\textsuperscript{o}~9, prop.~11).
\end{proof}

\begin{theorem}[19.6.4]
\label{0.19.6.4-fr}
Soient $k$ un corps, $A$ une $k$-algèbre locale noethérienne, $\mathfrak m$ son
idéal maximal, $K=A/\mathfrak m$ son corps résiduel, $A$ étant muni de la
topologie $\mathfrak m$-préadique. On suppose que $K$ soit une extension
séparable de $k$. Alors les conditions suivantes sont équivalentes~:
\begin{enumerate}
  \item[a)] $A$ est une $k$-algèbre formellement lisse.
  \item[b)] Le complété $\widehat A$ de $A$ est $k$-isomorphe à un anneau de
  séries formelles $K[[T_1,\ldots,T_n]]$ (dont la structure de $k$-algèbre est
  définie par l'homomorphisme composé
  \[
    k\xrightarrow{\varphi}K\xrightarrow{\psi}K[[T_1,\ldots,T_n]],
  \]
  où $\psi$ est l'injection canonique et $\varphi$ l'homomorphisme définissant
  la structure d'extension de $K$).
  \item[b')] $\widehat A$ est un anneau isomorphe à un anneau de séries
  formelles $K[[T_1,\ldots,T_n]]$.
  \item[c)] $A$ est un anneau local régulier.
\end{enumerate}
\end{theorem}

\begin{proof}
Le fait que $a)$ entraîne $b)$ est le cas particulier de ((19.5.4),
équivalence de $a)$ et $c)$), appliqué en y remplaçant $A$ par $k$, $B$ par
$A$ et $\mathfrak J$ par $\mathfrak m$, compte tenu de (19.6.1). Il est clair
que $b)$ entraîne $b')$ et $b')$ entraîne $c)$ par (17.1.1). Enfin si $c)$ est
vérifiée, il résulte de (17.1.1, $a)$) et de ((19.5.4), équivalence de $b)$ et
$a)$), appliqué à $K=A/\mathfrak m$ remplaçant $C$, que $A$ est une
$k$-algèbre formellement lisse.
\end{proof}

\begin{corollary}[19.6.5]
\label{0.19.6.5-fr}
Soient $k$ un corps, $A$ une $k$-algèbre locale noethérienne, $\mathfrak m$ son
idéal maximal, $A$ étant muni de la topologie $\mathfrak m$-préadique.
Supposons que $A$ soit une $k$-algèbre formellement lisse~; alors $A$ est
géométriquement régulier sur $k$, autrement dit (cf.~$\mathbf{IV}$, 6.7.6),
pour toute extension finie $k'$ de $k$, l'anneau semi-local
$A'=A\otimes_k k'$ est régulier (17.3.6). En outre, si $K$ est le corps
résiduel de $A$, $\widehat A$ est isomorphe à un anneau de séries formelles
$K[[T_1,\ldots,T_n]]$.
\end{corollary}

\begin{proof}
Comme $\widehat A'=\widehat A\otimes_k k'$, il résulte de (19.3.6) et de
(17.1.5) (appliqué aux anneaux locaux en les idéaux maximaux de $A'$) qu'on
peut se borner au cas où $A$ est complet~; alors $A'$ est une $k'$-algèbre
formellement lisse (19.3.5, (iii)) et par ailleurs est composée directe de
$k'$-algèbres locales, qui sont aussi formellement lisses (19.3.5, (v)). On est
donc ramené à prouver que $A$ est \emph{régulier}. Soit $k_0$ le sous-corps
premier de $k$~; comme $k$ est une extension séparable de $k_0$, c'est une
$k_0$-algèbre formellement lisse (19.6.1) et en vertu de l'hypothèse, il en est
de même de $A$ (19.3.5, (ii)). Comme le corps résiduel $K$ de $A$ est une
extension séparable de $k_0$, on peut alors appliquer (19.6.4) à $A$ et $k_0$,
d'où la conclusion.
\end{proof}

\oldpage[0\textsubscript{IV-1}]{103}

\begin{corollary}[19.6.5.1]
\label{0.19.6.5.1-fr}
Soient $A$ un anneau local noethérien, $k$ un sous-corps de $A$ tel que $A$
soit une $k$-algèbre formellement lisse (pour sa topologie préadique). Alors
tout corps $K$ tel que $k\subset K\subset A$ est une extension séparable de
$k$.
\end{corollary}

\begin{proof}
En effet, pour toute extension finie $k'$ de $k$, l'anneau $K\otimes_k k'$ 
s'identifie à un sous-anneau de $A'=A\otimes_k k'$~; comme $A'$ est un anneau
régulier, il est réduit, donc il en est de même de $K$, ce qui prouve que $K$
est une extension séparable de $k$ (Bourbaki, \emph{Alg.}, chap.~VIII,
\textsection7, n\textsuperscript{o}~3, th.~1).
\end{proof}

On notera que cela démontre que la condition de l'énoncé de (19.6.1) est
\emph{nécessaire}.

\begin{remark}[19.6.5.2]
\label{0.19.6.5.2-fr}
Si $K$ est une extension \emph{non séparable} de $k$, l'anneau de séries
formelles $K[[T_1,\ldots,T_n]]$ \emph{n'est donc pas} une $k$-algèbre
formellement lisse (pour la structure de $k$-algèbre usuelle rappelée dans
(19.6.4)). D'autre part, il y a des $k$-algèbres formellement lisses $A$ qui
sont des anneaux locaux noethériens complets dont le corps résiduel $K$ est une
extension \emph{non séparable} de $k$ (par exemple les complétés des localisés
de $B=k[T_1,\ldots,T_n]$ en des idéaux maximaux $\mathfrak n$ tels que
$B/\mathfrak n$ soit une extension algébrique finie non séparable de $k$)~; un
tel anneau ne peut donc être $k$-isomorphe à $K[[T_1,\ldots,T_n]]$ bien qu'il
lui soit isomorphe en vertu de (19.6.5).

On peut dans certains cas préciser le corollaire (19.6.5)~:
\end{remark}

\begin{proposition}[19.6.6]
\label{0.19.6.6-fr}
Soient $k$ un corps, $A$ une $k$-algèbre locale noethérienne, $\mathfrak m$ son
idéal maximal, $A$ étant muni de la topologie $\mathfrak m$-préadique.
Supposons que le corps résiduel $K=A/\mathfrak m$ de $A$ soit tel qu'il existe
une extension radicielle finie $k_0$ de $k$ pour laquelle le corps résiduel de
l'anneau local $K\otimes_k k_0$ soit une extension séparable de $k_0$ (nous
exprimerons plus tard cette propriété en disant que $K$ est de
\emph{multiplicité radicielle finie} sur $k$ ($\mathbf{IV}$, 4.7.4) et nous
verrons que cette condition est toujours satisfaite lorsque $K$ est une
\emph{extension de type fini} de $k$). Alors les conditions suivantes sont
équivalentes~:
\begin{enumerate}
  \item[a)] $A$ est une $k$-algèbre formellement lisse.
  \item[b)] Il existe une extension radicielle finie $k'$ de $k$ telle que
  $\widehat A\otimes_k k'$ soit $k'$-isomorphe à une algèbre de séries
  formelles $K'[[T_1,\ldots,T_n]]$, où $K'$ est une extension séparable de
  $k'$.
  \item[b')] Il existe une extension finie $k'$ de $k$ telle que l'anneau
  semi-local $\widehat A\otimes_k k'$ ait au moins un anneau local composant
  qui soit $k'$-isomorphe à une algèbre de séries formelles
  $K'[[T_1,\ldots,T_n]]$, où $K'$ est une extension séparable de $k'$.
  \item[c)] $A$ est géométriquement régulier sur $k$ (19.6.5).
\end{enumerate}
\end{proposition}

\begin{proof}
Notons d'abord que si $k'$ est une extension radicielle de $k$, il n'y a qu'un
seul idéal de $A'=A\otimes_k k'$ au-dessus de $\mathfrak m$, formé des éléments
dont une puissance $p^h$-ième ($p$ exposant caractéristique de $k$) est dans
$\mathfrak m$ pour un $h$ convenable (Bourbaki, \emph{Alg. comm.}, chap.~V,
\textsection2, n\textsuperscript{o}~3, lemme~4)~; $A'$ est donc un anneau
local, et il en est de même de
$K\otimes_k k'=(A\otimes_k k')/(\mathfrak m\otimes_k k')$~; en outre les corps
résiduels de ces deux anneaux sont identiques. Rappelons d'autre part que si
$K$ est une extension séparable de $k$, alors, pour toute extension finie
$k''$ de $k$, $K\otimes_k k''$ est composé direct de corps (Bourbaki,
\emph{Alg.}, chap.~VIII, \textsection7, n\textsuperscript{o}~3, cor.~1 du
th.~1), et par suite $\mathfrak m\otimes_k k''$ est le radical de
$A\otimes_k k''$, et les corps composants de $K\otimes_k k''$ sont les corps
résiduels aux idéaux maximaux de $A\otimes_k k''$~; en

\oldpage[0\textsubscript{IV-1}]{104}

outre, ces corps sont des extensions séparables de $k''$ puisque pour toute
extension $k_1$ de $k''$,
$K\otimes_k k_1=(K\otimes_k k'')\otimes_{k''}k_1$ est sans radical
(\emph{loc. cit.}, th.~1). Si l'on applique ces remarques au corps $K$ de
multiplicité radicielle finie sur $k$, on voit que pour \emph{toute} extension
finie $k''$ de $k_0$, $(K\otimes_k k'')_{\mathrm{red}}$ est séparable sur
$k''$.

Passons à la démonstration de (19.6.6). Il est trivial que $b)$ implique
$b')$. Montrons que $b')$ entraîne $a)$. Si l'on pose
$B'=K'[[T_1,\ldots,T_n]]$, l'hypothèse que $K'$ est séparable sur $k'$
entraîne, par les remarques du début, que pour toute extension finie $k''$ de
$k'$, les composants de l'anneau semi-local complet $B'\otimes_{k'}k''$
(égal à l'anneau de séries formelles
$(K'\otimes_{k'}k'')[[T_1,\ldots,T_n]]$) sont les anneaux
$K_i''[[T_1,\ldots,T_n]]$, où les $K_i''$ sont les corps composants de
$K'\otimes_{k'}k''$. Comme les composants locaux de $B'\otimes_{k'}k''$ sont
aussi des composants locaux de
\[
  \widehat A\otimes_k k''
  =(\widehat A\otimes_k k')\otimes_{k'}k'',
\]
on voit que l'hypothèse $b')$ est encore vérifiée lorsqu'on remplace $k'$ par
une quelconque de ses extensions finies. En particulier, on peut supposer que
$k'$ est \emph{quasi-galoisienne}, donc extension galoisienne d'une extension
radicielle $k'_0$ de $k$~; $A'=\widehat A\otimes_k k'$ peut alors s'écrire
$A'_0\otimes_{k'_0}k'$, où $A'_0=\widehat A\otimes_k k'_0$ est un anneau
local complet. On sait alors (Bourbaki, \emph{Alg.}, chap.~VIII,
\textsection8, prop.~4) que $A'$ est composé direct d'anneaux locaux
\emph{isomorphes} à $A'_0$ en tant que $k'_0$-algèbres. Or, un de ces anneaux
est par hypothèse $k'$-isomorphe à un anneau de séries formelles
$K'[[T_1,\ldots,T_n]]$, où $K'$ est une extension séparable de $k'$, donc
\emph{aussi} une extension séparable de $k'_0$~; compte tenu de (19.6.4),
$A'_0$ est donc une $k'_0$-algèbre formellement lisse. Mais comme $k'_0$ est
un $k$-module fidèlement plat, projectif et de type fini, on conclut de
(19.4.7) que $\widehat A$ est une $k$-algèbre formellement lisse.

On a déjà vu (19.6.5) que $a)$ implique $c)$~; il reste donc à vérifier que
$c)$ implique $b)$. Or, si $c)$ est vérifiée, alors, pour toute extension
radicielle finie $k'$ de $k$ contenant $k_0$, $A\otimes_k k'$ est un anneau
local régulier, dont le corps résiduel $K'$ est une extension séparable de
$k'$~; il résulte donc de (19.6.4) que son complété
$\widehat A\otimes_k k'$ est $k'$-isomorphe à un anneau de séries formelles
$K'[[T_1,\ldots,T_n]]$.
\end{proof}

\begin{remarks}[19.6.7]
\label{0.19.6.7-fr}
\begin{enumerate}
  \item[\rm(i)] On notera que l'hypothèse que $K$ est de multiplicité
  radicielle finie sur $k$ n'a été utilisée que dans la démonstration de
  l'implication $b')\Rightarrow a)$.
  \item[\rm(ii)] Nous prouverons plus tard (22.5.8) que $a)$ et $c)$ sont
  équivalentes, \emph{sans hypothèse} sur l'extension $K$ de $k$.
\end{enumerate}
\end{remarks}

\subsection{Cas des homomorphismes locaux; théorèmes d'existence et
d'unicité}
\label{subsection:0.19.7-fr}

Dans ce numéro, lorsqu'un anneau semi-local est considéré comme un anneau
topologique, il est toujours sous-entendu qu'il s'agit de sa topologie
$\mathfrak r$-préadique, où $\mathfrak r$ est son radical. Tout homomorphisme
local d'anneaux locaux est donc automatiquement \emph{continu}.

\begin{theorem}[19.7.1]
\label{0.19.7.1-fr}
Soient $A$, $B$ deux anneaux locaux noethériens, $\mathfrak m$, $\mathfrak n$
leurs idéaux maximaux respectifs, $k=A/\mathfrak m$ le corps résiduel de $A$~;
on suppose $A$ et $B$ munis respectivement des topologies
$\mathfrak m$-préadique et $\mathfrak n$-préadique. Soit
$\varphi:A\to B$ un homomorphisme local, et posons $B_0=B\otimes_A k$.
Les propriétés suivantes sont équivalentes~:
\begin{enumerate}
  \item[a)] $B$ est une $A$-algèbre \emph{formellement lisse}.

\oldpage[0\textsubscript{IV-1}]{105}

  \item[b)] $B$ est un $A$-module \emph{plat} et
  $B_0=B/\mathfrak mB$ (muni de la topologie quotient) est une $k$-algèbre
  \emph{formellement lisse}.
\end{enumerate}
\end{theorem}

\begin{proof}
La démonstration se fait en plusieurs étapes.

\medskip
\noindent\phantomsection\label{0.19.7.1.1-fr}\textbf{(19.7.1.1)}\enspace
Démontrons d'abord que $b)$ entraîne $a)$~; nous allons appliquer le critère
(19.4.7), avec $\mathfrak J=\mathfrak m$~; en vertu de la seconde hypothèse
dans $b)$, tout revient à montrer que $B$ est un $A$-module
\emph{formellement projectif}. L'hypothèse entraîne que pour tout $h>0$,
$B/\mathfrak m^hB$ est un $(A/\mathfrak m^h)$-module plat
($\mathbf 0_{\mathrm{III}}$, 10.2.1)~; comme les $\mathfrak m^h$ forment un
système fondamental de voisinages de 0 dans $A$ et que
$(B/\mathfrak m^hB)\otimes_{A/\mathfrak m^h}k=B_0$, on peut remplacer $A$ et
$B$ par $A/\mathfrak m^h$ et $B/\mathfrak m^hB$ respectivement, et par suite
supposer $A$ \emph{artinien} (donc discret). Comme $B_0$ est une $k$-algèbre
formellement lisse, c'est un anneau \emph{régulier} (19.6.5)~; soit
$(x_i^0)_{1\leq i\leq n}$ un système régulier de paramètres pour $B_0$
(17.1.6), et pour tout $i$, soit $x_i\in B$ tel que $x_i^0$ soit son image
dans $B_0=B/\mathfrak mB$~; comme les $x_i^0$ engendrent l'idéal maximal
$\mathfrak n_0=\mathfrak n/\mathfrak mB$ de $B_0$, les idéaux
\[
  \mathfrak Q_m=\sum_{i=1}^n(x_i^0)^mB_0 \qquad (m>0)
\]
forment un système fondamental de voisinages de 0 dans $B_0$, car
$\mathfrak Q_m$ est évidemment contenu dans $\mathfrak n_0^m$, et d'autre
part contient $\mathfrak n_0^{mn}$. Posons
$\mathfrak J_m=\sum_{i=1}^n x_i^mB$ pour tout $m>0$~; il est clair que
$\mathfrak n=\mathfrak J_1+\mathfrak mB$~; comme il existe un $h>0$ tel que
$\mathfrak m^h=0$, on a
$\mathfrak n^m\subset\mathfrak J_1^{m-h}\subset\mathfrak n^{m-h}$ pour
$m>h$, et comme on a vu que $\mathfrak J_m\supset\mathfrak J_1^{mn}$, on voit
que les $\mathfrak J_m$ forment un système fondamental de voisinages de 0 dans
$B$. Tout revient par suite à prouver que les $B/\mathfrak J_m$ sont des
$A$-modules \emph{libres}, et il revient au même de voir que ce sont des
$A$-modules \emph{plats} ($\mathbf 0_{\mathrm{III}}$, 10.1.3). Or,
l'hypothèse que $(x_i^0)$ est une suite $B_0$-régulière d'éléments de l'idéal
maximal de $B_0$ entraîne la même propriété pour la suite des $(x_i^0)^m$
($1\leq i\leq n$) pour tout $m>0$ (15.1.20)~; la conclusion résulte donc de
(15.1.16, $b)$ et $c)$).
\end{proof}

\begin{lemma}[19.7.1.2]
\label{0.19.7.1.2-fr}
Soient $A$ un anneau topologique, $B$, $C$ deux $A$-algèbres topologiques qui
sont des anneaux locaux noethériens. On suppose en outre que $C$ est complet
et que le corps résiduel $B/\mathfrak m$ de $B$ est un $A$-module de type fini.
Soit $E$ le produit tensoriel complété $B\widehat\otimes_A C$. Alors~:
\begin{enumerate}
  \item[(i)] $E$ est un anneau semi-local noethérien complet.
  \item[(ii)] L'idéal $\mathfrak mE$ est contenu dans le radical de $E$, et
  pour tout $h>0$, $E/\mathfrak m^hE$ est isomorphe à
  $(B/\mathfrak m^h)\otimes_A C$.
  \item[(iii)] Si $C$ est un $A$-module plat, $E$ est un $B$-module plat.
\end{enumerate}
\end{lemma}

\begin{proof}
Par définition, $E$ est le séparé complété du produit tensoriel $B\otimes_A C$
pour la topologie définie par les idéaux
$\operatorname{Im}(\mathfrak m^i\otimes_A C)+
\operatorname{Im}(B\otimes_A\mathfrak n'{}^j)$
($\mathbf 0_{\mathrm I}$, 7.7.5). Si l'on pose
\[
  \mathfrak r=\operatorname{Im}(\mathfrak m\otimes_A C)+
  \operatorname{Im}(B\otimes_A\mathfrak n'),
\]
on a
$\mathfrak r^{2i}\subset
\operatorname{Im}(\mathfrak m^i\otimes_A C)+
\operatorname{Im}(B\otimes_A\mathfrak n'{}^i)\subset\mathfrak r^i$, donc $E$
est aussi le séparé complété de $B\otimes_A C$ pour la topologie
$\mathfrak r$-préadique. Par hypothèse,
$(B\otimes_A C)/\mathfrak r=(B/\mathfrak m)\otimes_A(C/\mathfrak n')$ est un
$(C/\mathfrak n')$-module de type fini, donc un anneau artinien~; en outre,
$\mathfrak r/\mathfrak r^2$, étant un quotient de
\[
  ((\mathfrak m/\mathfrak m^2)\otimes_A C)
  \oplus(B\otimes_A(\mathfrak n'/\mathfrak n'{}^2)),
\]
est un $(B\otimes_A C)$-module de type fini~; appliquant
($\mathbf 0_{\mathrm I}$, 7.2.11), on voit donc que $E$ est un anneau
\emph{noethérien}~; en outre, $E/\mathfrak rE$, isomorphe à
$(B\otimes_A C)/\mathfrak r$, étant artinien, $E$ est \emph{semi-local}.
Notons maintenant que $E$, qui est isomorphe à
\[
  \varprojlim_{i,j}((B/\mathfrak m^i)\otimes_A(C/\mathfrak n'{}^j)),
\]
est aussi isomorphe, par le théorème de la double limite projective, à
\[
  \varprojlim_i\left(\varprojlim_j
  ((B/\mathfrak m^i)\otimes_A(C/\mathfrak n'{}^j))\right)~;
\]

\oldpage[0\textsubscript{IV-1}]{106}

mais $\varprojlim_j((B/\mathfrak m^i)\otimes_A(C/\mathfrak n'{}^j))$ est le
séparé complété de $(B/\mathfrak m^i)\otimes_A C$, et comme $C$ est complet,
ce n'est autre que $(B/\mathfrak m^i)\otimes_A C$ lui-même,
$B/\mathfrak m^i$ étant un $A$-module de type fini puisque
$\mathfrak m/\mathfrak m^i$ est un $(B/\mathfrak m)$-module de type fini
($\mathbf 0_{\mathrm I}$, 7.3.6). On a donc
\[
  E=\varprojlim_i((B/\mathfrak m^i)\otimes_A C).
\]
Pour tout entier $h>0$, $\mathfrak m^hE$, étant un idéal de $E$, est fermé
dans $E$ ($\mathbf 0_{\mathrm I}$, 7.3.5), donc complet, et d'autre part il est
évidemment dense dans
\[
  \varprojlim_i
  (\operatorname{Im}(\mathfrak m^h/\mathfrak m^{h+i})\otimes_A C),
\]
donc égal à cette dernière limite projective. En outre, tous les systèmes
projectifs considérés sont définis par des homomorphismes \emph{surjectifs}~;
il résulte donc de ($\mathbf 0_{\mathrm{III}}$, 13.2.2) que
$E/\mathfrak m^hE$ est isomorphe à
\[
  (B/\mathfrak m^h)\otimes_A C
  =(B\otimes_A C)/\operatorname{Im}(\mathfrak m^h\otimes_A C).
\]
En particulier, comme
$\operatorname{Im}(\mathfrak m\otimes_A C)\subset\mathfrak r$, cela montre
que $\mathfrak mE$ est contenu dans $\mathfrak rE$, donc dans le radical de
$E$. Enfin, l'hypothèse que $C$ est un $A$-module plat entraîne que
$(B/\mathfrak m^h)\otimes_A C=E/\mathfrak m^hE$ est un
$(B/\mathfrak m^h)$-module plat pour tout $h>0$~: comme $B$ et $E$ sont
noethériens et que $\mathfrak mE$ est contenu dans le radical de $E$, il
résulte de ($\mathbf 0_{\mathrm{III}}$, 10.2.2) que $E$ est un $B$-module
plat.
\end{proof}

\begin{lemma}[19.7.1.3]
\label{0.19.7.1.3-fr}
Soient $A$ un anneau local noethérien, $\mathfrak m$ son idéal maximal, $k$ son
corps résiduel, $B_0$ une $k$-algèbre~; on suppose que $B_0$ est un anneau
local noethérien, complet et régulier. Alors il existe une $A$-algèbre
topologique $B$ qui est un anneau local noethérien complet, un $A$-module plat,
et tel que $B_0$ soit $k$-isomorphe à
$B\otimes_A k=B/\mathfrak mB$.
\end{lemma}

\begin{proof}
Comme $\widehat A$ est plat sur $A$ et a même corps résiduel, on peut se borner
au cas où $A$ est \emph{complet}.

Soit $K$ le corps résiduel de $B_0$, et distinguons deux cas~:

\medskip
\noindent I) $K$ est une extension \emph{séparable} de $k$. En vertu de
(19.6.4), $B_0$ est $k$-isomorphe à un anneau de séries formelles
$K[[T_1,\ldots,T_n]]$. Lorsque $B_0=K$, le lemme a déjà été démontré
($\mathbf 0_{\mathrm{III}}$, 10.3.1)~; soit $C$ un anneau local noethérien
complet qui est un $A$-module plat et tel que $C\otimes_A k$ soit isomorphe à
$K$. Pour $n\geq1$, il suffit de prendre (avec la notation précédente)
$B=C[[T_1,\ldots,T_n]]$~; on sait en effet (Bourbaki,
\emph{Alg. comm.}, chap.~III, §~3, n\textsuperscript{o}~4, cor.~3 du th.~1)
que $B$ est un $C$-module plat, donc aussi un $A$-module plat, et d'autre part,
il est immédiat que $C[[T_1,\ldots,T_n]]\otimes_A k$ est isomorphe à
$(C/\mathfrak mC)[[T_1,\ldots,T_n]]=B_0$.

\medskip
\noindent II) $K$ est de \emph{caractéristique} $p>0$, et par suite il en est
de même de $k$. Notons $P$ le corps premier $\mathbf F_p$, et $W(P)$ l'anneau
local complet des nombres $p$-adiques $\mathbf Z_p$, qui est un anneau de
valuation discrète (donc \emph{régulier}), et a $P$ pour corps résiduel.
Montrons d'abord qu'il existe un homomorphisme continu d'anneaux $W(P)\to A$
faisant donc de $A$ une $W(P)$-\emph{algèbre topologique}. En effet, si
$j:\mathbf Z\to A$ est l'homomorphisme canonique, on a
$j(p\mathbf Z)\subset\mathfrak m$ par hypothèse, d'où
$j^{-1}(\mathfrak m)=p\mathbf Z$, et par suite $j$ se factorise en
$\mathbf Z\to\mathbf Z_{p\mathbf Z}\xrightarrow{f}A$, où $f$ est un
homomorphisme local, donc continu, qui (puisque $A$ est complet) se prolonge
par continuité en l'homomorphisme $W(P)\to A$ cherché.

Comme $k$ est extension séparable de $P$, le cas I) montre qu'il y a un
homomorphisme local $W(P)\to W(k)$, où $W(k)$ est un anneau local noethérien
complet et un $W(P)$-module plat, tel que $W(k)\otimes_{W(P)}P$ soit isomorphe
à $k$. D'ailleurs, comme l'uniformisante $p$ de $W(P)$ est un élément
$W(k)$-régulier par platitude ($\mathbf 0_{\mathrm I}$, 6.3.4) et comme

\oldpage[0\textsubscript{IV-1}]{107}

$W(k)/pW(k)=k$, $pW(k)$ est l'idéal maximal de $W(k)$, ce qui entraîne que ce
dernier anneau est un \emph{anneau de valuation discrète complet} (Bourbaki,
\emph{Alg. comm.}, chap.~VI, §~3, n\textsuperscript{o}~5, prop.~9). Par
(19.7.1.1) on voit en outre (puisque $k$ est séparable sur $P$, donc une
$P$-algèbre formellement lisse (19.6.1)) que $W(k)$ est une $W(P)$-algèbre
\emph{formellement lisse}. Le $W(P)$-homomorphisme continu $W(k)\to k$ se
factorise donc en $W(k)\to A\to k$ (19.3.11), ce qui permet de considérer $A$
comme une $W(k)$-algèbre topologique. Appliquant maintenant le cas I) à $B_0$
considéré comme $P$-algèbre et à $W(P)$, on voit qu'il existe une
$W(P)$-algèbre $B_P$ qui est un anneau local noethérien complet, un
$W(P)$-module plat, et telle que $B_P\otimes_{W(P)}P$ soit $P$-isomorphe à
$B_0$. Utilisant de nouveau le fait que $W(k)$ est une $W(P)$-algèbre
formellement lisse, on voit par (19.3.11) que l'homomorphisme composé
$W(k)\to k\to B_0$ se factorise en $W(k)\to B_P\to B_0$~; en outre, comme
$k=W(k)/pW(k)$, on a
$B_P\otimes_{W(k)}k=B_P/pB_P=B_P\otimes_{W(P)}P=B_0$. Montrons que $B_P$ est
un $W(k)$-module \emph{plat}~; comme $W(k)$ est un anneau de valuation discrète
dont $p$ est l'uniformisante, il suffit de vérifier que $B_P$ est un
$W(k)$-module \emph{sans torsion} ($\mathbf 0_{\mathrm I}$, 6.3.4), ou encore
que $p$ est un élément $B_P$-régulier, ce qui résulte de ce que $B_P$ est un
$W(P)$-module plat ($\mathbf 0_{\mathrm I}$, 6.3.4).

Posons maintenant
\[
  B=B_P\widehat\otimes_{W(k)}A
\]
et notons que le corps résiduel de $A$ étant égal à celui de $W(k)$, est
\emph{a fortiori} un $W(k)$-module de type fini. Il résulte donc tout d'abord
de (19.7.1.2) que $B$ est un anneau semi-local noethérien complet,
$\mathfrak mB$ étant contenu dans le radical de $B$~; en outre
$B/\mathfrak mB$ est $k$-isomorphe à
$B_P\otimes_{W(k)}k=B_0$, donc $B$ est en fait un anneau local. Comme $B_P$
est un $W(k)$-module plat, (19.7.1.2) montre enfin que $B$ est un $A$-module
plat. C.Q.F.D.
\end{proof}

\begin{lemma}[19.7.1.4]
\label{0.19.7.1.4-fr}
Soient $A$ un anneau, $\mathfrak J$ un idéal de $A$, $M$, $N$ deux $A$-modules
séparés pour la topologie $\mathfrak J$-préadique. On suppose en outre que $N$
est complet pour la topologie $\mathfrak J$-préadique et que $M$ soit un
$A$-module plat. Soit $u:N\to M$ un $A$-homomorphisme~; si
$u\otimes1:N\otimes_A(A/\mathfrak J)\to M\otimes_A(A/\mathfrak J)$ est
bijectif, alors $u$ est bijectif.
\end{lemma}

\begin{proof}
Les modules gradués associés étant pris relatifs aux filtrations
$\mathfrak J$-préadiques, il résulte des hypothèses sur $M$ et $N$ relativement
aux topologies $\mathfrak J$-préadiques qu'il suffit de prouver que
$\gr(u):\gr_\bullet(N)\to\gr_\bullet(M)$ est \emph{bijectif} (Bourbaki,
\emph{Alg. comm.}, chap.~III, §~2, n\textsuperscript{o}~8, cor.~3 du th.~1).
Or, on a un diagramme commutatif
\[
\xymatrix{
  \gr_0(N)\otimes_{A_0}\gr_\bullet(A)
    \ar[r]^-{\gr_0(u)\otimes1}\ar[d]_{\varphi_N} &
  \gr_0(M)\otimes_{A_0}\gr_\bullet(A)\ar[d]^{\varphi_M} \\
  \gr_\bullet(N)\ar[r]_-{\gr(u)} & \gr_\bullet(M)
}
\]

\oldpage[0\textsubscript{IV-1}]{108}

où $A_0=A/\mathfrak J=\gr_0(A)$, et $\varphi_M$ et $\varphi_N$ sont les
applications canoniques ($\mathbf 0_{\mathrm{III}}$, 10.1.1.2). Par
hypothèse, $\gr_0(u)$ est bijectif ainsi que $\varphi_M$
($\mathbf 0_{\mathrm{III}}$, 10.2.1), et $\varphi_N$ est surjectif~; on en
déduit d'abord que $\gr_0(u)\otimes1$ est bijectif, puis que $\varphi_N$ est
injectif, donc bijectif, et enfin que $\gr(u)$ est bijectif.
\end{proof}

\begin{lemma}[19.7.1.5]
\label{0.19.7.1.5-fr}
Soient $A$ un anneau noethérien, $\mathfrak J$ un idéal de $A$, $B$, $B'$ deux
$A$-algèbres qui sont des anneaux locaux noethériens, les homomorphismes
$A\to B$, $A\to B'$ étant continus pour la topologie $\mathfrak J$-préadique
sur $A$. On suppose que~:
\begin{enumerate}
  \item[$1^{\mathrm o}$] $B$ et $B'$ sont complets pour les topologies
  $\mathfrak J$-préadiques~;
  \item[$2^{\mathrm o}$] $B$ est une $A$-algèbre formellement lisse~;
  \item[$3^{\mathrm o}$] $B'$ est un $A$-module plat.
\end{enumerate}
Posons $A_0=A/\mathfrak J$, et soit
$u_0:B\otimes_A A_0\to B'\otimes_A A_0$ un $A_0$-isomorphisme~; alors il
existe un $A$-isomorphisme $u:B\to B'$ tel que $u_0=u\otimes1$ (ce qui
entraîne que $B'$ est une $A$-algèbre formellement lisse et $B$ un $A$-module
plat).
\end{lemma}

\begin{proof}
Posons $B_0=B\otimes_A A_0$, $B'_0=B'\otimes_A A_0$. Notons que si
$\mathfrak m$ et $\mathfrak m'$ sont les idéaux maximaux de $B$ et $B'$, les
topologies $\mathfrak J$-préadiques sur $B$ et $B'$ sont séparées puisque
$\mathfrak JB\subset\mathfrak m$ et $\mathfrak JB'\subset\mathfrak m'$~; en
outre, comme $\mathfrak JB'$ est fermé dans $B'$ pour la topologie
$\mathfrak J$-préadique, l'homomorphisme composé
$B\to B_0\xrightarrow{u_0}B'_0$, qui est continu pour les topologies
$\mathfrak J$-préadiques, se factorise en $B\xrightarrow{u}B'\to B'_0$, où
$u$ est un $A$-homomorphisme continu (19.3.10). On a évidemment
$u_0=u\otimes1$, et l'hypothèse que $u_0$ est bijectif entraîne qu'il en est de
même de $u$ en vertu de (19.7.1.4).
\end{proof}

\medskip
\noindent\textbf{(19.7.1.6)}\enspace
\emph{Fin de la démonstration.} --- Pour achever de prouver (19.7.1), il faut
montrer que $a)$ entraîne $b)$~; on sait déjà que $a)$ entraîne que $B_0$ est
une $k$-algèbre formellement lisse (19.3.5, (iii)), donc tout revient à prouver
que $B$ est un $A$-module plat. Il revient au même d'établir que $\widehat B$
est un $\widehat A$-module plat (Bourbaki, \emph{Alg. comm.}, chap.~III, §~5,
n\textsuperscript{o}~4, prop.~4), et l'on sait que $\widehat B$ est une
$\widehat A$-algèbre formellement lisse (19.3.6)~; on peut donc se borner au
cas où $A$ et $B$ sont complets. Comme $B_0$ est une $k$-algèbre formellement
lisse, c'est un anneau régulier (19.6.5) et complet
($\mathbf 0_{\mathrm I}$, 6.3.5)~; appliquant (19.7.1.3), on voit qu'il existe
une $A$-algèbre $B'$ qui est un anneau local noethérien complet et un
$A$-module plat, un homomorphisme local $A\to B'$ et un $k$-isomorphisme
$B_0\simeq B'_0=B'\otimes_A k$. Il suffit alors d'appliquer (19.7.1.5) en
prenant pour $\mathfrak J$ l'idéal maximal de $A$, pour obtenir que $B$ est
$A$-isomorphe à $B'$, donc est un $A$-module plat. C.Q.F.D.

\begin{theorem}[19.7.2]
\label{0.19.7.2-fr}
Soient $A$ un anneau local noethérien, $\mathfrak J$ un idéal contenu dans
l'idéal maximal de $A$, $A_0=A/\mathfrak J$, $B_0$ un anneau local noethérien
complet, $A_0\to B_0$ un homomorphisme local faisant de $B_0$ une
$A_0$-algèbre formellement lisse. Alors il existe un anneau local noethérien
complet $B$, un homomorphisme local $A\to B$ faisant de $B$ un $A$-module
plat, et un $A_0$-isomorphisme
$u:B\otimes_A A_0\xrightarrow{\sim}B_0$. Si $(B',u')$ est un couple
satisfaisant aux mêmes conditions que $(B,u)$, il existe un $A$-isomorphisme
$v:B\xrightarrow{\sim}B'$ rendant commutatif le diagramme
\[
\xymatrix{
  B\otimes_A A_0\ar[rr]^{v\otimes1}_{\sim}\ar[dr]^{\sim}_u &&
  B'\otimes_A A_0\ar[dl]^{u'}_{\sim} \\
  & B_0
}
\]
\end{theorem}

\oldpage[0\textsubscript{IV-1}]{109}

\begin{proof}
Soit $\mathfrak m$ l'idéal maximal de $A$, de sorte que
$\mathfrak m_0=\mathfrak m/\mathfrak J$ est l'idéal maximal de $A_0$, $A$ et
$A_0$ ayant le même corps résiduel $k$. Posons
$B_{00}=B_0\otimes_{A_0}k$~; comme $B_{00}$ est une $k$-algèbre formellement
lisse (19.3.5, (iii)), c'est un anneau local régulier (19.6.5)~; en appliquant
(19.7.1.3), on voit qu'il existe une $A$-algèbre topologique $B$ qui est un
anneau local noethérien complet, un $A$-module plat, et pour laquelle on a un
$k$-isomorphisme
$u_{00}:B\otimes_A k\xrightarrow{\sim}B_{00}$. Notons qu'en vertu de (19.7.1),
$B$ est une $A$-algèbre formellement lisse, donc
$B\otimes_A A_0=B/\mathfrak JB$ est une $A_0$-algèbre formellement lisse
(19.3.5, (iii)) et un anneau local noethérien complet~; en outre, $B_0$ est un
$A_0$-module plat en vertu de l'hypothèse et de (19.7.1)~; comme on a un
$k$-isomorphisme
\[
  u_{00}:B\otimes_A k=(B\otimes_A A_0)\otimes_{A_0}k
  \xrightarrow{\sim}B_0\otimes_{A_0}k=B_{00},
\]
on déduit de (19.7.1.5), appliqué en y remplaçant $A$ par $A_0$ et
$\mathfrak J$ par $\mathfrak m_0$, qu'il existe un $A_0$-isomorphisme
$u:B\otimes_A A_0\xrightarrow{\sim}B_0$ tel que $u_{00}=u\otimes1$. Quant à
l'assertion d'unicité, notons que les idéaux $\mathfrak J^hB$ (resp.
$\mathfrak J^hB'$) sont fermés dans $B$ (resp. $B'$)
($\mathbf 0_{\mathrm I}$, 7.3.5), donc $B$ et $B'$ sont séparés et complets
pour les topologies $\mathfrak J$-préadiques (Bourbaki, \emph{Top. gén.},
chap.~III, 3\textsuperscript{e} éd., §~3, n\textsuperscript{o}~5, cor.~2 de la
prop.~9)~; on a par hypothèse un $A_0$-isomorphisme
$v_0:B\otimes_A A_0\xrightarrow{\sim}B'\otimes_A A_0$ tel que
$u'\circ v_0=u$~; comme $B$ est une $A$-algèbre formellement lisse et $B'$ un
$A$-module plat, on peut appliquer (19.7.1.5), d'où l'existence du
$A$-isomorphisme $v$ répondant à la question.
\end{proof}

\begin{remarks}[19.7.3]
\label{0.19.7.3-fr}
\begin{enumerate}
  \item[(i)] On notera que l'assertion d'unicité dans (19.7.2) est encore
  valable si l'on suppose seulement que $B$ et $B'$ sont complets \emph{pour
  les topologies $\mathfrak J$-préadiques}. Nous ignorons si l'on peut
  améliorer de même l'assertion d'existence, autrement dit si l'on peut se
  dispenser de supposer l'anneau local $B_0$ \emph{complet} (pour sa topologie
  $\mathfrak n_0$-préadique, en désignant par $\mathfrak n_0$ son idéal
  maximal) en exigeant seulement que $B$ soit complet \emph{pour la topologie
  $\mathfrak J$-préadique}. Lorsque $A_0$ est complet pour la topologie
  $\mathfrak m$-préadique, on peut voir que ce problème revient au suivant~:
  si $B_0$ est un anneau local noethérien régulier (non nécessairement
  complet) contenant le corps premier $\mathbf F_p=\mathbf Z/p\mathbf Z$,
  existe-t-il pour tout $n>1$ une $(\mathbf Z/p^n\mathbf Z)$-algèbre plate $B$
  telle que $B/pB$ soit isomorphe à $B_0$~?
  \item[(ii)] On notera qu'en général, l'isomorphisme $v$ dont l'existence est
  affirmée dans (19.7.2) \emph{n'est pas unique} (cf.~(19.8.7)).
\end{enumerate}
\end{remarks}

\subsection{Algèbres de Cohen et $p$-anneaux de Cohen; application à la
structure des anneaux locaux complets}
\label{subsection:0.19.8-fr}

Les résultats de cette section sont des applications immédiates des théorèmes
de (19.7), mais méritent d'être explicités en raison de leur importance
pratique.

\begin{definition}[19.8.1]
\label{0.19.8.1-fr}
Soient $A$, $B$ deux anneaux locaux noethériens, $\mathfrak m$ l'idéal maximal
de $A$, $k=A/\mathfrak m$ son corps résiduel,
$\varphi:A\to B$ un homomorphisme local, faisant de $B$ une $A$-algèbre. On
dit que $B$ est une \emph{$A$-algèbre de Cohen} si elle vérifie les conditions
suivantes~:
\begin{enumerate}
  \item[(i)] $B$ est un anneau \emph{complet}.
  \item[(ii)] $B$ est un $A$-module \emph{plat}.
  \item[(iii)] $B\otimes_A k$ est un \emph{corps} (autrement dit,
  $\mathfrak mB$ est l'idéal maximal de $B$) qui est une \emph{extension
  séparable} de $k$.
\end{enumerate}
\end{definition}

\oldpage[0\textsubscript{IV-1}]{110}

\begin{theorem}[19.8.2]
\label{0.19.8.2-fr}
Soient $A$ un anneau local noethérien, $k$ son corps résiduel.
\begin{enumerate}
  \item[(i)] Si $B$ est une $A$-algèbre de Cohen, $B$ est une $A$-algèbre
  formellement lisse. Pour tout anneau local noethérien complet $C$, tout
  homomorphisme local $A\to C$ et tout idéal $\mathfrak J\neq C$ dans $C$,
  tout $A$-homomorphisme $B\to C/\mathfrak J$ se factorise donc en
  $B\xrightarrow{v}C\to C/\mathfrak J$, où $v$ est un $A$-homomorphisme
  (nécessairement local).
  \item[(ii)] Pour tout corps $K$, extension séparable de $k$, il existe une
  $A$-algèbre de Cohen $B$ telle que $B\otimes_A k$ soit $k$-isomorphe à $K$,
  et une telle $A$-algèbre est unique à isomorphisme près.
\end{enumerate}
\end{theorem}

\begin{proof}
Comme $K$ est une $k$-algèbre formellement lisse (19.6.1), l'assertion (i)
résulte de (19.7.1). Pour prouver (ii), on peut se borner au cas où $A$ est
\emph{complet}, car il revient au même de dire que $B$ est un $A$-module plat
ou un $\widehat A$-module plat ($\mathbf 0_{\mathrm{III}}$, 10.2.3), on a
$\mathfrak mB=\mathfrak m\widehat AB$ et $k$ est le corps résiduel de
$\widehat A$. Il suffit alors d'appliquer (19.7.2) en prenant
$\mathfrak J=\mathfrak m$ et $B_0=K$ (et utilisant (19.6.1)).
\end{proof}

\begin{definition}[19.8.3]
\label{0.19.8.3-fr}
On appelle \emph{anneau local premier} un anneau local de la forme
$\mathbf Z_{p\mathbf Z}$, où $p\mathbf Z$ est un idéal premier de $\mathbf Z$.
On appelle \emph{anneau local complet premier} le complété d'un anneau local
premier.
\end{definition}

Les anneaux locaux premiers sont donc de deux sortes~:
\begin{enumerate}
  \item[$1^{\mathrm o}$] Ceux qui correspondent aux idéaux maximaux
  $p\mathbf Z$ où $p\neq0$ est un nombre premier~; $\mathbf Z_{p\mathbf Z}$
  est un anneau de valuation discrète, dont le complété est \emph{l'anneau des
  entiers $p$-adiques} noté d'ordinaire $\mathbf Z_p$\footnote{Cette notation,
  actuellement universellement utilisée, est dans ce cas en conflit avec la
  notation $A_f$ adoptée dans ($\mathbf 0_{\mathrm I}$, 1.2.3)~: avec
  $A=\mathbf Z$ et $f=p$, $A_f$ signifie en effet l'anneau des nombres
  rationnels de la forme $k/p^n$ ($k\in\mathbf Z$, $n$ entier $\geq0$)~; nous
  éviterons toujours d'utiliser la notation $\mathbf Z_p$ pour désigner ce
  dernier anneau.}.
  \item[$2^{\mathrm o}$] Pour l'idéal premier $p\mathbf Z=(0)$,
  $\mathbf Z_{p\mathbf Z}$ est le corps des nombres rationnels $\mathbf Q$,
  identique à son complété (la topologie étant naturellement la topologie
  d'anneau local noethérien, donc ici la topologie discrète).
\end{enumerate}

La terminologie de (19.8.3), analogue à celle des «~corps premiers~», se
justifie de la même manière~: pour tout anneau local $A$, considérons
l'homomorphisme canonique $\varphi:\mathbf Z\to A$, et soit $p\mathbf Z$
l'image réciproque par cet homomorphisme de l'idéal maximal $\mathfrak m$ de
$A$~; $p\mathbf Z$ est un idéal premier de $\mathbf Z$ et l'homomorphisme
précédent se factorise donc en
$\mathbf Z\to\mathbf Z_{p\mathbf Z}\xrightarrow{\psi}A$~; d'ailleurs, comme
$\varphi$ est l'\emph{unique} homomorphisme de $\mathbf Z$ dans $A$, $p$ et
$\psi$ sont déterminés de façon \emph{unique}. Autrement dit, pour tout anneau
local $A$, il y a un \emph{unique homomorphisme} $\psi:P\to A$, où $P$ est un
anneau local premier~; si en outre $A$ est séparé et complet, on peut prolonger
par complétion cet homomorphisme, et il y a donc un \emph{unique
homomorphisme} $\widehat\psi:\widehat P\to A$, où $\widehat P$ est un anneau
local complet premier. D'ailleurs, par passage aux quotients, $\psi$ donne un
homomorphisme du corps résiduel $\mathbf F_p$ si $p>0$ (resp. $\mathbf Q$ si
$p=0$) dans le corps résiduel $k$ de $A$, et $p$ est donc la
\emph{caractéristique} de $k$.

Si l'on prend en particulier pour $A$ un anneau local premier (resp. complet
premier), on voit qu'il n'existe dans un tel anneau qu'\emph{un seul
endomorphisme}, savoir l'identité.

% FIN SOURCE ega4/ega0-19-fr.tex

%
% DEBUT SOURCE ega4/ega0-19b-fr.tex
\oldpage[0\textsubscript{IV-1}]{111}

\begin{definition}[19.8.4]
\label{0.19.8.4-fr}
Soient $A$ un anneau local, $P\to A$ l'unique homomorphisme d'un anneau local
premier $P$ dans $A$, $p$ la caractéristique des corps résiduels de $P$ et
$A$. On dit que $A$ est un \emph{anneau de Cohen} si c'est une $P$-algèbre de
Cohen, c'est-à-dire (19.8.1) si~:
\begin{enumerate}
  \item[$1^{\mathrm o}$] $A$ est noethérien et complet.
  \item[$2^{\mathrm o}$] $A$ est un $P$-module plat (ce qui revient aussi à
  dire que $A$ est un $\widehat P$-module plat (Bourbaki, \emph{Alg. comm.},
  chap.~III, §~5, n\textsuperscript{o}~4, prop.~4)).
  \item[$3^{\mathrm o}$] $A/pA$ est un corps (nécessairement séparable sur le
  corps résiduel de $P$, ce corps étant premier).
\end{enumerate}
\end{definition}

Si $p=0$, ces conditions équivalent à dire que $A$ est un \emph{corps de
caractéristique $0$}. Si $p>0$, on a nécessairement $pA\neq0$~; la condition
$3^{\mathrm o}$ signifie que $pA$ est l'idéal maximal $\mathfrak m$ de $A$~;
la condition $2^{\mathrm o}$ signifie que $p$ est $A$-régulier, puisque $P$
est un anneau de valuation discrète ($\mathbf 0_{\mathrm I}$, 6.3.4). Donc $A$
est un anneau régulier (17.1.1, $d)$) de dimension $1$, et par suite un anneau
de valuation discrète, complet en vertu de $1^{\mathrm o}$~; en résumé~:

\begin{proposition}[19.8.5]
\label{0.19.8.5-fr}
Les anneaux de Cohen sont les corps de caractéristique $0$ et les anneaux de
valuation discrète complets, dont le corps résiduel a une caractéristique
$p>0$, et dont l'idéal maximal est engendré par $p\cdot1$ ($1$ étant l'unité de
l'anneau).
\end{proposition}

On notera que dans le second cas, $p\cdot1\neq0$ puisque $p$ est $A$-régulier,
donc on peut identifier $p\cdot1$ à l'entier $p$~; l'homomorphisme canonique
$\mathbf Z_p\to A$ est injectif, et on identifie $p\cdot1$ à l'élément $p$ de
$\mathbf Z_p$~; on dit dans ce cas que $A$ est un \emph{$p$-anneau de Cohen}.

\begin{theorem}[19.8.6]
\label{0.19.8.6-fr}
\textup{(Cohen).} ---
\begin{enumerate}
  \item[(i)] Soient $W$ un anneau de Cohen, $C$ un anneau local noethérien
  complet, $\mathfrak J$ un idéal de $C$ distinct de $C$. Alors tout
  homomorphisme local $u:W\to C/\mathfrak J$ se factorise en
  $W\xrightarrow{v}C\to C/\mathfrak J$, où $v$ est un homomorphisme local.
  \item[(ii)] Soit $K$ un corps. Il existe un anneau de Cohen $W$ dont le corps
  résiduel est isomorphe à $K$. Si $W'$ est un second anneau de Cohen, $K'$ son
  corps résiduel, tout isomorphisme $u:K\xrightarrow{\sim}K'$ provient par
  passage aux quotients d'un isomorphisme $v:W\xrightarrow{\sim}W'$.
\end{enumerate}
\end{theorem}

\begin{proof}
Cela n'est autre que (19.8.2) appliqué au cas où $A$ est un anneau local
premier.
\end{proof}

\begin{remarks}[19.8.7]
\label{0.19.8.7-fr}
\begin{enumerate}
  \item[(i)] Lorsque $K$ est de caractéristique $0$, la partie (ii) de (19.8.6)
  devient triviale.
  \item[(ii)] L'homomorphisme $v$ de (19.8.6, (i)) \emph{n'est pas
  nécessairement déterminé de façon unique par $u$}, comme le montre déjà le
  cas où $W$ est un corps de caractéristique $0$, $\mathfrak J^2=0$ et $u$ est
  un isomorphisme (cf.~(21.5.5)). De même, dans (19.8.6, (ii)) l'isomorphisme
  $v$ \emph{n'est pas nécessairement déterminé de façon unique par $u$}
  (cf.~(21.5.5)).

  Toutefois, lorsque $K$ est \emph{parfait} et de caractéristique $p>0$, on
  verra (21.5.5) que dans (19.8.6, (ii)) l'isomorphisme $v$ est unique. On verra
  aussi plus tard que dans ce cas $W$ s'identifie à l'anneau $W_\infty(K)$ des
  \emph{vecteurs de Witt de longueur infinie} sur $K$.
  \item[(iii)] Dans (19.8.6, (i)), on peut affaiblir les hypothèses sur $C$ en
  utilisant (19.3.10) et (19.3.12).
\end{enumerate}
\end{remarks}

\begin{theorem}[19.8.8]
\label{0.19.8.8-fr}
\textup{(Cohen).} --- Soient $A$ un anneau local noethérien complet, $k$ son
corps résiduel.

\oldpage[0\textsubscript{IV-1}]{112}

\begin{enumerate}
  \item[(i)] Il existe un anneau de Cohen $W$ tel que $A$ soit isomorphe à un
  anneau quotient d'un anneau de séries formelles
  $W[[T_1,\ldots,T_n]]$ (et en particulier $A$ est isomorphe à un quotient d'un
  anneau local régulier complet (17.3.8)). Si $A$ contient un corps, il est
  isomorphe à un anneau quotient de $k[[T_1,\ldots,T_n]]$.
  \item[(ii)] Supposons en outre $A$ intègre. Alors il existe un sous-anneau
  $B$ de $A$, tel que~: $1^{\mathrm o}$ $B$ est isomorphe à un anneau de séries
  formelles sur un anneau $C$ qui est un corps ou un anneau de Cohen (ce qui
  entraîne que $B$ est un anneau local régulier et complet (17.3.8))~;
  $2^{\mathrm o}$ $B$ a même corps résiduel que $A$ et l'injection $B\to A$ est
  un homomorphisme local~; $3^{\mathrm o}$ $A$ est une $B$-algèbre finie.
\end{enumerate}
\end{theorem}

\begin{proof}
Soit $\mathfrak m$ l'idéal maximal de $A$. Il existe un anneau de Cohen $W$ dont
le corps résiduel est isomorphe à $k$ (19.8.6, (ii))~; on a donc un
homomorphisme local $W\to A/\mathfrak m$, qui par suite se factorise en
$W\xrightarrow{u}A\to A/\mathfrak m$, où $u$ est un homomorphisme local
(19.8.6, (i)). Pour toute famille finie $(x_i)_{1\leq i\leq n}$ d'éléments de
$\mathfrak m$, il existe alors un homomorphisme local
$v:W[[T_1,\ldots,T_n]]\to A$ prolongeant $u$ et tel que $v(T_i)=x_i$ pour tout
$i$ (Bourbaki, \emph{Alg. comm.}, chap.~III, §~4, n\textsuperscript{o}~5,
prop.~6). Lorsque $A$ contient un corps, il contient un corps premier $P$, dont
$k$ est une extension (nécessairement séparable), et par suite $A$ contient un
corps isomorphe à $k$ (19.6.2)~; on peut alors remplacer $W$ par $k$ dans la
définition précédente de $v$.

\medskip
\noindent(i)\enspace
Prenons d'abord pour les $x_i$ un système de générateurs de $\mathfrak m$. Comme
$W$ a même corps résiduel que $A$, et que les classes des $x_i$ dans l'anneau
gradué $\gr_\bullet(A)$ engendrent $\gr_\bullet(A)$ en tant que $k$-algèbre,
$\gr(v):\gr_\bullet(W[[T_1,\ldots,T_n]])\to\gr_\bullet(A)$ est surjectif~; on
en déduit que $v$ est lui-même surjectif (Bourbaki, \emph{Alg. comm.},
chap.~III, §~2, n\textsuperscript{o}~8, cor.~2 du th.~1). Rappelons que le cas
où $A$ contient un corps a déjà été vu et ne figure ici que pour mémoire
(19.6.3).

\medskip
\noindent(ii)\enspace
Si $A$ contient un corps, il contient un corps $k'$ isomorphe à $k$ comme on
l'a vu~; on considère alors un système de paramètres
$(y_j)_{1\leq j\leq m}$ de $A$ (16.3.6), on prend
$B=k[[T_1,\ldots,T_m]]$ et on considère l'homomorphisme local $w:B\to A$ qui
coïncide dans $k$ avec un isomorphisme $k\to k'$ et qui est tel que
$w(T_j)=y_j$ pour $1\leq j\leq m$. Si $A$ ne contient pas de corps, l'unique
homomorphisme $\mathbf Z_{p\mathbf Z}\to A$ (19.8.3) est nécessairement
injectif (sans quoi, comme $A$ est intègre, son noyau serait l'idéal maximal de
$\mathbf Z_{p\mathbf Z}$ et son image isomorphe à un corps)~; en outre, on a
alors $p>0$ par hypothèse, $\mathbf Z_{p\mathbf Z}$ étant un corps si $p=0$.
L'élément $p\cdot1$ de $A$ (identifié à $p$) est non diviseur de zéro dans $A$,
et est contenu dans $\mathfrak m$, donc (16.3.4 et 16.3.7) il existe une famille
$(z_i)_{1\leq i\leq m-1}$ qui, avec $p$, forme un système de paramètres de $A$.
L'anneau de Cohen $W$ considéré au début de la démonstration est alors un
anneau de valuation discrète de corps résiduel $k$, dans lequel $p$ engendre
l'idéal maximal (19.8.5), et l'homomorphisme unique $u:W\to A$ défini au début
applique $p$ sur lui-même. On prend alors $B=W[[T_1,\ldots,T_{m-1}]]$ et on
considère l'homomorphisme local $w:B\to A$ qui coïncide avec $u$ dans $W$ et
est tel que $w(T_i)=z_i$ pour $1\leq i\leq m-1$. Dans les deux cas, si
$\mathfrak n$ est l'idéal maximal de $B$, il est clair que $\mathfrak nA$ est
un idéal de définition de $A$~; comme en outre $B/\mathfrak n$ et
$A/\mathfrak m$ sont isomorphes, $A$ est un $B$-module quasi-fini
($\mathbf 0_{\mathrm I}$, 7.4.4), donc un $A$-module de type fini puisque $B$
est complet et $A$ séparé pour les topologies $\mathfrak n$-préadiques
($\mathbf 0_{\mathrm I}$, 7.4.1). D'autre part,

\oldpage[0\textsubscript{IV-1}]{113}

dans les deux cas, on a $\dim(B)=\dim(A)=m$~; dans le premier cas, cela résulte
de (17.1.4, (iii))~; dans le second, on voit directement que $p$ et les $T_i$
($1\leq i\leq m-1$) forment une suite $B$-régulière engendrant $\mathfrak n$,
ou on peut aussi utiliser le fait que ces éléments engendrent $\mathfrak n$ et
que l'on a $\dim(B)\geq\dim(A)$ par (16.3.10). Comme $A$ et $B$ sont intègres,
on tire finalement de (16.3.10) que $w$ est injectif, ce qui achève la
démonstration.
\end{proof}

\begin{corollary}[19.8.9]
\label{0.19.8.9-fr}
Soit $A$ un anneau local noethérien intègre complet contenant un corps $k_0$~;
soit $k$ le corps résiduel de $A$, et supposons que $k$ soit fini sur $k_0$.
Alors, dans la conclusion de (19.8.8, (ii)), on peut remplacer $1^{\mathrm o}$
et $2^{\mathrm o}$ par la condition que $B$ est de la forme
$k_0[[T_1,\ldots,T_n]]$, l'injection canonique $B\to A$ étant un
$k_0$-homomorphisme local (pour la structure usuelle de $k_0$-algèbre de $B$).
\end{corollary}

\begin{proof}
En effet, en reprenant la démonstration de (19.8.8, (ii)), on définit cette fois
$w:k_0[[T_1,\ldots,T_n]]\to A$ comme coïncidant dans $k_0$ avec l'identité et
appliquant $T_j$ sur $y_j$ pour $1\leq j\leq n$. L'hypothèse que $k$ est de
degré fini sur $k_0$ entraîne encore que $A$ est un $B$-module quasi-fini, donc
de type fini par ($\mathbf 0_{\mathrm I}$, 7.4.1), et on conclut comme dans
(19.8.8).
\end{proof}

\begin{corollary}[19.8.10]
\label{0.19.8.10-fr}
Soit $A$ un anneau local artinien dont l'idéal maximal $\mathfrak m$ est de
carré nul~; il existe alors un anneau local noethérien régulier $B$, d'idéal
maximal $\mathfrak n$, tel que $A$ soit isomorphe à $B/\mathfrak n^2$.
\end{corollary}

\begin{proof}
Soient $K$ le corps résiduel $A/\mathfrak m$ de $A$, $n$ le rang de
$\mathfrak m/\mathfrak m^2=\mathfrak m$ sur $K$. Si $A$ contient un corps, il
résulte de (19.6.3) que $A$ est isomorphe à $B/\mathfrak b$, où
$B=K[[T_1,\ldots,T_n]]$ et $\mathfrak b$ est contenu dans le carré de l'idéal
maximal $\mathfrak n$ de $B$~; mais comme
$\operatorname{long}(B/\mathfrak n^2)=n+1=\operatorname{long}(A)$, on a
nécessairement $\mathfrak b=\mathfrak n^2$.

Supposons ensuite que $A$ ne contienne pas de corps~; cela entraîne que $K$ est
de caractéristique $p>0$ et que $p\cdot1\neq0$ dans $A$ (19.6.3)~; donc
$p\cdot1$ est un élément de $\mathfrak m$, et il y a par suite $n-1$ autres
éléments $x_i$ ($2\leq i\leq n$) de $\mathfrak m$ formant avec $p\cdot1$ une
base de $\mathfrak m$ sur $K$. Soit $W$ un anneau de Cohen dont le corps
résiduel est isomorphe à $K$~; $W$ est un anneau de valuation discrète dont $p$
engendre l'idéal maximal~; on a vu dans la démonstration de (19.8.8) qu'il y a
un homomorphisme $u:W\to A$ appliquant $p$ sur lui-même et qui par passage aux
quotients donne l'identité sur $K$. On prend $B=W[[T_2,\ldots,T_n]]$ et l'on
considère l'homomorphisme local $w:B\to A$ qui coïncide avec $u$ dans $W$ et
est tel que $w(T_i)=x_i$ pour $2\leq i\leq n$. Il est clair que $w$ est
surjectif et que son noyau $\mathfrak b$ est contenu dans le carré de l'idéal
maximal $\mathfrak n=pB+BT_2+\cdots+BT_n$ de $B$~; comme
$\operatorname{long}(B/\mathfrak n^2)=n+1=\operatorname{long}(A)$, on a encore
$\mathfrak b=\mathfrak n^2$.
\end{proof}

\begin{proposition}[19.8.11]
\label{0.19.8.11-fr}
Soient $A$ un anneau local artinien, $\mathfrak m$ son idéal maximal, $k$ son
corps résiduel. Pour que $A$ soit isomorphe à un anneau quotient d'un anneau de
Cohen, il faut et il suffit que $\mathfrak m$ soit engendré par $p\cdot1$, où
$p$ est la caractéristique de $k$.
\end{proposition}

\begin{proof}
La condition est évidemment nécessaire (19.8.5). Pour voir qu'elle est
suffisante, on observe, comme au début de la démonstration de (19.8.8), qu'il
existe un anneau de Cohen $W$ dont le corps résiduel est isomorphe à $k$ et un
homomorphisme local $u:W\to A$. En outre, si l'on considère l'homomorphisme
composé $\mathbf Z_{p\mathbf Z}\to W\xrightarrow{u}A$ (qui est nécessairement
l'unique homomorphisme de $\mathbf Z_{p\mathbf Z}$ dans $A$), on voit que
l'image

\oldpage[0\textsubscript{IV-1}]{114}

par $u$ de l'élément $p\cdot1$ de $W$ est l'élément $p\cdot1$ de $A$~; comme
l'élément $p\cdot1$ de $W$ engendre l'idéal maximal de cet anneau, on déduit
aussitôt de l'hypothèse que
$\gr(u):\gr_\bullet(W)\to\gr_\bullet(A)$ est surjectif, et par suite il en est
de même de $u$ (Bourbaki, \emph{Alg. comm.}, chap.~III, §~2,
n\textsuperscript{o}~8, cor.~2 du th.~1).
\end{proof}

\subsection{Algèbres relativement formellement lisses}
\label{subsection:0.19.9-fr}

\begin{definition}[19.9.1]
\label{0.19.9.1-fr}
Soient $\Lambda$ un anneau topologique, $A$ une $\Lambda$-algèbre topologique,
$B$ une $A$-algèbre topologique. On dit que $B$ est une $A$-algèbre
\emph{formellement lisse relativement à $\Lambda$} si, pour toute $A$-algèbre
topologique discrète $C$, et tout idéal nilpotent $\mathfrak J$ de $C$, tout
$A$-homomorphisme continu $u_0:B\to C/\mathfrak J$ qui se factorise en
$B\xrightarrow{u}C\to C/\mathfrak J$, où $u$ est un $\Lambda$-homomorphisme
continu, se factorise aussi en $B\xrightarrow{v}C\to C/\mathfrak J$, où $v$
est un $A$-homomorphisme continu.
\end{definition}

Il résulte de cette définition que si $B$ est une $A$-algèbre \emph{formellement
lisse}, alors $B$ est aussi \emph{formellement lisse relativement à
$\Lambda$}, pour toute structure de $\Lambda$-algèbre topologique définie sur
$A$ (autrement dit, tout homomorphisme continu d'anneaux $\Lambda\to A$).

\begin{proposition}[19.9.2]
\label{0.19.9.2-fr}
Soient $\Lambda$ un anneau topologique, $A$ une $\Lambda$-algèbre topologique.
\begin{enumerate}
  \item[(i)] $A$ est une $\Lambda$-algèbre formellement lisse relativement à
  $\Lambda$.
  \item[(ii)] Si $B$ est une $A$-algèbre formellement lisse relativement à
  $\Lambda$ et $C$ une $B$-algèbre formellement lisse relativement à
  $\Lambda$, alors $C$ est une $A$-algèbre formellement lisse relativement à
  $\Lambda$.
  \item[(iii)] Soient $B$ une $A$-algèbre formellement lisse relativement à
  $\Lambda$, $A'$ une $A$-algèbre topologique, alors la $A'$-algèbre
  topologique $B\otimes_A A'$ est formellement lisse relativement à $\Lambda$.
  \item[(iv)] Soient $B$ une $A$-algèbre topologique, $S$ (resp. $T$) une
  partie multiplicative de $A$ (resp. $B$) telle que l'image canonique de $S$
  dans $B$ soit contenue dans $T$. Si $B$ est une $A$-algèbre formellement
  lisse relativement à $\Lambda$, alors $T^{-1}B$ est une $S^{-1}A$-algèbre
  formellement lisse relativement à $\Lambda$.
  \item[(v)] Soient $B_i$ ($1\leq i\leq n$) des $A$-algèbres topologiques.
  Pour que $\prod_{i=1}^n B_i$ soit une $A$-algèbre formellement lisse
  relativement à $\Lambda$, il faut et il suffit que chacune des $B_i$ le
  soit.
\end{enumerate}
\end{proposition}

\begin{proof}
L'assertion (i) est triviale, et la démonstration des autres est étroitement
calquée sur les démonstrations de (19.3.5)~; elle est donc laissée au lecteur.
\end{proof}

\begin{corollary}[19.9.3]
\label{0.19.9.3-fr}
Soient $A$ un anneau topologique, $A'$ et $B$ deux $A$-algèbres topologiques.
Alors la $A$-algèbre topologique $B'=B\otimes_A A'$ est une $B$-algèbre
formellement lisse relativement à $A$.
\end{corollary}

\begin{proof}
Cela résulte de (19.9.2, (i) et (iii)).
\end{proof}

\begin{proposition}[19.9.4]
\label{0.19.9.4-fr}
Soient $\Lambda$ un anneau topologique, $A$ une $\Lambda$-algèbre topologique,
$B$ une $A$-algèbre topologique. Les conditions suivantes sont équivalentes~:
\begin{enumerate}
  \item[a)] $B$ est une $A$-algèbre formellement lisse relativement à
  $\Lambda$.
  \item[b)] $\widehat B$ est une $A$-algèbre formellement lisse relativement à
  $\Lambda$.
  \item[c)] $\widehat B$ est une $\widehat A$-algèbre formellement lisse
  relativement à $\Lambda$.
  \item[d)] $\widehat B$ est une $\widehat A$-algèbre formellement lisse
  relativement à $\widehat\Lambda$.
\end{enumerate}
\end{proposition}

\begin{proof}
On laisse encore au lecteur la démonstration, calquée sur celle de (19.3.6).
\end{proof}

\oldpage[0\textsubscript{IV-1}]{115}

\begin{env}[19.9.5]
\label{0.19.9.5-fr}
De même, l'énoncé (19.3.8) est encore valable (avec la même démonstration) quand
on y remplace «~formellement lisse~» par «~formellement lisse relativement à
$\Lambda$~». Si dans l'énoncé de (19.3.10) on remplace «~formellement lisse~»
par «~formellement lisse relativement à $\Lambda$~» la conclusion est remplacée
par la suivante (la démonstration restant essentiellement inchangée)~: \emph{tout
$A$-homomorphisme continu $u_0:B\to C/\mathfrak J$ qui se factorise en
$B\xrightarrow{u}C\to C/\mathfrak J$, où $u$ est un $\Lambda$-homomorphisme
continu, se factorise aussi en $B\xrightarrow{v}C\to C/\mathfrak J$, où $v$
est un $A$-homomorphisme continu}.
\end{env}

\begin{env}[19.9.6]
\label{0.19.9.6-fr}
Les critères de lissité formelle (19.4.1) et (19.4.2) sont valables lorsqu'on y
remplace «~formellement lisse~» par «~formellement lisse relativement à
$\Lambda$~», les démonstrations restant pratiquement inchangées.
\end{env}

\begin{proposition}[19.9.7]
\label{0.19.9.7-fr}
Soient $\Lambda$ un anneau topologique, $A$ une $\Lambda$-algèbre topologique,
$B$ une $A$-algèbre topologique. Supposons que pour toute $A$-algèbre discrète
$C$ et tout idéal $\mathfrak J$ de $C$ tel que $\mathfrak J^2=0$, tout
$A$-homomorphisme continu $u_0:B\to C/\mathfrak J$ qui se factorise en
$B\xrightarrow{u}C\to C/\mathfrak J$, où $u$ est un $\Lambda$-homomorphisme
continu, se factorise aussi en $B\xrightarrow{v}C\to C/\mathfrak J$, où $v$
est un $A$-homomorphisme continu. Alors $B$ est une $A$-algèbre formellement
lisse relativement à $\Lambda$.
\end{proposition}

\begin{proof}
La démonstration de (19.4.3) se transcrit aussitôt.
\end{proof}

\begin{proposition}[19.9.8]
\label{0.19.9.8-fr}
Soient $\Lambda$ un anneau topologique, $A$ une $\Lambda$-algèbre topologique,
$B$ une $A$-algèbre topologique. Pour que $B$ soit une $A$-algèbre
formellement lisse relativement à $\Lambda$, il faut et il suffit que pour tout
$B$-module topologique discret $L$, annulé par un idéal ouvert de $B$, on ait
(cf.~18.4.2)
\[
  \operatorname{Exalcotop}_{A/\Lambda}(B,L)=0.
\]
\end{proposition}

\begin{proof}
Avec les notations de la démonstration de (19.4.4), il suffit de noter ici que
l'on peut supposer que l'extension $E_\Lambda$ de $B_\Lambda$ est
$\Lambda$-triviale~; le reste de la démonstration est alors inchangé.
\end{proof}

Lorsque $\Lambda$, $A$ et $B$ sont des anneaux discrets, le critère (19.9.8) se
réduit à
\begin{equation}
\tag{19.9.8.1}
\label{0.19.9.8.1-fr}
  \operatorname{Exalcom}_{A/\Lambda}(B,L)=0
  \quad\text{pour tout $B$-module $L$}~;
\end{equation}
autrement dit, \emph{toute $A$-extension commutative de $B$ par un $B$-module,
qui est $\Lambda$-triviale, est aussi $A$-triviale}.

\subsection{Algèbres formellement non ramifiées et algèbres formellement
étales}
\label{subsection:0.19.10-fr}
\label{0.19.10-fr}

\begin{env}[19.10.1]
\label{0.19.10.1-fr}
Nous aurons aussi à utiliser deux notions voisines de la notion d'algèbre
formellement lisse.
\end{env}

\begin{definition}[19.10.2]
\label{0.19.10.2-fr}
Soient $A$ un anneau topologique, $B$ une $A$-algèbre topologique. On dit que
$B$ est une $A$-algèbre \emph{formellement non ramifiée} si, pour tout
$A$-homomorphisme $p:E\to C$ de $A$-algèbres discrètes, dont le noyau est
nilpotent, et tout $A$-homomorphisme continu $u:B\to C$, il existe \emph{au plus
un} $A$-homomorphisme continu $v:B\to E$ tel que $u$ se factorise en
$B\xrightarrow{v}E\xrightarrow{p}C$. On dit que $B$ est une $A$-algèbre
\emph{formellement étale} si elle est formellement lisse et formellement non
ramifiée, autrement dit, si, pour tout $A$-homomorphisme surjectif $p:E\to C$
de $A$-algèbres discrètes, dont le noyau est nilpotent, et pour tout
$A$-homomorphisme continu $u:B\to C$, il existe \emph{un et un seul}
$A$-homomorphisme continu $v:B\to E$ tel que $u$ se factorise en
$B\xrightarrow{v}E\xrightarrow{p}C$.
\end{definition}

\oldpage[0\textsubscript{IV-1}]{116}

\begin{examples}[19.10.3]
\label{0.19.10.3-fr}
Nous nous bornerons aux anneaux \emph{discrets}.
\begin{enumerate}
  \item[(i)] Si l'homomorphisme structural $\rho:A\to B$ est \emph{surjectif},
  $B$ est une $A$-algèbre \emph{formellement non ramifiée}, car avec les
  notations de (19.10.2), l'homomorphisme composé
  $A\xrightarrow{j}B\xrightarrow{v}E$ doit être l'homomorphisme structural,
  donc $v$ est unique (s'il existe). Par contre, $B$ n'est formellement lisse
  que si $\rho$ est \emph{bijectif}.
  \item[(ii)] Soient $A$ un anneau, $S$ une partie multiplicative de $A$~;
  alors $B=S^{-1}A$ est une $A$-algèbre \emph{formellement étale}. En effet,
  soient $j:A\to B$ l'homomorphisme canonique, $\rho:A\to E$ un homomorphisme
  d'anneaux, $\mathfrak J$ un idéal nilpotent de $E$,
  $p:E\to C=E/\mathfrak J$ l'homomorphisme canonique. Si $u:B\to C$ est un
  $A$-homomorphisme, $u(s/1)$ est inversible dans $C$ pour $s\in S$ puisque
  $s/1$ est inversible dans $B$~; mais $u(s/1)=p(\rho(s))$, et comme
  $\mathfrak J$ est nilpotent, le fait que l'image de $\rho(s)$ dans
  $E/\mathfrak J$ soit inversible entraîne que $\rho(s)$ lui-même est
  inversible dans $E$ ($\mathbf 0_{\mathrm I}$, 7.1.12)~; on en conclut que
  $\rho$ se factorise d'une seule manière en
  $A\xrightarrow{j}B\xrightarrow{v}E$, et l'on a évidemment
  $u\circ j=(p\circ v)\circ j$, d'où $u=p\circ v$, ce qui prouve notre
  assertion.
  \item[(iii)] Une extension finie séparable $k'$ d'un corps $k$ est une
  $k$-algèbre formellement lisse (19.6.1)~; nous verrons plus loin que pour
  qu'une extension de \emph{type fini} $K$ d'un corps $k$ soit une $k$-algèbre
  formellement étale, il faut et il suffit qu'elle soit formellement non
  ramifiée, et cela équivaut aussi au fait que $K$ est une extension finie et
  séparable de $k$ (21.7.4).
  \item[(iv)] Une algèbre de polynômes $A[X_\alpha]_{\alpha\in I}$ sur un
  anneau $A$ est \emph{formellement lisse} sur $A$ (19.3.3) mais \emph{non
  formellement étale} si $I$ n'est pas vide, comme il résulte aussitôt des
  définitions.
\end{enumerate}
\end{examples}

\begin{remark}[19.10.4]
\label{0.19.10.4-fr}
Le même raisonnement que dans (19.4.3) montre que pour vérifier qu'une
$A$-algèbre topologique $B$ est formellement non ramifiée, on peut, dans la
définition (19.10.2), se borner au cas où le noyau de $p:E\to C$ est de
\emph{carré nul}~; on peut en outre se borner au cas où $u(B)\subset p(E)$
(sans quoi il n'y a aucun $v$ factorisant $u$) et, en remplaçant $C$ par $u(B)$
et $E$ par $p^{-1}(u(B))$, supposer $u$ et $p$ surjectifs.
\end{remark}

% FIN SOURCE ega4/ega0-19b-fr.tex

%
% DEBUT SOURCE ega4/ega0-20-fr.tex
\section{Dérivations et différentielles}
\label{section:0.20-fr}
\label{0.20-fr}

Les notions introduites dans ce paragraphe seront reprises sous forme
géométrique au chapitre~IV, §~16, et joueront un rôle important dans l'étude
des préschémas. Leur importance dans ce chapitre tient d'abord à leurs
relations avec la notion d'algèbre formellement lisse, et notamment aux
théorèmes (20.4.9), (20.5.7) et (20.5.12), qui seront traduits en langage
géométrique dans le paragraphe du chapitre~IV consacré aux morphismes lisses,
et sont d'un usage constant dans les applications. D'autre part, les notions
différentielles serviront à démontrer au §~22 d'importants critères de
régularité qui joueront un rôle essentiel dans l'étude approfondie des anneaux
locaux noethériens faite au §~7 du chapitre~IV.

\subsection{Dérivations et extensions d'algèbres}
\label{subsection:0.20.1-fr}

\oldpage[0\textsubscript{IV-1}]{117}
\begin{proposition}[20.1.1]
\label{0.20.1.1-fr}
Soient $A$ un anneau (non nécessairement commutatif), $B$, $E$, $C$ des
$A$-anneaux, $p:E\to C$ un $A$-homomorphisme dont le noyau $\mathfrak{Q}$
est de carré nul, $u:B\to C$ un $A$-homomorphisme. Supposons qu'il existe
un $A$-homomorphisme $v_0:B\to E$ tel que $u$ se factorise en
$B\xrightarrow{v_0}E\xrightarrow{p}C$. Alors l'ensemble des
$A$-homomorphismes $v:B\to E$ tels que $u=p\circ v$ est identique à
l'ensemble des applications $v_0+D$, où $D:B\to\mathfrak{Q}$ est une
application vérifiant les deux conditions suivantes :
\begin{enumerate}
  \item[\rm(i)] $D$ est un homomorphisme de $A$-bimodules.
  \item[\rm(ii)] Quels que soient $f$, $g$ dans $B$, on a
\end{enumerate}
\begin{equation}
\tag{20.1.1.1}
\label{0.20.1.1.1-fr}
  D(fg)=f.D(g)+D(f).g.
\end{equation}
\end{proposition}

\begin{proof}
Dire que $p(v(f))=p(v_0(f))$ pour $f\in B$ signifie que
$D(f)=v(f)-v_0(f)$ appartient à $\mathfrak{Q}$ ; si l'on écrit que
$v(fg)=v(f)v(g)$, on obtient la relation \sref{0.20.1.1.1-fr},
$\mathfrak{Q}$ étant de carré nul, et la condition (i) résulte de ce que $v$
et $v_0$ sont des $A$-homomorphismes.

Si $\rho:A\to B$ est l'homomorphisme structural, on tire de
\sref{0.20.1.1.1-fr} que l'on a
$D(\rho(a)f)=\rho(a)D(f)+D(\rho(a))f$ pour tout $a\in A$ ; mais on doit
avoir $D(\rho(a)f)=\rho(a)D(f)$ d'après (i), donc, en faisant $f=1$, il
vient
\[
  D(\rho(a))=0 \qquad\text{pour}\quad a\in A\,;
\]
réciproquement, si $D$ est nulle dans $\rho(A)$ et vérifie
\sref{0.20.1.1.1-fr}, elle vérifie aussi (i).
\end{proof}

\begin{definition}[20.1.2]
\label{0.20.1.2-fr}
Étant donnés un $A$-anneau $B$ et un $B$-bimodule $L$, on appelle
\emph{$A$-dérivation de $B$ dans $L$} une application $D:B\to L$ vérifiant
les conditions (i) et (ii) de \sref{0.20.1.1-fr}.
\end{definition}

Il résulte de \sref{0.20.1.1-fr} que le \emph{noyau} d'une $A$-dérivation
$D:B\to L$ est un sous-$A$-anneau de $B$.

On appelle parfois \emph{dérivation} de $B$ dans $L$ une
$\mathbf{Z}$-dérivation, c'est-à-dire une application \emph{additive} de $B$
dans $L$ vérifiant \sref{0.20.1.1.1-fr} ; on a donc $D(1)=0$ pour une telle
application. Si $B$ est une algèbre sur un \emph{corps premier} $P$, toute
$\mathbf{Z}$-dérivation de $B$ est une $P$-dérivation : c'est évident d'après
ce qui précède si $P$ est de caractéristique $>0$, et dans le cas contraire,
la relation $n.n^{-1}=1$ pour $n\in\mathbf{Z}$ donne
\[
  nD(n^{-1})=0
\]
donc $D(n^{-1})=0$ puisque $P$ est de caractéristique $0$.

Il résulte aussitôt de la définition \sref{0.20.1.2-fr} que si $D$, $D'$ sont
deux $A$-dérivations de $B$ dans $L$, il en est de même de $D-D'$. Autrement
dit, l'ensemble des $A$-dérivations de $B$ dans $L$ est muni d'une structure
de \emph{groupe additif} ; on note ce groupe
$\operatorname{D\acute{e}r}_A(B,L)$. Si $A$ est \emph{commutatif}, $B$ une
$A$-\emph{algèbre commutative}, et $L$ un $B$-module, alors, pour toute
$A$-dérivation $D$ de $B$ dans $L$ et tout $a\in A$, $aD$ est encore une
$A$-dérivation de $B$ dans $L$, autrement dit $\operatorname{D\acute{e}r}_A(B,L)$
est alors muni d'une structure de $A$-\emph{module}.

\oldpage[0\textsubscript{IV-1}]{118}
La prop. \sref{0.20.1.1-fr} s'interprète de la façon suivante :

\begin{corollary}[20.1.3]
\label{0.20.1.3-fr}
Étant donnés deux $A$-homomorphismes de $A$-anneaux $p:E\to C$,
$u:B\to C$ dont le premier a un noyau $\mathfrak{Q}$ de carré nul,
l'ensemble des $A$-homomorphismes $v:B\to E$ tels que $u=p\circ v$ est vide
ou est un espace homogène principal pour le groupe
$\operatorname{D\acute{e}r}_A(B,\mathfrak{Q})$.
\end{corollary}

En particulier :

\begin{corollary}[20.1.4]
\label{0.20.1.4-fr}
Soient $A$ un anneau, $C$ un $A$-anneau, $L$ un $C$-bimodule, $E$ une
$A$-extension de $C$ par $L$, $p:E\to C$ l'augmentation. L'application qui,
à toute dérivation $D\in\operatorname{D\acute{e}r}_A(C,L)$, associe l'application
$x\mapsto x+D(p(x))$ est un isomorphisme du groupe
$\operatorname{D\acute{e}r}_A(C,L)$ sur le groupe des $A$-équivalences de $E$ sur
elle-même.
\end{corollary}

\begin{proof}
Appliquons \sref{0.20.1.1-fr} pour $B=E$, $v_0=\mathrm{I}_E$ : l'ensemble
des $A$-homomorphismes $v:E\to E$ tels que $p\circ v=p$ est identique à
l'ensemble des applications $v=\mathrm{I}_E+D'$, où
$D'\in\operatorname{D\acute{e}r}_A(E,L)$. Dire qu'un tel $A$-homomorphisme $v$ est
une $A$-équivalence équivaut à dire que $v$ se réduit à l'injection canonique
dans $L$, ou encore que $D'(x)=0$ dans $L$ ; mais cela signifie aussi que
$D'$ se factorise en $E\xrightarrow{p}C\xrightarrow{D}L$, où $D$ est une
$A$-dérivation. D'où le corollaire.
\end{proof}

Pour les extensions triviales, \sref{0.20.1.3-fr} donne :

\begin{corollary}[20.1.5]
\label{0.20.1.5-fr}
Soient $A$ un anneau, $B$, $C$ deux $A$-anneaux, $L$ un $C$-bimodule,
$u:B\to C$ un $A$-homomorphisme ; l'application qui, à toute dérivation
$D\in\operatorname{D\acute{e}r}_A(B,L)$, associe l'application
$\alpha:x\mapsto(u(x),D(x))$ est une bijection sur l'ensemble des
$A$-homomorphismes $v:B\to D_C(L)$ (cf. 18.2.3) tels que $u$ se factorise en
$B\xrightarrow{v}D_C(L)\to C$.
\end{corollary}

Plus particulièrement :

\begin{corollary}[20.1.6]
\label{0.20.1.6-fr}
Soient $A$ un anneau, $B$ un $A$-anneau, $L$ un $B$-bimodule. Si, à toute
dérivation $D\in\operatorname{D\acute{e}r}_A(B,L)$, on associe : $1^\circ$ la
$A$-équivalence $(x,y)\mapsto(x,y+D(x))$ de l'extension $D_B(L)$ sur
elle-même ; $2^\circ$ le $A$-homomorphisme $x\mapsto(x,D(x))$ de $B$ dans
$D_B(L)$, inverse à droite de l'homomorphisme d'augmentation $D_B(L)\to B$,
on définit des correspondances biunivoques canoniques entre :
\begin{enumerate}
  \item[\rm(i)] l'ensemble $\operatorname{D\acute{e}r}_A(B,L)$ ;
  \item[\rm(ii)] l'ensemble des $A$-équivalences de $D_B(L)$ sur elle-même ;
  \item[\rm(iii)] l'ensemble des $A$-homomorphismes inverses à droite de
  l'homomorphisme d'augmentation $D_B(L)\to B$.
\end{enumerate}
\end{corollary}

On notera que la correspondance biunivoque ainsi établie entre (i) et (ii)
respecte les structures de groupe, et que celle qu'on en déduit entre (ii) et
(iii) n'est autre que la correspondance biunivoque déjà définie dans (18.3.8).

\begin{corollary}[20.1.7]
\label{0.20.1.7-fr}
Soient $A$ un anneau, $C$ un $A$-anneau, $L$ un $A$-bimodule, $E$ une
$A$-extension triviale de $C$ par $L$, $p:E\to C$ l'augmentation. L'ensemble
des $A$-dérivations $D$ de $E$ dans $L$ telles que $D(x)=x$ dans $L$ est
identique à l'ensemble des applications $x\mapsto x-w(p(x))$, où $w$
parcourt l'ensemble des $A$-homomorphismes inverses à droite de $p$.
\end{corollary}

\begin{proof}
Il suffit encore d'appliquer \sref{0.20.1.1-fr} pour $B=E$,
$v_0=\mathrm{I}_E$ ; si $v=\mathrm{I}_E-D$, la condition $D(x)=x$ pour
$x\in L$ équivaut à $v(x)=0$ pour $x\in L$, c'est-à-dire à $v=w\circ p$,
où $w:C\to E$ est un $A$-homomorphisme ; en outre la condition
$p\circ v=p$ équivaut à $p\circ w=\mathrm{I}_C$, autrement dit au fait que
$w$ est inverse à droite de $p$.
\end{proof}

\subsection{Propriétés fonctorielles des dérivations}
\label{subsection:0.20.2-fr}

\oldpage[0\textsubscript{IV-1}]{119}
\begin{env}[20.2.1]
\label{0.20.2.1-fr}
Soient $A$ un anneau, $B$ un $A$-anneau, $L$ un $B$-bimodule ; si $L'$ est
un second $B$-bimodule et $w:L\to L'$ un homomorphisme de $B$-bimodules, il
est clair que l'application $D\mapsto w\circ D$ est un homomorphisme de
groupes additifs
\begin{equation}
\tag{20.2.1.1}
\label{0.20.2.1.1-fr}
  w_0:\operatorname{D\acute{e}r}_A(B,L)\to\operatorname{D\acute{e}r}_A(B,L')
\end{equation}
et que si $w':L'\to L''$ est un second homomorphisme de $B$-bimodules, on a
$(w'\circ w)_0=w'_0\circ w_0$. Lorsque $A$ est commutatif, $B$ une
$A$-algèbre commutative et $L$ un $B$-module, \sref{0.20.2.1.1-fr} est un
homomorphisme de $A$-modules.

En second lieu, soient $B'$ un $A$-anneau, $v:B'\to B$ un
$A$-homomorphisme qui fait de $L$ un $B'$-bimodule ; alors l'application
$D\mapsto D\circ v$ est un homomorphisme de groupes additifs,
\begin{equation}
\tag{20.2.1.2}
\label{0.20.2.1.2-fr}
  v^0:\operatorname{D\acute{e}r}_A(B,L)\to\operatorname{D\acute{e}r}_A(B',L)
\end{equation}
comme il résulte de \sref{0.20.1.1.1-fr} ; si $v':B''\to B'$ est un second
$A$-homomorphisme, on a $(v\circ v')^0=v'^0\circ v^0$. Lorsque $A$, $B$ et
$B'$ sont commutatifs et $L$ un $B$-module, \sref{0.20.2.1.2-fr} est un
homomorphisme de $A$-modules.

Enfin, soit $u:A'\to A$ un homomorphisme d'anneaux faisant de $B$ un
$A'$-anneau ; toute $A$-dérivation $D\in\operatorname{D\acute{e}r}_A(B,L)$ est aussi
une $A'$-dérivation, d'où une injection canonique de groupes commutatifs
\begin{equation}
\tag{20.2.1.3}
\label{0.20.2.1.3-fr}
  u^0:\operatorname{D\acute{e}r}_A(B,L)\to\operatorname{D\acute{e}r}_{A'}(B,L)
\end{equation}
et si $u':A''\to A'$ est un second homomorphisme d'anneaux, on a
$(u\circ u')^0=u'^0\circ u^0$ ; lorsque $A$, $A'$ et $B$ sont commutatifs et
$L$ un $B$-module, \sref{0.20.2.1.3-fr} est un di-homomorphisme de modules
(relatif à $u$).

On peut encore dire que
\[
  (A,B,L)\leadsto\operatorname{D\acute{e}r}_A(B,L)
\]
est un \emph{foncteur covariant} de la catégorie $\boldsymbol{K}$ définie dans
(18.3.5) dans la catégorie $\boldsymbol{Ab}$ des groupes commutatifs, en
faisant correspondre, à tout triplet $(u,v,w)$ constituant un morphisme de
$\boldsymbol{K}$ l'homomorphisme $w_0\circ v^0\circ u^0$ ; la vérification de
la fonctorialité résulte de la commutativité des diagrammes
\[
\xymatrix{
  \operatorname{D\acute{e}r}_A(B,L) \ar[r]^{v^0} \ar[d]_{w_0} &
    \operatorname{D\acute{e}r}_A(B',L) \ar[d]^{w_0} \\
  \operatorname{D\acute{e}r}_A(B,L') \ar[r]^{v^0} &
    \operatorname{D\acute{e}r}_A(B',L')
}
\qquad
\xymatrix{
  \operatorname{D\acute{e}r}_A(B,L) \ar[r]^{u^0} \ar[d]_{w_0} &
    \operatorname{D\acute{e}r}_{A'}(B,L) \ar[d]^{w_0} \\
  \operatorname{D\acute{e}r}_A(B,L') \ar[r]^{u^0} &
    \operatorname{D\acute{e}r}_{A'}(B,L')
}
\]
pour tout homomorphisme $w:L\to L'$ de $B$-bimodules.
\end{env}

\oldpage[0\textsubscript{IV-1}]{120}
\begin{theorem}[20.2.2]
\label{0.20.2.2-fr}
Soient $u:A\to B$, $v:B\to C$ deux homomorphismes d'anneaux, $L$ un
$C$-bimodule. On a une suite exacte canonique de groupes commutatifs
\begin{align}
\tag{20.2.2.1}
\label{0.20.2.2.1-fr}
  0\to\operatorname{D\acute{e}r}_B(C,L)\xrightarrow{u^0}
  \operatorname{D\acute{e}r}_A(C,L)&\xrightarrow{v^0}
  \operatorname{D\acute{e}r}_A(B,L)\xrightarrow{\partial} \\
  &\to\operatorname{Exan}_B(C,L)\xrightarrow{u^1}
  \operatorname{Exan}_A(C,L)\xrightarrow{v^1}
  \operatorname{Exan}_A(B,L) \notag
\end{align}
où $u^0$, $v^0$ sont les homomorphismes (20.2.1.3) et (20.2.1.2)
respectivement, $u^1$, $v^1$ les homomorphismes définis dans (18.3.4.1) et
(18.3.3.1) respectivement, et où $\partial$ est défini comme suit : pour
toute $A$-dérivation $D$ de $B$ dans $L$, $\partial(D)$ est la classe de la
$B$-extension de $C$ par $L$ définie sur l'anneau $D_C(L)$ par le
$A$-homomorphisme $\alpha:x\mapsto(v(x),D(x))$ (cf.
\sref{0.20.1.5-fr}). En outre, la suite exacte \sref{0.20.2.2.1-fr} est
fonctorielle en $L$ (pour les homomorphismes définis dans (20.2.1.1) et
(18.3.1.1) respectivement).
\end{theorem}

\begin{proof}
Comme $D$ est une $A$-dérivation (et \emph{a fortiori} une
$\mathbf{Z}$-dérivation) de $B$ dans $L$, le $A$-homomorphisme
$\alpha:x\mapsto(v(x),D(x))$ définit bien sur $D_C(L)$ une structure de
$B$-extension, donc $\partial(D)$ est bien définie \sref{0.20.1.5-fr}. Il
faut vérifier l'exactitude en $5$ endroits :

\medskip
\noindent 1)\enspace
L'exactitude en $\operatorname{D\acute{e}r}_B(C,L)$ est triviale (cf.
\sref{0.20.2.1-fr}).

\medskip
\noindent 2)\enspace
Par définition \sref{0.20.2.1-fr} le noyau de $v^0$ est l'ensemble des
$A$-dérivations de $C$ dans $L$ qui s'annulent dans $v(B)$, c'est-à-dire
celles des $A$-dérivations qui sont aussi des $B$-dérivations
\sref{0.20.1.1-fr} ; d'où l'exactitude en $\operatorname{D\acute{e}r}_A(C,L)$.

\medskip
\noindent 3)\enspace
Le noyau de $\partial$ est formé des dérivations
$D\in\operatorname{D\acute{e}r}_A(B,L)$ pour lesquelles la $B$-extension définie par
$\alpha:x\mapsto(v(x),D(x))$ est $B$-triviale ; cela signifie (18.2.3) qu'il
existe un $B$-homomorphisme $z\mapsto(z,w(z))$ de $C$ dans $D_C(L)$ (la
structure de $B$-anneau de $D_C(L)$ étant définie par $\alpha$) ; mais un tel
homomorphisme, étant \emph{a fortiori} un $A$-homomorphisme, est de la forme
$z\mapsto(z,D'(z))$ où $D'\in\operatorname{D\acute{e}r}_A(C,L)$
\sref{0.20.1.6-fr} ; et en écrivant que c'est un $B$-homomorphisme, il vient
\[
  D(x)z+v(x)D'(z)=D'(v(x)z)=v(x)D'(z)+D'(v(x))z
\]
pour $x\in B$, $z\in C$, ce qui donne $D'(v(x))=D(x)$ ; le noyau de
$\partial$ est donc l'image de $v^0$.

\medskip
\noindent 4)\enspace
Le noyau de $u^1$ est l'ensemble des classes de $B$-extensions de $C$ par
$L$ qui sont $A$-triviales (18.3.7), donc (à équivalence près) de la forme
$D_C(L)$, où la structure de $A$-anneau est définie par l'homomorphisme
$t\mapsto(u(t),0)$. Or, toute structure de $B$-extension sur $D_C(L)$ est
définie par un homomorphisme $\alpha:x\mapsto(v(x),D(x))$, où $D$ est une
$\mathbf{Z}$-dérivation de $B$ dans $L$ \sref{0.20.1.5-fr} ; dire que la
structure de $A$-anneau de cette $B$-extension est déduite de sa structure de
$B$-anneau au moyen de $u:A\to B$ signifie que $D(u(t))=0$ pour $t\in A$,
donc que $D$ est une $A$-dérivation ou encore que la classe de la $B$-extension
considérée est de la forme $\partial(D)$ ; d'où l'exactitude en
$\operatorname{Exan}_B(C,L)$.

\medskip
\noindent 5)\enspace
Le noyau de $v^1$ est l'ensemble des classes des $A$-extensions $E$ de $C$
par $L$ qui se trivialisent sur $v(B)$, c'est-à-dire pour lesquelles il y a un
$A$-homomorphisme $w:B\to E$ tel que $v$ se factorise en
$B\xrightarrow{w}E\to C$ ; or un tel $A$-homomorphisme définit sur $E$ une
structure de $B$-extension dont la classe a pour image par $u^1$ la classe de
la $A$-extension donnée ; la réciproque étant triviale, on a prouvé
l'exactitude en $\operatorname{Exan}_A(C,L)$.

Enfin, la fonctorialité en $L$ découle trivialement des définitions.
\end{proof}

\oldpage[0\textsubscript{IV-1}]{121}
\begin{corollary}[20.2.3]
\label{0.20.2.3-fr}
Soient $A$, $B$, $C$ trois anneaux commutatifs, $u:A\to B$, $v:B\to C$
deux homomorphismes d'anneaux, $L$ un $C$-module. On a une suite exacte
canonique de $A$-modules
\begin{align}
\tag{20.2.3.1}
\label{0.20.2.3.1-fr}
  0\to\operatorname{D\acute{e}r}_B(C,L)\xrightarrow{u^0}
  \operatorname{D\acute{e}r}_A(C,L)&\xrightarrow{v^0}
  \operatorname{D\acute{e}r}_A(B,L)\xrightarrow{\partial} \\
  &\to\operatorname{Exalcom}_B(C,L)\xrightarrow{u^1}
  \operatorname{Exalcom}_A(C,L)\xrightarrow{v^1}
  \operatorname{Exalcom}_A(B,L) \notag
\end{align}
fonctorielle en $L$.
\end{corollary}

\begin{proof}
Le raisonnement est le même que dans \sref{0.20.2.2-fr} une fois qu'on a
vérifié que pour toute dérivation $D\in\operatorname{D\acute{e}r}_A(B,L)$,
$\partial(D)$ est bien la classe d'une $B$-extension de $C$ par $L$ qui est
une $B$-algèbre commutative ; mais cela résulte aussitôt de la commutativité
de $C$ et du fait que $L$ est un $C$-module.
\end{proof}

\begin{corollary}[20.2.4]
\label{0.20.2.4-fr}
Sous les hypothèses de \sref{0.20.2.2-fr} (resp.
\sref{0.20.2.3-fr}), on a une suite exacte canonique, fonctorielle en $L$,
\begin{equation}
\tag{20.2.4.1}
\label{0.20.2.4.1-fr}
  0\to\operatorname{D\acute{e}r}_B(C,L)\xrightarrow{u^0}
  \operatorname{D\acute{e}r}_A(C,L)\xrightarrow{v^0}
  \operatorname{D\acute{e}r}_A(B,L)\xrightarrow{\partial}
  \operatorname{Exan}_{B/A}(C,L)\to0
\end{equation}
(resp.
\begin{equation}
\tag{20.2.4.2}
\label{0.20.2.4.2-fr}
  0\to\operatorname{D\acute{e}r}_B(C,L)\xrightarrow{u^0}
  \operatorname{D\acute{e}r}_A(C,L)\xrightarrow{v^0}
  \operatorname{D\acute{e}r}_A(B,L)\xrightarrow{\partial}
  \operatorname{Exalcom}_{B/A}(C,L)\to0\text{).}
\end{equation}
Cela résulte de la définition de $\operatorname{Exan}_{B/A}(C,L)$ (resp.
$\operatorname{Exalcom}_{B/A}(C,L)$) ((18.3.7) et (18.4.2)).
\end{corollary}

\begin{remark}[20.2.5]
\label{0.20.2.5-fr}
Supposons que l'on ait un diagramme commutatif d'homomorphismes d'anneaux
\[
\xymatrix{
  A \ar[r] & B \ar[r] & C \\
  A' \ar[u] \ar[r] & B' \ar[u] \ar[r] & C' \ar[u]
}
\]
Alors on a un diagramme commutatif
\[
\xymatrix@C=0.5em{
  0 \ar[r] & \operatorname{D\acute{e}r}_B(C,L) \ar[r] \ar[d] &
  \operatorname{D\acute{e}r}_A(C,L) \ar[r] \ar[d] &
  \operatorname{D\acute{e}r}_A(B,L) \ar[r]^-\partial \ar[d] &
  \operatorname{Exan}_B(C,L) \ar[r] \ar[d] &
  \operatorname{Exan}_A(C,L) \ar[r] \ar[d] &
  \operatorname{Exan}_A(B,L) \ar[d] \\
  0 \ar[r] & \operatorname{D\acute{e}r}_{B'}(C',L) \ar[r] &
  \operatorname{D\acute{e}r}_{A'}(C',L) \ar[r] &
  \operatorname{D\acute{e}r}_{A'}(B',L) \ar[r]^-\partial &
  \operatorname{Exan}_{B'}(C',L) \ar[r] &
  \operatorname{Exan}_{A'}(C',L) \ar[r] &
  \operatorname{Exan}_{A'}(B',L)
}
\]
et de même pour les suites exactes \sref{0.20.2.3.1-fr},
\sref{0.20.2.4.1-fr} et \sref{0.20.2.4.2-fr}.
\end{remark}

\subsection{Dérivations continues dans les anneaux topologiques}
\label{subsection:0.20.3-fr}

\begin{env}[20.3.1]
\label{0.20.3.1-fr}
Étant donnés deux anneaux topologiques $A$, $B$ (linéairement topologisés
comme toujours), nous désignerons par $\operatorname{Hom.\ cont}(A,B)$
l'ensemble des homomorphismes continus de $A$ dans $B$. Étant donné un anneau
topologique $A$, la catégorie des $A$-anneaux topologiques se définit comme
celle des $A$-anneaux (18.1.4) en remplaçant partout « anneau » par « anneau
topologique » et « homomorphisme » par « homomorphisme continu » ; si $B$

\oldpage[0\textsubscript{IV-1}]{122}
et $C$ sont deux $A$-anneaux topologiques, on notera
$\operatorname{Hom.\ cont}_A(B,C)$ l'ensemble des $A$-homomorphismes continus
de $B$ dans $C$.

Soient $A$ un anneau topologique, $B$ un $A$-anneau topologique, $L$ un
$B$-bimodule topologique ; on note $\operatorname{D\acute{e}r.\ cont}_A(B,L)$
l'ensemble des $A$-dérivations \emph{continues} de $B$ dans $L$ ; il est clair
que c'est un \emph{sous-groupe} de $\operatorname{D\acute{e}r}_A(B,L)$ (et un
sous-$A$-module lorsque $A$ et $B$ sont commutatifs et $L$ un $B$-module).
\end{env}

\begin{env}[20.3.2]
\label{0.20.3.2-fr}
Il est immédiat que la prop. \sref{0.20.1.1-fr} subsiste quand on y remplace
« anneau » par « anneau topologique » et « homomorphisme » par
« homomorphisme continu ». De même, \sref{0.20.1.3-fr} et
\sref{0.20.1.4-fr} restent valables en remplaçant $\operatorname{D\acute{e}r}$ par
$\operatorname{D\acute{e}r.\ cont}$, « anneau » (resp. « bimodule ») par « anneau
topologique » (resp. « bimodule topologique »), « homomorphisme » par
« homomorphisme continu » ; il faut naturellement supposer dans
\sref{0.20.1.4-fr} que $p$ est continu, et remplacer « $A$-équivalences » par
« $A$-équivalences continues ». On a les mêmes résultats pour
\sref{0.20.1.5-fr} et \sref{0.20.1.6-fr}, à condition de prendre comme
topologie sur $D_C(L)$ (resp. $D_B(L)$) la topologie \emph{produit} (sur
$C\times L$, resp. $B\times L$).
\end{env}

\begin{proposition}[20.3.3]
\label{0.20.3.3-fr}
Soient $A$ un anneau topologique, $B$ un $A$-anneau topologique, $L$ un
$B$-bimodule topologique discret annulé par un idéal bilatère ouvert de $B$.
Si dans $B$ le carré de tout idéal bilatère ouvert est ouvert, on a
$\operatorname{D\acute{e}r.\ cont}_A(B,L)=\operatorname{D\acute{e}r}_A(B,L)$.
\end{proposition}

\begin{proof}
En effet, si $\mathfrak{R}$ est un idéal bilatère ouvert de $B$ annulant $L$,
et $D$ une $A$-dérivation de $B$ dans $L$, on a
$D(\mathfrak{R}^2)\subset\mathfrak{R}.D(\mathfrak{R})+
D(\mathfrak{R}).\mathfrak{R}=0$ \sref{0.20.1.1.1-fr}, donc $D$ est continue.
\end{proof}

\begin{env}[20.3.4]
\label{0.20.3.4-fr}
Tous les résultats de \sref{0.20.2.1-fr} restent valables lorsqu'on y remplace
« anneau » par « anneau topologique », « bimodule » par « bimodule
topologique », « homomorphisme » par « homomorphisme continu » et
$\operatorname{D\acute{e}r}$ par $\operatorname{D\acute{e}r.\ cont}$.
\end{env}

\begin{proposition}[20.3.5]
\label{0.20.3.5-fr}
Soient $A$ un anneau topologique, $B$ un $A$-anneau topologique, $L$ un
$B$-bimodule topologique discret annulé par un idéal bilatère ouvert de $B$.
On a alors un isomorphisme canonique
\begin{equation}
\tag{20.3.5.1}
\label{0.20.3.5.1-fr}
  \varinjlim\operatorname{D\acute{e}r}_{A/\mathfrak{J}}
  (B/\mathfrak{R},L)\xrightarrow{\sim}
  \operatorname{D\acute{e}r.\ cont}_A(B,L)
\end{equation}
où dans le premier membre la limite inductive est prise suivant l'ensemble
ordonné filtrant des couples $(\mathfrak{J},\mathfrak{R})$ d'idéaux bilatères
tels que $\mathfrak{R}.L=L.\mathfrak{R}=0$,
$\mathfrak{J}.B\subset\mathfrak{R}$,
$B.\mathfrak{J}\subset\mathfrak{R}$.
\end{proposition}

\begin{proof}
Comme $A/\mathfrak{J}$ et $B/\mathfrak{R}$ sont discrets, on a des
homomorphismes canoniques
$w_{\mathfrak{J},\mathfrak{R}}:
\operatorname{D\acute{e}r}_{A/\mathfrak{J}}(B/\mathfrak{R},L)\to
\operatorname{D\acute{e}r.\ cont}_A(B,L)$ formant un système inductif
\sref{0.20.3.4-fr}, d'où l'homomorphisme \sref{0.20.3.5.1-fr} par passage à
la limite inductive. Comme l'homomorphisme
$B/\mathfrak{R}'\to B/\mathfrak{R}$ est surjectif pour
$\mathfrak{R}\supset\mathfrak{R}'$, il résulte aussitôt de la définition que
l'homomorphisme
$\operatorname{D\acute{e}r}_{A/\mathfrak{J}}(B/\mathfrak{R},L)\to
\operatorname{D\acute{e}r}_{A/\mathfrak{J}}(B/\mathfrak{R}',L)$ (avec
$\mathfrak{R}.L=L.\mathfrak{R}=0$,
$\mathfrak{J}.B\subset\mathfrak{R}'$,
$B.\mathfrak{J}\subset\mathfrak{R}'$) est injectif, et il en est évidemment
de même de l'homomorphisme
$\operatorname{D\acute{e}r}_{A/\mathfrak{J}}(B/\mathfrak{R},L)\to
\operatorname{D\acute{e}r}_{A/\mathfrak{J}'}(B/\mathfrak{R},L)$ pour
$\mathfrak{J}'\subset\mathfrak{J}$ \sref{0.20.2.1-fr} ; on en conclut que
l'homomorphisme \sref{0.20.3.5.1-fr} est injectif. D'autre part, si $D$ est
une $A$-dérivation continue de $B$ dans $L$, son noyau contient un idéal
bilatère ouvert $\mathfrak{R}_0$ de $B$, et si $\mathfrak{J}_0$ est un idéal
bilatère ouvert de $A$ tel que $\mathfrak{J}_0B\subset\mathfrak{R}_0$ et
$B\mathfrak{J}_0\subset\mathfrak{R}_0$, il est clair que $D$ est l'image
canonique d'une $(A/\mathfrak{J}_0)$-dérivation de
$B/\mathfrak{R}_0$ dans $L$, donc \sref{0.20.3.5.1-fr} est surjectif.
\end{proof}

\oldpage[0\textsubscript{IV-1}]{123}
\begin{proposition}[20.3.6]
\label{0.20.3.6-fr}
Soient $u:A\to B$, $v:B\to C$ deux homomorphismes continus d'anneaux
topologiques, $L$ un $C$-bimodule discret annulé par un idéal bilatère ouvert
de $C$. On a une suite exacte canonique
\begin{align}
\tag{20.3.6.1}
\label{0.20.3.6.1-fr}
  0\to\operatorname{D\acute{e}r.\ cont}_B(C,L)\xrightarrow{u^0}
  \operatorname{D\acute{e}r.\ cont}_A(C,L)&\xrightarrow{v^0}
  \operatorname{D\acute{e}r.\ cont}_A(B,L)\xrightarrow{\partial} \\
  &\to\operatorname{Exantop}_B(C,L)\xrightarrow{u^1}
  \operatorname{Exantop}_A(C,L)\xrightarrow{v^1}
  \operatorname{Exantop}_A(B,L) \notag
\end{align}
où $\partial$ est défini par passage à la limite inductive à partir de
l'homomorphisme $\partial$ de \sref{0.20.2.2.1-fr} ; cette suite exacte est
fonctorielle en $L$ (dans la catégorie des $C$-bimodules discrets annulés par
des idéaux bilatères ouverts).
\end{proposition}

\begin{proof}
Cela résulte de l'exactitude du foncteur $\varinjlim$, à partir de
\sref{0.20.2.2-fr}.
\end{proof}

\begin{corollary}[20.3.7]
\label{0.20.3.7-fr}
Soient $A$, $B$, $C$ trois anneaux topologiques commutatifs, $u:A\to B$,
$v:B\to C$ deux homomorphismes continus, $L$ un $C$-module discret annulé
par un idéal ouvert de $C$. On a une suite exacte canonique de $A$-modules,
fonctorielle en $L$,
\begin{align}
\tag{20.3.7.1}
\label{0.20.3.7.1-fr}
  0\to\operatorname{D\acute{e}r.\ cont}_B(C,L)\xrightarrow{u^0}
  \operatorname{D\acute{e}r.\ cont}_A(C,L)&\xrightarrow{v^0}
  \operatorname{D\acute{e}r.\ cont}_A(B,L)\xrightarrow{\partial} \\
  &\to\operatorname{Exalcotop}_B(C,L)\xrightarrow{u^1}
  \operatorname{Exalcotop}_A(C,L)\xrightarrow{v^1}
  \operatorname{Exalcotop}_A(B,L). \notag
\end{align}
\end{corollary}

\begin{corollary}[20.3.8]
\label{0.20.3.8-fr}
Sous les hypothèses de \sref{0.20.3.5-fr} (resp.
\sref{0.20.3.6-fr}) on a une suite exacte canonique, fonctorielle en $L$
\begin{equation}
\tag{20.3.8.1}
\label{0.20.3.8.1-fr}
  0\to\operatorname{D\acute{e}r.\ cont}_B(C,L)\xrightarrow{u^0}
  \operatorname{D\acute{e}r.\ cont}_A(C,L)\xrightarrow{v^0}
  \operatorname{D\acute{e}r.\ cont}_A(B,L)\to
  \operatorname{Exantop}_{B/A}(C,L)\to0
\end{equation}
(resp.
\begin{equation}
\tag{20.3.8.2}
\label{0.20.3.8.2-fr}
  0\to\operatorname{D\acute{e}r.\ cont}_B(C,L)\xrightarrow{u^0}
  \operatorname{D\acute{e}r.\ cont}_A(C,L)\xrightarrow{v^0}
  \operatorname{D\acute{e}r.\ cont}_A(B,L)\to
  \operatorname{Exalcotop}_{B/A}(C,L)\to0\text{).}
\end{equation}
\end{corollary}

Nous laissons au lecteur le soin d'écrire les diagrammes analogues à ceux de
\sref{0.20.2.5-fr}.

\subsection{Parties principales et différentielles}
\label{subsection:0.20.4-fr}

Dans toute la suite de ce paragraphe et dans les trois suivants, tous les
anneaux sont supposés \emph{commutatifs}.

\begin{env}[20.4.1]
\label{0.20.4.1-fr}
Soient $A$ un anneau topologique, $B$ une $A$-algèbre topologique ; la
$A$-algèbre $B\otimes_A B$ sera munie de la topologie \emph{produit
tensoriel}, qui en fait une $A$-algèbre topologique ; nous désignerons par $p$
(ou $p_{B/A}$) le $A$-homomorphisme canonique surjectif
\begin{equation}
\tag{20.4.1.1}
\label{0.20.4.1.1-fr}
  p:B\otimes_A B\to B
\end{equation}
tel que $p(b\otimes b')=bb'$ ; il est immédiat que $p$ est continu. Le
\emph{noyau} de $p$ sera noté $\mathfrak{J}_{B/A}$ (ou seulement
$\mathfrak{J}$ s'il n'y a pas de confusion). On désignera par
$j_1:B\to B\otimes_A B$ et $j_2:B\to B\otimes_A B$ les deux
$A$-homomorphismes canoniques, tels que
\[
  j_1(b)=b\otimes1, \qquad j_2(b)=1\otimes b
\]
qui sont continus.
\end{env}

\oldpage[0\textsubscript{IV-1}]{124}
\begin{definition}[20.4.2]
\label{0.20.4.2-fr}
On appelle \emph{$B$-algèbre augmentée des parties principales d'ordre $1$ de
$B$ par rapport à $A$} et l'on note $\mathrm{P}^1_{B/A}$ la $A$-algèbre
topologique quotient
\begin{equation}
\tag{20.4.2.1}
\label{0.20.4.2.1-fr}
  \mathrm{P}^1_{B/A}=(B\otimes_A B)/\mathfrak{J}^2
\end{equation}
munie de la structure de $B$-algèbre topologique définie par l'homomorphisme
$\bar{j}_1:B\to\mathrm{P}^1_{B/A}$ (déduit de $j_1$ par composition avec
l'homomorphisme canonique $B\otimes_A B\to\mathrm{P}^1_{B/A}$), et de
l'augmentation de $B$-algèbre $\varepsilon:\mathrm{P}^1_{B/A}\to B$ (aussi
notée $\varepsilon_{B/A}$) déduite de $p$ par passage au quotient.

Comme $p(b\otimes1)=b$ par définition, il est clair que $\varepsilon$ est
bien une augmentation de $B$-algèbre.
\end{definition}

\begin{definition}[20.4.3]
\label{0.20.4.3-fr}
Le noyau de l'augmentation $\varepsilon:\mathrm{P}^1_{B/A}\to B$,
\begin{equation}
\tag{20.4.3.1}
\label{0.20.4.3.1-fr}
  \Omega^1_{B/A}=\mathfrak{J}_{B/A}/(\mathfrak{J}_{B/A})^2
\end{equation}
muni de la topologie induite par celle de $\mathrm{P}^1_{B/A}$, qui en fait
un $B$-module topologique, est appelé le $B$-module des
\emph{$1$-différentielles} (ou simplement des \emph{différentielles}) de $B$
par rapport à $A$.
\end{definition}

On notera que la topologie de $\Omega^1_{B/A}$ est aussi la topologie quotient
de la topologie induite sur $\mathfrak{J}_{B/A}$ par celle de
$B\otimes_A B$ (Bourbaki, \emph{Top. gén.}, chap.~III, 3\textsuperscript{e}
éd., §~2, n\textsuperscript{o}~7, prop.~20). Si $B$ est discret il en est de
même de $\Omega^1_{B/A}$. On note $\widehat{\Omega}^1_{B/A}$ le
\emph{séparé complété} du $B$-module topologique $\Omega^1_{B/A}$.

Tout anneau topologique $B$ pouvant être considéré comme une
$\mathbf{Z}$-algèbre topologique (où $\mathbf{Z}$ est muni de la topologie
discrète), on peut définir le $B$-module topologique $\Omega^1_{B/\mathbf{Z}}$,
qu'on appelle parfois aussi le $B$-module des \emph{différentielles absolues}
sur $B$ et que l'on note $\Omega^1_B$. Si $B$ est une algèbre topologique sur
un corps premier $P$ (discret), on a
$B\otimes_P B=B\otimes_{\mathbf{Z}}B$, donc
$\Omega^1_{B/P}=\Omega^1_B$.

\begin{lemma}[20.4.4]
\label{0.20.4.4-fr}
Soient $A$ un anneau, $B$ une $A$-algèbre. L'idéal $\mathfrak{J}_{B/A}$ de
$B\otimes_A B$ est engendré par les éléments $1\otimes s-s\otimes1$, où $s$
parcourt un ensemble de générateurs de la $A$-algèbre $B$.
\end{lemma}

\begin{proof}
Il est clair que pour tout $x\in B$, on a
$x\otimes1-1\otimes x\in\mathfrak{J}$ ; d'autre part, quels que soient $x$,
$y$ dans $B$, on a
$x\otimes y=xy\otimes1+(x\otimes1)(1\otimes y-y\otimes1)$. Si
$\sum_i(x_i\otimes y_i)\in\mathfrak{J}$, on a par définition
$\sum_i x_i y_i=0$, donc
\begin{equation}
\tag{20.4.4.1}
\label{0.20.4.4.1-fr}
  \sum_i(x_i\otimes y_i)=
  \sum_i(x_i\otimes1)(1\otimes y_i-y_i\otimes1)
\end{equation}
ce qui prouve que $\mathfrak{J}$ est l'idéal engendré par les éléments
$1\otimes x-x\otimes1$. En outre, si $x=st$, on a
\begin{equation}
\tag{20.4.4.2}
\label{0.20.4.4.2-fr}
  x\otimes1-1\otimes x=(s\otimes1)(t\otimes1-1\otimes t)
  +(s\otimes1-1\otimes s)(1\otimes t)
\end{equation}
ce qui termine aussitôt la démonstration par récurrence.
\end{proof}

\begin{proposition}[20.4.5]
\label{0.20.4.5-fr}
Soient $A$ un anneau topologique, $B$ une $A$-algèbre topologique. La
topologie de $\Omega^1_{B/A}$ est moins fine que la topologie déduite de celle
de $B$ (19.0.2) ; si dans $B$ le carré de tout idéal ouvert est ouvert, ces
deux topologies sont identiques.
\end{proposition}

\begin{proof}
La première assertion est triviale, la topologie produit tensoriel sur
$B\otimes_A B$ étant moins fine que la topologie déduite de celle de $B$ ;
\emph{a fortiori} la topologie induite sur
$\mathfrak{J}=\mathfrak{J}_{B/A}$

\oldpage[0\textsubscript{IV-1}]{125}
par celle de $B\otimes_A B$ est moins fine que la topologie sur $\mathfrak{J}$
déduite de celle de $B$. Pour démontrer la seconde assertion, écrivons
$M\otimes N$, par abus de notation, le sous-module
$\operatorname{Im}(M\otimes_A N)$ pour deux sous-$A$-modules $M$, $N$ de $B$.
Utilisant la relation
\[
\begin{aligned}
  (xy)\otimes z-x\otimes(yz)
  &=(x\otimes1)(1\otimes z)(y\otimes1-1\otimes y)\\
  &=xz.(y\otimes1-1\otimes y)
    -x.(z\otimes1-1\otimes z)(y\otimes1-1\otimes y)
\end{aligned}
\]
dans le $B$-module $B\otimes_A B$ (défini par $j_1$), on voit aussitôt,
compte tenu de \sref{0.20.4.4-fr}, que, si $\mathfrak{R}$ est un idéal de
$B$, on a
\[
  ((\mathfrak{R}^2\otimes B)+(B\otimes\mathfrak{R}^2))\cap\mathfrak{J}
  \subset(\mathfrak{R}\otimes\mathfrak{R})\cap\mathfrak{J}
  +\mathfrak{R}\mathfrak{J}+\mathfrak{J}^2
\]
et d'autre part, si $x_i$, $y_i$ sont des éléments de $\mathfrak{R}$ tels
que $\sum_i(x_i\otimes y_i)\in\mathfrak{J}$, il résulte de
\sref{0.20.4.4.1-fr} que l'on a
$\sum_i(x_i\otimes y_i)\in\mathfrak{R}\mathfrak{J}$, si bien que finalement
\begin{equation}
\tag{20.4.5.1}
\label{0.20.4.5.1-fr}
  (\mathfrak{R}^2\otimes B+B\otimes\mathfrak{R}^2)\cap\mathfrak{J}
  \subset\mathfrak{R}\mathfrak{J}+\mathfrak{J}^2.
\end{equation}
Or, on a un système fondamental de voisinages de $0$ dans $\mathfrak{J}$
(pour la topologie induite par celle de $B\otimes_A B$) en prenant pour
voisinages de $0$ les ensembles
$(\mathfrak{R}\otimes B+B\otimes\mathfrak{R})\cap\mathfrak{J}$, où
$\mathfrak{R}$ parcourt l'ensemble des idéaux ouverts. Comme la topologie de
$\Omega^1_{B/A}$ déduite de celle de $B$ est aussi le quotient par
$\mathfrak{J}^2$ de la topologie de $\mathfrak{J}$ déduite de celle de $B$,
l'hypothèse sur les idéaux ouverts de $B$ et la relation
\sref{0.20.4.5.1-fr} achèvent de prouver la seconde assertion.
\end{proof}

\begin{definition}[20.4.6]
\label{0.20.4.6-fr}
Soient $p_1$ et $p_2$ les applications composées
$B\xrightarrow{j_1}B\otimes_A B\to\mathrm{P}^1_{B/A}$,
$B\xrightarrow{j_2}B\otimes_A B\to\mathrm{P}^1_{B/A}$, qui sont des
$A$-homomorphismes continus tels que
$\varepsilon\circ p_1=\varepsilon\circ p_2=\mathrm{I}_B$. Le
$A$-homomorphisme continu de $A$-modules
\begin{equation}
\tag{20.4.6.1}
\label{0.20.4.6.1-fr}
  d_{B/A}=p_1-p_2:B\to\Omega^1_{B/A}
\end{equation}
est appelé la \emph{différentielle extérieure de $B$ relative à $A$} ; pour
tout $x\in B$, $d_{B/A}(x)$ (aussi noté $d(x)$ ou $dx$) est appelé la
\emph{différentielle de $x$} (relativement à $A$).
\end{definition}

Lorsque $A=\mathbf{Z}$, on écrit $d_B$ au lieu de $d_{B/\mathbf{Z}}$ ; si $B$
est une algèbre sur un corps premier $P$, on a $d_B=d_{B/P}$.

\begin{proposition}[20.4.7]
\label{0.20.4.7-fr}
Le $B$-module $\Omega^1_{B/A}$ est engendré par les éléments $d_{B/A}(x)$,
où $x$ parcourt un système de générateurs de la $A$-algèbre $B$.
\end{proposition}

\begin{proof}
Comme $d_{B/A}(x)$ est l'image canonique de $x\otimes1-1\otimes x$ dans
$\Omega^1_{B/A}$, la proposition est une conséquence immédiate de
\sref{0.20.4.4-fr}.
\end{proof}

\begin{theorem}[20.4.8]
\label{0.20.4.8-fr}
Soient $A$ un anneau topologique, $B$ une $A$-algèbre topologique.
\begin{enumerate}
  \item[\rm(i)] Il existe un isomorphisme unique de $B$-algèbres augmentées
  topologiques
\begin{equation}
\tag{20.4.8.1}
\label{0.20.4.8.1-fr}
  \varphi:\mathrm{P}^1_{B/A}\xrightarrow{\sim}D_B(\Omega^1_{B/A})
\end{equation}
qui se réduit à l'identité dans $\Omega^1_{B/A}$.
  \item[\rm(ii)] L'homomorphisme $d_{B/A}$ est une $A$-dérivation de $B$
  dans $\Omega^1_{B/A}$, ayant la propriété universelle suivante : pour tout
  $B$-module topologique $L$, l'application $u\mapsto u\circ d_{B/A}$ est
  un isomorphisme de $A$-modules
\begin{equation}
\tag{20.4.8.2}
\label{0.20.4.8.2-fr}
  \operatorname{Hom.\ cont}_B(\Omega^1_{B/A},L)
  \xrightarrow{\sim}\operatorname{D\acute{e}r.\ cont}_A(B,L).
\end{equation}
\end{enumerate}
\end{theorem}

\begin{proof}
\oldpage[0\textsubscript{IV-1}]{126}
(i) Il est immédiat que $\varphi$ (avec les notations de
\sref{0.20.4.6-fr}) est nécessairement l'application
$z\mapsto(\varepsilon(z),z-p_1(\varepsilon(z)))$, et l'isomorphisme
réciproque l'application $(b,x)\mapsto p_1(b)+x$ ; ces deux applications
étant continues, cela prouve la première assertion. On notera que cela
entraîne que la topologie de $B$ s'identifie (par $p_1$) au \emph{quotient}
de la topologie de $B\otimes_A B$ par $\mathfrak{J}_{B/A}$.

(ii) Le fait que $d_{B/A}$ soit une $A$-dérivation de $B$ résulte de la
définition \sref{0.20.4.6-fr} et de \sref{0.20.1.1-fr}.

Pour prouver la propriété universelle, rappelons que
$\operatorname{D\acute{e}r.\ cont}_A(B,L)$ s'identifie canoniquement à
l'ensemble $G$ des homomorphismes continus de $A$-algèbres
$u:B\to D_B(L)$ tels que le composé
$B\xrightarrow{u}D_B(L)\xrightarrow{q}B$ soit l'identité (où $q(b,0)=b$ est
la projection ; cf. \sref{0.20.1.6-fr} et \sref{0.20.3.2-fr}). D'autre part,
grâce à l'isomorphisme $\varphi$,
$\operatorname{Hom.\ cont}_B(\Omega^1_{B/A},L)$ s'identifie canoniquement à
l'ensemble des homomorphismes continus de $B$-algèbres
$v:\mathrm{P}^1_{B/A}\to D_B(L)$ tels que le composé
$\mathrm{P}^1_{B/A}\xrightarrow{v}D_B(L)\xrightarrow{q}B$ soit
l'augmentation $\varepsilon$. Comme on a $p_2=p_1-d_{B/A}$ par définition,
tout revient à prouver que tout $u\in G$ se factorise en
\[
\xymatrix{
  B \ar[r]^u \ar[d]_{p_2} & D_B(L) \\
  \mathrm{P}^1_{B/A} \ar[ur]_v &
}
\]
où $v$ est un $B$-homomorphisme continu. Or, on a déjà un homomorphisme
continu de $A$-algèbres $j:b\mapsto(b,0)$ de $B$ dans $D_B(L)$, qui
appartient à $G$ ; par définition du produit tensoriel topologique
d'algèbres topologiques ($\mathbf{0_I}$, 7.7.6), il existe donc un
$A$-homomorphisme continu d'algèbres
$w:B\otimes_A B\to D_B(L)$ rendant commutatif le diagramme
\[
\xymatrix{
  B\otimes_A B \ar[dr]^w & B \ar[l]_{j_1} \ar[d]^j \\
  B \ar[u]^{j_2} \ar[r]_u & D_B(L)
}
\]
On a donc par définition
$w(b\otimes1-1\otimes b)=j(b)-u(b)\in L$, et en vertu de
\sref{0.20.4.4-fr}, cela entraîne $w(\mathfrak{J})\subset L$ donc
$w(\mathfrak{J}^2)=0$ ; par suite $w$ se factorise en
\[
  B\otimes_A B\to\mathrm{P}^1_{B/A}\xrightarrow{v}D_B(L)
\]
où $v$ est un homomorphisme continu de $A$-algèbres ; d'ailleurs, comme
$v\circ p_1=j$ est un homomorphisme de $B$-algèbres, il en est de même de
$v$ par définition de la structure de $B$-algèbre de
$\mathrm{P}^1_{B/A}$ ; comme on a par définition $u=v\circ p_2$, cela
termine la démonstration.
\end{proof}

\begin{theorem}[20.4.9]
\label{0.20.4.9-fr}
Supposons que $B$ soit une $A$-algèbre topologique formellement lisse. Alors
le $B$-module topologique $\Omega^1_{B/A}$ est \emph{formellement projectif}.
\end{theorem}

\begin{proof}
En effet, l'hypothèse entraîne que $B\otimes_A B$ (muni de la structure de
$B$-algèbre topologique définie par $j_1$) est une $B$-algèbre topologique
formellement lisse \sref{0.19.3.5-fr}, (iii), et par suite aussi une
$A$-algèbre topologique formellement lisse \sref{0.19.3.5-fr}, (ii) ; comme
$B$ est topologiquement isomorphe à la $A$-algèbre quotient
$(B\otimes_A B)/\mathfrak{J}$ et est une $A$-algèbre formellement lisse, la
conclusion résulte de \sref{0.19.5.3-fr}.
\end{proof}

\oldpage[0\textsubscript{IV-1}]{127}
\begin{corollary}[20.4.10]
\label{0.20.4.10-fr}
Supposons en outre que dans $B$ le carré de tout idéal ouvert soit ouvert ;
alors, pour tout idéal ouvert $\mathfrak{R}$ de $B$,
$\Omega^1_{B/A}\otimes_B(B/\mathfrak{R})
=\Omega^1_{B/A}/\mathfrak{R}.\Omega^1_{B/A}$ est un
$(B/\mathfrak{R})$-module projectif.
\end{corollary}

\begin{proof}
En effet, la topologie de $\Omega^1_{B/A}$ est alors déduite de celle de
$B$ \sref{0.20.4.5-fr}, et il suffit d'appliquer \sref{0.19.2.4-fr}.
\end{proof}

\begin{corollary}[20.4.11]
\label{0.20.4.11-fr}
Soient $A$, $B$ deux anneaux locaux noethériens, $\rho:A\to B$ un
homomorphisme local, faisant de $B$ une $A$-algèbre topologique formellement
lisse (pour les topologies préadiques). Alors, pour tout idéal de définition
$\mathfrak{b}$ de $B$,
$\Omega^1_{B/A}\otimes_B(B/\mathfrak{b})$ est un
$(B/\mathfrak{b})$-module libre.
\end{corollary}

\begin{proof}
En effet, il résulte de \sref{0.20.4.10-fr} que ce module est projectif, et
comme $B/\mathfrak{b}$ est un anneau artinien, tout
$(B/\mathfrak{b})$-module projectif est libre
($\mathbf{0_{III}}$, 10.1.3).
\end{proof}

\begin{proposition}[20.4.12]
\label{0.20.4.12-fr}
Soient $A$ un anneau topologique, $B$ une $A$-algèbre topologique. Si
l'homomorphisme structural $A\to B$ est surjectif, on a
$\Omega^1_{B/A}=0$.
\end{proposition}

\begin{proof}
En effet, on a $\operatorname{D\acute{e}r}_A(B,L)=0$ pour tout $B$-module
$L$ en vertu de \sref{0.20.1.1-fr}, et la proposition découle donc aussitôt
de \sref{0.20.4.8-fr}.
\end{proof}

\begin{examples}[20.4.13]
\label{0.20.4.13-fr}
\begin{enumerate}
  \item[\rm(i)] Soient $A$ un anneau,
  $B=A[X_\alpha]_{\alpha\in I}$ une algèbre de polynômes sur $A$. Alors
  $\Omega^1_{B/A}$ est un $B$-module libre, dont les $dX_\alpha$ forment une
  base. En effet, les $dX_\alpha$ engendrent ce $B$-module
  \sref{0.20.4.7-fr}. D'autre part, si $L$ est un $B$-module libre, ayant
  une base $(z_\alpha)_{\alpha\in I}$ dont $I$ est l'ensemble d'indices, il
  existe un $A$-homomorphisme $u$ de $B$ dans $D_B(L)$ tel que
  $u(X_\alpha)=(X_\alpha,z_\alpha)$, pour tout $\alpha\in I$, donc
  \sref{0.20.1.5-fr} une $A$-dérivation $D$ de $B$ dans $L$ telle que
  $D(X_\alpha)=z_\alpha$ pour tout $\alpha$ ; en vertu de
  \sref{0.20.4.8.1-fr}, cela prouve que les $dX_\alpha$ sont linéairement
  indépendants.
  \item[\rm(ii)] Soient $A$ un anneau, $L$ un $A$-module, $B$ la
  $A$-algèbre $D_A(L)$ ; alors l'homomorphisme canonique
  $x\mapsto d_{B/A}x$ de $L$ dans $\Omega^1_{B/A}$ est \emph{bijectif}, car
  il est immédiat que les $B$-dérivations de $B=D_A(L)$ dans un $B$-module
  $M$ sont les applications de la forme $(b,x)\mapsto u(x)$, où
  $u\in\operatorname{Hom}_A(L,M)$ ; on conclut donc par
  \sref{0.20.4.8-fr}, (ii).
  \item[\rm(iii)] Supposons que $B$ soit le produit de deux $A$-algèbres
  topologiques $B_1$, $B_2$ (identifiées à des idéaux de $B$). Alors les
  images de $B_1\otimes_A B_2$ et de $B_2\otimes_A B_1$ par l'homomorphisme
  $p$ \sref{0.20.4.1.1-fr} sont nulles, d'où il résulte aussitôt que
  $\mathrm{P}^1_{B/A}$ s'identifie au \emph{produit}
  $\mathrm{P}^1_{B_1/A}\times\mathrm{P}^1_{B_2/A}$, et que le $B$-module
  $\Omega^1_{B/A}$ est \emph{somme directe} (topologique) des $B$-modules
  $\Omega^1_{B_1/A}$ et $\Omega^1_{B_2/A}$ (annulés respectivement par
  $B_2$ et $B_1$).
  \item[\rm(iv)] Soient $A$, $B$ deux anneaux intègres tels que $A\subset
  B$, que $A$ soit intégralement clos, $B$ entier sur $A$, et que le corps
  des fractions de $B$ soit une extension séparable de celui de $A$. Alors
  le $B$-module $\Omega^1_{B/A}$ est un \emph{module de torsion}. En effet,
  pour tout $x\in B$, le polynôme minimal de $x$ par rapport au corps des
  fractions $K$ de $A$ est un polynôme $f(T)$ appartenant à $A[T]$ ; comme
  $x$ est séparable sur $K$, on a $f'(x)\neq0$ et d'autre part on déduit de
  la relation $f(x)=0$ que $f'(x)d_{B/A}x=0$, d'où notre assertion en vertu
  de \sref{0.20.4.7-fr}.
\end{enumerate}
\end{examples}

\begin{remarks}[20.4.14]
\label{0.20.4.14-fr}
\begin{enumerate}
  \item[\rm(i)] On notera que le $B$-module $\Omega^1_{B/A}$, privé de sa
  topologie, est indépendant des topologies de $A$ et de $B$.
  \item[\rm(ii)] Nous introduirons plus tard l'« algèbre des parties
  principales d'ordre $n$ » de $B$ par rapport à $A$,
  $\mathrm{P}^n_{B/A}=(B\otimes_A B)/(\mathfrak{J}_{B/A})^{n+1}$, qui est à
  la base du « calcul différentiel d'ordre $n$ ».
\end{enumerate}
\end{remarks}

\subsection{Propriétés fonctorielles fondamentales de
\texorpdfstring{$\Omega^1_{B/A}$}{Omega1(B/A)}}
\label{subsection:0.20.5-fr}

\oldpage[0\textsubscript{IV-1}]{128}
\begin{env}[20.5.1]
\label{0.20.5.1-fr}
Dans tout ce numéro et le suivant, sauf mention expresse du contraire, les
anneaux et modules considérés sont supposés munis de la topologie
\emph{discrète}.
\end{env}

\begin{env}[20.5.2]
\label{0.20.5.2-fr}
Soient $A$ un anneau, $B$, $C$ deux $A$-algèbres, $u:B\to C$ un
$A$-homomorphisme ; on a un diagramme commutatif
\[
\xymatrix{
  B\otimes_A B \ar[r]^{u\otimes u} \ar[d]_{p_{B/A}} &
  C\otimes_A C \ar[d]^{p_{C/A}} \\
  B \ar[r]^u & C
}
\tag{20.5.2.1}
\label{0.20.5.2.1-fr}
\]
d'où par passage aux quotients, un $A$-homomorphisme d'algèbres
\begin{equation}
\tag{20.5.2.2}
\label{0.20.5.2.2-fr}
  u':\mathrm{P}^1_{B/A}\to\mathrm{P}^1_{C/A}
\end{equation}
tel que le diagramme
\[
\xymatrix{
  \mathrm{P}^1_{B/A} \ar[r]^{u'} & \mathrm{P}^1_{C/A} \\
  B \ar[u]^{p_1} \ar[r]^u & C \ar[u]_{p_1}
}
\]
soit commutatif ; comme $u\otimes u$ applique $\mathfrak{J}_{B/A}$ dans
$\mathfrak{J}_{C/A}$, on obtient, en restreignant $u'$ à
$\Omega^1_{B/A}$, une application
\begin{equation}
\tag{20.5.2.3}
\label{0.20.5.2.3-fr}
  u'':\Omega^1_{B/A}\to\Omega^1_{C/A}
\end{equation}
telle que le couple $(u'',u)$ soit un \emph{di-homomorphisme} pour la
structure de $B$-module de $\Omega^1_{B/A}$ et la structure de $C$-module de
$\Omega^1_{C/A}$ ; ce dernier fait permet d'en déduire canoniquement un
\emph{homomorphisme de $C$-modules}
\begin{equation}
\tag{20.5.2.4}
\label{0.20.5.2.4-fr}
  u_{C/B/A}:\Omega^1_{B/A}\otimes_B C\to\Omega^1_{C/A}.
\end{equation}

\oldpage[0\textsubscript{IV-1}]{129}
En outre, comme le diagramme
\[
\xymatrix{
  \mathrm{P}^1_{B/A} \ar[r]^{u'} & \mathrm{P}^1_{C/A} \\
  B \ar[u]^{p_2} \ar[r]^u & C \ar[u]_{p_2}
}
\tag{20.5.2.5}
\label{0.20.5.2.5-fr}
\]
est aussi commutatif, on en déduit que le diagramme
\[
\xymatrix{
  \Omega^1_{B/A} \ar[r]^{u''} & \Omega^1_{C/A} \\
  B \ar[u]^{d_{B/A}} \ar[r]^u & C \ar[u]_{d_{C/A}}
}
\tag{20.5.2.6}
\label{0.20.5.2.6-fr}
\]
est commutatif.

Enfin, si $w:C\to D$ est un second homomorphisme de $A$-algèbres, on a la
propriété de transitivité
\begin{equation}
\tag{20.5.2.7}
\label{0.20.5.2.7-fr}
  (w\circ u)_{D/B/A}
  =w_{D/C/A}\circ(u_{C/B/A}\otimes\mathrm{I}_D)
\end{equation}
comme il résulte de la définition.
\end{env}

\begin{env}[20.5.3]
\label{0.20.5.3-fr}
Soient maintenant $A$, $B$ deux anneaux, $v:A\to B$ un homomorphisme
d'anneaux, $C$ une $B$-algèbre qui devient une $A$-algèbre au moyen de $v$ ;
alors l'application canonique $v_0:C\otimes_A C\to C\otimes_B C$ est un
\emph{di-homomorphisme} surjectif d'algèbres (relatif à $v:A\to B$) tel que
le diagramme
\[
\xymatrix{
  C\otimes_A C \ar[r]^{v_0} \ar[d]_{p_{C/A}} &
  C\otimes_B C \ar[d]^{p_{C/B}} \\
  C \ar[r]^{\mathrm{I}_C} & C
}
\tag{20.5.3.1}
\label{0.20.5.3.1-fr}
\]
soit commutatif ; par passage aux quotients, on en déduit un
di-homomorphisme d'algèbres
\begin{equation}
\tag{20.5.3.2}
\label{0.20.5.3.2-fr}
  v':\mathrm{P}^1_{C/A}\to\mathrm{P}^1_{C/B}
\end{equation}
tel que le diagramme
\[
\xymatrix{
  \mathrm{P}^1_{C/A} \ar[r]^{v'} & \mathrm{P}^1_{C/B} \\
  C \ar[u]^{p_1} \ar[r]^{\mathrm{I}_C} & C \ar[u]_{p_1}
}
\]
soit commutatif. Comme $v_0$ applique $\mathfrak{J}_{C/A}$ dans
$\mathfrak{J}_{C/B}$ on obtient, en restreignant $v'$ à
$\Omega^1_{C/A}$, une application
\begin{equation}
\tag{20.5.3.3}
\label{0.20.5.3.3-fr}
  v_{C/B/A}:\Omega^1_{C/A}\to\Omega^1_{C/B}
\end{equation}
qui est un homomorphisme de $C$-modules.

\oldpage[0\textsubscript{IV-1}]{130}
En outre, comme le diagramme
\[
\xymatrix{
  \mathrm{P}^1_{C/A} \ar[r]^{v'} & \mathrm{P}^1_{C/B} \\
  C \ar[u]^{p_2} \ar[r]^{\mathrm{I}_C} & C \ar[u]_{p_2}
}
\tag{20.5.3.4}
\label{0.20.5.3.4-fr}
\]
est aussi commutatif, on en déduit que le diagramme
\[
\xymatrix{
  \Omega^1_{C/A} \ar[r]^{v_{C/B/A}} & \Omega^1_{C/B} \\
  C \ar[u]^{d_{C/A}} \ar[r]^{\mathrm{I}_C} & C \ar[u]_{d_{C/B}}
}
\tag{20.5.3.5}
\label{0.20.5.3.5-fr}
\]
est commutatif.

Enfin, si $s:A'\to A$ est un second homomorphisme d'anneaux, on a la
propriété de transitivité
\begin{equation}
\tag{20.5.3.6}
\label{0.20.5.3.6-fr}
  (v\circ s)_{C/B/A'}=v_{C/B/A}\circ s_{C/A/A'}.
\end{equation}
\end{env}

\begin{env}[20.5.4]
\label{0.20.5.4-fr}
Si maintenant on a un diagramme commutatif d'homomorphismes d'anneaux
\[
\xymatrix{
  B \ar[r]^v & B' \\
  A \ar[u]^f \ar[r]^u & A' \ar[u]_g
}
\]
on déduit de \sref{0.20.5.2.4-fr} et \sref{0.20.5.3.3-fr}, par composition,
un homomorphisme de $B'$-modules
\begin{equation}
\tag{20.5.4.1}
\label{0.20.5.4.1-fr}
  \Omega^1_{B/A}\otimes_B B'
  \xrightarrow{v_{B'/B/A}}\Omega^1_{B'/A}
  \xrightarrow{u_{B'/A'/A}}\Omega^1_{B'/A'}
\end{equation}
tel que le diagramme de $A'$-homomorphismes
\[
\xymatrix{
  \Omega^1_{B/A}\otimes_B B' \ar[r] & \Omega^1_{B'/A'} \\
  B' \ar[u]^{d_{B/A}\otimes1} \ar[r]^{\mathrm{I}_{B'}} &
  B' \ar[u]_{d_{B'/A'}}
}
\tag{20.5.4.2}
\label{0.20.5.4.2-fr}
\]
soit commutatif.

L'homomorphisme \sref{0.20.5.4.1-fr} correspond d'ailleurs à un
di-homomorphisme de $B$-modules
\begin{equation}
\tag{20.5.4.3}
\label{0.20.5.4.3-fr}
  \Omega^1_{B/A}\to\Omega^1_{B'/A'}.
\end{equation}
\end{env}

\begin{proposition}[20.5.5]
\label{0.20.5.5-fr}
Si $A'$, $B$ sont deux $A$-algèbres et $B'=B\otimes_A A'$,
l'homomorphisme canonique \sref{0.20.5.4.1-fr}
\begin{equation}
\tag{20.5.5.1}
\label{0.20.5.5.1-fr}
  \Omega^1_{B/A}\otimes_B B'\to\Omega^1_{B'/A'}
\end{equation}
est bijectif.
\end{proposition}

\begin{proof}
\oldpage[0\textsubscript{IV-1}]{131}
Le premier membre de \sref{0.20.5.5.1-fr} n'est autre alors que
$\Omega^1_{B/A}\otimes_A A'$. On peut écrire
$B'\otimes_{A'}B'=(B\otimes_A B)\otimes_A A'$ à un isomorphisme canonique
près et $p_{B'/A'}$ s'identifie alors à
$p_{B/A}\otimes\mathrm{I}_{A'}$ ; par suite (comme $p_{B/A}$ est surjectif)
$\mathfrak{J}_{B'/A'}=\operatorname{Im}(\mathfrak{J}_{B/A}\otimes_A A')$
et
$\mathfrak{J}^2_{B'/A'}
=\operatorname{Im}(\mathfrak{J}^2_{B/A}\otimes_A A')$, d'où
$\mathrm{P}^1_{B'/A'}=\mathrm{P}^1_{B/A}\otimes_A A'$ à un isomorphisme
canonique près, qui transforme en eux-mêmes les idéaux d'augmentation ;
comme le $A$-module $\mathrm{P}^1_{B/A}$ s'identifie canoniquement à la
somme directe de $B$ et de $\Omega^1_{B/A}$, on a
\begin{equation}
\tag{20.5.5.2}
\label{0.20.5.5.2-fr}
  \Omega^1_{B/A}\otimes_A A'=\Omega^1_{B'/A'}
\end{equation}
par le même isomorphisme, et on vérifie aussitôt que le composé de cet
isomorphisme et de l'isomorphisme canonique
$\Omega^1_{B/A}\otimes_B B'\xrightarrow{\sim}
\Omega^1_{B/A}\otimes_A A'$ n'est autre que \sref{0.20.5.5.1-fr}.
\end{proof}

\begin{env}[20.5.6]
\label{0.20.5.6-fr}
Les homomorphismes canoniques \sref{0.20.5.2.4-fr} et
\sref{0.20.5.3.3-fr} donnent, par fonctorialité, pour tout $C$-module $L$,
des homomorphismes canoniques
\begin{align}
\tag{20.5.6.1}
\label{0.20.5.6.1-fr}
  &\operatorname{Hom}_C(\Omega^1_{C/A},L)
  \to\operatorname{Hom}_C(\Omega^1_{B/A}\otimes_B C,L)
  =\operatorname{Hom}_B(\Omega^1_{B/A},L) \\
\tag{20.5.6.2}
\label{0.20.5.6.2-fr}
  &\operatorname{Hom}_C(\Omega^1_{C/B},L)
  \to\operatorname{Hom}_C(\Omega^1_{C/A},L).
\end{align}
Compte tenu de \sref{0.20.4.8.2-fr} et des diagrammes commutatifs
\sref{0.20.5.2.6-fr} et \sref{0.20.5.3.5-fr}, ces homomorphismes ne sont
autres (à une identification canonique près) que les homomorphismes
\sref{0.20.2.1.2-fr} et \sref{0.20.2.1.3-fr} respectivement.
\end{env}

\begin{theorem}[20.5.7]
\label{0.20.5.7-fr}
Soient $u:A\to B$, $v:B\to C$ deux homomorphismes d'anneaux.
\begin{enumerate}
  \item[\rm(i)] La suite de $C$-modules
\begin{equation}
\tag{20.5.7.1}
\label{0.20.5.7.1-fr}
  \Omega^1_{B/A}\otimes_B C
  \xrightarrow{v_{C/B/A}}\Omega^1_{C/A}
  \xrightarrow{u_{C/B/A}}\Omega^1_{C/B}\to0
\end{equation}
est exacte.
  \item[\rm(ii)] Pour que $v_{C/B/A}$ soit inversible à gauche, il faut et
  il suffit que $C$ soit une $B$-algèbre formellement lisse relativement à
  $A$ (pour les topologies discrètes (cf. \sref{0.19.9.1-fr})) ; en
  particulier, il suffit pour cela que $C$ soit une $B$-algèbre formellement
  lisse (pour les topologies discrètes).
\end{enumerate}
\end{theorem}

\begin{proof}
(i) L'exactitude de la suite \sref{0.20.2.4.2-fr} montre tout d'abord,
compte tenu de \sref{0.20.5.6-fr}, que la suite
\[
  0\to\operatorname{Hom}_C(\Omega^1_{C/B},L)
  \to\operatorname{Hom}_C(\Omega^1_{C/A},L)
  \to\operatorname{Hom}_C(\Omega^1_{B/A}\otimes_B C,L)
\]
est exacte pour tout $C$-module $L$. On sait que cela implique l'exactitude
de la suite \sref{0.20.5.7.1-fr} (Bourbaki, \emph{Alg.}, chap.~II,
3\textsuperscript{e} éd., §~2, n\textsuperscript{o}~1, th.~1).

(ii) En vertu de l'exactitude de \sref{0.20.5.7.1-fr}, dire que
$v_{C/B/A}$ est inversible à gauche signifie que la suite
\begin{equation}
\tag{20.5.7.2}
\label{0.20.5.7.2-fr}
  0\to\Omega^1_{B/A}\otimes_B C
  \xrightarrow{v_{C/B/A}}\Omega^1_{C/A}
  \xrightarrow{u_{C/B/A}}\Omega^1_{C/B}\to0
\end{equation}
est \emph{exacte et scindée} ; on sait (Bourbaki, \emph{loc. cit.},
n\textsuperscript{o}~1, prop.~1) que cela équivaut à dire que pour tout
$C$-module $L$, la suite
\[
  0\to\operatorname{Hom}_C(\Omega^1_{C/B},L)
  \to\operatorname{Hom}_C(\Omega^1_{C/A},L)
  \to\operatorname{Hom}_C(\Omega^1_{B/A}\otimes_B C,L)\to0
\]
est exacte ; compte tenu de \sref{0.20.5.6-fr} et
\sref{0.20.2.4.2-fr}, cette condition équivaut à
\[
  \operatorname{Exalcom}_{B/A}(C,L)=0
\]
pour tout $C$-module $L$, et la conclusion résulte donc de
\sref{0.19.9.8.1-fr}.
\end{proof}

\oldpage[0\textsubscript{IV-1}]{132}
Notons en outre que si l'on a un diagramme commutatif d'homomorphismes
d'anneaux
\[
\xymatrix{
  A' \ar[r] & B' \ar[r] & C' \\
  A \ar[u] \ar[r] & B \ar[u] \ar[r] & C \ar[u]
}
\]
on en déduit un diagramme commutatif
\[
\xymatrix{
  \Omega^1_{B/A}\otimes_B C \ar[r] \ar[d] &
  \Omega^1_{C/A} \ar[r] \ar[d] &
  \Omega^1_{C/B} \ar[r] \ar[d] & 0 \\
  \Omega^1_{B'/A'}\otimes_{B'} C' \ar[r] &
  \Omega^1_{C'/A'} \ar[r] &
  \Omega^1_{C'/B'} \ar[r] & 0
}
\tag{20.5.7.3}
\label{0.20.5.7.3-fr}
\]
où les flèches verticales proviennent des di-homomorphismes
\sref{0.20.5.4.3-fr}.

\begin{corollary}[20.5.8]
\label{0.20.5.8-fr}
Supposons que l'homomorphisme $v:B\to C$ fasse de $C$ une $B$-algèbre
formellement étale (pour les topologies discrètes
\sref{0.19.10.2-fr}) ; alors l'homomorphisme \sref{0.20.5.3.3-fr}
\[
  v_{C/B/A}:\Omega^1_{B/A}\otimes_B C\to\Omega^1_{C/A}
\]
est bijectif.
\end{corollary}

\begin{proof}
En effet, si $C$ est une $B$-algèbre formellement non ramifiée pour les
topologies discrètes, il résulte de \sref{0.19.10.4-fr},
\sref{0.20.4.8-fr} et \sref{0.20.1.1-fr} que l'on a
$\operatorname{Hom}_C(\Omega^1_{C/B},L)=0$ pour tout $C$-module $L$, donc
$\Omega^1_{C/B}=0$ (cf. (20.7.4)) ; d'autre part, si $C$ est une
$B$-algèbre formellement lisse pour les topologies discrètes, la suite
\sref{0.20.5.7.2-fr} est exacte ; d'où le corollaire.
\end{proof}

\begin{corollary}[20.5.9]
\label{0.20.5.9-fr}
Soient $A$ un anneau, $B$ une $A$-algèbre, $S$ une partie multiplicative de
$B$ ; alors l'homomorphisme canonique
\begin{equation}
\tag{20.5.9.1}
\label{0.20.5.9.1-fr}
  S^{-1}\Omega^1_{B/A}\to\Omega^1_{S^{-1}B/A}
\end{equation}
est bijectif.
\end{corollary}

\begin{proof}
Il suffit d'appliquer \sref{0.20.5.8-fr} à $C=S^{-1}B$, qui est une
$B$-algèbre formellement étale pour les topologies discrètes
\sref{0.19.10.3-fr}, (ii).
\end{proof}

Compte tenu de \sref{0.20.5.5-fr}, on peut donc écrire
\begin{equation}
\tag{20.5.9.2}
\label{0.20.5.9.2-fr}
  \Omega^1_{S^{-1}B/S^{-1}A}
  =S^{-1}\Omega^1_{B/A}
  =\Omega^1_{S^{-1}B/A},
\end{equation}
à des isomorphismes canoniques près.

\begin{corollary}[20.5.10]
\label{0.20.5.10-fr}
Si $k$ est un corps et $K=k(X_\alpha)_{\alpha\in I}$ une extension pure de
$k$, les $dX_\alpha$ forment une base du $K$-espace vectoriel
$\Omega^1_{K/k}$.
\end{corollary}

\begin{proof}
Comme $K$ est le corps des fractions de l'anneau de polynômes
$k[X_\alpha]_{\alpha\in I}$, cela résulte de \sref{0.20.4.13-fr}, (i) et de
\sref{0.20.5.9-fr}.
\end{proof}

\oldpage[0\textsubscript{IV-1}]{133}
\begin{env}[20.5.11]
\label{0.20.5.11-fr}
Soient $A$ un anneau, $B$ une $A$-algèbre, $\mathfrak{R}$ un idéal de $B$,
$C$ la $A$-algèbre quotient $B/\mathfrak{R}$, et considérons
l'homomorphisme composé de $A$-modules
\begin{equation}
\tag{20.5.11.1}
\label{0.20.5.11.1-fr}
  \delta:\mathfrak{R}\longrightarrow B\xrightarrow{d_{B/A}}\Omega^1_{B/A}
\end{equation}
où la première flèche est l'injection canonique; comme
$d(xy)=xdy+ydx$, on voit que l'on a
$\delta(\mathfrak{R}^2)\subset\mathfrak{R}.\Omega^1_{B/A}$, d'où, par
passage aux quotients, un homomorphisme de $A$-modules
\begin{equation}
\tag{20.5.11.2}
\label{0.20.5.11.2-fr}
  \delta_{C/B/A}:\mathfrak{R}/\mathfrak{R}^2\longrightarrow
  \Omega^1_{B/A}\otimes_B C=
  \Omega^1_{B/A}/(\mathfrak{R}.\Omega^1_{B/A}).
\end{equation}
Mais en fait, $\delta_{C/B/A}$ est un homomorphisme de $C$-modules, car
pour $x\in B$ et $y\in\mathfrak{R}$, on a
$ydx\in\mathfrak{R}.\Omega^1_{B/A}$, donc
$d(xy)\equiv xdy\pmod{\mathfrak{R}.\Omega^1_{B/A}}$, ce qui prouve
d'abord que (20.5.11.2) est un homomorphisme de $B$-modules, et comme
$\mathfrak{R}$ annule les deux membres, cela établit notre assertion.

Si $B'$ est une seconde $A$-algèbre, $u:B\to B'$ un $A$-homomorphisme,
$\mathfrak{R}'$ un idéal de $B'$ tel que
$u(\mathfrak{R})\subset\mathfrak{R}'$ et $C'=B'/\mathfrak{R}'$ l'algèbre
quotient, on a un diagramme commutatif
\begin{equation}
\tag{20.5.11.3}
\label{0.20.5.11.3-fr}
\xymatrix{
  \mathfrak{R}/\mathfrak{R}^2 \ar[r]^-{\delta_{C/B/A}} \ar[d] &
  \Omega^1_{B/A}\otimes_B C \ar[d] \\
  \mathfrak{R}'/\mathfrak{R}'^2 \ar[r]_-{\delta_{C'/B'/A}} &
  \Omega^1_{B'/A}\otimes_{B'} C'
}
\end{equation}
où les flèches verticales proviennent de $u$ (20.5.2.4).
\end{env}

\begin{theorem}[20.5.12]
\label{0.20.5.12-fr}
Soient $A$ un anneau, $B$ une $A$-algèbre, $C$ la $A$-algèbre quotient
$B/\mathfrak{R}$, $u:B\to C$ l'homomorphisme canonique.
\begin{enumerate}
  \item[\rm(i)] On a une suite exacte de $C$-modules
\begin{equation}
\tag{20.5.12.1}
\label{0.20.5.12.1-fr}
  \mathfrak{R}/\mathfrak{R}^2
  \xrightarrow{\delta_{C/B/A}}\Omega^1_{B/A}\otimes_B C
  \xrightarrow{u_{C/B/A}}\Omega^1_{C/A}\longrightarrow 0,
\end{equation}
où $u_{C/B/A}$ et $\delta_{C/B/A}$ sont définis par (20.5.2.4) et
(20.5.11.2) respectivement.
  \item[\rm(ii)] Si l'on pose $E=B/\mathfrak{R}^2$, l'homomorphisme
canonique (20.5.2.4)
\[
  \Omega^1_{B/A}\otimes_B C\longrightarrow
  \Omega^1_{E/A}\otimes_E C
\]
est bijectif.
  \item[\rm(iii)] Les trois conditions suivantes sont équivalentes :
  \begin{enumerate}
    \item[\rm a)] $\delta_{C/B/A}$ est inversible à gauche.
    \item[\rm b)] Toute $A$-extension de $C$ par un $C$-module $L$, dont
l'image réciproque par $u:B\to C$ est $A$-triviale, est elle-même
$A$-triviale.
    \item[\rm c)] La $A$-algèbre $E=B/\mathfrak{R}^2$ est une extension
$A$-triviale de $C$ par $\mathfrak{R}/\mathfrak{R}^2$.
  \end{enumerate}
  \item[\rm(iv)] Il y a correspondance biunivoque canonique entre les
inverses à gauche de $\delta_{C/B/A}$ et les inverses à droite de
l'homomorphisme canonique $E\to C$.
\end{enumerate}
\end{theorem}

\begin{proof}
\oldpage[0\textsubscript{IV-1}]{134}
(i) Comme $u$ est surjectif, on a $\operatorname{D\acute{e}r}_B(C,L)=0$ pour tout $C$-module
$L$ en vertu de (20.1.1). La suite exacte (20.2.3.1) devient donc
\begin{equation}
\tag{20.5.12.2}
\label{0.20.5.12.2-fr}
  0\longrightarrow\operatorname{D\acute{e}r}_A(C,L)\xrightarrow{u^0}\operatorname{D\acute{e}r}_A(B,L)
  \xrightarrow{\partial}\operatorname{Exalcom}_B(C,L)
  \xrightarrow{v^1}\operatorname{Exalcom}_A(C,L)
  \xrightarrow{u^1}\operatorname{Exalcom}_A(B,L),
\end{equation}
où $v$ est l'homomorphisme $A\to B$. Rappelons d'autre part (18.3.8) que
$\operatorname{Exalcom}_B(C,L)$ s'identifie canoniquement à
$\operatorname{Hom}_C(\mathfrak{R}/\mathfrak{R}^2,L)$; on déduit donc de (20.5.12.2)
et (20.4.8) la suite exacte
\begin{equation}
\tag{20.5.12.3}
\label{0.20.5.12.3-fr}
\begin{split}
  0\longrightarrow\operatorname{Hom}_C(\Omega^1_{C/A},L)
  \longrightarrow\operatorname{Hom}_C(\Omega^1_{B/A}\otimes_B C,L)
  \xrightarrow{\varphi}\operatorname{Hom}_C(\mathfrak{R}/\mathfrak{R}^2,L)\\
  \xrightarrow{\psi}\operatorname{Ker}(u^1)\longrightarrow 0
\end{split}
\end{equation}
avec $\varphi=\partial\circ\eta^{-1}$ et $\psi=v^1\circ\eta$. Revenant
aux définitions de $\partial$ (20.2.2) et de $\eta$ (18.3.8), on constate
aussitôt que $\varphi$ est précisément l'homomorphisme
$\operatorname{Hom}(\delta_{C/B/A},1_L)$. L'existence de la suite exacte formée des
4 premiers termes de (20.5.12.3) montre donc que la suite (20.5.12.1)
est exacte (Bourbaki, \emph{Alg.}, chap. II, 3\textsuperscript{e} éd.,
\S~2, n\textsuperscript{o}~1, th.~1).

\medskip
(ii) Appliquons à $B$ et à l'idéal $\mathfrak{R}^2$ la suite exacte
(20.5.12.1), ce qui donne
\begin{equation}
\tag{20.5.12.4}
\label{0.20.5.12.4-fr}
  \mathfrak{R}^2/\mathfrak{R}^4\longrightarrow
  \Omega^1_{B/A}\otimes_B E\longrightarrow\Omega^1_{E/A}\longrightarrow 0,
\end{equation}
d'où, en tensorisant par $C$ (considéré comme $E$-algèbre), la suite
exacte
\[
  \mathfrak{R}^2/\mathfrak{R}^4\otimes_E C\longrightarrow
  \Omega^1_{B/A}\otimes_B C\longrightarrow
  \Omega^1_{E/A}\otimes_E C\longrightarrow 0.
\]
Or, si $x,y$ sont deux éléments de $\mathfrak{R}$, et $\zeta$ la classe
de $xy$ mod. $\mathfrak{R}^4$, l'image
$\delta'(\zeta\otimes1)$ est par définition
$d_{B/A}(xy)\otimes1=(xd_{B/A}(y)+yd_{B/A}(x))\otimes1$; mais comme les
images de $x$ et de $y$ dans $C$ sont nulles, on a aussi
$d_{B/A}(xy)\otimes1=0$ dans $\Omega^1_{B/A}\otimes_B C$, ce qui prouve
notre assertion.

\medskip
(iii) Dire que $\delta_{C/B/A}$ est inversible à gauche revient, compte
tenu de l'exactitude de (20.5.12.1), à dire que la suite
\begin{equation}
\tag{20.5.12.5}
\label{0.20.5.12.5-fr}
  0\longrightarrow\mathfrak{R}/\mathfrak{R}^2
  \xrightarrow{\delta_{C/B/A}}\Omega^1_{B/A}\otimes_B C
  \xrightarrow{u_{C/B/A}}\Omega^1_{C/A}\longrightarrow0
\end{equation}
est \emph{exacte et scindée}, et il revient au même (Bourbaki,
\emph{loc. cit.}) de dire que $\operatorname{Ker}(u^1)=0$ dans la suite exacte
(20.5.12.3) pour tout $L$, ce qui montre l'équivalence des conditions
a) et b) (cf. (18.3.6.2)).

Le fait que b) entraîne c) provient de ce que l'image réciproque par
$u:B\to C$ de la $A$-extension $E$ de $C$ par
$\mathfrak{R}/\mathfrak{R}^2$ est $A$-triviale, $u$ étant composée des
homomorphismes canoniques $B\to E\to C$ (18.1.6). Inversement, c)
entraîne b), car toute $B$-extension de $C$ par $L$ est $B$-équivalente
à $E\oplus_{\mathfrak{R}/\mathfrak{R}^2}L$ (18.3.8), autrement dit sa
classe est l'image de la classe de $E$ (considéré comme $B$-extension)
par l'homomorphisme

\oldpage[0\textsubscript{IV-1}]{135}
$w_*:\operatorname{Exan}_B(C,\mathfrak{R}/\mathfrak{R}^2)
\to\operatorname{Exan}_B(C,L)$ correspondant à un $B$-homomorphisme
$w:\mathfrak{R}/\mathfrak{R}^2\to L$. Le fait que c) entraîne b) résulte
alors de la commutativité du diagramme (18.3.6.5)
\[
\xymatrix{
  \operatorname{Exan}_B(C,\mathfrak{R}/\mathfrak{R}^2)
    \ar[r]^{w_*} \ar[d]_{u^*} &
  \operatorname{Exan}_B(C,L) \ar[d]^{u^*} \\
  \operatorname{Exan}_A(C,\mathfrak{R}/\mathfrak{R}^2)
    \ar[r]_{w_*} & \operatorname{Exan}_A(C,L).
}
\]

\medskip
(iv) On a vu (20.1.7) que les inverses à droite de $E\to C$
correspondent biunivoquement de façon canonique à l'ensemble des éléments
$D\in\operatorname{D\acute{e}r}_A(E,\mathfrak{R}/\mathfrak{R}^2)$ tels que $D(x)=x$ dans
$\mathfrak{R}/\mathfrak{R}^2$, donc aussi, par (20.4.8), à l'ensemble des
$E$-homomorphismes
$h:\Omega^1_{E/A}\to\mathfrak{R}/\mathfrak{R}^2$ tels que le composé
\[
  \mathfrak{R}/\mathfrak{R}^2\xrightarrow{d_{E/A}}
  \Omega^1_{E/A}\xrightarrow{h}\mathfrak{R}/\mathfrak{R}^2
\]
soit l'identité. Par tensorisation avec $C$, on en déduit (puisque
$\mathfrak{R}/\mathfrak{R}^2$ est un $C$-module) que le composé
\[
  \mathfrak{R}/\mathfrak{R}^2\xrightarrow{\delta_{C/E/A}}
  \Omega^1_{E/A}\otimes_E C\xrightarrow{h\otimes1}
  \mathfrak{R}/\mathfrak{R}^2
\]
est l'identité; or, puisque $\mathfrak{R}/\mathfrak{R}^2$ est un
$C$-module, $h\mapsto h\otimes1$ est un isomorphisme de l'ensemble
$\operatorname{Hom}_E(\Omega^1_{E/A},\mathfrak{R}/\mathfrak{R}^2)$ sur
$\operatorname{Hom}_C(\Omega^1_{E/A}\otimes_E C,\mathfrak{R}/\mathfrak{R}^2)$; et par
ailleurs (ii) prouve que l'on peut identifier canoniquement
$\Omega^1_{E/A}\otimes_E C$ et $\Omega^1_{B/A}\otimes_B C$,
$\delta_{C/E/A}$ s'identifiant alors à $\delta_{C/B/A}$. C.Q.F.D.
\end{proof}

\begin{example}[20.5.13]
\label{0.20.5.13-fr}
Soient $B=A[X_\alpha]_{\alpha\in I}$ une algèbre de polynômes sur $A$,
$\mathfrak{R}$ un idéal de $B$, $(P_\lambda)$ un système de générateurs
de $\mathfrak{R}$ et $C=B/\mathfrak{R}$; on sait que
$\Omega^1_{B/A}$ est un $B$-module libre dont les $dX_\alpha$ forment une
base (20.4.13, (i)), donc les $dX_\alpha$ forment aussi une base du
$C$-module libre $\Omega^1_{B/A}\otimes_B C$. D'autre part, il résulte
aussitôt de la définition que l'image de
$\mathfrak{R}/\mathfrak{R}^2$ par $\delta_{C/B/A}$ est le sous-$C$-module
engendré par les
\[
  dP_\lambda=\sum_\alpha
  \frac{\partial P_\lambda}{\partial X_\alpha}\,dX_\alpha.
\]
On en conclut que $\Omega^1_{C/A}$ est isomorphe au quotient du
$C$-module libre ayant les $dX_\alpha$ pour base, par le sous-$C$-module
engendré par les $dP_\lambda$, ce qui donne une description d'un module
de différentielles d'une algèbre \emph{quelconque}, toute $A$-algèbre
$C$ pouvant s'obtenir de la façon précédente.
\end{example}

\begin{corollary}[20.5.14]
\label{0.20.5.14-fr}
Si $C$ est une $A$-algèbre formellement lisse (pour les topologies
discrètes), la suite
\begin{equation}
\tag{20.5.14.1}
\label{0.20.5.14.1-fr}
  0\longrightarrow\mathfrak{R}/\mathfrak{R}^2
  \xrightarrow{\delta_{C/B/A}}\Omega^1_{B/A}\otimes_B C
  \xrightarrow{u_{C/B/A}}\Omega^1_{C/A}\longrightarrow0
\end{equation}
est exacte et scindée.
\end{corollary}

\begin{proof}
En effet, toute $A$-extension de $C$ par un $C$-module est alors triviale
(19.4.4.1).
\end{proof}

\oldpage[0\textsubscript{IV-1}]{136}
\begin{remark}[20.5.15]
\label{0.20.5.15-fr}
Soit $u:A\to B$ un homomorphisme \emph{surjectif} d'anneaux; alors, pour
tout homomorphisme d'anneaux $v:B\to C$, l'homomorphisme canonique
\begin{equation}
\tag{20.5.15.1}
\label{0.20.5.15.1-fr}
  u_{C/B/A}:\Omega^1_{C/A}\longrightarrow\Omega^1_{C/B}
\end{equation}
est \emph{bijectif}; cela résulte en effet de la suite exacte (20.5.7.1),
puisque $\Omega^1_{B/A}=0$ (20.4.12).
\end{remark}

\subsection{Modules d'imperfection et homomorphismes caractéristiques}
\label{subsection:0.20.6-fr}
\label{0.20.6-fr}

\begin{definition}[20.6.1]
\label{0.20.6.1-fr}
Étant donnés deux homomorphismes d'anneaux $u:A\to B$, $v:B\to C$, on
appelle \emph{module d'imperfection de la $B$-algèbre $C$ relativement à
$A$}, et on note $\Upsilon_{C/B/A}$, le $C$-module noyau de
l'homomorphisme
$v_{C/B/A}:\Omega^1_{B/A}\otimes_B C\to\Omega^1_{C/A}$.
On a donc par définition (cf. (20.5.7)) la suite exacte
\begin{equation}
\tag{20.6.1.1}
\label{0.20.6.1.1-fr}
  0\longrightarrow\Upsilon_{C/B/A}\longrightarrow
  \Omega^1_{B/A}\otimes_B C\xrightarrow{v_{C/B/A}}
  \Omega^1_{C/A}\xrightarrow{u_{C/B/A}}\Omega^1_{C/B}
  \longrightarrow0.
\end{equation}
\end{definition}

Lorsque $A=\mathbb{Z}$ (les modules $\Omega^1_{B/A}$ et $\Omega^1_{C/A}$
étant donc les modules de différentielles « absolues » $\Omega^1_B$ et
$\Omega^1_C$), on écrit $\Upsilon_{C/B}$ au lieu de
$\Upsilon_{C/B/\mathbb{Z}}$; lorsque $B$ et $C$ sont des algèbres sur un
corps premier $P$, on a $\Upsilon_{C/B/P}=\Upsilon_{C/B}$.

Soient $R$, $S$ des parties multiplicatives de $B$ et $C$ respectivement,
telles que l'image de $R$ soit contenue dans $S$. Il résulte alors de la
suite exacte (20.6.1.1) et de (20.5.9) que l'on a
\begin{equation}
\tag{20.6.1.2}
\label{0.20.6.1.2-fr}
  \Upsilon_{S^{-1}C/R^{-1}B/A}=S^{-1}\Upsilon_{C/B/A}.
\end{equation}

\begin{proposition}[20.6.2]
\label{0.20.6.2-fr}
Si $C$ est une $B$-algèbre formellement lisse relativement à $A$ (et en
particulier si $C$ est une $B$-algèbre formellement lisse), on a
$\Upsilon_{C/B/A}=0$.
\end{proposition}

\begin{proof}
Cela résulte de (20.5.7, (ii)).
\end{proof}

\begin{proposition}[20.6.3]
\label{0.20.6.3-fr}
Soient $k$ un corps, $K$ une extension de $k$. Pour que $K$ soit une
extension séparable de $k$, il faut et il suffit que $\Upsilon_{K/k}=0$
(autrement dit, l'homomorphisme canonique
$\Omega^1_k\otimes_k K\to\Omega^1_K$ est injectif ou encore (20.4.8),
toute dérivation de $k$ dans $K$ se prolonge en une dérivation de $K$
dans lui-même).
\end{proposition}

\begin{proof}
En effet, il est équivalent de dire que $K$ est séparable sur $k$ ou une
$k$-algèbre formellement lisse (19.6.1); d'autre part, si $P$ est le
corps premier de $k$, $K$ est séparable sur $P$, donc une $P$-algèbre
formellement lisse, et il revient donc au même de dire que $K$ est une
$k$-algèbre formellement lisse ou une $k$-algèbre formellement lisse
relativement à $P$. Enfin, dire que $K$ est une $k$-algèbre formellement
lisse relativement à $P$ équivaut d'après (20.5.7, (ii)) à dire que
l'homomorphisme
$v_{K/k/P}:\Omega^1_k\otimes_k K\to\Omega^1_K$ est inversible à gauche;
mais comme $K$ est un corps, cette dernière condition équivaut à dire
que le noyau de $v_{K/k/P}$, c'est-à-dire $\Upsilon_{K/k}$, est nul.
\end{proof}

\begin{env}[20.6.4]
\label{0.20.6.4-fr}
Considérons un diagramme commutatif
\begin{equation}
\tag{20.6.4.1}
\label{0.20.6.4.1-fr}
\xymatrix{
  A' \ar[r]^{u'} & B' \ar[r]^{v'} & C' \\
  A \ar[u]^f \ar[r]_u & B \ar[u]^g \ar[r]_v & C \ar[u]^h
}
\end{equation}

\oldpage[0\textsubscript{IV-1}]{137}
d'homomorphismes d'anneaux commutatifs. La commutativité du diagramme
correspondant (20.5.7.3) entraîne l'existence d'un unique
$C$-homomorphisme
\begin{equation}
\tag{20.6.4.2}
\label{0.20.6.4.2-fr}
  \Upsilon_{C/B/A}\longrightarrow\Upsilon_{C'/B'/A'}
\end{equation}
déduit canoniquement de (20.6.4.1) et rendant commutatif le diagramme
\begin{equation}
\tag{20.6.4.3}
\label{0.20.6.4.3-fr}
\xymatrix{
  0 \ar[r] & \Upsilon_{C/B/A} \ar[r] \ar[d] &
  \Omega^1_{B/A}\otimes_B C \ar[r]^{v_{C/B/A}} \ar[d] &
  \Omega^1_{C/A} \ar[r]^{u_{C/B/A}} \ar[d] &
  \Omega^1_{C/B} \ar[r] \ar[d] & 0 \\
  0 \ar[r] & \Upsilon_{C'/B'/A'} \ar[r] &
  \Omega^1_{B'/A'}\otimes_{B'} C' \ar[r]_{v_{C'/B'/A'}} &
  \Omega^1_{C'/A'} \ar[r]_{u_{C'/B'/A'}} &
  \Omega^1_{C'/B'} \ar[r] & 0
}
\end{equation}

La donnée de l'homomorphisme (20.6.4.2) équivaut d'ailleurs à celle d'un
$C'$-homomorphisme
\begin{equation}
\tag{20.6.4.4}
\label{0.20.6.4.4-fr}
  \Upsilon_{C/B/A}\otimes_C C'\longrightarrow\Upsilon_{C'/B'/A'}
\end{equation}
qui, composé avec l'homomorphisme canonique
$\Upsilon_{C/B/A}\to\Upsilon_{C/B/A}\otimes_C C'$, redonne (20.6.4.2).
Il est clair que (20.6.4.2) jouit d'une propriété évidente de
transitivité, permettant de dire que $\Upsilon_{C/B/A}$ est un
\emph{foncteur} en le triplet $(A,B,C)$.
\end{env}

\begin{env}[20.6.5]
\label{0.20.6.5-fr}
Il nous sera commode pour la suite, sous les conditions de (20.6.1),
d'introduire un \emph{complexe} (de chaînes) \emph{de $C$-modules}
$\mathrm{K}_\bullet(C/B/A)$ dont les termes sont nuls sauf en degrés 0
et 1, où on prend
\begin{equation}
\tag{20.6.5.1}
\label{0.20.6.5.1-fr}
\left\{
\begin{aligned}
  \mathrm{K}_0(C/B/A)&=\Omega^1_{C/A},\\
  \mathrm{K}_1(C/B/A)&=\Omega^1_{B/A}\otimes_B C,
\end{aligned}
\right.
\end{equation}
l'opérateur différentiel $\mathrm{K}_1\to\mathrm{K}_0$ étant
$v_{C/B/A}$. Cela permet d'écrire (à des isomorphismes canoniques près)
$\Omega^1_{C/B}$ et $\Upsilon_{C/B/A}$ comme les \emph{modules
d'homologie} de ce complexe
\begin{equation}
\tag{20.6.5.2}
\label{0.20.6.5.2-fr}
  \mathrm{H}_0(\mathrm{K}_\bullet(C/B/A))=\Omega^1_{C/B},\qquad
  \mathrm{H}_1(\mathrm{K}_\bullet(C/B/A))=\Upsilon_{C/B/A}.
\end{equation}
\end{env}

De même :
\begin{proposition}[20.6.6]
\label{0.20.6.6-fr}
Sous les hypothèses de (20.6.1), pour tout $C$-module $L$, on a des
$C$-isomorphismes canoniques
\begin{align}
\tag{20.6.6.1}
\label{0.20.6.6.1-fr}
  \mathrm{H}^0(\mathrm{K}_\bullet(C/B/A),L)&\simeq\operatorname{D\acute{e}r}_B(C,L),\\
\tag{20.6.6.2}
\label{0.20.6.6.2-fr}
  \mathrm{H}^1(\mathrm{K}_\bullet(C/B/A),L)&\simeq
  \operatorname{Exalcom}_{B/A}(C,L).
\end{align}
\end{proposition}

\begin{proof}
En effet, le complexe de cochaînes
$\operatorname{Hom}_C(\mathrm{K}_\bullet(C/B/A),L)$ n'est autre, en vertu de (20.4.8)
et (20.5.6), que le complexe
\[
  \cdots\longrightarrow0\longrightarrow\operatorname{D\acute{e}r}_A(C,L)
  \longrightarrow\operatorname{D\acute{e}r}_A(B,L)\longrightarrow0\longrightarrow\cdots
\]
\oldpage[0\textsubscript{IV-1}]{138}
où l'opérateur différentiel est $v^0$ (avec les notations de (20.2.1)).
La proposition résulte alors de la suite exacte (20.2.4.2) et de la
définition des modules de cohomologie
\begin{equation}
\tag{20.6.6.3}
\label{0.20.6.6.3-fr}
  \mathrm{H}^i(\mathrm{K}_\bullet,L)=
  \mathrm{H}^i(\operatorname{Hom}_C(\mathrm{K}_\bullet,L)).
\end{equation}
\end{proof}

Si on a un diagramme commutatif d'homomorphismes d'anneaux
\[
\xymatrix{
  A' \ar[r] & B' \ar[r] & C' \\
  A \ar[u] \ar[r] & B \ar[u] \ar[r] & C \ar[u]
}
\]
les di-homomorphismes (20.5.4.3) définissent un di-homomorphisme de
complexes de modules
\begin{equation}
\tag{20.6.6.4}
\label{0.20.6.6.4-fr}
  \mathrm{K}_\bullet(C/B/A)\longrightarrow
  \mathrm{K}_\bullet(C'/B'/A')
\end{equation}
et les di-homomorphismes qu'on en déduit pour l'homologie ou la
cohomologie s'identifient, par les formules (20.6.5.2), (20.6.6.1) et
(20.6.6.2), aux di-homomorphismes déjà définis dans (20.5.4.3),
(20.6.4.2), (20.2.1) et (18.3.7.2) (compte tenu, pour le dernier, de la
définition de l'opérateur $\partial$ dans la suite exacte (20.2.4.2)).

\begin{env}[20.6.7]
\label{0.20.6.7-fr}
On sait que pour un complexe $\mathrm{K}_\bullet$ de $C$-modules et un
$C$-module $L$, on a des homomorphismes canoniques
\[
  \alpha_i:\mathrm{H}^i(\mathrm{K}_\bullet,L)\longrightarrow
  \operatorname{Hom}_C(\mathrm{H}_i(\mathrm{K}_\bullet),L)
\]
(M, IV, 6). Ici, l'homomorphisme canonique
\[
  \alpha_1:\mathrm{H}^1(\mathrm{K}_\bullet(C/B/A),L)
  \longrightarrow
  \operatorname{Hom}_C(\mathrm{H}_1(\mathrm{K}_\bullet(C/B/A)),L)
\]
se définit immédiatement comme obtenu par passage au quotient par l'image
de $\operatorname{Hom}_C(\mathrm{K}_0,L)$ de l'homomorphisme de restriction
\[
  \operatorname{Hom}_C(\mathrm{K}_1(C/B/A),L)\longrightarrow
  \operatorname{Hom}_C(\mathrm{H}_1(\mathrm{K}_\bullet(C/B/A)),L)
\]
puisque $\mathrm{H}_1(\mathrm{K}_\bullet(C/B/A))$ n'est autre que le
noyau de $\mathrm{K}_1\to\mathrm{K}_0$; tenant compte de (20.6.6.2) et
(20.6.5.2), on obtient donc un $C$-homomorphisme canonique
\begin{equation}
\tag{20.6.7.1}
\label{0.20.6.7.1-fr}
  \operatorname{Exalcom}_{B/A}(C,L)\longrightarrow
  \operatorname{Hom}_C(\Upsilon_{C/B/A},L)
\end{equation}
qui s'explicite de la façon suivante : en vertu de (20.2.4.2), toute
$B$-extension de $C$ par $L$ qui est $A$-triviale provient de la donnée
d'une $A$-dérivation $D$ de $B$ dans $L$, donc (20.4.8) d'un
$C$-homomorphisme $f$ de $\Omega^1_{B/A}\otimes_B C$ dans $L$; on fait
correspondre à la classe de cette extension la \emph{restriction} de $f$
à $\Upsilon_{C/B/A}$, qui ne dépend effectivement que de cette classe et
non du choix de $D$.
\end{env}

\begin{definition}[20.6.8]
\label{0.20.6.8-fr}
Soient $u:A\to B$, $v:B\to C$ deux homomorphismes d'anneaux. Pour toute
$B$-extension $E$ de $C$ par un $C$-module $L$, qui est $A$-triviale
(18.3.7), on appelle \emph{homomorphisme caractéristique de $E$}, et on
note $\chi_E$, le $C$-homomorphisme $\Upsilon_{C/B/A}\to L$, image de la
classe de $E$ par l'homomorphisme canonique (20.6.7.1).
\end{definition}

\oldpage[0\textsubscript{IV-1}]{139}
On peut définir l'homomorphisme $\chi_E$ d'une autre manière :
\begin{proposition}[20.6.9]
\label{0.20.6.9-fr}
Soit $E$ une $B$-extension de $C$ par un $C$-module $L$, qui est
$A$-triviale (18.3.7); alors le diagramme
\begin{equation}
\tag{20.6.9.1}
\label{0.20.6.9.1-fr}
\xymatrix{
  0 \ar[r] & \Upsilon_{C/B/A} \ar[r] \ar[d]^{\chi_E} &
  \Omega^1_{B/A}\otimes_B C \ar[r]^{v_{C/B/A}}
    \ar[d]^{q_{E/B/A}\otimes1_C} &
  \Omega^1_{C/A} \ar[d]^{1} \\
  0 \ar[r] & L \ar[r]_{\delta_{C/E/A}} &
  \Omega^1_{E/A}\otimes_E C \ar[r]_{p_{C/E/A}} &
  \Omega^1_{C/A} \ar[r] & 0
}
\end{equation}
où $q:B\to E$ définit la structure de $B$-extension de $E$ et
$p:E\to C$ est l'homomorphisme d'augmentation, est commutatif et ses
lignes sont exactes.
\end{proposition}

\begin{proof}
La ligne inférieure du diagramme est la suite (20.5.12.5) relative aux
deux homomorphismes $A\xrightarrow{q\circ u}E$ et
$p:E\to C=E/L$; comme $p$ est surjectif et que $E$ est une extension
$A$-triviale, cette suite est exacte et scindée en vertu de (20.5.12,
(iii)). La commutativité du carré de droite de (20.6.9.1) résulte de la
relation $v=p\circ q$ (20.5.2.7); l'image par
$q_{E/B/A}\otimes1_C$ du noyau $\Upsilon_{C/B/A}$ de $v_{C/B/A}$ est
donc contenue dans le noyau $L$ de $p_{C/E/A}$. D'autre part soit
$h:C\to E$ un $A$-homomorphisme inverse à droite de $p$, et soit
$j:L\to E$ l'injection canonique, de sorte que l'on a
$q(b)=h(v(b))+j(D(b))$ pour $b\in B$, où $D$ est la $A$-dérivation de
$B$ dans $L$ définissant la $B$-extension $E$; on peut écrire
$D=f\circ d_{B/A}$, où $f:\Omega^1_{B/A}\to L$ est un
$B$-homomorphisme. En vertu de (20.5.2.6), on a
\[
  q_{E/B/A}(d_{B/A}(b))=d_{E/A}(q(b))
  =d_{E/A}(h(v(b)))+d_{E/A}(j(f(d_{B/A}(b))))
\]
pour $b\in B$. Soit alors
$z=\sum_i d_{B/A}(b_i)\otimes c_i$, où $b_i\in B$ et $c_i\in C$, un
élément de $\Omega^1_{B/A}\otimes_B C$; on a
\[
  (q_{E/B/A}\otimes1_C)(z)=
  \sum_i d_{E/A}(h(v(b_i)))\otimes c_i+
  \sum_i d_{E/A}(j(f(d_{B/A}(b_i))))\otimes c_i.
\]
Dans la première somme, on a
$d_{E/A}(h(v(b_i)))=h_{E/C/A}(d_{C/A}(v(b_i)))$, donc cette somme est
$(h_{E/C/A}\otimes1_C)(v_{C/B/A}(z))$ en vertu de (20.5.2.6). Si l'on
prend $z\in\Upsilon_{C/B/A}$, cette somme est donc nulle, et il reste,
par définition de $\delta_{C/E/A}$,
\[
  (q_{E/B/A}\otimes1_C)(z)=
  \sum_i\delta_{C/E/A}(c_i f(d_{B/A}(b_i)))
  =\delta_{C/E/A}(f(z)),
\]
ce qui démontre la commutativité du carré de gauche dans (20.6.9.1).
\end{proof}

On notera que lorsque l'on suppose seulement que $\delta_{C/E/A}$ est
\emph{injectif} (et non nécessairement inversible à gauche), cette
interprétation permettrait encore de définir $\chi_E$ comme la
restriction de $q_{E/B/A}\otimes1_C$ à $\Upsilon_{C/B/A}$.

\begin{corollary}[20.6.10]
\label{0.20.6.10-fr}
Si $\mathfrak{R}$ est un idéal de $B$, $C=B/\mathfrak{R}$,
$E=B/\mathfrak{R}^2$, et si $E$ est une $B$-extension $A$-triviale
(18.3.7) de $C$ par $\mathfrak{R}/\mathfrak{R}^2$, alors
l'homomorphisme caractéristique
$\chi_E:\Upsilon_{C/B/A}\to\mathfrak{R}/\mathfrak{R}^2$ est bijectif.
\end{corollary}

\begin{proof}
\oldpage[0\textsubscript{IV-1}]{140}
En effet, dans le diagramme (20.6.9.1), les deux flèches verticales de
droite sont des homomorphismes bijectifs (20.5.12, (ii)).
\end{proof}

\begin{theorem}[20.6.11]
\label{0.20.6.11-fr}
Soient $u:A\to B$, $v:B\to C$ deux homomorphismes d'anneaux, $L$ un
$C$-module. Supposons vérifiée l'une des conditions suivantes :
\begin{enumerate}
  \item[\rm(i)] $L$ est un $C$-module injectif.
  \item[\rm(ii)] $\Upsilon_{C/B/A}$ est facteur direct du $C$-module
$\Omega^1_{B/A}\otimes_B C$ et
$u_{C/B/A}:\Omega^1_{C/A}\to\Omega^1_{C/B}$ est inversible à droite.
\end{enumerate}
Alors l'homomorphisme canonique (20.6.7.1)
\[
  \operatorname{Exalcom}_{B/A}(C,L)\longrightarrow
  \operatorname{Hom}_C(\Upsilon_{C/B/A},L)
\]
est bijectif.

En particulier, si $\Omega^1_{C/B}$ et $\Omega^1_{C/A}$ sont des
$C$-modules projectifs, l'homomorphisme canonique (20.6.7.1) est
bijectif.
\end{theorem}

\begin{proof}
Le fait que chacune des conditions (i), (ii) entraîne que (20.6.7.1) est
bijectif résulte dans les deux cas de la définition de $\alpha_1$. On
notera d'ailleurs que la condition (ii) est \emph{nécessaire et
suffisante} pour que l'homomorphisme (20.6.7.1) soit bijectif pour
\emph{tout} $C$-module $L$ (Bourbaki, \emph{Alg.}, chap. II,
3\textsuperscript{e} éd., \S~2, n\textsuperscript{o}~1, prop.~1). Si on
suppose que $\Omega^1_{C/B}$ et $\Omega^1_{C/A}$ sont des $C$-modules
projectifs, alors, dans la suite exacte (20.6.1.1),
$\operatorname{Ker}(u_{C/B/A})$ est un $C$-module projectif, car la suite exacte
\[
  0\longrightarrow\operatorname{Ker}(u_{C/B/A})\longrightarrow\Omega^1_{C/A}
  \longrightarrow\Omega^1_{C/B}\longrightarrow0
\]
est scindée, $\Omega^1_{C/B}$ étant projectif; comme
$\operatorname{Ker}(u_{C/B/A})=\operatorname{Im}(v_{C/B/A})$, la suite exacte
\[
  0\longrightarrow\Upsilon_{C/B/A}\longrightarrow
  \Omega^1_{B/A}\otimes_B C\longrightarrow
  \operatorname{Im}(v_{C/B/A})\longrightarrow0
\]
est scindée.
\end{proof}

\begin{corollary}[20.6.12]
\label{0.20.6.12-fr}
Supposons que $C$ soit une $A$-algèbre formellement lisse. Il existe alors
un homomorphisme canonique
\begin{equation}
\tag{20.6.12.1}
\label{0.20.6.12.1-fr}
  \operatorname{Exalcom}_B(C,L)\longrightarrow
  \operatorname{Hom}_C(\Upsilon_{C/B/A},L).
\end{equation}
En outre, cet homomorphisme est bijectif si l'une des conditions (i),
(ii) de (20.6.11) est vérifiée.
\end{corollary}

\begin{proof}
En effet, il résulte de l'hypothèse sur $C$ que
$\operatorname{Exalcom}_B(C,L)=\operatorname{Exalcom}_{B/A}(C,L)$, et
l'homomorphisme (20.6.12.1) n'est autre que (20.6.7.1).
\end{proof}

\begin{corollary}[20.6.13]
\label{0.20.6.13-fr}
Si $C$ est une $A$-algèbre formellement lisse et si $\Omega^1_{C/B}$ est
un $C$-module projectif, l'homomorphisme (20.6.12.1) est bijectif.
\end{corollary}

\begin{proof}
En effet, on sait alors que $\Omega^1_{C/A}$ est un $C$-module projectif
(20.4.9), les topologies étant discrètes.
\end{proof}

\begin{env}[20.6.14]
\label{0.20.6.14-fr}
Les notations restant les mêmes, supposons maintenant en outre que $A$,
$B$, $C$ soient des $\Lambda$-algèbres et $u$, $v$ des
$\Lambda$-homomorphismes, ce qui revient à se donner trois homomorphismes
d'anneaux
\[
  \Lambda\xrightarrow{s}A\xrightarrow{u}B\xrightarrow{v}C.
\]

\oldpage[0\textsubscript{IV-1}]{141}
On a donc, en dehors du module d'imperfection $\Upsilon_{C/B/A}$, les
modules d'imperfection $\Upsilon_{B/A/\Lambda}$,
$\Upsilon_{C/A/\Lambda}$ et $\Upsilon_{C/B/\Lambda}$, et on a déjà défini
des homomorphismes canoniques de $C$-modules (20.6.4.2)
\begin{align}
\tag{20.6.14.1}
\label{0.20.6.14.1-fr}
  u':\Upsilon_{C/A/\Lambda}&\longrightarrow\Upsilon_{C/B/\Lambda},\\
\tag{20.6.14.2}
\label{0.20.6.14.2-fr}
  s':\Upsilon_{C/B/\Lambda}&\longrightarrow\Upsilon_{C/B/A}.
\end{align}

Comme dans le diagramme commutatif (20.5.7.3)
\begin{equation}
\tag{20.6.14.3}
\label{0.20.6.14.3-fr}
\xymatrix{
  \Omega^1_{A/\Lambda}\otimes_A B \ar[r] \ar[d] &
  \Omega^1_{B/\Lambda} \ar[r] \ar[d] &
  \Omega^1_{B/A} \ar[r] \ar[d] & 0 \\
  \Omega^1_{A/\Lambda}\otimes_A C \ar[r] &
  \Omega^1_{C/\Lambda} \ar[r] & \Omega^1_{C/A} \ar[r] & 0
}
\end{equation}
la ligne inférieure est formée de $C$-modules, on en déduit par
tensorisation un diagramme commutatif
\begin{equation}
\tag{20.6.14.4}
\label{0.20.6.14.4-fr}
\xymatrix{
  \Omega^1_{A/\Lambda}\otimes_A C \ar[r] \ar@{=}[d] &
  \Omega^1_{B/\Lambda}\otimes_B C \ar[r] \ar[d] &
  \Omega^1_{B/A}\otimes_B C \ar[r] \ar[d] & 0 \\
  \Omega^1_{A/\Lambda}\otimes_A C \ar[r] &
  \Omega^1_{C/\Lambda} \ar[r] & \Omega^1_{C/A} \ar[r] & 0
}
\end{equation}
où la première ligne est encore exacte et la flèche verticale de gauche
est l'identité; si l'on pose
\begin{equation}
\tag{20.6.14.5}
\label{0.20.6.14.5-fr}
  \Upsilon^C_{B/A/\Lambda}=
  \operatorname{Ker}(u_{B/A/\Lambda}\otimes1_C)=
  \operatorname{Ker}(\Omega^1_{A/\Lambda}\otimes_A C\longrightarrow
       \Omega^1_{B/\Lambda}\otimes_B C)
\end{equation}
on voit, compte tenu de la définition de $\Upsilon_{C/A/\Lambda}$, que
l'on a un unique $C$-homomorphisme
\begin{equation}
\tag{20.6.14.6}
\label{0.20.6.14.6-fr}
  v':\Upsilon^C_{B/A/\Lambda}\longrightarrow\Upsilon_{C/A/\Lambda}
\end{equation}
rendant commutatif le diagramme
\begin{equation}
\tag{20.6.14.7}
\label{0.20.6.14.7-fr}
\xymatrix{
  0 \ar[r] & \Upsilon^C_{B/A/\Lambda} \ar[r] \ar[d]^{v'} &
  \Omega^1_{A/\Lambda}\otimes_A C \ar[r] \ar@{=}[d] &
  \Omega^1_{B/\Lambda}\otimes_B C \ar[r] \ar[d] &
  \Omega^1_{B/A}\otimes_B C \ar[r] \ar[d] & 0 \\
  0 \ar[r] & \Upsilon_{C/A/\Lambda} \ar[r] &
  \Omega^1_{A/\Lambda}\otimes_A C \ar[r] &
  \Omega^1_{C/\Lambda} \ar[r] & \Omega^1_{C/A} \ar[r] & 0
}
\end{equation}
dont les lignes sont exactes.

Lorsque $\Lambda=\mathbb{Z}$, on écrira $\Upsilon^C_{B/A}$ au lieu de
$\Upsilon^C_{B/A/\mathbb{Z}}$. Si $\Lambda$ est un corps premier, on a
$\Upsilon^C_{B/A/\Lambda}=\Upsilon^C_{B/A}$.
\end{env}

\begin{env}[20.6.15]
\label{0.20.6.15-fr}
Pour étudier les relations entre les modules précédents, nous introduirons
d'une part le complexe de $B$-modules
$\mathrm{K}_\bullet(B/A/\Lambda)$, d'autre part, les complexes
\oldpage[0\textsubscript{IV-1}]{142}
de $C$-modules $\mathrm{K}_\bullet(C/A/\Lambda)$ et
$\mathrm{K}_\bullet(C/B/\Lambda)$ (20.6.5), et en outre les complexes de
$C$-modules suivants. On posera tout d'abord
\begin{equation}
\tag{20.6.15.1}
\label{0.20.6.15.1-fr}
  \mathrm{K}^C_\bullet(B/A/\Lambda)=
  \mathrm{K}_\bullet(B/A/\Lambda)\otimes_B C.
\end{equation}
D'autre part, nous désignerons par $\mathrm{T}_\bullet(C/B/\Lambda)$ le
complexe de $C$-modules dont les termes sont nuls sauf en degrés 0 et 1,
et où
\begin{equation}
\tag{20.6.15.2}
\label{0.20.6.15.2-fr}
  \mathrm{T}_0(C/B/\Lambda)=\mathrm{T}_1(C/B/\Lambda)=
  \Omega^1_{B/\Lambda}\otimes_B C
\end{equation}
l'opérateur différentiel étant l'identité, de sorte que ce complexe est
homotope à 0; posons enfin
\begin{equation}
\tag{20.6.15.3}
\label{0.20.6.15.3-fr}
  \mathrm{K}'_\bullet(C/A/\Lambda)=
  \mathrm{K}_\bullet(C/A/\Lambda)\oplus
  \mathrm{T}_\bullet(C/B/\Lambda).
\end{equation}
En vertu du caractère trivial de $\mathrm{T}_\bullet$, il est clair que
l'on a
\begin{equation}
\tag{20.6.15.4}
\label{0.20.6.15.4-fr}
  \mathrm{H}^i(\mathrm{K}'_\bullet,L)=
  \mathrm{H}^i(\mathrm{K}_\bullet,L)
  \quad\text{et}\quad
  \mathrm{H}_i(\mathrm{K}'_\bullet,L)=
  \mathrm{H}_i(\mathrm{K}_\bullet,L)
\end{equation}
pour tout $C$-module $L$ et tout $i$.
\end{env}
\begin{env}[20.6.16]
\label{0.20.6.16-fr}
Définissons maintenant une suite exacte de complexes, \emph{scindée} en chaque
degré
\begin{equation}
\tag{20.6.16.1}
\label{0.20.6.16.1-fr}
  0\to\mathrm{K}^C_\bullet(B/A/\Lambda)
  \xrightarrow{j}\mathrm{K}'_\bullet(C/A/\Lambda)
  \xrightarrow{p}\mathrm{K}_\bullet(C/B/\Lambda)\to0
\end{equation}
de la façon suivante : désignons pour un moment par
\[
  f:\Omega^1_{A/\Lambda}\otimes_A C\to
    \Omega^1_{B/\Lambda}\otimes_B C,
  \qquad
  g:\Omega^1_{B/\Lambda}\otimes_B C\to\Omega^1_{C/\Lambda}
\]
les homomorphismes canoniques $u_{B/A/\Lambda}\otimes1_C$ et
$v_{C/B/\Lambda}$ respectivement, dont le composé est
$g\circ f=(v\circ u)_{C/A/\Lambda}$ (cf.\ \sref{0.20.6.14.4-fr}).
On prendra $j_1(x)=(x,f(x))$, $p_1(y,z)=z-f(y)$,
$j_0(x)=(g(x),x)$, $p_0(y,z)=g(z)-y$ de sorte que
$\operatorname{Im}(j_1)=\operatorname{Ker}(p_1)$ est le graphe de $f$,
supplémentaire de $\{0\}\oplus\mathrm{T}_1$, et
$\operatorname{Im}(j_0)=\operatorname{Ker}(p_0)$ est le graphe de $g$,
supplémentaire de $\mathrm{K}_1(C/A/\Lambda)\oplus\{0\}$ ; la vérification
de la commutativité du diagramme
\[
\xymatrix{
  0 \ar[r] & \mathrm{K}^C_1(B/A/\Lambda) \ar[r]^{j_1} \ar[d] &
  \mathrm{K}'_1(C/A/\Lambda) \ar[r]^{p_1} \ar[d] &
  \mathrm{K}_1(C/B/\Lambda) \ar[r] \ar[d] & 0 \\
  0 \ar[r] & \mathrm{K}^C_0(B/A/\Lambda) \ar[r]_{j_0} &
  \mathrm{K}'_0(C/A/\Lambda) \ar[r]_{p_0} &
  \mathrm{K}_0(C/B/\Lambda) \ar[r] & 0
}
\]
où les flèches verticales sont les opérateurs différentiels, est immédiate.
\end{env}

\begin{theorem}[20.6.17]
\label{0.20.6.17-fr}
On a une suite exacte de $C$-modules
\begin{equation}
\tag{20.6.17.1}
\label{0.20.6.17.1-fr}
  0\to\Upsilon^C_{B/A/\Lambda}\xrightarrow{v'}
  \Upsilon_{C/A/\Lambda}\xrightarrow{u'}
  \Upsilon_{C/B/\Lambda}\xrightarrow{\partial}
  \Omega^1_{B/A}\otimes_B C\xrightarrow{v_{C/B/A}}
  \Omega^1_{C/A}\xrightarrow{u_{C/B/A}}\Omega^1_{C/B}\to0
\end{equation}
\oldpage[0\textsubscript{IV-1}]{143}
où l'opérateur bord $\partial$ est le composé
\begin{equation}
\tag{20.6.17.2}
\label{0.20.6.17.2-fr}
  \Upsilon_{C/B/\Lambda}\xrightarrow{s'}\Upsilon_{C/B/A}
  \longrightarrow\Omega^1_{B/A}\otimes_B C
\end{equation}
la seconde flèche étant l'injection canonique.
\end{theorem}

\begin{proof}
En écrivant la suite exacte d'homologie pour la suite exacte de complexes
\sref{0.20.6.16.1-fr}, on obtient \sref{0.20.6.17.1-fr}, l'homologie étant
nulle en degrés distincts de 0 et de 1 ; le fait que les homomorphismes de
cette suite exacte qui proviennent par fonctorialité de $j$ et de $p$ sont
bien ceux de l'énoncé est immédiat. Reste à vérifier que $\partial$ est égal
à \sref{0.20.6.17.2-fr} ; or un élément $z\in\Upsilon_{C/B/\Lambda}$ est
image par $p_1$ de $(0,z)$, d'où on déduit aussitôt que $\partial(z)$ est
l'image de $z$ par l'homomorphisme canonique
\[
  s_{B/A/\Lambda}\otimes1_C:
  \Omega^1_{B/\Lambda}\otimes_B C\to\Omega^1_{B/A}\otimes_B C.
\]
Notre assertion résulte de la commutativité du diagramme
\begin{equation}
\tag{20.6.17.3}
\label{0.20.6.17.3-fr}
\xymatrix{
  \Upsilon_{C/B/\Lambda} \ar[r]^{s'} \ar[dd] & \Upsilon_{C/B/A} \ar[dd] \\
  & \\
  \Omega^1_{B/\Lambda}\otimes_B C \ar[r] &
  \Omega^1_{B/A}\otimes_B C
}
\end{equation}
(cf.\ \sref{0.20.6.4.3-fr}).
\end{proof}

\begin{corollary}[20.6.18]
\label{0.20.6.18-fr}
\begin{enumerate}
  \item[\rm(i)] La suite de $C$-modules
\begin{equation}
\tag{20.6.18.1}
\label{0.20.6.18.1-fr}
  0\to\Upsilon^C_{B/A/\Lambda}\xrightarrow{v'}
  \Upsilon_{C/A/\Lambda}\xrightarrow{u'}
  \Upsilon_{C/B/\Lambda}\xrightarrow{s'}\Upsilon_{C/B/A}\to0
\end{equation}
est exacte.
  \item[\rm(ii)] Si $B$ est une $A$-algèbre formellement lisse relativement
  à $\Lambda$, on a $\Upsilon^C_{B/A/\Lambda}=0$.
\end{enumerate}
\end{corollary}

\begin{proof}
(i) Dans \sref{0.20.6.17.1-fr}, l'image de $\Upsilon_{C/B/\Lambda}$ par
$\partial$ est le noyau de $v_{C/B/A}$, donc $\Upsilon_{C/B/A}$ par
définition.

(ii) L'hypothèse entraîne que la suite
\[
  0\to\Omega^1_{A/\Lambda}\otimes_A B\to
  \Omega^1_{B/\Lambda}\to\Omega^1_{B/A}\to0
\]
est exacte et scindée \sref{0.20.5.7-fr} ; par tensorisation avec $C$, la
suite
\[
  0\to\Omega^1_{A/\Lambda}\otimes_A C\to
  \Omega^1_{B/\Lambda}\otimes_B C\to
  \Omega^1_{B/A}\otimes_B C\to0
\]
reste donc exacte, d'où notre assertion.
\end{proof}

\begin{corollary}[20.6.19]
\label{0.20.6.19-fr}
Soient $K$ un corps, $E$, $F$ deux extensions de $K$ telles que
$K\subset E\subset F$.
\begin{enumerate}
  \item[\rm(i)] Si $F$ est une extension séparable de $E$, on a
  $\Upsilon_{F/E/K}=0$ (autrement dit, l'homomorphisme canonique
  $\Omega^1_{E/K}\otimes_E F\to\Omega^1_{F/K}$ est injectif).
  \item[\rm(ii)] Inversement, si $F$ est une extension séparable de $K$ et
  si l'on a $\Upsilon_{F/E/K}=0$, alors $F$ est une extension séparable de
  $E$.
\end{enumerate}
\end{corollary}

\begin{proof}
Si $P$ est le corps premier de $K$, on a en effet la suite exacte
\sref{0.20.6.18-fr}
\[
  \Upsilon_{F/K/P}\to\Upsilon_{F/E/P}\to\Upsilon_{F/E/K}\to0.
\]
Si $\Upsilon_{F/E/P}=\Upsilon_{F/E}=0$, on a donc
$\Upsilon_{F/E/K}=0$, d'où (i) en vertu de \sref{0.20.6.3-fr} ;
inversement, si $\Upsilon_{F/E/K}=0$ et
$\Upsilon_{F/K/P}=\Upsilon_{F/K}=0$, on a
$\Upsilon_{F/E/P}=0$, d'où (ii) en vertu de \sref{0.20.6.3-fr}.
\end{proof}

\oldpage[0\textsubscript{IV-1}]{144}
\begin{corollary}[20.6.20]
\label{0.20.6.20-fr}
\begin{enumerate}
  \item[\rm(i)] Si $K$ est une extension algébrique séparable de $k$, on a
  $\Omega^1_{K/k}=0$.
  \item[\rm(ii)] Soient $k$ un corps de caractéristique 0, $K$ une extension
  de $k$. Pour que $\Omega^1_{K/k}=0$, il faut et il suffit que $K$ soit une
  extension algébrique de $k$. En particulier, pour qu'un corps $K$ de
  caractéristique 0 soit tel que $\Omega^1_K=0$, il faut et il suffit que
  $K$ soit une extension algébrique de $\mathbb{Q}$ (cf.\ (21.4.4) et
  (21.7.5)).
\end{enumerate}
\end{corollary}

\begin{proof}
(i) Pour tout $x\in K$, soit $f$ le polynôme minimal de $x$ sur $k$ ; comme
$f'(x)\ne0$ et $d_{K/k}(f(x))=f'(x)d_{K/k}(x)=0$, on a
$d_{K/k}(x)=0$ et notre assertion résulte de \sref{0.20.4.7-fr}.

(ii) Il existe une extension pure $L$ de $k$ telle que
$k\subset L\subset K$ et que $K$ soit une extension algébrique de $L$.
Comme $K$ est séparable sur $L$, il résulte de \sref{0.20.6.19-fr}, (i) que
la suite \sref{0.20.5.7.2-fr}
\[
  0\to\Omega^1_{L/k}\otimes_L K\to\Omega^1_{K/k}
  \to\Omega^1_{K/L}\to0
\]
est exacte et de (i) que $\Omega^1_{K/L}=0$. La relation
$\Omega^1_{K/k}=0$ est donc équivalente à $\Omega^1_{L/k}=0$, et comme
$L$ est une extension pure de $k$, il résulte de \sref{0.20.5.10-fr} que
la relation $\Omega^1_{L/k}=0$ équivaut à $L=k$.
\end{proof}

\begin{remarks}[20.6.21]
\label{0.20.6.21-fr}
\begin{enumerate}
  \item[\rm(i)] Comme $\Upsilon_{B/A/\Lambda}$ est le noyau de
  \[
    u_{B/A/\Lambda}:\Omega^1_{A/\Lambda}\otimes_A B
    \to\Omega^1_{B/\Lambda}
  \]
  et $\Upsilon^C_{B/A/\Lambda}$ le noyau de
  $u_{B/A/\Lambda}\otimes1_C$, on a un homomorphisme canonique
\begin{equation}
\tag{20.6.21.1}
\label{0.20.6.21.1-fr}
  \Upsilon_{B/A/\Lambda}\otimes_B C\to\Upsilon^C_{B/A/\Lambda}.
\end{equation}
Cet homomorphisme est bijectif lorsque la suite
\begin{equation}
\tag{20.6.21.2}
\label{0.20.6.21.2-fr}
  0\to\Upsilon_{B/A/\Lambda}\otimes_B C\to
  \Omega^1_{A/\Lambda}\otimes_A C\to
  \Omega^1_{B/\Lambda}\otimes_B C\to
  \Omega^1_{B/A}\otimes_B C\to0
\end{equation}
est exacte, ce qui se produit dans les cas suivants :
\begin{enumerate}
  \item[$1^\circ$] $C$ est un $B$-module \emph{plat}.
  \item[$2^\circ$] Les $B$-modules $\Omega^1_{B/\Lambda}$ et
  $\Omega^1_{B/A}$ sont \emph{plats} : en effet, il en est alors de même de
  $\operatorname{Ker}(\Omega^1_{B/\Lambda}\to\Omega^1_{B/A})$
  ($\mathbf{0_I}$, 6.1.2), et la suite \sref{0.20.6.21.2-fr} est alors
  exacte en vertu de ($\mathbf{0_I}$, 6.1.2).
\end{enumerate}
  \item[\rm(ii)] Considérons un diagramme commutatif d'homomorphismes
  d'anneaux
\[
\xymatrix{
  \Lambda' \ar[r] & A' \ar[r] & B' \ar[r] & C' \\
  \Lambda \ar[u] \ar[r] & A \ar[u] \ar[r] & B \ar[u] \ar[r] & C \ar[u]
}
\]
Alors les définitions de \sref{0.20.6.16-fr} montrent que l'on a un
diagramme commutatif de complexes (où les flèches verticales proviennent de
\sref{0.20.6.6.4-fr})
\[
\xymatrix{
  0 \ar[r] & \mathrm{K}^C_\bullet(B/A/\Lambda) \ar[r] \ar[d] &
  \mathrm{K}'_\bullet(C/A/\Lambda) \ar[r] \ar[d] &
  \mathrm{K}_\bullet(C/B/\Lambda) \ar[r] \ar[d] & 0 \\
  0 \ar[r] & \mathrm{K}^{C'}_\bullet(B'/A'/\Lambda') \ar[r] &
  \mathrm{K}'_\bullet(C'/A'/\Lambda') \ar[r] &
  \mathrm{K}_\bullet(C'/B'/\Lambda') \ar[r] & 0
}
\]
\oldpage[0\textsubscript{IV-1}]{145}
d'où, par passage à l'homologie, un \emph{diagramme commutatif}
\begin{equation}
\tag{20.6.21.3}
\label{0.20.6.21.3-fr}
\xymatrix@C=0.7em{
  0 \ar[r] & \Upsilon^C_{B/A/\Lambda} \ar[r] \ar[d] &
  \Upsilon_{C/A/\Lambda} \ar[r] \ar[d] &
  \Upsilon_{C/B/\Lambda} \ar[r] \ar[d] &
  \Omega^1_{B/A}\otimes_B C \ar[r] \ar[d] &
  \Omega^1_{C/A} \ar[r] \ar[d] &
  \Omega^1_{C/B} \ar[r] \ar[d] & 0 \\
  0 \ar[r] & \Upsilon^{C'}_{B'/A'/\Lambda'} \ar[r] &
  \Upsilon_{C'/A'/\Lambda'} \ar[r] &
  \Upsilon_{C'/B'/\Lambda'} \ar[r] &
  \Omega^1_{B'/A'}\otimes_{B'} C' \ar[r] &
  \Omega^1_{C'/A'} \ar[r] & \Omega^1_{C'/B'} \ar[r] & 0
}
\end{equation}
On a un diagramme commutatif analogue pour \sref{0.20.6.18.1-fr}.
\end{enumerate}
\end{remarks}

\begin{proposition}[20.6.22]
\label{0.20.6.22-fr}
Soient $s:\Lambda\to A$, $u:A\to B$ deux homomorphismes d'anneaux,
$\mathfrak{R}$ un idéal de $B$, $C$ l'anneau quotient $B/\mathfrak{R}$.
Supposons que $E=B/\mathfrak{R}^2$ soit une $B$-extension
$\Lambda$-triviale de $C$ par $\mathfrak{R}/\mathfrak{R}^2$. On a alors une
suite exacte
\begin{equation}
\tag{20.6.22.1}
\label{0.20.6.22.1-fr}
  0\to\Upsilon^C_{B/A/\Lambda}\xrightarrow{v'}
  \Upsilon_{C/A/\Lambda}\xrightarrow{\chi_E}
  \mathfrak{R}/\mathfrak{R}^2\xrightarrow{\delta_{C/B/A}}
  \Omega^1_{B/A}\otimes_B C\xrightarrow{v_{C/B/A}}
  \Omega^1_{C/A}\to0.
\end{equation}
\end{proposition}

\begin{proof}
En effet, comme $v:B\to C$ est surjectif, on a $\Omega^1_{C/B}=0$
\sref{0.20.4.12-fr}. En outre, il résulte de \sref{0.20.6.10-fr} que
$\Upsilon_{C/B/\Lambda}$ s'identifie canoniquement à
$\mathfrak{R}/\mathfrak{R}^2$. Il suffit alors d'appliquer les suites exactes
\sref{0.20.6.17.1-fr} et \sref{0.20.6.18.1-fr}, en notant que l'on a un
diagramme commutatif
\[
\xymatrix{
  \Upsilon_{C/A/\Lambda} \ar[r] \ar[d]_{\chi_E} &
  \Omega^1_{A/\Lambda}\otimes_A C \ar[d] \\
  \Upsilon_{C/E/\Lambda}=\Upsilon_{C/B/\Lambda}=\mathfrak{R}/\mathfrak{R}^2
    \ar[r]^-{\delta_{C/B/A}} \ar[d]_{\ell} &
  \Omega^1_{B/A}\otimes_B C=\Omega^1_{E/A}\otimes_E C \ar[d] \\
  \mathfrak{R}/\mathfrak{R}^2 \ar[r]^{\delta_{C/B/A}} &
  \Omega^1_{B/A}\otimes_B C
}
\]
et utilisant la commutativité du diagramme \sref{0.20.6.17.3-fr}.
\end{proof}

On a ainsi précisé le \emph{noyau} de $\delta_{C/B/A}$ dans le cas où il
existe un anneau $\Lambda$ tel que $A$ soit une $\Lambda$-algèbre et que
$B/\mathfrak{R}^2$ soit $\Lambda$-triviale ; on observera que c'est le cas
en particulier lorsque $C$ est une \emph{$\Lambda$-algèbre formellement
lisse} (pour la topologie discrète).

Supposons en outre qu'on ait un diagramme commutatif d'homomorphismes
d'anneaux
\[
\xymatrix{
  \Lambda' \ar[r] & A' \ar[r] & B' \\
  \Lambda \ar[u] \ar[r] & A \ar[u] \ar[r] & B \ar[u]^f
}
\]
\oldpage[0\textsubscript{IV-1}]{146}
que $\mathfrak{R}'$ soit un idéal de $B'$ tel que
$f(\mathfrak{R})\subset\mathfrak{R}'$, et que
$E'=B'/\mathfrak{R}'^2$ soit une $B'$-extension $\Lambda'$-triviale de
$C'=B'/\mathfrak{R}'$ par $\mathfrak{R}'/\mathfrak{R}'^2$. On a alors un
diagramme commutatif
\begin{equation}
\tag{20.6.22.2}
\label{0.20.6.22.2-fr}
\xymatrix@C=0.7em{
  0 \ar[r] & \Upsilon^{C'}_{B'/A'/\Lambda'} \ar[r] &
  \Upsilon_{C'/A'/\Lambda'} \ar[r] &
  \mathfrak{R}'/\mathfrak{R}'^2 \ar[r] &
  \Omega^1_{B'/A'}\otimes_{B'} C' \ar[r] &
  \Omega^1_{C'/A'} \ar[r] & 0 \\
  0 \ar[r] & \Upsilon^C_{B/A/\Lambda} \ar[r] \ar[u] &
  \Upsilon_{C/A/\Lambda} \ar[r] \ar[u] &
  \mathfrak{R}/\mathfrak{R}^2 \ar[r] \ar[u] &
  \Omega^1_{B/A}\otimes_B C \ar[r] \ar[u] &
  \Omega^1_{C/A} \ar[r] \ar[u] & 0
}
\end{equation}
comme il résulte de \sref{0.20.6.21.3-fr} et \sref{0.20.5.11.3-fr}.

\begin{corollary}[20.6.23]
\label{0.20.6.23-fr}
Sous les hypothèses de \sref{0.20.6.22-fr}, supposons en outre que $B$ soit
une $A$-algèbre formellement lisse. Alors on a une suite exacte
\begin{equation}
\tag{20.6.23.1}
\label{0.20.6.23.1-fr}
  0\to\Upsilon_{C/A/\Lambda}\xrightarrow{\chi_E}
  \mathfrak{R}/\mathfrak{R}^2\xrightarrow{\delta_{C/B/A}}
  \Omega^1_{B/A}\otimes_B C\to\Omega^1_{C/A}\to0.
\end{equation}
\end{corollary}

\begin{proof}
Cela résulte en effet de \sref{0.20.6.18-fr}, (ii).
\end{proof}

\begin{env}[20.6.24]
\label{0.20.6.24-fr}
Lorsque les hypothèses de \sref{0.20.6.22-fr} sont satisfaites, on dit encore
que l'homomorphisme caractéristique $\chi_E$ est
l'\emph{homomorphisme caractéristique de la $A$-algèbre $B$, relativement à
$\Lambda$ et à l'idéal $\mathfrak{R}$} (en supprimant ces dernières
précisions lorsqu'il n'y a pas de confusion à craindre) ; on le notera
parfois $\chi_B$ ou $\chi_{B/A}$.
\end{env}

\begin{proposition}[20.6.25]
\label{0.20.6.25-fr}
Soient $s:\Lambda\to A$, $u:A\to B$, $v:B\to C$ trois homomorphismes
d'anneaux, $L$ un $C$-module. On a alors une suite exacte
\begin{equation}
\tag{20.6.25.1}
\label{0.20.6.25.1-fr}
\begin{split}
  0\to\operatorname{D\acute{e}r}_B(C,L)\xrightarrow{u^1}
  \operatorname{D\acute{e}r}_A(C,L)\xrightarrow{v^1}
  \operatorname{D\acute{e}r}_A(B,L)\xrightarrow{\partial} \\
  \to\operatorname{Exalcom}_{B/\Lambda}(C,L)\xrightarrow{u^1}
  \operatorname{Exalcom}_{A/\Lambda}(C,L)\xrightarrow{v^1}
  \operatorname{Exalcom}_{A/\Lambda}(B,L)\to0
\end{split}
\end{equation}
où $u^1$, $v^1$ sont les homomorphismes définis dans (18.3.6.4) et
(18.3.6.2), et $\partial$ est défini comme dans \sref{0.20.2.2-fr}.
\end{proposition}

\begin{proof}
En effet, comme la suite exacte \sref{0.20.6.16.1-fr} est scindée, on en
déduit une suite exacte
\[
  0\to\operatorname{Hom}_C(\mathrm{K}_\bullet(C/B/\Lambda),L)\to
  \operatorname{Hom}_C(\mathrm{K}'_\bullet(C/A/\Lambda),L)\to
  \operatorname{Hom}_C(\mathrm{K}^C_\bullet(B/A/\Lambda),L)\to0.
\]
Si on applique à ce complexe la suite exacte de cohomologie, en tenant compte
de \sref{0.20.6.15.4-fr}, et de \sref{0.20.6.6-fr}, on obtient
\sref{0.20.6.25.1-fr} car on a
\[
  \operatorname{Hom}_C(\mathrm{K}^C_\bullet(B/A/\Lambda),L)
  =\operatorname{Hom}_B(\mathrm{K}_\bullet(B/A/\Lambda),L)
\]
par définition ; l'identification de $u^1$ et $v^1$ avec les homomorphismes
de (18.3.4.2) résulte de \sref{0.20.6.6.4-fr}.
\end{proof}

\begin{remark}[20.6.26]
\label{0.20.6.26-fr}
Dans ce numéro, les complexes $\mathrm{K}_\bullet(C/B/A)$ sont apparus comme
un artifice technique destiné à simplifier l'exposition de certains
comportements fonctoriels. En réalité, ces complexes, considérés comme objets
de la catégorie des
\oldpage[0\textsubscript{IV-1}]{147}
complexes de $C$-modules « à homotopie près » (c'est-à-dire où les morphismes
sont les classes d'homomorphismes homotopes) sont des invariants remarquables,
plus fins que le couple formé de $\Omega^1_{C/B}$ et de
$\Upsilon_{C/B/A}$. Lorsque $k$ est un corps premier, et $C$ une $k$-algèbre
formellement lisse (par exemple un anneau régulier de type fini sur une
extension de $k$ (cf.\ ($\mathbf{IV}$, 6.8.6))), on peut montrer que le
complexe $\mathrm{K}_\bullet(C/A/k)$ peut se décrire uniquement en faisant
intervenir $C$ et $A$ (à l'exclusion de $k$) : on exprime $C$ comme quotient
d'une algèbre de polynômes $B$ sur $A$ par un idéal $\mathfrak{Q}$, et on
considère le complexe $\mathrm{F}_\bullet(C/A)$ à 2 termes non nuls
\[
  \cdots\to0\to\mathfrak{Q}/\mathfrak{Q}^2\to
  \Omega^1_{B/A}\otimes_B C\to0\to\cdots
\]
(dont l'homologie coïncide bien avec celle de
$\mathrm{K}_\bullet(C/A/k)$ en vertu de \sref{0.20.6.23.1-fr}). Ces
complexes $\mathrm{F}_\bullet(C/A)$, qui du point de vue de l'algèbre
homologique jouent le rôle d'un fibré conormal pour $\operatorname{Spec}(C)$
au-dessus de $\operatorname{Spec}(A)$, tiendront une place importante dans
les chapitres de cet ouvrage consacrés à la dualité des faisceaux cohérents et
au théorème de Riemann--Roch.
\end{remark}

\subsection{Généralisations aux anneaux topologiques}
\label{subsection:0.20.7-fr}

\begin{env}[20.7.1]
\label{0.20.7.1-fr}
Il résulte aussitôt des définitions que si, dans \sref{0.20.5.2-fr} et
\sref{0.20.5.3-fr}, les anneaux $A$, $B$, $C$ sont supposés être des anneaux
topologiques et les homomorphismes d'anneaux $u$, $v$ continus, alors les
homomorphismes $u_{C/B/A}$ et $v_{C/B/A}$ sont \emph{continus} (il faut
naturellement prendre sur $\Omega^1_{B/A}\otimes_B C$ la topologie produit
tensoriel).

En outre, $u_{C/B/A}:\Omega^1_{C/A}\to\Omega^1_{C/B}$ est un
\emph{morphisme strict surjectif} de $C$-modules topologiques ; en effet, il en
est ainsi de l'homomorphisme canonique $C\otimes_A C\to C\otimes_B C$, compte
tenu de la définition de la topologie produit tensoriel ; par passage aux
quotients on en déduit que l'homomorphisme correspondant
$\mathrm{P}^1_{C/A}\to\mathrm{P}^1_{C/B}$ est un morphisme strict surjectif,
et $u_{C/B/A}$ est la restriction de ce dernier à $\Omega^1_{C/A}$.

Dans \sref{0.20.5.5-fr}, si $A'$ et $B$ sont des $A$-algèbres topologiques,
et si $B'$ et $\Omega^1_{B/A}\otimes_B B'$ sont munis des topologies produits
tensoriels, l'homomorphisme canonique \sref{0.20.5.5.1-fr} est un isomorphisme
topologique, en tenant compte de ce que $\mathrm{P}^1_{B/A}$ est somme directe
topologique de $B$ et de $\Omega^1_{B/A}$ \sref{0.20.4.8-fr}.
\end{env}

\begin{proposition}[20.7.2]
\label{0.20.7.2-fr}
Soient $u:A\to B$, $v:B\to C$ deux homomorphismes continus d'anneaux
topologiques. Pour que l'homomorphisme continu
$v_{C/B/A}:\Omega^1_{B/A}\otimes_B C\to\Omega^1_{C/A}$ soit formellement
inversible à gauche (cf.\ \sref{0.19.1.5-fr}), il faut et il suffit que $C$
soit une $B$-algèbre formellement lisse relativement à $A$
\sref{0.19.9.1-fr} (et \emph{a fortiori} il suffit que $C$ soit une
$B$-algèbre formellement lisse).
\end{proposition}

\begin{proof}
Dire que $v_{C/B/A}$ est formellement inversible à gauche signifie en effet,
vu que les topologies de $\Omega^1_{C/A}$ et $\Omega^1_{B/A}\otimes_B C$
sont moins fines que celles déduites de la topologie de $C$
\sref{0.20.4.5-fr}, que pour tout $C$-module \emph{discret} $L$, annulé par un
idéal ouvert de $C$, l'homomorphisme canonique
\[
  \operatorname{Hom.\ cont}_C(\Omega^1_{C/A},L)\to
  \operatorname{Hom.\ cont}_C(\Omega^1_{B/A}\otimes_B C,L)
\]
\oldpage[0\textsubscript{IV-1}]{148}
est \emph{surjectif} \sref{0.19.1.5-fr} ; comme
$\operatorname{Hom.\ cont}_C(\Omega^1_{B/A}\otimes_B C,L)=
\operatorname{Hom.\ cont}_B(\Omega^1_{B/A},L)$ par définition de la topologie
produit tensoriel, il revient au même, en vertu de \sref{0.20.4.8-fr}, de dire
que l'homomorphisme canonique
\[
  \operatorname{D\acute{e}r.\ cont}_A(C,L)\to
  \operatorname{D\acute{e}r.\ cont}_A(B,L)
\]
est surjectif. Mais la suite exacte \sref{0.20.3.8.2-fr} montre que cette
condition équivaut à $\operatorname{Exalcotop}_{B/A}(C,L)=0$, c'est-à-dire
précisément au fait que $C$ est formellement lisse relativement à $A$
\sref{0.19.9.8-fr}.
\end{proof}

\begin{corollary}[20.7.3]
\label{0.20.7.3-fr}
Supposons que dans $B$ et dans $C$, le carré d'un idéal ouvert soit ouvert.
Pour que $C$ soit une $B$-algèbre formellement lisse relativement à $A$, il
faut et il suffit que si l'on désigne par $(\mathfrak{R}_\lambda)$ un système
fondamental de voisinages de $0$ formé d'idéaux de $C$, alors, pour tout
$\lambda$, l'homomorphisme
\begin{equation}
\tag{20.7.3.1}
\label{0.20.7.3.1-fr}
  v_{C/B/A}\otimes 1_{C/\mathfrak{R}_\lambda}:
  \Omega^1_{B/A}\otimes_B(C/\mathfrak{R}_\lambda)\to
  \Omega^1_{C/A}\otimes_C(C/\mathfrak{R}_\lambda)
\end{equation}
soit inversible à gauche.
\end{corollary}

\begin{proof}
On sait en effet alors que la topologie de $\Omega^1_{B/A}$ (resp.\
$\Omega^1_{C/A}$) est déduite de celle de $B$ (resp.\ de $C$)
\sref{0.20.4.5-fr} ; on en conclut aussitôt que la topologie de
$\Omega^1_{B/A}\otimes_B C$ est aussi déduite de celle de $C$ ; le corollaire
résulte alors de \sref{0.20.7.2-fr} et \sref{0.19.1.7-fr}.
\end{proof}

\begin{proposition}[20.7.4]
\label{0.20.7.4-fr}
Soient $A$ un anneau topologique, $B$ une $A$-algèbre topologique. Pour que
$B$ soit formellement non ramifiée \sref{0.19.10.2-fr}, il faut et il suffit
que le séparé complété $\widehat{\Omega}^1_{B/A}$ soit nul.
\end{proposition}

\begin{proof}
En effet, il résulte aussitôt de \sref{0.19.10.4-fr} et
\sref{0.20.1.1-fr} que, pour que $B$ soit formellement non ramifiée, il faut et
il suffit que pour tout idéal ouvert $\mathfrak{R}$ de $B$, tout idéal ouvert
$\mathfrak{H}$ de $A$ tel que $\mathfrak{H}B\subset\mathfrak{R}$ et tout
$(B/\mathfrak{R})$-module $L$, on ait
$\operatorname{D\acute{e}r}_{A/\mathfrak{H}}(B/\mathfrak{R},L)=0$,
c'est-à-dire $\operatorname{D\acute{e}r.\ cont}_A(B,L)=0$ pour tout
$B$-module discret $L$ annulé par un idéal ouvert de $B$ ; en vertu de
\sref{0.20.4.8-fr}, cela équivaut à
$\operatorname{Hom.\ cont}_B(\Omega^1_{B/A},L)=0$ pour un tel $B$-module,
d'où aussitôt la proposition.

Lorsque $B$ est \emph{discret}, la condition de l'énoncé de
\sref{0.20.7.4-fr} équivaut donc à $\Omega^1_{B/A}=0$.
\end{proof}

\begin{corollary}[20.7.5]
\label{0.20.7.5-fr}
Soient $A$ un anneau, $\mathfrak{m}$ un idéal de $A$, $B$ une $A$-algèbre ;
munissons $A$ de la topologie $\mathfrak{m}$-préadique, $B$ de la topologie
déduite de celle de $A$, et posons $A_0=A/\mathfrak{m}$,
$B_0=B/\mathfrak{m}B=B\otimes_A A_0$. Alors, pour que $B$ soit une
$A$-algèbre formellement non ramifiée, il faut et il suffit que
$\Omega^1_{B_0/A_0}=0$ (ou encore que $B_0$ soit une $A_0$-algèbre
formellement non ramifiée).
\end{corollary}

\begin{proof}
En effet, on a
$\Omega^1_{B/A}\otimes_B B_0=\Omega^1_{B/A}/\mathfrak{m}.\Omega^1_{B/A}
=\Omega^1_{B_0/A_0}$ en vertu de \sref{0.20.5.5-fr} ; écrire que tout
sous-module ouvert de $\Omega^1_{B/A}$ est égal à $\Omega^1_{B/A}$ équivaut
par \sref{0.20.4.5-fr} à écrire que
$\mathfrak{m}.\Omega^1_{B/A}=\Omega^1_{B/A}$, d'où la conclusion.
\end{proof}

\begin{proposition}[20.7.6]
\label{0.20.7.6-fr}
Soient $u:A\to B$, $v:B\to C$ deux homomorphismes continus d'anneaux
topologiques, et supposons que $v$ fasse de $C$ une $B$-algèbre formellement
étale ; alors $v_{C/B/A}:\Omega^1_{B/A}\otimes_B C\to\Omega^1_{C/A}$ est un
bimorphisme formel \sref{0.19.1.2-fr}.
\end{proposition}

\begin{proof}
En effet, pour tout $C$-module discret $L$ annulé par un idéal ouvert de $C$,
on a alors \sref{0.20.7.4-fr}
$\operatorname{D\acute{e}r.\ cont}_B(C,L)=0$ et par suite
\sref{0.20.3.6.1-fr} l'homomorphisme canonique
\oldpage[0\textsubscript{IV-1}]{149}
$\operatorname{D\acute{e}r.\ cont}_A(C,L)\to
\operatorname{D\acute{e}r.\ cont}_A(B,L)$ est injectif ; il est par ailleurs
surjectif en vertu de \sref{0.20.7.2-fr}, donc il est bijectif ; il en résulte
que l'image de $v_{C/B/A}$ est nécessairement dense dans $\Omega^1_{C/A}$,
sans quoi le quotient de $\Omega^1_{C/A}$ par l'adhérence de
$\operatorname{Im}(v_{C/B/A})$ serait séparé et $\neq0$ et aurait donc un
quotient discret $L\neq0$, contrairement à ce qu'on vient de voir (compte tenu
de \sref{0.20.4.8-fr}). Par suite $v_{C/B/A}$, qui est un monomorphisme formel
en vertu de \sref{0.20.7.2-fr}, est aussi un épimorphisme formel
\sref{0.19.1.2-fr}, donc un bimorphisme formel.
\end{proof}

\begin{corollary}[20.7.7]
\label{0.20.7.7-fr}
Supposons que dans $B$ et dans $C$ le carré de tout idéal ouvert soit ouvert.
Si $C$ est une $B$-algèbre formellement étale, alors, pour tout idéal ouvert
$\mathfrak{R}_\lambda$ de $C$, l'homomorphisme \sref{0.20.7.3.1-fr} est
bijectif.
\end{corollary}

\begin{proof}
En effet, en vertu de \sref{0.19.1.1-fr}, cet homomorphisme est surjectif, et
il est injectif par \sref{0.20.7.3-fr}.
\end{proof}

\begin{proposition}[20.7.8]
\label{0.20.7.8-fr}
Soient $u:A\to B$ un homomorphisme continu d'anneaux topologiques,
$\mathfrak{R}$ un idéal de $B$, $C$ l'anneau topologique quotient
$B/\mathfrak{R}$, $v:B\to C$ l'homomorphisme canonique. Alors :
\begin{enumerate}
  \item[\rm(i)] Dans la suite exacte \sref{0.20.5.12.1-fr}
\[
  \mathfrak{R}/\mathfrak{R}^2\xrightarrow{\delta_{C/B/A}}
  \Omega^1_{B/A}\otimes_B C\xrightarrow{v_{C/B/A}}\Omega^1_{C/A}\to0
\]
l'homomorphisme $\delta_{C/B/A}$ est continu et l'homomorphisme $v_{C/B/A}$
est un morphisme strict de $C$-modules topologiques.
  \item[\rm(ii)] Pour que $\delta_{C/B/A}$ soit formellement inversible à
gauche \sref{0.19.1.5-fr}, il faut et il suffit que pour tout $C$-module
discret $L$ annulé par un idéal ouvert de $C$, l'homomorphisme canonique
\begin{equation}
\tag{20.7.8.1}
\label{0.20.7.8.1-fr}
  \operatorname{Exalcotop}_A(C,L)\to\operatorname{Exalcotop}_A(B,L)
\end{equation}
soit injectif.
\end{enumerate}
\end{proposition}

\begin{proof}
(i) La première assertion est évidente. D'autre part, l'homomorphisme
canonique $v\otimes v:B\otimes_A B\to C\otimes_A C$ est un morphisme strict
par définition de la topologie produit tensoriel, et on en déduit aussitôt
(cf.\ \sref{0.20.5.2-fr}) qu'il en est de même de $v_{C/B/A}$.

(ii) Dire que $\delta_{C/B/A}$ est formellement inversible à gauche signifie
que pour tout $C$-module discret $L$ annulé par un idéal ouvert de $C$,
l'homomorphisme canonique
\[
  \operatorname{Hom.\ cont}_C(\Omega^1_{B/A}\otimes_B C,L)\to
  \operatorname{Hom.\ cont}_C(\mathfrak{R}/\mathfrak{R}^2,L)
\]
est surjectif. Or, compte tenu de \sref{0.18.4.3-fr} et de
\sref{0.20.4.8-fr}, cela revient à dire que l'homomorphisme canonique
\[
  \operatorname{D\acute{e}r.\ cont}_A(B,L)\to
  \operatorname{Exalcotop}_B(C,L)
\]
est surjectif, et la conclusion résulte donc de la suite exacte
\sref{0.20.3.6.1-fr}.

Le fait que \sref{0.20.7.8.1-fr} soit injectif s'exprime encore de la manière
suivante, compte tenu de la définition des deux membres
\sref{0.18.4.1-fr} : pour tout idéal ouvert $\mathfrak{M}$ de $A$, tout idéal
ouvert $\mathfrak{N}$ de $B$ tel que $\mathfrak{M}B\subset\mathfrak{N}$, et
toute $(A/\mathfrak{M})$-extension $E$ de
$B/(\mathfrak{R}+\mathfrak{N})$ par un
$(B/(\mathfrak{R}+\mathfrak{N}))$-module $L$, telle que l'image réciproque de
$E$ par l'homomorphisme canonique
$B/\mathfrak{N}\to B/(\mathfrak{R}+\mathfrak{N})$ soit
$(A/\mathfrak{M})$-triviale, il existe un idéal ouvert
$\mathfrak{M}'\subset\mathfrak{M}$ de $A$,
\oldpage[0\textsubscript{IV-1}]{150}
un idéal ouvert $\mathfrak{N}'\subset\mathfrak{N}$ de $B$ tels que
$\mathfrak{M}'B\subset\mathfrak{N}'$ et que l'image réciproque de $E$ par
l'homomorphisme canonique
$B/(\mathfrak{R}+\mathfrak{N}')\to B/(\mathfrak{R}+\mathfrak{N})$ soit
$(A/\mathfrak{M}')$-triviale.
\end{proof}

\begin{corollary}[20.7.9]
\label{0.20.7.9-fr}
Si la $A$-algèbre topologique $C=B/\mathfrak{R}$ est formellement lisse,
l'homomorphisme canonique $\delta_{C/B/A}$ est formellement inversible à
gauche.
\end{corollary}

\begin{env}[20.7.10]
\label{0.20.7.10-fr}
Dans \sref{0.20.6.1-fr}, lorsque $A$, $B$, $C$ sont des anneaux topologiques,
$u$ et $v$ des homomorphismes continus, on munit $\Upsilon_{C/B/A}$ de la
topologie induite par celle de $\Omega^1_{B/A}\otimes_B C$ ; les
homomorphismes \sref{0.20.6.4.2-fr} et \sref{0.20.6.4.4-fr} sont alors
continus pourvu qu'il en soit de même de ceux du diagramme
\sref{0.20.6.4.1-fr}. En outre, si, dans \sref{0.20.6.7-fr}, on suppose que
$L$ est un $C$-module discret annulé par un idéal ouvert de $C$, on déduit,
par passage à la limite inductive, un $C$-homomorphisme canonique
\begin{equation}
\tag{20.7.10.1}
\label{0.20.7.10.1-fr}
  \operatorname{Exalcotop}_{B/A}(C,L)\to
  \operatorname{Hom.\ cont}_C(\Upsilon_{C/B/A},L)
\end{equation}
compte tenu de \sref{0.18.5.3.1-fr} : pour tout idéal ouvert
$\mathfrak{M}$ de $A$, tout idéal ouvert $\mathfrak{N}$ de $B$ tel que
$\mathfrak{M}B\subset\mathfrak{N}$, tout idéal ouvert $\mathfrak{P}$ de $C$
tel que $\mathfrak{N}C\subset\mathfrak{P}$ et que $\mathfrak{P}.L=0$, toute
$(B/\mathfrak{N})$-extension $E$ de $C/\mathfrak{P}$ par $L$ qui est
$(A/\mathfrak{M})$-triviale provient de la donnée d'un $C$-homomorphisme
continu $f$ de $\Omega^1_{B/A}\otimes_B C$ dans $L$, et l'homomorphisme
\sref{0.20.7.10.1-fr} fait correspondre à l'image dans
$\operatorname{Exalcotop}_{B/A}(C,L)$ de la classe de $E$, la restriction de
$f$ à $\Upsilon_{C/B/A}$, dit \emph{homomorphisme caractéristique de $E$} et
noté encore $\chi_E$.
\end{env}

\begin{proposition}[20.7.11]
\label{0.20.7.11-fr}
Supposons que la topologie de $C$ soit telle que le carré d'un idéal ouvert
soit ouvert. Si $C$ est une $A$-algèbre topologique formellement lisse et si
$\Omega^1_{C/B}$ est un $C$-module formellement projectif, on a un
isomorphisme canonique
\begin{equation}
\tag{20.7.11.1}
\label{0.20.7.11.1-fr}
  \operatorname{Exalcotop}_B(C,L)\xrightarrow{\sim}
  \operatorname{Hom.\ cont}_C(\Upsilon_{C/B/A},L)
\end{equation}
pour tout $C$-module discret $L$ annulé par un idéal ouvert de $C$.
\end{proposition}

\begin{proof}
On a en effet, dans la suite exacte \sref{0.20.3.7.1-fr},
$\operatorname{Exalcotop}_A(C,L)=0$ \sref{0.19.4.4-fr}, donc
$\operatorname{Exalcotop}_B(C,L)=\operatorname{Exalcotop}_{B/A}(C,L)$ et
l'homomorphisme \sref{0.20.7.8.1-fr} n'est autre que
\sref{0.20.7.10.1-fr} ; le fait qu'il est bijectif se déduit de
\sref{0.20.6.13-fr} par passage à la limite inductive, compte tenu de ce que
la topologie de $\Omega^1_{C/B}$ est alors déduite de celle de $C$
\sref{0.20.4.5-fr} et de \sref{0.19.2.4-fr}.
\end{proof}

\begin{env}[20.7.12]
\label{0.20.7.12-fr}
Dans \sref{0.20.6.14-fr} on munit encore $\Upsilon^C_{B/A/\Lambda}$ de la
topologie induite par celle de $\Omega^1_{A/\Lambda}\otimes_A C$ et alors
l'homomorphisme \sref{0.20.6.14.6-fr} est continu, lorsque les anneaux
considérés sont topologiques et les homomorphismes d'anneaux continus.
\end{env}

\begin{env}[20.7.13]
\label{0.20.7.13-fr}
Si, dans \sref{0.20.6.23-fr}, on suppose que les anneaux $\Lambda$, $A$, $B$,
$C$ sont topologiques, les homomorphismes $s$, $u$, $v$ continus et $L$ un
$C$-module discret annulé par un idéal ouvert de $C$, alors on peut passer à
la limite inductive comme dans \sref{0.20.3.6-fr}, et on a une suite exacte
\begin{equation}
\tag{20.7.13.1}
\label{0.20.7.13.1-fr}
\begin{split}
  0\to\operatorname{D\acute{e}r.\ cont}_B(C,L)\to
  \operatorname{D\acute{e}r.\ cont}_A(C,L)\to
  \operatorname{D\acute{e}r.\ cont}_A(B,L)\to \\
  \to\operatorname{Exalcotop}_{B/\Lambda}(C,L)\to
  \operatorname{Exalcotop}_{A/\Lambda}(C,L)\to
  \operatorname{Exalcotop}_{A/\Lambda}(B,L)\to0.
\end{split}
\end{equation}
\end{env}

\begin{env}[20.7.14]
\label{0.20.7.14-fr}
Soient $A$ un anneau topologique, $B$ une $A$-algèbre topologique,
$\mathfrak{M}'$ un idéal ouvert de $A$, $\mathfrak{N}'$ un idéal ouvert de
$B$ tels que $\mathfrak{M}'B\subset\mathfrak{N}'$ ; posons
$A'=A/\mathfrak{M}'$,
\oldpage[0\textsubscript{IV-1}]{151}
$B'=B/\mathfrak{N}'$ ; le noyau de l'homomorphisme
$B\otimes_A B\to B'\otimes_{A'}B'$ est
$\mathfrak{U}'=\operatorname{Im}(\mathfrak{N}'\otimes B+B\otimes\mathfrak{N}')$,
d'où résulte aussitôt que le noyau de l'homomorphisme
\begin{equation}
\tag{20.7.14.1}
\label{0.20.7.14.1-fr}
  \varphi_{(\mathfrak{M}',\mathfrak{N}')}:
  \Omega^1_{B/A}\to\Omega^1_{B'/A'}
\end{equation}
est $((\mathfrak{J}\cap\mathfrak{U}')+\mathfrak{J}^2)/\mathfrak{J}^2$ ; par
ailleurs, l'homomorphisme \sref{0.20.7.14.1-fr} est surjectif, comme il
résulte de \sref{0.20.4.7-fr}. Si $\mathfrak{M}''$ (resp.\
$\mathfrak{N}''$) est un second idéal ouvert de $A$ (resp.\ $B$) tel que
$\mathfrak{M}''\subset\mathfrak{M}'$, $\mathfrak{N}''\subset\mathfrak{N}'$
et $\mathfrak{M}''B\subset\mathfrak{N}''$, et si l'on pose
$A''=A/\mathfrak{M}''$, $B''=B/\mathfrak{N}''$, on a de même un
homomorphisme surjectif
\[
  \varphi_{(\mathfrak{M}'',\mathfrak{N}''),(\mathfrak{M}',\mathfrak{N}')}:
  \Omega^1_{B''/A''}\to\Omega^1_{B'/A'}
\]
et ces homomorphismes forment évidemment un \emph{système projectif}. Si
l'on remarque que les
$((\mathfrak{J}\cap\mathfrak{U}')+\mathfrak{J}^2)/\mathfrak{J}^2$ forment
un système fondamental de voisinages de $0$ dans $\Omega^1_{B/A}$, on voit
donc que le \emph{séparé complété} $\widehat{\Omega}^1_{B/A}$ du $B$-module
topologique $\Omega^1_{B/A}$ est donné, à un isomorphisme canonique près, par
\begin{equation}
\tag{20.7.14.2}
\label{0.20.7.14.2-fr}
  \widehat{\Omega}^1_{B/A}=\varprojlim(\Omega^1_{B'/A'}).
\end{equation}

En outre, l'homomorphisme canonique
$j:\Omega^1_{B/A}\to\widehat{\Omega}^1_{B/A}$ est limite projective du
système projectif des $\varphi_{(\mathfrak{M}',\mathfrak{N}')}$, donc
$\varphi_{(\mathfrak{M}',\mathfrak{N}')}$ se factorise en
$\Omega^1_{B/A}\xrightarrow{j}\widehat{\Omega}^1_{B/A}\to
\Omega^1_{B'/A'}$, et comme il est surjectif on en conclut que
l'homomorphisme canonique
\begin{equation}
\tag{20.7.14.3}
\label{0.20.7.14.3-fr}
  \widehat{\Omega}^1_{B/A}\to\Omega^1_{B'/A'}
\end{equation}
est surjectif pour tout couple $(\mathfrak{M}',\mathfrak{N}')$.

On peut d'ailleurs dans ce qui précède remplacer partout $A'$ par $A$.

Enfin, si $L$ est un $B$-module topologique \emph{séparé et complet}, tout
$B$-homomorphisme continu de $\Omega^1_{B/A}$ dans $L$ se prolonge de façon
unique en un $\widehat B$-homomorphisme continu de
$\widehat{\Omega}^1_{B/A}$ dans $L$, et réciproquement un tel homomorphisme
redonne par composition avec $\Omega^1_{B/A}\to\widehat{\Omega}^1_{B/A}$ un
$B$-homomorphisme continu, si bien que l'on a un isomorphisme canonique
\[
  \operatorname{Hom.\ cont}_B(\Omega^1_{B/A},L)\xrightarrow{\sim}
  \operatorname{Hom.\ cont}_{\widehat B}(\widehat{\Omega}^1_{B/A},L).
\]
Plus particulièrement, si $L$ est un $B$-module discret annulé par un idéal
ouvert de $B$, on voit que l'isomorphisme canonique
\sref{0.20.4.8.2-fr} peut aussi s'écrire
\begin{equation}
\tag{20.7.14.4}
\label{0.20.7.14.4-fr}
  \operatorname{Hom.\ cont}_{\widehat B}(\widehat{\Omega}^1_{B/A},L)
  \xrightarrow{\sim}\operatorname{D\acute{e}r.\ cont}_A(B,L).
\end{equation}
\end{env}

\begin{proposition}[20.7.15]
\label{0.20.7.15-fr}
Soient $A$ un anneau discret, $B$ une $A$-algèbre topologique adique
($\mathbf{0}_{\mathrm I}$, 7.1.9), $\mathfrak{n}$ un idéal de définition de
$B$, $B_0=B/\mathfrak{n}$. On suppose que $\Omega^1_{B_0/A}$ et
$\mathfrak{n}/\mathfrak{n}^2$ sont des $B_0$-modules de type fini. Alors
$\widehat{\Omega}^1_{B/A}$ est un $\widehat B$-module de type fini.
\end{proposition}

\begin{proof}
Comme le carré de tout idéal ouvert de $B$ est ouvert, la topologie de
$\Omega^1_{B/A}$ est la topologie $\mathfrak{n}$-préadique
\sref{0.20.4.5-fr}. Compte tenu de l'hypothèse que $B$ est un anneau adique,
il suffit donc, en vertu de ($\mathbf{0}_{\mathrm I}$, 7.2.7 et 7.2.9), de
prouver
\oldpage[0\textsubscript{IV-1}]{152}
que $\Omega^1_{B/A}/\mathfrak{n}.\Omega^1_{B/A}=
\Omega^1_{B/A}\otimes_B B_0$ est un $B_0$-module de type fini. Mais cela
résulte de l'hypothèse et de la suite exacte \sref{0.20.5.12.1-fr}
\[
  \mathfrak{n}/\mathfrak{n}^2\to
  \Omega^1_{B/A}\otimes_B B_0\to\Omega^1_{B_0/A}\to0.
\]
\end{proof}

\begin{env}[20.7.16]
\label{0.20.7.16-fr}
La proposition \sref{0.20.7.15-fr} s'applique par exemple lorsque $A$ est un
corps $k$, $B=k'[[T_1,\ldots,T_n]]$ une algèbre de séries formelles munie de
sa topologie usuelle, $k'$ une extension \emph{finie} de $k$ (cf.\
\sref{0.21.9.2-fr}). On notera que le corps des fractions $K$ de $B$ a un
degré de transcendance infini sur $k$ ; lorsque $k$ est de caractéristique
$0$, on en déduit aussitôt (à l'aide de \sref{0.20.4.13-fr}, (i) et de
\sref{0.20.5.9-fr} notamment, en utilisant aussi le fait que toute dérivation
d'un corps de caractéristique $0$ se prolonge à toute extension) que
$\Omega^1_{B/k}$ \emph{n'est pas un $B$-module de type fini}.
\end{env}

\begin{env}[20.7.17]
\label{0.20.7.17-fr}
Soient $A$, $B$, $C$ trois anneaux topologiques, $u:A\to B$, $v:B\to C$ deux
homomorphismes continus ; remplaçant $A$, $B$, $C$ par des quotients par des
idéaux ouverts $A'=A/\mathfrak{M}'$, $B'=B/\mathfrak{N}'$,
$C'=C/\mathfrak{P}'$, avec $u(\mathfrak{M}')\subset\mathfrak{N}'$,
$v(\mathfrak{N}')\subset\mathfrak{P}'$, de sorte que l'on a des
homomorphismes $u':A'\to B'$, $v':B'\to C'$, on en déduit des homomorphismes
canoniques $u'_{C'/B'/A'}$, $v'_{C'/B'/A'}$ qui, en vertu de
\sref{0.20.5.4-fr}, forment des systèmes projectifs, et donnent par suite,
par passage à la limite, des homomorphismes canoniques, prolongements aux
séparés complétés des homomorphismes de la suite exacte
\sref{0.20.5.7.1-fr}
\[
  \Omega^1_{B/A}\otimes_B C\xrightarrow{v_{C/B/A}}\Omega^1_{C/A}
  \xrightarrow{u_{C/B/A}}\Omega^1_{C/B}\to0:
\]
\begin{align}
  \widehat v_{C/B/A}:\widehat{\Omega}^1_{B/A}
  \widehat\otimes_{\widehat B}\widehat C&\to\widehat{\Omega}^1_{C/A},
  \tag{20.7.17.1}\label{0.20.7.17.1-fr}\\
  \widehat u_{C/B/A}:\widehat{\Omega}^1_{C/A}&\to
  \widehat{\Omega}^1_{C/B}. \tag{20.7.17.2}\label{0.20.7.17.2-fr}
\end{align}
et dans la suite
\begin{equation}
\tag{20.7.17.3}
\label{0.20.7.17.3-fr}
  \widehat{\Omega}^1_{B/A}\widehat\otimes_{\widehat B}\widehat C
  \xrightarrow{\widehat v_{C/B/A}}\widehat{\Omega}^1_{C/A}
  \xrightarrow{\widehat u_{C/B/A}}\widehat{\Omega}^1_{C/B}\to0
\end{equation}
le composé de deux homomorphismes consécutifs est $0$, mais la suite
\emph{n'est pas nécessairement exacte}. Toutefois, si $B$ et $C$ sont
\emph{métrisables}, l'homomorphisme $\widehat u_{C/B/A}$ est \emph{surjectif},
et $\operatorname{Im}(\widehat v_{C/B/A})$ est \emph{dense} dans
$\operatorname{Ker}(\widehat u_{C/B/A})$ : cela résulte aussitôt
(cf.\ ($\mathbf{0}_{\mathrm I}$, 7.3.1)) de ce que, si
$(\mathfrak{M}_k)$ (resp.\ $(\mathfrak{N}_k)$) est une suite décroissante
d'idéaux de $B$ (resp.\ $C$), formant un système fondamental de voisinages de
$0$, et si l'on pose $B_k=B/\mathfrak{M}_k$, $C_k=C/\mathfrak{N}_k$, les
homomorphismes de transition
$\Omega^1_{C_{k+1}/B_{k+1}}\to\Omega^1_{C_k/B_k}$,
$\Omega^1_{C_{k+1}/A}\to\Omega^1_{C_k/A}$ et
$\Omega^1_{B_{k+1}/A}\to\Omega^1_{B_k/A}$ sont surjectifs
\sref{0.20.7.14-fr}.
\end{env}

\begin{proposition}[20.7.18]
\label{0.20.7.18-fr}
Soient $A$, $B$, $C$ trois anneaux topologiques, $u:A\to B$, $v:B\to C$
deux homomorphismes continus. On suppose $B$ et $C$ admissibles
($\mathbf{0}_{\mathrm I}$, 7.1.2) et métrisables. Pour que l'homomorphisme
canonique
\[
  \widehat v_{C/B/A}:\widehat{\Omega}^1_{B/A}
  \widehat\otimes_{\widehat B}\widehat C\to\widehat{\Omega}^1_{C/A}
\]
admette un inverse à gauche qui soit un $C$-homomorphisme continu, il faut et
il suffit que $C$ soit une $B$-algèbre formellement lisse relativement à $A$.
\end{proposition}

\begin{proof}
La condition est nécessaire en vertu de \sref{0.20.7.2-fr} et
\sref{0.19.1.6-fr}. Pour voir qu'elle est suffisante, notons que le $C$-module
topologique
$L=\widehat{\Omega}^1_{B/A}\widehat\otimes_{\widehat B}\widehat C$ est
métrisable et
\oldpage[0\textsubscript{IV-1}]{153}
complet en vertu de l'hypothèse, donc $E=D_C(L)$, muni de la topologie produit,
est métrisable et complet ; il est en outre admissible, car si $\mathfrak{R}$
est un idéal de définition dans $C$, la suite des
$(\mathfrak{R}\oplus L)^n=\mathfrak{R}^n\oplus\mathfrak{R}^{n-1}L$ tend vers
$0$. Comme l'application composée
\[
  D:B\xrightarrow{d_{B/A}}\Omega^1_{B/A}\to
  \widehat{\Omega}^1_{B/A}\widehat\otimes_{\widehat B}\widehat C=L
\]
est une $A$-dérivation continue de $B$ dans $L$, le $A$-homomorphisme continu
\[
  f:x\longmapsto(v(x),D(x))
\]
de $B$ dans $E$ définit sur $E$ une structure de $B$-extension. Comme $L$ est
un idéal fermé dans $E$, il résulte de \sref{0.19.9.5-fr} et de l'hypothèse
que l'application identique $C\to E/L$ (qui est un $B$-homomorphisme) se
factorise en $C\xrightarrow{g}E\to E/L$, où $g$ est un homomorphisme continu
tel que $g\circ v=f$ ; par suite \sref{0.20.1.3-fr}, $g$ est de la forme
$y\mapsto(y,D'(y))$, où $D'$ est une $B$-dérivation continue de $C$ dans $L$,
autrement dit $D'\circ v=D$. Compte tenu de \sref{0.20.7.14.4-fr}, on a
$D'=h\circ\widehat d_{C/A}$, où
$h:\widehat{\Omega}^1_{C/A}\to L$ est un $C$-homomorphisme continu ; mais on
a $\widehat d_{C/A}\circ v=\widehat v_{C/B/A}\circ D$ par définition, et
comme l'image de $B$ par $D$ engendre (topologiquement) le $C$-module $L$
\sref{0.20.4.7-fr}, la relation $D'\circ v=D$ donne bien
$h\circ\widehat v_{C/B/A}=1_L$.
\end{proof}

\begin{corollary}[20.7.19]
\label{0.20.7.19-fr}
Sous les hypothèses de \sref{0.20.7.18-fr}, si l'on suppose en outre que dans
$B$ et $C$ le carré d'un idéal ouvert soit ouvert, alors, pour que $C$ soit
une $B$-algèbre formellement lisse relativement à $A$, il faut et il suffit
que $\widehat v_{C/B/A}$ soit inversible à gauche.
\end{corollary}

\begin{proof}
En effet, les topologies de
$\widehat{\Omega}^1_{B/A}\widehat\otimes_{\widehat B}\widehat C$ et de
$\widehat{\Omega}^1_{C/A}$ sont alors déduites de celle de $C$
\sref{0.20.4.5-fr}, et tout $C$-homomorphisme de l'un dans l'autre est
nécessairement continu.
\end{proof}

\begin{env}[20.7.20]
\label{0.20.7.20-fr}
Soient $A$ un anneau topologique, $B$ une $A$-algèbre topologique
\emph{métrisable et complète}, $\mathfrak{R}$ un idéal \emph{fermé} de $B$,
$C=B/\mathfrak{R}$ l'anneau topologique quotient, qui est métrisable et
complet. Soit $(\mathfrak{M}_k)$ un système fondamental décroissant de
voisinages de $0$ dans $B$ formé d'idéaux, et posons
$B_k=B/\mathfrak{M}_k$,
$\mathfrak{R}_k=(\mathfrak{R}+\mathfrak{M}_k)/\mathfrak{M}_k$,
$C_k=B_k/\mathfrak{R}_k$. On a un système projectif d'homomorphismes
\[
  \delta_{C_k/B_k/A}:\mathfrak{R}_k/\mathfrak{R}_k^2\to
  \Omega^1_{B_k/A}\otimes_{B_k}C_k
\]
\sref{0.20.5.11.3-fr}, d'où l'on déduit en passant à la limite un
homomorphisme canonique
\[
  \widehat\delta_{C/B/A}:\mathfrak{R}/\mathfrak{R}^2\to
  \widehat{\Omega}^1_{B/A}\widehat\otimes_{\widehat B}\widehat C
\]
et en raisonnant comme dans \sref{0.20.7.17-fr}, on voit que l'homomorphisme
canonique
$\widehat u_{C/B/A}:\widehat{\Omega}^1_{B/A}
\widehat\otimes_{\widehat B}\widehat C\to\widehat{\Omega}^1_{C/A}$ est
\emph{surjectif} et que $\operatorname{Im}(\widehat\delta_{C/B/A})$ est
\emph{dense} dans $\operatorname{Ker}(\widehat u_{C/B/A})$.
\end{env}

% FIN SOURCE ega4/ega0-20-fr.tex

%
% DEBUT SOURCE ega4/ega0-21-fr.tex
\section{Différentielles dans les anneaux de caractéristique p}
Les résultats du présent paragraphe, de nature plus spéciale et technique que
ceux des §§19, 20 et 22, ne serviront qu'exceptionnellement dans le cours du
chap. IV. Leur rôle principal est ici dans la démonstration de trois théorèmes
du §22 (22.3.3, 22.5.8 et 22.7.3), dont le premier et le dernier interviennent
de façon essentielle dans la théorie «fine» des anneaux locaux noethériens du
chap. IV, §7.
\oldpage[0\textsubscript{IV-1}]{154}
\subsection{Systèmes de $p$-générateurs et $p$-bases}
\label{0.21.1-fr}

\begin{env}[21.1.1]
\label{0.21.1.1-fr}
Étant donné un nombre $p$ qui est, soit $0$, soit un nombre premier, nous
dirons qu'un anneau $A$ est de caractéristique $p$ s'il existe un homomorphisme
d'anneaux $P\to A$, où $P$ est le corps premier de caractéristique $p$~; on
notera que cet homomorphisme est alors unique, le composé
$\mathbf Z\to P\to A$ étant l'unique homomorphisme de $\mathbf Z$ dans $A$.
Si $A\neq0$, il revient au même de dire que $A$ contient un corps de
caractéristique $p$, l'image de $P$ étant nécessairement un corps isomorphe à
$P$ (et d'ailleurs le seul sous-corps de $A$ isomorphe à $P$).
\end{env}

\begin{env}[21.1.2]
\label{0.21.1.2-fr}
Si $p>0$, dire que $A$ est de caractéristique $p$ équivaut à dire que, dans
$A$, on a $p\cdot1=0$, ou encore $pA=0$. Si $p=0$, dire que $A$ est de
caractéristique $p$ équivaut à dire que pour tout entier $n\neq0$, $n\cdot1$
est inversible dans $A$. Si $A\neq0$, il ne peut exister qu'un seul $p$
(premier ou $0$) tel que $A$ soit de caractéristique $p$~; cela résulte de
l'identité de Bézout $ap+bq=1$ pour deux nombres premiers distincts $p,q$.
Par contre, l'anneau réduit à $0$ est de caractéristique $p$ pour tout $p$.
\end{env}

\begin{env}[21.1.3]
\label{0.21.1.3-fr}
Si $A$ est de caractéristique $p$, il en est de même de toute algèbre sur
$A$. En particulier, pour tout idéal premier $\mathfrak p$ de $A$, le corps
résiduel de $A$ en $\mathfrak p$ est de caractéristique $p$. Réciproquement,
si $p=0$ et si, pour tout idéal maximal $\mathfrak m$ de $A$, le corps
résiduel de $A$ en $\mathfrak m$ est de caractéristique $0$, il en est de même
de $A$, car pour tout entier $n\neq0$, $n\cdot1$ est alors inversible dans
tous les $A_{\mathfrak m}$, donc aussi dans $A$. Par contre, si $p>0$, un
anneau local peut avoir son corps résiduel de caractéristique $p$ sans être
lui-même de caractéristique $p$, comme le montre l'exemple de l'anneau
premier $\mathbf Z_{p\mathbf Z}$ (19.8.3) ou de l'anneau local artinien
$\mathbf Z/p^n\mathbf Z$ pour $n\geq2$, qui ne contiennent pas de corps.

Notons enfin que, pour un anneau (même réduit), les corps résiduels en ses
idéaux premiers peuvent avoir des caractéristiques différentes, comme le
montre l'exemple de $\mathbf Z$.
\end{env}

\begin{env}[21.1.4]
\label{0.21.1.4-fr}
Dans toute la suite de ce paragraphe, nous supposons fixé un nombre premier
$p$ et tous les anneaux seront supposés de caractéristique $p$, sauf mention
expresse du contraire. Pour un tel anneau $A$, l'application $x\mapsto x^p$
est un endomorphisme de $A$, que l'on note $F_A$~; si $A$ est réduit, $F_A$
est injectif. On pose $A^p=F_A(A)$ (sous-anneau de $A$ formé des $x^p$,
$x\in A$)~; on peut naturellement considérer $A$ comme une $A^p$-algèbre.

On peut aussi considérer $A$ comme une $A$-algèbre au moyen de
l'homomorphisme $F_A:x\mapsto x^p$ de $A$ dans $A$~; autrement dit, il
s'agit de l'algèbre $A^{(p)}$ pour laquelle le produit $\lambda\cdot x$ de
$x\in A$ par un scalaire $\lambda\in A$ est le produit $\lambda^p x$ dans
l'anneau $A$. Il est clair que, pour tout homomorphisme d'anneaux $u:A\to B$,
le couple $(u,u)$ est un di-homomorphisme d'algèbres $A^{(p)}\to B^{(p)}$.
Pour tout $A$-module $E$, on pose $E^{(p)}=E\otimes_A A^{(p)}$, où
$A^{(p)}$ est considéré comme un $(A,A)$-bimodule, la structure de
$A$-module à gauche étant celle qu'on vient de définir et la structure de
$A$-module à droite définie par la multiplication dans $A$~; ainsi, pour
$x\in E$, $\alpha,\beta\in A$, on a
\[
  \alpha(x\otimes\beta)=x\otimes(\alpha\beta),\qquad
  (\alpha x)\otimes\beta=x\otimes\alpha^p\beta.
\]
\end{env}

\oldpage[0\textsubscript{IV-1}]{155}
Posant $x^{(p)}=x\otimes1$, on a donc $(\alpha x)^{(p)}=\alpha^p x^{(p)}$.
Lorsque $F_A$ est injectif, on peut identifier $A^{(p)}$ à l'anneau $A$
considéré comme algèbre sur son sous-anneau $A^p$.

\begin{proposition}[21.1.5]
\label{0.21.1.5-fr}
Soient $A$, $B$ deux anneaux, $u:A\to B$ un homomorphisme.
\begin{enumerate}
  \item[(i)] Toute $A$-dérivation de $B$ dans un $B$-module $L$ est nulle
  dans $A[B^p]$ (donc est une $A[B^p]$-dérivation).
  \item[(ii)] Pour toute sous-$A$-algèbre $A'$ de $A[B^p]$, si
  $j:A'\to B$ est l'injection canonique, l'homomorphisme canonique
  \[
  j_{B/A'/A}:\Omega^1_{B/A}\longrightarrow\Omega^1_{B/A'}
  \]
  est bijectif.
  \item[(iii)] Supposons qu'il existe un entier $s\geq0$ tel que
  $u(A)\supset B^{p^s}$. Alors, dans l'anneau $B\otimes_A B$,
  $\mathfrak J_{B/A}$ est un nilidéal.
\end{enumerate}
\end{proposition}

\begin{proof}
(i) Par récurrence à partir de (20.1.1.1), on déduit, pour toute
$A$-dérivation $D:B\to L$, que $D(x^n)=nx^{n-1}D(x)$, d'où en particulier
$D(x^p)=0$.

(ii) Vu (20.4.8), l'assertion (ii) n'est qu'une traduction de (i), cette
dernière s'écrivant $\operatorname{Der}_{A[B^p]}(B,L)=\operatorname{Der}_A(B,L)$
pour tout $B$-module $L$.

(iii) Pour tout $x\in B$, on a
$(x\otimes1-1\otimes x)^{p^s}=x^{p^s}\otimes1-1\otimes x^{p^s}=0$,
puisque $x^{p^s}\in u(A)$. La conclusion résulte de ce que les éléments
$1\otimes x-x\otimes1$ engendrent $\mathfrak J_{B/A}$ (20.4.4).
\end{proof}

Il résulte aussitôt de (21.1.5) que, pour tout couple d'homomorphismes
d'anneaux $A\to B\to C$, on a
\begin{equation}
\tag{21.1.5.1}
\label{0.21.1.5.1-fr}
\Upsilon_{C/B/A}=\Upsilon_{C/B/A'}
\end{equation}
pour toute sous-$A$-algèbre $A'\subset A[B^p]$ de $B$.

D'autre part, (21.1.5) montre en particulier que, pour les modules de
différentielles «absolues»,
\begin{equation}
\tag{21.1.5.2}
\label{0.21.1.5.2-fr}
\Omega^1_A=\Omega^1_{A/A^p}.
\end{equation}

\begin{corollary}[21.1.6]
\label{0.21.1.6-fr}
Supposons que $B$ soit une $A$-algèbre de type fini et qu'il existe un
entier $s\geq0$ tel que $u(A)\supset B^{p^s}$. Si $\Omega^1_{B/A}=0$,
$u$ est surjectif.
\end{corollary}

\begin{proof}
\oldpage[0\textsubscript{IV-1}]{156}
D'après (21.1.5, (iii)), $\mathfrak J_{B/A}$ est un nilidéal~; en outre,
en vertu de (20.4.4) et de l'hypothèse, $\mathfrak J_{B/A}$ est un idéal de
type fini, donc il est nilpotent~; comme la relation $\Omega^1_{B/A}=0$
signifie que $\mathfrak J_{B/A}=\mathfrak J_{B/A}^2$, on en conclut que
$\mathfrak J_{B/A}=0$, ou encore que $\pi:B\otimes_A B\to B$ est bijectif.
Par ailleurs, tout élément de $B$ ayant sa puissance $p^s$-ième dans $u(A)$,
$B$ est entier sur $A$, et étant de type fini, c'est une $A$-algèbre finie.
On est ainsi ramené à prouver le lemme suivant (où l'on ne suppose pas que
les anneaux considérés sont de caractéristique $p$) :
\end{proof}

\begin{lemma}[21.1.6.1]
\label{0.21.1.6.1-fr}
Soient $R$ un anneau, $S$ une $R$-algèbre finie~; si l'homomorphisme
canonique $\pi:S\otimes_R S\to S$ est bijectif, alors l'homomorphisme
structural $u:R\to S$ est surjectif.
\end{lemma}

\begin{proof}
Il suffit de montrer que, pour tout idéal maximal $\mathfrak m$ de $R$, si
l'on pose $T=R-\mathfrak m$, l'homomorphisme
$u_{\mathfrak m}:R_{\mathfrak m}\to T^{-1}S$ est surjectif (Bourbaki,
\emph{Alg. comm.}, chap.~II, §~3, n$^\circ$~3, th.~1)~; or l'hypothèse
entraîne que l'homomorphisme
$(T^{-1}S)\otimes_{R_{\mathfrak m}}(S/\mathfrak mS)\to T^{-1}S$
est bijectif (0$_I$, 1.3.4), et comme $T^{-1}S$ est une
$R_{\mathfrak m}$-algèbre finie, on voit qu'on peut se borner au cas où
$R$ est un anneau local.

Désignant encore par $\mathfrak m$ son idéal maximal, il suffit de prouver que
$u\otimes1:R/\mathfrak m\to S/\mathfrak mS$ est surjectif, en vertu du
lemme de Nakayama ($S$ étant un $R$-module de type fini)~; comme
l'homomorphisme canonique
$(S/\mathfrak mS)\otimes_{R/\mathfrak m}(S/\mathfrak mS)\to S/\mathfrak mS$
est bijectif et que $S/\mathfrak mS$ est une $(R/\mathfrak m)$-algèbre finie,
on est finalement ramené au cas où $R$ est un corps~; mais alors les rangs de
$S\otimes_R S$ et de $S$ sur $R$ ne peuvent être égaux que si $S$ est de rang
$0$ ou $1$, donc si $u:R\to S$ est surjectif.
\end{proof}

\begin{proposition}[21.1.7]
\label{0.21.1.7-fr}
Soient $A$ un anneau, $B$ une $A$-algèbre, $(x_\alpha)_{\alpha\in I}$ une
famille d'éléments de $B$. Considérons les propriétés suivantes~:
\begin{enumerate}
  \item[a)] $B=A[B^p,(x_\alpha)]$, autrement dit la $A[B^p]$-algèbre
  $B$ est engendrée par la famille $(x_\alpha)_{\alpha\in I}$.
  \item[b)] Le $A[B^p]$-module $B$ est engendré par les monômes
  $\prod_\alpha x_\alpha^{n(\alpha)}$, où $(n(\alpha))_{\alpha\in I}$ est
  une famille d'entiers à support fini telle que $0\leq n(\alpha)<p$ pour
  tout $\alpha\in I$.
  \item[c)] Le $B$-module $\Omega^1_{B/A}$ est engendré par les
  $d_{B/A}(x_\alpha)$ $(\alpha\in I)$.
\end{enumerate}
Alors les propriétés a) et b) sont équivalentes et entraînent c)~; si, en
outre, $B$ est une $A[B^p]$-algèbre de type fini, c) est équivalente à a) et b).
\end{proposition}

\begin{proof}
Il est clair que b) entraîne a), et inversement a) entraîne b), car tout
monôme $\prod_\alpha x_\alpha^{m(\alpha)}$, où les $m(\alpha)$ sont des
entiers $\geq0$ (formant une famille à support fini), peut s'écrire
$(\prod_\alpha x_\alpha^{q(\alpha)})^p\prod_\alpha x_\alpha^{r(\alpha)}$,
en divisant chaque $m(\alpha)$ par $p$, ce qui donne
$m(\alpha)=p q(\alpha)+r(\alpha)$ avec $0\leq r(\alpha)<p$.
Le fait que a) entraîne c) résulte de (21.1.5, (ii)) et de (20.4.7).

Supposons inversement c) vérifiée et que $B$ soit une $A[B^p]$-algèbre de
type fini~; soit $B'$ la sous-$A[B^p]$-algèbre de $B$ engendrée par les
$x_\alpha$~; dans la suite exacte (20.5.7.1)
\[
\Omega^1_{B'/A[B^p]}\otimes_{B'}B\longrightarrow
\Omega^1_{B/A[B^p]}\longrightarrow\Omega^1_{B/B'}\longrightarrow0
\]
l'hypothèse entraîne que la flèche de gauche est surjective (compte tenu de
(21.1.5, (ii)))~; on a donc $\Omega^1_{B/B'}=0$, et comme
$B^p\subset B'\subset B$ et que $B$ est une $B'$-algèbre de type fini,
on a nécessairement $B'=B$ par (21.1.6).
\end{proof}

\begin{remark}[21.1.8]
\label{0.21.1.8-fr}
Lorsque $B$ est un corps, nous démontrerons l'équivalence des propriétés a),
b) et c) sans hypothèse de finitude (21.4.5).
\end{remark}

\begin{definition}[21.1.9]
\label{0.21.1.9-fr}
Soient $A$ un anneau (de caractéristique $p$), $B$ une $A$-algèbre. On dit
qu'une famille $(x_\alpha)_{\alpha\in I}$ d'éléments de $B$ est $p$-libre sur
$A$ (resp. un système de $p$-générateurs de $B$ sur $A$, resp. une $p$-base
de $B$ sur $A$) si la famille des monômes
$\prod_\alpha x_\alpha^{n(\alpha)}$ ($0\leq n(\alpha)<p$,
$(n(\alpha))_{\alpha\in I}$ à support fini) est une famille libre (resp. un
système de générateurs, resp. une base) dans le $A[B^p]$-module $B$.
\end{definition}

On a des définitions correspondantes pour un ensemble $M$ d'éléments de $B$,
en considérant la famille définie par l'injection canonique $M\to B$.
Lorsqu'on prend pour $A$ le corps premier $\mathbf F_p$ (auquel cas
$A[B^p]=B^p$), on omet la mention de $A$ dans la définition précédente (ou
l'on dit encore qu'une famille est «absolument» $p$-libre, resp. un système
«absolu» de $p$-générateurs, resp. une $p$-base «absolue»).

\oldpage[0\textsubscript{IV-1}]{157}
\begin{lemma}[21.1.10]
\label{0.21.1.10-fr}
Soient $C$ une sous-$A$-algèbre de $B$ telle que $B^p\subset C$, $M$ une
partie de $C$, $N$ une partie de $B$.
\begin{enumerate}
  \item[(i)] Si $M$ est un système de $p$-générateurs de $C$ sur $A$ et $N$
  un système de $p$-générateurs de $B$ sur $C$, alors $M\cup N$ est un système
  de $p$-générateurs de $B$ sur $A$.
  \item[(ii)] Supposons que $M$ soit une $p$-base de $C$ sur $A$~; alors,
  pour que $N$ soit $p$-libre sur $C$, il faut et il suffit que $M\cup N$
  soit $p$-libre sur $A$.
\end{enumerate}
\end{lemma}

\begin{proof}
On peut se borner au cas où $B^p\subset A\subset C\subset B$. Alors (i)
est un cas particulier du fait que si $P$ (resp. $Q$) est un système de
générateurs du $A$-module $C$ (resp. du $C$-module $B$), l'ensemble des
produits $xy$, où $x\in P$ et $y\in Q$, est un système de générateurs du
$A$-module $B$.

Conservant les mêmes notations, si $P$ est une base du $A$-module $C$,
dire que $Q$ est une famille libre sur $C$ signifie que la relation
$\sum_{\lambda,\mu}a_{\lambda\mu}x_\lambda y_\mu=0$, où
$a_{\lambda\mu}\in A$, $x_\lambda\in P$, $y_\mu\in Q$ (les $x_\lambda$
(resp. les $y_\mu$) étant deux à deux distincts) équivaut à
$a_{\lambda\mu}=0$ pour tout couple $(\lambda,\mu)$, ou encore à
$\sum_\lambda a_{\lambda\mu}x_\lambda=0$ pour tout $\mu$~; d'où
l'assertion (ii).
\end{proof}

\subsection{$p$-bases et lissité formelle}
\label{0.21.2-fr}

\begin{theorem}[21.2.1]
\label{0.21.2.1-fr}
Soient $B$ un anneau, $A$ un sous-anneau de $B$ tel que $B^p\subset A\subset B$,
$(x_\alpha)_{\alpha\in I}$ une $p$-base de $B$ sur $A$.
Soient $E$ une $A$-algèbre, $u:A\to E$ l'homomorphisme structural,
$q:E\to B$ un $A$-homomorphisme surjectif, $\mathfrak R$ son noyau, et
supposons que $\mathfrak R^p=0$. Alors~:
\begin{enumerate}
  \item[(i)] Pour qu'il existe un $A$-homomorphisme $v:B\to E$ inverse à
  droite de l'homomorphisme $q:E\to B$, il faut et il suffit que
  $u(q(z)^p)=z^p$ pour tout $z\in E$.
  \item[(ii)] Lorsque la condition de (i) est satisfaite, pour toute famille
  $(z_\alpha)_{\alpha\in I}$ d'éléments de $E$ telle que
  $q(z_\alpha)=x_\alpha$ pour tout $\alpha\in I$, il existe un
  $A$-homomorphisme $v:B\to E$ et un seul tel que $v(x_\alpha)=z_\alpha$
  pour tout $\alpha\in I$, et $v$ est inverse à droite de $q$.
\end{enumerate}
\end{theorem}

\begin{proof}
S'il existe un $A$-homomorphisme $v:B\to E$, on doit avoir $v(a)=u(a)$ pour
tout $a\in A\subset B$, donc $v(q(z)^p)=u(q(z)^p)$ pour tout $z\in E$,
puisque $B^p\subset A$. Mais on a $v(q(z)^p)=(v(q(z)))^p$ et, par définition,
$v(q(z))\equiv z\pmod{\mathfrak R}$ si $q\circ v=1_B$, donc
$(v(q(z)))^p=z^p$ puisque $\mathfrak R^p=0$~; d'où la nécessité de (i).
La suffisance de la condition (i) résultera de (ii).

Or, sous les hypothèses de (ii), l'unicité de $v$ est évidente puisque les
monômes $\prod_\alpha x_\alpha^{n(\alpha)}$ engendrent le $A$-module $B$~;
comme ces monômes forment en outre une \emph{base} du $A$-module $B$, il existe
une application $A$-\emph{linéaire} $v$ de $B$ dans $E$ telle que
$v(\prod_\alpha x_\alpha^{n(\alpha)})=\prod_\alpha z_\alpha^{n(\alpha)}$
pour toute famille $(n(\alpha))_{\alpha\in I}$ à support fini avec
$0\leq n(\alpha)<p$ pour tout $\alpha$. Reste à voir que $v$ est un
homomorphisme d'anneaux. Or, on peut écrire
$(\prod_\alpha x_\alpha^{m(\alpha)})(\prod_\alpha x_\alpha^{n(\alpha)})
=a\prod_\alpha x_\alpha^{r(\alpha)}$, où
$r(\alpha)=m(\alpha)+n(\alpha)$ si $m(\alpha)+n(\alpha)<p$,
$r(\alpha)=m(\alpha)+n(\alpha)-p$ dans le cas contraire, et $a\in A$
\oldpage[0\textsubscript{IV-1}]{158}
est le produit des $x_\alpha^p$ pour les $\alpha\in I$ tels que
$m(\alpha)+n(\alpha)\geq p$~; il s'agit de voir que
$u(a)=\prod_\alpha z_\alpha^p$~; mais comme $x_\alpha=q(z_\alpha)$, cela
résulte de l'hypothèse.
\end{proof}

\begin{corollary}[21.2.2]
\label{0.21.2.2-fr}
Soient $B$ un anneau, $A$ un sous-anneau de $B$ tel que $B^p\subset A\subset B$.
S'il existe une $p$-base de $B$ sur $A$, $B$ est une $A$-algèbre formellement
lisse relativement à $B^p$ (pour les topologies discrètes).
\end{corollary}

\begin{proof}
Soient $E$ une $A$-extension de $B$ par un $B$-module $L$, $q:E\to B$
l'augmentation, $u:A\to E$ l'homomorphisme structural. Dire que $E$ est
$B^p$-triviale signifie qu'il existe un homomorphisme d'anneaux $v:B\to E$
tel que $q(v(b))=b$ et $v(b^p)=u(b^p)$ pour tout $b\in B$. On en déduit,
pour $z\in E$, que $u(q(z)^p)=v(q(z)^p)=(v(q(z)))^p$~; mais, en vertu de
la relation $L^2=0$, on a aussi $L^p=0$, et comme $v(q(z))-z\in L$, on a
$(v(q(z)))^p=z^p$. La condition de (21.2.1), (i), est donc satisfaite et
$E$ est $B^p$-triviale.
\end{proof}

\begin{corollary}[21.2.3]
\label{0.21.2.3-fr}
Soient $B$ un anneau, $A$ un sous-anneau de $B$ tel que $B^p\subset A\subset B$,
et $\left(x_\alpha\right)_{\alpha\in I}$ une $p$-base de $B$ sur $A$. Soit
$L$ un $B$-module. Alors~:
\begin{enumerate}
  \item[(i)] Pour qu'une dérivation $D$ de $A$ dans $L$ se prolonge en une
  dérivation de $B$ dans $L$, il faut et il suffit que $D$ s'annule dans $B^p$.
  \item[(ii)] Si $D$ s'annule dans $B^p$, alors, pour toute famille
  $(y_\alpha)_{\alpha\in I}$ d'éléments de $L$, il existe une dérivation
  $D'$ de $B$ dans $L$ et une seule prolongeant $D$ et telle que
  $D'(x_\alpha)=y_\alpha$ pour tout $\alpha\in I$.
\end{enumerate}
\end{corollary}

\begin{proof}
Étant donnée une dérivation $D$ de $A$ dans $L$, considérons l'anneau
$E=D_B(L)$ et l'homomorphisme $u:A\to E$ défini par $u(a)=(a,D(a))$~; un
$A$-homomorphisme $v:B\to E$ inverse à droite de l'homomorphisme canonique
$q:E\to B$ est alors de la forme $x\mapsto(x,D'(x))$, où $D'$ est une
dérivation de $B$ dans $L$ prolongeant $D$ (20.1.5)~; comme
$u(q(z)^p)=(q(z)^p,D(q(z)^p))$ pour $z\in E$, le corollaire résulte aussitôt
de (21.2.1) appliqué à $E$.
\end{proof}

Il résulte de (21.2.3) que la suite
\begin{equation}
\tag{21.2.3.1}
\label{0.21.2.3.1-fr}
0\to\operatorname{Der}_A(B,L)\to\operatorname{Der}_{B^p}(B,L)
\to\operatorname{Der}_{B^p}(A,L)\to0
\end{equation}
(cf. (20.2.2.1)) est exacte, et que l'application
\begin{equation}
\tag{21.2.3.2}
\label{0.21.2.3.2-fr}
D'\longmapsto(D'(x_\alpha))_{\alpha\in I}
\end{equation}
est un isomorphisme du $B$-module $\operatorname{Der}_A(B,L)$ sur le
$B$-module produit $L^I$ (en faisant $D=0$ dans (21.2.3)).

\begin{corollary}[21.2.4]
\label{0.21.2.4-fr}
Sous les hypothèses de (21.2.3), la suite
\begin{equation}
\tag{21.2.4.1}
\label{0.21.2.4.1-fr}
0\to\Omega^1_{A/B^p}\otimes_A B\to\Omega^1_{B/B^p}
\to\Omega^1_{B/A}\to0
\end{equation}
est exacte et scindée, et la famille $(d_{B/A}(x_\alpha))_{\alpha\in I}$ est
une base du $B$-module $\Omega^1_{B/A}$.
\end{corollary}

\begin{proof}
Cela résulte aussitôt de (21.2.3) et de la formule
$\operatorname{Der}_A(B,L)=\operatorname{Hom}_B(\Omega^1_{B/A},L)$ (20.4.8).
\end{proof}

\begin{corollary}[21.2.5]
\label{0.21.2.5-fr}
Soient $A$ un anneau, $B$ une $A$-algèbre, $(x_\alpha)_{\alpha\in I}$ une
$p$-base de $B$ sur $A$. Alors $(d_{B/A}(x_\alpha))_{\alpha\in I}$ est une
base du $B$-module $\Omega^1_{B/A}$.
\end{corollary}

\begin{proof}
En utilisant (21.1.5), (ii), on se ramène au cas où $A=A[B^p]$ et il suffit
alors d'appliquer (21.2.4).
\end{proof}

\oldpage[0\textsubscript{IV-1}]{159}
\begin{env}[21.2.6]
\label{0.21.2.6-fr}
Soient $A$, $B$ deux anneaux, $u:A\to B$ un homomorphisme d'anneaux~; on a
(avec les notations de (21.1.4)) le diagramme commutatif
\begin{equation}
\tag{21.2.6.1}
\label{0.21.2.6.1-fr}
\xymatrix{
 A^{(p)} \ar[r]^{u} & B^{(p)}\\
 A \ar[u]^{F_A} \ar[r]^{u} & B \ar[u]_{F_B}
}
\end{equation}
On en déduit canoniquement un homomorphisme de $B$-algèbres
\begin{equation}
\tag{21.2.6.2}
\label{0.21.2.6.2-fr}
u\otimes1:A^{(p)}\otimes_A B\longrightarrow B^{(p)}
\end{equation}
dont l'image est l'anneau $A[B^p]$ (qui est une $B$-algèbre pour
l'homomorphisme $F_B:B\to B^p$).
\end{env}

\begin{theorem}[21.2.7]
\label{0.21.2.7-fr}
Soient $A$ un anneau, $B$ une $A$-algèbre, $u:A\to B$ l'homomorphisme
structural. On suppose vérifiées les conditions suivantes~:
\begin{enumerate}
  \item[(i)] L'homomorphisme (21.2.6.2) déduit de $u$ est injectif.
  \item[(ii)] $B$ admet une $p$-base $(x_\alpha)_{\alpha\in I}$ par rapport à $A$.
\end{enumerate}
Alors $B$ est une $A$-algèbre formellement lisse (pour les topologies
discrètes). De façon plus précise, munissons $A$ et $B$ des topologies
discrètes~; soient $E$ une $A$-algèbre topologique admissible
(0$_I$, 7.1.2), $\mathfrak R$ un idéal de définition de $E$, $C=E/\mathfrak R$,
$v:B\to C$ un $A$-homomorphisme, $q:E\to C$ l'augmentation. Alors, pour
toute famille $(z_\alpha)_{\alpha\in I}$ d'éléments de $E$ telle que
$q(z_\alpha)=v(x_\alpha)$ pour tout $\alpha\in I$, il existe un
$A$-homomorphisme $w:B\to E$ et un seul tel que $v=q\circ w$ et
$w(x_\alpha)=z_\alpha$ pour tout $\alpha\in I$.
\end{theorem}

\begin{proof}
Considérons d'abord le cas où $E$ est discret, donc $\mathfrak R$ nilpotent.
On peut se borner au cas où $v$ est surjectif ; en outre, le raisonnement de
(19.4.3) permet de supposer $\mathfrak R^2=0$ et par suite $\mathfrak R^p=0$.
Enfin, en considérant l'image réciproque par $v$ de l'extension $E$ de $C$
par $\mathfrak R$, on peut se borner au cas où $C=B$ et où $v$ est l'identité
(19.4.4). Comme l'homomorphisme d'anneaux $F_E:z\rightsquigarrow z^p$
s'annule dans $\mathfrak R$, il se factorise en
$E\xrightarrow q B\xrightarrow{F'}E$,
et $F'$, considéré comme homomorphisme de $B$ dans $E^{(p)}$, est un
\emph{$A$-homomorphisme} par définition de la structure de $A$-algèbre de $E^{(p)}$
à l'aide de l'homomorphisme composé
$A\xrightarrow r E\xrightarrow{F_E}E^{(p)}$, où $r:A\to E$ est l'homomorphisme
structural de la $A$-algèbre $E$ ; d'ailleurs
$r:A^{(p)}\to E^{(p)}$ est aussi un $A$-homomorphisme. Il y a donc un unique
$A$-homomorphisme $f:A^{(p)}\otimes_A B\to E^{(p)}$ tel que les composés de $f$
et des homomorphismes canoniques $A^{(p)}\to A^{(p)}\otimes_A B$ et
$B\to A^{(p)}\otimes_A B$ soient respectivement $r$ et $F'$. Or, par hypothèse,
on peut identifier $A^{(p)}\otimes_A B$ à $A[B^p]$, les homomorphismes canoniques
$A^{(p)}\to A^{(p)}\otimes_A B$ et $B\to A^{(p)}\otimes_A B$ s'identifiant
respectivement à $u$ et à $F_B$. On peut donc maintenant considérer $E$ comme
une $A[B^p]$-algèbre au moyen de l'homomorphisme $f$, et, par construction, on a
$f(q(z)^p)=z^p$ pour tout $z\in E$ ; on est donc dans les conditions
d'application de (21.2.1), d'où le théorème dans ce cas.

Pour passer au cas général, considérons un système fondamental
$(\mathfrak J_\lambda)$ d'idéaux ouverts de $E$, et posons
$E_\lambda=E/\mathfrak J_\lambda$,
$\mathfrak R_\lambda=(\mathfrak R+\mathfrak J_\lambda)/\mathfrak J_\lambda$,
$C_\lambda=E_\lambda/\mathfrak R_\lambda=E/(\mathfrak R+\mathfrak J_\lambda)$ ;
comme $E$ est admissible, on a $E=\varprojlim E_\lambda$ (0$_I$, 7.2.2).
Pour tout couple $(\alpha,\lambda)$, désignons par $z_{\alpha\lambda}$
l'image canonique de $z_\alpha$ dans $E_\lambda$, par
$p_\lambda:C\to C_\lambda$ l'homomorphisme canonique,
\oldpage[0\textsubscript{IV-1}]{160}
par $q_\lambda$ l'homomorphisme canonique $E_\lambda\to C_\lambda$.
Comme par hypothèse $\mathfrak R_\lambda$ est nilpotent dans $E_\lambda$,
la première partie de la démonstration prouve l'existence et l'unicité d'un
$A$-homomorphisme $w_\lambda:B\to E_\lambda$ tel que
$p_\lambda\circ v=q_\lambda\circ w_\lambda$ et
$w_\lambda(x_\alpha)=z_{\alpha\lambda}$ pour tout $\alpha$.
L'unicité des $w_\lambda$ montre alors aussitôt que $(w_\lambda)$ est un
système projectif d'homomorphismes, et $w=\varprojlim w_\lambda$ répond
à la question.
\end{proof}

\begin{remark}[21.2.8]
\label{0.21.2.8-fr}
La vérification de l'existence et de l'unicité de l'homomorphisme $w:B\to E$
tel que $w(x_\alpha)=z_\alpha$ pour tout $\alpha$ peut se déduire directement
du fait que $B$ est une $A$-algèbre formellement lisse et que les
$d_{B/A}(x_\alpha)$ forment une base du $B$-module $\Omega^1_{B/A}$ (21.2.4),
sans faire intervenir le fait qu'il s'agit d'une $p$-base (de sorte que le
résultat est valable sans supposer les anneaux de caractéristique $p$).
En effet, on peut se borner au cas où $\mathfrak R^2=0$~; comme
$\operatorname{Der}_A(B,\mathfrak R)=\operatorname{Hom}_B(\Omega^1_{B/A},\mathfrak R)$
(20.4.8), il existe une $A$-dérivation $D$, et une seule, de $B$ dans
$\mathfrak R$ telle que $D(x_\alpha)=z'_\alpha$ pour toute famille
$(z'_\alpha)$ d'éléments de $\mathfrak R$~; la conclusion résulte de (20.1.1).
\end{remark}

\subsection{$p$-bases et modules d'imperfection}
\label{0.21.3-fr}

\begin{env}[21.3.1]
\label{0.21.3.1-fr}
Soient $A$, $B$ deux anneaux (de caractéristique $p$), $i:A\to B$,
$j:B\to A$ deux homomorphismes d'anneaux tels que
\begin{align}
\tag{21.3.1.1}\label{0.21.3.1.1-fr}
j(i(a))&=a^p\quad\text{pour tout }a\in A,\\
\tag{21.3.1.2}\label{0.21.3.1.2-fr}
i(j(b))&=b^p\quad\text{pour tout }b\in B.
\end{align}
Le plus souvent, $i$ sera injectif, de sorte que $A$ s'identifiera par
$i:A\to B$ à un sous-anneau de $B$~; une fois cette identification faite,
l'existence de $j$ implique que $B^p\subset A$, et $j$ s'identifie alors à $F_B$.
\end{env}

\begin{env}[21.3.2]
\label{0.21.3.2-fr}
Si $h:A^p\to A$ est l'injection canonique, on a, d'après (21.3.1.1) et
(21.3.1.2), un diagramme commutatif
\begin{equation}
\tag{21.3.2.1}\label{0.21.3.2.1-fr}
\xymatrix{
 B \ar[r]^{j} & A\\
 A \ar[u]^{i} \ar[r]^{F_A} & A^p \ar[u]_{h}
}
\end{equation}
si bien que le couple $(j,F_A)$ peut être considéré comme un di-homomorphisme
de la $A$-algèbre $B$ (pour $i$) dans la $A^p$-algèbre $A$ (pour $h$).
On en déduit un homomorphisme canonique de $A$-modules
\begin{equation}
\tag{21.3.2.2}\label{0.21.3.2.2-fr}
\pi_{B/A}:\Omega^1_{B/A}\otimes_B A_{[j]}\longrightarrow
\Omega^1_{A/A^p}=\Omega^1_A
\end{equation}
(cf. (20.5.4)~; l'identification de $\Omega^1_{A/A^p}$ et du module des
différentielles absolues $\Omega^1_A$ provient de (21.1.5), (ii)).
\begin{align}
\tag{21.3.2.3}\label{0.21.3.2.3-fr}
\Theta_{B/A}&=\Omega^1_{B/A}\otimes_B A_{[j]},\\
\tag{21.3.2.4}\label{0.21.3.2.4-fr}
\Xi_{B/A}&=\operatorname{Ker}(\pi_{B/A}).
\end{align}
\oldpage[0\textsubscript{IV-1}]{161}
De sorte que l'on a la suite exacte
\begin{equation}
\tag{21.3.2.5}\label{0.21.3.2.5-fr}
0\to\Xi_{B/A}\to\Theta_{B/A}\xrightarrow{\pi_{B/A}}\Omega^1_A
\end{equation}
(on remarquera que ces notations peuvent prêter à confusion puisque
$\Theta_{B/A}$ et $\Xi_{B/A}$ dépendent non seulement de $A$ et de $B$,
mais aussi de $i$ et de $j$).
\end{env}

\begin{env}[21.3.3]
\label{0.21.3.3-fr}
Comme, en vertu de (21.3.1.2), on a $F_B=i\circ j$, on peut écrire pour
tout $B$-module $M$ (cf. (21.1.4))
\begin{equation}
\tag{21.3.3.1}\label{0.21.3.3.1-fr}
M^{(p)}=M\otimes_B B^{(p)}=(M\otimes_B A_{[j]})\otimes_A B_{[i]},
\end{equation}
d'où en particulier
\begin{equation}
\tag{21.3.3.2}\label{0.21.3.3.2-fr}
(\Omega^1_{B/A})^{(p)}=\Theta_{B/A}\otimes_A B_{[i]},
\end{equation}
et l'on déduit donc de (21.3.2.2) un homomorphisme canonique de $B$-modules
\begin{equation}
\tag{21.3.3.3}\label{0.21.3.3.3-fr}
\pi_{B/A}\otimes1_B:(\Omega^1_{B/A})^{(p)}\longrightarrow
\Omega^1_A\otimes_A B_{[i]}.
\end{equation}
Compte tenu de (20.5.4.2), l'image de cet homomorphisme est le $B$-module
engendré par les éléments $d_A(j(b))\otimes1$ pour $b\in B$. Leur image
par l'homomorphisme canonique $i_{B/A}:\Omega^1_A\otimes_A B\to\Omega^1_B$
déduit de $i$ (20.5.2) est donc dans le sous-$B$-module de $\Omega^1_B$
engendré par les $d_B(i(j(b)))$ pour $b\in B$~; en vertu de (21.3.1.2),
cette image est nulle. Autrement dit, on a une suite d'homomorphismes
\begin{equation}
\tag{21.3.3.4}\label{0.21.3.3.4-fr}
0\to\Xi_{B/A}\otimes_A B_{[i]}\to(\Omega^1_{B/A})^{(p)}
\to\Upsilon_{B/A}
\end{equation}
qui n'est pas nécessairement exacte, mais où le composé de deux
homomorphismes consécutifs est nul.
\end{env}

\begin{proposition}[21.3.4]
\label{0.21.3.4-fr}
Si $B$ est un $A$-module plat (pour $i$), la suite (21.3.3.4) est exacte.
\end{proposition}
\begin{proof}
Cela résulte de la définition de la platitude. La proposition s'applique en
particulier lorsque $B^p\subset A\subset B$, $i$ et $j$ étant respectivement
l'injection canonique et $F_B$, et que $B$ admet une $p$-base sur $A$ (de sorte
que $B$ est alors un $A$-module libre). Mais même dans ce cas le noyau
$\Xi_{B/A}$ n'est pas nécessairement nul. Toutefois~:
\end{proof}

\begin{proposition}[21.3.5]
\label{0.21.3.5-fr}
Soient $B$ un anneau réduit, $A$ un sous-anneau de $B$ tel que
$B^p\subset A\subset B$. Supposons qu'il existe une $p$-base
$(x_\lambda)_{\lambda\in L}$ de $B$ sur $A$ et une $p$-base
$(y_\mu)_{\mu\in M}$ de $A$ sur $B^p$~; alors l'homomorphisme canonique
(21.3.3.4)
\begin{equation}
\tag{21.3.5.1}\label{0.21.3.5.1-fr}
(\Omega^1_{B/A})^{(p)}\longrightarrow\Upsilon_{B/A}
\end{equation}
est bijectif, et les éléments $d_A(x_\lambda^p)\otimes1$ forment une base
du $B$-module $\Upsilon_{B/A}$.
\end{proposition}

\begin{proof}
\oldpage[0\textsubscript{IV-1}]{162}
Comme $B$ est réduit, $x\mapsto x^p$ est un isomorphisme de structure de
$B$ sur $B^p$~; par transport de structure au moyen de cet isomorphisme,
on voit que $(x_\lambda^p)_{\lambda\in L}$ et les $y_\mu$ forment une $p$-base
de $B^p$ sur $A^p$ (21.1.10)~; on en conclut par (21.1.5), (ii), et (21.2.5)
que les $d_A(x_\lambda^p)$ et les $d_A(y_\mu)$ forment une base du
$A$-module $\Omega^1_A$. Or l'image de $d_A(x_\lambda^p)\otimes1$ par
$i_{B/A}$ est $d_B(x_\lambda^p)=0$~; d'autre part, l'image de $d_A(y_\mu)$
par $i_{B/A}$ est $d_B(i(y_\mu))$, et comme les $y_\mu$ et les $x_\lambda$
forment une $p$-base de $B$ sur $B^p$ (21.1.10), les $d_B(y_\mu)$ sont une
partie d'une base de $\Omega^1_B$ (21.2.5), donc sont linéairement indépendants
sur $B$. On en déduit que le noyau de $i_{B/A}$ a pour base les
$d_A(x_\lambda^p)\otimes1$~; comme ce sont les images par (21.3.5.1) des
éléments correspondants de $(\Omega^1_{B/A})^{(p)}$, et que ces derniers
forment une base du $B$-module $(\Omega^1_{B/A})^{(p)}$ (21.2.5),
l'homomorphisme (21.3.5.1) de $B$-modules est bijectif.
\end{proof}

\begin{remarks}[21.3.6]
\label{0.21.3.6-fr}
\begin{enumerate}
  \item[(i)] Soient $A'$, $B'$ deux anneaux de caractéristique $p$, et
  supposons donnés deux homomorphismes $i':A'\to B'$, $j':B'\to A'$ vérifiant
  (21.3.1.1) et (21.3.1.2), ainsi que des homomorphismes d'anneaux
  $f:A\to A'$, $g:B\to B'$ rendant commutatif le diagramme
  \[
  \xymatrix{
  A'\ar[r]^{i'}&B'\ar[r]^{j'}&A'\\
  A\ar[u]^f\ar[r]^i&B\ar[u]^g\ar[r]^j&A\ar[u]^f
  }
  \]
  Alors l'homomorphisme canonique $\Omega^1_{B/A}\to\Omega^1_{B'/A'}$
  (20.5.4.3) donne un di-homomorphisme canonique
  $\Theta_{B/A}\to\Theta_{B'/A'}$ rendant commutatif le diagramme
  \[
  \xymatrix{
  0\ar[r]&\Xi_{B/A}\ar[r]\ar[d]&\Theta_{B/A}\ar[r]^{\pi_{B/A}}\ar[d]&\Omega^1_A\ar[d]\\
  0\ar[r]&\Xi_{B'/A'}\ar[r]&\Theta_{B'/A'}\ar[r]^{\pi_{B'/A'}}&\Omega^1_{A'}
  }
  \]
  \item[(ii)] Soit $A'$ une $A$-algèbre quelconque et posons
  $B'=B\otimes_A A'$~; alors $i'=i\otimes1:A'\to B'$ et
  $j'=j\otimes1:B'\to A'$ vérifient (21.3.1.1) et (21.3.1.2), et l'on a
  $\Omega^1_{B'/A'}=\Omega^1_{B/A}\otimes_A A'=\Omega^1_{B/A}\otimes_B B'$
  (20.5.5)~; on en déduit un $A'$-isomorphisme canonique
  \begin{equation}
  \tag{21.3.6.1}\label{0.21.3.6.1-fr}
  \Theta_{B/A}\otimes_A A'\xrightarrow{\sim}\Theta_{B'/A'}
  \end{equation}
  et aussi un $B'$-homomorphisme canonique
  \begin{equation}
  \tag{21.3.6.2}\label{0.21.3.6.2-fr}
  (\Omega^1_{B/A})^{(p)}\otimes_B B'\xrightarrow{\sim}
  (\Omega^1_{B'/A'})^{(p)}.
  \end{equation}
\end{enumerate}
\end{remarks}

\subsection{Cas des extensions de corps}
\label{0.21.4-fr}

\begin{env}[21.4.0]
\label{0.21.4.0-fr}
Soient $K$ un corps de caractéristique $p>0$, $k$ un sous-corps de $K$~;
alors l'anneau $k[K^p]$ est égal au corps $k(K^p)$ puisque $k$ est algébrique
sur $K^p$. On pourra donc appliquer les résultats des numéros précédents en
remplaçant partout $A$, $B$ et $A[B^p]$ par $k$, $K$ et $k(K^p)$.
\end{env}

\begin{lemma}[21.4.1]
\label{0.21.4.1-fr}
Soient $k$ un corps de caractéristique $p>0$, $K$ une extension de $k$.
Pour qu'un élément $x\in K$ soit $p$-libre sur $k$, il faut et il suffit que
$x\notin k(K^p)$.
\end{lemma}

\begin{proof}
En effet, $x$ est racine du polynôme $X^p-x^p$ de $k(K^p)[X]$, et l'on sait
(Bourbaki, \emph{Alg.}, chap.~V, §~8, n$^\circ$~1, prop.~1) que si
$x\notin k(K^p)$ ce polynôme est irréductible, de sorte que les éléments
$1,x,\ldots,x^{p-1}$ forment une base du $k(K^p)$-module $k(K^p)(x)$.
\end{proof}

\begin{theorem}[21.4.2]
\label{0.21.4.2-fr}
Soient $k$ un corps de caractéristique $p>0$, $K$ une extension de $k$,
$S$ un système de $p$-générateurs de $K$ sur $k$, $L\subset S$ une partie
$p$-libre sur $k$. Il existe alors une $p$-base $B$ de $K$ sur $k$ telle que
$L\subset B\subset S$. En particulier, toute extension de $k$ admet une
$p$-base sur $k$.
\end{theorem}

\begin{proof}
On peut se borner au cas où $K^p\subset k$. En vertu du théorème de Zorn,
il existe dans $K$ une partie $B$ telle que $L\subset B\subset S$, $p$-libre
sur $k$ et maximale parmi les parties de $S$ ayant ces propriétés. Il suffit
de voir que le sous-corps $K'$ de $K$ engendré par $k$ et $B$ est égal à $K$.
Dans le cas contraire, il existerait $x\in S$ non dans $K'$~; comme
$K^p\subset k$, on a $K'=K'(K^p)$~; donc $x$ serait $p$-libre sur $K'$
(21.4.1), et par suite $B\cup\{x\}$ serait $p$-libre sur $k$ (21.1.10),
contrairement à la définition de $B$. La dernière assertion de (21.4.2)
s'obtient en prenant $L=\varnothing$, $S=K$.
\end{proof}

\begin{corollary}[21.4.3]
\label{0.21.4.3-fr}
Soient $k$ un corps de caractéristique $p>0$, $K$ une extension de $k$.
Pour qu'une famille $(x_\alpha)$ d'éléments de $K$ soit $p$-libre sur $k$,
il faut et il suffit que, pour tout $\alpha$, $x_\alpha$ n'appartienne pas au
corps $K_\alpha$ engendré par $k(K^p)$ et par les $x_\beta$ d'indice
$\beta\neq\alpha$.
\end{corollary}

\begin{proof}
La condition est nécessaire en vertu de (21.1.10). Inversement, supposons-la
remplie~; on peut se borner au cas où $(x_\alpha)$ est un système de
$p$-générateurs de $K$. Il existe alors une sous-famille de $(x_\alpha)$ qui
est une $p$-base de $K$ (21.4.2), mais cette famille ne peut être distincte de
$(x_\alpha)$, sans quoi, en vertu de l'hypothèse, ce ne serait pas une famille
de $p$-générateurs de $K$.
\end{proof}

\oldpage[0\textsubscript{IV-1}]{163}
\begin{corollary}[21.4.4]
\label{0.21.4.4-fr}
Soient $k$ un corps de caractéristique $p>0$, $K$ une extension de $k$.
Pour que $\Omega^1_{K/k}=0$, il faut et il suffit que $K=k(K^p)$.
En particulier, pour que $\Omega^1_K=0$, il faut et il suffit que $K$ soit
un corps parfait.
\end{corollary}

\begin{proof}
En effet, si $(x_\alpha)$ est une $p$-base de $K$ sur $k$, les
$d_{K/k}(x_\alpha)$ forment une base de l'espace vectoriel $\Omega^1_{K/k}$
(21.2.5).
\end{proof}

\begin{theorem}[21.4.5]
\label{0.21.4.5-fr}
Soient $k$ un corps de caractéristique $p>0$, $K$ une extension de $k$,
$(x_\alpha)$ une famille d'éléments de $K$. Pour que $(x_\alpha)$ soit
$p$-libre sur $k$ (resp. une famille de $p$-générateurs de $K$ sur $k$,
resp. une $p$-base de $K$ sur $k$), il faut et il suffit que la famille
$(d_{K/k}(x_\alpha))$ soit une famille libre (resp. un système de
générateurs, resp. une base) dans l'espace vectoriel $\Omega^1_{K/k}$.
\end{theorem}

\begin{proof}
\oldpage[0\textsubscript{IV-1}]{164}
Soit $K'$ le sous-corps (égal au sous-anneau) de $K$ engendré par
$k(K^p)$ et les $x_\alpha$~; compte tenu de (21.4.3), il suffit de voir que,
pour chaque $\alpha$, $x_\alpha$ n'appartient pas à $K_\alpha$. Or, dans la
suite exacte
\[
\Omega^1_{K_\alpha/k(K^p)}\otimes_{K_\alpha}K\longrightarrow
\Omega^1_{K/k(K^p)}\longrightarrow\Omega^1_{K/K_\alpha}\longrightarrow0,
\]
les images par la flèche de gauche des $d_{K_\alpha/k}(x_\beta)$ pour
$\beta\neq\alpha$ sont les $d_{K/k}(x_\beta)$, donc engendrent un sous-espace
vectoriel de $\Omega^1_{K/k(K^p)}$ ne contenant pas $d_{K/k}(x_\alpha)$~; comme
ce sous-espace est le noyau de la flèche de droite, on voit que
$d_{K/K_\alpha}(x_\alpha)\neq0$, donc $x_\alpha\notin K_\alpha$.

Si maintenant $(x_\alpha)$ est une famille $p$-libre, elle est une sous-famille
d'une $p$-base de $K$ sur $k$ (21.4.2)~; les $d_{K/k}(x_\alpha)$ font donc
partie d'une base de l'espace vectoriel $\Omega^1_{K/k}$ (21.2.5), et sont
par suite linéairement indépendants sur $K$. Réciproquement, si les
$d_{K/k}(x_\alpha)$ sont linéairement indépendants sur $K$, les deux
paragraphes précédents montrent que $(x_\alpha)$ est $p$-libre sur $k$.
\end{proof}

\begin{corollary}[21.4.6]
\label{0.21.4.6-fr}
Soient $k$ un corps de caractéristique $p>0$, $K$ une extension de $k$,
$x$ un élément de $K$. Les trois conditions suivantes sont équivalentes~:
\begin{enumerate}
  \item[a)] $x\notin k(K^p)$.
  \item[b)] $d_{K/k}(x)\neq0$.
  \item[c)] L'élément $x$ est $p$-libre sur $k$.
\end{enumerate}
\end{corollary}

\begin{proposition}[21.4.7]
\label{0.21.4.7-fr}
Soient $L$ un corps de caractéristique $p>0$, $K$ un sous-corps de $L$ tel
que $L^p\subset K\subset L$. Alors la suite d'espaces vectoriels sur $L$
\begin{equation}
\tag{21.4.7.1}\label{0.21.4.7.1-fr}
0\to\Omega^1_{K/L^p}\otimes_KL\to\Omega^1_{L/L^p}
\to\Omega^1_{L/K}\to0
\end{equation}
est exacte~; en d'autres termes, on a $\Upsilon_{L/K/L^p}=0$.
\end{proposition}
\begin{proof}
C'est un cas particulier de (21.2.4), compte tenu de (21.4.2).
\end{proof}

\begin{proposition}[21.4.8]
\label{0.21.4.8-fr}
Sous les hypothèses de (21.4.7), l'homomorphisme canonique
\begin{equation}
\tag{21.4.8.1}\label{0.21.4.8.1-fr}
(\Omega^1_{L/K})^{(p)}\longrightarrow\Upsilon_{L/K}
\end{equation}
est bijectif~; si $(x_\lambda)$ est une $p$-base de $L$ sur $K$, les éléments
$d_K(x_\lambda^p)\otimes1$ forment une base de l'espace vectoriel $\Upsilon_{L/K}$.
\end{proposition}
\begin{proof}
C'est un cas particulier de (21.3.5).
\end{proof}

\subsection{Application~: critères de séparabilité}
\label{0.21.5-fr}

Dans ce numéro et les deux suivants, on ne suppose plus que les anneaux
considérés soient de caractéristique $p>0$.

\begin{env}[21.5.1]
\label{0.21.5.1-fr}
Notons d'abord que le critère (21.2.7) permet de démontrer une partie du
théorème de Cohen sur les extensions séparables (19.6.1), savoir que si $k$
est de caractéristique $p>0$ et si $K$ est une extension séparable de $k$,
alors $K$ est une $k$-algèbre formellement lisse. En effet, $K$ admet une
$p$-base sur $k$ (21.4.2), et d'autre part, il résulte du critère de MacLane
(Bourbaki, \emph{Alg.}, chap.~V, §~8, n$^\circ$~2, prop.~3) que, dans une
clôture algébrique de $K$, $k^{p^{-1}}$ et $K$ sont linéairement disjoints sur
$k$, et par suite que l'homomorphisme canonique
$k^{p^{-1}}\otimes_k K\to k^{p^{-1}}(K)$ est bijectif, ce qui est précisément
la condition (i) de (21.2.7), après transport de structure par
l'isomorphisme $x\mapsto x^p$.
\end{env}
\begin{proposition}[21.5.2]
\label{0.21.5.2-fr}
\textnormal{(MacLane).}
Soient $A$, $B$ deux anneaux de valuation discrète complets, $\mathfrak m$,
$\mathfrak n$ leurs idéaux maximaux respectifs, $K=A/\mathfrak m$ et
$L=B/\mathfrak n$ leurs corps résiduels, $u:A\to B$ un homomorphisme tel que
$Bu(\mathfrak m)=\mathfrak n$, $u_0=u\otimes 1:K\to L$ l'homomorphisme
correspondant pour les corps résiduels. On considère les conditions suivantes~:
\begin{enumerate}
  \item[\rm a)] $L$ est une extension séparable de $K$ (pour $u_0$).
  \item[\rm b)] Pour tout anneau de valuation discrète complet $B'$ d'idéal
  maximal $\mathfrak n'$ et de corps résiduel $L'$,
\oldpage[0\textsubscript{IV-1}]{165}
  tout homomorphisme
  $u':A\to B'$ tel que $B'u'(\mathfrak m)=\mathfrak n'$ et tout
  $K$-isomorphisme $\sigma:L\to L'$ (relatif à $u_0$ et
  $u'_0=u'\otimes1:K\to L'$), il existe un isomorphisme $w:B\to B'$ tel que
  $u'=w\circ u$ et que $w_0=w\otimes1:L\to L'$ soit égal à $\sigma$.
  \item[\rm b$'$)] Pour tout homomorphisme $u':A\to B$ tel que
  $Bu'(\mathfrak m)=\mathfrak n$ et tel que $u'_0=u'\otimes1:K\to L$ soit égal
  à $u_0$, il existe un automorphisme $w$ de $B$ tel que $u'=w\circ u$ et que
  $w_0=w\otimes1:L\to L$ soit l'identité.
  \item[\rm c)] Désignant par $u_1:A/\mathfrak m^2\to B/\mathfrak n^2$
  l'homomorphisme déduit de $u$ par passage aux quotients, alors, pour tout
  homomorphisme local $u'_1:A/\mathfrak m^2\to B/\mathfrak n^2$ tel que
  $\operatorname{gr}_0(u'_1)=\operatorname{gr}_0(u_1)(=u_0)$, il existe un
  automorphisme $w_1$ de $B/\mathfrak n^2$ tel que
  $\operatorname{gr}_0(w_1)$ soit l'identité et que $u'_1=w_1\circ u_1$.
  \item[\rm c$'$)] Pour tout homomorphisme local $u'_1:A/\mathfrak m^2\to
  B/\mathfrak n^2$ tel que $\operatorname{gr}_0(u'_1)=\operatorname{gr}_0(u_1)$
  et $\operatorname{gr}_1(u'_1)=\operatorname{gr}_1(u_1)$
  (homomorphisme de $\mathfrak m/\mathfrak m^2$ dans
  $\mathfrak n/\mathfrak n^2$), il existe un automorphisme $w_1$ de
  $B/\mathfrak n^2$ tel que $\operatorname{gr}_0(w_1)$ soit l'identité et que
  $u'_1=w_1\circ u_1$.
\end{enumerate}
 Alors on a les implications $c)\Leftrightarrow c')\Leftrightarrow a)\Rightarrow b)
 \Rightarrow b')$. Si de plus $A$ est un anneau de Cohen, les cinq conditions
précédentes sont équivalentes.

\end{proposition}
\begin{proof}
Les implications $b)\Rightarrow b')$ et $c)\Rightarrow c')$ sont triviales.
Montrons que $a)$ implique $b)$.
L'homomorphisme $u$ fait de $B$ un $A$-module sans torsion, puisqu'il
transforme par hypothèse une uniformisante de $A$ en une uniformisante de $B$~;
il en résulte que $B$ est un $A$-module plat (0$_{\rm I}$, 6.3.4), donc une
$A$-algèbre de Cohen (19.8.1)~; le fait que $a)$ implique $b)$ est alors
 conséquence de (19.8.2, (i)) appliquée à $G=B'$, $\mathfrak J=\mathfrak n'$,
 $B'$ étant considéré comme $A$-algèbre pour $u'$ et l'homomorphisme
$B\to B'/\mathfrak n'$ étant le composé
$B\to B/\mathfrak n\to B'/\mathfrak n'$, qui est un $A$-homomorphisme en
 vertu de l'hypothèse $u'_0=\sigma\circ u_0$. Le même raisonnement prouve que
$a)$ entraîne $c)$, en prenant cette fois $G=B/\mathfrak n^2$,
$\mathfrak J=\mathfrak n/\mathfrak n^2$, $G$ étant considéré comme
$A$-algèbre pour $u'_1$.
Prouvons en second lieu que $c')$ implique $a)$. Les deux homomorphismes
$u_1$ et $u'_1$ sont tels que $u'_1=u_1+D$, où $D$ est une dérivation de
$A/\mathfrak m^2$ dans $\mathfrak n/\mathfrak n^2$ (20.1.1)~; d'ailleurs,
l'hypothèse $\operatorname{gr}_1(u'_1)=\operatorname{gr}_1(u_1)$ signifie
que $D$ est nulle dans $\mathfrak m/\mathfrak m^2$, et peut par suite être
considérée comme une dérivation de $K=A/\mathfrak m$ dans
$\mathfrak n/\mathfrak n^2$, et $\mathfrak n/\mathfrak n^2$ s'identifie
(par choix d'une uniformisante de $B$) au $K$-module $L$ (pour $u_0$).
D'autre part, les conditions imposées à $w_1$ entraînent que
 $\operatorname{gr}_1(w_1)$ est aussi l'identité (puisque
 $u_1(\mathfrak m/\mathfrak m^2)$ engendre
$\mathfrak n/\mathfrak n^2$)~; $w_1$ est donc de la forme
$x\mapsto x+D'(x)$, où $D'$ est cette fois une dérivation de
$B/\mathfrak n^2$ dans le $(B/\mathfrak n^2)$-module
$\mathfrak n/\mathfrak n^2$~; comme $w_1$ est l'identité dans
$\mathfrak n/\mathfrak n^2$, $D'$ peut encore (par l'identification
précédente de $\mathfrak n/\mathfrak n^2$ et de $L$) être considérée comme
une dérivation de $L$ dans $L$~; enfin, la relation $u'_1=w_1\circ u_1$
signifie que $D'$ prolonge $D$. Notons d'autre part que toute dérivation $D$
de $K$ dans $L$ correspond à un homomorphisme $u'_1$ vérifiant les conditions
de $c')$ (20.1.1)~; la condition $c')$ signifie donc que toute dérivation de
$K$ dans $L$ se prolonge en une dérivation de $L$ dans lui-même, c'est-à-dire
(20.6.5) que $L$ est séparable sur $K$.
Prouvons enfin que, lorsque $A$ est un anneau de Cohen, $b')$ entraîne $c')$.
Prouvons d'abord que sous les hypothèses de $c')$, il existe un homomorphisme
$u':A\to B$ qui vérifie les hypothèses de $b')$ et qui, par passage aux quotients,
donne $u'_1:A/\mathfrak m^2\to B/\mathfrak n^2$. En effet, cela résulte de
(19.8.6, (i)) appliqué à l'homomorphisme composé
$v':A\to A/\mathfrak m^2\xrightarrow{u'_1}B/\mathfrak n^2$~;
\oldpage[0\textsubscript{IV-1}]{166}
ce dernier se factorise donc en $A\xrightarrow{u'}B\to B/\mathfrak n^2$ et
les hypothèses sur $\operatorname{gr}_0(u'_1)$ et
$\operatorname{gr}_1(u'_1)$ entraînent
$\operatorname{gr}_0(u')=\operatorname{gr}_0(u)=u_0$ et
$\operatorname{gr}_1(u')=\operatorname{gr}_1(u)$, donc l'image par $u'$ d'une
uniformisante de $A$ est une uniformisante de $B$, et l'on a par suite
$Bu'(\mathfrak m)=\mathfrak n$. Il suffit alors, pour obtenir un automorphisme
$w_1$ répondant à la question, de prendre l'automorphisme déduit par passage
aux quotients de l'automorphisme $w$ de $B$ fourni par l'application de $b')$
à l'homomorphisme $u'$.
\end{proof}
\begin{remarks}[21.5.3]
\label{0.21.5.3-fr}
\begin{enumerate}
  \item[(i)] Les propriétés différentielles des corps permettent de résoudre la
  question de l'unicité du corps de représentants dans un anneau local
  noethérien complet $C$ (19.8.7, (ii)). Soient $\mathfrak J$ l'idéal maximal
  de $C$, $K=C/\mathfrak J$ le corps résiduel de $C$~; on peut se borner au
  cas $\mathfrak J\ne0$. Supposons qu'il existe un homomorphisme $u:K\to C$
  qui, composé avec l'augmentation $C\to K$, donne l'identité~; alors, pour
  que cet homomorphisme soit unique, il faut et il suffit que
  $\Omega^1_K=0$. En effet, la condition $\Omega^1_K=0$ entraîne que $K$ est
  formellement non ramifié sur son corps premier (20.7.4)~; s'il existait un
  second homomorphisme $v\ne u$ répondant à la question, il y aurait un plus
  grand entier $n$ tel que $\mathfrak J^n$ contienne l'ensemble des
  $u(x)-v(x)$ pour $x\in K$~; par passage au quotient, $u$ et $v$ donneraient
  deux homomorphismes distincts $u_1,v_1$ de $K$ dans
  $C/\mathfrak J^{n+1}$, dont les composés avec
  $C/\mathfrak J^{n+1}\to C/\mathfrak J^n$ seraient égaux, ce qui contredit
  la définition (19.10.2). Inversement, supposons que $\Omega^1_K\ne0$~; il
  existe alors une dérivation $D\ne0$ de $K$ dans
  $\mathfrak J/\mathfrak J^2$ (20.4.8), donc un homomorphisme
  $v_1:K\to C/\mathfrak J^2$ tel que, si $u_1:K\to C/\mathfrak J^2$ est
  obtenu par passage au quotient à partir de $u$, on ait $v_1=u_1+D$
  (20.1.1). Si $k$ est le corps premier de $K$, $C$ est une $k$-algèbre et
  $K$ une $k$-algèbre formellement lisse (19.6.1), et les restrictions de
  $u_1$ et $v_1$ à $k$ coïncident, donc $v_1$ se factorise en
  $K\xrightarrow{v}C\to C/\mathfrak J^2$ et l'on a $v\ne u$.
  Rappelons ((20.6.20) et (21.4.4)) que la condition $\Omega^1_K=0$ signifie
  que $K$ est parfait s'il est de caractéristique $\ne0$, et extension
  algébrique de $\mathbf Q$ s'il est de caractéristique $0$.
  \item[(ii)] De la même manière, soient $W$ un anneau de Cohen, $C$ un anneau
  local noethérien complet, soit $\mathfrak J\ne0$ un idéal de $C$ contenu dans
  l'idéal maximal~; pour que la factorisation $W\to C\to C/\mathfrak J$
  dans (19.8.6, (i)) soit unique, il faut et il suffit que le corps résiduel
  $K$ de l'anneau de Cohen $W$ soit tel que $\Omega^1_K=0$. En effet, si
  $\Omega^1_K\ne0$ et si $\mathfrak m$ désigne l'idéal maximal de $W$, il
  suffit de composer une dérivation $D\ne0$ de $K$ dans
  $\mathfrak J/\mathfrak m\mathfrak J$ avec l'augmentation $W\to K$ pour
  obtenir une dérivation $D'\ne0$ de $W$ dans
  $\mathfrak J/\mathfrak m\mathfrak J$ et l'on termine le raisonnement comme
  dans (i) en formant à l'aide de $D'$ un homomorphisme
  $v_1:W\to C/\mathfrak m\mathfrak J$ distinct de l'homomorphisme
  $u_1:W\to C/\mathfrak m\mathfrak J$ déduit de $u$ par passage au quotient~;
  on achève le raisonnement en invoquant cette fois (19.8.6, (i)). Si au
  contraire $\Omega^1_K=0$, l'unicité de $v$ résulte déjà de (i) lorsque
  $W=K$ est un corps de caractéristique $0$. Dans le cas contraire, on a
  $\widehat{\Omega}^1_W=0$~; en effet, on a alors $\mathfrak m=pW$, $p$ étant
  la caractéristique de $K$ (19.8.5), et l'homomorphisme canonique
  $\mathfrak m/\mathfrak m^2\to\Omega^1_W/\mathfrak m\Omega^1_W$ (20.5.11.2)
  est par suite nul. La suite exacte (20.5.12.1) appliquée à $W$ et à
  $K=W/\mathfrak m$ entraîne alors que $\Omega^1_W=\mathfrak m\Omega^1_W$,
  d'où notre assertion. Mais alors (20.7.4) $W$ est formellement non ramifié
  (pour sa topologie adique) sur $\mathbf Z$, et l'unicité de $v$ se prouve
  comme dans (i).
\end{enumerate}
\end{remarks}

\oldpage[0\textsubscript{IV-1}]{167}
\subsection{Corps admissibles pour une extension}
\label{section:0.21.6-fr}

\begin{env}[21.6.1]
\label{0.21.6.1-fr}
Étant donnés quatre corps $k_0\subset k\subset K\subset L$, il résulte de
(20.6.16) et (20.6.17) que l'on a une suite exacte
\begin{equation}
\tag{21.6.1.1}
\label{0.21.6.1.1-fr}
0\to\Upsilon_{K/k/k_0}\otimes_KL\xrightarrow{v}\Upsilon_{L/k/k_0}
\xrightarrow{u}\Upsilon_{L/K/k_0}\xrightarrow{s}\Upsilon_{L/K/k}\to0.
\end{equation}

Lorsqu'on laisse fixes $k_0$, $K$ et $L$ et qu'on fait «varier» le corps
intermédiaire $k$ entre $k_0$ et $K$, on a évidemment
$\Upsilon_{L/K/k_0}=\Upsilon_{L/K/k}$ lorsque $k=k_0$. Lorsque
l'homomorphisme canonique $s$ de (21.6.1.1) est encore bijectif, on dit que
$k$ est un corps $k_0$-admissible pour l'extension $L$ de $K$. L'intérêt de
l'existence, sous certaines conditions, de tels corps $k$, qui soient pourtant
«suffisamment proches» de $K$ (par exemple tels que $[K:k]$ soit fini), est
qu'ils permettent de remplacer dans certaines questions les modules de
différentielles $\Omega^1_{K/k_0}$ et $\Omega^1_{L/k_0}$ (qui peuvent être
«trop grands», par exemple lorsque $k_0$ est le corps premier) par
$\Omega^1_{K/k}$ et $\Omega^1_{L/k}$, plus aisément maniables.

Lorsque $k_0$ est le \emph{corps premier}, on dira «corps admissible» au lieu
de «corps $k_0$-admissible».
 
On posera
\end{env}

\begin{env}[21.6.1.2]
\label{0.21.6.1.2-fr}
\[
\Delta(L/K,k/k_0)=\operatorname{Coker}
(\Upsilon_{K/k/k_0}\otimes_KL\to\Upsilon_{L/k/k_0})
\simeq\operatorname{Ker}(\Upsilon_{L/K/k_0}\to\Upsilon_{L/K/k})
\]
(espace vectoriel sur $L$)~; son rang sera noté $d(L/K,k/k_0)$ et appelé
\emph{défaut de $k_0$-admissibilité de $k$ pour l'extension $L$ de $K$}
(il est évidemment nul si et seulement si $k$ est $k_0$-admissible pour cette
extension). Lorsque $k_0$ est le corps premier, on écrira
$\Delta(L/K,k)$ et $d(L/K,k)$ au lieu de $\Delta(L/K,k/k_0)$ et
$d(L/K,k/k_0)$.
\end{env}

\begin{proposition}[21.6.2]
\label{0.21.6.2-fr}
Soient $k_0\subset k\subset K\subset L$ quatre corps.
\begin{enumerate}
  \item[\rm (i)] Les conditions suivantes sont équivalentes~:
  \begin{enumerate}
    \item[\rm a)] Le corps $k$ est $k_0$-admissible pour l'extension $L$ de $K$
    (autrement dit, l'homomorphisme $\Upsilon_{L/K/k_0}\to
    \Upsilon_{L/K/k}$ est injectif, donc bijectif).
    \item[\rm b)] L'homomorphisme canonique
    $u:\Upsilon_{K/k/k_0}\to\Upsilon_{L/k/k_0}$ est nul.
    \item[\rm c)] L'homomorphisme canonique
    $v:\Upsilon_{K/k/k_0}\otimes_KL\to\Upsilon_{L/k/k_0}$ est surjectif
    (donc bijectif).
    \item[\rm d)] On a $d(L/K,k/k_0)=0$ (ou
    $\Delta(L/K,k/k_0)=0$).
  \end{enumerate}
  \item[\rm (ii)] Les conditions équivalentes de (i) sont vérifiées lorsque
  l'on est dans l'un des cas suivants~: $\alpha)$ $L$ est séparable sur $k$~;
  $\beta)$ $L$ est séparable sur $K$~; $\gamma)$ on a
  $k\subset k_0(K^p)$, en désignant par $p$ l'exposant caractéristique de $k_0$.
\end{enumerate}
\end{proposition}

\begin{proof}
\begin{enumerate}
  \item[(i)] Les assertions résultent trivialement de l'exactitude de la suite
  (21.6.1.1).
  \item[(ii)] Si $L$ est séparable sur $k$, on a
  $\Upsilon_{L/K/k_0}=0$ (20.6.19), donc la condition b) de (i) est remplie~;
  si $L$ est séparable sur $K$, on a
  $\Upsilon_{L/K/k_0}=0$ (20.6.19), donc la condition a) de (i) est remplie~;
  enfin, si l'on a $k\subset k_0(K^p)$, il en résulte que
  $\Upsilon_{L/K/k_0}=\Upsilon_{L/K/k}$ en vertu de (21.1.5.1), et la
  condition a) de (i) est vérifiée.
\end{enumerate}
\end{proof}
\oldpage[0\textsubscript{IV-1}]{168}
\begin{env}[21.6.3]
\label{0.21.6.3-fr}
Supposons que l'on ait un diagramme commutatif de monomorphismes de corps
\[
\xymatrix{
 k'_0\ar[r]&k'\ar[r]&K'\ar[r]&L'\\
 k_0\ar[u]\ar[r]&k\ar[u]\ar[r]&K\ar[u]\ar[r]&L\ar[u]
}
\]
Il résulte alors de (20.6.17, (ii)) que l'on a un homomorphisme canonique
\begin{equation}
\tag{21.6.3.1}
\label{0.21.6.3.1-fr}
\Delta(L/K,k/k_0)\longrightarrow\Delta(L'/K',k'/k'_0)
\end{equation}
avec une propriété de transitivité évidente, de sorte que l'on peut dire que
$\Delta(L/K,k/k_0)$ est un \emph{foncteur} en le quadruplet
$(k_0,k,K,L)$.
\end{env}

\begin{proposition}[21.6.4]
\label{0.21.6.4-fr}
\begin{enumerate}
  \item[\rm (i)] Soient $k\subset k'\subset k''\subset K\subset L$ cinq corps.
  On a une suite exacte d'homomorphismes canoniques
\begin{equation}
\tag{21.6.4.1}
\label{0.21.6.4.1-fr}
0\to\Delta(L/K,k'/k)\to\Delta(L/K,k''/k)\to\Delta(L/K,k''/k')\to0
\end{equation}
et par suite l'égalité
\begin{equation}
\tag{21.6.4.2}
\label{0.21.6.4.2-fr}
d(L/K,k''/k)=d(L/K,k''/k')+d(L/K,k'/k).
\end{equation}
  \item[\rm (ii)] Soient $k_0\subset k\subset K\subset L\subset M$ cinq corps.
  On a une suite exacte d'homomorphismes canoniques
\begin{equation}
\tag{21.6.4.3}
\label{0.21.6.4.3-fr}
0\to\Delta(L/K,k/k_0)\otimes_LM\to\Delta(M/K,k/k_0)
\to\Delta(M/L,k/k_0)\to0
\end{equation}
et par suite l'égalité
\begin{equation}
\tag{21.6.4.4}
\label{0.21.6.4.4-fr}
d(M/K,k/k_0)=d(M/L,k/k_0)+d(L/K,k/k_0).
\end{equation}
\end{enumerate}
\end{proposition}

\begin{proof}
\begin{enumerate}
  \item[(i)] Considérons le diagramme commutatif
\[
\xymatrix{
0\ar[r]&\Upsilon_{K/k'/k}\otimes_KL\ar[r]\ar[d]&
\Upsilon_{K/k''/k}\otimes_KL\ar[r]\ar[d]&
\Upsilon_{K/k''/k'}\otimes_KL\ar[r]\ar[d]&0\\
0\ar[r]&\Upsilon_{L/k'/k}\ar[r]\ar[d]&
\Upsilon_{L/k''/k}\ar[r]\ar[d]&
\Upsilon_{L/k''/k'}\ar[r]\ar[d]&0\\
0\ar[r]&\Delta(L/K,k'/k)\ar[r]&\Delta(L/K,k''/k)\ar[r]&
\Delta(L/K,k''/k')\ar[r]&0
}
\]
où l'on peut considérer les trois lignes comme des \emph{complexes}
$T_1,T_2,T_3$ respectivement~; la suite exacte (21.6.1.1) et la définition
de $\Delta$ (21.6.1.2) montrent que l'on a une suite exacte de complexes
$0\to T_1\to T_2\to T_3\to0$~; appliquons-lui la suite exacte de
\oldpage[0\textsubscript{IV-1}]{169}
cohomologie, et notons que, en vertu de l'exactitude de (21.6.1.1), la
cohomologie de $T_1$ et celle de $T_2$ sont nulles sauf en un seul et même
degré, pour lequel les modules de cohomologie sont tous deux égaux à
$\Upsilon_{k''/k'/k}\otimes_{k''}L$~; comme $T_1$ et $T_2$ ont ainsi même
cohomologie, celle de $T_3$ est nécessairement nulle, ce qui prouve (i).
  \item[(ii)] Considérons de même le diagramme commutatif
\[
\xymatrix{
0\ar[r]&\Delta(L/K,k/k_0)\otimes_LM\ar[r]\ar[d]&
\Delta(M/K,k/k_0)\ar[r]\ar[d]&\Delta(M/L,k/k_0)\ar[r]&0\\
0\ar[r]&\Upsilon_{L/K/k_0}\otimes_LM\ar[r]\ar[d]&
\Upsilon_{M/K/k_0}\ar[r]\ar[d]&\Upsilon_{M/L/k_0}\ar[r]&0\\
0\ar[r]&\Upsilon_{L/K/k}\otimes_LM\ar[r]&
\Upsilon_{M/K/k}\ar[r]&\Upsilon_{M/L/k}\ar[r]&0
}
\]
où de nouveau on considère les trois lignes comme des \emph{complexes}
$T'_1,T'_2,T'_3$~; la suite exacte (21.6.1.1) et la définition de $\Delta$
(21.6.1.2) donnent ici une suite exacte de complexes
$0\to T'_1\to T'_2\to T'_3\to0$ à laquelle on applique encore la suite exacte
de cohomologie~; cette fois, en vertu de l'exactitude de (21.6.1.1), la
cohomologie de $T'_2$ et celle de $T'_3$ sont nulles sauf en un seul et même
degré, pour lequel les modules de cohomologie sont tous deux égaux à
$\Upsilon_{M/L/k_0}$~; on conclut donc ici que la cohomologie de $T'_1$ est
nulle, ce qui établit (ii).
\end{enumerate}
\end{proof}

\begin{corollary}[21.6.5]
\label{0.21.6.5-fr}
\begin{enumerate}
  \item[(i)] Étant donnés cinq corps $k\subset k'\subset k''\subset K\subset L$,
  pour que $k''$ soit $k$-admissible pour l'extension $L$ de $K$, il faut et
  il suffit que $k'$ soit $k$-admissible et que $k''$ soit $k'$-admissible
  pour l'extension $L$ de $K$.
  \item[(ii)] Étant donnés cinq corps
  $k_0\subset k\subset K\subset L\subset M$, pour que $k$ soit
  $k_0$-admissible pour l'extension $M$ de $K$, il faut et il suffit qu'il le
  soit pour l'extension $L$ de $K$ et pour l'extension $M$ de $L$.
\end{enumerate}

Cela résulte aussitôt des relations (21.6.4.2) et (21.6.4.4), les valeurs de
$d$ étant $\geqslant0$.
\end{corollary}

\begin{corollary}[21.6.6]
\label{0.21.6.6-fr}
Soient $k_0\subset k\subset K\subset L$ quatre corps, et supposons que $k$
soit $k_0$-admissible pour l'extension $L$ de $K$. Alors, si $k'_0$, $k'$,
$K'$, $L'$ sont quatre corps tels que
$k_0\subset k'_0\subset k'\subset k\subset K\subset K'\subset L'\subset L$,
$k'$ est $k'_0$-admissible pour l'extension $L'$ de $K'$.
\end{corollary}

\subsection{L'égalité de Cartier}
\label{subsection:0.21.7-fr}

Le résultat suivant traduit en termes de différentielles un théorème de
MacLane sur les dérivations~:

\oldpage[0\textsubscript{IV-1}]{170}
\begin{theorem}[21.7.1]
\label{0.21.7.1-fr}
\textnormal{(Cartier).}
Soient $K$ un corps, $L$ une extension de type fini de $K$. Alors
$\Omega^1_{L/K}$ et $\Upsilon_{L/K}$ sont des $L$-espaces vectoriels de rang
fini, et l'on a
\begin{equation}
\tag{21.7.1.1}
\label{0.21.7.1.1-fr}
\operatorname{rg}\Omega^1_{L/K}-\operatorname{rg}\Upsilon_{L/K}
=\deg.\operatorname{tr}_K L.
\end{equation}
\end{theorem}

\begin{proof}
Si $L'$ est un corps tel que $K\subset L'\subset L$, on a la suite exacte
(20.6.15.1) (appliquée à $\Lambda=P$, corps premier de $K$, $A=K$,
$B=L'$, $C=L$)
\[
0\to\Upsilon_{L'/K}\otimes_{L'}L\to\Upsilon_{L/K}
\to\Omega^1_{L'/K}\otimes_{L'}L\to\Omega^1_{L/K}
\to\Omega^1_{L/L'}\to0
\]
d'où (0$_{\rm III}$, 11.10.2)
\[
\operatorname{rg}_L\Omega^1_{L/K}-\operatorname{rg}_L\Upsilon_{L/K}
=(\operatorname{rg}_{L'}\Omega^1_{L/L'}-\operatorname{rg}_{L'}\Upsilon_{L/L'})
+(\operatorname{rg}_{L'}\Omega^1_{L'/K}-\operatorname{rg}_{L'}\Upsilon_{L'/K}).
\]
Comme par ailleurs $\deg.\operatorname{tr}_K L=\deg.\operatorname{tr}_K L'
+\deg.\operatorname{tr}_{L'}L$, on voit, par récurrence sur le nombre de
générateurs de l'extension $L$, que l'on est ramené à prouver (21.7.1.1)
lorsque $L=K(x)$. Distinguons alors trois cas~:
\begin{enumerate}
  \item[\rm a)] $x$ est transcendant sur $K$~; comme $L$ est séparable sur
  $K$, on a alors $\Upsilon_{L/K}=0$ (20.6.3)~; d'autre part, (20.5.9) et
  (20.4.13) montrent que $\Omega^1_{L/K}$ est de rang $1$, d'où
  (21.7.1.1) dans ce cas.
  \item[\rm b)] $L$ est une extension algébrique séparable de $K$, de sorte
  que l'on a encore $\Upsilon_{L/K}=0$ (20.6.3). D'autre part, on a
  $\Omega^1_{L/K}=0$ par (20.6.20), d'où encore (21.7.1.1).
  \item[\rm c)] $L$ est une extension algébrique inséparable de $K$~; le
  raisonnement du début montre que l'on peut se borner au cas où $L^p\subset K$~;
  il résulte donc de (21.4.8) que l'on a
  $\operatorname{rg}_L\Upsilon_{L/K}=\operatorname{rg}_L\Omega^1_{L/K}$,
  d'où encore (21.7.1.1).
\end{enumerate}
\end{proof}

\begin{corollary}[21.7.2]
\label{0.21.7.2-fr}
Soient $K$ un corps, $L$ une extension de type fini de $K$, $k$ un sous-corps
de $K$. Alors $\Omega^1_{L/K}$ et $\Upsilon_{L/K/k}$ sont des espaces
vectoriels de rang fini sur $L$, et l'on a
\begin{equation}
\tag{21.7.2.1}
\label{0.21.7.2.1-fr}
\operatorname{rg}\Omega^1_{L/K}-\operatorname{rg}\Upsilon_{L/K/k}
=\deg.\operatorname{tr}_K L+d(L/K,k).
\end{equation}
Par suite, on a l'inégalité
\begin{equation}
\tag{21.7.2.2}
\label{0.21.7.2.2-fr}
\operatorname{rg}\Omega^1_{L/K}-\operatorname{rg}\Upsilon_{L/K/k}
\geqslant\deg.\operatorname{tr}_K L.
\end{equation}
En outre, pour que les deux membres de (21.7.2.2) soient égaux, il faut et
il suffit que $k$ soit un corps admissible pour l'extension $L$ de $K$.
\end{corollary}

\begin{proof}
En effet, comme l'homomorphisme
$\Upsilon_{L/K}\to\Upsilon_{L/K/k}$ est surjectif, on a par définition
(21.6.1.2)
$\operatorname{rg}\Upsilon_{L/K}-\operatorname{rg}\Upsilon_{L/K/k}
=d(L/K,k)$, et le corollaire résulte donc aussitôt de (21.7.1) et de
(21.6.2).
\end{proof}

\begin{corollary}[21.7.3]
\label{0.21.7.3-fr}
Soient $k\subset K\subset L$ trois corps tels que $L$ soit une extension de
type fini de $k$. Alors on a
\begin{equation}
\tag{21.7.3.1}
\label{0.21.7.3.1-fr}
\operatorname{rg}_L\Omega^1_{L/k}-\operatorname{rg}_K\Omega^1_{K/k}
=\deg.\operatorname{tr}_K L+d(L/K,k).
\end{equation}
\oldpage[0\textsubscript{IV-1}]{171}
Par suite, on a l'inégalité
\begin{equation}
\tag{21.7.3.2}
\label{0.21.7.3.2-fr}
\operatorname{rg}_L\Omega^1_{L/k}-\operatorname{rg}_K\Omega^1_{K/k}
\geqslant\deg.\operatorname{tr}_K L.
\end{equation}
L'égalité étant atteinte si et seulement si $k$ est un corps admissible pour
l'extension $L$ de $K$.
\end{corollary}

\begin{proof}
On a en effet la suite exacte (20.6.1.1)
\[
0\to\Upsilon_{L/K/k}\to\Omega^1_{K/k}\otimes_KL\to\Omega^1_{L/k}
\to\Omega^1_{L/K}\to0
\]
donc le premier membre de (21.7.3.1) est égal à
\[
\operatorname{rg}_L\Omega^1_{L/k}-\operatorname{rg}_L\Upsilon_{L/K/k}
\]
et il suffit d'appliquer (21.7.2).
\end{proof}

L'intérêt de ce dernier corollaire est qu'il ne fait intervenir que des
modules de différentielles, à l'exclusion de modules d'imperfection. Nous
verrons d'ailleurs plus bas (21.8.6) que pour toute extension de type fini
$L$ de $K$, il existe un sous-corps $k$ de $K$ tel que $[K:k]$ soit fini et
qui est admissible pour l'extension $L$ de $K$, de sorte que les deux membres
de (21.7.3.2) sont alors égaux.

\begin{corollary}[21.7.4]
\label{0.21.7.4-fr}
Soit $K$ une extension de type fini d'un corps $k$. Les conditions suivantes
sont équivalentes~:
\begin{enumerate}
  \item[\rm a)] $K$ est une extension finie séparable de $k$.
  \item[\rm b)] $K$ est une $k$-algèbre formellement non ramifiée (19.10.2).
  \item[\rm c)] $K$ est une $k$-algèbre formellement étale (19.10.2).
  \item[\rm d)] On a $\Omega^1_{K/k}=0$.
\end{enumerate}
\end{corollary}

\begin{proof}
En effet, on sait (les topologies étant discrètes) que les conditions b) et
d) sont équivalentes (20.7.4) et que a) entraîne que $K$ est une
$k$-algèbre formellement lisse (19.6.1), donc la conjonction de a) et b)
équivaut à c). Tout revient à voir que d) entraîne a). Or, en vertu de
(21.7.1.1), la relation $\Omega^1_{K/k}=0$ entraîne que $K$ est algébrique
(donc fini) sur $k$ et que $\Upsilon_{K/k}=0$, c'est-à-dire (20.6.3) que
$K$ est séparable sur $k$.
\end{proof}

\begin{remark}[21.7.5]
\label{0.21.7.5-fr}
En vertu de (20.6.3.2) et de (0$_{\rm III}$, 11.10.2), le premier membre de
(21.7.1.1) n'est autre que la \emph{caractéristique d'Euler-Poincaré} du
complexe $K_\bullet(C/B/A)$ introduit dans (20.6.3), pour $C=L$, $B=K$ et
$A=P$ (corps premier de $K$). Dans le chapitre consacré à la théorie des
intersections et au théorème de Riemann--Roch, un rôle important sera joué
également par les caractéristiques d'Euler-Poincaré généralisées (à valeurs
dans des groupes de classes de $\mathcal O_X$-Modules) de complexes
généralisant les complexes $F_\bullet(C/A)$ considérés dans (20.6.22).
\end{remark}

\subsection{Critères d'admissibilité}
\label{subsection:0.21.8-fr}

Nous revenons à nos conventions antérieures et supposons donc que tous les
corps considérés dans ce numéro sont de caractéristique $p>0$.

\begin{lemma}[21.8.1]
\label{0.21.8.1-fr}
Soient $K$ un corps, $k$ un sous-corps de $K$, $(k_\alpha)_{\alpha\in I}$ une
famille filtrante décroissante de sous-corps de $K$ telle que
$k=\bigcap_{\alpha\in I}k_\alpha$. Soient $V$ un espace vectoriel sur $K$,
$(a_i)_{1\leqslant i\leqslant n}$ une famille finie de vecteurs de $V$~; si la
famille $(a_i)$ est libre sur $k$, il existe un indice $\gamma$ tel qu'elle
soit aussi libre sur $k_\gamma$.
\end{lemma}

\oldpage[0\textsubscript{IV-1}]{172}
\begin{proof}
Soit $r$ le rang de la famille $(a_i)_{1\leqslant i\leqslant n}$ sur $K$, et
raisonnons par récurrence sur $n-r$~; la proposition est évidente pour $n=r$,
car alors la famille $(a_i)_{1\leqslant i\leqslant r}$ est libre sur $K$, donc
sur tout sous-corps de $K$. Supposons par exemple que
$(a_i)_{1\leqslant i\leqslant r}$ soit libre sur $K$, et écrivons
$a_{r+1}=\sum_{i=1}^r\lambda_i a_i$ avec $\lambda_i\in K$~; la famille
$(a_i)_{1\leqslant i\leqslant r+1}$ étant libre sur $k$, les $\lambda_i$ ne
peuvent tous appartenir à $k$~; supposons par exemple que
$\lambda_{i_0}\notin k$. Alors il existe un indice $\beta$ tel que
$\lambda_{i_0}\notin k_\beta$~; on en conclut que la famille $(a_i)_{1\leqslant i\leqslant r+1}$ est libre sur $k_\beta$~; en effet, comme la
famille $(a_i)_{1\leqslant i\leqslant r}$ est libre sur tout sous-corps de
$K$, si la famille $(a_i)_{1\leqslant i\leqslant r+1}$ n'était pas libre sur
$k_\beta$, $a_{r+1}$ serait égal à une combinaison linéaire des $a_i$ à
coefficients dans $k_\beta$, et comme ces coefficients sont nécessairement
les $\lambda_i$, on aboutit à une contradiction. Il suffit maintenant
d'appliquer l'hypothèse de récurrence en remplaçant $K$ par $k_\beta$ et la
famille $(k_\alpha)_{\alpha\in I}$ par la sous-famille des $k_\alpha$
contenus dans $k_\beta$.
\end{proof}

\begin{lemma}[21.8.2]
\label{0.21.8.2-fr}
Soient $K$ un corps, $p$ son exposant caractéristique, $k_0$ un sous-corps de
$K$, $(k_\alpha)_{\alpha\in I}$ une famille filtrante décroissante de
sous-corps de $K$ telle que $\bigcap_{\alpha\in I}k_\alpha(K^p)=k_0(K^p)$.
Si $(x_i)_{1\leqslant i\leqslant n}$ est une famille finie d'éléments de $K$
qui est $p$-libre sur $k_0$, il existe un indice $\alpha$ tel que $(x_i)$ soit
$p$-libre sur $k_\alpha$.

En effet, dire que la famille $(x_i)$ est $p$-libre sur un sous-corps $k$
de $K$ signifie que la famille finie des monômes
$\prod_{i=0}^n x_i^{m(i)}$ avec $0\leqslant m(i)<p$ est libre sur $k(K^p)$~;
il suffit donc d'appliquer le lemme (21.8.1) à cette famille de monômes dans
l'espace vectoriel $V=K$, et aux sous-corps $k_\alpha(K^p)$ et $k_0(K^p)$
de $K$.
\end{lemma}

\begin{theorem}[21.8.3]
\label{0.21.8.3-fr}
Soient $K$ un corps de caractéristique $p>0$, $k_0$ un sous-corps de $K$,
$(k_\alpha)_{\alpha\in I}$ une famille filtrante décroissante de sous-corps
de $K$ contenant $k_0$. Les conditions suivantes sont équivalentes~:
\begin{enumerate}
  \item[\rm a)] $\bigcap_{\alpha\in I}k_\alpha(K^p)=k_0(K^p)$.
  \item[\rm b)] Pour toute extension $L$ de $K$ telle que
  $\Upsilon_{L/K/k_\alpha}$ soit un $L$-espace vectoriel de rang fini
  (ce qui a lieu en particulier si $L$ est une extension de type fini en
  vertu de (21.7.2)), il existe $\alpha\in I$ tel que $k_\alpha$ soit
  $k_0$-admissible pour l'extension $L$ de $K$.
  \item[\rm b$'$)] Pour toute extension $L=K(x)$ de $K$, avec $x^p\in K$,
  il existe $\alpha\in I$ tel que $k_\alpha$ soit $k_0$-admissible pour
  l'extension $L$ de $K$.
  \item[\rm c)] L'application canonique
\begin{equation}
\tag{21.8.3.1}
\label{0.21.8.3.1-fr}
\Omega^1_{K/k_0}\longrightarrow\varprojlim_{\alpha}\Omega^1_{K/k_\alpha}
\end{equation}
est injective.
\end{enumerate}
\end{theorem}

\begin{proof}
L'application canonique (21.8.3.1) est bien entendu obtenue par passage à la
limite projective dans le système projectif d'homomorphismes
$\Omega^1_{K/k_0}\to\Omega^1_{K/k_\alpha}$ (20.5.3.3). Nous allons prouver
le théorème suivant le schéma logique $c)\Rightarrow b)\Rightarrow b')
\Rightarrow a)\Rightarrow c)$.

Dire que $k_\alpha$ est $k_0$-admissible pour l'extension $L$ de $K$ signifie
que l'homomorphisme
\oldpage[0\textsubscript{IV-1}]{173}
canonique $\Upsilon_{L/K/k_0}\to\Upsilon_{L/K/k_\alpha}$ est injectif~; or,
si $N_\alpha$ est le noyau de cet homomorphisme, les $N_\alpha$ forment un
système filtrant décroissant de sous-espaces vectoriels de
$\Upsilon_{L/K/k_0}$, et comme $\Upsilon_{L/K/k_0}$ est par hypothèse de rang
fini, il revient au même de dire que l'un des $N_\alpha$ est nul ou que leur
intersection est nulle. Mais cette intersection n'est autre que le noyau de
l'homomorphisme limite projective
$\Upsilon_{L/K/k_0}\to\varprojlim_\alpha\Upsilon_{L/K/k_\alpha}$. Or, on a
le diagramme commutatif
\[
\xymatrix{
\Upsilon_{L/K/k_0}\ar[r]\ar[d]&\varprojlim\Upsilon_{L/K/k_\alpha}\ar[d]\\
\Omega^1_{K/k_0}\otimes_KL\ar[r]&
\varprojlim(\Omega^1_{K/k_\alpha}\otimes_KL)
}
\]
où la flèche verticale de gauche est injective par définition. Pour prouver
que $c)$ entraîne $b)$, il suffit donc de montrer que l'application canonique
\begin{equation}
\tag{21.8.3.2}
\label{0.21.8.3.2-fr}
\Omega^1_{K/k_0}\otimes_KV\longrightarrow
\varprojlim_{\alpha}(\Omega^1_{K/k_\alpha}\otimes_KV)
\end{equation}
est injective pour tout espace vectoriel $V$ sur $K$ (et en particulier pour
$V=L$). Or cela est évident si $V=K^n$ puisque alors
$W\otimes_KV=W^n$ pour tout espace vectoriel $W$ sur $K$ et que les produits
et limites projectives commutent. D'autre part, pour tout élément $z$ du
premier membre de (21.8.3.2), il existe un sous-espace $V'$ de $V$ de rang
fini tel que $z\in\Omega^1_{K/k_0}\otimes_KV'$ (identifié canoniquement à un
sous-espace de $\Omega^1_{K/k_0}\otimes_KV$)~; si $z\ne0$, son image dans
$\varprojlim_\alpha(\Omega^1_{K/k_\alpha}\otimes_KV')$ n'est donc pas nulle~;
comme le foncteur $\varprojlim$ est exact à gauche, l'image de $z$ dans le
second membre de (21.8.3.2) est donc aussi $\ne0$, ce qui achève de montrer
que $c)$ implique $b)$.

Il est trivial que $b)$ entraîne $b')$~; montrons que $b')$ entraîne $a)$.
Avec les notations de $b')$, on peut supposer que $L\ne K$. Il résulte alors
de (21.4.7) que l'on a $\operatorname{rg}\Upsilon_{L/K}\leqslant1$. Si
$\Upsilon_{L/K/k_\alpha}=0$, il n'y a rien à démontrer en vertu de
(21.6.1.1). Sinon, $\Upsilon_{L/K/k_\alpha}$ s'identifie canoniquement à
$\Upsilon_{L/K}$ et, si l'on pose $a=x^p$, $\Upsilon_{L/K}$ a une base
formée du seul élément $d_K(a)\otimes1$ (21.4.7)~;
$\Upsilon_{L/K/k_\alpha}$ a donc une base formée du seul élément
$d_{K/k_\alpha}(a)\otimes1$, et dire qu'un sous-corps $k_\alpha$ est
$k_0$-admissible pour l'extension $L$ de $K$ signifie que l'on a
$d_{K/k_\alpha}(a)\ne0$, ou encore (21.4.5) que
$a\notin k_\alpha(K^p)$. Or, pour tout $a\notin k_0(K^p)$, on a
$d_{K/k_0}(a)\ne0$ et $K(a^{1/p})\ne K$, donc on peut appliquer $b')$, qui
prouve l'existence d'un $\alpha$ tel que $a\notin k_\alpha(K^p)$~;
autrement dit $b')$ implique $a)$.

Reste à montrer que $a)$ implique $c)$. Soit $(x_\mu)_{\mu\in M}$ une
$p$-base de $K$ sur $k_0$~; alors, les $d_{K/k_0}(x_\mu)$ forment une base du
$K$-espace vectoriel $\Omega^1_{K/k_0}$ (21.4.5), et la condition $c)$
signifie que pour toute partie finie $J$ de l'ensemble d'indices $M$, il
existe un $\alpha\in I$ tel que les $d_{K/k_\alpha}(x_\mu)$ pour $\mu\in J$
soient linéairement indépendants dans $\Omega^1_{K/k_\alpha}$~; mais cela
signifie aussi (21.4.5) que les $x_\mu$ pour $\mu\in J$ forment une famille
$p$-libre sur $k_\alpha$, et l'existence d'un $\alpha$ ayant cette propriété
découle de (21.8.2) et de l'hypothèse $a)$.
\end{proof}
\oldpage[0\textsubscript{IV-1}]{174}

\begin{corollary}[21.8.4]
\label{0.21.8.4-fr}
Soient $K$ un corps de caractéristique $p>0$, $k_0$ un sous-corps de $K$,
$(k_\alpha)_{\alpha\in I}$ une famille filtrante décroissante de sous-corps
de $K$, contenant $k_0$, telle que
\[
  \bigcap_{\alpha\in I}k_\alpha(K^p)=k_0(K^p).
\]
Soient $L$ une extension de $K$ telle que $\Upsilon_{L/K/k_0}$ soit un
$L$-espace vectoriel de rang fini~; alors, pour tout corps $K'$ tel que
$K\subset K'\subset L$, il existe $\alpha\in I$ tel que $k_\alpha$ soit
$k_0$-admissible pour l'extension $L$ de $K'$.
\end{corollary}

\begin{proof}
Il suffit d'appliquer (21.8.3) et (21.6.6).
\end{proof}

\begin{corollary}[21.8.5]
\label{0.21.8.5-fr}
Soient $K$ un corps de caractéristique $p>0$, $k_0$ un sous-corps de $K$,
$(k_\alpha)_{\alpha\in I}$ une famille filtrante décroissante de sous-corps
de $K$ contenant $k_0$, telle que
\[
  \bigcap_{\alpha\in I}k_\alpha(K^p)=k_0(K^p).
\]
Alors, pour toute extension $L$ de $K$ telle que $\Upsilon_{L/K/k_0}$ soit
un $L$-espace vectoriel de rang fini, on a
\[
  \bigcap_{\alpha\in I}k_\alpha(L^p)=k_0(L^p).
\]
\end{corollary}

\begin{proof}
Supposons en effet que $M$ soit une extension de $L$ telle que
$\Upsilon_{M/L/k_0}$ soit un $M$-espace vectoriel de rang fini. La suite
exacte (21.6.1.1)
\[
  0\longrightarrow \Upsilon_{L/K/k_0}\otimes_L M
  \longrightarrow \Upsilon_{M/K/k_0}
  \longrightarrow \Upsilon_{M/L/k_0}
\]
et l'hypothèse montrent alors que $\Upsilon_{M/K/k_0}$ est aussi un
$M$-espace vectoriel de rang fini. Il existe donc un indice $\alpha$ tel que
$k_\alpha$ soit $k_0$-admissible pour l'extension $M$ de $K$ (21.8.3), donc
aussi pour l'extension $M$ de $L$ (21.6.6)~; cela ayant lieu pour toute
extension $M$ de $L$ telle que $\Upsilon_{M/L/k_0}$ soit de rang fini, le
corollaire résulte de l'équivalence de $a)$ et $b)$ dans (21.8.3).
\end{proof}

\begin{corollary}[21.8.6]
\label{0.21.8.6-fr}
Soient $K$ un corps, $p$ son exposant caractéristique, $k_0$ un sous-corps de
$K$. Si $L$ est une extension de $K$ telle que $\Upsilon_{L/K/k_0}$ soit un
$L$-espace vectoriel de rang fini, il existe un sous-corps $k$ de $K$,
contenant $k_0(K^p)$, tel que $[K:k]$ soit fini, et qui soit
$k_0$-admissible pour l'extension $L$ de $K$.
\end{corollary}

\begin{proof}
Il suffit en effet, en vertu de (21.8.3), de construire une famille filtrante
décroissante $(k_\alpha)_{\alpha\in I}$ de sous-corps de $K$, contenant
$K^p$ et $k_0$, pour lesquels $[K:k_\alpha]<+\infty$ et
$\bigcap_\alpha k_\alpha=k_0(K^p)$. Pour cela on considère une $p$-base
$(x_\lambda)_{\lambda\in J}$ de $K$ sur $k_0$ et, pour toute partie finie
$H$ de $J$, on considère le sous-corps $k_H$ de $K$ engendré par $k_0(K^p)$
et les $x_\lambda$ d'indice $\lambda\in J-H$~; il résulte de cette définition
que $(x_\lambda)_{\lambda\in H}$ est une $p$-base de $K$ sur $k_H$, et on en
conclut aussitôt que les $k_H$ vérifient les conditions voulues.
\end{proof}

\begin{remarks}[21.8.7]
\label{0.21.8.7-fr}
\begin{enumerate}
  \item[(i)] On a déjà vu (21.7.2) que si $L$ est une extension de
  \emph{type fini} de $K$, $\Upsilon_{L/K/k_0}$ est de rang fini pour tout
  sous-corps $k_0$ de $K$. Il en est de même si $L$ est une extension
  \emph{séparable} de $K$, car en vertu de (20.6.19), on a
  $\Upsilon_{L/K/k_0}=0$. Enfin, si $L$ est une \emph{extension de type fini
  d'une extension séparable $L_0$ de $K$}, le même raisonnement que dans
  (21.8.5) montre que $\Upsilon_{L/K/k_0}$ est encore de rang fini (et en fait
  est isomorphe à un sous-espace de $\Upsilon_{L/L_0/k_0}$).

  \oldpage[0\textsubscript{IV-1}]{175}
  \item[(ii)] Dans l'énoncé de (21.8.5), et par suite aussi dans celui de
  (21.8.3, $b)$), on ne peut omettre l'hypothèse que $\Upsilon_{L/K/k_0}$ est
  de rang fini sur $L$. Prenons pour $k_0$ le corps premier $\mathbf F_p$,
  pour $K$ un corps tel que $[K:K^p]$ soit infini dénombrable (par exemple le
  corps de fractions rationnelles $\mathbf F_p(X_1,\ldots,X_n,\ldots)$ à une
  infinité d'indéterminées)~; le procédé de construction de (21.8.6) montre
  aussitôt qu'il existe une suite infinie strictement décroissante $(K_n)$ de
  sous-corps de $K$ telle que $K_0=K$ et $\bigcap_nK_n=K^p$. Nous allons
  construire une suite croissante d'extensions finies $M_n$ de $K$, telle que
  si $M$ est la réunion des $M_n$, l'extension $L=M^{1/p}$ de $K$ mette
  (21.8.5) en défaut. Pour cela, soit $x$ un élément de $K-K^p$, posons
  $M_0=K^p$, et $M_n=K^p(a_1x,a_2x,\ldots,a_nx)$ pour $n\geqslant1$, où les
  $a_n$ sont construits par récurrence de façon que $a_n\in K_n$ et
  $a_{n+1}\notin M_n(x)$ pour tout $n$~: cela est possible, car $M_n(x)$ est
  de degré fini sur $K^p$, tandis qu'il n'en est pas de même de $K_n$. On en
  conclut aussitôt que $x\notin M=L^p$, mais comme
  $x=(a_nx)a_n^{-1}$, on a $x\in K_n(M)=K_n(L^p)$ pour tout $n$.
\end{enumerate}
\end{remarks}

\begin{proposition}[21.8.8]
\label{0.21.8.8-fr}
Soit $k$ un corps de caractéristique $p>0$, et soient
$A=k[[T_1,\ldots,T_r]]$ l'anneau des séries formelles à $r$ indéterminées sur
$k$, $K=k((T_1,\ldots,T_r))$ son corps des fractions. Alors il existe une
famille filtrante décroissante $(A_\alpha)$ de sous-anneaux noethériens de $A$
telle que $A$ soit un $A_\alpha$-module libre de type fini pour tout $\alpha$
et que, si $K_\alpha$ est le corps des fractions de $A_\alpha$, on ait
$\bigcap_\alpha K_\alpha=K^p$.
\end{proposition}

\begin{proof}
On peut écrire $K^p=k^p((T_1^p,\ldots,T_r^p))$~; on a vu dans la démonstration
de (21.8.6) qu'il existe une famille décroissante $(k_\alpha)$ de sous-corps
de $k$ telle que $[k:k_\alpha]$ soit fini pour tout $\alpha$ et
$\bigcap_\alpha k_\alpha=k^p$~; il est clair que si l'on pose
$A_\alpha=k_\alpha[[T_1^p,\ldots,T_r^p]]$, $A$ est un $A_\alpha$-module
libre de type fini~; tout revient donc à prouver la relation
\begin{equation}
\tag{21.8.8.1}
\label{0.21.8.8.1-fr}
  \bigcap_\alpha k_\alpha((T_1^p,\ldots,T_r^p))
  =k^p((T_1^p,\ldots,T_r^p)).
\end{equation}
Comme $\bigcap_\alpha k_\alpha=k^p$, cela va résulter des deux lemmes suivants.
\end{proof}

\begin{lemma}[21.8.8.2]
\label{0.21.8.8.2-fr}
Si $k'$ est une extension d'un corps $k$, on a
\[
  k'[[T_1,\ldots,T_r]]\cap k((T_1,\ldots,T_r))=k[[T_1,\ldots,T_r]].
\]
\end{lemma}

\begin{proof}
En effet, posons $C=k[[T_1,\ldots,T_r]]$, $D=k'[[T_1,\ldots,T_r]]$~; comme
$k((T_1,\ldots,T_r))$ est le corps des fractions de $C$, il suffira de
prouver que $D$ est un $C$-module \emph{fidèlement plat} (Bourbaki,
\emph{Alg. comm.}, chap.~I, \S~3, n\textsuperscript{o}~5, prop.~10). Or, $C$
et $D$ sont des anneaux locaux noethériens, et si $\mathfrak m$ est l'idéal
maximal de $C$, on a
$D/\mathfrak m^jD=(C/\mathfrak m^j)\otimes_k k'$, donc
$D/\mathfrak m^jD$ est un $(C/\mathfrak m^j)$-module plat~; il suffit donc
d'appliquer ($\mathbf 0_{\mathrm{III}}$, 10.2.1) et
($\mathbf 0_{\mathrm I}$, 6.6.2).
\end{proof}

\begin{lemma}[21.8.8.3]
\label{0.21.8.8.3-fr}
Soient $k$ un corps, $(k_\alpha)$ une famille filtrante décroissante de
sous-corps de $k$, et posons $k_0=\bigcap_\alpha k_\alpha$. On suppose qu'il
existe une puissance $q$ de l'exposant caractéristique de $k$ telle que
$k^q\subset k_0$. Alors on a
\begin{equation}
\tag{21.8.8.4}
\label{0.21.8.8.4-fr}
  \bigcap_\alpha k_\alpha((T_1,\ldots,T_r))=k_0((T_1,\ldots,T_r)).
\end{equation}
\end{lemma}

\begin{proof}
\oldpage[0\textsubscript{IV-1}]{176}
Il suffit de prouver qu'un élément $f\ne0$ du premier membre de (21.8.8.4)
appartient au second membre. Soient $\gamma$ un indice,
$g\in k_\gamma[[T_1,\ldots,T_r]]$ tel que
$gf\in k_\gamma[[T_1,\ldots,T_r]]$~; quitte à remplacer $g$ par $g^q$, on
peut supposer que $g\in k_0[[T_1,\ldots,T_r]]$. Alors, pour tout
$\alpha\geqslant\gamma$, on a
\[
  gf\in k_\gamma[[T_1,\ldots,T_r]]\cap k_\alpha((T_1,\ldots,T_r))
  =k_\alpha[[T_1,\ldots,T_r]]
\]
en vertu du lemme (21.8.8.2). Mais il est clair que l'intersection des
anneaux $k_\alpha[[T_1,\ldots,T_r]]$ n'est autre que
$k_0[[T_1,\ldots,T_r]]$, et l'on a donc bien
$f\in k_0((T_1,\ldots,T_r))$.
\end{proof}

\subsection{Modules de différentielles complétés dans les anneaux de séries formelles}
\label{subsection:0.21.9-fr}

Dans ce numéro, les corps ne sont plus nécessairement supposés être de
caractéristique $>0$.

\begin{lemma}[21.9.1]
\label{0.21.9.1-fr}
Soient $k$ un corps, $A$ un anneau local noethérien complet qui est une
$k$-algèbre, $K$ le corps résiduel de $A$. Si $K$ est une extension de type
fini de $k$, le $A$-module $\widehat\Omega^1_{A/k}$ est de type fini.
\end{lemma}

\begin{proof}
En vertu de (20.7.15), il suffit de prouver que $\Omega^1_{K/k}$ est un
$K$-espace vectoriel de rang fini. Or, par hypothèse, $K$ est le corps des
fractions d'une $k$-algèbre de type fini $B$~; comme $\Omega^1_{B/k}$ est un
$B$-module de type fini en vertu de (20.4.7), la conclusion résulte de
(20.5.9).
\end{proof}

\begin{proposition}[21.9.2]
\label{0.21.9.2-fr}
Soient $k_0$ un corps, $A$ un anneau local noethérien complet qui est une
$k_0$-algèbre formellement lisse, $\mathfrak m$ son idéal maximal,
$K=A/\mathfrak m$ son corps résiduel (de sorte qu'en tant qu'anneau $A$ est
isomorphe à un anneau de séries formelles $K[[T_1,\ldots,T_r]]$ (19.6.5)).
Si $K$ est une extension de type fini de $k_0$, alors
$\widehat\Omega^1_{A/k_0}$ est un $A$-module libre de rang égal à
$\dim(A)+\deg.\operatorname{tr}_{k_0}K$. En outre, pour tout sous-corps $k$
de $k_0$ tel que $\Omega^1_{k_0/k}$ soit de rang fini,
$\widehat\Omega^1_{A/k}$ est un $A$-module libre de rang égal à
\[
  \dim(A)+\deg.\operatorname{tr}_{k_0}K
  +\operatorname{rg}_{k_0}(\Omega^1_{k_0/k}).
\]
\end{proposition}

\begin{proof}
Il est clair que, si $P$ est le sous-corps premier de $k_0$, $K$ est une
$P$-algèbre formellement lisse (19.6.1). Notons d'autre part que l'on a
$\Upsilon^K_{A/k_0/P}=0$~: en effet, cela revient à voir que l'homomorphisme
\[
  u_{A/k_0/P}\otimes1_K:
  \Omega^1_{k_0/P}\otimes_{k_0}K
  \longrightarrow \Omega^1_{A/P}\otimes_A K
\]
est injectif (20.6.14.5), $u$ désignant l'homomorphisme structural
$k_0\to A$. Mais comme $A$ est une $k_0$-algèbre formellement lisse par
hypothèse, il résulte de (20.7.2) que $u_{A/k_0/P}$ est formellement
inversible à gauche, donc (19.1.7) $u_{A/k_0/P}\otimes1_K$ est inversible à
gauche et \emph{a fortiori} injectif, ce qui prouve notre assertion.
Appliquant la suite exacte (20.6.22.1), où $\Lambda,A,B,C$ sont remplacés par
$P,k_0,A,K$, il vient la suite exacte
\[
  0\longrightarrow \Upsilon^K_{A/k_0}
  \longrightarrow \mathfrak m/\mathfrak m^2
  \longrightarrow \Omega^1_{A/k_0}\otimes_A K
  \longrightarrow \Omega^1_{K/k_0}\longrightarrow0,
\]
de sorte que l'on a
\[
\begin{split}
  \operatorname{rg}_K(\Omega^1_{A/k_0}\otimes_AK)
  &={\operatorname{rg}}_K(\mathfrak m/\mathfrak m^2)
    +\bigl(\operatorname{rg}_K(\Omega^1_{K/k_0})
    -\operatorname{rg}_K(\Upsilon_{K/k_0})\bigr)\\
  &={\operatorname{rg}}_K(\mathfrak m/\mathfrak m^2)
    +\deg.\operatorname{tr}_{k_0}K
\end{split}
\]
en vertu de (21.7.1). D'autre part, $\widehat\Omega^1_{A/k_0}$ est un
$A$-module formellement projectif (20.4.9), et comme sa topologie est la
topologie $\mathfrak m$-préadique (20.4.5),
$\widehat\Omega^1_{A/k_0}/\mathfrak m^{j+1}\widehat\Omega^1_{A/k_0}$ est un
$(A/\mathfrak m^{j+1})$-module projectif pour tout $j$ (19.2.4), donc libre
($\mathbf 0_{\mathrm{III}}$, 10.1.2) et de rang

\oldpage[0\textsubscript{IV-1}]{177}
$n=\operatorname{rg}_K(\Omega^1_{A/k_0}\otimes_AK)$. En vertu de (21.9.1),
$\widehat\Omega^1_{A/k_0}$ est un $A$-module de type fini, et en outre ce
$A$-module est plat en vertu de ($\mathbf 0_{\mathrm{III}}$, 10.2.1)~; il est
donc libre ($\mathbf 0_{\mathrm{III}}$, 10.1.3) de rang $n$, d'où la première
assertion.

Comme par hypothèse $\Omega^1_{k_0/k}$ est de rang fini sur $k_0$, le produit
tensoriel complété
$\Omega^1_{k_0/k}\widehat\otimes_{k_0}A$ s'identifie à
$\Omega^1_{k_0/k}\otimes_{k_0}A$~; comme $A$ est une $k_0$-algèbre
formellement lisse, il résulte de (20.7.18) que l'homomorphisme
\[
  \widehat u_{A/k_0/k}:
  \Omega^1_{k_0/k}\otimes_{k_0}A
  \longrightarrow\widehat\Omega^1_{A/k}
\]
admet un inverse à gauche qui est un $A$-homomorphisme continu, ce qui
implique en particulier que son image est fermée dans
$\widehat\Omega^1_{A/k}$~; la suite exacte
\[
  0\longrightarrow \Omega^1_{k_0/k}\otimes_{k_0}A
  \xrightarrow{\widehat u_{A/k_0/k}}\widehat\Omega^1_{A/k}
  \longrightarrow\widehat\Omega^1_{A/k_0}\longrightarrow0
\]
est donc exacte et scindée en vertu de (20.7.17) et de ce qui précède~; ce
qui achève de prouver la proposition.
\end{proof}

\begin{corollary}[21.9.3]
\label{0.21.9.3-fr}
Soient $k$ un corps, $A$ l'anneau de séries formelles
$k[[T_1,\ldots,T_r]]$, muni de sa structure usuelle de $k$-algèbre. Alors
$\widehat\Omega^1_{A/k}$ est un $A$-module libre de rang $r=\dim(A)$.
\end{corollary}

\begin{proof}
Il suffit de noter que $A$ est une $k$-algèbre formellement lisse (19.3.4),
et d'appliquer (21.9.2) avec $k_0=k=K$.
\end{proof}

\begin{lemma}[21.9.4]
\label{0.21.9.4-fr}
Soient $k$ un corps de caractéristique $p>0$, $A$ une $k$-algèbre qui est un
anneau local noethérien complet, $B$ une sous-$k$-algèbre de $A$ isomorphe à
une $k$-algèbre topologique de la forme $k[[T_1,\ldots,T_r]]$ et telle que
$A$ soit une $B$-algèbre de type fini. Si $B_1$ est la sous-algèbre
$k[[T_1^p,\ldots,T_r^p]]$ de $B$, $\widehat\Omega^1_{A/k}$ s'identifie
canoniquement à $\Omega^1_{A/B_1}$.
\end{lemma}

\begin{proof}
Toute dérivation continue de $B_1$ dans un $A$-module topologique $P$, qui
est restriction d'une $k$-dérivation continue de $A$ dans $P$, est \emph{nulle},
car elle l'est dans le sous-anneau $k[T_1^p,\ldots,T_r^p]$ des polynômes, et
ce dernier est dense dans $B_1$. La suite exacte (20.3.7.1) montre donc que
l'homomorphisme canonique
\[
  \operatorname{D\acute{e}r.\,cont}_{B_1}(A,P)
  \longrightarrow \operatorname{D\acute{e}r.\,cont}_k(A,P)
\]
est bijectif~; la conclusion résulte de (20.7.14.4), compte tenu de ce que
$\Omega^1_{A/B_1}$ est un $A$-module de type fini (20.4.7), donc séparé et
complet puisque $A$ est un anneau local noethérien complet.
\end{proof}

\begin{proposition}[21.9.5]
\label{0.21.9.5-fr}
Soient $A$ un anneau local noethérien complet intègre, $A_0$ un sous-anneau
de $A$ isomorphe à un anneau de séries formelles $k[[T_1,\ldots,T_r]]$ sur
un corps $k$, tel que $A$ soit une $A_0$-algèbre finie et que le corps des
fractions $E$ de $A$ soit une extension séparable du corps des fractions
$L_0$ de $A_0$. Alors on a
\begin{equation}
\tag{21.9.5.1}
\label{0.21.9.5.1-fr}
  \operatorname{rg}_A\widehat\Omega^1_{A/k}
  =\operatorname{rg}_E(\widehat\Omega^1_{A/k}\otimes_AE)=\dim(A).
\end{equation}
\end{proposition}

\begin{proof}
Notons que si $\mathfrak m_0$ est l'idéal maximal de $A_0$, la topologie de
$A$ est la topologie $\mathfrak m_0$-adique puisque $A$ est une
$A_0$-algèbre finie (Bourbaki, \emph{Alg. comm.}, chap.~IV, \S~2,
n\textsuperscript{o}~5, cor.~3 de la prop.~9) et induit sur $A_0$ la topologie
$\mathfrak m_0$-adique (Bourbaki, \emph{Alg. comm.}, chap.~III, \S~3,
n\textsuperscript{o}~4, th.~3). On sait (21.9.1) que
$\widehat\Omega^1_{A/k}$ est un $A$-module de type fini et,

\oldpage[0\textsubscript{IV-1}]{178}
en vertu de (21.9.3) $\widehat\Omega^1_{A_0/k}$ est un $A_0$-module libre de
rang $r=\dim(A_0)=\dim(A)$ (16.1.5)~; le produit tensoriel
$\widehat\Omega^1_{A_0/k}\otimes_{A_0}A$ est donc un $A$-module libre de rang
$r$ identique à son séparé complété~; enfin $\widehat\Omega^1_{A/A_0}$ est
un $A$-module de type fini (20.4.7), donc identique à son séparé complété
($\mathbf 0_{\mathrm I}$, 7.3.6). Cela étant, dans la suite d'homomorphismes
\begin{equation}
\tag{21.9.5.2}
\label{0.21.9.5.2-fr}
  \widehat\Omega^1_{A_0/k}\otimes_{A_0}A
  \xrightarrow{u}\widehat\Omega^1_{A/k}
  \xrightarrow{v}\widehat\Omega^1_{A/A_0}\longrightarrow0
\end{equation}
on sait que $v$ est surjectif et que $\operatorname{Im}(u)$ est dense dans
$\operatorname{Ker}(v)$ (20.7.17)~; mais tout sous-$A_0$-module de
$\widehat\Omega^1_{A/k}$ est fermé pour la topologie $\mathfrak m_0$-adique
($\mathbf 0_{\mathrm I}$, 7.3.5), donc la suite (21.9.5.2) est exacte. Notons
d'autre part que puisque $A$ est entier sur $A_0$, $A_0$ intégralement clos
et $E$ séparable sur $L_0$, $\widehat\Omega^1_{A/A_0}$ est un $A$-module de
torsion (20.4.13, (iv))~; tensorisant la suite exacte (21.9.5.2) par $E$, il
vient une suite exacte
\[
  \widehat\Omega^1_{A_0/k}\otimes_{A_0}E
  \longrightarrow \widehat\Omega^1_{A/k}\otimes_AE\longrightarrow0
\]
d'où $\operatorname{rg}_E(\widehat\Omega^1_{A/k}\otimes_AE)\leqslant r$. Il
reste donc à montrer qu'il existe, dans $\widehat\Omega^1_{A/k}\otimes_AE$,
$r$ éléments linéairement indépendants sur $E$. Or, soient
$D_i=\partial/\partial T_i$ les dérivations canoniques de $A_0$ dans lui-même
$(1\leqslant i\leqslant r)$~; elles se prolongent de façon unique en des
dérivations (encore notées $D_i$) de $E$ dans lui-même, puisque $E$ est
extension séparable finie de $L_0$~; par restriction à $A$, ces dérivations
donnent des dérivations de $A$ dans $E$, et il est immédiat que ces
dérivations prennent leurs valeurs dans un même sous-$A$-module de type fini
de $E$~; comme elles sont continues dans $A_0$ et que la topologie de $A$ est
la topologie $\mathfrak m_0$-adique, on a ainsi formé $r$ dérivations
continues de $A$ dans un $A$-module topologique, qui sont évidemment
linéairement indépendantes puisque $D_i(T_j)=\delta_{ij}$~; ceci achève de
prouver la proposition.
\end{proof}

La proposition suivante est un analogue «~formel~» de (21.7.3)~:

\begin{proposition}[21.9.6]
\label{0.21.9.6-fr}
Soient $k_0\subset k\subset K_0$ trois corps de caractéristique $p>0$, tels
que $[K_0:k]<+\infty$~; on pose
\[
  L_0=K_0((T_1,\ldots,T_r)),\qquad
  M_0=K_0((T_1^p,\ldots,T_r^p)),
\]
\[
  L=k((T_1,\ldots,T_r)),\qquad
  M=k((T_1^p,\ldots,T_r^p)),\qquad
  N=k((T_1^{p^2},\ldots,T_r^{p^2})).
\]
Soit $E$ une extension de type fini de $L_0$.
\begin{enumerate}
  \item[(i)] Si $\Upsilon_{K_0/k_0}$ est de rang fini, on a
  \begin{equation}
  \tag{21.9.6.1}
  \label{0.21.9.6.1-fr}
    \operatorname{rg}_E\Omega^1_{E/M}
    -\operatorname{rg}_{K_0}\Omega^1_{K_0/k}
    =(r+\deg.\operatorname{tr}_{L_0}E)
      +d(E/K_0,k_0)+d(E/M_0,N/k_0)
  \end{equation}
  (notations de (21.6.1)), et en particulier on a
  \begin{equation}
  \tag{21.9.6.2}
  \label{0.21.9.6.2-fr}
    \operatorname{rg}_E\Omega^1_{E/M}
    -\operatorname{rg}_{K_0}\Omega^1_{K_0/k}
    \geqslant r+\deg.\operatorname{tr}_{L_0}E+d(E/K_0,k_0).
  \end{equation}
  En outre, si les deux membres de (21.9.6.2) sont égaux, ils le restent
  lorsqu'on remplace $k$ par un corps $k'$ tel que $k_0\subset k'\subset k$.

  \item[(ii)] On suppose en outre que $[k_0:k_0^p]<+\infty$. Soit
  $(k_\alpha)$ une famille filtrante décroissante de sous-corps de $K_0$
  contenant $k_0$, tels que $[K_0:k_\alpha]<+\infty$ pour tout $\alpha$ et
  que
  $\bigcap_\alpha k_\alpha(K_0^p)=k_0(K_0^p)$~; alors il existe un $\alpha$
  tel que, pour $k=k_\alpha$, les deux membres de (21.9.6.2) soient égaux.
\end{enumerate}
\end{proposition}

\begin{proof}
\begin{enumerate}
  \oldpage[0\textsubscript{IV-1}]{179}
  \item[(i)] On sait (21.1.5, (ii)) que
  $\operatorname{rg}_E\Omega^1_{E/M}$ ne change pas quand on remplace $M$ par
  $M(L_0^p)$~; comme
  $L_0^p=K_0^p((T_1^p,\ldots,T_r^p))$ et que
  $[k(K_0^p):k]<+\infty$, on a
  \[
    M(L_0^p)=k(K_0^p)((T_1^p,\ldots,T_r^p))~;
  \]
  de même $\Omega^1_{K_0/k}$ ne change pas quand on y remplace $k$ par
  $k(K_0^p)$, si bien que l'on peut \emph{supposer que $K_0^p\subset k$}, ce
  que nous ferons dans toute la suite. Nous introduirons aussi les corps
  \[
    k'_0=k_0(K_0^p),\qquad
    N'=k'_0((T_1^{p^2},\ldots,T_r^{p^2}))
  \]
  de sorte que l'on a le diagramme de corps
  \[
    \xymatrix{
      K_0\ar[r] & M_0\ar[r] & L_0\ar[r] & E\\
      k\ar[u]\ar[r] & N\ar[u]\ar[r] & M\ar[u]\ar[r] & L\ar[u]\\
      k_0\ar[u]\ar[r] & k'_0\ar[u]\ar[r] & N'\ar[u] &
    }
  \]

  Nous nous proposons d'évaluer la différence $\delta$ du premier et du
  second membre de (21.9.6.2).

  Il résulte de (21.7.3) que l'on a
  \[
    \operatorname{rg}_E\Omega^1_{E/M}
    -\operatorname{rg}_{M_0}\Omega^1_{M_0/M}
    =\deg.\operatorname{tr}_{M_0}E+d(E/M_0,M)
  \]
  et comme $L_0$ est une extension finie de $M_0$, on a
  $\deg.\operatorname{tr}_{L_0}E=\deg.\operatorname{tr}_{M_0}E$. D'autre
  part, une $p$-base de $K_0$ sur $k=k(K_0^p)$ est aussi une $p$-base de
  $M_0$ sur $M$, donc (21.4.5), on a
  \[
    \operatorname{rg}_{M_0}\Omega^1_{M_0/M}
    =\operatorname{rg}_{K_0}\Omega^1_{K_0/k}
  \]
  si bien que l'on a
  \begin{equation}
  \tag{21.9.6.3}
  \label{0.21.9.6.3-fr}
    \delta=d(E/M_0,M)-d(E/K_0,k_0)-r.
  \end{equation}

  Notons maintenant le lemme classique~:

  \begin{lemma}[21.9.6.4]
  \label{0.21.9.6.4-fr}
  Pour tout corps $k$, le corps de séries formelles
  $K=k((T_1,\ldots,T_r))$ est une extension séparable de $k$.
  \end{lemma}

  \begin{proof}
  Rappelons rapidement la démonstration de ce lemme pour être complet. Il
  suffit (Bourbaki, \emph{Alg.}, chap.~VIII, \S~7,
  n\textsuperscript{o}~3, démonstration du th.~1) de prouver que pour toute
  extension finie $k'$ de $k$, $K\otimes_k k'$ est sans élément nilpotent~;
  mais si l'on pose $A=k[[T_1,\ldots,T_r]]$ et $S=A-\{0\}$,
  $K\otimes_k k'$ est égal à $S^{-1}(A\otimes_k k')$, et
  $A\otimes_k k'$ s'identifie canoniquement à l'anneau intègre
  $A'=k'[[T_1,\ldots,T_r]]$ et $A$ à un sous-anneau de $A'$, d'où la
  conclusion.
  \end{proof}

  Utilisant ce lemme et (21.6.2, (ii)), on a $d(L_0/K_0,k_0)=0$, et la
  formule (21.6.4.4) donne donc
  \[
    d(E/K_0,k_0)=d(E/L_0,k_0).
  \]

  \oldpage[0\textsubscript{IV-1}]{180}
  Les formules (21.6.4.4) et (21.6.4.2) donnent d'autre part
  \[
    d(E/M_0,M)=d(E/L_0,M)+d(L_0/M_0,M),
  \]
  \[
    d(L_0/M_0,M)=d(L_0/M_0,M/N)+d(L_0/M_0,N),
  \]
  d'où
  \[
    \delta=d(E/L_0,M)+d(L_0/M_0,N)-d(E/L_0,k_0)
    +d(L_0/M_0,M/N)-r.
  \]

  Nous allons montrer que l'on a la relation
  \begin{equation}
  \tag{21.9.6.5}
  \label{0.21.9.6.5-fr}
    d(L_0/M_0,M/N)=r.
  \end{equation}
  Pour cela, remarquons que, par définition, on a
  \[
    d(L_0/M_0,M/N)
    =\operatorname{rg}_{L_0}\Upsilon_{L_0/M/N}
    -\operatorname{rg}_{M_0}\Upsilon_{M_0/M/N}.
  \]
  Compte tenu des deux suites exactes
  \[
  \begin{aligned}
    0&\longrightarrow\Upsilon_{L_0/M/N}
      \longrightarrow\Omega^1_{M/N}\otimes_ML_0
      \longrightarrow\Omega^1_{L_0/N}
      \longrightarrow\Omega^1_{L_0/M}\longrightarrow0,\\
    0&\longrightarrow\Upsilon_{M_0/M/N}
      \longrightarrow\Omega^1_{M/N}\otimes_MM_0
      \longrightarrow\Omega^1_{M_0/N}
      \longrightarrow\Omega^1_{M_0/M}\longrightarrow0,
  \end{aligned}
  \]
  il vient
  \[
  \begin{split}
    d(L_0/M_0,M/N)={}&\operatorname{rg}_{L_0}\Omega^1_{L_0/M}
      -\operatorname{rg}_{L_0}\Omega^1_{L_0/N}\\
      &-\operatorname{rg}_{M_0}\Omega^1_{M_0/M}
      +\operatorname{rg}_{M_0}\Omega^1_{M_0/N}.
  \end{split}
  \]
  Or, on a $N(L_0^p)\supset M$, donc
  $\Omega^1_{L_0/N}=\Omega^1_{L_0/M}$ (21.1.5, (ii)). D'autre part, les
  $T_i^p$ $(1\leqslant i\leqslant r)$ forment une $p$-base de $M$ sur $N$,
  et une $p$-base de $K_0$ sur $k$ est aussi une $p$-base de $M_0$ sur $M$,
  donc, comme $M_0^p\subset N$,
  \[
    \operatorname{rg}_{M_0}\Omega^1_{M_0/N}
    =r+\operatorname{rg}_{K_0}\Omega^1_{K_0/k}
    =r+\operatorname{rg}_{M_0}\Omega^1_{M_0/M}
  \]
  en vertu de (21.4.5) et (21.1.10), ce qui prouve (21.9.6.5).

  On a encore, par (21.6.4.2),
  \[
    d(E/L_0,M)=d(E/L_0,k_0)+d(E/L_0,M/k_0)
  \]
  et
  \[
    d(L_0/M_0,N)=d(L_0/M_0,N/k)+d(L_0/M_0,k).
  \]
  Mais comme $L_0$ est séparable sur $K_0$ (21.9.6.4), on a
  $d(L_0/K_0,k)=0$, et par suite aussi $d(L_0/M_0,k)=0$ ((21.6.2) et
  (21.6.4)). Il vient donc
  \[
    \delta=d(E/L_0,M/k_0)+d(L_0/M_0,N/k).
  \]
  Appliquons encore deux fois (21.6.4.2) au premier terme du second membre de
  cette formule~; on obtient
  \[
    d(E/L_0,M/k_0)=d(E/L_0,M/N)+d(E/L_0,N/k)+d(E/L_0,k/k_0).
  \]
  Comme $M\subset N(L_0^p)$, on a $d(E/L_0,M/N)=0$ (21.6.2), donc (par
  (21.6.4.4))
  \[
    \delta=d(E/L_0,k/k_0)+d(E/M_0,N/k).
  \]
  Mais $M_0$ et $L_0$ sont des extensions séparables de $K_0$ (21.9.6.4),
  donc
  \[
    d(L_0/K_0,k/k_0)=d(M_0/K_0,k/k_0)=0
  \]
  (21.6.2), et appliquant encore deux fois (21.6.4.4),

  \oldpage[0\textsubscript{IV-1}]{181}
  on a
  \[
    d(E/L_0,k/k_0)=d(E/K_0,k/k_0)=d(E/M_0,k/k_0),
  \]
  d'où, par une dernière application de (21.6.4.2),
  \[
    \delta=d(E/M_0,N/k_0),
  \]
  ce qui prouve (21.9.6.1). En outre, lorsque $k$ est remplacé par une
  sous-extension $k'$ de $k_0$, $N$ est remplacé par une sous-extension $N'$,
  donc $d(E/M_0,N'/k_0)\leqslant d(E/M_0,N/k_0)$ (21.6.4.2), ce qui démontre
  la dernière assertion de (i).

  \item[(ii)] Pour tout $\alpha$, posons
  $N_\alpha=k_\alpha((T_1^{p^2},\ldots,T_r^{p^2}))$. En vertu de (21.9.6.1)
  et de (21.8.3), il s'agit de montrer (compte tenu de ce que
  $K_0^p\subset k_\alpha$ et de ce que $E$ est une extension de type fini de
  $M_0$) que l'on a
  \[
    \bigcap_\alpha N_\alpha=k_0(M_0^p).
  \]
  Or, on a
  \begin{equation}
  \tag{21.9.6.6}
  \label{0.21.9.6.6-fr}
    \bigcap_\alpha N_\alpha=N'
    =(k_0(K_0^p))((T_1^{p^2},\ldots,T_r^{p^2}))
  \end{equation}
  en vertu de (21.8.8.3) et de la relation
  $\bigcap_\alpha k_\alpha=k_0(K_0^p)$ (21.8.6). D'autre part, on a
  $k_0(M_0^p)=k_0(K_0^p((T_1^{p^2},\ldots,T_r^{p^2})))$. Pour achever la
  démonstration de (21.9.6, (ii)), il reste donc à prouver le

  \begin{lemma}[21.9.6.7]
  \label{0.21.9.6.7-fr}
  Soient $k_0$ un corps de caractéristique $p>0$ tel que
  $[k_0:k_0^p]<+\infty$, $K_0$ une extension de $k_0$. Alors on a
  \begin{equation}
  \tag{21.9.6.8}
  \label{0.21.9.6.8-fr}
    (k_0(K_0^p))((T_1,\ldots,T_r))
    =k_0\bigl(K_0^p((T_1,\ldots,T_r))\bigr).
  \end{equation}
  \end{lemma}

  \begin{proof}
  Comme $k_0^p\subset K_0^p$, on a
  \[
    k_0^p\bigl(K_0^p((T_1,\ldots,T_r))\bigr)
    =K_0^p((T_1,\ldots,T_r)).
  \]
  Si $(c_i)_{1\leqslant i\leqslant n}$ est une base de $k_0$ sur $k_0^p$,
  il est immédiat que $(c_i)$ est aussi un système de générateurs de chacun
  des deux membres de (21.9.6.8) sur $K_0^p((T_1,\ldots,T_r))$.
  \end{proof}
\end{enumerate}
\end{proof}

\begin{remark}[21.9.7]
\label{0.21.9.7-fr}
L'assertion de (21.9.6, (ii)) \emph{ne s'étend pas} au cas où
$[k_0:k_0^p]$ est infini. Prenons en effet pour $k_0$ le corps
$\mathbf F_p(X_n)_{n\geqslant0}$ des fractions rationnelles à une infinité
d'indéterminées, de sorte que les $X_n$ forment une $p$-base de $k_0$ sur
$\mathbf F_p$, et par suite $[k_0:k_0^p]=+\infty$. Dans l'énoncé de (21.9.6),
prenons $K_0=k_0$ (donc nécessairement $k=k_0$), $L_0=k_0((T))$, pour $E$
l'extension $L_0(y)$, où $y$ est racine du polynôme
\[
  Y^p-\sum_{n=0}^{\infty}X_nT^{np}
\]
de $L_0[Y]$. Alors la différence $\delta$ des deux membres de (21.9.6.2) est
\emph{non nulle}. En effet, $E$ est \emph{séparable} sur $k_0$, sans quoi
$y^p$ serait élément de $k_0(L_0^p)$, et l'on voit aussitôt que cela n'est
pas possible (dû au fait qu'il y a une infinité de $X_n$ algébriquement
indépendants)~; on a donc $d(E/K_0,k_0)=0$, la formule (21.9.6.3) se réduit à
$\delta=d(E/M_0,M_0)-1$, et il est clair que
$d(E/M_0,M_0)=\operatorname{rg}_E\Upsilon_{E/M_0}$. Tout revient à vérifier
que ce rang est $2$. Or on a $E^p\subset M_0$, une $p$-base de $E$ sur $L_0$
est formée du seul élément $y$, et $L_0$ a évidemment une $p$-base sur $M_0$
formée du seul élément $T$~; donc, puisque $E$ est extension algébrique de
$M_0$, notre assertion résulte de (21.7.1) et (21.4.5).
\end{remark}

\begin{proposition}[21.9.8]
\label{0.21.9.8-fr}
Soient $k_0$ un corps (de caractéristique quelconque), $K_0$ une extension
séparable de $k_0$, $A$ un anneau local noethérien complet, qui est une
$k_0$-algèbre et dont le corps résiduel

\oldpage[0\textsubscript{IV-1}]{182}
est une extension finie de $K_0$. Pour tout idéal premier $\mathfrak p$ de
$A$, soit $k(\mathfrak p)$ le corps résiduel de $A$ en $\mathfrak p$. Alors,
pour tout corps $k$ tel que $k_0\subset k\subset K_0$ et tel que
$[K_0:k]<+\infty$, on a
\begin{equation}
\tag{21.9.8.1}
\label{0.21.9.8.1-fr}
\begin{split}
  &\operatorname{rg}_{k(\mathfrak p)}
    \bigl(\widehat\Omega^1_{(A/\mathfrak p)/k}
      \otimes_{A/\mathfrak p}k(\mathfrak p)\bigr)
    -\operatorname{rg}_{K_0}\Omega^1_{K_0/k}\\
  &\hspace{35mm}\geqslant
    \dim(A/\mathfrak p)
    +\operatorname{rg}_{k(\mathfrak p)}\Upsilon_{k(\mathfrak p)/k_0}.
\end{split}
\end{equation}
Supposons en outre que $[k_0:k_0^p]<+\infty$ (où $p$ est l'exposant
caractéristique de $k_0$). Alors on peut trouver un corps $k$ tel que
$k_0(K_0^p)\subset k\subset K_0$ et $[K_0:k]<+\infty$, pour lequel les deux
membres de (21.9.8.1) sont égaux.
\end{proposition}

\begin{proof}
Comme $A/\mathfrak p$ a même corps résiduel $K$ que $A$ et est complet, on
peut se borner au cas où $A$ est intègre et $\mathfrak p=0$~; on posera
alors $E=k(\mathfrak p)$, corps des fractions de $A$. Comme $K_0$ est
formellement lisse sur $k_0$ (19.6.1), il existe un $k_0$-monomorphisme
$K_0\to A$ faisant de $A$ une $K_0$-algèbre~; on sait en outre qu'il existe
une sous-$K_0$-algèbre $A_0$ de $A$, $K_0$-isomorphe à une algèbre de séries
formelles $K_0[[T_1,\ldots,T_r]]$ et telle que $A$ soit une $A_0$-algèbre
\emph{finie} (19.8.9), ce qui entraîne que $E$ est une extension finie du
corps des fractions $L_0=K_0((T_1,\ldots,T_r))$ de $A_0$. On a alors
$\Upsilon_{K_0/k_0}=0$ puisque $K_0$ est séparable sur $k_0$, et
$\deg.\operatorname{tr}_{L_0}E=0$~; en outre, comme $A$ est une
$A_0$-algèbre finie on a $\dim(A)=\dim(A_0)$ (16.1.5).

Si $p=1$, on a
$\operatorname{rg}_E(\widehat\Omega^1_{A/k}\otimes_AE)=r$ par (21.9.5) et
(17.1.4, (iii)), et $\Omega^1_{K_0/k}=0$ puisque $K_0$ est extension finie
séparable de $k$ (21.7.4)~; d'autre part $\Upsilon_{E/k_0}=0$, $E$ étant
extension séparable de $k_0$~; les deux membres de (21.9.8.1) sont donc égaux
dans ce cas.

Supposons maintenant que $p>1$~; si $B=k[[T_1,\ldots,T_r]]$, $A$ est une
$B$-algèbre finie, donc, en posant $B_1=k[[T_1^p,\ldots,T_r^p]]$,
$\widehat\Omega^1_{A/k}$ s'identifie à $\Omega^1_{A/B_1}$ (21.9.4), et en
désignant par $M$ le corps des fractions $k((T_1^p,\ldots,T_r^p))$, il
résulte de (20.5.9) que $\widehat\Omega^1_{A/k}\otimes_AE$ s'identifie à
$\Omega^1_{E/M}$~; l'inégalité (21.9.8.1) n'est autre alors que (21.9.6.2),
et la dernière assertion du corollaire résulte de (21.9.6, (ii)).
\end{proof}

% FIN SOURCE ega4/ega0-21-fr.tex

%
% DEBUT SOURCE ega4/ega0-22-fr.tex
\section{Critères différentiels de lissité formelle et de régularité}
\label{section:0.22-fr}

Les principaux résultats de ce paragraphe sont~:
\begin{enumerate}
  \item[a)] Les critères (22.2.2) et (22.5.8), qui donnent des conditions
  nécessaires et suffisantes pour qu'un anneau local noethérien complet $A$
  contenant un corps $k$ soit formellement lisse sur $k$~; le critère
  (22.5.8) s'énonce sans faire intervenir de notions différentielles, et
  complète (19.6.6)~; l'énoncé (22.2.2), de nature un peu plus technique,
  permet de donner une \emph{classification}, pour $k$ fixé, des anneaux
  locaux $A$ considérés (22.2.5)~: ils sont donnés par leur dimension $n$ et
  leur corps résiduel $K$, sous la seule condition
  $n\geqslant\operatorname{rg}_K\Upsilon_{K/k}$.

  \item[b)] Le théorème (22.3.3), et le corollaire (22.7.6) du critère
  jacobien de Nagata (22.7.3), qui donnent des conditions permettant
  d'affirmer qu'un anneau \emph{localisé} d'un anneau local noethérien complet
  $A$ contenant un corps $k$ est géométriquement régulier sur $k$~; ils
  joueront un rôle important dans l'étude des propriétés «~fines~» des
  anneaux locaux faite au chap.~IV, §~7. On ne possède pour le moment aucune
  démonstration de (22.7.6) indépendante du critère jacobien de Nagata.

  \item[c)] Le critère jacobien de Zariski (22.6.7) et ses variantes, qui
  sont des conséquences faciles de (20.7.8).
\end{enumerate}

\oldpage[0\textsubscript{IV-1}]{183}

\subsection{Relèvement de la lissité formelle}
\label{subsection:0.22.1-fr}

\begin{theorem}[22.1.1]
\label{0.22.1.1-fr}
Soient $s:\Lambda\to A$, $u:A\to B$ deux homomorphismes d'anneaux,
$\mathfrak J$ un idéal de $B$, $C$ l'anneau quotient $B/\mathfrak J$. On
suppose que $C$ est une $\Lambda$-algèbre formellement lisse pour la topologie
discrète.
\begin{enumerate}
  \item[(i)] Supposons vérifiées les conditions suivantes~:
  \begin{enumerate}
    \item[$\alpha$)] $\mathfrak J/\mathfrak J^2$ est un $C$-module projectif.
    \item[$\beta$)] L'homomorphisme canonique
    $\mathbf S_C^*(\mathfrak J/\mathfrak J^2)\to
    \operatorname{gr}_{\mathfrak J}^*(B)$ (19.5.2) est bijectif.
    \item[$\gamma$)] L'homomorphisme caractéristique
    $\chi:\Upsilon_{C/B/\Lambda}\to\mathfrak J/\mathfrak J^2$ (20.6.20)
    est injectif.
    \item[$\delta$)] Le $C$-module
    $\Omega^1_{B/A}\otimes_BC$ est projectif.
  \end{enumerate}
  Alors $B$, muni de la topologie $\mathfrak J$-préadique, est une
  $\Lambda$-algèbre formellement lisse.

  \item[(ii)] Inversement, supposons que $B$ soit une $\Lambda$-algèbre
  formellement lisse pour la topologie $\mathfrak J$-préadique, et que $A$
  soit une $\Lambda$-algèbre formellement lisse pour la topologie discrète.
  Alors les conditions $\alpha$), $\beta$), $\gamma$), $\delta$) de (i) sont
  vérifiées.
\end{enumerate}
\end{theorem}

\begin{proof}
En vertu de la suite exacte (20.6.22.1) (qui est applicable, puisque
$B/\mathfrak J^2$ est une extension $\Lambda$-triviale de $C$ par
$\mathfrak J/\mathfrak J^2$, $C$ étant supposée être une $\Lambda$-algèbre
formellement lisse (19.4.4)), la condition $\gamma$) équivaut à
$\Upsilon^C_{B/A/\Lambda}=0$, ou encore à la suivante~:
\begin{enumerate}
  \item[$\gamma'$)] L'homomorphisme
  \[
    u_{B/A/\Lambda}\otimes 1_C:
    \Omega^1_{A/\Lambda}\otimes_A C
    \longrightarrow \Omega^1_{B/\Lambda}\otimes_B C
  \]
  est injectif.
\end{enumerate}

(i) Les conditions $\alpha$) et $\beta$) entraînent que $B$ est une
$\Lambda$-algèbre formellement lisse pour la topologie
$\mathfrak J$-préadique (19.5.4). Il revient donc au même de dire que $B$ est
(pour la topologie $\mathfrak J$-préadique) une $A$-algèbre formellement
lisse, ou une $A$-algèbre formellement lisse \emph{relativement à $\Lambda$}.
Pour voir que $B$ est une $A$-algèbre formellement lisse, il suffit donc, en
vertu de (20.7.2), de prouver que l'homomorphisme continu de $B$-modules
topologiques
\[
  u_{B/A/\Lambda}:\Omega^1_{A/\Lambda}\otimes_A B
  \longrightarrow\Omega^1_{B/\Lambda}
\]
est \emph{formellement inversible à gauche}. Or, puisque $B$ est une
$\Lambda$-algèbre formellement lisse, le $B$-module topologique
$\Omega^1_{B/\Lambda}$ est formellement projectif (20.4.9) et sa topologie
est déduite de celle de $B$ (20.4.5). En vertu de (19.1.9), il suffit donc,
pour que $u_{B/A/\Lambda}$ soit formellement inversible à gauche, que
$u_{B/A/\Lambda}\otimes 1_C$ soit inversible à gauche. Mais en vertu de
$\gamma'$), cette dernière application est injective, donc on a la suite
exacte (20.6.14.7)
\[
  0\longrightarrow\Omega^1_{A/\Lambda}\otimes_A C
  \longrightarrow\Omega^1_{B/\Lambda}\otimes_B C
  \longrightarrow\Omega^1_{B/A}\otimes_B C\longrightarrow0.
\]
Enfin, en vertu de $\delta$), le $C$-module
$\Omega^1_{B/A}\otimes_BC$ est projectif, donc la suite exacte précédente
est scindée, ce qui achève de prouver (i).

\oldpage[0\textsubscript{IV-1}]{184}

(ii) Si $B$ est une $\Lambda$-algèbre formellement lisse pour la topologie
$\mathfrak J$-préadique, et $A$ une $\Lambda$-algèbre formellement lisse pour
la topologie discrète, $B$ est une $A$-algèbre formellement lisse pour la
topologie $\mathfrak J$-préadique (19.3.5, (ii))~; les conditions $\alpha$) et
$\beta$) résultent donc de (19.5.6). En outre, (20.7.2) montre que
$u_{B/A/\Lambda}$ est formellement inversible à gauche, donc
$u_{B/A/\Lambda}\otimes 1_C$ est inversible à gauche (19.1.7), et \emph{a
fortiori} injectif. Enfin, $\Omega^1_{B/A}$ est un $B$-module formellement
projectif (20.4.9), et par suite $\Omega^1_{B/A}\otimes_BC$ est un $C$-module
projectif (19.2.4).
\end{proof}

\begin{corollary}[22.1.2]
\label{0.22.1.2-fr}
Soient $s:\Lambda\to A$, $u:A\to B$ deux homomorphismes d'anneaux,
$\mathfrak m$ un idéal maximal de $B$, $K=B/\mathfrak m$ le corps quotient.
On suppose que les $\Lambda$-algèbres $A$ et $K$ soient formellement lisses
pour la topologie discrète. Alors les conditions suivantes sont
équivalentes~:
\begin{enumerate}
  \item[a)] $B$ est une $A$-algèbre formellement lisse pour la topologie
  $\mathfrak m$-préadique.
  \item[b)] L'homomorphisme canonique
  \begin{equation}
  \tag{22.1.2.1}
  \label{0.22.1.2.1-fr}
    \mathbf S_K^*(\mathfrak m/\mathfrak m^2)
    \longrightarrow\operatorname{gr}_{\mathfrak m}^*(B)
  \end{equation}
  est bijectif, et l'homomorphisme caractéristique
  \begin{equation}
  \tag{22.1.2.2}
  \label{0.22.1.2.2-fr}
    \chi:\Upsilon_{K/B/\Lambda}\longrightarrow\mathfrak m/\mathfrak m^2
  \end{equation}
  est injectif.
\end{enumerate}
\end{corollary}

\begin{proof}
Cela résulte aussitôt de (22.1.1), les conditions $\alpha$) et $\delta$)
étant ici trivialement vérifiées, puisque $K$ est un corps.
\end{proof}

\begin{remark}[22.1.3]
\label{0.22.1.3-fr}
Supposons vérifiées les hypothèses de (22.1.2) et en outre supposons que
$\mathfrak m/\mathfrak m^2$ soit un $K$-espace vectoriel de rang \emph{fini}.
Alors, pour que $B$ soit une $A$-algèbre formellement lisse pour la topologie
$\mathfrak m$-préadique, il suffit que les conditions suivantes soient
satisfaites~:

\begin{env}[22.1.3.1]
\label{0.22.1.3.1-fr}
Étant donnés une algèbre de polynômes $E=K[T_1,\ldots,T_n]$, l'idéal maximal
$\mathfrak n$ de $E$ engendré par les $T_i$ $(1\leqslant i\leqslant n)$, et
un idéal $\mathfrak b\supset\mathfrak n^k$ dans $E$, tout
$\Lambda$-homomorphisme $v:B\to E/\mathfrak b$ qui, par passage aux
quotients, donne l'identité $K\to K$, se factorise en
\[
  B\xrightarrow{w}E/\mathfrak n^k\longrightarrow E/\mathfrak b,
\]
où $w$ est un $\Lambda$-homomorphisme.
\end{env}

\begin{env}[22.1.3.2]
\label{0.22.1.3.2-fr}
Si $F$ est l'extension triviale $D_K(K)$ («~algèbre des nombres duaux sur
$K$~»), $\rho:A\to F$ un homomorphisme d'anneaux tel que le composé
$A\xrightarrow{\rho}F\to K$ soit l'homomorphisme canonique, alors tout
$\Lambda$-homomorphisme $v:B\to K$ qui, par passage aux quotients, donne
l'identité $K\to K$, se factorise en
\[
  B\xrightarrow{w}F\longrightarrow K,
\]
où $w$ est un $\Lambda$-homomorphisme (pour la structure de $A$-algèbre sur
$F$ définie par $\rho$).
\end{env}

En effet, on a vu (19.5.5) que la condition (22.1.3.1) entraîne que
l'homomorphisme canonique (22.1.2.1) est bijectif. En vertu de (22.1.2), il
suffit ensuite de voir que l'homomorphisme canonique
\[
  u_{B/A/\Lambda}\otimes1_K:
  \Omega^1_{A/\Lambda}\otimes_AK
  \longrightarrow\Omega^1_{B/\Lambda}\otimes_BK
\]
est injectif, ou encore (puisqu'il s'agit d'espaces vectoriels sur $K$) que
pour tout $K$-espace vectoriel $L$, l'homomorphisme canonique
$\operatorname{D\acute{e}r}_\Lambda(B,L)\to\operatorname{D\acute{e}r}_\Lambda(A,L)$ est
\emph{surjectif} (en vertu de (20.4.8))~; il est clair que pour cela, il faut
et il suffit que l'homomorphisme canonique
$\operatorname{D\acute{e}r}_\Lambda(B,K)\to\operatorname{D\acute{e}r}_\Lambda(A,K)$ soit
surjectif, d'où notre assertion, en vertu de (20.1.5).

On notera que les deux conditions (22.1.3.1) et (22.1.3.2) sont entraînées
par la suivante~:

\begin{env}[22.1.3.3]
\label{0.22.1.3.3-fr}
Pour toute $\Lambda$-algèbre locale artinienne $E$ de corps résiduel égal à
$K$, et tout idéal nilpotent $\mathfrak r$ de $E$, tout
$\Lambda$-homomorphisme $B\to E/\mathfrak r$ qui, par passage aux quotients,
donne l'identité $K\to K$, se factorise en
$B\xrightarrow{u}E\to E/\mathfrak r$, où $u$ est un
$\Lambda$-homomorphisme.
\end{env}
\end{remark}

\oldpage[0\textsubscript{IV-1}]{185}

Ce résultat se généralise comme suit~:

\begin{proposition}[22.1.4]
\label{0.22.1.4-fr}
Soit $\varphi:A\to B$ un homomorphisme local d'anneaux locaux noethériens,
$k$ (resp. $K$) le corps résiduel de $A$ (resp. $B$). Pour que $B$ soit une
$A$-algèbre formellement lisse (pour les topologies préadiques), il faut et
il suffit que pour tout anneau local artinien $C$, de corps résiduel égal à
$K$, tout homomorphisme local $A\to C$ (faisant de $C$ une $A$-algèbre
topologique), et tout idéal nilpotent $\mathfrak J$ de $C$, tout
$A$-homomorphisme local $B\to C/\mathfrak J$ tel que l'homomorphisme
$K\to K$ obtenu par passage aux quotients soit l'identité, se factorise en
\[
  B\xrightarrow{f}C\longrightarrow C/\mathfrak J,
\]
où $f$ est un $A$-homomorphisme local.
\end{proposition}

\begin{proof}
La condition étant nécessaire par définition (19.3.1), montrons qu'elle est
suffisante. Comme les $\widehat A$-homomorphismes locaux de $\widehat B$ dans
une $\widehat A$-algèbre locale discrète proviennent par extension des
$A$-homomorphismes locaux de $B$ dans cette algèbre, on peut se borner au cas
où $A$ et $B$ sont \emph{complets}, en vertu de (19.3.6). Lorsque $A$ est un
\emph{corps} $k$, la proposition résulte de (22.1.3.3), où l'on prend pour
$\Lambda$ le sous-corps premier de $k$, et de (19.6.1).

Dans le cas général, posons $B_0=B\otimes_Ak=B/\mathfrak mB$
($\mathfrak m$ idéal maximal de $A$), qui est complet. Si $C_0$ est une
$k$-algèbre locale artinienne de corps résiduel égal à $K$, $\mathfrak J_0$
un idéal nilpotent de $C_0$, tout $k$-homomorphisme local
$g_0:B_0\to C_0/\mathfrak J_0$ donnant l'identité sur $K$ par passage aux
quotients, donne par composition un $A$-homomorphisme local
\[
  g:B\longrightarrow B_0\xrightarrow{g_0}C_0/\mathfrak J_0
\]
qui par hypothèse se factorise donc en
$B\xrightarrow{f}C_0\to C_0/\mathfrak J_0$, où $f$ est un
$A$-homomorphisme local~; mais comme $C_0$ est une $k$-algèbre, $f$ se
factorise en $B\to B_0\xrightarrow{f_0}C_0$, où $f_0$ est un
$k$-homomorphisme local. On voit donc que les hypothèses de l'énoncé sont
encore remplies lorsqu'on y substitue $B_0$ et $k$ à $B$ et $A$
respectivement~; donc $B_0$ est une $k$-algèbre \emph{formellement lisse}
d'après ce que l'on vient de voir.

Appliquant (19.7.2) et (19.7.1), on voit qu'il existe un anneau local
noethérien complet $B'$, un homomorphisme local $A\to B'$ faisant de $B'$ un
$A$-module \emph{plat} et une $A$-algèbre \emph{formellement lisse}, et un
$k$-isomorphisme
\[
  u'_0:B'_0=B'\otimes_Ak\xrightarrow{\sim}B_0.
\]
Soit $v'_0$ l'isomorphisme réciproque de $u'_0$~; nous nous proposons de
montrer qu'il existe un $A$-homomorphisme local $v:B\to B'$ tel que $v'_0$
provienne de $v$ par passage aux quotients. Pour cela, désignons par
$\mathfrak n$ et $\mathfrak n'$ les idéaux maximaux respectifs de $B$ et
$B'$, et, pour tout $j$, soit
\[
  w_j:B/(\mathfrak n^{j+1}+\mathfrak mB)
  \longrightarrow B'/(\mathfrak n'^{j+1}+\mathfrak mB')
\]
l'homomorphisme déduit de $v'_0$ par passage aux quotients~; il suffira de
former pour tout $j$ un $A$-homomorphisme local
$v_j:B/\mathfrak n^{j+1}\to B'/\mathfrak n'^{j+1}$ tel que les diagrammes
\[
\xymatrix@C=4em@R=3em{
  B/\mathfrak n^{j+1} \ar[r]^{v_j} \ar[d] &
  B'/\mathfrak n'^{j+1} \ar[d] \\
  B/\mathfrak n^j \ar[r]_{v_{j-1}} & B'/\mathfrak n'^j
}
\qquad
\xymatrix@C=4em@R=3em{
  B/\mathfrak n^{j+1} \ar[r]^{v_j} \ar[d] &
  B'/\mathfrak n'^{j+1} \ar[d] \\
  B/(\mathfrak n^{j+1}+\mathfrak mB) \ar[r]_{w_j} &
  B'/(\mathfrak n'^{j+1}+\mathfrak mB')
}
\]
soient commutatifs~; $B$ et $B'$ étant complets, l'homomorphisme
$v=\varprojlim v_j$ répondra à la question.

\oldpage[0\textsubscript{IV-1}]{186}

Or, la formation par récurrence de $v_j$ résulte de l'hypothèse sur $B$ et
du même raisonnement que dans (19.3.10.1), en utilisant le lemme
(19.3.10.2) et le fait que
\[
  \mathfrak n'^j+(\mathfrak n'^{j+1}+\mathfrak mB')
  =\mathfrak n'^j+\mathfrak mB'.
\]
Il résulte alors de (19.7.1.4) que $v$ est un $A$-\emph{isomorphisme}, car
$B'$ est un $A$-module plat, et $B$ est complet pour la topologie
$\mathfrak m$-préadique, les idéaux $\mathfrak m^iB$ étant fermés dans $B$
pour la topologie $\mathfrak n$-adique ($\mathbf 0_{\mathrm I}$, 7.3.5).
C.Q.F.D.
\end{proof}

\subsection{Caractérisation différentielle des algèbres locales formellement
lisses sur un corps}
\label{subsection:0.22.2-fr}

\begin{env}[22.2.1]
\label{0.22.2.1-fr}
Soient $k$ un corps, $P$ son sous-corps premier, $A$ une $k$-algèbre qui est
un anneau \emph{local}, $\mathfrak m$ son idéal maximal, $K=A/\mathfrak m$
son corps résiduel. Comme $K$ est séparable sur $P$, $K$ est une
$P$-algèbre formellement lisse pour la topologie discrète (19.6.1), et par
suite la $P$-extension $A/\mathfrak m^2$ de $K$ par
$\mathfrak m/\mathfrak m^2$ est $P$-\emph{triviale}, et l'homomorphisme
caractéristique
\begin{equation}
\tag{22.2.1.1}
\label{0.22.2.1.1-fr}
  \chi_{A/k}:\Upsilon_{K/k}\longrightarrow\mathfrak m/\mathfrak m^2
\end{equation}
est donc défini (20.6.23).
\end{env}

\begin{theorem}[22.2.2]
\label{0.22.2.2-fr}
Sous les conditions de (22.2.1), pour que $A$ soit une $k$-algèbre
formellement lisse pour la topologie $\mathfrak m$-préadique, il faut et il
suffit que les conditions suivantes soient satisfaites~:
\begin{enumerate}
  \item[(i)] L'homomorphisme canonique
  $\mathbf S_K^*(\mathfrak m/\mathfrak m^2)\to
  \operatorname{gr}_{\mathfrak m}^*(A)$ est bijectif.
  \item[(ii)] L'homomorphisme caractéristique
  $\chi_{A/k}:\Upsilon_{K/k}\to\mathfrak m/\mathfrak m^2$ est injectif.
\end{enumerate}
\end{theorem}

\begin{proof}
C'est un cas particulier de (22.1.2), où l'on remplace $\Lambda$ par $P$,
$A$ par $k$ et $B$ par $A$~; $K$ et $k$, étant séparables sur $P$, sont en
effet des $P$-algèbres formellement lisses (19.6.1).
\end{proof}

\begin{remark}[22.2.3]
\label{0.22.2.3-fr}
La condition (ii) de (22.2.2) peut être remplacée par une quelconque des
suivantes~:
\begin{enumerate}
  \item[a)] L'homomorphisme canonique
  $\Omega^1_k\otimes_kK\to\Omega^1_A\otimes_AK$ est injectif.
  \item[b)] Pour tout $n$, l'homomorphisme canonique
  \[
    \Omega^1_k\otimes_k(A/\mathfrak m^{n+1})
    \longrightarrow
    \Omega^1_A\otimes_A(A/\mathfrak m^{n+1})
  \]
  est inversible à gauche.
  \item[c)] L'homomorphisme canonique
  $\Omega^1_k\widehat\otimes_kA\to\widehat\Omega^1_A$ est inversible à
  gauche.
  \item[d)] $A$ est une $k$-algèbre formellement lisse (pour la topologie
  $\mathfrak m$-préadique) \emph{relativement au corps premier}.
\end{enumerate}
En effet, on a déjà remarqué dans (22.1.1) que l'injectivité de $\chi_{A/k}$
est équivalente à $a)$, et que la conjonction de $a)$ et de la condition (i)
de (22.2.2) équivaut à celle de $d)$ et de (i). Enfin, on a aussi vu dans
(22.1.1) que $b)$ et $d)$ sont équivalentes, et l'équivalence de $b)$ et
$c)$ résulte de (19.1.8).
\end{remark}
\begin{env}[22.2.4]
\label{0.22.2.4-fr}
La principale application de (22.2.2) a trait aux $k$-algèbres locales
noethériennes et \emph{complètes} (cas où $\mathfrak m/\mathfrak m^2$ est un
$K$-espace vectoriel de rang \emph{fini}). De façon précise, soient $k$ un
corps, $K$ une extension de $k$, $V$ un $K$-espace vectoriel de

\oldpage[0\textsubscript{IV-1}]{187}

rang fini~; nous considérons les triplets $(A,\alpha,\beta)$, où $A$ est une
$k$-algèbre locale noethérienne complète \emph{formellement lisse} d'idéal
maximal $\mathfrak m$, $\alpha:A/\mathfrak m\xrightarrow{\sim}K$ un
isomorphisme de $k$-algèbres,
$\beta:\mathfrak m/\mathfrak m^2\xrightarrow{\sim}V$ un di-isomorphisme
d'espaces vectoriels (pour l'isomorphisme $\alpha$). Étant donnés deux tels
triplets $(A,\alpha,\beta)$, $(A',\alpha',\beta')$, une \emph{équivalence} du
premier sur le second est un $k$-isomorphisme $u:A\xrightarrow{\sim}A'$ tel
que (si $\mathfrak m'$ est l'idéal maximal de $A'$), les diagrammes
\[
\xymatrix@C=4em@R=3em{
  A/\mathfrak m \ar[r]^{\operatorname{gr}^0(u)} \ar[d]_{\alpha} &
  A'/\mathfrak m' \ar[d]^{\alpha'} \\
  K \ar[r]_{1_K} & K
}
\qquad
\xymatrix@C=4em@R=3em{
  \mathfrak m/\mathfrak m^2 \ar[r]^{\operatorname{gr}^1(u)}
    \ar[d]_{\beta} &
  \mathfrak m'/\mathfrak m'^2 \ar[d]^{\beta'} \\
  V \ar[r]_{1_V} & V
}
\]
sont commutatifs. Il est immédiat que les \emph{classes d'équivalence} des
triplets $(A,\alpha,\beta)$ forment un \emph{ensemble} que nous noterons
$\operatorname{FL}(K/k,V)$. Remarquons maintenant qu'à tout triplet
$(A,\alpha,\beta)$ est associé le $K$-homomorphisme composé d'espaces
vectoriels
\[
  \Upsilon_{K/k}\xrightarrow{\chi_{A/k}}
  \mathfrak m/\mathfrak m^2\xrightarrow{\beta}V
\]
(où $\chi_{A/k}$ est considéré comme un di-homomorphisme relatif à
$\alpha^{-1}$)~; la définition précédente prouve (par (20.6.22.2)) que ce
$K$-homomorphisme \emph{ne dépend que de la classe d'équivalence} de
$(A,\alpha,\beta)$ dans $\operatorname{FL}(K/k,V)$~; autrement dit on a
défini une application canonique
\begin{equation}
\tag{22.2.4.1}
\label{0.22.2.4.1-fr}
  \operatorname{FL}(K/k,V)\longrightarrow
  \operatorname{Hom}_K(\Upsilon_{K/k},V).
\end{equation}
\end{env}

\begin{proposition}[22.2.5]
\label{0.22.2.5-fr}
Soient $k$ un corps, $K$ une extension de $k$, $V$ un $K$-espace vectoriel
de rang fini. L'application canonique (22.2.4.1) est une bijection de
$\operatorname{FL}(K/k,V)$ sur l'ensemble des $K$-homomorphismes injectifs
de $\Upsilon_{K/k}$ dans $V$.
\end{proposition}

\begin{proof}
(i) Pour montrer que (22.2.4.1) est injective, il faut prouver ce qui suit~:
soient $A$, $A'$ deux $k$-algèbres locales noethériennes complètes
formellement lisses, $\mathfrak m$, $\mathfrak m'$ leurs idéaux maximaux
respectifs, $\alpha:A/\mathfrak m\xrightarrow{\sim}K$,
$\alpha':A'/\mathfrak m'\xrightarrow{\sim}K$ deux $k$-isomorphismes~;
supposons donnés deux $k$-isomorphismes
\[
  v^0:A/\mathfrak m\xrightarrow{\sim}A'/\mathfrak m',\qquad
  v^1:\mathfrak m/\mathfrak m^2\xrightarrow{\sim}
      \mathfrak m'/\mathfrak m'^2
\]
tels que les diagrammes
\begin{equation}
\tag{22.2.5.1}
\label{0.22.2.5.1-fr}
\xymatrix@C=4em@R=3em{
  A/\mathfrak m \ar[r]^{v^0} \ar[dr]_{\alpha} &
  A'/\mathfrak m' \ar[dl]^{\alpha'} \\
  & K
}
\qquad
\xymatrix@C=4em@R=3em{
  \mathfrak m/\mathfrak m^2 \ar[r]^{v^1} &
  \mathfrak m'/\mathfrak m'^2 \\
  \Upsilon_{K/k} \ar[u]^{\chi_A} \ar[ur]_{\chi_{A'}} &
}
\end{equation}
soient commutatifs. Il s'agit de prouver qu'il existe un $k$-isomorphisme
$u:A\xrightarrow{\sim}A'$ tel que $\operatorname{gr}^0(u)=v^0$ et
$\operatorname{gr}^1(u)=v^1$. Remarquons en premier lieu que puisque $K$ est
un corps, l'homomorphisme canonique (20.6.7.1) est \emph{bijectif} (20.6.11),
donc il existe déjà un $k$-isomorphisme
$v:A/\mathfrak m^2\xrightarrow{\sim}A'/\mathfrak m'^2$ tel que
$\operatorname{gr}^0(v)=v^0$ et $\operatorname{gr}^1(v)=v^1$. Notons
maintenant que puisque $A'$ est métrisable et complète, et $A$ une
$k$-algèbre formellement lisse, l'homomorphisme composé
$A\to A/\mathfrak m^2\xrightarrow{v}A'/\mathfrak m'^2$ se factorise en
$A\xrightarrow{u}A'\to A'/\mathfrak m'^2$, où $u$ est

\oldpage[0\textsubscript{IV-1}]{188}

un $k$-homomorphisme (19.3.11), et l'on a donc déjà
$\operatorname{gr}^0(u)=v^0$ et $\operatorname{gr}^1(u)=v^1$. Or, on a le
diagramme commutatif
\[
\xymatrix@C=5em@R=3em{
  \mathbf S_K^*(\mathfrak m/\mathfrak m^2) \ar[r] \ar[d]_{w} &
  \operatorname{gr}_{\mathfrak m}^*(A) \ar[d]^{\operatorname{gr}^*(u)} \\
  \mathbf S_K^*(\mathfrak m'/\mathfrak m'^2) \ar[r] &
  \operatorname{gr}_{\mathfrak m'}^*(A')
}
\]
où les flèches horizontales sont les homomorphismes canoniques, et $w$
provient de $v^0$ et $v^1$~; comme $v^0$ et $v^1$ sont bijectifs, il en est
de même de $w$, et puisque $A$ et $A'$ sont des $k$-algèbres formellement
lisses, on déduit de (22.2.2, (i)) que $\operatorname{gr}(u)$ est bijectif~;
mais comme $A$ et $A'$ sont séparés et complets, cela entraîne que $u$
lui-même est \emph{bijectif} (Bourbaki, \emph{Alg. comm.}, chap.~III, §~2,
n\textsuperscript{o}~8, cor.~3 du th.~1).

(ii) Reste à montrer que l'image de (22.2.4.1) est l'ensemble des
homomorphismes \emph{injectifs} de $\Upsilon_{K/k}$ dans $V$~; on sait déjà
par (22.2.2, (ii)) qu'elle est contenue dans cet ensemble~; il suffira alors
de prouver le lemme suivant.
\end{proof}

\begin{lemma}[22.2.5.2]
\label{0.22.2.5.2-fr}
Pour tout $K$-homomorphisme $h:\Upsilon_{K/k}\to V$, il existe une
$k$-algèbre $A$ qui est un anneau local noethérien régulier et complet,
d'idéal maximal $\mathfrak m$, de corps résiduel $K$, et un $K$-isomorphisme
$\beta:\mathfrak m/\mathfrak m^2\xrightarrow{\sim}V$ tels que
$h=\beta\circ\chi_{A/k}$.
\end{lemma}

\begin{proof}
En effet, si ce lemme est prouvé, le fait que $A$ est régulier entraîne que
la condition (i) de (22.2.2) est satisfaite (17.1.1)~; si en outre $h$ est
injectif, (22.2.2) montrera que $A$ est une $k$-algèbre formellement lisse
dont la classe aura $h$ pour image (par contre si $h$ n'est pas injectif, la
$k$-algèbre $A$ dont (22.2.5.2) prouve l'existence n'est pas formellement
lisse).

Pour prouver (22.2.5.2), soient $B$ la $K$-algèbre $\mathbf S_K^*(V)$,
$\mathfrak n$ l'idéal d'augmentation de $B$ et $A=\widehat B$ la
$K$-algèbre complétée de $B$ pour la topologie $\mathfrak n$-préadique (de
sorte que $A$ est isomorphe à l'algèbre de séries formelles
$K[[T_1,\ldots,T_n]]$ si $n=\operatorname{rg}_KV$). En vertu de (20.6.11)
il existe sur $A/\mathfrak m^2$ (où $\mathfrak m=\mathfrak nA$) une structure
de $k$-extension de $K$ par $V$ (\emph{distincte} en général de la structure
$k$-triviale définie par l'injection canonique $k\to K\to A$) telle que $h$
soit l'homomorphisme caractéristique de cette extension. Si
$f:k\to A/\mathfrak m^2$ est l'homomorphisme structural, le fait que $k$
soit séparable sur le corps premier permet (19.6.1) de factoriser $f$ en
$k\xrightarrow{g}A\to A/\mathfrak m^2$, et la $k$-algèbre $A$ ainsi définie
répond à la question (20.6.22.2).
\end{proof}

\begin{theorem}[22.2.6]
\label{0.22.2.6-fr}
Soient $k$ un corps, $K$ une extension de $k$. Pour qu'il existe une
$k$-algèbre locale noethérienne $A$, d'idéal maximal $\mathfrak m$, telle que
$A/\mathfrak m=K$ et que $A$ soit formellement lisse (pour sa topologie
$\mathfrak m$-préadique), il faut et il suffit que $\Upsilon_{K/k}$ soit un
$K$-espace vectoriel de rang fini. La $k$-algèbre $A$ est alors déterminée
(à un $k$-isomorphisme près donnant par passage aux quotients l'identité
$K\to K$) par sa dimension, qui est soumise à la seule condition d'être
$\geqslant\operatorname{rg}_K\Upsilon_{K/k}$.
\end{theorem}

\begin{proof}
Comme $\mathfrak m/\mathfrak m^2$ est de rang fini sur $K$, la nécessité de
la condition résulte de (22.2.2, (ii))~; inversement, (22.2.5) montre que la
condition est suffisante et que l'on peut prendre pour
$\mathfrak m/\mathfrak m^2$ un $K$-espace vectoriel de rang quelconque
$\geqslant\operatorname{rg}_K(\Upsilon_{K/k})$.
\end{proof}

\oldpage[0\textsubscript{IV-1}]{189}

En particulier, l'homomorphisme identique de $\Upsilon_{K/k}$ associe
\emph{canoniquement} à $K$ une $k$-algèbre noethérienne complète,
formellement lisse, pour laquelle $\Upsilon_{K/k}$ est isomorphe à
$\mathfrak m/\mathfrak m^2$.

\begin{remarks}[22.2.7]
\label{0.22.2.7-fr}
\begin{enumerate}
  \item[(i)] Soit $\Lambda$ un anneau local noethérien, de corps résiduel
  $k$~; il résulte de (19.7.1) et (19.7.2) que la détermination des
  $\Lambda$-algèbres $A$ qui sont des anneaux locaux noethériens complets et
  qui sont formellement lisses est \emph{équivalente} au même problème où
  $\Lambda$ est remplacé par $k$, et est donc en principe entièrement
  résolue par (22.2.5), qui ramène la question à la structure des \emph{corps
  résiduels} des $k$-algèbres cherchées, en tant qu'extensions de $k$.

  \item[(ii)] Le fait que $\Upsilon_{K/k}$ soit de rang fini sur $K$
  \emph{n'entraîne pas} que $K$ soit «~de multiplicité radicielle finie~» sur
  $k$ (cf. (19.6.6)). Par exemple, soient $k_0$ un corps parfait de
  caractéristique $p>0$, $k=k_0(T)$ le corps des fractions rationnelles à une
  indéterminée sur $k_0$ et $K=k^{p^{-\infty}}$ la plus petite extension
  parfaite de $k$. On a alors $\Omega^1_K=\Omega^1_{K/k}=0$ (21.4.4), donc
  par définition (20.6.1.1)
  $\Upsilon_{K/k}=\Omega^1_k\otimes_kK$, et comme $k^p=k_0(T^p)$, $T$ est
  une $p$-base de $k$ sur $k^p$, donc (21.4.5), $\Omega^1_k$ est de rang $1$
  sur $k$, et par suite $\Upsilon_{K/k}$ de rang $1$ sur $K$. Mais pour tout
  $n$ le corps résiduel de l'anneau local $K\otimes_kk^{p^{-n}}$ est
  isomorphe à $K$ donc n'est pas séparable sur $k^{p^{-n}}$.
\end{enumerate}
\end{remarks}

Le théorème (22.2.2) se généralise comme suit~:

\begin{proposition}[22.2.8]
\label{0.22.2.8-fr}
Sous les conditions de (22.2.1), supposons que la caractéristique $p$ de $k$
soit $>0$~; soit $(k_\alpha)$ une famille filtrante décroissante de
sous-corps de $k$, telle que
$\bigcap_\alpha k_\alpha(k^p)=k^p$. Alors la condition (ii) de (22.2.2) peut
se remplacer par l'une quelconque des suivantes~:
\begin{enumerate}
  \item[a)] Il existe $\alpha_0$ tel que l'homomorphisme canonique
  \[
    \Omega^1_{k/k_\alpha}\otimes_kK
    \longrightarrow\Omega^1_{A/k_\alpha}\otimes_AK
  \]
  soit injectif pour tout $\alpha\geqslant\alpha_0$.
  \item[b)] Il existe $\alpha_0$ tel que $A$ soit une $k$-algèbre
  formellement lisse (pour la topologie $\mathfrak m$-préadique) relativement
  à $k_\alpha$, pour tout $\alpha\geqslant\alpha_0$.
  \item[c)] Pour tout $\alpha$ et tout entier $n\geqslant0$,
  l'homomorphisme canonique
  \[
    \Omega^1_{k/k_\alpha}\otimes_k(A/\mathfrak m^{n+1})
    \longrightarrow
    \Omega^1_{A/k_\alpha}\otimes_A(A/\mathfrak m^{n+1})
  \]
  est inversible à gauche.
\end{enumerate}
\end{proposition}

\begin{proof}
En effet, si $A$ est une $k$-algèbre formellement lisse, elle l'est \emph{a
fortiori} relativement à tout sous-corps de $k$, ce qui entraîne $c)$ en
vertu de (20.7.2) et (19.1.7)~; il est clair que $c)$ implique $b)$, en vertu
de (20.7.2), et que $b)$ implique $a)$. Enfin, compte tenu de (22.2.3), il
reste à prouver que $a)$ entraîne que l'homomorphisme canonique
\[
  u:\Omega^1_k\otimes_kK\longrightarrow\Omega^1_A\otimes_AK
\]
est injectif. Or, on a pour tout $\alpha$ un diagramme commutatif
\[
\xymatrix@C=5em@R=3em{
  \Omega^1_k\otimes_kK \ar[r]^{u} \ar[d]_{v_\alpha} &
  \Omega^1_A\otimes_AK \ar[d] \\
  \Omega^1_{k/k_\alpha}\otimes_kK \ar[r]_{u_\alpha} &
  \Omega^1_{A/k_\alpha}\otimes_AK
}
\]

\oldpage[0\textsubscript{IV-1}]{190}

et l'hypothèse que $u_\alpha$ est injectif pour
$\alpha\geqslant\alpha_0$ implique que $\operatorname{Ker}(u)$ est contenu
dans l'\emph{intersection} des $\operatorname{Ker}(v_\alpha)$. Mais on a vu
(21.8.3) que cette intersection est nulle en vertu de l'hypothèse
$\bigcap_\alpha k_\alpha(k^p)=k^p$, ce qui achève la démonstration.
\end{proof}

\begin{proposition}[22.2.9]
\label{0.22.2.9-fr}
Sous les hypothèses de (22.2.1), les conditions suivantes sont
équivalentes~:
\begin{enumerate}
  \item[a)] $A/\mathfrak m^2$ est une $k$-extension $k$-triviale de $K$ par
  $\mathfrak m/\mathfrak m^2$.
  \item[b)] $\chi_{A/k}=0$.
  \item[c)] L'homomorphisme canonique (20.5.11.2)
  \[
    \delta_{K/A/k}:\mathfrak m/\mathfrak m^2
    \longrightarrow\Omega^1_{A/k}\otimes_AK
  \]
  est injectif (en d'autres termes, la suite
  \begin{equation}
  \tag{22.2.9.1}
  \label{0.22.2.9.1-fr}
    0\longrightarrow\mathfrak m/\mathfrak m^2
    \longrightarrow\Omega^1_{A/k}\otimes_AK
    \longrightarrow\Omega^1_{K/k}\longrightarrow0
  \end{equation}
  est exacte).
\end{enumerate}
\end{proposition}

\begin{proof}
L'équivalence de $b)$ et $c)$ résulte de la suite exacte (20.6.22.1)~;
l'équivalence de $c)$ et $a)$ résulte de (20.5.12), la suite (22.2.9.1)
étant scindée si elle est exacte, puisqu'il s'agit d'espaces vectoriels.

On notera que si $K$ est \emph{séparable sur $k$}, les conditions
équivalentes de (22.2.9) sont satisfaites, en vertu de la définition de
$\chi_{A/k}$ (22.2.1) et de (20.6.3).
\end{proof}

\begin{corollary}[22.2.10]
\label{0.22.2.10-fr}
Sous les hypothèses de (22.2.1), considérons un sous-corps $k_0$ de $k$. Les
conditions suivantes sont équivalentes~:
\begin{enumerate}
  \item[a)] $A/\mathfrak m^2$ est une $k_0$-extension $k_0$-triviale de $K$
  par $\mathfrak m/\mathfrak m^2$.
  \item[b)] L'homomorphisme composé
  \[
    \Upsilon_{K/k_0}\longrightarrow\Upsilon_{K/k}
    \xrightarrow{\chi_{A/k}}\mathfrak m/\mathfrak m^2
  \]
  est nul.
  \item[c)] L'homomorphisme caractéristique
  $\chi_{A/k}:\Upsilon_{K/k}\to\mathfrak m/\mathfrak m^2$ se factorise en
  \[
    \Upsilon_{K/k}\longrightarrow\Upsilon_{K/k/k_0}
    \xrightarrow{\chi'}\mathfrak m/\mathfrak m^2.
  \]
\end{enumerate}
En outre, lorsque ces conditions sont remplies, l'homomorphisme $\chi'$ est
déterminé de façon unique et coïncide avec l'homomorphisme caractéristique
de la $k_0$-extension $k_0$-triviale $A/\mathfrak m^2$ de $K$ par
$\mathfrak m/\mathfrak m^2$.
\end{corollary}

\begin{proof}
Comme le composé dans $b)$ n'est autre que $\chi_{A/k_0}$, l'équivalence de
$a)$ et $b)$ résulte de (22.2.9)~; d'autre part, $b)$ et $c)$ sont
équivalentes en vertu de la suite exacte (20.6.18.1)
\[
  \Upsilon_{K/k_0}\longrightarrow\Upsilon_{K/k}
  \longrightarrow\Upsilon_{K/k/k_0}\longrightarrow0,
\]
et comme l'homomorphisme canonique
$\Upsilon_{K/k}\to\Upsilon_{K/k/k_0}$ est surjectif, cela implique l'unicité
de $\chi'$~; le fait que $\chi'$ est l'homomorphisme caractéristique de la
$k_0$-extension $A/\mathfrak m^2$ résulte de (20.6.22.2).
\end{proof}

\begin{corollary}[22.2.11]
\label{0.22.2.11-fr}
Sous les hypothèses de (22.2.1), soit $p$ l'exposant caractéristique de $k$,
et soit $(k_\alpha)$ une famille filtrante décroissante de sous-corps de $k$
telle que $\bigcap_\alpha k_\alpha(k^p)=k^p$. Si $\Upsilon_{K/k}$ est de rang
fini sur $K$, il existe un indice $\alpha$ tel que $A/\mathfrak m^2$ soit une
$k_\alpha$-extension $k_\alpha$-triviale de $K$ par
$\mathfrak m/\mathfrak m^2$.
\end{corollary}

\begin{proof}
On sait en effet (21.8.3) qu'il existe alors un indice $\alpha$ tel que
l'homomorphisme canonique
$\Upsilon_{K/k}\to\Upsilon_{K/k/k_\alpha}$ soit bijectif, donc la condition
$b)$ de (22.2.10) (où $k_\alpha$ remplace $k_0$) est satisfaite.
\end{proof}
\oldpage[0\textsubscript{IV-1}]{191}

\subsection{Application aux relations entre certains anneaux locaux et leurs complétés}
\label{subsection:0.22.3-fr}

\begin{lemma}[22.3.1]
\label{0.22.3.1-fr}
Soient $A_0$ un anneau semi-local noethérien, $A$ une $A_0$-algèbre finie
(qui est donc un anneau semi-local noethérien, cf. Bourbaki,
\emph{Alg. comm.}, chap.~IV, §~2, n\textsuperscript{o}~5, cor.~3 de la
prop.~9). Si $\widehat A_0$ et $\widehat A$ sont les complétés de $A_0$ et
$A$ pour leurs topologies préadiques respectives, alors, lorsqu'on munit
$A_0$, $A$ et $\widehat A$ des topologies discrètes, $\widehat A$ est une
$A$-algèbre formellement lisse relativement à $A_0$ (19.9.1).
\end{lemma}

\begin{proof}
On sait en effet que
$\widehat A=A\otimes_{A_0}\widehat A_0$ (Bourbaki,
\emph{Alg. comm.}, chap.~IV, §~2, n\textsuperscript{o}~5, cor.~3 de la
prop.~9 et chap.~III, §~3, n\textsuperscript{o}~4, th.~1), donc la proposition
résulte de (19.9.3).
\end{proof}

\begin{proposition}[22.3.2]
\label{0.22.3.2-fr}
Soient $A$ un anneau local noethérien régulier, $K$ son corps des fractions,
$p$ la caractéristique de $K$. On suppose vérifiée l'une des hypothèses
suivantes~:
\begin{enumerate}
  \item[(i)] $p=0$.
  \item[(ii)] $p>0$ et il existe une famille filtrante décroissante
  $(A_\alpha)$ de sous-anneaux noethériens de $A$, tels que $A$ soit une
  $A_\alpha$-algèbre finie pour tout $\alpha$, que
  $A^{p^{h(\alpha)}}\subset A_\alpha$ pour un entier $h(\alpha)$ convenable,
  et que, si $K_\alpha$ désigne le corps des fractions de $A_\alpha$, on ait
  \[
    \bigcap_\alpha K_\alpha(K^p)=K^p.
  \]
\end{enumerate}
Soient alors $B$ une $A$-algèbre finie intègre et contenant $A$,
$\mathfrak n$ un idéal premier de $B$, $C$ l'anneau local $B_\mathfrak n$,
$\mathfrak q$ un idéal premier du complété $\widehat C$ tel que
$\mathfrak q\cap C=0$, de sorte que l'anneau local
$\widehat C_\mathfrak q$ est une algèbre sur le corps des fractions $L$ de
$B$. Alors $\widehat C_\mathfrak q$ est une $L$-algèbre formellement lisse
pour sa topologie $\mathfrak q$-adique, et par suite un anneau
géométriquement régulier sur $L$ (19.6.5).
\end{proposition}

\begin{proof}
Distinguons deux cas.

I) Supposons que $\mathfrak n$ contienne l'idéal maximal $\mathfrak m$ de
$A$, et par suite $\mathfrak n\cap A=\mathfrak m$. Alors $\mathfrak n$ est
maximal (Bourbaki, \emph{Alg. comm.}, chap.~V, §~2,
n\textsuperscript{o}~1, prop.~1) et si $\mathfrak r$ est le radical de
l'anneau semi-local $B$, $\widehat C$ est le séparé complété de $B$ pour la
topologie $\mathfrak n$-préadique, et un des composants de l'anneau
semi-local $\widehat B$, complété de $B$ pour la topologie
$\mathfrak r$-préadique (Bourbaki, \emph{Alg. comm.}, chap.~III, §~3,
n\textsuperscript{o}~4, prop.~8)~; on peut donc écrire
$\widehat C_\mathfrak q=\widehat B_{\mathfrak q'}$, où $\mathfrak q'$ est
l'idéal premier de $\widehat B$ image réciproque de $\mathfrak q$. Notons
maintenant que $B$ est un sous-anneau de $\widehat B$ et l'hypothèse
$\mathfrak q\cap C=0$ entraîne $\mathfrak q'\cap B=0$, donc
$\widehat B_{\mathfrak q'}$ est aussi localisé de l'anneau
$\widehat B\otimes_B L$ en un de ses idéaux premiers $\mathfrak q''$, dont
$\mathfrak q'$ est l'image réciproque dans $\widehat B$. En outre on a
$\widehat B=\widehat A\otimes_A B$ (Bourbaki, \emph{Alg. comm.}, chap.~IV,
§~2, n\textsuperscript{o}~5, cor. de la prop.~9 et chap.~III, §~3,
n\textsuperscript{o}~4, th.~1), donc
$\widehat B\otimes_B L=\widehat A\otimes_A L$~; enfin, si $\mathfrak p$ est
l'idéal premier de $\widehat A$, image réciproque de $\mathfrak q''$, le
fait que $A\to B$ est injectif et la commutativité du diagramme
\[
  \xymatrix{
    B \ar[r] & \widehat B\\
    A \ar[u] \ar[r] & \widehat A \ar[u]
  }
\]
entraînent que $\mathfrak p\cap A=0$, donc
$(\widehat A\otimes_A L)_{\mathfrak q''}$ est aussi localisé de l'anneau
\[
  \widehat A_\mathfrak p\otimes_A L
  =(\widehat A_\mathfrak p\otimes_A K)\otimes_K L
\]
en un de ses idéaux premiers.

\oldpage[0\textsubscript{IV-1}]{192}

Pour montrer que cet anneau local est une $L$-algèbre formellement lisse,
il suffit donc de prouver que $\widehat A_\mathfrak p\otimes_A K$ est une
$K$-algèbre formellement lisse (19.3.5, (iii) et (iv))~; d'ailleurs
$\widehat A_\mathfrak p\otimes_A K=\widehat A_\mathfrak p$ puisque
$\mathfrak p\cap A=0$. On va appliquer les critères (22.2.2) et (22.2.8).
En premier lieu, comme $A$ est régulier, il en est de même de $\widehat A$
(17.1.5), donc de $\widehat A_\mathfrak p$ (17.3.2), et la condition (i) de
(22.2.2) est vérifiée (17.1.1). Il suffit par suite (lorsque $p>0$) de
vérifier la condition b) de (22.2.8). Or, (22.3.1) montre que, pour tout
$\alpha$, $\widehat A$ est une $A$-algèbre formellement lisse relativement
à $A_\alpha$ (pour les topologies discrètes)~; par suite
\[
  \widehat A\otimes_{A_\alpha}K_\alpha
  =\widehat A\otimes_A(A\otimes_{A_\alpha}K_\alpha)
\]
est une $(A\otimes_{A_\alpha}K_\alpha)$-algèbre formellement lisse
relativement à $K_\alpha$, pour les topologies discrètes (19.9.2, (iii)).
Mais comme $A$ est une $A_\alpha$-algèbre finie, on a
$A\otimes_{A_\alpha}K_\alpha=K$~; comme $\widehat A_\mathfrak p$ est un
anneau de fractions de $\widehat A\otimes_A K$, on en conclut (19.9.2,
(iv)) que $\widehat A_\mathfrak p$ est une $K$-algèbre formellement lisse
relativement à $K_\alpha$, pour la topologie discrète, et \emph{a fortiori}
pour sa topologie $\mathfrak p$-adique ((19.3.8) et (19.9.5)). Mais cela est
précisément la condition b) de (22.2.8), donc la proposition est démontrée
dans ce cas pour $p>0$. Lorsque $p=0$, le corps résiduel de
$\widehat A_\mathfrak p$ est séparable sur $K$, et comme
$\widehat A_\mathfrak p$ est un anneau régulier, c'est une $K$-algèbre
formellement lisse en vertu de (19.6.4).

II) Cas général. Soit $\mathfrak s$ l'idéal premier $\mathfrak n\cap A$ de
$A$, et posons $S=A-\mathfrak s$, de sorte que $S^{-1}A=A_\mathfrak s$, et
l'on a $C=(S^{-1}B)_{S^{-1}\mathfrak n}$, où cette fois
$S^{-1}\mathfrak n$ contient l'idéal maximal
$\mathfrak sA_\mathfrak s$ de $A_\mathfrak s$. Comme $A_\mathfrak s$ est
régulier (17.3.2) et $S^{-1}B$ une $A_\mathfrak s$-algèbre finie intègre
contenant $A_\mathfrak s$, tout revient à vérifier la condition (ii) pour
$A_\mathfrak s$ lorsque $p>0$~: or, posons
$\mathfrak s_\alpha=\mathfrak s\cap A_\alpha$ et considérons l'anneau local
$(A_\alpha)_{\mathfrak s_\alpha}$. Pour tout $t\notin\mathfrak s$, on a par
hypothèse, en posant $q=p^{h(\alpha)}$, $t^q\in A_\alpha$ et
$t^q\notin\mathfrak s$, donc $t^q\notin\mathfrak s_\alpha$~; si l'on pose
$S_\alpha=A_\alpha-\mathfrak s_\alpha$, on voit donc que l'on a
$A_\mathfrak s=S_\alpha^{-1}A$, et l'hypothèse entraîne que
$A_\mathfrak s$ est une $(A_\alpha)_{\mathfrak s_\alpha}$-algèbre finie~;
en outre, $K_\alpha$ est aussi le corps des fractions de
$(A_\alpha)_{\mathfrak s_\alpha}$, ce qui achève la démonstration.
\end{proof}

\begin{theorem}[22.3.3]
\label{0.22.3.3-fr}
Soient $A$ un anneau local noethérien complet, $\mathfrak p$ un idéal
premier de $A$, $B$ l'anneau localisé $A_\mathfrak p$, $\mathfrak q$ un
idéal premier du complété $\widehat B$, $\mathfrak r=\mathfrak q\cap B$.
Alors l'anneau localisé $\widehat B_\mathfrak q$ de $\widehat B$ est une
$B_\mathfrak r$-algèbre formellement lisse pour les topologies préadiques.
\end{theorem}

\begin{proof}
L'idéal premier $\mathfrak r$ de $B$ est de la forme
$\mathfrak n_\mathfrak p$, où $\mathfrak n\subset\mathfrak p$ est un idéal
premier de $A$, et l'on a $B_\mathfrak r=A_\mathfrak n$~; soit $k$ le corps
résiduel de $B_\mathfrak r$, qui est aussi le corps des fractions de l'anneau
local intègre et complet $A'=A/\mathfrak n$. Comme $\widehat B$ est un
$B$-module plat $(\mathbf 0_\mathrm I,~7.3.5)$,
$\widehat B_\mathfrak q$ est un $B_\mathfrak r$-module plat
$(\mathbf 0_\mathrm I,~6.3.2)$, et pour prouver que
$\widehat B_\mathfrak q$ est une $B_\mathfrak r$-algèbre formellement
lisse, il suffit de prouver que
\[
  \widehat B_\mathfrak q\otimes_{B_\mathfrak r}k
  =\widehat B_\mathfrak q/\mathfrak r\widehat B_\mathfrak q
\]
est une $k$-algèbre formellement lisse (19.7.1). Or
$\widehat B_\mathfrak q/\mathfrak r\widehat B_\mathfrak q
=(\widehat B/\mathfrak r\widehat B)_{\mathfrak q'}$, où $\mathfrak q'$ est
l'image canonique de $\mathfrak q$ dans
$\widehat B/\mathfrak r\widehat B$~; on a par définition
$\mathfrak q'\cap(B/\mathfrak r)=0$ et
$B/\mathfrak r=A_{\mathfrak p'}$, où
$\mathfrak p'=\mathfrak p/\mathfrak n$~; en outre
$\widehat B/\mathfrak r\widehat B$ est le complété de $B/\mathfrak r$.

On est donc ramené à démontrer le corollaire suivant~:
\end{proof}

\begin{corollary}[22.3.4]
\label{0.22.3.4-fr}
Soient $A$ un anneau local noethérien complet et intègre, $K$ son corps des
fractions, $\mathfrak p$ un idéal premier de $A$, $B$ l'anneau localisé
$A_\mathfrak p$. Alors, pour tout idéal premier $\mathfrak q$

\oldpage[0\textsubscript{IV-1}]{193}

du complété $\widehat B$, tel que $\mathfrak q\cap B=0$, l'anneau local
$\widehat B_\mathfrak q$ est une $K$-algèbre formellement lisse pour sa
topologie $\mathfrak q$-adique, donc un anneau géométriquement régulier sur
$K$.
\end{corollary}

\begin{proof}
En effet, il résulte de (19.8.9) qu'il existe un sous-anneau $A_0$ de $A$
qui est un anneau de séries formelles sur un corps de caractéristique $p>0$
ou sur un anneau de Cohen (on rappelle que les anneaux de Cohen ont un corps
des fractions de caractéristique $0$), et est tel que $A$ soit une
$A_0$-algèbre finie. Or, on sait que l'anneau $A_0$ est régulier (17.3.8) et
vérifie l'une des hypothèses (i), (ii) de (22.3.2), en vertu de (21.8.8)~;
il suffit donc d'appliquer (22.3.2), en remplaçant $B$ par $A$ et $A$ par
$A_0$.
\end{proof}

\subsection{Résultats préliminaires sur les extensions finies d'anneaux locaux dont l'idéal maximal est de carré nul}
\label{subsection:0.22.4-fr}
\label{0.22.4-fr}

\begin{proposition}[22.4.1]
\label{0.22.4.1-fr}
Soient $A$ un anneau local dont l'idéal maximal $\mathfrak m$ est de carré
nul, $K=A/\mathfrak m$ son corps résiduel~; désignons par $V$ l'idéal
$\mathfrak m$ considéré comme $K$-espace vectoriel. Soient $T_i$
$(1\leqslant i\leqslant r)$ des indéterminées, et pour chaque $i$, soit
$F_i$ un polynôme unitaire de $A[T_i]$ de degré $d_i$~; désignons par $A'$ le
quotient de l'anneau de polynômes $A[T_1,\ldots,T_r]$ par l'idéal
$\mathfrak b$ engendré par les $r$ polynômes $F_i$. Supposons que $A'$ soit
un anneau local, et soient $\mathfrak m'$ son idéal maximal,
$K'=A'/\mathfrak m'$ son corps résiduel. Alors~:
\begin{enumerate}
  \item[(i)] Si l'on pose $d=\prod_{i=1}^r d_i$, $A'$ est un $A$-module
  libre de rang $d$, et l'on a $[K':K]\leqslant d$.
  \item[(ii)] Pour que $[K':K]=d$, il faut et il suffit que
  $A'\otimes_AK=A'/\mathfrak mA'$ soit un corps (nécessairement isomorphe à
  $K'$), autrement dit, que $\mathfrak m'=\mathfrak mA'$. On a dans ce cas
  $\mathfrak m'^2=0$, et si $V'$ désigne l'idéal $\mathfrak m'$ considéré
  comme $K'$-espace vectoriel, l'homomorphisme canonique
  $V\otimes_KK'\to V'$ est bijectif (et par suite
  $\operatorname{rg}_{K'}V'=\operatorname{rg}_K V$).
\end{enumerate}
\end{proposition}

\begin{proof}
Sans supposer $A'$ local, il est évident, par division euclidienne et
récurrence sur $r$, que si $t_i$ est la classe de $T_i$ mod.~$\mathfrak b$,
les monômes
$M_\nu=t_1^{\nu_1}t_2^{\nu_2}\cdots t_r^{\nu_r}$, où
$0\leqslant\nu_i\leqslant d_i-1$, forment une base du $A$-module $A'$. Si
$A'$ est supposé local, $K'$ est un quotient de
$A'/\mathfrak mA'=A'\otimes_AK$, donc
$[K':K]\leqslant[A'\otimes_AK:K]=d$, et il ne peut y avoir égalité que si
$\mathfrak m'=\mathfrak mA'$, ce qui démontre (i) et la première assertion
de (ii). D'autre part, il est clair que lorsque
$\mathfrak m'=\mathfrak mA'$, on a $\mathfrak m'^2=0$, et
$V\otimes_KK'=\mathfrak m\otimes_AA'=\mathfrak mA'=V'$ puisque $A'$ est
un $A$-module libre.
\end{proof}

\begin{remark}[22.4.2]
\label{0.22.4.2-fr}
Lorsque $A$ est artinien (autrement dit $\operatorname{rg}_K(V)$ fini),
l'hypothèse que $A'$ est local entraîne toujours que
$\operatorname{rg}_{K'}(\mathfrak m'/\mathfrak m'^2)
\geqslant\operatorname{rg}_K(V)$, comme nous le verrons plus loin (22.5.2).
On notera que ces deux rangs peuvent être égaux sans que l'on ait
$\mathfrak m'=\mathfrak mA'$.
\end{remark}

\begin{proposition}[22.4.3]
\label{0.22.4.3-fr}
Les hypothèses sur $A$ et les notations étant celles de (22.4.1), supposons
que les $F_i$ soient de la forme
\begin{equation}
\label{0.22.4.3.1-fr}
  F_i(T_i)=T_i^{d_i}+\sum_{1\leqslant k\leqslant d_i}
  c_{ik}T_i^{d_i-k},
  \tag{22.4.3.1}
\end{equation}
où $d_i\geqslant2$ pour tout $i$, et $c_{ik}\in\mathfrak m$
(«~polynômes d'Eisenstein~»). Alors~:
\begin{enumerate}
  \item[(i)] $A'$ est un anneau local, et son corps résiduel
  $K'=A'/\mathfrak m'$ est isomorphe à $K$.

  \oldpage[0\textsubscript{IV-1}]{194}

  \item[(ii)] Le noyau de l'homomorphisme canonique
  $V\to V'=\mathfrak m'/\mathfrak m'^2$ est l'idéal $\mathfrak n$ de $A$
  engendré par les $\xi_i=c_{i,d_i}$ («~termes constants~» des $F_i$), et
  le conoyau de cet homomorphisme a pour base sur $K'$ les images des $T_i$~;
  en particulier, on a
  \begin{equation}
  \label{0.22.4.3.2-fr}
    \operatorname{rg}_{K'}(V')=\operatorname{rg}_K(V)+r-r',
    \tag{22.4.3.2}
  \end{equation}
  où $r'=\operatorname{rg}_K(\mathfrak n)$. Lorsque
  $\operatorname{rg}_K(V)$ est fini (autrement dit lorsque $A$ est
  artinien), pour que
  $\operatorname{rg}_{K'}(V')=\operatorname{rg}_K(V)$, il faut et il suffit
  que les $\xi_i$ soient linéairement indépendants sur $K$.
\end{enumerate}
\end{proposition}

\begin{proof}
(i) Il est clair que $A'\otimes_AK=A'/\mathfrak mA'$ est isomorphe à
$K[T_1,\ldots,T_r]/\mathfrak b'$, où $\mathfrak b'$ est l'idéal engendré
par les $r$ polynômes $T_i^{d_i}$~; c'est donc un anneau local de corps
résiduel $K$, dont l'idéal maximal est engendré par les classes des $T_i$
mod.~$\mathfrak b'$, d'où (i).

(ii) Soit $t_i$ $(1\leqslant i\leqslant r)$ la classe de $T_i$
mod.~$\mathfrak b$~; il résulte de (i) que l'idéal maximal $\mathfrak m'$ de
$A'$ est donné par
\[
  \mathfrak m'=\mathfrak mA'+\sum_{i=1}^r t_iA',
\]
d'où, compte tenu de $\mathfrak m^2=0$,
\[
  \mathfrak m'^2=\sum_{i,j}t_it_jA'+\sum_i t_i\mathfrak mA'.
\]
On en conclut que $A'/\mathfrak m'^2$ est isomorphe à l'anneau
$A[T_1,\ldots,T_r]/\mathfrak c$, où $\mathfrak c$ est l'idéal engendré par
les éléments
\[
  F_i\ (1\leqslant i\leqslant r),\qquad
  T_iT_j\ (1\leqslant i\leqslant r,\ 1\leqslant j\leqslant r),\qquad
  xT_i\ (1\leqslant i\leqslant r,\ x\in\mathfrak m).
\]
Comme $d_i\geqslant2$ et que les $T_i^2$ appartiennent à
$\mathfrak c$, les hypothèses faites entraînent que $\mathfrak c$ est aussi
engendré par les éléments
\[
  \xi_i\ (1\leqslant i\leqslant r),\qquad
  T_iT_j\ (1\leqslant i\leqslant r,\ 1\leqslant j\leqslant r),\qquad
  T_ix\ (1\leqslant i\leqslant r,\ x\in\mathfrak m).
\]
Soient
\[
  \mathfrak c_1=\sum_{i,j}T_iT_jA[T_1,\ldots,T_r],\qquad
  \mathfrak c_2=\mathfrak c_1+\sum_i\mathfrak mT_iA[T_1,\ldots,T_r],
\]
de sorte que
$\mathfrak c=\mathfrak c_2+\sum_i\xi_iA[T_1,\ldots,T_r]$. Il est clair que
$A[T_1,\ldots,T_r]/\mathfrak c_1$ est le $A$-module somme directe de $A$ et
des $At_i^{(1)}$, où $t_i^{(1)}$ est la classe de $T_i$
mod.~$\mathfrak c_1$. On en déduit que
$A[T_1,\ldots,T_r]/\mathfrak c_2$ est somme directe de $A$ et des
$Kt_i^{(2)}$, où $t_i^{(2)}$ est la classe de $T_i$
mod.~$\mathfrak c_2$. Enfin $A[T_1,\ldots,T_r]/\mathfrak c$ est somme
directe de $A/\mathfrak n$ et des $Kt_i'$, où $t_i'$ est la classe de $T_i$
mod.~$\mathfrak c$, et $\mathfrak n$ est l'idéal de $A$ engendré par les
$\xi_i$~; dans l'identification de $A'/\mathfrak m'^2$ à
$A[T_1,\ldots,T_r]/\mathfrak c$, l'image de $V$ s'identifie à
$\mathfrak m/\mathfrak n$, et ce qui précède prouve donc (ii).
\end{proof}

\begin{proposition}[22.4.4]
\label{0.22.4.4-fr}
Les hypothèses sur $A$ et les notations étant celles de (22.4.1), supposons
que $K$ soit de caractéristique $p>0$ (donc $p\cdot1\in\mathfrak m$ dans
$A$) et que les $F_i$ soient de la forme
\begin{equation}
\label{0.22.4.4.1-fr}
  F_i(T_i)=T_i^p+p\sum_{k=1}^{p-1}c_{ik}T_i^{p-k}-f_i
  \tag{22.4.4.1}
\end{equation}
avec $c_{ik}\in A$ et $f_i\in A$ pour tout $i$ et tout $k$ tel que
$1\leqslant k\leqslant p-1$~; en outre si $a_i$ est la classe de $f_i$
mod.~$\mathfrak m$, on suppose que la famille $(a_i)_{1\leqslant i\leqslant r}$
est $p$-libre sur $K^p$ («~absolument $p$-libre~»). Alors~:
\begin{enumerate}
  \item[(i)] $A'$ est un anneau local, d'idéal maximal $\mathfrak mA'$,
  dont le corps résiduel $K'$ est isomorphe à l'extension de $K$ obtenue par
  adjonction des $a_i^{1/p}$ $(1\leqslant i\leqslant r)$.

  \oldpage[0\textsubscript{IV-1}]{195}

  \item[(ii)] Les éléments $d_Af_i\otimes1_{K'}$ forment une base du noyau
  $N=\Upsilon^{K'}_{A'/A}$ de l'homomorphisme canonique
  \[
    \Omega^1_A\otimes_AK'\longrightarrow
    \Omega^1_{A'}\otimes_{A'}K'.
  \]
\end{enumerate}
\end{proposition}

\begin{proof}
(i) Cette fois $A'\otimes_AK=A'/\mathfrak mA'$ est isomorphe à
$K[T_1,\ldots,T_r]/\mathfrak b'$, où $\mathfrak b'$ est l'idéal engendré
par les polynômes $T_i^p-a_i$~; l'hypothèse faite sur les $a_i$ entraîne que
cet anneau quotient est un corps (21.1.8), d'où les assertions de (i).

Avant de démontrer (ii), nous établirons le

\begin{lemma}[22.4.4.2]
\label{0.22.4.4.2-fr}
Soient $\Lambda$ un anneau, $A$ une $\Lambda$-algèbre, $\mathfrak J$ un
idéal de $A$ tel que l'anneau $C=A/\mathfrak J$ soit de caractéristique
$p>0$ (ce qui équivaut à dire que $p\cdot1\in\mathfrak J$ dans $A$). Soit
$\pi$ l'image canonique de $p\cdot1$ dans
$\mathfrak J/\mathfrak J^2$. Si $C$ est une $(\Lambda/p\Lambda)$-algèbre
formellement lisse pour les topologies discrètes, on a une suite exacte
canonique
\begin{equation}
\label{0.22.4.4.3-fr}
  0\longrightarrow(\mathfrak J/\mathfrak J^2)/C\cdot\pi
  \longrightarrow\Omega^1_{A/\Lambda}\otimes_AC
  \longrightarrow\Omega^1_{C/\Lambda}\longrightarrow0.
  \tag{22.4.4.3}
\end{equation}
\end{lemma}

\begin{proof}[Démonstration du lemme]
Posons en effet $A'=A/pA$, $\mathfrak J'=\mathfrak J/pA$, de sorte que
$C=A'/\mathfrak J'$. Comme
$A'=A\otimes_\Lambda(\Lambda/p\Lambda)$, on a
$\Omega^1_{A/\Lambda}\otimes_AA'=\Omega^1_{A'/\Lambda'}$ à un
isomorphisme canonique près, $\Lambda'$ désignant l'anneau
$\Lambda/p\Lambda$ (20.5.5). Appliquant alors la suite exacte (20.5.14.1) à
la $\Lambda'$-algèbre formellement lisse $C=A'/\mathfrak J'$, on obtient la
suite exacte
\[
  0\longrightarrow\mathfrak J'/\mathfrak J'^2
  \longrightarrow\Omega^1_{A'/\Lambda'}\otimes_{A'}C
  \longrightarrow\Omega^1_{C/\Lambda'}\longrightarrow0
\]
et comme l'homomorphisme $\Lambda\to\Lambda'$ est surjectif, on a
$\Omega^1_{C/\Lambda'}=\Omega^1_{C/\Lambda}$ à un isomorphisme canonique
près (20.5.15)~; d'où la suite exacte (22.4.4.3) en vertu de ce qui précède.
\end{proof}

(ii) Comme on a $\mathfrak m'^2=0$, nous désignerons encore par $V'$ l'idéal
$\mathfrak m'$ considéré comme $K'$-espace vectoriel, et nous poserons
\begin{equation}
\label{0.22.4.4.4-fr}
  V_0=V/K\cdot p=\mathfrak m/(\mathfrak m^2+pA),\qquad
  V_0'=V'/K'\cdot p=\mathfrak m'/(\mathfrak m'^2+pA').
  \tag{22.4.4.4}
\end{equation}
Appliquant la suite exacte (22.4.4.3) au cas où $\Lambda=\mathbf Z$, à la
$\mathbf Z$-algèbre $A$ (resp. $A'$) et à son idéal $\mathfrak m$
(resp. $\mathfrak m'$), il vient, puisque $K$ (resp. $K'$) est séparable sur
le corps premier $\mathbf Z/p\mathbf Z$, donc une
$(\mathbf Z/p\mathbf Z)$-algèbre formellement lisse (19.6.1), les suites
exactes
\begin{equation}
\label{0.22.4.4.5-fr}
  0\longrightarrow V_0\longrightarrow\Omega^1_A\otimes_AK
  \longrightarrow\Omega^1_K\longrightarrow0,\qquad
  0\longrightarrow V_0'\longrightarrow\Omega^1_{A'}\otimes_{A'}K'
  \longrightarrow\Omega^1_{K'}\longrightarrow0.
  \tag{22.4.4.5}
\end{equation}
Considérons alors le diagramme (22.4.4.6) où les deux colonnes médianes sont
la seconde suite (22.4.4.5) et la première tensorisée par $K'$~; comme la
première suite exacte (22.4.4.5) est formée d'espaces vectoriels sur $K$,
les deux colonnes médianes de (22.4.4.6) sont exactes~; en outre, le
diagramme formé par ces colonnes et les flèches horizontales qui les
connectent est commutatif, en vertu de la démonstration de (22.4.4.2) et de
la commutativité du diagramme (20.5.11.3). La dernière ligne de (22.4.4.6)
est la suite exacte obtenue à partir de (20.6.1.1) en y remplaçant $A$, $B$,
$C$ par $\mathbf Z$, $K$ et $K'$~; la ligne médiane est obtenue par
tensorisation avec $K'$ de la suite exacte de $K$-modules déduite de
(20.6.1.1) en y remplaçant $A$, $B$, $C$ par $\mathbf Z$, $A$, $K$
respectivement. Le diagramme formé de ces deux lignes et des flèches
verticales qui les connectent est commutatif en vertu de la commutativité de
(20.6.4.3). Notons d'autre part qu'on est sous les conditions d'application
de (22.4.1, (ii)), donc la ligne supérieure du diagramme (22.4.4.6) est
exacte, les colonnes extrêmes du diagramme

\oldpage[0\textsubscript{IV-1}]{196}

\begin{equation}
\label{0.22.4.4.6-fr}
  \xymatrix{
    && 0 \ar[d] & 0 \ar[d] && \\
    & 0 \ar[r] & V_0\otimes_KK' \ar[r] \ar[d] & V_0' \ar[r] \ar[d] & 0 & \\
    0 \ar[r] & N \ar[r] \ar[d] & \Omega^1_A\otimes_AK' \ar[r] \ar[d] &
      \Omega^1_{A'}\otimes_{A'}K' \ar[r] \ar[d] &
      \Omega^1_{A'/A}\otimes_{A'}K' \ar[r] \ar[d] & 0 \\
    0 \ar[r] & \Upsilon_{K'/K} \ar[r] \ar[d] &
      \Omega^1_K\otimes_KK' \ar[r] \ar[d] & \Omega^1_{K'} \ar[r] \ar[d] &
      \Omega^1_{K'/K} \ar[r] \ar[d] & 0 \\
    & 0 & 0 & 0 & 0 &
  }
  \tag{22.4.4.6}
\end{equation}
sont donc formées respectivement des noyaux et des conoyaux des flèches
horizontales médianes, ce qui achève de prouver que le diagramme est
commutatif. En outre~:

\begin{lemma}[22.4.4.7]
\label{0.22.4.4.7-fr}
Les homomorphismes
\[
  N\longrightarrow\Upsilon_{K'/K}\qquad\text{et}\qquad
  \Omega^1_{A'/A}\otimes_{A'}K'\longrightarrow\Omega^1_{K'/K}
\]
du diagramme (22.4.4.6) sont bijectifs.
\end{lemma}

\begin{proof}
En effet, le diagramme du serpent (Bourbaki, \emph{Alg. comm.}, chap.~I,
§~1, n\textsuperscript{o}~4, prop.~2) appliqué aux deux colonnes médianes de
(22.4.4.6) donne une suite exacte
\[
  0\longrightarrow N\longrightarrow\Upsilon_{K'/K}\longrightarrow0
  \longrightarrow\Omega^1_{A'/A}\otimes_{A'}K'
  \longrightarrow\Omega^1_{K'/K}\longrightarrow0.
\]
Or, si $t_i$ est l'image canonique de $T_i$ dans $A'$, la relation
$F_i(t_i)=0$, jointe au fait que dans $A'$ on a
$p\cdot1\in\mathfrak mA'\subset\mathfrak m'$, montre aussitôt que l'on a
$d_{A'}f_i\in\mathfrak m'\Omega^1_{A'}$, donc
$d_{A'}f_i\otimes1_{K'}=0$~; cela prouve que les
$d_Af_i\otimes1_{K'}$ appartiennent à $N$. En outre, leurs images dans
$\Upsilon_{K'/K}$ sont les $d_Ka_i\otimes1_{K'}$, qui forment une base de
$\Upsilon_{K'/K}$ sur $K'$ (21.4.7)~; compte tenu de (22.4.4.7), cela achève
de démontrer (22.4.4).
\end{proof}
\end{proof}
\begin{env}[22.4.5]
\label{0.22.4.5-fr}
Considérons maintenant deux anneaux $A$, $B$ \emph{de caractéristique} $p>0$, et
deux homomorphismes
\[
  A\xrightarrow{i}B\xrightarrow{j}A
\]

\oldpage[0\textsubscript{IV-1}]{197}

vérifiant les conditions de (21.3.1). Soit en outre $\mathfrak m$ un idéal
\emph{de carré nul} dans $A$~; posons
\[
  K=A/\mathfrak m,\qquad
  B_{(K)}=B\otimes_AK=B/Bi(\mathfrak m)
\]
et soient
\[
  \varphi:A\longrightarrow K,\qquad
  \psi:B\longrightarrow B_{(K)}
\]
les homomorphismes canoniques.

Comme $\varphi$ est surjectif, $\Omega^1_{B_{(K)}/A}$ s'identifie
canoniquement à $\Omega^1_{B_{(K)}/K}$ (20.5.15). Notons que, puisque
$p\geqslant2$, pour tout $x\in\mathfrak m$, on a $j(i(x))=x^p=0$, autrement
dit $j(i(\mathfrak m))=0$~; par passage aux quotients on déduit donc de $j$ un
homomorphisme
\[
  j':B_{(K)}\longrightarrow A
\]
tel que, si $i'=\psi\circ i:A\longrightarrow B_{(K)}$, on ait
\[
  j'(i'(x))=x^p\qquad\text{et}\qquad i'(j'(y))=y^p
  \qquad\text{pour }x\in A\text{ et }y\in B_{(K)}.
\]
En posant, suivant les notations de (21.3.2),
\[
  \Theta_{B_{(K)}/A}
  =\Omega^1_{B_{(K)}/A}\otimes_{B_{(K)}}A_{[j']}
  =\Omega^1_{B_{(K)}/K}\otimes_{B_{(K)}}A_{[j']},
\]
on a donc (21.3.2.2) un homomorphisme canonique
\begin{equation}
\label{0.22.4.5.1-fr}
  \pi_{B_{(K)}/A}:\Theta_{B_{(K)}/A}\longrightarrow\Omega^1_A.
  \tag{22.4.5.1}
\end{equation}

D'autre part, on déduit de $i$ et $j$, par passage aux quotients (tenant compte
de ce que $j(i(\mathfrak m))=0$), des homomorphismes
\[
  K\xrightarrow{i_0}B_{(K)}\xrightarrow{j_0}K
\]
de sorte que l'on a le diagramme commutatif
\[
  \xymatrix{
    A \ar[r]^i \ar[d]_\varphi & B \ar[r]^j \ar[d]^\psi & A \ar[d]^\varphi \\
    K \ar[r]_{i_0} & B_{(K)} \ar[r]_{j_0} \ar[ur]_{j'} & K
  }
\]
Il est clair en outre que l'on a
\[
  j_0(i_0(s))=s^p\qquad\text{et}\qquad i_0(j_0(t))=t^p
  \qquad\text{pour }s\in K\text{ et }t\in B_{(K)}.
\]
On définit donc de même
\[
  \Theta_{B_{(K)}/K}
  =\Omega^1_{B_{(K)}/K}\otimes_{B_{(K)}}K_{[j_0]}
  =\Theta_{B_{(K)}/A}\otimes_AK,
\]
donc, en tensorisant (22.4.5.1) par $K$, on obtient un
$K$-\emph{homomorphisme canonique}
\begin{equation}
\label{0.22.4.5.2-fr}
  \pi'_{B_{(K)}/K}=\pi_{B_{(K)}/A}\otimes1_K:
  \Theta_{B_{(K)}/K}\longrightarrow\Omega^1_A\otimes_AK.
  \tag{22.4.5.2}
\end{equation}

\oldpage[0\textsubscript{IV-1}]{198}

Il résulte aussitôt des définitions que le diagramme
\[
  \xymatrix{
    \Theta_{B_{(K)}/K} \ar[r]^{\pi'_{B_{(K)}/K}} \ar[dr]_{\pi_{B_{(K)}/K}}
      & \Omega^1_A\otimes_AK \ar[d]^{\varphi_{K/A}} \\
    & \Omega^1_K
  }
\]
est commutatif, $\pi_{B_{(K)}/K}$ étant l'homomorphisme canonique
correspondant à $i_0$ et $j_0$ (21.3.2.2). Par \emph{restriction} à
$\Xi_{B_{(K)}/K}=\operatorname{Ker}(\pi_{B_{(K)}/K})$, l'homomorphisme
$\pi'_{B_{(K)}/K}$ définit donc un \emph{homomorphisme canonique}
\begin{equation}
\label{0.22.4.5.3-fr}
  \Xi_{B_{(K)}/K}\longrightarrow
  \operatorname{Ker}(\Omega^1_A\otimes_AK\longrightarrow\Omega^1_K)=\Upsilon_{K/A}.
  \tag{22.4.5.3}
\end{equation}
\end{env}

\begin{env}[22.4.6]
\label{0.22.4.6-fr}
Les hypothèses et notations étant celles de (22.4.5), supposons en outre que
$\mathfrak m$ soit un idéal \emph{maximal de carré nul}, donc $K$ un
\emph{corps}, et notons $V$ l'idéal $\mathfrak m$ considéré comme $K$-espace
vectoriel. Comme $K$ est une algèbre formellement lisse sur son corps premier
$\mathbf F_p$ (21.5.2) et que $A$ est une algèbre sur $\mathbf F_p$ (de sorte
que $\Omega^1_A=\Omega^1_{A/\mathbf F_p}$), on déduit de (20.5.14) que l'on a
une \emph{suite exacte}
\begin{equation}
\label{0.22.4.6.1-fr}
  0\longrightarrow V\longrightarrow\Omega^1_A\otimes_AK
  \longrightarrow\Omega^1_K\longrightarrow0.
  \tag{22.4.6.1}
\end{equation}
On a donc par (22.4.5.3) un homomorphisme noté
\begin{equation}
\label{0.22.4.6.2-fr}
  \chi'_{B/A}:\Xi_{B_{(K)}/K}\longrightarrow V
  \tag{22.4.6.2}
\end{equation}
que nous appellerons \emph{homomorphisme caractéristique} de la $A$-algèbre
$B$ (relative à $i$, $j$ et $\mathfrak m$). Il résulte des définitions que pour
que $\chi'_{B/A}$ soit \emph{injectif}, il faut et il suffit que
$\pi'_{B_{(K)}/K}$ soit injectif, le noyau de ce dernier étant évidemment
contenu dans $\Xi_{B_{(K)}/K}$.
\end{env}

\begin{proposition}[22.4.7]
\label{0.22.4.7-fr}
Soient $A$ un anneau local artinien dont l'idéal maximal $\mathfrak m$ est de
carré nul, $K=A/\mathfrak m$ son corps résiduel, $V$ l'idéal $\mathfrak m$
considéré comme $K$-espace vectoriel~; on suppose que $A$ est de
caractéristique $p>0$~; soient
\begin{equation}
\label{0.22.4.7.1-fr}
  F_i(T_i)=T_i^p-f_i\qquad\text{avec }f_i\in A\quad(1\leqslant i\leqslant r)
  \tag{22.4.7.1}
\end{equation}
et désignons par $B$ l'anneau quotient de $A[T_1,\ldots,T_r]$ par l'idéal
engendré par les $r$ polynômes $F_i$. Alors~:
\begin{enumerate}
  \item[\rm(i)] $B$ est un anneau local.
  \item[\rm(ii)] Si $\mathfrak m'$ est l'idéal maximal de $B$,
  $K'=B/\mathfrak m'$ son corps résiduel, $V'$ le $B$-module
  $\mathfrak m'/\mathfrak m'^2$ considéré comme $K'$-espace vectoriel, alors,
  pour que $\operatorname{rg}_K(V)=\operatorname{rg}_{K'}(V')$, il faut et il suffit que l'homomorphisme
  caractéristique de $K$-espaces vectoriels
  \begin{equation}
  \label{0.22.4.7.2-fr}
    \chi'_{B/A}:\Xi_{B_{(K)}/K}\longrightarrow V
    \tag{22.4.7.2}
  \end{equation}
  (cf. (22.4.6.2)) soit injectif.
\end{enumerate}
\end{proposition}

\begin{proof}
Désignons encore par $a_i$ $(1\leqslant i\leqslant r)$ la classe de $f_i$
mod.~$\mathfrak m$~; les $a_i$ ne sont plus nécessairement $p$-indépendants
sur $K^p$, mais on peut toujours supposer que la sous-famille

\oldpage[0\textsubscript{IV-1}]{199}

$(a_i)_{s+1\leqslant i\leqslant r}$ est $p$-libre sur $K^p$ et forme une
$p$-base sur $K^p$ du corps $K^p(a_1,\ldots,a_r)$ de sorte que l'on a
\[
  a_i\in K^p(a_{s+1},\ldots,a_r)\qquad\text{pour }1\leqslant i\leqslant s.
\]
Désignons par $A''$ l'anneau quotient de $A[T_{s+1},\ldots,T_r]$ par l'idéal
engendré par les $F_i$ d'indice $i\geqslant s+1$~; alors, par (22.4.4), $A''$
est un anneau local dont l'idéal maximal est $\mathfrak m''=\mathfrak mA''$
(donc de carré nul) et le corps résiduel
$K''=K(a_{s+1}^{1/p},\ldots,a_r^{1/p})$. On a donc
$a_i\in K''^p$ pour $1\leqslant i\leqslant s$ et par suite, dans $A''$,
l'élément $f_i$, pour $1\leqslant i\leqslant s$, peut s'écrire
\begin{equation}
\label{0.22.4.7.3-fr}
  f_i=g_i^p+h_i
  \tag{22.4.7.3}
\end{equation}
où $g_i\in A''$ et $h_i\in\mathfrak m''$ (comme $\mathfrak m''^2=0$ et
$p\mathfrak m''=0$, il est immédiat que $h_i$ est déterminé de façon
\emph{unique} par ces conditions). Remplaçant $T_i$ par $T_i+g_i$, on voit donc
que si l'on pose
\[
  G_i(T_i)=T_i^p+\sum_{k=1}^{p-1}\binom pk g_i^{p-k}T_i^k-h_i
  \qquad(1\leqslant i\leqslant s),
\]
$B$ est l'anneau quotient de $A''[T_1,\ldots,T_s]$ par l'idéal engendré par les
$G_i$~; or, tous les coefficients de $G_i$ sauf le coefficient dominant
appartiennent à $\mathfrak m''$, donc on est dans la situation de (22.4.3),
$B$ est un anneau local, et son corps résiduel $K'$ est \emph{isomorphe} à
$K''$, ce qui prouve (i) (qui d'ailleurs résulte directement de ce que
$B^p\subset A$, et par suite il n'y a qu'un seul idéal de $B$ au-dessus de
$\mathfrak m$).

Notons maintenant que pour que (22.4.7.2) soit injectif, il faut et il suffit
qu'il en soit ainsi de
\[
  \pi'_{B_{(K)}/K}:\Omega^1_{B_{(K)}/K}\otimes_{B_{(K)}}K
  \longrightarrow\Omega^1_A\otimes_AK
\]
(par (22.4.6)). Or $B_{(K)}=B/\mathfrak mB$ est l'algèbre quotient de
$C=K[T_1,\ldots,T_r]$ par l'idéal engendré par les polynômes
$\overline F_i=T_i^p-a_i$, et comme $d_{C/K}\overline F_i=0$ il résulte de
(20.5.13) que $\Omega^1_{B_{(K)}/K}$ est un $B_{(K)}$-module libre ayant pour
\emph{base} les $d_{B_{(K)}/K}t_i$, où $t_i$ est l'image canonique de $T_i$
$(1\leqslant i\leqslant r)$. Mais par définition (22.4.5) l'image canonique par
$j':B_{(K)}\longrightarrow A$ d'une classe mod.~$\mathfrak mB$ d'un élément
$x\in B$ est l'élément $x^p\in A$~; on en déduit aussitôt, puisque $t_i^p=f_i$,
que l'image de $d_{B_{(K)}/K}t_i\otimes1$ par $\pi'_{B_{(K)}/K}$ est
$d_Af_i\otimes1$ dans $\Omega^1_A\otimes_AK$~; la condition que
$\pi'_{B_{(K)}/K}$ soit injectif est donc équivalente au fait que les
$d_Af_i\otimes1_K$ sont \emph{linéairement indépendants sur $K$}, ou encore
que leurs images $d_Af_i\otimes1_{K'}$ sont \emph{linéairement indépendants
sur $K'$} dans $\Omega^1_A\otimes_AK'$.

Or, si l'on applique (22.4.4) à la $A$-algèbre $A''$, on voit que le noyau de
l'homomorphisme canonique
$\Omega^1_A\otimes_AK'\longrightarrow\Omega^1_{A''}\otimes_{A''}K'$ a pour
base les $d_Af_i\otimes1_{K'}$ pour $s+1\leqslant i\leqslant r$. La condition
précédente équivaut donc encore à dire que les $d_{A''}f_i\otimes1_{K'}$ sont
linéairement indépendants dans $\Omega^1_{A''}\otimes_{A''}K'$ pour
$1\leqslant i\leqslant s$. Or, on a $d_{A''}f_i=d_{A''}h_i$ par (22.4.7.3)~;
d'autre part, si $V''$ est le $K'$-espace vectoriel $\mathfrak m''$, on a
$\operatorname{rg}_{K'}(V'')=\operatorname{rg}_K(V)$ (22.4.1, (ii)), et la condition
$\operatorname{rg}_{K'}(V')=\operatorname{rg}_K(V)$ est donc équivalente à
$\operatorname{rg}_{K'}(V'')=\operatorname{rg}_{K'}(V')$. Mais il résulte de (22.4.3) que cette dernière
relation est équivalente au fait que les classes $h_i$

\oldpage[0\textsubscript{IV-1}]{200}

dans $V''$ sont \emph{linéairement indépendantes sur $K'$}. Or, dans
$\Omega^1_{A''}\otimes_{A''}K'$, les $d_{A''}h_i\otimes1_{K'}$ sont les images
canoniques (20.5.11.2) des $h_i$ par l'injection
$V''\longrightarrow\Omega^1_{A''}\otimes_{A''}K'$ (cf. (22.4.6.1)), et cela
achève de prouver (22.4.7).
\end{proof}

On a en outre prouvé~:

\begin{corollary}[22.4.7.4]
\label{0.22.4.7.4-fr}
Pour que $\operatorname{rg}_K(V)=\operatorname{rg}_{K'}(V')$, il faut et il suffit que les éléments
$d_Af_i\otimes1_K$ $(1\leqslant i\leqslant r)$ soient linéairement
indépendants dans le $K$-espace vectoriel $\Omega^1_A\otimes_AK$.
\end{corollary}

\begin{remark}[22.4.8]
\label{0.22.4.8-fr}
On peut définir les homomorphismes (22.4.5.2) et (22.4.6.2) dans des conditions
un peu différentes. Soient $A$ un anneau, $\mathfrak m$ un idéal \emph{de carré
nul} dans $A$, $K=A/\mathfrak m$, $p$ un nombre premier tel que
$p\cdot1\in\mathfrak m$ (autrement dit $K$ est de caractéristique $p$, mais non
nécessairement $A$). Soit d'autre part $B$ une $A$-algèbre qui est un $A$-module
\emph{fidèlement plat} (de sorte que $A\subset B$), et supposons que l'on ait
\begin{equation}
\label{0.22.4.8.1-fr}
  y^p\in A+pB\subset A+\mathfrak mB
  \tag{22.4.8.1}
\end{equation}
pour tout $y\in B$. Si $y^p=x+pz$ avec $x\in A$, $z\in B$, la \emph{classe}
mod.~$(A\cap pB)$ de $x$ est bien déterminée, et comme $A\cap pB=pA$ puisque
$B$ est un $A$-module fidèlement plat, on a ainsi défini une application
canonique $j:B\longrightarrow A/pA$, telle que, pour $x\in A$, $j(x)$ soit la
classe de $x^p$ mod.~$pA$. D'ailleurs, si $y'\in B$ et
$y'^p=x'+pz'$ avec $x'\in A$, $z'\in B$, il est immédiat que les éléments
$(y+y')^p-x-x'$ et $(yy')^p-xx'$ appartiennent à $pB$ en vertu du fait que les
coefficients binomiaux $\binom pi$ sont des multiples de $p$ pour
$1\leqslant i\leqslant p-1$. L'application $j$ est donc un
\emph{homomorphisme d'anneaux}. Par composition, on déduit de $j$ un
homomorphisme $j':B\xrightarrow{j}A/pA\longrightarrow A/\mathfrak m=K$ et par
suite un homomorphisme $j_0:B_{(K)}\longrightarrow K$ tel que $j'$ soit le
composé $B\longrightarrow B_{(K)}=B\otimes_AK\xrightarrow{j_0}K$~; en outre,
si $i_0:K\longrightarrow B_{(K)}$ est l'application canonique, $i_0$ et $j_0$
vérifient les conditions de (21.3.1) pour les anneaux \emph{de caractéristique
$p$}, $K$ et $B_{(K)}$. Cela étant, pour définir un $K$-homomorphisme
$\pi'_{B_{(K)}/K}:\Omega^1_{B_{(K)}/K}\otimes_{B_{(K)}}K_{[j_0]}
\longrightarrow\Omega^1_A\otimes_AK$, il suffit (20.4.8) de définir une
$K$-dérivation
\begin{equation}
\label{0.22.4.8.2-fr}
  D_0:B_{(K)}\longrightarrow\Omega^1_A\otimes_AK
  \tag{22.4.8.2}
\end{equation}
où le second membre est considéré comme un $B_{(K)}$-module au moyen de $j_0$.
Pour cela, il suffit de définir une $A$-dérivation
\begin{equation}
\label{0.22.4.8.3-fr}
  D:B\longrightarrow\Omega^1_A\otimes_AK
  \tag{22.4.8.3}
\end{equation}
qui \emph{s'annule dans} $\mathfrak mB$ (le second membre étant considéré comme
$B$-module au moyen de $j'$). Or, si l'on a la relation $y^p=x+pz$ avec
$y\in B$, $x\in A$, $z\in B$, l'élément $d_Ax\otimes1_K$ de
$\Omega^1_A\otimes_AK$ ne dépend que de la classe $j(y)$ mod.~$pA$ de $x$
(donc seulement de $y$), car pour tout $t\in A$ on a
$d_A(pt)=p\,d_A(t)$, donc $d_A(pt)\otimes1=0$ dans
$\Omega^1_A\otimes_AK=\Omega^1_A/\mathfrak m\Omega^1_A$, puisque
$p\cdot1\in\mathfrak m$. On peut donc prendre
\[
  D(y)=d_Ax\otimes1_K.
\]
Le fait que $D$ est une dérivation résulte aisément de ce que $j$ est un
homomorphisme d'anneaux~; en outre, pour $t\in A$, on a
$d_A(t^px)=t^p d_A(x)+pt^{p-1}x d_A(t)$ et l'on en tire que
$D(ty)=tD(y)$, donc $D$ est une $A$-dérivation~; enfin, si l'on a de plus
$t\in\mathfrak m$, on a

\oldpage[0\textsubscript{IV-1}]{201}

$D(ty)=tD(y)=0$ puisque $\Omega^1_A\otimes_AK$ est annulé par $\mathfrak m$.
On a ainsi défini l'homomorphisme (22.4.5.2) cherché, et il est immédiat de
vérifier que l'homomorphisme composé
\[
  \Theta_{B_{(K)}/K}\xrightarrow{\pi'_{B_{(K)}/K}}
  \Omega^1_A\otimes_AK\longrightarrow\Omega^1_K
\]
est encore l'homomorphisme canonique $\pi_{B_{(K)}/K}$ (relatif à $i_0$ et
$j_0$). On a donc aussi un homomorphisme canonique analogue à (22.4.5.3).

Si maintenant $K$ est un \emph{corps}, $A/pA$ est une $\mathbf F_p$-algèbre,
et l'on a la suite exacte (20.5.14)
\[
  0\longrightarrow V/pA\longrightarrow
  \Omega^1_{A/pA}\otimes_{A/pA}K\longrightarrow\Omega^1_K\longrightarrow0
\]
et d'autre part on a la suite exacte (20.5.12.1)
\[
  pA\longrightarrow\Omega^1_A\otimes_A(A/pA)
  \longrightarrow\Omega^1_{A/pA}\longrightarrow0
\]
qui montre, en tensorisant avec $K$, que les $K$-espaces vectoriels
$\Omega^1_A\otimes_AK$ et
$\Omega^1_{A/pA}\otimes_{A/pA}K$ sont isomorphes. L'analogue de (22.4.6.2) est
donc ici un homomorphisme
\begin{equation}
\label{0.22.4.8.4-fr}
  \Xi_{B_{(K)}/K}\longrightarrow V/pA=\mathfrak m/(\mathfrak m^2+pA).
  \tag{22.4.8.4}
\end{equation}

Ceci dit, \emph{le critère (22.4.7), où l'on remplace $V$ et $V'$ par $V/pA$
et $V'/pB$ respectivement, reste valable en supposant seulement que $K$ soit
de caractéristique $p>0$} (autrement dit, que $p\cdot1\in\mathfrak m$ dans
$A$), \emph{et que les $F_i(T_i)$ soient de la forme plus générale
(22.4.4.1)}~: en effet, $B$ est un $A$-module libre (donc fidèlement plat), et
il suffit de reprendre le raisonnement de (22.4.7), en y remplaçant (22.4.6.2)
par (22.4.8.4) et $V''$ par $V''/pA''$.
\end{remark}

\subsection{Algèbres géométriquement régulières et algèbres formellement lisses}
\label{subsection:0.22.5-fr}

\begin{proposition}[22.5.1]
\label{0.22.5.1-fr}
Soient $A$ un anneau local régulier, $A'$ un anneau local contenant $A$ et qui
est une $A$-algèbre finie, $\mathfrak m$ (resp. $\mathfrak m'$) l'idéal maximal
de $A$ (resp. $A'$), $K=A/\mathfrak m$ (resp. $K'=A'/\mathfrak m'$) son corps
résiduel. Pour que $A'$ soit un anneau régulier, il faut et il suffit que l'on
ait
\begin{equation}
\label{0.22.5.1.1-fr}
  \operatorname{rg}_K(\mathfrak m/\mathfrak m^2)
  =\operatorname{rg}_{K'}(\mathfrak m'/\mathfrak m'^2).
  \tag{22.5.1.1}
\end{equation}
\end{proposition}

\begin{proof}
En effet, le premier membre de (22.5.1.1) est $\dim(A)$ (17.1.1) et comme
$A\subset A'$ et que $A'$ est une $A$-algèbre finie,
$\dim(A')=\dim(A)$ (16.1.5)~; l'égalité (22.5.1.1) est donc nécessaire et
suffisante pour que $A'$ soit régulier (17.1.1).
\end{proof}

\begin{remark}[22.5.2]
\label{0.22.5.2-fr}
\begin{enumerate}
  \item[\rm(i)] Soient $A_1=A/\mathfrak m^2$,
  $A'_1=A'\otimes_AA_1=A'/\mathfrak m^2A'$~; comme
  $\mathfrak m^2A'\subset\mathfrak m'^2$,
  $\mathfrak m'_1=\mathfrak m'/\mathfrak m^2A'$ est l'idéal maximal de $A'_1$
  et $\mathfrak m'_1/\mathfrak m_1'^2$ est un $K'$-espace vectoriel isomorphe
  à $\mathfrak m'/\mathfrak m'^2$~; la condition de régularité de $A'$ ne
  dépend donc que de la structure de la $A_1$-algèbre $A'_1$, ce qui va nous
  permettre ci-dessous d'appliquer les résultats préliminaires de (22.4) sur
  les algèbres finies sur des anneaux artiniens.
  \item[\rm(ii)] Il résulte de la démonstration de (22.5.1) et de (16.2.6.1)
  que l'on a en tout cas
  \begin{equation}
  \label{0.22.5.2.1-fr}
    \operatorname{rg}_K(\mathfrak m/\mathfrak m^2)
    \leqslant\operatorname{rg}_{K'}(\mathfrak m'/\mathfrak m'^2).
    \tag{22.5.2.1}
  \end{equation}

  \oldpage[0\textsubscript{IV-1}]{202}

  \item[\rm(iii)] Supposons que $A$ et $A'$ vérifient les hypothèses générales
  de (22.4.1). On sait (19.8.10) qu'il existe un anneau local \emph{régulier}
  $B$, d'idéal maximal $\mathfrak n$, tel que $A$ soit isomorphe à
  $B/\mathfrak n^2$~; désignons par $G_i\in B[T_i]$ un polynôme unitaire dont
  l'image canonique dans $A[T_i]$ soit $F_i$ $(1\leqslant i\leqslant r)$~;
  alors, si $B'$ est l'anneau quotient de $B[T_1,\ldots,T_r]$ par l'idéal
  engendré par les $G_i$, il est clair que $B'$ est un $B$-module libre de rang
  $d$ et que $A'=B'\otimes_BA=B'/\mathfrak n^2B'$~; comme $A'$ est supposé être
  un anneau local, il en est de même de $B'$, et si $\mathfrak n'$ est l'idéal
  maximal de $B'$, $B'/\mathfrak n'$ est isomorphe à $K'$ et
  $\mathfrak n'/\mathfrak n'^2$ isomorphe à $\mathfrak m'/\mathfrak m'^2$ en
  tant qu'espace vectoriel sur $K'$~; l'inégalité (22.5.2.1) montre donc que
  l'on a, sous les hypothèses de (22.4.1),
  \begin{equation}
  \label{0.22.5.2.2-fr}
    \operatorname{rg}_K(V)\leqslant\operatorname{rg}_{K'}(\mathfrak m'/\mathfrak m'^2).
    \tag{22.5.2.2}
  \end{equation}
\end{enumerate}
\end{remark}

\begin{corollary}[22.5.3]
\label{0.22.5.3-fr}
Soient $A$ un anneau local régulier, $\mathfrak m$ son idéal maximal,
$T_i$ $(1\leqslant i\leqslant r)$ des indéterminées, et pour chaque $i$, soit
\begin{equation}
\label{0.22.5.3.1-fr}
  F_i(T_i)=T_i^{d_i}+\sum_{1\leqslant k\leqslant d_i}c_{ik}T_i^{d_i-k}
  \tag{22.5.3.1}
\end{equation}
un polynôme unitaire de $A[T_i]$, tel que $d_i\geqslant2$ pour tout $i$ et
$c_{ik}\in\mathfrak m$. Soit $A'$ le quotient de l'anneau de polynômes
$A[T_1,\ldots,T_r]$ par l'idéal engendré par les $r$ polynômes $F_i$. Alors~:
\begin{enumerate}
  \item[\rm(i)] $A'$ est un anneau local, et son corps résiduel $K'$ est
  isomorphe au corps résiduel $K$ de $A$.
  \item[\rm(ii)] Pour que $A'$ soit un anneau régulier, il faut et il suffit
  que les classes $\xi_i$ mod.~$\mathfrak m^2$ des éléments $c_{i,d_i}$
  $(1\leqslant i\leqslant r)$ soient linéairement indépendantes sur $K$.
\end{enumerate}
\end{corollary}

\begin{proof}
En effet, avec les notations de (22.5.2, (i)), si $\overline F_i$ désigne le
polynôme de $A_1[T_i]$ obtenu en appliquant aux coefficients de $F_i$
l'homomorphisme $A\longrightarrow A_1$, $A'_1$ se définit à partir de $A_1$ et
des polynômes $\overline F_i$ comme dans (22.4.1), et la conclusion résulte
donc de (22.5.1) et de (22.4.3).
\end{proof}

\begin{theorem}[22.5.4]
\label{0.22.5.4-fr}
Soient $A$ un anneau local régulier, $K$ son corps résiduel~; on suppose que
$K$ est de caractéristique $p>0$~; soient
\begin{equation}
\label{0.22.5.4.1-fr}
  F_i(T_i)=T_i^p+p\sum_{k=1}^{p-1}c_{ik}T_i^{p-k}-f_i
  \tag{22.5.4.1}
\end{equation}
avec $c_{ik}\in A$ et $f_i\in A$ $(1\leqslant i\leqslant r)$~; soit $B$
l'anneau quotient de $A[T_1,\ldots,T_r]$ par l'idéal engendré par les $r$
polynômes $F_i$. Alors $B$ est un anneau local, et les conditions suivantes
sont équivalentes~:
\begin{enumerate}
  \item[a)] $B$ est régulier~;
  \item[b)] les éléments $d_Af_i\otimes1_K$ $(1\leqslant i\leqslant r)$ sont
  linéairement indépendants dans le $K$-espace vectoriel
  $\Omega^1_A\otimes_AK$~;
  \item[c)] si $B'=B/\mathfrak m^2B$ et si l'on pose
  $\Xi_{B_{(K)}/K}=\Xi_{B'_{(K)}/K}$ (22.4.5), l'homomorphisme
  caractéristique $\Xi_{B_{(K)}/K}\longrightarrow\mathfrak m/\mathfrak m^2$
  (22.4.6.2) est injectif.
\end{enumerate}
\end{theorem}

\begin{proof}
En effet, $B'$ se définit encore à partir de $A_1=A/\mathfrak m^2$ et des
polynômes $\overline F_i$ obtenus en appliquant aux $F_i$ l'homomorphisme
$A\longrightarrow A_1$. Compte tenu de ce que $\Omega^1_A\otimes_AK$ est
canoniquement isomorphe à $\Omega^1_{A_1}\otimes_{A_1}K$ (20.5.12, (ii)), le
théorème résulte de (22.5.1) et de (22.4.7.4), en utilisant la remarque
(22.4.8) lorsque $A_1$ n'est pas de caractéristique $p$.
\end{proof}

\oldpage[0\textsubscript{IV-1}]{203}

\begin{env}[22.5.5]
\label{0.22.5.5-fr}
Soient $k$ un corps de caractéristique $p>0$, $A$ un anneau local, $k'$ une
extension \emph{radicielle finie} de $k$ telle que $k'^p\subset k$,
$\mathfrak m$ l'idéal maximal de $A$, $K=A/\mathfrak m$ son corps résiduel~;
rappelons que $A'=A\otimes_kk'$ est un anneau local dont le corps résiduel est
le même que celui de l'anneau local $K\otimes_kk'$ (Bourbaki,
\emph{Alg. comm.}, chap.~V, §~2, n\textsuperscript{o}~3, lemme~4). Notons que
l'on a $A'_{(K)}=A'\otimes_AK=K\otimes_kk'$, et par suite (21.3.6)
\begin{equation}
\label{0.22.5.5.1-fr}
  \Theta_{A'_{(K)}/K}=\Theta_{k'/k}\otimes_kK
  \tag{22.5.5.1}
\end{equation}
à un isomorphisme canonique près. En outre, on sait (21.4.7) que l'on a
$\Xi_{k'/k}=0$, autrement dit l'homomorphisme canonique (21.3.2.2)
\[
  \pi_{k'/k}:\Theta_{k'/k}\longrightarrow\Omega^1_k
\]
est injectif, si bien que l'on a une \emph{injection canonique}
\begin{equation}
\label{0.22.5.5.2-fr}
  \Theta_{A'_{(K)}/K}\longrightarrow\Omega^1_k\otimes_kK
  \tag{22.5.5.2}
\end{equation}
qui s'explicite de la façon suivante (compte tenu de (22.5.5.1))~: pour tout
$x'\in k'$, à l'élément
$d_{A'_{(K)}/K}(x'\otimes1_K)\otimes1_K$ de
$\Omega^1_{A'_{(K)}/K}\otimes_{A'_{(K)}}K$, on fait correspondre l'élément
$d_k(x'^p)\otimes1_K$ de $\Omega^1_k\otimes_kK$.
\end{env}

\begin{lemma}[22.5.6]
\label{0.22.5.6-fr}
Avec les notations de (22.5.5)~:
\begin{enumerate}
  \item[\rm(i)] Le noyau de l'homomorphisme canonique
  $\pi_{A'_{(K)}/K}$ (21.3.2.2) s'identifie par (22.5.5.2) à
  $(\Theta_{k'/k}\otimes_kK)\cap\Upsilon_{K/k}$.
  \item[\rm(ii)] Soient $A_1=A/\mathfrak m^2$,
  $A'_1=A_1\otimes_kk'=A'/\mathfrak m^2A'$~; alors l'homomorphisme
  caractéristique $\chi'_{A'_1/A_1}$ (22.4.6.2) s'identifie à la restriction de
  l'homomorphisme caractéristique
  $\chi_{A/k}:\Upsilon_{K/k}\longrightarrow V=\mathfrak m/\mathfrak m^2$
  (20.6.24) à $(\Theta_{k'/k}\otimes_kK)\cap\Upsilon_{K/k}$.
\end{enumerate}
\end{lemma}

\begin{proof}
(i) En vertu de (20.6.9), appliqué en remplaçant $A$ par le corps premier, $B$
par $k$, $C$ par $K$ et $E$ par $A_1=A/\mathfrak m^2$ (extension de $K$ par
$V=\mathfrak m/\mathfrak m^2$), il suffit (cf. (22.5.5)) de vérifier que si
$x'\in k'$ et si $z\in A_1$ est un élément dont la classe mod.
$\mathfrak m/\mathfrak m^2$ est $x'^p\in k$, alors l'image de
$d_k(x'^p)\otimes1_K$ par l'homomorphisme canonique
$\Omega^1_k\otimes_kK\longrightarrow\Omega^1_{A_1}\otimes_{A_1}K$ est
$d_{A_1}z\otimes1_K$, ce qui résulte des définitions.

(ii) Comme $A'_{(K)}=K\otimes_kk'=(A'_1)_{(K)}$ (en vertu de la relation
$K=A_1/(\mathfrak m/\mathfrak m^2)$), $\Theta_{A'_{(K)}/K}$ s'identifie à
$\Theta_{(A'_1)_{(K)}/K}$ et la vérification résulte du même calcul que dans
(i) et de la définition (22.4.5.2) de $\pi'_{(A'_1)_{(K)}/K}$.
\end{proof}

\begin{proposition}[22.5.7]
\label{0.22.5.7-fr}
Soient $k$ un corps de caractéristique $p>0$, $A$ une $k$-algèbre qui est un
anneau local régulier, $\mathfrak m$ son idéal maximal, $K=A/\mathfrak m$ son
corps résiduel~; soit $V=\mathfrak m/\mathfrak m^2$, considéré comme
$K$-espace vectoriel. Soient $k'$ une extension finie de $k$ telle que
$k'^p\subset k$, et $A'=A\otimes_kk'$~; pour que l'anneau local $A'$ soit
régulier, il faut et il suffit que la restriction de l'homomorphisme
caractéristique $\chi_{A/k}:\Upsilon_{K/k}\longrightarrow V$ (20.6.24) à
$(\Theta_{k'/k}\otimes_kK)\cap\Upsilon_{K/k}$ soit injective.
\end{proposition}

\begin{proof}
Il s'agit de prouver que la condition de l'énoncé est équivalente à
(22.5.1.1), en désignant par $\mathfrak m'$ l'idéal maximal de $A'$, par
$K'=A'/\mathfrak m'$ son corps résiduel. Avec les notations de la remarque
(22.5.2, (i)), on a $A'_1=A_1\otimes_kk'$. Or, soit
$(a_i)_{1\leqslant i\leqslant r}$ une $p$-base de $k'$ sur $k$, de sorte que
$a_i^p=b_i\in k$, et $k'$ est isomorphe au quotient de l'anneau de

\oldpage[0\textsubscript{IV-1}]{204}

polynômes $k[T_1,\ldots,T_r]$ par l'idéal engendré par les polynômes
$F_i=T_i^p-b_i$ $(1\leqslant i\leqslant r)$~; on en déduit que $A'_1$ est
isomorphe au quotient de l'anneau de polynômes $A_1[T_1,\ldots,T_r]$ par
l'idéal engendré par les $F_i$. On peut alors appliquer le critère (22.4.7), et
la condition de l'énoncé équivaut donc au fait que $\chi'_{A'_1/A_1}$ soit
injectif~; on conclut à l'aide de (22.5.6).
\end{proof}

Nous pouvons maintenant démontrer la réciproque de (19.6.5)~:

\begin{theorem}[22.5.8]
\label{0.22.5.8-fr}
Soient $k$ un corps, $p$ son exposant caractéristique, $A$ une $k$-algèbre
locale noethérienne. Les conditions suivantes sont équivalentes~:
\begin{enumerate}
  \item[a)] $A$ est une $k$-algèbre formellement lisse (pour sa topologie
  préadique).
  \item[b)] $A$ est géométriquement régulier sur $k$.
  \item[b')] Pour toute extension finie $k'$ de $k$ telle que $k'^p\subset k$,
  $A'=A\otimes_kk'$ est un anneau régulier.
\end{enumerate}
\end{theorem}

\begin{proof}
Compte tenu de (19.6.6), on peut se borner au cas $p>1$.

Le fait que a) implique b) n'est autre que (19.6.5), et b) entraîne
trivialement b'). Montrons que b') entraîne que les conditions de (22.2.2) sont
satisfaites~; c'est évident pour la première (17.1.1) puisque b') entraîne
d'abord que $A$ est régulier. D'autre part, soit $(x_\alpha)_{\alpha\in I}$ une
$p$-base de $k^{1/p}$ sur $k$~; on sait que les éléments $d_k(x_\alpha^p)$
forment une base du $k$-espace vectoriel $\Omega^1_k$ (21.4.5), donc les
éléments $d_k(x_\alpha^p)\otimes1_K$ forment une base du $K$-espace vectoriel
$\Omega^1_k\otimes_kK$. On en conclut (cf. (22.5.5)) que lorsque $k'$ parcourt
l'ensemble des sous-extensions de $k^{1/p}$, finies sur $k$, la famille des
sous-espaces $\Theta_{k'/k}\otimes_kK$ de $\Omega^1_k\otimes_kK$ est filtrante
croissante et a pour réunion $\Omega^1_k\otimes_kK$. Il résulte alors de
(22.5.7) que la condition b') entraîne que $\chi_{A/k}$ est injectif, ce qui
n'est autre que la condition (ii) de (22.2.2).
\end{proof}

\begin{corollary}[22.5.9]
\label{0.22.5.9-fr}
Soient $k$ un corps, $A$ une $k$-algèbre locale noethérienne formellement
lisse. Alors, pour tout idéal premier $\mathfrak p$ de $A$, l'anneau local
$A_{\mathfrak p}$ est une $k$-algèbre formellement lisse (pour la topologie
$\mathfrak p$-préadique).
\end{corollary}

\begin{proof}
En effet, avec les notations de (22.5.8, b')), il suffit de montrer que
l'anneau local $A_{\mathfrak p}\otimes_kk'$ est régulier~; mais comme c'est un
anneau de fractions de l'anneau local $A\otimes_kk'$ et que ce dernier est
régulier par hypothèse, il en est de même de $A_{\mathfrak p}\otimes_kk'$
(17.3.6).
\end{proof}

\begin{remark}[22.5.10]
\label{0.22.5.10-fr}
\begin{enumerate}
  \item[\rm(i)] La conclusion de (22.5.8) est en défaut lorsque, dans la
  condition b'), on se borne aux extensions \emph{monogènes} $k'=k(x)$ de $k$
  (avec $x^p\in k$). Il se peut en effet que pour toutes ces extensions $k'$,
  on ait
  $\Upsilon_{K/k}\cap(\Theta_{k'/k}\otimes_kK)=0$, bien que
  $\Upsilon_{K/k}\neq0$~; cela signifie que pour tout $x\in k^{1/p}$ tel que
  $x\notin k$, on doit avoir
  $d_k(x^p)\otimes1_K\notin\Upsilon_{K/k}$, ou encore $d_K(x^p)\neq0$,
  c'est-à-dire $x^p\notin K^p$~; autrement dit, on doit avoir
  $K^p\cap k=k^p$, bien que $K$ soit une extension \emph{inséparable} de $k$.
  Or, on construit aisément des exemples de telles extensions~: partons d'un
  corps parfait $k_0$ de caractéristique $p>0$, soient $X$, $Y$, $Z$ trois
  indéterminées et posons $k=k_0(X,Y)$ et
  $K=k(X^{1/p}+ZY^{1/p})$~; on vérifie aisément que $K$ vérifie les conditions
  précédentes, donc $\Upsilon_{K/k}\neq0$. Appliquons maintenant le lemme
  (22.2.5.2) en prenant $V=K$ et pour $h$ l'homomorphisme nul~; $A$ est alors
  un anneau de valuation discrète ayant $K$ pour corps résiduel, avec
  $\chi_{A/k}=0$, et n'est donc pas une $k$-algèbre formellement lisse par
  (22.2.2)~; toutefois $A\otimes_kk'$ est un anneau régulier pour toute
  extension \emph{monogène} $k'\subset k^{1/p}$ de $k$.

  \oldpage[0\textsubscript{IV-1}]{205}

  \item[\rm(ii)] Le corollaire (22.5.9) amène à considérer la question
  suivante~: si $A$ est une $k$-algèbre noethérienne, à quelles conditions
  l'ensemble des idéaux premiers $\mathfrak p\in\operatorname{Spec}(A)$ tels que
  $A_{\mathfrak p}$ soit une $k$-algèbre formellement lisse est-il
  \emph{ouvert}~? Nous en aborderons plus tard certains cas particuliers.
\end{enumerate}
\end{remark}

\subsection{Critère jacobien de Zariski}
\label{subsection:0.22.6-fr}

\begin{theorem}[22.6.1]
\label{0.22.6.1-fr}
\textnormal{(critère jacobien de lissité formelle).}
Soient $A$, $B$ deux anneaux topologiques,
$u:A\longrightarrow B$ un homomorphisme continu faisant de $B$ une
$A$-algèbre formellement lisse, $\mathfrak J$ un idéal de $B$ (non
nécessairement fermé), $C=B/\mathfrak J$ la $A$-algèbre topologique quotient.
Pour que $C$ soit une $A$-algèbre formellement lisse, il faut et il suffit que
l'homomorphisme canonique (cf. (20.5.11.2))
\begin{equation}
\label{0.22.6.1.1-fr}
  \delta_{C/B/A}:\mathfrak J/\mathfrak J^2
  \longrightarrow\Omega^1_{B/A}\otimes_BC
  \tag{22.6.1.1}
\end{equation}
soit \emph{formellement inversible à gauche} (cf. (19.1.5)).
\end{theorem}

\begin{proof}
En effet, dire que $B$ (resp. $C$) est une $A$-algèbre formellement lisse
signifie (19.4.4) que l'on a
$\operatorname{Exalcotop}_A(B,L)=0$ (resp.
$\operatorname{Exalcotop}_A(C,L)=0$) pour tout $B$-module (resp. $C$-module)
discret $L$ annulé par un idéal ouvert de $B$ (resp. $C$). Comme par hypothèse
$\operatorname{Exalcotop}_A(B,L)=0$ pour tout $C$-module discret $L$ annulé
par un idéal ouvert de $C$, dire que $C$ est une $A$-algèbre formellement lisse
signifie donc que l'homomorphisme canonique
$\operatorname{Exalcotop}_A(C,L)\longrightarrow
\operatorname{Exalcotop}_A(B,L)$ est injectif, et le théorème résulte donc de
(20.7.8).
\end{proof}

Rappelons (19.5.3) que lorsque $B$ et $C$ sont des $A$-algèbres formellement
lisses, $\mathfrak J/\mathfrak J^2$ est un $C$-module topologique
\emph{formellement projectif}~; en outre, les homomorphismes canoniques
$\varphi_n:\mathbf S_C^n(\mathfrak J/\mathfrak J^2)
\longrightarrow\mathfrak J^n/\mathfrak J^{n+1}$ sont des
\emph{bimorphismes formels} lorsque $B$ est supposé être un anneau
\emph{préadmissible}.
\begin{corollary}[22.6.2]
\label{0.22.6.2-fr}
Les hypothèses et notations étant celles de (22.6.1), supposons de plus que
dans $B$ le carré de tout idéal ouvert soit ouvert, et que sur $\mathfrak J$ la
topologie induite par celle de $B$ soit identique à la topologie déduite de
celle de $B$ (19.0) (on notera que ces deux conditions sont vérifiées lorsque
$B$ est noethérien et sa topologie préadique ($0_{\mathrm I}$, 7.3.2), ou si la
topologie de $B$ est la topologie $\mathfrak J$-préadique). Soit
$(\mathfrak b_\lambda)$ un système fondamental d'idéaux ouverts de $B$ et
posons $B_\lambda=B/\mathfrak b_\lambda$ pour tout $\lambda$. Alors~:
\begin{enumerate}
  \item[\rm(i)] Les conditions suivantes sont équivalentes~:
  \begin{enumerate}
    \item[a)] $C$ est une $A$-algèbre formellement lisse.
    \item[b)] Pour tout $\lambda$, l'homomorphisme de $B_\lambda$-modules
    \begin{equation}
    \label{0.22.6.2.1-fr}
      (\mathfrak J/\mathfrak J^2)\otimes_BB_\lambda
      \longrightarrow\Omega^1_{B/A}\otimes_BB_\lambda
      \tag{22.6.2.1}
    \end{equation}
    déduit de $\delta_{C/B/A}$ par tensorisation, est inversible à gauche
    (autrement dit un isomorphisme sur un facteur direct de
    $\Omega^1_{B/A}\otimes_BB_\lambda$).
  \end{enumerate}
  \item[\rm(ii)] Supposons de plus que l'anneau topologique $C$ soit
  préadmissible ($0_{\mathrm I}$, 7.1.2) et

  \oldpage[0\textsubscript{IV-1}]{206}

  soit $\mathfrak L$ un idéal de définition de $C$~; alors les conditions a)
  et b) de (i) équivalent aussi à~:
  \begin{enumerate}
    \item[c)] L'homomorphisme de $(C/\mathfrak L)$-modules
    \[
      (\mathfrak J/\mathfrak J^2)\otimes_C(C/\mathfrak L)
      \longrightarrow\Omega^1_{B/A}\otimes_B(C/\mathfrak L)
    \]
    est inversible à gauche.
  \end{enumerate}
\end{enumerate}
\end{corollary}

\begin{proof}
Les hypothèses entraînent que sur $\Omega^1_{B/A}$ la topologie est déduite de
celle de $B$ (20.4.5), donc (i) résulte de (22.6.1) et de (19.1.7). D'autre
part, rappelons (20.4.9) que puisque $B$ est une $A$-algèbre formellement
lisse, $\Omega^1_{B/A}$ est un $B$-module formellement projectif, donc
$\Omega^1_{B/A}\otimes_BB_\lambda$ est un $B_\lambda$-module projectif pour
tout $\lambda$ (19.2.4)~; l'équivalence de c) avec a) et b) lorsque $C$ est
préadmissible résulte alors de (19.1.9).
\end{proof}

En particulier~:

\begin{corollary}[22.6.3]
\label{0.22.6.3-fr}
Soient $A$ un anneau, $B$ une $A$-algèbre, $\mathfrak J$ un idéal de $B$,
$C=B/\mathfrak J$~; on suppose $A$, $B$, $C$ munis des topologies discrètes
et que $B$ est une $A$-algèbre formellement lisse. Pour que $C$ soit une
$A$-algèbre formellement lisse, il faut et il suffit que l'homomorphisme
canonique (22.6.1.1) soit inversible à gauche.
\end{corollary}

On notera que puisque toute $A$-algèbre $C$ est quotient d'une algèbre de
polynômes $A[T_\alpha]_{\alpha\in I}$, et que cette dernière est formellement
lisse (19.3.3), le critère (22.6.3) permet en principe de reconnaître si $C$
est une $A$-algèbre formellement lisse.

\begin{proposition}[22.6.4]
\label{0.22.6.4-fr}
Soient $A$ un anneau, $B$ une $A$-algèbre formellement lisse (pour les
topologies discrètes), $\mathfrak J$ un idéal de $B$, $C=B/\mathfrak J$~; on
suppose que $\mathfrak J/\mathfrak J^2$ est un $C$-module de type fini. Soient
$\mathfrak p$ un idéal premier de $C$, $k(\mathfrak p)$ le corps résiduel de
$C_{\mathfrak p}$, $\mathfrak p'$ l'image réciproque de $\mathfrak p$ dans
$B$, $\mathfrak q$ l'idéal premier de $A$ image réciproque de $\mathfrak p'$.
Les conditions suivantes sont équivalentes~:
\begin{enumerate}
  \item[a)] $C_{\mathfrak p}$ est une $A_{\mathfrak q}$-algèbre formellement
  lisse (pour les topologies discrètes).
  \item[a')] $C_{\mathfrak p}$ est une $A$-algèbre formellement lisse (pour
  les topologies discrètes).
  \item[b)] L'homomorphisme canonique
  \begin{equation}
  \label{0.22.6.4.1-fr}
    (\mathfrak J/\mathfrak J^2)\otimes_Ck(\mathfrak p)
    \longrightarrow\Omega^1_{B/A}\otimes_Bk(\mathfrak p)
    \tag{22.6.4.1}
  \end{equation}
  est injectif.
  \item[c)] Il existe $f\in C-\mathfrak p$ tel que $C_f$ soit une
  $A$-algèbre formellement lisse.
\end{enumerate}
Si de plus $B$ est noethérien, les conditions précédentes sont encore
équivalentes à~:
\begin{enumerate}
  \item[d)] $C_{\mathfrak p}$ est une $A_{\mathfrak q}$-algèbre formellement
  lisse pour la topologie $\mathfrak p$-préadique sur $C_{\mathfrak p}$ et la
  topologie discrète (ou la topologie $\mathfrak q$-préadique) sur
  $A_{\mathfrak q}$.
\end{enumerate}
\end{proposition}

\begin{proof}
Comme $A_{\mathfrak q}$ est une $A$-algèbre formellement lisse (19.3.5, (iv)
et (ii)), on sait déjà que a) et a') sont équivalentes (19.3.5, (ii)). On a
$C_{\mathfrak p}=B_{\mathfrak p'}/\mathfrak J B_{\mathfrak p'}$, et
$B_{\mathfrak p'}$ est une $A$-algèbre formellement lisse pour la topologie
discrète (19.3.5, (iv))~; en outre
$\Omega^1_{B_{\mathfrak p'}/A}=(\Omega^1_{B/A})_{\mathfrak p'}$ (20.5.9).
L'équivalence de a') et b) résulte alors de l'application de (22.6.3) à
$B_{\mathfrak p'}$ et $C_{\mathfrak p}$, en tenant compte de ce que
$\Omega^1_{B/A}$ est un $B$-module projectif (20.4.9) et (19.2.11), et que
$\mathfrak J/\mathfrak J^2$ est un $B$-module de type fini, et en utilisant
(19.1.12). L'application de (19.1.12) prouve aussi l'équivalence de b) et c).
Enfin, notons que puisque $B_{\mathfrak p'}$ est une
$A_{\mathfrak q}$-algèbre formellement lisse pour les topologies

\oldpage[0\textsubscript{IV-1}]{207}

discrètes (19.3.5, (iv)), $B_{\mathfrak p'}$, muni de la topologie
$\mathfrak p'$-préadique, est encore une $A_{\mathfrak q}$-algèbre
formellement lisse lorsque l'on prend sur $A_{\mathfrak q}$ la topologie
discrète ou la topologie $\mathfrak q$-préadique (19.3.8). Pour montrer
l'équivalence de b) et d) lorsque $B$ est noethérien, appliquons (22.6.2) à
l'anneau $A_{\mathfrak q}$ discret (ou $\mathfrak q$-préadique) et à l'anneau
$B_{\mathfrak p'}$, muni de la topologie $\mathfrak p'$-préadique~; comme
$C_{\mathfrak p}$ est préadmissible et que $\mathfrak pC_{\mathfrak p}$ est
un idéal de définition pour sa topologie, on peut invoquer l'équivalence de
c) et a) dans (22.6.2).
\end{proof}

\begin{corollary}[22.6.5]
\label{0.22.6.5-fr}
Sous les hypothèses générales de (22.6.4), l'ensemble des
$\mathfrak p\in\operatorname{Spec}(C)$ tels que $C_{\mathfrak p}$ soit une $A$-algèbre
formellement lisse (ou une $A_{\mathfrak q}$-algèbre formellement lisse, en
désignant par $\mathfrak q$ l'image réciproque de $\mathfrak p$ dans $A$)
pour les topologies discrètes, est ouvert dans $\operatorname{Spec}(C)$.
\end{corollary}

Cela résulte de la forme c) de (22.6.4). On notera que lorsque $B$ est
noethérien, on peut remplacer les topologies discrètes par les topologies
préadiques.

\begin{corollary}[22.6.6]
\label{0.22.6.6-fr}
Sous les hypothèses générales de (22.6.4), les conditions suivantes sont
équivalentes~:
\begin{enumerate}
  \item[a)] $C$ est une $A$-algèbre formellement lisse (pour les topologies
  discrètes).
  \item[b)] Pour tout $\mathfrak p\in\operatorname{Spec}(C)$ (ou seulement pour tout idéal
  maximal $\mathfrak p$ de $C$), $C_{\mathfrak p}$ est une $A$-algèbre
  formellement lisse (ou une $A_{\mathfrak q}$-algèbre formellement lisse, en
  désignant par $\mathfrak q$ l'image réciproque de $\mathfrak p$ dans $A$)
  pour les topologies discrètes.
\end{enumerate}
En outre, lorsque $B$ est noethérien, on peut remplacer dans b) les topologies
discrètes par les topologies préadiques sur $A_{\mathfrak q}$ et
$C_{\mathfrak p}$.
\end{corollary}

\begin{proof}
La dernière assertion résulte de (22.6.4). En vertu de (22.6.3), la condition
a) équivaut au fait que l'homomorphisme (22.6.1.1) est inversible à gauche~;
l'équivalence de a) et b) résulte donc de (19.1.14), compte tenu du fait que
$\Omega^1_{B/A}$ est un $B$-module projectif et $\mathfrak J/\mathfrak J^2$
un $C$-module de type fini.
\end{proof}

\begin{proposition}[22.6.7]
\label{0.22.6.7-fr}
Soient $k_0$ un corps, $k$ une extension séparable de $k_0$, $C$ une
$k$-algèbre de type fini.
\begin{enumerate}
  \item[\rm(i)] Soit $\mathfrak p$ un idéal premier de $C$. Les conditions
  suivantes sont équivalentes~:
  \begin{enumerate}
    \item[a)] $C_{\mathfrak p}$ est une $k_0$-algèbre formellement lisse pour
    la topologie discrète.
    \item[b)] $C_{\mathfrak p}$ est une $k_0$-algèbre formellement lisse pour
    la topologie $\mathfrak p$-préadique.
    \item[c)] Il existe $f\in C-\mathfrak p$ tel que $C_f$ soit une
    $k_0$-algèbre formellement lisse pour la topologie discrète.
    \item[d)] $C_{\mathfrak p}$ est un anneau géométriquement régulier
    (19.6.5) sur $k_0$.
  \end{enumerate}
  En outre l'ensemble des $\mathfrak p\in\operatorname{Spec}(C)$ ayant l'une quelconque de
  ces propriétés est ouvert dans $\operatorname{Spec}(C)$.
  \item[\rm(ii)] Les conditions suivantes sont équivalentes~:
  \begin{enumerate}
    \item[a)] $C$ est une $k_0$-algèbre formellement lisse pour la topologie
    discrète.
    \item[b)] Tout $\mathfrak p\in\operatorname{Spec}(C)$ vérifie les conditions
    équivalentes de (i).
    \item[c)] Tout idéal maximal $\mathfrak m$ de $C$ vérifie les conditions
    équivalentes de (i).
    \item[d)] Pour toute extension $k'$ de $k_0$, $C\otimes_{k_0}k'$ est un
    anneau régulier (17.3.6)~; on dit encore que $C$ est un anneau
    géométriquement régulier sur $k_0$.
  \end{enumerate}

  \oldpage[0\textsubscript{IV-1}]{208}

  \item[\rm(iii)] Soient $B=k[T_1,\ldots,T_r]$ une algèbre de polynômes,
  $\mathfrak J$ un idéal de $B$ tels que $C$ soit isomorphe à
  $B/\mathfrak J$. Soit $\mathfrak q$ un idéal premier de $B$ contenant
  $\mathfrak J$ et soit $\mathfrak p=\mathfrak q/\mathfrak J$. Les conditions
  de (i) sont alors encore équivalentes à la suivante («~critère jacobien de
  Zariski~»)~:
  \begin{enumerate}
    \item[e)] Il existe un système d'éléments
    $(f_i)_{1\leq i\leq m}$ de $\mathfrak J$, engendrant l'idéal
    $\mathfrak J_{\mathfrak q}$ dans $B_{\mathfrak q}$ et des
    $k_0$-dérivations $D_j$ de $B$ dans lui-même $(1\leq j\leq m)$ tels que
    $\det(D_j(f_i))\notin\mathfrak q$.
  \end{enumerate}
\end{enumerate}
\end{proposition}

\begin{proof}
On sait que $C$ est toujours isomorphe à une $k$-algèbre de la forme
$B/\mathfrak J$. Comme $B$ est une $k$-algèbre formellement lisse (19.3.3) et
que $k$, étant séparable sur $k_0$, est une $k_0$-algèbre formellement lisse
(19.6.1), $B$ est une $k_0$-algèbre formellement lisse (19.3.5, (ii)).
L'équivalence des conditions a), b), c) de (i) résulte donc de (22.6.4), et
celle de a), b), c) dans (ii) résulte de (22.6.6)~; l'équivalence de a) et d)
dans (i) résulte de (22.5.8)~; comme tout localisé de $C\otimes_{k_0}k'$ est
aussi un localisé de $C_{\mathfrak p}\otimes_{k_0}k'$ pour un idéal premier
$\mathfrak p\in\operatorname{Spec}(C)$ convenable, l'équivalence de d) et b) dans (ii)
résulte de l'équivalence de d) et a) dans (i). Enfin, comme
$\Omega^1_{B/k_0}$ est un $B$-module projectif, l'équivalence du critère de
Zariski e) et des autres conditions de (i) découle de (19.1.12) et de
(22.6.3), compte tenu de (20.4.8).
\end{proof}

Zariski s'intéressait en fait à un critère différentiel de régularité pour les
anneaux locaux $C_{\mathfrak p}$. Comme il revient au même de dire qu'un anneau
local noethérien contenant un corps est régulier ou est formellement lisse en
tant qu'algèbre sur son sous-corps premier (19.6.4), on obtient aussitôt un tel
critère en remplaçant $k_0$ par le sous-corps premier de $k$ dans (22.6.7)~; en
particulier, on obtient le résultat suivant, que nous retrouverons plus tard
(IV, 6.12.5) comme cas particulier de résultats plus généraux de Nagata~:

\begin{corollary}[22.6.8]
\label{0.22.6.8-fr}
\textnormal{(Zariski).}
Soit $C$ une algèbre de type fini sur un corps. Alors l'ensemble des
$\mathfrak p\in\operatorname{Spec}(C)$ tels que $C_{\mathfrak p}$ soit un anneau local
régulier est ouvert dans $\operatorname{Spec}(C)$.
\end{corollary}

\begin{remark}[22.6.9]
\label{0.22.6.9-fr}
\begin{enumerate}
  \item[\rm(i)] Les énoncés de (22.6.7) sont encore valables si au lieu de
  supposer $k$ séparable sur $k_0$, on suppose seulement qu'il est de
  \emph{multiplicité radicielle finie}, autrement dit (19.6.6) qu'il existe
  une extension radicielle finie $k'_0$ de $k_0$ telle que le corps résiduel de
  l'anneau artinien local $k\otimes_{k_0}k'_0$ soit une extension séparable
  $k'$ de $k'_0$. Il existe alors un $k'_0$-monomorphisme
  $k'\longrightarrow k\otimes_{k_0}k'_0$ qui, composé avec l'homomorphisme
  canonique $k\otimes_{k_0}k'_0\longrightarrow k'$, donne l'identité, puisque
  $k'$ est une $k'_0$-algèbre formellement lisse (19.6.1)~; on en conclut que
  $C\otimes_{k_0}k'_0$ est muni d'une structure de $k'$-algèbre de type fini,
  et comme $k'$ est séparable sur $k'_0$, on peut appliquer (22.6.7) à cette
  $k'$-algèbre~; on en conclut notre assertion en appliquant (19.4.7). Comme
  nous n'aurons pas à nous servir de cette généralisation, nous en laissons le
  détail au lecteur. Nous ignorons par contre si les résultats de (22.6.7) sont
  valables sans \emph{aucune} hypothèse sur l'extension $k$ de $k_0$.
  \item[\rm(ii)] Sous les hypothèses générales de (22.6.7), soit $(k_\alpha)$
  une famille filtrante décroissante de sous-corps de $k$ contenant $k_0$ et
  telle que
  $\bigcap_\alpha k_\alpha(k^p)=k_0(k^p)$ (où $p$ est l'exposant
  caractéristique de $k_0$). Utilisant une méthode de décompte des dimensions
  due à Nagata, on peut montrer que, dans le critère jacobien (22.6.7, (iii)),
  on peut se borner à des dérivations $D_j$ de $B$ qui soient des
  $k_\alpha$-dérivations pour un $\alpha$ convenable.

  \oldpage[0\textsubscript{IV-1}]{209}

  L'intérêt de ce résultat est qu'il y a toujours de telles familles
  $(k_\alpha)$ pour lesquelles $[k:k_\alpha]$ soit fini (21.8.6). Nous ne
  démontrerons pas ce raffinement du critère de Zariski, dont nous n'aurons pas
  à faire usage. Dans (22.7), on va donner, pour les anneaux locaux complets,
  une variante (due aussi à Nagata) du critère de Zariski, qui se démontre
  essentiellement par la même méthode (avec des difficultés techniques un peu
  plus grandes).
\end{enumerate}
\end{remark}

\subsection{Le critère jacobien de Nagata}
\label{subsection:0.22.7-fr}

\begin{env}[22.7.1]
\label{0.22.7.1-fr}
Le critère jacobien de Nagata est l'analogue du critère jacobien de Zariski,
mais pour des anneaux quotients d'anneaux de \emph{séries formelles} sur un
corps. Nous en donnerons, comme Nagata [31], deux versions, présentées ici
comme critères de lissité formelle.
\end{env}

\begin{proposition}[22.7.2]
\label{0.22.7.2-fr}
Soient $k$ un corps, $A=k[[T_1,\ldots,T_r]]$ un anneau de séries formelles sur
$k$, muni de sa structure usuelle de $k$-algèbre, $\mathfrak q$ un idéal de
$A$, $B=A/\mathfrak q$. Soit $\mathfrak p$ un idéal premier de $A$ contenant
$\mathfrak q$. On suppose qu'il existe une sous-$k$-algèbre $C'$ de
$C=A/\mathfrak p$, isomorphe à une algèbre de séries formelles
$k[[X_1,\ldots,X_s]]$, telle que $C$ soit finie sur $C'$ et que le corps des
fractions $L$ de $C$ soit une extension séparable de
$L'=k((X_1,\ldots,X_s))$ (hypothèse toujours satisfaite si $k$ est de
caractéristique 0). Alors les conditions suivantes sont équivalentes~:
\begin{enumerate}
  \item[a)] $B_{\mathfrak p}=A_{\mathfrak p}/\mathfrak qA_{\mathfrak p}$ est
  une $k$-algèbre formellement lisse pour la topologie
  $\mathfrak p$-préadique.
  \item[b)] Il existe des $k$-dérivations $D_i$ $(1\leq i\leq m)$ de $A$ dans
  lui-même, et des éléments $f_i$ $(1\leq i\leq m)$ de $\mathfrak q$, tels
  que les images des $f_i$ dans $\mathfrak qA_{\mathfrak p}$ engendrent cet
  idéal de $A_{\mathfrak p}$, et que l'on ait
  $\det(D_i(f_j))\notin\mathfrak p$.
  \item[c)] $B_{\mathfrak p}$ est un anneau régulier.
\end{enumerate}
\end{proposition}

\begin{proof}
Le corps résiduel de $B_{\mathfrak p}$ est $k(\mathfrak p)$~; comme
$k((X_1,\ldots,X_s))$ est séparable sur $k$ (21.9.6.4), $k(\mathfrak p)$ est
séparable sur $k$ par hypothèse, et l'on sait déjà que dans ces conditions les
propriétés a) et c) sont équivalentes (19.6.4). D'ailleurs
$A_{\mathfrak p}$ est un anneau régulier ((17.1.4) et (17.3.2)), et comme son
corps résiduel $k(\mathfrak p)$ est séparable sur $k$, $A_{\mathfrak p}$ est
formellement lisse sur $k$ pour la topologie $\mathfrak p$-préadique
(19.6.4). Il résulte donc de (22.6.2, (ii)) que la condition a) équivaut au
fait que l'homomorphisme de $L$-espaces vectoriels
\begin{equation}
\label{0.22.7.2.1-fr}
  j_0:(\mathfrak q/\mathfrak q^2)\otimes_AL
  \longrightarrow\Omega^1_{A_{\mathfrak p}/k}\otimes_{A_{\mathfrak p}}L
  =\widehat\Omega^1_{A/k}\otimes_AL
  \tag{22.7.2.1}
\end{equation}
soit injectif.

Considérons d'autre part l'homomorphisme composé
\begin{equation}
\label{0.22.7.2.2-fr}
  j:(\mathfrak q/\mathfrak q^2)\otimes_AL
  \xrightarrow{j_0}\Omega^1_{A/k}\otimes_AL
  \longrightarrow\widehat\Omega^1_{A/k}\otimes_AL
  \tag{22.7.2.2}
\end{equation}
et montrons que la condition b) équivaut à dire que $j$ est injectif. Notons en
effet que
\[
  (\mathfrak q/\mathfrak q^2)\otimes_AL
  =(\mathfrak qA_{\mathfrak p}/\mathfrak q^2A_{\mathfrak p})
  \otimes_{A_{\mathfrak p}}L~;
\]
la condition b) signifie que si $F_i=j(f'_i\otimes1)$, où $f'_i$ est l'image
de $f_i$ dans
$\mathfrak qA_{\mathfrak p}/\mathfrak q^2A_{\mathfrak p}$, la matrice
$(\langle F_i,D_j\otimes1\rangle)$ est inversible~; donc cela entraîne que les
$F_i$ sont linéairement indépendants, et \emph{a fortiori} qu'il en est de
même des $f'_i\otimes1$~; mais comme ces derniers engendrent
$(\mathfrak qA_{\mathfrak p}/\mathfrak q^2A_{\mathfrak p})
\otimes_{A_{\mathfrak p}}L$, on en conclut que $j$ est injectif. Inversement,
supposons $j$ injectif, et prenons les $f_i$ tels que les $f'_i\otimes1$
forment

\oldpage[0\textsubscript{IV-1}]{210}

une base de
$(\mathfrak qA_{\mathfrak p}/\mathfrak q^2A_{\mathfrak p})
\otimes_{A_{\mathfrak p}}L$~; alors les $F_i$ forment une base de l'image de
$j$~; mais on sait (21.9.3) que $\widehat\Omega^1_{A/k}$ est un $A$-module
libre de rang $r$ et les $k$-dérivations de $A$ dans lui-même engendrent son
dual~; le fait que les $F_i$ soient linéairement indépendants entraîne donc
l'existence de $k$-dérivations $D_j$ de $A$ dans lui-même telles que la matrice
$(\langle F_i,D_j\otimes1\rangle)$ soit inversible, autrement dit la condition
b).

Ces remarques montrent déjà que b) entraîne a). Pour prouver inversement que
la condition a) entraîne que $j$ est injectif, considérons le diagramme
commutatif
\begin{equation}
\label{0.22.7.2.3-fr}
\xymatrix{
  (\mathfrak q/\mathfrak q^2)\otimes_AL
    \ar[r]^j \ar[d]_\alpha &
  \widehat\Omega^1_{A/k}\otimes_AL \ar@{=}[d] \\
  (\mathfrak p/\mathfrak p^2)\otimes_AL
    \ar[r]_i &
  \widehat\Omega^1_{A/k}\otimes_AL
}
\tag{22.7.2.3}
\end{equation}
et notons le lemme suivant~:
\end{proof}

\begin{lemma}[22.7.2.4]
\label{0.22.7.2.4-fr}
Soient $R$ un anneau local régulier, $\mathfrak m$ son idéal maximal,
$K=R/\mathfrak m$ son corps résiduel. Si $\mathfrak n$ est un idéal de $R$ tel
que $R/\mathfrak n$ soit régulier, alors l'homomorphisme canonique
\begin{equation}
\label{0.22.7.2.5-fr}
  \alpha:(\mathfrak n/\mathfrak n^2)\otimes_RK
  \longrightarrow(\mathfrak m/\mathfrak m^2)\otimes_RK
  \tag{22.7.2.5}
\end{equation}
est injectif.
\end{lemma}

\begin{proof}
En effet, on sait (17.1.6) que $\mathfrak n$ est engendré par une suite
$(x_i)_{1\leq i\leq t}$ faisant partie d'un système régulier de paramètres
$(x_i)_{1\leq i\leq n}$ de $R$~; les images des $x_i$ $(1\leq i\leq t)$
forment une base de $(\mathfrak n/\mathfrak n^2)\otimes_RK$, et leurs images
par $\alpha$ font partie de la base de
$(\mathfrak m/\mathfrak m^2)\otimes_RK$ formée des images des $x_i$ pour
$1\leq i\leq n$, d'où le lemme.
\end{proof}

\begin{proof}[Suite de la démonstration de la proposition~{\rm(22.7.2)}]
Ce lemme et le diagramme (22.7.2.3) ramènent donc (en vertu de l'hypothèse a))
à prouver que $i$ est injectif. Or, dans la suite
\begin{equation}
\label{0.22.7.2.6-fr}
  \mathfrak p/\mathfrak p^2\xrightarrow{h}
  \widehat\Omega^1_{A/k}\widehat\otimes_A(A/\mathfrak p)
  \xrightarrow{g}\widehat\Omega^1_{(A/\mathfrak p)/k}\longrightarrow0
  \tag{22.7.2.6}
\end{equation}
on sait (20.7.20) que $g$ est surjectif et que l'image de $h$ est dense dans
le noyau de $g$ pour la topologie $\mathfrak m$-adique ($\mathfrak m$ étant
l'idéal maximal de $A$)~; comme $\widehat\Omega^1_{A/k}$ est un $A$-module de
type fini, tous les sous-modules du $(A/\mathfrak p)$-module
\[
  \widehat\Omega^1_{A/k}\otimes_A(A/\mathfrak p)
  =\widehat\Omega^1_{A/k}\widehat\otimes_A(A/\mathfrak p)
\]
sont fermés pour la topologie $\mathfrak m$-adique ($0_{\mathrm I}$, 7.3.5),
donc la suite (22.7.2.6) est exacte. Tensorisant par $L$, il vient la suite
exacte
\[
  (\mathfrak p/\mathfrak p^2)\otimes_AL\xrightarrow{i}
  \widehat\Omega^1_{A/k}\otimes_AL\longrightarrow
  \widehat\Omega^1_{(A/\mathfrak p)/k}\otimes_{A/\mathfrak p}L
  \longrightarrow0.
\]

Or, l'hypothèse de séparabilité faite sur $L$ entraîne, en vertu de (21.9.5),
que $\widehat\Omega^1_{(A/\mathfrak p)/k}\otimes_{A/\mathfrak p}L$ est de
rang $s$ sur $L$~; comme $\widehat\Omega^1_{A/k}\otimes_AL$ est de rang $r$
sur $L$ (21.9.3), on a $\operatorname{rg}_L(\operatorname{Im}i)=r-s$. Mais comme
$A/\mathfrak p$ est finie sur une sous-algèbre isomorphe à
$k[[T_1,\ldots,T_s]]$, on a $\dim(A/\mathfrak p)=s$ ((17.1.4) et (16.1.5))~;
donc
\[
  r-s=\dim(A)-\dim(A/\mathfrak p)=\dim(A_{\mathfrak p})
\]

\oldpage[0\textsubscript{IV-1}]{211}

par (16.5.11). Or, comme $A_{\mathfrak p}$ est régulier (17.3.2),
$\dim(A_{\mathfrak p})$ est égal au rang sur $L$ de
$\mathfrak pA_{\mathfrak p}/\mathfrak p^2A_{\mathfrak p}$ (17.1.1), donc au
rang sur $L$ de $(\mathfrak p/\mathfrak p^2)\otimes_AL$, ce qui achève de
prouver que $i$ est injectif.
\end{proof}

\begin{theorem}[22.7.3]
\label{0.22.7.3-fr}
Soient $k$ un corps, $A$ un anneau local noethérien complet de corps résiduel
$K$. On suppose que~:
\begin{enumerate}
  \item[$1^\circ$] $[k:k^p]<+\infty$ (où $p$ est l'exposant caractéristique
  de $k$)~;
  \item[$2^\circ$] $K$ est une extension finie d'une extension séparable
  $K_0$ de $k$~;
  \item[$3^\circ$] $A$ est muni d'une structure de $K_0$-algèbre formellement
  lisse (pour la topologie préadique).
\end{enumerate}
Soient $\mathfrak q$ un idéal de $A$, $B=A/\mathfrak q$, $\mathfrak p$ un
idéal premier de $A$ contenant $\mathfrak q$. Les conditions suivantes sont
équivalentes~:
\begin{enumerate}
  \item[a)] $B_{\mathfrak p}$ est une $k$-algèbre formellement lisse (pour la
  topologie $\mathfrak p$-préadique).
  \item[b)] Il existe des $k$-dérivations $D_i$ de $A$ dans lui-même
  $(1\leq i\leq m)$ et des éléments $f_i$ $(1\leq i\leq m)$ de
  $\mathfrak q$, tels que les images des $f_i$ dans
  $\mathfrak qA_{\mathfrak p}$ engendrent cet idéal de $A_{\mathfrak p}$, et
  que l'on ait $\det(D_i(f_j))\notin\mathfrak p$.
  \item[b')] Il existe une sous-extension $k'$ de $K_0$, contenant $K_0^p$,
  telle que $[K_0:k']<+\infty$, des $k'$-dérivations $D_i$ de $A$ dans
  lui-même $(1\leq i\leq m)$ et des éléments $f_i$ de $\mathfrak q$
  $(1\leq i\leq m)$, tels que les images des $f_i$ dans
  $\mathfrak qA_{\mathfrak p}$ engendrent cet idéal de $A_{\mathfrak p}$, et
  que l'on ait $\det(D_i(f_j))\notin\mathfrak p$.
\end{enumerate}
\end{theorem}

\begin{proof}
Distinguons deux cas, suivant que $p=1$ ou $p>1$.

\emph{A)} $k$ \emph{est de caractéristique} $0$. Comme $K$ est alors une
extension séparable de $K_0$, il résulte de (19.6.4) que $A$ est
$K_0$-isomorphe à un anneau de séries formelles $K[[T_1,\ldots,T_r]]$ muni de
sa structure usuelle de $K$-algèbre. Mais alors, compte tenu de (19.8.8, (ii))
appliqué à $A/\mathfrak p$, on voit que les conditions générales de (22.7.2)
sont vérifiées en y remplaçant $k$ par $K$. En outre, en vertu de (19.6.4), il
revient au même de dire que $B_{\mathfrak p}$ est une $k$-algèbre
formellement lisse ou une $K$-algèbre formellement lisse, les deux conditions
étant équivalentes au fait que $B_{\mathfrak p}$ est un anneau régulier. On
peut donc appliquer les conclusions de (22.7.2), et il est immédiat que cela
prouve l'équivalence de a), b) et b') (avec $k'=K_0$ dans b')).

\emph{B)} $k$ \emph{est de caractéristique} $p>0$. Comme $A$ est une
$K_0$-algèbre formellement lisse et que $K_0$ est séparable sur $k$, il
résulte de (19.3.5, (ii)) et (19.6.1) que $A$ est une $k$-algèbre
formellement lisse~; en vertu de (22.5.9), $A_{\mathfrak p}$ est aussi une
$k$-algèbre formellement lisse pour la topologie $\mathfrak p$-préadique. Il
résulte donc encore de (22.6.2, (ii)) que la condition a) équivaut au fait que
l'homomorphisme $j_0$ défini dans (22.7.2.1) (où $L=k(\mathfrak p)$) est
injectif. Donc (19.1.12, c)) montre donc déjà que b) implique a)~; comme par
ailleurs b') implique trivialement b), tout revient à montrer que l'hypothèse
que $j_0$ est injectif entraîne b').

Désignons d'abord par $k'$ une sous-extension quelconque de $K_0$ telle que
$[K_0:k']<+\infty$. Notons que $\widehat\Omega^1_{A/k'}$ est un $A$-module
libre de type fini. Par le raisonnement de (22.7.2), il suffit de montrer que
(pour un choix convenable de $k'$), l'homomorphisme composé
\begin{equation}
\label{0.22.7.3.1-fr}
  j':(\mathfrak q/\mathfrak q^2)\otimes_AL
  \xrightarrow{j_0}\Omega^1_{A/k}\otimes_AL
  \longrightarrow\widehat\Omega^1_{A/k'}\otimes_AL
  \tag{22.7.3.1}
\end{equation}

\oldpage[0\textsubscript{IV-1}]{212}

est injectif. Considérons encore le diagramme commutatif
\begin{equation}
\label{0.22.7.3.2-fr}
\xymatrix{
  (\mathfrak q/\mathfrak q^2)\otimes_AL
    \ar[r]^{j_0} \ar[d]_\alpha &
  \Omega^1_{A/k}\otimes_AL \ar[r] \ar@{=}[d] &
  \widehat\Omega^1_{A/k'}\otimes_AL \ar@{=}[d] \\
  (\mathfrak p/\mathfrak p^2)\otimes_AL
    \ar[r]_{i_0} &
  \Omega^1_{A/k}\otimes_AL \ar[r] &
  \widehat\Omega^1_{A/k'}\otimes_AL
}
\tag{22.7.3.2}
\end{equation}

Comme $j_0$ est injectif, on a $\operatorname{Ker}(i_0\circ\alpha)=0$. Mais comme
$A_{\mathfrak p}$ est un anneau régulier et que l'hypothèse a) implique que
$B_{\mathfrak p}=A_{\mathfrak p}/\mathfrak qA_{\mathfrak p}$ est aussi un
anneau régulier (19.6.5), le lemme (22.7.2.4) montre que $\alpha$ est injectif,
d'où $\alpha^{-1}(\operatorname{Ker} i_0)=0$. Si $i'$ est le composé des homomorphismes de la
seconde ligne de (22.7.3.2), on a de même
$\operatorname{Ker}(j')=\alpha^{-1}(\operatorname{Ker}(i'))$, et tout revient donc à montrer que, pour un
choix convenable de $k'$, on a
\begin{equation}
\label{0.22.7.3.3-fr}
  \operatorname{Ker}(i')=\operatorname{Ker}(i).
  \tag{22.7.3.3}
\end{equation}

Mais le même raisonnement que dans (22.7.2) montre que l'on a une suite exacte
\[
  (\mathfrak p/\mathfrak p^2)\otimes_AL\xrightarrow{i'}
  \widehat\Omega^1_{A/k'}\otimes_AL\longrightarrow
  \widehat\Omega^1_{(A/\mathfrak p)/k'}\otimes_{A/\mathfrak p}L
  \longrightarrow0.
\]
Or, les hypothèses faites permettent d'appliquer (21.9.8), donc il existe une
sous-extension $k'$ de $K_0$ contenant $K_0^p$, telle que
$[K_0:k']<+\infty$, et pour laquelle on a
\[
  \operatorname{rg}_L\bigl(\widehat\Omega^1_{(A/\mathfrak p)/k'}
  \otimes_{A/\mathfrak p}L\bigr)
  =\dim(A/\mathfrak p)+\operatorname{rg}_L\Upsilon_{L/k}
  +\operatorname{rg}_{K_0}\Omega^1_{K_0/k'}.
\]
Mais on a
\[
  \dim(A/\mathfrak p)=\dim(A)-\dim(A_{\mathfrak p})
\]
par (16.5.11),
\[
  \dim(A_{\mathfrak p})
  =\operatorname{rg}_L\bigl((\mathfrak p/\mathfrak p^2)\otimes_AL\bigr)
\]
puisque $A_{\mathfrak p}$ est régulier (17.1.1), et enfin
\[
  \dim(A)+\operatorname{rg}_{K_0}\Omega^1_{K_0/k'}
  =\operatorname{rg}_L\bigl(\widehat\Omega^1_{A/k'}\otimes_AL\bigr)
\]
en vertu de (21.9.2)~; on obtient donc
\begin{equation}
\label{0.22.7.3.4-fr}
  \operatorname{rg}_L(\operatorname{Ker}(i'))=\operatorname{rg}_L\Upsilon_{L/k}.
  \tag{22.7.3.4}
\end{equation}

D'autre part, le corps résiduel de $A_{\mathfrak p}$ étant formellement lisse
sur le corps premier de $k$, on a une suite exacte (20.6.22.1)
\[
  \Upsilon_{L/k}\xrightarrow{\chi_{A_{\mathfrak p}/k}}
  (\mathfrak pA_{\mathfrak p})/(\mathfrak pA_{\mathfrak p})^2
  \xrightarrow{i}\Omega^1_{A_{\mathfrak p}/k}
  \otimes_{A_{\mathfrak p}}L,
\]
et comme $A_{\mathfrak p}$ est une $k$-algèbre formellement lisse pour la
topologie $\mathfrak p$-préadique, il suit de (22.2.2) que
$\chi_{A_{\mathfrak p}/k}$ est un homomorphisme injectif, donc $\operatorname{Ker}(i)$ est
isomorphe à $\Upsilon_{L/k}$, et compte tenu de (22.7.3.4) et de ce que
$\operatorname{Ker}(i)\subset\operatorname{Ker}(i')$, cela achève de prouver (22.7.3.3) et par suite le
théorème.
\end{proof}

\oldpage[0\textsubscript{IV-1}]{213}

\begin{remark}[22.7.4]
\label{0.22.7.4-fr}
On a en fait montré (en utilisant (21.9.6)) que la condition b') est encore
équivalente aux autres conditions de (22.7.3) lorsque l'on assujettit $k'$ à
être un des corps d'une famille filtrante décroissante $(k_\alpha)$ de
sous-corps de $K_0$ contenant $K_0^p$, tels que
$[K_0:k_\alpha]<+\infty$ pour tout $\alpha$ et que
$\bigcap_\alpha k_\alpha=K_0^p$.
\end{remark}

\begin{corollary}[22.7.5]
\label{0.22.7.5-fr}
Sous les hypothèses générales de (22.7.3), l'ensemble des idéaux premiers
$\mathfrak n\in\operatorname{Spec}(B)$ tels que $B_{\mathfrak n}$ soit une
$k$-algèbre formellement lisse (pour la topologie $\mathfrak n$-préadique)
est ouvert dans $\operatorname{Spec}(B)$.
\end{corollary}

\begin{proof}
Avec les notations de (22.7.3), il suffit de voir que si $B_{\mathfrak p}$
est formellement lisse sur $k$, il en est de même de $B_{\mathfrak p'}$, pour
tout $\mathfrak p'\in V(\mathfrak q)$ assez voisin de $\mathfrak p$ dans
$\operatorname{Spec}(A)$. Or, si les $k$-dérivations $D_i$ et les éléments
$f_i$ de $\mathfrak q$ vérifient le critère b) de (22.7.3), on a aussi
$f=\det(D_if_j)\notin\mathfrak p'$ pour $\mathfrak p'\in D(f)$~; d'autre
part, il existe un $g\notin\mathfrak p$ tel que les images des $f_i$ dans
$\mathfrak qA_{\mathfrak p'}$ engendrent $\mathfrak qA_{\mathfrak p'}$ pour
$\mathfrak p'\in D(g)$ (Bourbaki, \emph{Alg. comm.}, chap.~II, §~5,
n\textsuperscript{o}~1, prop.~2), ce qui achève de prouver le corollaire.
\end{proof}

\begin{corollary}[22.7.6]
\label{0.22.7.6-fr}
\textnormal{(Nagata).} Soit $B$ un anneau local noethérien complet contenant
un corps. Alors l'ensemble des $\mathfrak n\in\operatorname{Spec}(B)$ tels
que $B_{\mathfrak n}$ soit régulier est ouvert dans $\operatorname{Spec}(B)$.
\end{corollary}

\begin{proof}
En effet, $B$ est une algèbre sur un corps premier $k$, donc (19.8.8, (ii))
de la forme $A/\mathfrak q$, où
$A=K_0[[T_1,\ldots,T_n]]$ est un anneau de séries formelles et $K_0$ une
extension de $k$~; d'autre part, dire que $B_{\mathfrak n}$ est régulier
équivaut à dire que c'est une $k$-algèbre formellement lisse pour la topologie
$\mathfrak n$-préadique (19.6.4). Toutes les conditions d'application de
(22.7.5) sont donc réalisées.
\end{proof}

\begin{remarks}[22.7.7]
\label{0.22.7.7-fr}
\begin{enumerate}
  \item[(i)] Nous verrons plus loin, avec Nagata, en utilisant (22.7.6), que
  l'on peut étendre la conclusion de (22.7.6) au cas d'un anneau local
  noethérien complet quelconque (IV, 6.12.7).

  \item[(ii)] La conclusion du théorème (22.7.3) n'est plus nécessairement
  exacte lorsqu'on ne suppose plus que $[k:k^p]$ soit fini. Prenons par
  exemple $k=\mathbf F_p(X_n)_{n\geqslant0}$,
  $A=k[[T,U]]$, $\mathfrak q$ égal à l'idéal principal $Af$, où
  \[
    f=U^p-\sum_{n=0}^{\infty}X_nT^{np}~;
  \]
  $A$ est un anneau factoriel dans lequel $f$ est un élément extrémal, car il
  est extrémal dans l'anneau de polynômes $k((T))[U]$, donc aussi dans
  l'anneau de séries formelles $k((T))[[U]]$ (Bourbaki,
  \emph{Alg. comm.}, chap.~VII, §~3), et \emph{a fortiori} dans $A$.
  Prenons $\mathfrak p=\mathfrak q$, de sorte que $B_{\mathfrak p}$ est le
  corps désigné par $E$ dans (21.9.7), qui est séparable sur $k$
  (\emph{loc. cit.}), donc une $k$-algèbre formellement lisse. Mais on a ici,
  avec les notations de (22.7.3), $k=K_0=K$, et les $k$-dérivations de $A$
  dans lui-même sont les combinaisons linéaires de $\partial/\partial T$ et
  $\partial/\partial U$~; comme on a
  $\partial f/\partial T=\partial f/\partial U=0$, les critères b) et b') de
  (22.7.3) ne sont pas vérifiés.
\end{enumerate}
\end{remarks}

% FIN SOURCE ega4/ega0-22-fr.tex

%
% DEBUT SOURCE ega4/ega0-23-fr.tex
\section{Anneaux japonais}
\label{section:0.23-fr}

Les résultats de ce paragraphe seront complétés dans IV, (7.6) et (7.7).

\subsection{Anneaux japonais}
\label{subsection:0.23.1-fr}

\begin{definition}[23.1.1]
\label{0.23.1.1-fr}
On dit qu'un anneau intègre $A$ est un anneau japonais si, pour toute
extension finie $K'$ de son corps des fractions $K$, la fermeture intégrale
$A'$ de $A$ dans $K'$ est

\oldpage[0\textsubscript{IV-1}]{214}

un $A$-module de type fini (autrement dit, une $A$-algèbre finie). On dit
qu'un anneau $A$ est universellement japonais si toute $A$-algèbre intègre de
type fini est un anneau japonais.
\end{definition}

Il est clair que si $A$ est un anneau japonais, tout anneau de fractions
$S^{-1}A$ est alors un anneau japonais, car (avec les notations précédentes)
$S^{-1}A'$ est la fermeture intégrale de $S^{-1}A$ dans $K'$, et est
évidemment un $S^{-1}A$-module de type fini.

Un anneau intègre noethérien, même un \emph{anneau de valuation discrète},
n'est pas nécessairement un anneau japonais [30].

\begin{proposition}[23.1.2]
\label{0.23.1.2-fr}
Soient $A$ un anneau noethérien intègre, $K$ son corps des fractions. Si, pour
toute extension radicielle finie $K'$ de $K$, la fermeture intégrale de $A$
dans $K'$ est un $A$-module de type fini, alors $A$ est un anneau japonais.
\end{proposition}

\begin{proof}
Comme $A$ est noethérien, il suffit, pour vérifier que $A$ est un anneau
japonais, de voir que pour toute extension quasi-galoisienne finie $L$ de
$K$, la fermeture intégrale $B$ de $A$ dans $L$ est un $A$-module de type
fini. Or $L$ est une extension galoisienne de la plus grande extension
radicielle $K'$ de $K$ contenue dans $L$~; et si $A'$ est la fermeture
intégrale de $A$ dans $K'$, $B$ est la fermeture intégrale de $A'$ dans $L$.
Mais on sait que dans une extension séparable de $K'$, la fermeture intégrale
de $A'$ est un $A'$-module de type fini (Bourbaki, \emph{Alg. comm.},
chap.~V, §~1, n\textsuperscript{o}~6, cor.~1 de la prop.~20), d'où la
proposition.
\end{proof}

Il résulte de (23.1.2) que lorsque $K$ est de caractéristique $0$, dire que
$A$ est un anneau japonais signifie que sa \emph{clôture intégrale} $A'$ est
un $A$-module de type fini.

\begin{theorem}[23.1.3]
\label{0.23.1.3-fr}
\textnormal{(Tate).} Soient $A$ un anneau intègre noethérien, $x\neq0$ un
élément de $A$. Supposons vérifiées les conditions suivantes~:
\begin{enumerate}
  \item[(i)] $A$ est intégralement clos.
  \item[(ii)] $\mathfrak p=xA$ est premier et $A$ est séparé et complet pour
  la topologie $\mathfrak p$-préadique.
  \item[(iii)] $A/xA$ est un anneau japonais.
\end{enumerate}
Alors $A$ est un anneau japonais.
\end{theorem}

Nous utiliserons le lemme suivant~:

\begin{lemma}[23.1.3.1]
\label{0.23.1.3.1-fr}
Soient $A$ un anneau, $x$ un élément de $A$ non diviseur de $0$ dans $A$ et
tel que $xA=\mathfrak p$ soit premier~; alors, pour tout entier $n>0$,
l'image réciproque de $x^nA_{\mathfrak p}$ dans $A$ est $x^nA$.
\end{lemma}

\begin{proof}
En effet, supposons que $b\in A$ soit un élément tel que
$b/1=x^na/s$ dans $A_{\mathfrak p}$, où $a\in A$ et
$s\notin\mathfrak p$~; il existe donc $s'\notin\mathfrak p$ tel que
$s'sb=x^nas'$, d'où $s'sb\in\mathfrak p$, et comme
$s's\notin\mathfrak p$, cela entraîne $b\in\mathfrak p$, autrement dit
$b=xb'$ avec $b'\in A$~; puisque $x$ est non diviseur de $0$, on en conclut
$s'sb'=x^{n-1}as$ et il suffit de raisonner par récurrence sur $n$.
\end{proof}

\begin{proof}
Pour démontrer le théorème, on peut se borner au cas où le corps des fractions
$K$ de $A$ est de caractéristique $p>0$, puisque $A$ est intégralement clos.
Soit $K'$ une extension radicielle finie de $K$, de sorte qu'il existe une
puissance $q=p^e$ telle que $K'^q\subset K$~; en remplaçant $K'$ par une
extension radicielle plus grande, on peut même supposer qu'il existe $y\in K'$
tel que $y^q=x$. Soit $A'$ la fermeture intégrale de $A$ dans $K'$~; comme
$A$ est intégralement clos, $A'$ est l'ensemble des $x'\in K'$ tels que
$x'^q\in A'\cap K=A$. Posons $V=A_{\mathfrak p}$~; l'idéal maximal
$\mathfrak m=xV$ de $V$ étant principal, on sait que l'anneau local
noethérien $V$ est un \emph{anneau de valuation discrète} (17.1.4)~; comme
$V$ est intégralement clos, le même

\oldpage[0\textsubscript{IV-1}]{215}

raisonnement que ci-dessus montre que la fermeture intégrale $V'$ de $V$ dans
$K'$ est l'ensemble des $x'\in K'$ tels que $x'^q\in V$~; on sait
(Bourbaki, \emph{Alg. comm.}, chap.~VI, §~8, n\textsuperscript{o}~6,
prop.~6, et chap.~V, §~2, n\textsuperscript{o}~3, lemme~4) que $V'$ est un
\emph{anneau de valuation discrète}, dont l'idéal maximal $\mathfrak m'$ est
l'ensemble des $x'\in K'$ tels que $x'^q\in\mathfrak m$ et en outre le corps
résiduel $V'/\mathfrak m'$ est une extension \emph{finie} de
$V/\mathfrak m$ (\emph{loc. cit.}, chap.~VI, §~8,
n\textsuperscript{o}~1, lemme~2). Montrons que, pour tout entier $n>0$, on a
\begin{equation}
\label{0.23.1.3.2-fr}
  \mathfrak m'^n\cap A'=y^nA'.
  \tag{23.1.3.2}
\end{equation}

En effet, comme $y^q=x$, on a $y\in\mathfrak m'$, donc
$y^nA'\subset\mathfrak m'^n\cap A'$. Inversement, soit
$x'\in\mathfrak m'^n\cap A'$, et posons $x'=z'y^n$ avec $z'\in K'$. On a
$x'^q\in A'$~; d'autre part, $x'$ est somme de produits
$t'_1\cdots t'_n$ avec $t'_i\in\mathfrak m'$, donc
$t_i'^q\in\mathfrak m$, et l'on en conclut que
$x'^q\in\mathfrak m^n\cap A$. Or le lemme (23.1.3.1) prouve que
$\mathfrak m^n\cap A=x^nA$, donc on a $z'^qx^n\in x^nA$, ou encore
$z'^q\in A$ puisque $x\neq0$, ce qui établit (23.1.3.2).

Montrons en second lieu que sur $A'$ la topologie $xA'$-préadique est
\emph{séparée}~; cette topologie est aussi la topologie $yA'$-préadique de
$A'$, et il résulte de (23.1.3.2) que cette topologie est induite par la
topologie $\mathfrak m'$-préadique de $V'$ qui est séparée puisque $V'$ est
un anneau de valuation discrète.

Prouvons ensuite que $A'/xA'$ est un $A$-module de type fini. Comme $x=y^q$
et que $y^kA'/y^{k+1}A'$ est isomorphe à $A'/yA'$, on peut se borner à
montrer que $A'/yA'$ est un $A$-module de type fini~; mais la formule
(23.1.3.2) appliquée pour $n=1$ montre que $A'/yA'$ est un sous-anneau de
$V'/\mathfrak m'$, c'est-à-dire du corps résiduel de $V'$, qui est une
extension finie du corps résiduel $V/\mathfrak m$ de $V$. Or
$V/\mathfrak m$ est le corps des fractions de $A/\mathfrak p$, et comme
$A'$ est entier sur $A$, $A'/yA'$ est entier sur $A/\mathfrak p$, donc
contenu dans la fermeture intégrale de $A/\mathfrak p$ dans
$V'/\mathfrak m'$~; comme par hypothèse $A/\mathfrak p$ est un anneau
japonais noethérien, $A'/yA'$ est un $(A/\mathfrak p)$-module de type fini,
et \emph{a fortiori} un $A$-module de type fini.

Cela étant, la topologie $xA'$-préadique de $A'$ étant séparée, $A'$ est un
sous-anneau de son complété $\widehat{A'}$ pour cette topologie~; mais puisque
$A'/xA'$ est un $A$-module de type fini et que $A$ est séparé et complet pour
la topologie $xA$-préadique, $\widehat{A'}$ est un $A$-module de type fini en
vertu de (0\textsubscript{I}, 7.2.9), donc il en est de même de $A'$, C.Q.F.D.
\end{proof}

\begin{corollary}[23.1.4]
\label{0.23.1.4-fr}
Soit $A$ un anneau japonais noethérien et intégralement clos. Alors tout
anneau de séries formelles $A[[T_1,\ldots,T_r]]$ est un anneau japonais.
\end{corollary}

\begin{proof}
On sait que pour tout anneau noethérien intégralement clos $B$, l'anneau de
séries formelles $B[[T]]$ est noethérien et intégralement clos (Bourbaki,
\emph{Alg. comm.}, chap.~V, §~1, n\textsuperscript{o}~4, prop.~14)~; on peut
donc, par récurrence sur $n$, se borner à prouver que $A[[T]]$ est un anneau
japonais. Or l'élément $x=T$ vérifie toutes les conditions de (23.1.3), d'où
la conclusion.
\end{proof}

\begin{theorem}[23.1.5]
\label{0.23.1.5-fr}
\textnormal{(Nagata).} Tout anneau local noethérien intègre complet $A$ est
un anneau japonais.
\end{theorem}

\begin{proof}
On sait (19.8.8, (ii)) qu'il existe un sous-anneau $B$ de $A$ qui est un
anneau local complet et régulier, tel que $A$ soit une $B$-algèbre finie~; il
suffit évidemment de prouver

\oldpage[0\textsubscript{IV-1}]{216}

que $B$ est un anneau japonais, autrement dit, on peut se borner au cas où
$A$ est en outre \emph{régulier}. Raisonnons par récurrence sur
$n=\dim(A)$, le théorème étant évident pour $n=0$. Supposons donc $n>0$, et,
en désignant par $\mathfrak m$ l'idéal maximal de $A$, soit $x$ un élément de
$\mathfrak m-\mathfrak m^2$~; on sait (17.1.8) que $A/xA$ est un anneau
régulier, et en outre ((17.1.7) et (16.3.4)) que
$\dim(A/xA)=n-1$~; de plus, l'idéal $xA$ est fermé dans $A$
(0\textsubscript{I}, 7.3.5), donc $A/xA$ est complet. En vertu de l'hypothèse
de récurrence, $A/xA$ est donc un anneau japonais, et il résulte alors de
(23.1.3) qu'il en est de même de $A$.
\end{proof}

\begin{corollary}[23.1.6]
\label{0.23.1.6-fr}
Soient $A$ un anneau local noethérien intègre complet, $K$ son corps des
fractions, $K'$ une extension finie de $K$~; alors la fermeture intégrale
$A'$ de $A$ dans $K'$ est un anneau local complet.
\end{corollary}

\begin{proof}
On sait déjà par (23.1.5) que $A'$ est une $A$-algèbre finie, donc un anneau
semi-local noethérien complet, et par suite un \emph{composé direct}
d'anneaux locaux~; mais comme $A'$ est intègre, c'est nécessairement un
anneau local.
\end{proof}

\begin{proposition}[23.1.7]
\label{0.23.1.7-fr}
Soient $A$ un anneau local noethérien intègre, $K$ son corps des fractions,
$K'$ une extension finie de $K$, $A'$ la fermeture intégrale de $A$ dans
$K'$. On suppose que le complété $\widehat A$ est réduit, et l'on désigne par
$R$ son anneau total des fractions.
\begin{enumerate}
  \item[(i)] Si l'anneau $R\otimes_KK'$ est réduit (ce qui aura lieu en
  particulier lorsque $K'$ est une extension séparable de $K$, $R$ étant
  composé direct de corps, extensions de $K$ (Bourbaki, \emph{Alg.},
  chap.~VIII, §~7, n\textsuperscript{o}~3, cor.~1 du th.~1)), alors $A'$ est
  un $A$-module de type fini.
  \item[(ii)] Si en particulier $R$ est une $K$-algèbre séparable, alors $A$
  est un anneau japonais.
\end{enumerate}
\end{proposition}

\begin{proof}
\emph{(i)} Posons pour simplifier $A_1=\widehat A$,
$A'_1=A'\otimes_AA_1$, $K'_1=K'\otimes_AA_1$. Comme $A_1$ est un
$A$-module plat (0\textsubscript{I}, 7.3.5), $A'_1$ s'identifie à un
sous-anneau de $K'_1$ et est évidemment entier sur $A_1$. D'autre part, si
l'on pose $K_1=K\otimes_AA_1$, on peut écrire
$K'_1=K'\otimes_KK_1$~; comme $A$ est intègre et que $A_1$ est un $A$-module
plat, tout élément $\neq0$ de $A$ est $A_1$-régulier
(0\textsubscript{I}, 6.3.4)~; comme $K=S^{-1}A$, où $S=A-\{0\}$, on a
$K_1=S^{-1}A_1$, et la remarque précédente prouve que $K_1$ s'identifie à un
sous-anneau de l'anneau total des fractions $R$ de $A_1$. Puisque tout
$K$-module est plat, $K'_1$ s'identifie à un sous-anneau de $K'\otimes_KR$,
qui par hypothèse est réduit et est une $R$-algèbre finie. Désignons par
$\mathfrak q_i$ ($1\leqslant i\leqslant n$) les idéaux premiers minimaux de
$A_1$, par $B_i$ l'anneau intègre $A_1/\mathfrak q_i$, par $L_i$ le corps des
fractions de $B_i$, de sorte que $R$ est composé direct des $L_i$, et
$K'\otimes_KR$ composé direct des $K'\otimes_KL_i$~; ces dernières
$K$-algèbres, étant réduites, sont des composées directes de corps,
extensions finies des $L_i$. Comme $B_i$ est un anneau local intègre
\emph{complet}, le théorème de Nagata (23.1.5) prouve que la fermeture
intégrale de $B_i$ dans $K'\otimes_KL_i$ est un $B_i$-module de type fini, et
\emph{a fortiori} un $A_1$-module de type fini~; comme tout élément de
$K'\otimes_KL_i$ qui est entier sur $A_1$ est aussi entier sur $B_i$ (étant
annulé par $\mathfrak q_i$), on en conclut finalement que la fermeture
intégrale de $A_1$ dans $K'\otimes_KR$ est un $A_1$-module de type fini. Mais
comme $A'_1$ est contenu dans cette fermeture intégrale et que $A_1$ est
noethérien, $A'_1$ est aussi un $A_1$-module de type fini. Enfin, comme $A_1$
est un $A$-module fidèlement plat (0\textsubscript{I}, 7.3.5), on en conclut
que $A'$ est un $A$-module de type fini (Bourbaki, \emph{Alg. comm.},
chap.~I, §~3, n\textsuperscript{o}~6, prop.~11).

\oldpage[0\textsubscript{IV-1}]{217}

\emph{(ii)} L'hypothèse entraîne que les $L_i$ sont des extensions séparables
de $K$, donc $K'\otimes_KL_i$ est réduit (Bourbaki, \emph{Alg.}, chap.~VIII,
§~7, n\textsuperscript{o}~3, th.~1) et l'on peut appliquer (i) à toute
extension finie $K'$ de $K$, ce qui démontre notre assertion.
\end{proof}

\subsection{Clôture intégrale d'un anneau local noethérien intègre}
\label{subsection:0.23.2-fr}

\begin{env}[23.2.1]
\label{0.23.2.1-fr}
Dans ce qui suit, pour tout anneau $A$, nous noterons $A_{\mathrm{red}}$ pour
simplifier, le quotient de $A$ par son nilradical (de sorte que si
$X=\operatorname{Spec}(A)$, on a
$X_{\mathrm{red}}=\operatorname{Spec}(A_{\mathrm{red}})$).

Nous dirons qu'un anneau \emph{local} $A$ est \emph{unibranche} si l'anneau
$A_{\mathrm{red}}$ est intègre et si la clôture intégrale de
$A_{\mathrm{red}}$ est un anneau \emph{local}, ce qui généralise la définition
donnée dans (III, 4.3.6). Nous dirons que $A$ est \emph{géométriquement
unibranche} s'il est unibranche et si le corps résiduel de l'anneau local,
clôture intégrale de $A_{\mathrm{red}}$, est une extension \emph{radicielle}
de celui de $A$. Il est clair qu'un anneau local intégralement clos est
géométriquement unibranche~; il résulte de (23.1.6) qu'un anneau local
noethérien intègre \emph{complet} est unibranche.
\end{env}

\begin{lemma}[23.2.2]
\label{0.23.2.2-fr}
\textnormal{(*)}\footnote{Dans la suite de ce numéro, nous utilisons des
notions et résultats exposés au chap.~IV, §§~2 et 5~; comme les résultats de
ce numéro ne sont pas utilisés avant le chap.~IV, §~6, cela n'entraîne pas de
cercle vicieux.} Soient $A$ un anneau intègre, $A'$ sa clôture intégrale~; pour
que le morphisme canonique
$f:\operatorname{Spec}(A')\to\operatorname{Spec}(A)$ soit radiciel, il faut et
il suffit que pour tout $\mathfrak p\in\operatorname{Spec}(A)$,
$A_{\mathfrak p}$ soit géométriquement unibranche~; lorsqu'il en est ainsi,
$f$ est un homéomorphisme.
\end{lemma}

\begin{proof}
En effet, pour tout $\mathfrak p\in\operatorname{Spec}(A)$, la clôture
intégrale de $A_{\mathfrak p}$ est $A'_{\mathfrak p}$ (Bourbaki,
\emph{Alg. comm.}, chap.~V, §~1, n\textsuperscript{o}~5, prop.~16), et tous
les idéaux premiers de $A'_{\mathfrak p}$ au-dessus de
$\mathfrak pA_{\mathfrak p}$ sont maximaux (\emph{loc. cit.}, §~2,
n\textsuperscript{o}~1, prop.~1). Dire que $f$ est injectif signifie donc que
pour tout $\mathfrak p\in\operatorname{Spec}(A)$, $A'_{\mathfrak p}$ est un
anneau \emph{local}, c'est-à-dire que les $A_{\mathfrak p}$ sont
unibranches. Dire que tout $A_{\mathfrak p}$ est géométriquement unibranche
signifie alors que $f$ est radiciel en vertu de (I, 3.5.8). Lorsqu'il en est
ainsi, $f$ est surjectif et fermé (II, 6.1.10), donc un homéomorphisme.
\end{proof}

\begin{lemma}[23.2.3]
\label{0.23.2.3-fr}
Soient $A$ un anneau intègre, $K$ son corps des fractions, $B$ un anneau
noethérien, $\varphi:A\to B$ un homomorphisme d'anneau faisant de $B$ un
$A$-module plat. Soit $K'$ une extension de $K$~; posons
$B_1=B_{\mathrm{red}}$, et soit $R$ l'anneau total des fractions de $B_1$.
Alors, pour toute sous-$A$-algèbre $A'$ de $K'$, l'homomorphisme canonique
\begin{equation}
\label{0.23.1.8.1-fr}
  (A'\otimes_AB)_{\mathrm{red}}
  \longrightarrow(K'\otimes_AR)_{\mathrm{red}}
  \tag{23.1.8.1}
\end{equation}
est injectif.
\end{lemma}

\begin{proof}
On peut en effet considérer l'homomorphisme canonique
$h:A'\otimes_AB\to K'\otimes_AR$ comme composé des homomorphismes suivants
\[
  A'\otimes_AB\xrightarrow{u}K'\otimes_AB
  \xrightarrow{v}K'\otimes_AB_1\xrightarrow{w}K'\otimes_AR.
\]
Comme $B$ est un $A$-module plat, $u$ est injectif~; de même, $K'$ étant un
$K$-module plat et $K$ un $A$-module plat, $K'$ est un $A$-module plat, et
$w$ est donc injectif puisque $B_1$ est réduit, donc l'homomorphisme
$B_1\to R$ injectif. Enfin, pour la même raison, si $\mathfrak N$ est

\oldpage[0\textsubscript{IV-1}]{218}

le nilradical de $B$, le noyau de $v$ est $K'\otimes_A\mathfrak N$, donc il est
contenu dans le nilradical de $K'\otimes_AB$~; on en conclut aussitôt que si
l'image par $h=w\circ v\circ u$ d'un élément $x\in A'\otimes_AB$ est
nilpotent, $x$ est nilpotent, ce qui prouve le lemme.
\end{proof}

\begin{proposition}[23.2.4]
\label{0.23.2.4-fr}
Soient $A$ un anneau intègre, $K$ son corps des fractions, $B$ un anneau
noethérien, $\mathfrak q_i$ ($1\leqslant i\leqslant n$) ses idéaux premiers
minimaux, $\varphi:A\to B$ un homomorphisme d'anneaux. On suppose que les
$B/\mathfrak q_i$ sont des anneaux japonais, et que $B$ soit un $A$-module
plat. Soit $K'$ une extension finie de $K$, et soit $(C_\lambda)$ la famille
filtrante croissante des sous-anneaux de $K'$ qui sont des $A$-algèbres
finies et admettent $K'$ pour corps des fractions~; la réunion $A'$ des
$C_\lambda$ est donc la fermeture intégrale de $A$ dans $K'$. Alors~:
\begin{enumerate}
  \item[(i)] Il existe un indice $\alpha$ tel que pour $\lambda\geqslant
  \alpha$, l'homomorphisme canonique
  \begin{equation}
  \label{0.23.2.4.1-fr}
    (C_\alpha\otimes_AB)_{\mathrm{red}}
    \longrightarrow(C_\lambda\otimes_AB)_{\mathrm{red}}
    \tag{23.2.4.1}
  \end{equation}
  soit bijectif.
  \item[(ii)] Si en outre $B$ est un $A$-module fidèlement plat, le morphisme
  $\operatorname{Spec}(A')\to\operatorname{Spec}(C_\alpha)$ est radiciel.
\end{enumerate}
\end{proposition}

\begin{proof}
\emph{(i)} Soit $B_1=B_{\mathrm{red}}$, et soit $R$ l'anneau total des
fractions de $B_1$~; comme $A$ est intègre et $B$ un $A$-module plat, tout
élément $\neq0$ de $A$ est $B$-régulier (0\textsubscript{I}, 6.3.4), donc
l'homomorphisme composé $A\xrightarrow{\varphi}B\to B_1$ se prolonge en un
homomorphisme $K\to R$~; on peut alors écrire
$K'\otimes_AR=(K'\otimes_AK)\otimes_KR$, et comme $K'\otimes_AK$
s'identifie à $K'$, on a $K'\otimes_AR=K'\otimes_KR$. Comme $R$ est composé
direct des corps des fractions $L_i$ des $B/\mathfrak q_i$,
$K'\otimes_AR$ est composé direct des $K'\otimes_KL_i$, et par suite
$(K'\otimes_AR)_{\mathrm{red}}$ est composé direct d'un nombre fini
d'extensions finies des $L_i$. En vertu de l'hypothèse, la fermeture
intégrale de $B/\mathfrak q_i$ dans une extension finie de $L_i$ est un
$(B/\mathfrak q_i)$-module de type fini, donc un $B$-module de type fini~;
on en conclut que la fermeture intégrale $B'$ de $B$ dans
$(K'\otimes_AR)_{\mathrm{red}}$ est un $B$-module de type fini. En vertu de
(23.2.3), les $(C_\lambda\otimes_AB)_{\mathrm{red}}$ s'identifient
canoniquement à des sous-anneaux de $(K'\otimes_AR)_{\mathrm{red}}$ qui sont
des $B$-algèbres finies, donc contenues dans $B'$. Comme $B$ est noethérien
et $B'$ un $B$-module de type fini, la famille filtrante des
$(C_\lambda\otimes_AB)_{\mathrm{red}}$ admet un \emph{plus grand élément}
$(C_\alpha\otimes_AB)_{\mathrm{red}}$, ce qui prouve (i).

\emph{(ii)} Comme $B$ est un $A$-module fidèlement plat, il suffit
(IV, 2.6.1, (v)) de montrer que le morphisme
$\operatorname{Spec}(A'\otimes_AB)\to
\operatorname{Spec}(C_\alpha\otimes_AB)$ est radiciel, ou, ce qui revient au
même (I, 5.1.6), que le morphisme
$\operatorname{Spec}(A'\otimes_AB)_{\mathrm{red}}\to
\operatorname{Spec}(C_\alpha\otimes_AB)_{\mathrm{red}}$ est radiciel~; mais
on a $A'\otimes_AB=\varinjlim_\lambda(C_\lambda\otimes_AB)$, donc
(IV, 5.13.2)
$(A'\otimes_AB)_{\mathrm{red}}=
\varinjlim_\lambda(C_\lambda\otimes_AB)_{\mathrm{red}}$, et la conclusion
résulte de (i).
\end{proof}

\begin{corollary}[23.2.5]
\label{0.23.2.5-fr}
Soient $A$ un anneau local noethérien intègre, $K$ son corps des fractions,
$K'$ une extension finie de $K$. Soit $(C_\lambda)$ la famille filtrante
croissante des sous-anneaux de $K'$ qui sont des $A$-algèbres finies (donc
des anneaux noethériens semi-locaux (Bourbaki, \emph{Alg. comm.}, chap.~IV,
§~2, n\textsuperscript{o}~5, cor.~3 de la prop.~9)) et admettent $K'$ pour
corps des fractions. Alors il existe $\alpha$ tel que l'homomorphisme
$(\widehat C_\alpha)_{\mathrm{red}}\to
(\widehat C_\lambda)_{\mathrm{red}}$ soit un isomorphisme pour
$\lambda\geqslant\alpha$, et si $A'$ est la fermeture intégrale de $A$ dans
$K'$, le morphisme
$\operatorname{Spec}(A')\to\operatorname{Spec}(C_\alpha)$ est radiciel
(cf. (23.2.2)).
\end{corollary}

\oldpage[0\textsubscript{IV-1}]{219}

\begin{proof}
On applique (23.2.4) en prenant $B=\widehat A$~; comme les
$B/\mathfrak q_i$ sont des anneaux locaux noethériens complets, ce sont des
anneaux japonais (23.1.5) et $B$ est un $A$-module fidèlement plat
(0\textsubscript{I}, 7.3.5)~; en outre on a
$C_\lambda\otimes_AB=\widehat C_\lambda$ (Bourbaki, \emph{Alg. comm.},
chap.~III, §~3, n\textsuperscript{o}~4, th.~3 et chap.~IV, §~2,
n\textsuperscript{o}~5, cor.~3 de la prop.~9).
\end{proof}

\begin{corollary}[23.2.6]
\label{0.23.2.6-fr}
Sous les hypothèses de (23.2.5), la fermeture intégrale $A'$ de $A$ dans
$K'$ est un anneau semi-local~; si $\mathfrak m$ est l'idéal maximal de $A$,
alors, pour tout idéal maximal $\mathfrak m'$ de $A'$,
$A'/\mathfrak m'$ est une extension finie de $A/\mathfrak m$.
\end{corollary}

\begin{proof}
Le fait que l'homomorphisme
$(\widehat C_\alpha)_{\mathrm{red}}\to
(\widehat C_\lambda)_{\mathrm{red}}$ soit bijectif pour
$\lambda\geqslant\alpha$ entraîne que le nombre des idéaux maximaux de
$C_\lambda$ est constant pour $\lambda\geqslant\alpha$ et que si
$\mathfrak m_\alpha$ est un idéal maximal de $C_\alpha$ et
$\mathfrak m_\lambda$ l'unique idéal maximal de $C_\lambda$ au-dessus de
$\mathfrak m_\alpha$, les corps $C_\alpha/\mathfrak m_\alpha$ et
$C_\lambda/\mathfrak m_\lambda$ sont canoniquement isomorphes. La conclusion
résulte de ce que $\mathfrak m'=\varinjlim\mathfrak m_\lambda$ et
$A'/\mathfrak m'=\varinjlim(C_\lambda/\mathfrak m_\lambda)$
(0\textsubscript{III}, 10.3.1.1.3).
\end{proof}

\begin{proposition}[23.2.7]
\label{0.23.2.7-fr}
\textnormal{(Y. Mori).} Soient $A$ un anneau local noethérien intègre, $K$ son
corps des fractions, $A'$ la clôture intégrale de $A$. Alors $A'$ est un
anneau de Krull (Bourbaki, \emph{Alg. comm.}, chap.~VII, §~1) semi-local~;
autrement dit, il existe une famille $(V_\lambda)$ d'anneaux de valuation
discrète ayant $K$ pour corps des fractions, tels que
$A'=\bigcap_\lambda V_\lambda$ et que tout $x\in K$ appartienne à tous les
$V_\lambda$ sauf pour un nombre fini d'entre eux. En outre, il existe un
sous-anneau $C$ de $K$ qui est une $A$-algèbre finie, et telle que les
$V_\lambda$ soient les clôtures intégrales des anneaux
$C_{\mathfrak p_\lambda}$, où $(\mathfrak p_\lambda)$ est la famille des
idéaux premiers de hauteur $1$ de l'anneau local $C$.
\end{proposition}

Prouvons d'abord le lemme suivant~:

\begin{lemma}[23.2.7.1]
\label{0.23.2.7.1-fr}
Soient $A$, $B$ deux anneaux, $\varphi:A\to B$ un homomorphisme faisant de $B$
un $A$-module fidèlement plat. On suppose que $A$ est intègre~; alors, si
$C=B_{\mathrm{red}}$, l'homomorphisme composé
$\psi:A\to B\to B_{\mathrm{red}}=C$ est injectif et se prolonge en un
homomorphisme injectif du corps des fractions $K$ de $A$ dans l'anneau total
des fractions $R$ de $C$~; en outre, si $C'$ est la fermeture intégrale de
$C$ dans $R$, $C'\cap K$ est la clôture intégrale de $A$.
\end{lemma}

\begin{proof}
Comme $A$ est intègre et $\varphi$ injectif, l'intersection de $\varphi(A)$
et du nilradical $\mathfrak N$ de $B$ est réduite à $0$, et comme
$C=B/\mathfrak N$, $\psi$ est injectif. Tout $a\neq0$ dans $A$ étant non
diviseur de zéro dans $A$ ne l'est pas non plus dans $B$ par platitude
(0\textsubscript{I}, 6.3.4)~; on en déduit que $a$ n'est pas non plus diviseur
de $0$ dans $C$, car si l'on avait $ax\in\mathfrak N$ pour un
$x\notin\mathfrak N$ dans $B$, on en tirerait $a^nx^n=0$ pour un entier $n$,
ce qui contredit ce qui précède puisque $x^n\neq0$. On peut donc prolonger
$\psi$ en un homomorphisme injectif de $K$ dans $R$. Pour prouver la dernière
assertion, notons qu'il est clair que la clôture intégrale $A'$ de $A$ est
contenue dans $C'\cap K$. Inversement, soit $x\in C'\cap K$~; $x$ est donc
entier sur $C$, et \emph{a fortiori} sur $B$, autrement dit $B[x]$ est une
$B$-algèbre finie~; d'ailleurs, $K$ s'identifie à un sous-anneau de
$B\otimes_AK$, et le sous-anneau $B[x]$ de $B\otimes_AK$ s'identifie à
$B\otimes_AA[x]$ par platitude~; on en conclut que $A[x]$ est un $A$-module
de type fini (Bourbaki, \emph{Alg. comm.}, chap.~I\textsuperscript{er}, §~3,
n\textsuperscript{o}~6, prop.~11), donc que $x\in A'$.
\end{proof}

Ce lemme étant établi, nous allons l'appliquer à l'injection canonique de $A$
dans son complété $\widehat A$, qui est un $A$-module fidèlement plat~;
$B=(\widehat A)_{\mathrm{red}}=\widehat A/\mathfrak N$ (où $\mathfrak N$ est
le nilradical de $\widehat A$) est un anneau local noethérien complet réduit,
dont l'anneau total des

\oldpage[0\textsubscript{IV-1}]{220}

fractions $R$ est donc composé direct d'un nombre fini de corps $L_i$~;
l'image canonique $B_i$ de $B$ dans $L_i$ est un anneau local noethérien
\emph{intègre et complet}, dont $L_i$ est le corps de fractions, et la
fermeture intégrale $B'$ de $B$ dans $R$ est \emph{composée directe} des
$B'_i$, où $B'_i$ est la clôture intégrale de $B_i$. Mais en vertu de
(23.1.5), $B'_i$ est une $B_i$-algèbre \emph{finie}, donc un anneau local
\emph{noethérien} intégralement clos. On sait (Bourbaki,
\emph{Alg. comm.}, chap.~VII, §~1, n\textsuperscript{o}~3, cor. du th.~2) que
$B'_i$ est un anneau de Krull. Or, pour tout $i$, on a un homomorphisme
$\varphi_i:K\to L_i$, qui est injectif, et par suite
$\varphi_i^{-1}(B'_i)$ est un anneau de Krull dans $K$. Mais comme
$A'=B'\cap K$ en vertu du lemme et que
$B'=\bigcap_i\varphi_i^{-1}(B'_i)$, $A'$ est intersection d'un nombre fini
d'anneaux de Krull et est par suite un anneau de Krull.

Pour prouver la dernière assertion de la proposition, notons qu'il existe une
$A$-algèbre finie $C\subset K$ telle que le morphisme
$\operatorname{Spec}(A')\to\operatorname{Spec}(C)$ soit radiciel et un
homéomorphisme (23.2.4)~; pour tout idéal premier
$\mathfrak p'\in\operatorname{Spec}(A')$, $\mathfrak p'$ est donc le seul
idéal premier de $A'$ au-dessus de $\mathfrak p=\mathfrak p'\cap C$ et
$A'_{\mathfrak p'}$ la clôture intégrale de $C_{\mathfrak p}$~; d'autre part
le fait que $\operatorname{Spec}(A')\to\operatorname{Spec}(C)$ est un
homéomorphisme entraîne que l'application
$\mathfrak p'\mapsto\mathfrak p'\cap C$ est une bijection de l'ensemble des
idéaux premiers de hauteur $1$ de $A'$ sur l'ensemble des idéaux premiers de
hauteur $1$ de $C$~; d'où la conclusion, compte tenu de Bourbaki,
\emph{Alg. comm.}, chap.~VII, §~1, n\textsuperscript{o}~6, th.~4.

\begin{remarks}[23.2.8]
\label{0.23.2.8-fr}
\begin{enumerate}
  \item[(i)] On aura soin de noter, dans l'application de (23.2.6) et
  (23.2.7), que l'anneau $A'$ \emph{n'est pas nécessairement un anneau
  noethérien}, comme le montre un exemple de Nagata avec $\dim(A)=3$ [30].

  \item[(ii)] Les conclusions de (23.2.7) sont encore valables si $A'$ est
  la fermeture intégrale de $A$ dans une extension finie $K'$ de $K$~; il
  suffit en effet de considérer une $A$-algèbre finie $B$ dont $K'$ est le
  corps des fractions~; $B$ est un anneau semi-local noethérien, donc
  intersection d'un nombre fini d'anneaux locaux noethériens
  $B_{\mathfrak m_i}$ ($\mathfrak m_i$ idéaux maximaux de $B$), et sa clôture
  intégrale (qui est égale à $A'$) est intersection des
  $A'_{\mathfrak m_i}$ (Bourbaki, \emph{Alg. comm.}, chap.~II, §~3,
  n\textsuperscript{o}~3, cor.~4 du th.~1) qui sont les clôtures intégrales
  des $B_{\mathfrak m_i}$, donc des anneaux de Krull par (23.1.13)~; par
  suite $A'$ est un anneau de Krull.

  \item[(iii)] On peut montrer ([30], 3.3.10) que pour tout anneau noethérien
  intègre $A$ (non nécessairement local), la clôture intégrale de $A$ est un
  anneau de Krull.
\end{enumerate}
\end{remarks}

\begin{corollary}[23.2.9]
\label{0.23.2.9-fr}
Soient $A$ un anneau noethérien intègre, $A'$ sa clôture intégrale. Supposons
que pour tout anneau $C$ tel que $A\subset C\subset A'$ et tel que $C$ soit
une $A$-algèbre finie, et pour tout idéal premier $\mathfrak q$ de hauteur $1$
dans $C$, $\mathfrak q\cap A$ soit un idéal premier de hauteur $1$ dans $A$.
Soit $P$ l'ensemble des idéaux premiers de hauteur $1$ dans $A$~; on a alors
\[
  A'=\bigcap_{\mathfrak p\in P}A'_{\mathfrak p}.
\]
\end{corollary}

\begin{proof}
Comme $A'$ est un $A$-module sans torsion, on a
$A'=\bigcap A'_{\mathfrak m}$, où $\mathfrak m$ parcourt l'ensemble des idéaux
maximaux de $A$ (Bourbaki, \emph{Alg. comm.}, chap.~II, §~3,
n\textsuperscript{o}~3, cor.~4 du th.~1)~;

\oldpage[0\textsubscript{IV-1}]{221}

il suffira donc de prouver que $A'_{\mathfrak m}$ est l'intersection des
$A'_{\mathfrak p}$ pour les idéaux premiers $\mathfrak p\subset\mathfrak m$
de hauteur $1$. Comme $A'_{\mathfrak m}$ est la clôture intégrale de
$A_{\mathfrak m}$, on voit qu'on peut se borner à démontrer le corollaire pour
$A_{\mathfrak m}$. Or, il y a alors, en vertu de (23.2.7), une
$A_{\mathfrak m}$-algèbre finie $B$ contenue dans $A'_{\mathfrak m}$, telle
que $A'_{\mathfrak m}$ soit l'intersection des clôtures intégrales des anneaux
$B_{\mathfrak q}$ où $\mathfrak q$ parcourt l'ensemble des idéaux premiers de
hauteur $1$ de $B$. Mais $B$ est de la forme $C_{\mathfrak m}$, où $C$ est
une $A$-algèbre finie~; donc, pour tout idéal premier $\mathfrak q$ de hauteur
$1$ dans $B$, $\mathfrak q\cap A$ est par hypothèse un idéal premier de
hauteur $1$ dans $A$~; comme la clôture intégrale de $B_{\mathfrak q}$
contient celle de $A_{\mathfrak q\cap A}$ et que cette dernière contient
$A'_{\mathfrak m}$, il en résulte bien que $A'_{\mathfrak m}$ est
l'intersection des $A'_{\mathfrak p}$ pour $\mathfrak p\in P$ et
$\mathfrak p\subset\mathfrak m$, ce qui achève la démonstration.
\end{proof}

\begin{remarks}[23.2.10]
\label{0.23.2.10-fr}
\begin{enumerate}
  \item[(i)] Nous verrons (IV, 5.6.3 et 5.6.10) que l'hypothèse de (23.2.9)
  est vérifiée lorsque $A$ est \emph{universellement caténaire}, en particulier
  (IV, 5.10.17 et 5.11.2) lorsque $A$ est \emph{quotient d'un anneau
  régulier}~; enfin, elle est évidemment vérifiée si
  $\operatorname{Spec}(A')\to\operatorname{Spec}(A)$ est un homéomorphisme,
  en particulier si ce morphisme est radiciel (23.2.2).

  \item[(ii)] La conclusion de (23.2.9) n'est plus nécessairement exacte
  lorsqu'on omet l'hypothèse sur les sous-$A$-algèbres finies de $A'$. En
  effet, on peut définir un anneau local noethérien intègre $A$, de dimension
  $2$, dont la clôture intégrale $A'$ est une $A$-algèbre finie, mais tel qu'il
  y ait un idéal premier $\mathfrak p'$ de $A'$ de hauteur $1$ au-dessus de
  l'idéal maximal $\mathfrak m$ (de hauteur $2$) de $A$, et tel en outre que
  pour tout idéal premier $\mathfrak p$ de hauteur $1$ dans $A$,
  $A_{\mathfrak p}$ soit intégralement clos (IV, 5.6.11)~; l'intersection de
  ces anneaux $A_{\mathfrak p}$ ne peut alors être égale à $A'$ puisque $A'$
  est un anneau de Krull (Bourbaki, \emph{Alg. comm.}, chap.~VII, §~1,
  n\textsuperscript{o}~5, prop.~9).
\end{enumerate}
\end{remarks}

\begin{flushright}
\emph{(A suivre.)}
\end{flushright}

% FIN SOURCE ega4/ega0-23-fr.tex

%
% DEBUT SOURCE ega4/ega4-1-fr.tex
\cleardoublepage
\markboth{CHAPITRE IV}{CHAPITRE IV}
\oldpage[IV-1]{222}
\thispagestyle{plain}

\begin{center}
{\Large CHAPITRE IV}

\vspace{1.5cm}
{\LARGE\bfseries ÉTUDE LOCALE DES SCHÉMAS\\[0.3em]
ET DES MORPHISMES DE SCHÉMAS}

\vspace{1.5cm}
{\large\bfseries Sommaire}
\end{center}

\medskip
\begin{tabular}{@{}r p{0.84\textwidth}}
§ 1. & Conditions de finitude relatives. Ensembles constructibles dans les préschémas.\\
§ 2. & Changement de base et platitude.\\
§ 3. & Cycles premiers associés et décompositions primaires.\\
§ 4. & Changement du corps de base dans les préschémas.\\
§ 5. & Dimension et profondeur dans les préschémas.\\
§ 6. & Morphismes plats de préschémas localement noethériens.\\
§ 7. & Application aux relations entre un anneau local noethérien et son complété.
Anneaux excellents.\\
§ 8. & Limites projectives de préschémas.\\
§ 9. & Propriétés constructibles.\\
§ 10. & Préschémas de Jacobson.\\
§ 11.\footnotemark[1] & Propriétés topologiques des morphismes plats de présentation finie.
Critères locaux de platitude.\\
§ 12. & Étude des fibres des morphismes plats de présentation finie.\\
§ 13. & Morphismes équidimensionnels.\\
§ 14. & Morphismes universellement ouverts.\\
§ 15. & Étude des fibres d'un morphisme universellement ouvert.\\
§ 16. & Invariants différentiels. Morphismes différentiellement lisses.\\
§ 17. & Morphismes lisses, non ramifiés, étales.\\
§ 18. & Compléments sur les morphismes étales. Anneaux locaux henséliens.\\
§ 19. & Immersions régulières et transversalement régulières.\\
§ 20. & Sections hyperplanes~; projections génériques.\\
§ 21. & Prolongements infinitésimaux.\\
\end{tabular}
\footnotetext[1]{L'ordre et le contenu des §§ 11 à 21 ne sont donnés qu'à titre
indicatif et pourront être modifiés avant leur publication.}

\clearpage
\oldpage[IV-1]{223}

Les sujets traités dans ce Chapitre appellent les remarques suivantes~:

\medskip
\noindent
\emph{a)} Leur caractère commun est d'être relatifs à des propriétés
\emph{locales} de préschémas ou de morphismes, i.e. considérées en un point,
ou aux points d'une fibre, ou dans un voisinage (non spécifié) d'un point ou
d'une fibre. Ces propriétés sont généralement de nature \emph{topologique},
\emph{différentielle} ou \emph{dimensionnelle} (i.e. faisant intervenir les
notions de \emph{dimension} et de \emph{profondeur}) et sont liées aux
propriétés des \emph{anneaux locaux} aux points considérés. Un problème type
est de mettre en relation, pour un morphisme donné $f:X\to Y$ et un point
$x\in X$, les propriétés de $X$ en $x$ avec celles de $Y$ en $y=f(x)$ et de
la fibre $X_y=f^{-1}(y)$ en $x$. Un autre est de déterminer la nature
topologique (par exemple la constructibilité, ou le fait d'être ouvert ou
fermé) de l'ensemble des points $x\in X$ en lesquels $X$ possède une certaine
propriété ou pour lesquels la fibre $X_{f(x)}$ passant par $x$ possède une
certaine propriété en $x$. De même, on s'intéressera à la nature topologique
de l'ensemble des points $y\in Y$ tels que $X$ ait une certaine propriété en
tous les points de la fibre $X_y$, ou tels que cette fibre ait une certaine
propriété.

\medskip
\noindent
\emph{b)} Les notions les plus importantes pour les chapitres suivants sont
celle de \emph{morphisme plat de présentation finie} et celles de
\emph{morphisme lisse} et de \emph{morphisme étale}, qui en sont des cas
particuliers. Leur étude détaillée (jointe à celle de questions connexes)
commence à proprement parler au §~11.

\medskip
\noindent
\emph{c)} Les §§~1 à 10 peuvent être considérés comme de nature préliminaire
et développent trois types de techniques, utilisées non seulement dans les
autres paragraphes du chapitre mais aussi bien entendu dans les chapitres
suivants~:

\medskip
\noindent
\emph{c 1)} Aux §§~1 à 4 sont envisagés divers aspects de la notion de
\emph{changement de base}, surtout en relation avec des conditions de
\emph{finitude} ou de \emph{platitude}~; on y amorce la technique de
\emph{descente} sous ses aspects les plus élémentaires (les questions
d'«~effectivité~» liées à cette technique seront étudiées dans le chap.~V).

\medskip
\noindent
\emph{c 2)} Les §§~5 à 7 sont centrés sur ce que l'on peut appeler des
\emph{techniques noethériennes}, les préschémas considérés étant toujours
supposés localement noethériens, tandis qu'au contraire aucune condition de
finitude n'est généralement imposée aux \emph{morphismes}~; cela tient
essentiellement à ce que les notions de dimension et de profondeur ne sont
guère maniables que pour les anneaux locaux noethériens. Rappelons que le §~7
constitue une théorie «~fine~» des anneaux locaux noethériens, assez peu
utilisée dans la suite du chapitre.

\medskip
\noindent
\emph{c 3)} Les §§~8 à 10 donnent entre autres le moyen d'\emph{éliminer les
hypothèses noethériennes} sur les préschémas étudiés, en leur substituant des
hypothèses convenables de \emph{finitude} (la «~présentation finie~») sur les
\emph{morphismes} considérés~: l'avantage de cette substitution est que ces
dernières hypothèses sont \emph{stables par changement de base}, ce qui n'est
pas le cas pour les hypothèses noethériennes sur les préschémas. La technique
permettant ce passage repose, d'une part, sur l'utilisation de la notion de
\emph{limite projective} de préschémas, grâce à laquelle on peut ramener une
question à la même question sous des hypothèses \emph{noethériennes}~; d'autre
part, sur l'emploi systématique des \emph{ensembles constructibles}, qui ont
le double intérêt de se conserver par image réciproque (par un morphisme

\oldpage[IV-1]{224}

quelconque) et par image directe (par un morphisme de présentation finie), et
d'avoir des propriétés topologiques maniables dans les préschémas localement
noethériens. Les mêmes techniques permettent même souvent de se ramener à des
anneaux noethériens plus particuliers, par exemple les
$\mathbf Z$-\emph{algèbres de type fini}, et c'est là que les propriétés des
anneaux «~excellents~» (étudiées au §~7) interviennent de façon décisive.
Indépendamment d'ailleurs de la question de l'élimination des hypothèses
noethériennes, les techniques des §§~8 à 10, de nature très élémentaire, sont
d'usage constant dans presque toutes les applications.

\setcounter{section}{0}

\section{Conditions de finitude relatives. Ensembles constructibles dans les préschémas}
\label{section:IV.1-fr}

Nous allons, dans ce paragraphe, reprendre, en les complétant, l'exposé des
«~conditions de finitude~» pour un morphisme de préschémas $f:X\to Y$, donné
dans (I, 6.3 et 6.6). Il y a essentiellement deux notions de «~finitude~» de
nature \emph{globale} sur $X$, celle de morphisme \emph{quasi-compact}
(définie dans (I, 6.6.1)) et celle de morphisme \emph{quasi-séparé}~; il y a
d'autre part deux notions de «~finitude~» de nature \emph{locale} sur $X$,
celle de morphisme \emph{localement de type fini} (définie dans (I, 6.6.2))
et celle de morphisme \emph{localement de présentation finie}. En combinant
ces notions locales aux notions globales précédentes, on obtient les notions
de morphisme \emph{de type fini} (définie dans (I, 6.3.1)) et de morphisme
\emph{de présentation finie}. Pour la commodité du lecteur, nous donnerons de
nouveau, dans ce paragraphe, les principales propriétés énoncées dans (I, 6.3
et 6.6), en renvoyant bien entendu à ces numéros du chapitre I pour leurs
démonstrations.

Aux n\textsuperscript{os}~(1.8) et (1.9), nous complétons, dans le cadre des
préschémas, et en faisant usage des notions de finitude précédentes, les
résultats sur les ensembles constructibles donnés dans
(0\textsubscript{III}, §~9).

\subsection{Morphismes quasi-compacts}
\label{subsection:IV.1.1-fr}

\begin{definition}[1.1.1]
\label{IV.1.1.1-fr}
On dit qu'un morphisme de préschémas $f:X\to Y$ est \emph{quasi-compact} si
l'application continue $f$ de l'espace topologique $X$ dans l'espace
topologique $Y$ est quasi-compacte (0\textsubscript{III}, 9.1.1), autrement
dit si l'image réciproque $f^{-1}(U)$ de tout ouvert quasi-compact $U$ de $Y$
est quasi-compacte (cf. (I, 6.6.1)).
\end{definition}

Si $\mathfrak B$ est une base de la topologie de $Y$ formée d'ouverts
affines, pour que $f$ soit quasi-compact, il faut et il suffit que pour tout
$V\in\mathfrak B$, $f^{-1}(V)$ soit \emph{réunion finie d'ouverts affines}.
Par exemple, si $Y$ est affine et $X$ quasi-compact, \emph{tout} morphisme
$f:X\to Y$ est quasi-compact (I, 6.6.1).

Si $f:X\to Y$ est un morphisme quasi-compact, il est clair que pour tout
ouvert $V$ de $Y$, la restriction de $f$ à $f^{-1}(V)$ est un morphisme
quasi-compact $f^{-1}(V)\to V$. Inversement, si $(U_\alpha)$ est un
recouvrement ouvert de $Y$ et $f:X\to Y$ un morphisme tel que les restrictions
$f^{-1}(U_\alpha)\to U_\alpha$ soient quasi-compactes, alors $f$ est
quasi-compact. Par

\oldpage[IV-1]{225}

suite, si $f:X\to Y$ est un $S$-morphisme de $S$-préschémas, et s'il existe
un recouvrement ouvert $(S_\lambda)$ de $S$ tel que les restrictions
$g^{-1}(S_\lambda)\to h^{-1}(S_\lambda)$ de $f$ (où $g$ et $h$ sont les
morphismes structuraux) soient quasi-compactes, alors $f$ est quasi-compact.

\begin{proposition}[1.1.2]
\label{IV.1.1.2-fr}
\begin{enumerate}
  \item[(i)] Une immersion $X\to Y$ est quasi-compacte si elle est fermée, ou si
  l'espace sous-jacent à $Y$ est localement noethérien ou si l'espace
  sous-jacent à $X$ est noethérien.
  \item[(ii)] Le composé de deux morphismes quasi-compacts est quasi-compact.
  \item[(iii)] Si $f:X\to Y$ est un $S$-morphisme quasi-compact, il en est de même de
  $f_{(S')}:X_{(S')}\to Y_{(S')}$ pour toute extension $g:S'\to S$ du
  préschéma de base.
  \item[(iv)] Si $f:X\to X'$ et $g:Y\to Y'$ sont deux $S$-morphismes
  quasi-compacts,
  \[
    f\times_S g:X\times_S Y\to X'\times_S Y'
  \]
  est quasi-compact.
  \item[(v)] Si le composé $g\circ f$ de deux morphismes $f:X\to Y$, $g:Y\to Z$
  est quasi-compact, et si $g$ est séparé, ou l'espace sous-jacent à $X$
  localement noethérien, $f$ est quasi-compact.
  \item[(vi)] Pour qu'un morphisme $f$ soit quasi-compact, il faut et il suffit que
  $f_{\mathrm{red}}$ le soit.
\end{enumerate}
\end{proposition}

Pour la démonstration, voir (I, 6.6.4). On notera que l'assertion (vi) est
aussi conséquence de la proposition plus générale~:

\begin{proposition}[1.1.3]
\label{IV.1.1.3-fr}
Soient $f:X\to Y$, $g:Y\to Z$ deux morphismes. Si $g\circ f$ est
quasi-compact et $f$ surjectif, $g$ est quasi-compact.
\end{proposition}

\begin{proof}
En effet, si $V$ est un ouvert quasi-compact dans $Z$,
$f^{-1}(g^{-1}(V))$ est quasi-compact par hypothèse, et l'on a
$g^{-1}(V)=f(f^{-1}(g^{-1}(V)))$, puisque $f$ est surjectif~; donc
$g^{-1}(V)$ est quasi-compact.
\end{proof}

\begin{corollary}[1.1.4]
\label{IV.1.1.4-fr}
Soient $f:X\to Y$, $g:Y'\to Y$ deux morphismes, $X'=X\times_Y Y'$,
$f'=f_{(Y')}:X'\to Y'$~; supposons $g$ quasi-compact et surjectif (resp.
surjectif). Alors, pour que $f$ soit quasi-compact (resp. dominant), il suffit
que $f'$ le soit.
\end{corollary}

\begin{proof}
Si $g':X'\to X$ est la projection canonique, $g'$ est un morphisme surjectif
(I, 3.5.2), et l'on a $f\circ g'=g\circ f'$~; si $f'$ est quasi-compact
(resp. dominant), il en est de même de $g\circ f'$ puisque $g$ est
quasi-compact (1.1.2) (resp. surjectif)~; comme $f\circ g'$ est quasi-compact
(resp. dominant) et que $g'$ est surjectif, on en déduit que $f$ est
quasi-compact (1.1.3) (resp. dominant).
\end{proof}

\oldpage[IV-1]{226}

Pour abréger, nous appellerons \emph{point maximal} d'un préschéma $X$ tout
point générique d'une composante irréductible de $X$~; si
$X=\operatorname{Spec}(A)$ est affine, les points maximaux de $X$ sont les
\emph{idéaux premiers minimaux} de $A$.

\begin{proposition}[1.1.5]
\label{IV.1.1.5-fr}
Soit $f:X\to Y$ un morphisme quasi-compact. Les conditions suivantes sont
équivalentes~:
\begin{enumerate}
  \item[a)] $f$ est dominant~;
  \item[b)] pour tout point maximal $y$ de $Y$, on a
  $f^{-1}(y)\neq\varnothing$~;
  \item[c)] pour tout point maximal $y$ de $Y$, $f^{-1}(y)$ contient un point
  maximal de $X$.
\end{enumerate}
\end{proposition}

L'équivalence de \emph{a)} et \emph{c)} a été démontrée dans (I, 6.6.5). Il
est clair que \emph{c)} entraîne \emph{b)}~; d'autre part, si $X'$ est une
composante irréductible de $X$ rencontrant $f^{-1}(y)$, $f(X')$ est
irréductible et contient $y$, donc est contenu dans la composante irréductible
$Y'=\overline{\{y\}}$ de $Y$ (0\textsubscript{I}, 2.1.6)~; si $x$ est le
point générique de $X'$, on a donc nécessairement $f(x)=y$
(0\textsubscript{I}, 2.1.5)~; donc \emph{b)} entraîne \emph{c)}.

\begin{proposition}[1.1.6]
\label{IV.1.1.6-fr}
Soient $f':X'\to Y$, $f'':X''\to Y$ deux morphismes, $X$ le préschéma somme
$X'\coprod X''$, $f:X\to Y$ le morphisme égal à $f'$ dans $X'$ et à $f''$
dans $X''$. Pour que $f$ soit quasi-compact, il faut et il suffit que $f'$ et
$f''$ le soient.
\end{proposition}

Cela résulte aussitôt de la définition.

\subsection{Morphismes quasi-séparés}
\label{subsection:IV.1.2-fr}

\begin{definition}[1.2.1]
\label{IV.1.2.1-fr}
Soient $X$, $Y$ deux préschémas. On dit qu'un morphisme $f:X\to Y$ est
\emph{quasi-séparé} (ou que $X$ est \emph{quasi-séparé sur} $Y$) si le
morphisme diagonal
$\Delta_f=(1_X,1_X)_Y:X\to X\times_Y X$ est quasi-compact. On dit qu'un
préschéma $X$ est \emph{quasi-séparé} s'il est quasi-séparé sur
$\operatorname{Spec}(\mathbf Z)$.
\end{definition}

Par définition, tout morphisme séparé est quasi-séparé, $\Delta_f$ étant une
immersion fermée (1.1.2, (i))~; un \emph{schéma} $X$ est un préschéma
quasi-séparé, étant séparé sur $\operatorname{Spec}(\mathbf Z)$ (I, 5.4.1).

\begin{proposition}[1.2.2]
\label{IV.1.2.2-fr}
\begin{enumerate}
  \item[(i)] Tout monomorphisme de préschémas $f:X\to Y$ (en particulier toute
  immersion) est quasi-séparé.
  \item[(ii)] Si $f:X\to Y$ et $g:Y\to Z$ sont deux morphismes quasi-séparés,
  $g\circ f:X\to Z$ est quasi-séparé.
  \item[(iii)] Soient $X$, $Y$ deux $S$-préschémas, $f:X\to Y$ un $S$-morphisme
  quasi-séparé. Alors, pour tout changement de base $g:S'\to S$, le morphisme
  $f_{(S')}:X_{(S')}\to Y_{(S')}$ est quasi-séparé.
  \item[(iv)] Si $f:X\to Y$, $f':X'\to Y'$ sont deux $S$-morphismes quasi-séparés,
  alors
  \[
    f\times_S f':X\times_S X'\to Y\times_S Y'
  \]
  est quasi-séparé.
  \item[(v)] Si $f:X\to Y$, $g:Y\to Z$ sont deux morphismes tels que $g\circ f$
  soit quasi-séparé, alors $f$ est quasi-séparé.
  \item[(vi)] Si $f:X\to Y$ est un morphisme quasi-séparé,
  $f_{\mathrm{red}}:X_{\mathrm{red}}\to Y_{\mathrm{red}}$ est quasi-séparé.
\end{enumerate}
\end{proposition}

\begin{proof}
En vertu de (I, 5.5.12), il suffit de démontrer (i), (ii) et (iii).

L'assertion (i) est triviale, tout monomorphisme étant séparé (I, 5.5.1).
Pour démontrer (iii), on se ramène aussitôt au cas où $Y=S$ (I, 3.3.11) et
l'assertion résulte de ce que
$\Delta_{f_{(S')}}=(\Delta_f)_{(S')}$ (I, 5.3.4) et de (1.1.2, (iii)). Pour
démontrer (ii), considérons les projections $p$ et $q$ de $X\times_Y X$ sur
$X$~; si $i=(p,q)_Z$, on sait que le diagramme
\[
\xymatrix{
  X\times_Y X \ar[r]^-{i} \ar[d]_-{\pi} & X\times_Z X
    \ar[d]^-{f\times_Z f} \\
  Y \ar[r]_-{\Delta_g} & Y\times_Z Y
}
\]
\oldpage[IV-1]{227}
(où $\pi$ est le morphisme structural) est commutatif et identifie
$X\times_Y X$ au produit des $(Y\times_Z Y)$-préschémas $Y$ et
$X\times_Z X$ (I, 5.3.5). Si $g$ est quasi-séparé, $\Delta_g$ est
quasi-compact, donc il en est de même de $i$ (1.1.2, (iii))~; si de plus $f$
est quasi-séparé, $\Delta_f$ est quasi-compact, donc il en est de même de
$i\circ\Delta_f$ (1.1.2, (ii)) qui est égal à $\Delta_{g\circ f}$.
\end{proof}

\begin{corollary}[1.2.3]
\label{IV.1.2.3-fr}
\begin{enumerate}
  \item[(i)] Soit $X$ un préschéma quasi-séparé. Alors tout morphisme $f:X\to Y$
  est quasi-séparé.
  \item[(ii)] Soit $Y$ un préschéma quasi-séparé. Pour qu'un morphisme $f:X\to Y$
  soit quasi-séparé, il faut et il suffit que le préschéma $X$ soit
  quasi-séparé.
\end{enumerate}
\end{corollary}

Cela résulte aussitôt de (1.2.2, (ii) et (v)) appliqué à $X$, $Y$ et
$\operatorname{Spec}(\mathbf Z)$.

\begin{proposition}[1.2.4]
\label{IV.1.2.4-fr}
Soient $f:X\to Y$ un morphisme, $g:Y\to Z$ un morphisme quasi-séparé. Si
$g\circ f$ est quasi-compact, il en est de même de $f$.
\end{proposition}

\begin{proof}
On sait que $f$ est le morphisme composé
$X\xrightarrow{\Gamma_f}X\times_ZY\xrightarrow{\mathrm{pr}_2}Y$
(I, 5.3.13)~; d'autre part, $\mathrm{pr}_2$ s'identifie à
$(g\circ f)\times_Z1_Y$ (I, 3.3.4) et si $g\circ f$ est quasi-compact, il en
est de même de $\mathrm{pr}_2$ (1.1.2, (iv)). Enfin, on a le diagramme
commutatif
\[
\xymatrix{
  X \ar[r]^-{\Gamma_f} \ar[d]_-{f} & X\times_ZY
    \ar[d]^-{f\times 1_Y} \\
  Y \ar[r]_-{\Delta_g} & Y\times_ZY
}
\tag{1.2.4.1}
\label{IV.1.2.4.1-fr}
\]
qui identifie $X$ au produit des $(Y\times_ZY)$-préschémas $Y$ et
$X\times_ZY$ (I, 5.3.7). Comme par hypothèse $\Delta_g$ est un morphisme
quasi-compact, il en est de même de $\Gamma_f$ (1.1.2, (iii))~; la conclusion
résulte donc de (1.1.2, (ii)).
\end{proof}

\begin{proposition}[1.2.5]
\label{IV.1.2.5-fr}
Soient $f:X\to Y$, $g:Y'\to Y$ deux morphismes,
$X'=X\times_Y Y'$, $f'=f_{(Y')}:X'\to Y'$. Si $g$ est quasi-compact et
surjectif et $f'$ quasi-séparé, alors $f$ est quasi-séparé.
\end{proposition}

\begin{proof}
On a en effet $(X'\times_{Y'}X')=(X\times_YX)_{(Y')}$ et
$\Delta_{f'}=(\Delta_f)_{(Y')}$ (I, 5.3.4)~; comme la projection
$X'\times_{Y'}X'\to X\times_YX$ est quasi-compacte (1.1.2, (iii)) et
surjective (I, 3.5.2), on peut, en vertu de (I, 3.3.11), appliquer (1.1.4),
qui montre que puisque $\Delta_{f'}$ est un morphisme quasi-compact, il en est
de même de $\Delta_f$.
\end{proof}

\begin{proposition}[1.2.6]
\label{IV.1.2.6-fr}
Soient $f:X\to Y$ un morphisme, $(U_\alpha)$ un recouvrement de $Y$ par des
ouverts tels que les préschémas induits $U_\alpha$ soient quasi-séparés. Pour
que $f$ soit quasi-séparé, il faut et il suffit que chacun des préschémas
induits par $X$ sur les ouverts $f^{-1}(U_\alpha)$ soit quasi-séparé.
\end{proposition}

\begin{proof}
L'image réciproque dans $X\times_YX$ de $U_\alpha$ est
$X_\alpha\times_{U_\alpha}X_\alpha$, en posant
$X_\alpha=f^{-1}(U_\alpha)$, et la restriction
$X_\alpha\to X_\alpha\times_{U_\alpha}X_\alpha$ de $\Delta_f$ n'est autre
que $\Delta_{f_\alpha}$, en désignant par $f_\alpha$ la restriction
$X_\alpha\to U_\alpha$ de $f$~; en vertu de la définition (1.2.1) et du
caractère local de la notion de morphisme quasi-compact (1.1.1), pour que $f$
soit quasi-séparé, il faut et il suffit que chacun des morphismes $f_\alpha$
le soit. Mais comme par hypothèse le morphisme
$U_\alpha\to\operatorname{Spec}(\mathbf Z)$ est quasi-séparé, il revient au
même de dire que $f_\alpha$ est quasi-séparé ou que le composé
$X_\alpha\xrightarrow{f_\alpha}U_\alpha\to\operatorname{Spec}(\mathbf Z)$
l'est (1.2.2, (ii) et (v)), d'où la proposition.
\end{proof}

\oldpage[IV-1]{228}

Pour vérifier qu'un morphisme est quasi-séparé, on est donc ramené à donner
des critères pour qu'un préschéma soit quasi-séparé~:

\begin{proposition}[1.2.7]
\label{IV.1.2.7-fr}
Soient $X$ un préschéma, $(U_\alpha)$ un recouvrement de $X$ formé d'ouverts
quasi-compacts. Les propriétés suivantes sont équivalentes~:
\begin{enumerate}
  \item[a)] $X$ est quasi-séparé.
  \item[b)] Pour tout ouvert quasi-compact $U$ de $X$, l'immersion canonique
  $U\to X$ est quasi-compacte (autrement dit, $U$ est rétrocompact dans $X$).
  \item[b')] L'intersection de deux ouverts quasi-compacts de $X$ est
  quasi-compacte.
  \item[c)] Pour tout couple d'indices $(\alpha,\beta)$, l'intersection
  $U_\alpha\cap U_\beta$ est quasi-compacte.
\end{enumerate}
\end{proposition}

\begin{proof}
Les propriétés \emph{b)} et \emph{b')} sont trivialement équivalentes par
définition d'un morphisme quasi-compact. Comme un ouvert quasi-compact est
réunion finie d'ouverts affines, pour deux ouverts quasi-compacts $U$, $V$ de
$X$, $U\times_{\mathbf Z}V$ est un ouvert quasi-compact de
$X\times_{\mathbf Z}X$ (I, 3.2.7), dont l'image réciproque par $\Delta_X$ est
$U\cap V$, donc \emph{a)} entraîne \emph{b')}. Il est trivial que \emph{b')}
entraîne \emph{c)}~; enfin, si \emph{c)} est vérifiée, les
$U_\alpha\times_{\mathbf Z}U_\beta$ forment un recouvrement de
$X\times_{\mathbf Z}X$ par des ouverts quasi-compacts, et l'on sait que pour
que le morphisme $\Delta_X:X\to X\times_{\mathbf Z}X$ soit quasi-compact, il
suffit que les images réciproques $U_\alpha\cap U_\beta$ des
$U_\alpha\times_{\mathbf Z}U_\beta$ par $\Delta_X$ soient quasi-compactes
(1.1.1)~; donc \emph{c)} entraîne \emph{a)}.
\end{proof}

\begin{corollary}[1.2.8]
\label{IV.1.2.8-fr}
Tout préschéma $X$ dont l'espace sous-jacent est localement noethérien (par
exemple un préschéma localement noethérien) est quasi-séparé~; tout morphisme
$f:X\to Y$ est alors quasi-séparé.
\end{corollary}

\begin{proposition}[1.2.9]
\label{IV.1.2.9-fr}
Soient $X'$, $X''$ deux préschémas, $X=X'\coprod X''$ leur somme,
$f':X'\to Y$, $f'':X''\to Y$ deux morphismes, $f:X\to Y$ le morphisme qui
coïncide avec $f'$ dans $X'$ et avec $f''$ dans $X''$. Pour que $f$ soit
quasi-séparé il faut et il suffit que $f'$ et $f''$ le soient.
\end{proposition}

\begin{proof}
En effet, $X\times_YX$ est somme des quatre préschémas
$X'\times_YX'$, $X'\times_YX''$, $X''\times_YX'$ et $X''\times_YX''$, et
$\Delta_f$ est le morphisme qui coïncide avec $\Delta_{f'}$ dans $X'$ et avec
$\Delta_{f''}$ dans $X''$~; la proposition résulte donc aussitôt des
définitions.
\end{proof}

\subsection{Morphismes localement de type fini}
\label{subsection:IV.1.3-fr}

\begin{env}[1.3.1]
\label{IV.1.3.1-fr}
Rappelons que si $A$ est un anneau, une $A$-algèbre (commutative) $B$ est dite
\emph{de type fini} si elle est engendrée par un nombre fini d'éléments de
$B$, ou, ce qui revient au même, si elle est isomorphe au quotient d'une
algèbre de polynômes $A[T_1,\ldots,T_n]$ par un idéal de cette algèbre. Il est
clair que pour toute $A$-algèbre (commutative) $A'$,
$B\otimes_AA'$ est alors une $A'$-algèbre de type fini. Si $B$ est une
$A$-algèbre de type fini et $C$ une $B$-algèbre de type fini, $C$ est une
$A$-algèbre de type fini, car si $C$ est un quotient de
$B[T_1,\ldots,T_n]$ et $B$ un quotient de $A[S_1,\ldots,S_m]$,
$B[T_1,\ldots,T_n]=B\otimes_AA[T_1,\ldots,T_n]$ est un quotient de
$A[S_1,\ldots,S_m,T_1,\ldots,T_n]$, donc il en est de même de $C$.
\end{env}

\begin{definition}[1.3.2]
\label{IV.1.3.2-fr}
Soient $f:X\to Y$ un morphisme de préschémas, $x$ un point de $X$,
$y=f(x)$. On dit que $f$ est \emph{de type fini en $x$} s'il existe un
voisinage ouvert affine $V$ de $y$ et un voisinage ouvert affine $U$ de $x$
tels que $f(U)\subset V$ et que l'anneau $A(U)$ soit une $A(V)$-algèbre de
type fini. On dit que $f$ est \emph{localement de type fini} s'il est de type
fini en tout point de $X$.
\end{definition}

\oldpage[IV-1]{229}

\begin{proposition}[1.3.3]
\label{IV.1.3.3-fr}
Si $Y$ est un préschéma localement noethérien et $f:X\to Y$ un morphisme
localement de type fini, alors $X$ est localement noethérien.
\end{proposition}

Pour la démonstration, voir (I, 6.3.7).

\begin{proposition}[1.3.4]
\label{IV.1.3.4-fr}
\begin{enumerate}
  \item[(i)] Toute immersion locale est localement de type fini.
  \item[(ii)] Si deux morphismes $f:X\to Y$, $g:Y\to Z$ sont localement de
  type fini, il en est de même de $g\circ f$.
  \item[(iii)] Si $f:X\to Y$ est un $S$-morphisme localement de type fini,
  $f_{(S')}:X_{(S')}\to Y_{(S')}$ est localement de type fini pour toute
  extension $S'\to S$ du préschéma de base.
  \item[(iv)] Si $f:X\to X'$ et $g:Y\to Y'$ sont deux $S$-morphismes
  localement de type fini, $f\times_Sg$ est localement de type fini.
  \item[(v)] Si le composé $g\circ f$ de deux morphismes est localement de
  type fini, $f$ est localement de type fini.
  \item[(vi)] Si un morphisme $f$ est localement de type fini, il en est de
  même de $f_{\mathrm{red}}$.
\end{enumerate}
\end{proposition}

Pour la démonstration, voir (I, 6.6.6).

\begin{corollary}[1.3.5]
\label{IV.1.3.5-fr}
Soit $f:X\to Y$ un morphisme localement de type fini. Pour tout morphisme
$Y'\to Y$ tel que $Y'$ soit localement noethérien, $X\times_YY'$ est
localement noethérien.
\end{corollary}

\begin{proof}
En effet, $f_{(Y')}:X\times_YY'\to Y'$ est localement de type fini par
(1.3.4, (iii)), et il suffit d'appliquer (1.3.3).
\end{proof}

\begin{proposition}[1.3.6]
\label{IV.1.3.6-fr}
Soit $\varphi:A\to B$ un homomorphisme d'anneaux. Pour que le morphisme
correspondant $f:\operatorname{Spec}(B)\to\operatorname{Spec}(A)$ soit localement de type fini, il faut et
il suffit que $B$ soit une $A$-algèbre de type fini.
\end{proposition}

Pour la démonstration, voir (I, 6.3.3), en tenant compte de ce que $f$ est
quasi-compact et de (I, 6.6.3).

\begin{proposition}[1.3.7]
\label{IV.1.3.7-fr}
Pour qu'un morphisme localement de type fini $f:X\to Y$ soit surjectif, il
faut et il suffit que, pour tout corps algébriquement clos $\Omega$,
l'application $X(\Omega)\to Y(\Omega)$ correspondant à $f$ (I, 3.4.1) soit
surjective.
\end{proposition}

Pour la démonstration, voir (I, 6.3.10).

\begin{env}[1.3.8]
\label{IV.1.3.8-fr}
Soit $A$ un anneau. On dit qu'une $A$-algèbre $B$ est \emph{essentiellement de
type fini} si $B$ est $A$-isomorphe à une $A$-algèbre de la forme $S^{-1}C$,
où $C$ est une $A$-algèbre \emph{de type fini} et $S$ une partie multiplicative
de $C$.
\end{env}

\par\medskip\noindent
\begin{proposition}[1.3.9]
\label{IV.1.3.9-fr}
\begin{enumerate}
  \item[(i)] Si $B$ est une $A$-algèbre essentiellement de type fini et $C$
  une $B$-algèbre essentiellement de type fini, alors $C$ est une
  $A$-algèbre essentiellement de type fini.
  \item[(ii)] Soient $B$ une $A$-algèbre essentiellement de type fini, $A'$
  une $A$-algèbre~; alors $B'=B\otimes_AA'$ est une $A'$-algèbre
  essentiellement de type fini.
\end{enumerate}
\end{proposition}

\begin{proof}
\begin{enumerate}
  \item[(i)] Soient $B=S'^{-1}B'$ et $C=T'^{-1}C'$, où $B'$ (resp. $C'$) est
  une $A$-algèbre (resp. une $B$-algèbre) de type fini et $S'$ (resp. $T'$)
  une partie multiplicative de $B'$ (resp. $C'$)~; alors $C'$ est de la forme
  $S'^{-1}C''$, où $C''$ est une $B'$-algèbre de type fini, donc $C$ est aussi
  un anneau de fractions de $C''$~; comme $C''$ est une $A$-algèbre de type
  fini, cela prouve notre assertion.
  \item[(ii)] Si $B$ est anneau de fractions d'une $A$-algèbre de type fini
  $C$, $B'$ est un anneau de fractions de la $A'$-algèbre de type fini
  $C\otimes_AA'$, d'où (ii).
\end{enumerate}
\end{proof}

\oldpage[IV-1]{230}

\begin{env}[1.3.10]
\label{IV.1.3.10-fr}
Si $B$ est une $A$-algèbre \emph{locale} essentiellement de type fini, $B$ est
de la forme $C_{\mathfrak q}$, où $C$ est une $A$-algèbre de type fini et
$\mathfrak q$ un idéal premier de $C$ (Bourbaki, \emph{Alg. comm.}, chap. II,
\S~5, prop. 11). Soit $\mathfrak p$ l'image réciproque dans $A$ de l'idéal
$\mathfrak q$~; si l'on pose $S=A-\mathfrak p$, $C_{\mathfrak q}$ est aussi un
anneau local en un idéal premier de $S^{-1}C$~; comme $S^{-1}C$ est une
algèbre de type fini sur $A_{\mathfrak p}=S^{-1}A$, on voit que $B$ est aussi
une $A_{\mathfrak p}$-algèbre \emph{essentiellement de type fini}, et en outre
l'homomorphisme $A_{\mathfrak p}\to B$ est \emph{local}.
\end{env}

\begin{proposition}[1.3.11]
\label{IV.1.3.11-fr}
Si $B$ est une $A$-algèbre locale essentiellement de type fini, $B$ est
$A$-isomorphe à une $A$-algèbre quotient d'une $A$-algèbre de la forme
$C_{\mathfrak q}$, où $C$ est une algèbre de polynômes
$A[T_1,\ldots,T_n]$ et $\mathfrak q$ un idéal premier de $C$.
\end{proposition}

\begin{proof}
En effet, par définition, $B$ est isomorphe à $C'_{\mathfrak q'}$, où $C'$ est
une $A$-algèbre de type fini et $\mathfrak q'$ un idéal premier de $C'$~; mais
$C'=C/\mathfrak b$, où $C=A[T_1,\ldots,T_n]$ et $\mathfrak b$ est un idéal de
$C$~; donc $\mathfrak q'=\mathfrak q/\mathfrak b$, où $\mathfrak q$ est un
idéal premier de $C$, et $C'_{\mathfrak q'}$ est isomorphe à
$C_{\mathfrak q}/\mathfrak bC_{\mathfrak q}$.
\end{proof}

\subsection{Morphismes localement de présentation finie}
\label{subsection:IV.1.4-fr}

\begin{env}[1.4.1]
\label{IV.1.4.1-fr}
Étant donné un anneau $A$, nous dirons qu'une $A$-algèbre $B$ (commutative)
est \emph{de présentation finie} si elle est isomorphe au quotient
$A[T_1,\ldots,T_n]/\mathfrak a$ d'une algèbre de polynômes sur $A$ par un
idéal \emph{de type fini} $\mathfrak a$ de $A[T_1,\ldots,T_n]$. Pour toute
$A$-algèbre (commutative) $A'$, $B\otimes_AA'$ est alors une $A'$-algèbre de
présentation finie, étant isomorphe au quotient de $A'[T_1,\ldots,T_n]$ par
l'image canonique dans cet anneau de $\mathfrak a\otimes_AA'$, qui est
évidemment un $A'[T_1,\ldots,T_n]$-module de type fini. Si $B$ est une
$A$-algèbre de présentation finie et $C$ une $B$-algèbre de présentation
finie, alors $C$ est une $A$-algèbre de présentation finie. En effet, soient
$B=A[S_1,\ldots,S_m]/\mathfrak a$ et
$C=B[T_1,\ldots,T_n]/\mathfrak b$, où $\mathfrak a$ (resp. $\mathfrak b$) est
un idéal de type fini de $A[S_1,\ldots,S_m]$ (resp.
$B[T_1,\ldots,T_n]$)~; l'anneau $B[T_1,\ldots,T_n]$ est isomorphe à
$A[S_1,\ldots,S_m,T_1,\ldots,T_n]/\mathfrak a'$, où $\mathfrak a'$ est
l'image canonique de
$\mathfrak a\otimes_AA[T_1,\ldots,T_n]$, et est donc un idéal de type fini de
$A[S_1,\ldots,S_m,T_1,\ldots,T_n]$~; l'idéal $\mathfrak b$ de
$B[T_1,\ldots,T_n]$ est de la forme $\mathfrak b'/\mathfrak a'$, où
$\mathfrak b'$ est un idéal de $A[S_1,\ldots,S_m,T_1,\ldots,T_n]$~; comme
$\mathfrak a'$ et $\mathfrak b'/\mathfrak a'$ sont des modules de type fini
sur $A[S_1,\ldots,S_m,T_1,\ldots,T_n]$, il en est de même de $\mathfrak b'$,
et $C$, isomorphe à
$A[S_1,\ldots,S_m,T_1,\ldots,T_n]/\mathfrak b'$, est donc bien une
$A$-algèbre de présentation finie. Notons enfin que si $A$ est
\emph{noethérien}, il en est de même de $A[T_1,\ldots,T_n]$, donc toute
$A$-algèbre de type fini est alors de présentation finie.
\end{env}

\begin{definition}[1.4.2]
\label{IV.1.4.2-fr}
Soient $f:X\to Y$ un morphisme de préschémas, $x$ un point de $X$, $y=f(x)$.
On dit que $f$ est \emph{de présentation finie en $x$} s'il existe un voisinage
ouvert affine $V$ de $y$ et un voisinage ouvert affine $U$ de $x$ tels que
$f(U)\subset V$ et que l'anneau $A(U)$ soit une $A(V)$-algèbre de présentation
finie. On dit que $f$ est \emph{localement de présentation finie} s'il est de
présentation finie en tout point de $X$.
\end{definition}

Si $Y$ est localement noethérien, il revient au même de dire que $f$ est
localement de type fini ou localement de présentation finie.

\begin{proposition}[1.4.3]
\label{IV.1.4.3-fr}
\begin{enumerate}
  \item[(i)] Tout isomorphisme local est localement de présentation finie.
  \item[(ii)] Si deux morphismes $f:X\to Y$, $g:Y\to Z$ sont localement de
  présentation finie, il en est de même de $g\circ f$.

\oldpage[IV-1]{231}

  \item[(iii)] Si $f:X\to Y$ est un $S$-morphisme localement de présentation
  finie, $f_{(S')}:X_{(S')}\to Y_{(S')}$ est localement de présentation finie
  pour toute extension $S'\to S$ du préschéma de base.
  \item[(iv)] Si $f:X\to X'$ et $g:Y\to Y'$ sont deux $S$-morphismes
  localement de présentation finie, $f\times_Sg$ est localement de
  présentation finie.
  \item[(v)] Si le composé $g\circ f$ de deux morphismes est localement de
  présentation finie et si $g$ est localement de type fini, alors $f$ est
  localement de présentation finie.
\end{enumerate}
\end{proposition}

\begin{proof}
L'assertion (i) est triviale. Pour démontrer (iii), on peut se limiter au cas
où $S=Y$ (I, 3.3.11)~; soient $x'$ un point de $X_{(S')}$, $s'$ et $x$ ses
projections sur $S'$ et $X$ respectivement, $s$ la projection commune de $s'$
et $x$ sur $S$. Par hypothèse, il y a un voisinage ouvert affine $V$ (resp.
$U$) de $s$ (resp. $x$) tels que l'image de $U$ soit contenue dans $V$ et que
$A(U)$ soit une $A(V)$-algèbre de présentation finie~; soit $V'$ un voisinage
ouvert affine de $s'$ dont l'image dans $S$ soit contenue dans $V$~; alors
$U\times_VV'$ est un voisinage ouvert affine de $x'$ dont l'image dans $S'$
est contenue dans $V'$ et dont l'anneau $A(U)\otimes_{A(V)}A(V')$ est une
$A(V')$-algèbre de présentation finie (1.4.1).

Pour démontrer (ii), considérons un point $x\in X$, et soient $y=f(x)$,
$z=g(f(x))$. Il y a par hypothèse un voisinage ouvert affine $W=\operatorname{Spec}(C)$
(resp. $V=\operatorname{Spec}(B)$) de $z$ (resp. $y$) tels que $g(V)\subset W$ et que $B$
soit une $C$-algèbre de présentation finie. D'autre part, il y a un voisinage
ouvert affine $V'=\operatorname{Spec}(B')$ (resp. $U=\operatorname{Spec}(A)$) de $y$ (resp. $x$) tels que
$f(U)\subset V'$ et que $A$ soit une $B'$-algèbre de présentation finie. Soit
$s\in B$ tel que l'ouvert $D(s)=\operatorname{Spec}(B_s)$ dans $V$ soit un voisinage de $y$
contenu dans $V'$, et soit $U'=f^{-1}(D(s))\subset U$~; on a
$U'=\operatorname{Spec}(A\otimes_{B'}B_s)$, et $A\otimes_{B'}B_s$ est une $B_s$-algèbre de
présentation finie (1.4.1)~; d'aute part $B_s=B[T]/(1-sT)$ est une
$B$-algèbre de présentation finie par définition~; par suite (1.4.1),
$A\otimes_{B'}B_s$ est une $C$-algèbre de présentation finie, ce qui prouve
(ii).

On sait que (iv) résulte de (i), (ii) et (iii) (I, 3.5.1). Pour prouver (v),
considérons le diagramme commutatif
\[
\xymatrix{
X \ar[r]^{\Gamma_f} \ar[d] & X\times_ZY \ar[d]^{f\times 1} \\
Y \ar[r]_{\Delta_g} & Y\times_ZY
}
\]
qui identifie $X$ au produit des $(Y\times_ZY)$-préschémas $Y$ et
$X\times_ZY$ (I, 5.3.7, $\Delta_g$ étant le morphisme diagonal). Si l'on sait
que $\Delta_g$ est localement de présentation finie, il en résultera par
(iii) qu'il en est de même de $\Gamma_f$. D'autre part, on a la factorisation
de $f$ en
$X\xrightarrow{\Gamma_f}X\times_ZY\xrightarrow{p_2}Y$ (I, 5.3.13), où la
projection $p_2$ est égale à $(g\circ f)\times_Z1_Y$~; comme $g\circ f$ est
supposé localement de présentation finie, il en est de même de $p_2$ par
(iv), et il résultera finalement de (ii) que $f$ est aussi localement de
présentation finie. On voit donc qu'on est ramené à démontrer le
\end{proof}

\begin{corollary}[1.4.3.1]
\label{IV.1.4.3.1-fr}
Soit $g:Y\to Z$ un morphisme localement de type fini~; alors le morphisme
diagonal $\Delta_g:Y\to Y\times_ZY$ est localement de présentation finie.
\end{corollary}

\oldpage[IV-1]{232}

\begin{proof}
On peut se borner au cas où $Z=\operatorname{Spec}(A)$, $Y=\operatorname{Spec}(B)$, $B$ étant une
$A$-algèbre de type fini~; le morphisme $\Delta_g$ correspond alors à
l'\emph{homomorphisme d'augmentation} $\delta:B\otimes_AB\to B$ tel que
$\delta(x\otimes y)=xy$. Compte tenu de la définition (1.4.2), il suffit de
montrer que le noyau $\mathfrak J$ de $\delta$ est un idéal \emph{de type
fini}, ce qui résulte de (0, 20.4.4).
\end{proof}

\begin{proposition}[1.4.4]
\label{IV.1.4.4-fr}
Soient $A$ un anneau, $B$ une $A$-algèbre, $B'$ une algèbre de polynômes
$A[T_1,\ldots,T_n]$, $u:B'\to B$ un homomorphisme surjectif de
$A$-algèbres. Pour que $B$ soit une $A$-algèbre de présentation finie, il faut
et il suffit que le noyau $\mathfrak J$ de $u$ soit un idéal de type fini de
$B'$.
\end{proposition}

\begin{proof}
La condition est suffisante par définition (1.4.1). Pour voir qu'elle est
nécessaire, remarquons que le morphisme $g:\operatorname{Spec}(B')\to\operatorname{Spec}(A)$ est
localement de type fini~; si $B$ est une $A$-algèbre de présentation finie,
le morphisme $f:\operatorname{Spec}(B)\to\operatorname{Spec}(B')$ correspondant à $u$ est tel que
$g\circ f:\operatorname{Spec}(B)\to\operatorname{Spec}(A)$ soit localement de présentation finie, donc il
résulte de (1.4.3, (v)) que $f$ est aussi localement de présentation finie.
Or, pour démontrer que l'idéal $\mathfrak J$ est de type fini, il suffit de
prouver que $\widetilde{\mathfrak J}$ est un $\widetilde{B'}$-Module de type
fini (II, 6.1.4.1). Revenant à la définition (1.4.2), on est finalement ramené
à prouver le
\end{proof}

\begin{lemma}[1.4.4.1]
\label{IV.1.4.4.1-fr}
Soient $\mathfrak a$ un idéal d'un anneau $C$, $s$ un élément de $C$, tel que
$C_s/\mathfrak a_s$ soit une $C$-algèbre de présentation finie~; alors
$\mathfrak a_s$ est un idéal de type fini dans l'anneau $C_s$.
\end{lemma}

\begin{proof}
Par hypothèse, il existe une $C$-algèbre de polynômes
$C'=C[T_1,\ldots,T_n]$ et un $C$-homomorphisme surjectif
$v:C'\to C_s/\mathfrak a_s$ dont le noyau $\mathfrak b$ est de type fini.
Soit $p:C_s\to C_s/\mathfrak a_s$ l'homomorphisme canonique~; pour tout $i$,
il y a un $t_i\in C$ tel que $p(t_i/s^k)=v(T_i)$ (on peut supposer que
l'exposant $k$ est le même pour tous les indices $i$). Considérons alors la
$C_s$-algèbre $D=C_s[T_1,\ldots,T_n]$, et l'homomorphisme de $C_s$-algèbres
$w:D\to C_s$ tel que $w(T_i)=t_i/s^k$~; $w$ est évidemment surjectif, et il
en est par suite de même de $v'=p\circ w:D\to C_s/\mathfrak a_s$. D'autre
part, tout polynôme de $D$ peut s'écrire $P/s^m$, où
$P\in C'=C[T_1,\ldots,T_n]$, et l'on a
$v'(P/s^m)=(1/s^m)v(P)$~; comme $1/s$ est inversible dans $C_s$, la relation
$v'(P/s^m)=0$ dans $C_s/\mathfrak a_s$ équivaut à $v(P)=0$, et par suite le
noyau $\mathfrak b'$ de $v'$ est engendré par l'image canonique de
$\mathfrak b$ dans $D$, et \emph{a fortiori} est de type fini dans l'anneau
$D$. Mais
$\mathfrak b'=v'^{-1}(0)=w^{-1}(p^{-1}(0))=w^{-1}(\mathfrak a_s)$, et comme
$w$ est surjectif, $\mathfrak a_s=w(\mathfrak b')$, ce qui prouve que
$\mathfrak a_s$ est un idéal de type fini de $C_s$.
\end{proof}

\begin{corollary}[1.4.5]
\label{IV.1.4.5-fr}
Soient $X$, $Y$ deux préschémas, $j:X\to Y$ une immersion, $U$ un ouvert de
$Y$ tel que $j(X)$ soit fermé dans $U$, $\mathcal J$ l'idéal quasi-cohérent de
$\mathcal O_U$ qui définit le sous-préschéma fermé de $Y$ associé à $j$. Pour
que $j$ soit localement de présentation finie, il faut et il suffit que
$\mathcal J$ soit un $\mathcal O_U$-Module de type fini.
\end{corollary}

\par\medskip\noindent
\begin{proposition}[1.4.6]
\label{IV.1.4.6-fr}
Soit $\varphi:A\to B$ un homomorphisme d'anneaux. Pour que le morphisme
correspondant $f:\operatorname{Spec}(B)\to\operatorname{Spec}(A)$ soit localement de présentation finie,
il faut et il suffit que $B$ soit une $A$-algèbre de présentation finie.
\end{proposition}

\begin{proof}
La condition étant trivialement suffisante, il reste à prouver qu'elle est
nécessaire. Si $f$ est localement de présentation finie, il résulte déjà de
(1.3.6) que $B$ est une $A$-algèbre de type fini, donc il existe un
homomorphisme surjectif $u:B'=A[T_1,\ldots,T_n]\to B$ de $A$-algèbres~; il
résulte de (1.4.3, (v)) que l'immersion
$j:\operatorname{Spec}(B)\to\operatorname{Spec}(B')$ est localement

\oldpage[IV-1]{233}

de présentation finie~; donc (1.4.5), si $\mathfrak J$ est le noyau de $u$,
le $\widetilde{B'}$-Module $\widetilde{\mathfrak J}$ est de type fini, et par
suite (II, 6.1.4.1) $\mathfrak J$ est un idéal de type fini.
\end{proof}

\begin{proposition}[1.4.7]
\label{IV.1.4.7-fr}
Soit $\rho:A\to B$ un homomorphisme d'anneaux faisant de $B$ une $A$-algèbre
finie. Pour que $B$ soit une $A$-algèbre de présentation finie, il faut et il
suffit que $B$ soit un $A$-module de présentation finie.
\end{proposition}

Nous utiliserons le

\begin{lemma}[1.4.7.1]
\label{IV.1.4.7.1-fr}
Si $B$ est une $A$-algèbre finie, il existe un $A$-homomorphisme surjectif
d'algèbres $u:B'\to B$, où $B'$ est une $A$-algèbre finie et de présentation
finie qui est un $A$-module libre.
\end{lemma}

\begin{proof}
En effet, il y a un système fini $(a_i)_{1\leq i\leq m}$ de générateurs de la
$A$-algèbre $B$ tel que chaque $a_i$ vérifie une relation $F_i(a_i)=0$, où
$F_i\in A[X]$ est un polynôme unitaire de degré $>0$~; la $A$-algèbre quotient
$B'_i=A[X]/(F_i)$ est donc libre de rang fini sur $A$~; soit $c_i$ la classe de
$X$ dans $B'_i$. Il existe un $A$-homomorphisme d'algèbres $u_i:B'_i\to B$
tel que $u_i(c_i)=a_i$~; il suffit alors de prendre pour $B'$ le produit
tensoriel $B'_1\otimes_AB'_2\otimes_A\cdots\otimes_AB'_m$ et de considérer
l'homomorphisme $u_1\otimes u_2\otimes\cdots\otimes u_m:B'\to B$, qui est
surjectif par construction. En outre, si $B''=A[T_1,\ldots,T_m]$ et si $b'_i$
désigne l'image canonique de $c_i$ dans $B'$, le $A$-homomorphisme d'algèbres
$v:B''\to B'$ tel que $v(T_i)=b'_i$ ($1\leq i\leq m$) est surjectif, et son
noyau $\mathfrak b'$ est l'idéal de $B''$ engendré par les $F_i(T_i)$, donc de
type fini~; ceci prouve que $B'$ est une $A$-algèbre de présentation finie.
\end{proof}

\begin{proof}
Ce lemme étant démontré, gardons les mêmes notations, et soit
$w=v\circ u:B''\to B$. Si $\mathfrak b$ (resp. $\mathfrak a$) est le noyau de
$w$ (resp. $u$), on a $\mathfrak a=v(\mathfrak b)$ puisque $v$ est surjectif.
Or (1.4.4), si $B$ est une $A$-algèbre de présentation finie, $\mathfrak b$
est un idéal de type fini de $B''$, donc $\mathfrak a$ est un idéal de type
fini de $B'$, donc un $A$-module de type fini puisque $B'$ est une $A$-algèbre
finie~; comme $B'$ est un $A$-module libre, $B$ est un $A$-module de
présentation finie. Inversement, si $B'$ est un $A$-module de présentation
finie, $\mathfrak a$ est un $A$-module de type fini (Bourbaki,
\emph{Alg. comm.}, chap. I, \S~2, n\textsuperscript{o}~8, lemme 9) et
\emph{a fortiori} un idéal de type fini de $B'$~; par suite, $B$ est par
définition une $B'$-algèbre de présentation finie, et comme $B'$ est une
$A$-algèbre de présentation finie, $B$ est une $A$-algèbre de présentation
finie.
\end{proof}

\subsection{Morphismes de type fini}
\label{subsection:IV.1.5-fr}

\begin{proposition}[1.5.1]
\label{IV.1.5.1-fr}
Soient $f:X\to Y$ un morphisme, $(U_\alpha)$ un recouvrement de $Y$ formé
d'ouverts affines. Les conditions suivantes sont équivalentes~:
\begin{enumerate}
  \item[a)] $f$ est localement de type fini et quasi-compact.
  \item[b)] Pour tout $\alpha$, $f^{-1}(U_\alpha)$ est réunion finie d'ouverts
  affines $V_{\alpha i}$ tels que chaque anneau $A(V_{\alpha i})$ soit une
  $A(U_\alpha)$-algèbre de type fini.
  \item[c)] Pour tout ouvert affine $U$ de $Y$, $f^{-1}(U)$ est réunion finie
  d'ouverts affines $V_j$ tels que les anneaux $A(V_j)$ soient des
  $A(U)$-algèbres de type fini.
\end{enumerate}
\end{proposition}

Pour la démonstration, voir (I, 6.3.2 et 6.6.3).

\begin{definition}[1.5.2]
\label{IV.1.5.2-fr}
On dit qu'un morphisme $f$ est \emph{de type fini} s'il vérifie les conditions
équivalentes de la proposition (1.5.1).
\end{definition}

\oldpage[IV-1]{234}

\begin{proposition}[1.5.3]
\label{IV.1.5.3-fr}
Soit $f:X\to Y$ un morphisme de type fini~; si $Y$ est noethérien, il en est
de même de $X$.
\end{proposition}

Pour la démonstration, voir (I, 6.3.7).

\begin{proposition}[1.5.4]
\label{IV.1.5.4-fr}
\begin{enumerate}
  \item[(i)] Soit $j:X\to Y$ un morphisme d'immersion. Si $j$ est
  quasi-compact (ce qui est le cas si $j$ est une immersion fermée, ou si
  l'espace sous-jacent à $X$ est noethérien, ou si l'espace sous-jacent à $Y$
  est localement noethérien), $j$ est de type fini.
  \item[(ii)] Le composé de deux morphismes de type fini est de type fini.
  \item[(iii)] Si $f:X\to Y$ est un $S$-morphisme de type fini,
  $f_{(S')}:X_{(S')}\to Y_{(S')}$ est de type fini pour toute extension
  $S'\to S$ du préschéma de base.
  \item[(iv)] Si $f:X\to X'$ et $g:Y\to Y'$ sont deux $S$-morphismes de type
  fini, $f\times_Sg$ est de type fini.
  \item[(v)] Si le composé $g\circ f$ de deux morphismes est de type fini et
  si $g$ est quasi-séparé ou si $X$ est noethérien, $f$ est de type fini.
  \item[(vi)] Si un morphisme $f$ est de type fini, il en est de même de
  $f_{\mathrm{red}}$.
\end{enumerate}
\end{proposition}

Cela résulte aussitôt (sauf le cas (v) où $X$ est noethérien, démontré dans
(I, 6.3.6)) des définitions et de (1.1.2), (1.2.4) et (1.3.4).

\begin{corollary}[1.5.5]
\label{IV.1.5.5-fr}
Soit $f:X\to Y$ un morphisme de type fini. Pour tout morphisme $Y'\to Y$ tel
que $Y'$ soit noethérien, $X\times_YY'$ est noethérien.
\end{corollary}

Cela résulte de (1.5.4, (iii)) et de (1.5.3).

\begin{corollary}[1.5.6]
\label{IV.1.5.6-fr}
Soit $X$ un préschéma de type fini sur un préschéma localement noethérien $S$.
Alors tout $S$-morphisme $f:X\to Y$ est de type fini.
\end{corollary}

Pour la démonstration, voir (I, 6.3.9).

\begin{proposition}[1.5.7]
\label{IV.1.5.7-fr}
Soit $\varphi:A\to B$ un homomorphisme d'anneaux. Pour que le morphisme
correspondant $f:\operatorname{Spec}(B)\to\operatorname{Spec}(A)$ soit de type fini, il faut et il suffit
que $\varphi$ fasse de $B$ une $A$-algèbre de type fini.
\end{proposition}

Pour la démonstration, voir (I, 6.3.3).

\subsection{Morphismes de présentation finie}
\label{subsection:IV.1.6-fr}

\begin{definition}[1.6.1]
\label{IV.1.6.1-fr}
Soient $X$, $Y$ deux préschémas. On dit qu'un morphisme $f:X\to Y$ est
\emph{de présentation finie} s'il vérifie les trois conditions suivantes~:
\begin{enumerate}
  \item[(i)] $f$ est localement de présentation finie~;
  \item[(ii)] $f$ est quasi-compact (ce qui, lorsque (i) est vérifiée,
  équivaut à dire que $f$ est de type fini (1.5.2))~;
  \item[(iii)] $f$ est quasi-séparé.
\end{enumerate}
On dit alors aussi que $X$ est \emph{de présentation finie sur $Y$}, ou un
\emph{$Y$-préschéma de présentation finie}.
\end{definition}

La condition (iii) est automatiquement vérifiée si $f$ est séparé, ou si $X$
est localement noethérien (1.2.8)~; lorsque $Y$ est localement noethérien, il
revient au même de dire que $f$ est \emph{de présentation finie} ou que $f$
est \emph{de type fini} (cette dernière condition entraînant que $X$ est
localement noethérien (1.3.3)).

\oldpage[IV-1]{235}

\begin{proposition}[1.6.2]
\label{IV.1.6.2-fr}
\begin{enumerate}
  \item[(i)] Toute immersion quasi-compacte et localement de présentation
  finie (en particulier toute immersion ouverte quasi-compacte) est de
  présentation finie.
  \item[(ii)] Le composé de deux morphismes de présentation finie est de
  présentation finie.
  \item[(iii)] Soient $X$, $Y$ deux $S$-préschémas, $f:X\to Y$ un
  $S$-morphisme de présentation finie. Alors, pour tout morphisme $S'\to S$,
  le morphisme $f_{(S')}:X_{(S')}\to Y_{(S')}$ est de présentation finie.
  \item[(iv)] Soient $f:X\to Y$, $f':X'\to Y'$ deux $S$-morphismes de
  présentation finie~; alors
  $f\times_Sf':X\times_SX'\to Y\times_SY'$ est de présentation finie.
  \item[(v)] Si le composé $g\circ f$ de deux morphismes est de présentation
  finie et si $g$ est quasi-séparé et localement de présentation finie, $f$
  est de présentation finie.
\end{enumerate}
\end{proposition}

Cela résulte immédiatement de (1.1.2), (1.2.2), (1.2.4) et (1.4.3).

Il résulte en particulier de (1.6.2, (iii)) que si $f$ est un morphisme de
présentation finie et $U$ un ouvert de $Y$, la restriction
$f^{-1}(U)\to U$ de $f$ est un morphisme de présentation finie. Inversement,
soit $(U_\alpha)$ un recouvrement de $Y$ par des ouverts affines, et supposons
que les restrictions $f^{-1}(U_\alpha)\to U_\alpha$ de $f$ soient des
morphismes de présentation finie~; il en résulte alors que $f$ est de
présentation finie, car $f$ est évidemment localement de présentation finie,
quasi-compact en vertu de (1.1.1) et quasi-séparé en vertu de (1.2.6). On peut
donc dire que la propriété pour un morphisme $f:X\to Y$ d'être de présentation
finie est \emph{locale sur $Y$}.

Si $X$ est un préschéma quasi-séparé, tout morphisme $f:X\to Y$ est
quasi-séparé (1.2.3). Donc, si $f$ est quasi-compact et localement de
présentation finie, $f$ est alors de présentation finie.

En particulier~:

\begin{corollary}[1.6.3]
\label{IV.1.6.3-fr}
Soit $\varphi:A\to B$ un homomorphisme d'anneaux. Pour que le morphisme
correspondant $f:\operatorname{Spec}(B)\to\operatorname{Spec}(A)$ soit de présentation finie, il faut et
il suffit que $\varphi$ fasse de $B$ une $A$-algèbre de présentation finie.
\end{corollary}

La nécessité de la condition résulte en effet de (1.4.6).

\begin{remark}[1.6.4]
\label{IV.1.6.4-fr}
Dans la définition (1.6.1), la condition (iii) n'est pas conséquence des deux
autres. Soit par exemple $Y$ un schéma affine dont l'espace sous-jacent ne
soit pas noethérien, et soit $U$ un ouvert non quasi-compact dans $Y$. Soit
$X$ le préschéma obtenu en recollant deux préschémas $Y_1$, $Y_2$ isomorphes à
$Y$ le long des ouverts $U_1$, $U_2$ correspondant à $U$ (0$_{\mathrm I}$,
4.1.7), de sorte que $X$ est réunion de deux ouverts isomorphes respectivement
à $Y_1$ et $Y_2$ (et que nous identifions à ces derniers), avec
$Y_1\cap Y_2=U$. En outre, il y a un morphisme $f:X\to Y$ qui coïncide dans
$Y_i$ avec l'isomorphisme donné $Y_i\to Y$ ($i=1,2$). Il est clair que $f$
vérifie la condition (i) de (1.6.1), étant un isomorphisme local~; il vérifie
aussi (ii), l'image réciproque d'un ouvert quasi-compact de $Y$ étant la
réunion de ses images dans $Y_1$ et $Y_2$~; mais comme $Y_1\cap Y_2$ n'est pas
quasi-compact, $f$ n'est pas quasi-séparé (1.2.7).
\end{remark}

\begin{proposition}[1.6.5]
\label{IV.1.6.5-fr}
Soient $X'$, $X''$ deux préschémas, $X=X'\amalg X''$ leur somme,
$f':X'\to Y$, $f'':X''\to Y$ deux morphismes, $f:X\to Y$ le morphisme qui
coïncide avec $f'$ dans $X'$ et avec $f''$ dans $X''$. Pour que $f$ soit de
présentation finie, il faut et il suffit que $f'$ et $f''$ le soient.
\end{proposition}

\par\medskip\noindent
\begin{proof}
Il suffit de montrer que pour que $f$ possède une des trois propriétés de la
définition

\oldpage[IV-1]{236}

(1.6.1), il faut et il suffit que $f'$ et $f''$ la possèdent~; c'est évident
pour la propriété d'être localement de présentation finie, qui est locale sur
$X$~; pour la propriété d'être quasi-compact, cela a été vu en (1.1.6) et pour
la propriété d'être quasi-séparé, en (1.2.9).
\end{proof}

\subsection{Améliorations de résultats antérieurs}
\label{subsection:IV.1.7-fr}

Nous allons donner dans ce numéro une liste de propositions démontrées dans
les chapitres précédents, dont l'énoncé peut être amélioré à l'aide des
nouvelles conditions de finitude introduites ci-dessus.

\begin{env}[1.7.1]
\label{IV.1.7.1-fr}
Dans les énoncés de (I, 6.4.2, 6.4.3 et 6.4.9), on peut, au lieu de supposer
que $X$ et $Y$ sont de type fini sur le corps $K$, supposer seulement qu'ils
sont \emph{localement de type fini} sur $K$~; pour (I, 6.4.2), il suffit
d'observer que tout point d'un préschéma $X$ qui est fermé dans tout ouvert
affine de $X$ est fermé dans $X$, et un schéma affine localement de type fini
sur $K$ est \emph{ipso facto} de type fini (1.5.1). Pour (I, 6.4.9), on utilise
(1.3.7) et (1.3.4, (v)).
\end{env}

\begin{env}[1.7.2]
\label{IV.1.7.2-fr}
Dans (I, 6.5.1, (ii)), on peut au lieu de supposer que $S$ est localement
noethérien et $Y$ de type fini sur $S$, supposer seulement que $Y$ est
\emph{localement de présentation finie} sur $S$, comme le montre aussitôt la
démonstration (en appliquant la définition (1.4.2)). De même dans (I, 6.5.4,
(ii)) et (I, 6.5.5, (ii)) il suffit de supposer que $f$ est un $S$-morphisme,
$Y$ étant localement de présentation finie sur $S$.
\end{env}

\begin{env}[1.7.3]
\label{IV.1.7.3-fr}
Dans (I, 7.1.11), pour la deuxième assertion, au lieu de supposer $S$
localement noethérien et $Y$ de type fini sur $S$, on peut seulement supposer
$Y$ localement de présentation finie sur $S$ (la démonstration dépendant de
(I, 6.5.1, (ii)))~; de même dans (I, 7.1.12 à 7.1.15).
\end{env}

\begin{env}[1.7.4]
\EGAReaderTarget{ega4.c04.statement.1-7-4}{IV.1.7.4-fr}
Dans (I, 9.2.2), on peut remplacer les trois hypothèses par la seule hypothèse
(entraînée par chacune des autres) que $f$ est \emph{quasi-compact et
quasi-séparé}, en vertu de (1.2.6) et (1.2.7). On notera que dans (I, 9.2.1),
l'hypothèse signifie exactement que $f$ est quasi-compact et quasi-séparé.
\end{env}

\begin{env}[1.7.5]
\label{IV.1.7.5-fr}
Dans (I, 9.3.1, 9.3.2 et 9.3.3), on peut remplacer les deux hypothèses sur $X$
par l'hypothèse (moins forte) que $X$ est un préschéma \emph{quasi-compact et
quasi-séparé}, la démonstration utilisant exactement ces deux propriétés (en
vertu de (1.2.7)).
\end{env}

\begin{env}[1.7.6]
\label{IV.1.7.6-fr}
Dans (I, 9.3.4 et 9.3.5), il suffit de supposer $X$ quasi-compact,
$\mathcal F$ et $\mathcal J$ quasi-cohérents et de type fini.
\end{env}

\begin{env}[1.7.7]
\label{IV.1.7.7-fr}
Dans (I, 9.4.7), on peut affaiblir l'hypothèse b) en supposant seulement $X$
\emph{quasi-séparé}, puisqu'il suffit seulement que l'immersion canonique
$U\to X$ soit quasi-compacte (1.2.7). On peut faire la même modification dans
(I, 9.4.8, 9.4.9 et 9.4.10).
\end{env}

\begin{env}[1.7.8]
\label{IV.1.7.8-fr}
Dans (I, 9.5.1), les conditions indiquées entre parenthèses dans l'énoncé
peuvent être remplacées par la condition que $f$ soit quasi-compact et
quasi-séparé.
\end{env}

\begin{env}[1.7.9]
\label{IV.1.7.9-fr}
Dans (I, 9.6.5), on peut remplacer les hypothèses sur $X$ par l'hypothèse
moins forte que $X$ est quasi-compact et quasi-séparé, comme le montre la
démonstration,

\oldpage[IV-1]{237}

qui n'utilise que (I, 9.4.7). Dans (I, 9.6.6), l'hypothèse «~$X$ est un
schéma quasi-compact~» peut être remplacée par «~$X$ est un préschéma
quasi-compact et quasi-séparé~».
\end{env}

\begin{env}[1.7.10]
\label{IV.1.7.10-fr}
Dans (II, 1.7.15), on peut supposer seulement que $Y$ est quasi-compact et
quasi-séparé, la démonstration n'utilisant que (I, 9.6.5).
\end{env}

\begin{env}[1.7.11]
\label{IV.1.7.11-fr}
Dans la seconde assertion de (II, 3.4.5), on peut substituer aux deux
hypothèses sur $Y$ l'hypothèse plus faible que $Y$ est quasi-compact et
quasi-séparé. De même dans (II, 3.4.8).
\end{env}

\begin{env}[1.7.12]
\label{IV.1.7.12-fr}
Dans (II, 3.8.5), on peut de même supposer seulement que $Y$ est quasi-compact
et quasi-séparé, cette hypothèse intervenant par (I, 9.4.7).
\end{env}

\begin{env}[1.7.13]
\label{IV.1.7.13-fr}
Dans (II, 4.4.1, 4.4.6 et 4.4.7), il suffit de supposer que $Y$ est
\emph{quasi-compact et quasi-séparé}, compte tenu de (1.2.3) dans la
démonstration de (II, 4.4.7). De même, dans (II, 4.4.10.1), il suffit de
supposer $Z$ quasi-compact et quasi-séparé.
\end{env}

\begin{env}[1.7.14]
\label{IV.1.7.14-fr}
Dans (II, 4.5.2) et (II, 4.5.5), il suffit de supposer $X$
\emph{quasi-compact et quasi-séparé}, cette hypothèse intervenant par (I,
9.3.1 et 9.3.2).
\end{env}

\begin{env}[1.7.15]
\label{IV.1.7.15-fr}
Dans (II, 4.6.8), il suffit encore de supposer $X$ \emph{quasi-compact et
quasi-séparé}, cette hypothèse intervenant par (I, 9.4.7).
\end{env}

\begin{env}[1.7.16]
\label{IV.1.7.16-fr}
Dans (II, 5.1.2), il suffit de supposer que $X$ est quasi-compact et
quasi-séparé (l'hypothèse intervenant par (II, 4.5.2 et 4.5.5)). De même dans
(II, 5.1.9), où on utilise (I, 9.6.5).
\end{env}

\begin{env}[1.7.17]
\label{IV.1.7.17-fr}
Dans (II, 5.2.1), il suffit toujours de supposer $X$ quasi-compact et
quasi-séparé (l'hypothèse intervenant par (II, 4.5.2)).
\end{env}

\begin{env}[1.7.18]
\label{IV.1.7.18-fr}
Dans (II, 5.2.2), il faut supposer $f$ quasi-compact et $X$ quasi-séparé~; la
démonstration actuelle est en effet insuffisante (avec les hypothèses faites),
car du fait que pour tout $\mathcal O_X$-Module quasi-cohérent $\mathcal F$ on
a $\mathrm R^1f_*(\mathcal F)=0$, il ne s'ensuit pas nécessairement, pour un
ouvert affine $U$ de $Y$, que la même propriété soit valable pour les
$(\mathcal O_X|f^{-1}(U))$-Modules quasi-cohérents, si on ne sait pas qu'un tel
Module est la \emph{restriction} d'un $\mathcal O_X$-Module quasi-cohérent (la
même remarque s'appliquant aux conditions b) et c'))~; lorsque $f$ est
quasi-compact et $X$ quasi-séparé, ce prolongement est possible en vertu de
(I, 9.4.7), modifié plus haut (1.7.7).
\end{env}

\begin{env}[1.7.19]
\label{IV.1.7.19-fr}
Dans (II, 5.3.2), il suffit de supposer $Y$ \emph{quasi-compact et
quasi-séparé} (ce qui intervient par (II, 4.4.6 et 4.4.7)). De même dans (II,
5.5.3, (ii)).
\end{env}

\begin{env}[1.7.20]
\label{IV.1.7.20-fr}
Dans (III, 1.2.2, 1.2.3 et 1.2.4), on peut remplacer les deux hypothèses sur
$X$ par l'hypothèse moins forte que $X$ est un préschéma \emph{quasi-compact
et quasi-séparé}, la démonstration n'utilisant que (I, 9.3.3).
\end{env}

\begin{env}[1.7.21]
\label{IV.1.7.21-fr}
Dans (III, 1.4.10 à 1.4.14), on peut remplacer «~séparé~» par
«~quasi-séparé~», cette hypothèse servant uniquement à pouvoir appliquer (I,
9.2.2) (voir ci-dessus (1.7.4)). De même, dans (III, 1.4.15), il suffit de
supposer que le morphisme $u$ est \emph{quasi-séparé et quasi-compact}.
\end{env}

\begin{env}[1.7.22]
\label{IV.1.7.22-fr}
Dans (III, 2.1.3), on peut encore remplacer les hypothèses par l'hypothèse
moins forte que $X$ est un préschéma quasi-compact et quasi-séparé, le
raisonnement étant le même que dans (III, 1.2.2).
\end{env}

\subsection{Morphismes de présentation finie et ensembles constructibles}
\label{subsection:IV.1.8-fr}

\oldpage[IV-1]{238}

\begin{env}[1.8.1]
\label{IV.1.8.1-fr}
Soit $X$ un préschéma~; on sait que tout ouvert quasi-compact dans $X$ est
réunion finie d'ouverts affines, et réciproquement. Pour qu'un ouvert $U$ de
$X$ soit \emph{rétrocompact} dans $X$ (0\textsubscript{III}, 9.1.1), il faut
et il suffit donc que pour tout \emph{ouvert affine} $V$ de $X$, $U\cap V$
soit quasi-compact. Lorsque $X$ est \emph{quasi-séparé} (1.2.1), tout ouvert
quasi-compact dans $X$ est rétrocompact dans $X$ (1.2.7), donc toute partie
localement constructible de $X$ est rétrocompacte dans $X$
(0\textsubscript{III}, 9.1.13)~; si \emph{de plus} $X$ est
\emph{quasi-compact}, il y a identité entre parties constructibles et parties
localement constructibles dans $X$ (0\textsubscript{III}, 9.1.12), et entre
parties ouvertes constructibles et parties ouvertes quasi-compactes
(0\textsubscript{III}, 9.1.5).
\end{env}

\begin{proposition}[1.8.2]
\label{IV.1.8.2-fr}
Soient $X$, $Y$ deux préschémas, $f:X\to Y$ un morphisme. Pour toute partie
constructible (resp. localement constructible) $Z$ de $Y$, $f^{-1}(Z)$ est
une partie constructible (resp. localement constructible) de $X$.
\end{proposition}

\begin{proof}
Supposons $Z$ localement constructible~; pour tout $x\in X$, il y aura un
voisinage ouvert affine $V$ de $f(x)$ tel que $Z\cap V$ soit constructible
dans $V$~; si la proposition est démontrée pour les parties constructibles,
$f^{-1}(Z)\cap f^{-1}(V)$ sera constructible dans $f^{-1}(V)$, donc
$f^{-1}(Z)$ sera localement constructible. Il suffit par suite de considérer
le cas où $Z$ est constructible, et compte tenu de (0\textsubscript{III},
9.1.3), on est ramené au cas où $Z$ est ouvert et \emph{rétrocompact} dans
$Y$, autrement dit tel que l'injection canonique $Z\to Y$ soit un morphisme
quasi-compact~; il suffit alors de voir que $f^{-1}(Z)$ est
\emph{rétrocompact dans $X$}, autrement dit tel que l'injection canonique
$f^{-1}(Z)\to X$ soit un morphisme quasi-compact~; mais cela résulte de
(1.1.2, (iii)), puisque $f^{-1}(Z)=Z\times_YX$ (I, 4.4.1).
\end{proof}

\begin{lemma}[1.8.3]
\label{IV.1.8.3-fr}
Soient $X$ un préschéma quasi-compact et quasi-séparé, $Z$ une partie
constructible de $X$. Il existe alors un schéma affine $X'$ et un morphisme
de présentation finie $f:X'\to X$ tels que $f(X')=Z$.
\end{lemma}

\begin{proof}
Comme $X$ est quasi-compact, il est réunion finie d'ouverts affines $X_i$
($i\in I$), et comme $X$ est quasi-séparé, il résulte de (1.2.7) que chacune
des immersions canoniques $g_i:X_i\to X$ est de présentation finie~; on en
conclut que si $f_i:X'_i\to X_i$ est un morphisme de présentation finie, il
en est de même de $g_i\circ f_i:X'_i\to X$ (1.6.2), et si $X'$ est le
préschéma somme des $X'_i$, le morphisme $h:X'\to X$ qui coïncide avec
$g_i\circ f_i$ dans chaque $X'_i$ est lui aussi de présentation finie
(1.6.5). Comme $Z\cap X_i$ est constructible dans $X_i$
(0\textsubscript{III}, 9.1.8), on voit qu'on est ramené à prouver le lemme
lorsque $X_i$ est \emph{affine} d'anneau $A$. Comme $Z$ est alors réunion
finie d'ensembles de la forme $U\cap\complement V$, où $U$ et $V$ sont des
ouverts quasi-compacts (0\textsubscript{III}, 9.1.3), on voit, en considérant
encore une somme convenable de préschémas, qu'on peut se ramener au cas où
$Z=U\cap\complement V$~; comme d'ailleurs $U$ est réunion finie d'ouverts
affines, on peut même se borner au cas où $U$ est \emph{affine}. Comme $V$
est quasi-compact, il est réunion finie d'ouverts de la forme $X_{f_i}$, où
$f_i\in A$~; soit $\mathfrak J$ l'idéal de $A$ engendré par les $f_i$, et
soit $Z'$ le sous-schéma fermé affine de $X$ qu'il définit (et qui est par
construction de présentation finie sur $X$)~; on a par définition
$V=\complement Z'$, donc $U\cap\complement V=U\cap Z'$. Considérons le
schéma affine $X'=Z'\times_XU$

\oldpage[IV-1]{239}

et soit $f:X'\to X$ le morphisme structural~; on vient de voir que
l'immersion canonique $Z'\to X$ est de présentation finie, et il en est de
même de l'immersion ouverte $U\to X$, qui est quasi-compacte puisque $U$ et
$X$ sont affines~; on conclut donc de (1.6.2, (iv)) que $f$ est de
présentation finie, et l'on a $f(X')=Z'\cap U$ (I, 3.4.8), ce qui achève la
démonstration.
\end{proof}

\begin{theorem}[1.8.4]
\label{IV.1.8.4-fr}
\textup{(Chevalley).} --- Soit $f:X\to Y$ un morphisme de présentation
finie. Pour toute partie localement constructible $Z$ de $X$, $f(Z)$ est
localement constructible dans $Y$.
\end{theorem}

\begin{proof}
Soient $y\in Y$ et $V$ un voisinage ouvert affine de $y$~; comme $f$ est
quasi-compact et quasi-séparé, il en est de même de sa restriction
$f^{-1}(V)\to V$, donc $f^{-1}(V)$ est un préschéma quasi-compact et
quasi-séparé (1.2.3)~; comme $Z\cap f^{-1}(V)$ est constructible dans
$f^{-1}(V)$ (1.8.1), on voit qu'on peut se borner au cas où $Y$ est affine
et $Z$ constructible~; $X$ est alors quasi-compact et quasi-séparé, donc
(1.8.3), il existe un morphisme de présentation finie $g:X'\to X$ tel que
$Z=g(X')$~; comme $f\circ g$ est de présentation finie, on voit qu'on est
ramené au cas où $Z=X$~; autrement dit, il suffira de prouver le
\end{proof}

\begin{lemma}[1.8.4.1]
\label{IV.1.8.4.1-fr}
Soient $Y$ un schéma affine, $f:X\to Y$ un morphisme quasi-compact et
localement de présentation finie~; alors $f(X)$ est une partie constructible
de $Y$.
\end{lemma}

\begin{proof}
Comme $X$ est quasi-compact, il est réunion finie d'ouverts affines~; on peut
donc se borner au cas où $Y=\operatorname{Spec}(A)$,
$X=\operatorname{Spec}(B)$, $B$ étant une $A$-algèbre de présentation finie
(1.4.6). Or, on a le
\end{proof}

\begin{lemma}[1.8.4.2]
\label{IV.1.8.4.2-fr}
Soient $A_0$ un anneau, $(A_\lambda)_{\lambda\in L}$ un système inductif de
$A_0$-algèbres, $A=\varinjlim A_\lambda$~; si $B$ est une $A$-algèbre de
présentation finie, il existe un $\lambda$ et une $A_\lambda$-algèbre de
présentation finie $B_\lambda$ tels que $B$ soit isomorphe à
$B_\lambda\otimes_{A_\lambda}A$.
\end{lemma}

\begin{proof}
Par hypothèse, $B$ est isomorphe à une algèbre de la forme
$A[T_1,\ldots,T_n]/\mathfrak J$, où les $T_i$ sont des indéterminées et
$\mathfrak J$ un idéal de type fini~; soit $(F_j)$ ($1\leq j\leq m$) un
système de générateurs de $\mathfrak J$. Soit
$\varphi_\lambda:A_\lambda\to A$ l'homomorphisme canonique. Comme $L$ est
filtrant, il existe $\lambda$ tel que chacun des coefficients de chacun des
polynômes $F_j$ soit l'image par $\varphi_\lambda$ d'un élément de
$A_\lambda$. Autrement dit, il existe un système de polynômes
$(F_{j\lambda})_{1\leq j\leq m}$ de $A_\lambda[T_1,\ldots,T_n]$ tels que
$F_j=\varphi_\lambda(F_{j\lambda})$ pour $1\leq j\leq m$. Si
$\mathfrak J_\lambda$ est l'idéal de $A_\lambda[T_1,\ldots,T_n]$ engendré
par les $F_{j\lambda}$, $\mathfrak J$ est l'image de
$\mathfrak J_\lambda\otimes_{A_\lambda}A$ dans
$A[T_1,\ldots,T_n]=A_\lambda[T_1,\ldots,T_n]\otimes_{A_\lambda}A$~;
l'anneau $B_\lambda=A_\lambda[T_1,\ldots,T_n]/\mathfrak J_\lambda$ répond à
la question.
\end{proof}

On appliquera ici ce lemme en considérant $A$ comme limite inductive de ses
sous-$\mathbf Z$-algèbres \emph{de type fini}. On voit donc qu'il existe une
telle sous-$\mathbf Z$-algèbre $A_0$ de $A$, et une $A_0$-algèbre de
présentation finie $B_0$ telles que $B$ soit isomorphe à
$B_0\otimes_{A_0}A$~; posons $X_0=\operatorname{Spec}(B_0)$,
$Y_0=\operatorname{Spec}(A_0)$, de sorte que l'on a $X=X_0\times_{Y_0}Y$,
$f$ étant la projection $X\to Y$~; soient $f_0:X_0\to Y_0$,
$g_0:Y\to Y_0$ les morphismes structuraux~; il résulte de (I, 3.4.8) que
l'on a $f(X)=g_0^{-1}(f_0(X_0))$~; compte tenu de (1.8.2), il suffit donc de
montrer que $f_0(X_0)$ est constructible. Autrement dit, on est finalement
ramené à démontrer le

\begin{corollary}[1.8.5]
\label{IV.1.8.5-fr}
Soient $Y$ un préschéma noethérien, $f:X\to Y$ un morphisme de type fini.
Pour toute partie constructible $Z$ de $X$, $f(Z)$ est une partie
constructible de $Y$.
\end{corollary}

\oldpage[IV-1]{240}

\begin{proof}
On se ramène comme ci-dessus à prouver que $f(X)$ est constructible.
Appliquons le critère (0\textsubscript{III}, 9.2.3)~: il faut prouver que pour
toute partie fermée irréductible $T$ de $Y$, $T\cap f(X)=f(f^{-1}(T))$ est
rare dans $T$ ou contient un ouvert non vide de $T$~; en prenant un
sous-préschéma de $Y$ ayant $T$ pour espace sous-jacent (I, 5.2.1) et en
considérant son image réciproque par $f$, on voit finalement (compte tenu de
(1.5.4, (iii)) qu'on est ramené à prouver que si l'on suppose $Y$
\emph{irréductible et $f$ dominant}, $f(X)$ contient un ouvert non vide de
$Y$. On peut en outre supposer $Y$ affine~; alors $X$ est réunion finie
d'ouverts affines $V_i$, et comme $f(X)$ est dense dans $Y$, il en est de
même d'un au moins des $f(V_i)$. On peut donc aussi supposer $X$ affine~; si
$(X_j)$ est la famille (finie) des composantes irréductibles de $X$, on voit
comme ci-dessus qu'un au moins des $f(X_j)$ est dense dans $Y$~; on peut donc
supposer $X$ irréductible. Enfin, en remplaçant $f$ par $f_{\mathrm{red}}$
(1.5.4, (vi)) on peut supposer $X$ et $Y$ \emph{intègres}. Alors
$X=\operatorname{Spec}(B)$, $Y=\operatorname{Spec}(A)$, où $A$ est un
anneau intègre noethérien, $B$ une $A$-algèbre intègre de type fini contenant
$A$ (I, 1.2.7). Il suffit de montrer qu'il existe $g\in A$ tel que, pour tout
$y\in D(g)$ (c'est-à-dire tel que $g\notin\mathfrak p_y$) l'idéal
$\mathfrak p_y$ soit l'intersection de $A$ et d'un idéal premier de $B$, car
cela montrera que $D(g)\subset f(X)$. Finalement il suffit de prouver le
\end{proof}

\begin{lemma}[1.8.5.1]
\label{IV.1.8.5.1-fr}
Soient $A$ un anneau intègre, $B$ une $A$-algèbre intègre de type fini. Il
existe $g\in A$ tel que tout homomorphisme de $A$ dans un corps
algébriquement clos $\Omega$, non nul en $g$, se prolonge en un homomorphisme
de $B$ dans $\Omega$.
\end{lemma}

\begin{proof}
Or c'est là un résultat classique d'algèbre commutative (Bourbaki,
\emph{Alg. comm.}, chap. V, \S~3, n\textsuperscript{o}~1, cor. 3 du th. 1).
\end{proof}

\begin{corollary}[1.8.6]
\label{IV.1.8.6-fr}
Soient $Y$ un préschéma irréductible, $\eta$ son point générique,
$f:X\to Y$ un morphisme localement de type fini. Si $f^{-1}(\eta)$ n'est pas
vide, il existe un voisinage ouvert $U$ de $\eta$ dans $Y$ tel que
$f^{-1}(y)$ soit non vide pour tout $y\in U$. Si $\mathscr F$ est un
$\mathcal O_X$-Module quasi-cohérent de type fini, et si
$\mathscr F_\eta=\mathscr F\otimes_{\mathcal O_Y}k(\eta)$ n'est pas nul,
alors il existe un voisinage ouvert $U'$ de $\eta$ dans $Y$ tel que
$\mathscr F_y=\mathscr F\otimes_{\mathcal O_Y}k(y)$ ne soit pas nul pour
tout $y\in U'$.
\end{corollary}

\begin{proof}
Si $p_y$ est la projection canonique de la fibre
$f^{-1}(y)=X\times_Y\operatorname{Spec}(k(y))$ dans $X$, on a
$\mathscr F_y=p_y^*(\mathscr F)$, donc
$\operatorname{Supp}(\mathscr F_y)=p_y^{-1}(\operatorname{Supp}(\mathscr F))
=\operatorname{Supp}(\mathscr F)\cap f^{-1}(y)$ en vertu de (I, 9.1.13.1) et
(I, 3.6.1), puisque $\mathscr F$ est de type fini~; en outre
$\operatorname{Supp}(\mathscr F)$ est fermé dans $X$ (0\textsubscript{I},
5.2.2), et si $Z$ est un sous-préschéma fermé de $X$ ayant pour espace
sous-jacent $\operatorname{Supp}(\mathscr F)$ (I, 5.2.1), et $j$ l'immersion
canonique $Z\to X$, $f\circ j$ est localement de type fini (1.3.4)~; cela
prouve que la première assertion entraîne la seconde. Pour prouver la
première assertion, remarquons d'abord que l'on peut supposer $X$ et $Y$
réduits (1.3.4, (vi)), autrement dit $Y$ intègre. Soit $\xi$ un point de
$f^{-1}(\eta)$~; on peut remplacer $X$ et $Y$ par des voisinages ouverts
affines de $\xi$ et $\eta$ respectivement, et par suite supposer que
$X=\operatorname{Spec}(B)$, $Y=\operatorname{Spec}(A)$, $B$ étant une
$A$-algèbre de type fini. En outre, si $Z$ est le sous-préschéma fermé réduit
de $X$ ayant $\overline{\{\xi\}}$ pour ensemble sous-jacent (I, 5.2.1), on
peut, comme ci-dessus, remplacer $X$ par $Z$, autrement dit supposer $B$
\emph{intègre} (I, 5.1.4). Comme le morphisme $f$ est alors dominant,
l'homomorphisme correspondant $\varphi:A\to B$ est injectif (I, 1.2.7)~;
donc le corollaire résulte finalement du lemme (1.8.5.1).
\end{proof}

\oldpage[IV-1]{241}

\begin{proposition}[1.8.7]
\label{IV.1.8.7-fr}
Soient $S$ un préschéma, $g:X\to S$, $h:Y\to S$ deux morphismes de
présentation finie, $f:X\to Y$ un $S$-morphisme. Pour tout $s\in S$, posons
$X_s=g^{-1}(s)=X\times_S\operatorname{Spec}(k(s))$,
$Y_s=h^{-1}(s)=Y\times_S\operatorname{Spec}(k(s))$,
$f_s=f\times 1:X_s\to Y_s$. Alors l'ensemble des $s\in S$ tels que $f_s$
soit surjectif (resp. radiciel) est localement constructible.
\end{proposition}

\begin{proof}
Soit $E$ l'ensemble des $s\in S$ tels que $f_s$ soit surjectif~; l'ensemble
$S-E$ est égal à $h(Y-f(X))$~; or $f$ est de présentation finie (1.6.2, (v))
donc $f(X)$ est localement constructible dans $Y$ (1.8.4)~; comme $h$ est de
présentation finie, $h(Y-f(X))$ est localement constructible dans $S$, donc
aussi $E$.

Pour montrer que l'ensemble $E'$ des $s\in S$ tels que $f_s$ soit radiciel
est localement constructible, nous utiliserons le
\end{proof}

\begin{lemma}[1.8.7.1]
\EGAReaderTarget{ega4.c04.lem.1-8-7-1}{IV.1.8.7.1-fr}
Soient $U$, $V$ deux préschémas~; pour qu'un morphisme $f:U\to V$ soit
radiciel, il faut et il suffit que le morphisme diagonal
$\Delta_f:U\to U\times_VU$ soit surjectif. Par suite, tout morphisme radiciel
est séparé.
\end{lemma}

\begin{proof}
En effet, $\Delta_f$ étant une immersion (I, 5.3.9), est un morphisme
localement de type fini (1.3.4)~; dire que $\Delta_f$ est surjectif signifie
donc que pour tout corps algébriquement clos $K$, l'application
correspondante $U(K)\to(U\times_VU)(K)=U(K)\times_{V(K)}U(K)$ est surjective
(1.3.7 et I, 3.4.2.1)~; mais par définition d'un produit fibré d'ensembles,
cela signifie que l'application $U(K)\to V(K)$ correspondant à $f$ est
injective, et cela équivaut à dire que $f$ est radiciel (I, 3.5.5).

Cela étant, $\Delta_{f_s}$ se déduit de $\Delta_f:X\to X\times_YX$ par le
changement de base $\operatorname{Spec}(k(s))\to S$ (I, 5.3.4). Comme $f$
est de présentation finie, il en est de même du morphisme structural
$X\times_YX\to Y$ (1.6.2, (iv)), donc aussi de $\Delta_f$ (1.6.2, (v))~; il
suffit donc d'appliquer la première partie de la proposition à $\Delta_f$,
en utilisant le lemme (1.8.7.1).
\end{proof}

\subsection{Ensembles pro-constructibles et ind-constructibles}
\label{subsection:IV.1.9-fr}

\begin{lemma}[1.9.1]
\label{IV.1.9.1-fr}
Soit $(A_\alpha)_{\alpha\in I}$ un système inductif d'anneaux dont l'ensemble
d'indices $I$ est filtrant à droite, et soit $A=\varinjlim A_\alpha$. Pour que
$A=0$, il faut et il suffit qu'il existe un indice $\gamma$ tel que
$A_\gamma=0$ (et l'on a alors $A_\beta=0$ pour tout $\beta\geq\gamma$).
\end{lemma}

\begin{proof}
En effet, pour tout $\alpha\in I$, l'homomorphisme canonique
$\varphi_\alpha:A_\alpha\to A$ transforme l'élément unité en élément unité~;
dire que $A=0$ signifie que $\varphi_\alpha(1)=0$, ou encore
$\varphi_\alpha(1)=\varphi_\alpha(0)$~; on sait que cela équivaut à dire
qu'il existe un indice $\gamma\geq\alpha$ tel que l'homomorphisme
$\varphi_{\gamma\alpha}:A_\alpha\to A_\gamma$ soit tel que
$\varphi_{\gamma\alpha}(1)=\varphi_{\gamma\alpha}(0)$, et cette dernière
relation équivaut à $A_\gamma=0$, d'où le lemme.
\end{proof}

\begin{proposition}[1.9.2]
\label{IV.1.9.2-fr}
Soient $B$ un anneau, $(A_\alpha)_{\alpha\in I}$ un système inductif de
$B$-algèbres dont l'ensemble d'indices $I$ est filtrant à droite, et soit
$A=\varinjlim A_\alpha$~; posons $Y=\operatorname{Spec}(B)$,
$X_\alpha=\operatorname{Spec}(A_\alpha)$, $X=\operatorname{Spec}(A)$, et
soient $u:X\to Y$, $u_\alpha:X_\alpha\to Y$ les morphismes correspondant aux
homomorphismes structuraux $B\to A$, $B\to A_\alpha$. Alors~:
\end{proposition}

\begin{env}[1.9.2.i]
\label{IV.1.9.2.i-fr}
\emph{Pour que $X=\varnothing$, il faut et il suffit qu'il existe un
$\gamma\in I$ tel que $X_\gamma=\varnothing$, et l'on a alors
$X_\alpha=\varnothing$ pour $\alpha\geq\gamma$.}
\end{env}

\oldpage[IV-1]{242}

\begin{env}[1.9.2.ii]
\label{IV.1.9.2.ii-fr}
\emph{On a}
\[
  u(X)=\bigcap_{\alpha\in I}u_\alpha(X_\alpha). \tag{1.9.2.1}
  \label{IV.1.9.2.1-fr}
\]
\end{env}

L'assertion (i) n'est autre que la traduction de (1.9.1). Pour démontrer
(ii), notons que, puisque $u$ se factorise en
$X\to X_\alpha\xrightarrow{u_\alpha}Y$, le premier membre de (1.9.2.1) est
contenu dans le second. Inversement, soit $y\in Y-u(X)$~; on a
$u^{-1}(y)=\varnothing$, et $u^{-1}(y)$ est l'espace sous-jacent au
$k(y)$-préschéma $\operatorname{Spec}(A\otimes_Bk(y))$~; comme dans la
catégorie des $B$-modules le foncteur $\varinjlim$ commute au produit
tensoriel, $A\otimes_Bk(y)$ est limite inductive du système inductif d'anneaux
$A_\alpha\otimes_Bk(y)$~; le lemme (1.9.1) montre qu'il existe $\alpha$ tel
que $A_\alpha\otimes_Bk(y)=0$, c'est-à-dire
$u_\alpha^{-1}(y)=\varnothing$, donc $y\notin u_\alpha(X_\alpha)$. C.Q.F.D.

\begin{proposition}[1.9.3]
\label{IV.1.9.3-fr}
Soient $X$ un préschéma quasi-compact et quasi-séparé, $E$ une partie de $X$,
$(X_\alpha)_{\alpha\in I}$ un recouvrement ouvert affine de $X$, et pour
tout $\alpha\in I$, posons $E_\alpha=E\cap X_\alpha$. Les conditions
suivantes sont équivalentes~:
\begin{enumerate}
  \item[a)] Il existe un morphisme quasi-compact $f:X'\to X$ tel que
  $E=f(X')$.
  \item[a$'$)] Pour tout $\alpha$, il existe un morphisme quasi-compact
  $f_\alpha:X'_\alpha\to X_\alpha$ tel que
  $E_\alpha=f_\alpha(X'_\alpha)$.
  \item[a$''$)] Pour tout $\alpha$, il existe un schéma affine $X''_\alpha$
  et un morphisme $f'_\alpha:X''_\alpha\to X_\alpha$ tel que
  $f'_\alpha(X''_\alpha)=E_\alpha$.
  \item[b)] $E$ est intersection de parties constructibles de $X$.
  \item[b$'$)] $F=X-E$ est réunion de parties constructibles de $X$.
\end{enumerate}
\end{proposition}

\begin{proof}
Il est clair que $b)$ et $b')$ sont équivalentes. La condition $a)$ entraîne
$a')$ en prenant pour $X'_\alpha$ le préschéma induit sur l'ouvert
$f^{-1}(X_\alpha)$ de $X'$ et pour $f_\alpha$ la restriction de $f$. Pour
voir que $a')$ entraîne $a'')$, il suffit de remarquer que $X'_\alpha$ est
réunion d'ouverts affines $X''_{\alpha\lambda}$ ($\lambda\in L$)~; soit
$X''_\alpha$ le préschéma \emph{somme} des $X''_{\alpha\lambda}$
($\lambda\in L$), qui est un schéma affine~; si
$g_\alpha:X''_\alpha\to X'_\alpha$ est le morphisme coïncidant avec
l'identité dans chacun des $X''_{\alpha\lambda}$, il est clair que si l'on
pose $f'_\alpha=f_\alpha\circ g_\alpha$, on a
$E_\alpha=f'_\alpha(X''_\alpha)$ puisque $g_\alpha$ est surjectif. Pour
montrer que $a'')$ entraîne $a)$, notons qu'il y a une partie finie $J$ de
$I$ telle que les $X_\alpha$ d'indices $\alpha\in J$ recouvrent $X$~; soit
$X'$ le préschéma somme des $X''_\alpha$ pour $\alpha\in J$, et soit
$f:X'\to X$ le morphisme coïncidant avec $j_\alpha\circ f'_\alpha$ pour tout
$\alpha\in J$, où $j_\alpha:X_\alpha\to X$ est l'injection canonique. Comme
$E$ est réunion des $E\cap X_\alpha$ pour $\alpha\in J$, on a bien
$E=f(X')$ et il reste à voir que $f$ est un morphisme quasi-compact~; mais
l'hypothèse que $X$ est quasi-séparé entraîne que $j_\alpha$ est
quasi-compact (1.2.7), donc il en est de même de
$j_\alpha\circ f'_\alpha$ (1.1.1 et 1.1.2) et finalement de $f$ (1.1.6).

Reste donc à prouver l'équivalence de $a'')$ et de $b)$. Prouvons en premier
lieu que $a'')$ entraîne $b)$~: considérons une partie finie $J$ de $I$ telle
que les $X_\alpha$ pour $\alpha\in J$ recouvrent $X$~; il suffira de montrer
que, pour tout $\alpha\in J$, $E_\alpha$ est intersection de parties
constructibles $F_{\alpha\lambda}$ de $X_\alpha$ ($\lambda\in L_\alpha$)~;
en effet, tout $F_{\alpha\lambda}$ est aussi une partie constructible de $X$
en vertu de (0\textsubscript{III}, 9.1.8, (ii)), car l'hypothèse que $X$ est
quasi-séparé entraîne que tout $X_\alpha$ est rétrocompact dans $X$ (1.2.7)~;
$E$ étant réunion des $E_\alpha$ pour $\alpha\in J$, est intersection des
réunions (finies) $\bigcup_{\alpha\in J}F_{\alpha,\lambda(\alpha)}$, pour
tous les choix des $\lambda(\alpha)\in L_\alpha$~; ces réunions étant
constructibles dans $X$, cela établit notre assertion. On peut donc se borner
au cas

\oldpage[IV-1]{243}

où $X=\operatorname{Spec}(A)$ et où $E=f(X')$, où
$X'=\operatorname{Spec}(A')$ est affine, $f:X'\to X$ un morphisme
correspondant à un homomorphisme d'anneaux $A\to A'$. Notons maintenant que
$A'$ est limite inductive de la famille $(A'_\lambda)$ de ses
sous-$A$-algèbres \emph{de type fini}, ordonnée par inclusion~; si
$X'_\lambda=\operatorname{Spec}(A'_\lambda)$, $f$ se factorise en
$X'\to X'_\lambda\xrightarrow{f_\lambda}X$ et il résulte de (1.9.2, (ii))
que l'on a $f(X')=\bigcap_\lambda f_\lambda(X'_\lambda)$~; si l'on démontre
que $f_\lambda(X'_\lambda)$ est intersection de parties constructibles de
$X$, il en sera de même de $f(X')$. On peut donc se borner au cas où $A'$ est
une $A$-algèbre \emph{de type fini}. Or, on a le lemme suivant~:
\end{proof}

\begin{lemma}[1.9.3.1]
\label{IV.1.9.3.1-fr}
Soient $A$ un anneau, $A'$ une $A$-algèbre de type fini. Alors $A'$ est
$A$-isomorphe à une limite inductive filtrante $\varinjlim A''_\mu$ où les
$A''_\mu$ sont des $A$-algèbres de présentation finie (1.4.1).
\end{lemma}

\begin{proof}
En effet, on peut écrire $A'=B/\mathfrak J$, avec
$B=A[T_1,\ldots,T_n]$ et $\mathfrak J$ un idéal de $B$. Mais $\mathfrak J$
est limite inductive des idéaux \emph{de type fini} $\mathfrak J'_\mu$ de
$B$, contenus dans $\mathfrak J$, ordonnés par inclusion~; comme dans la
catégorie des $B$-modules le foncteur $\varinjlim$ est exact, on en conclut
que $A'=\varinjlim B/\mathfrak J'_\mu$ à un $A$-isomorphisme près.
\end{proof}

Le même raisonnement que ci-dessus permet alors de se ramener au cas où $A'$
est une $A$-algèbre de présentation finie~; mais alors $f(X')$ est
constructible dans $X$ en vertu du théorème de Chevalley (1.8.5), ce qui
prouve $b)$.

Démontrons enfin que $b)$ entraîne $a'')$. Si $E$ est intersection d'une
famille $(G_\lambda)_{\lambda\in L}$ de parties constructibles de $X$,
chaque $E_\alpha$ est intersection des $G_\lambda\cap X_\alpha$, qui sont
constructibles dans $X_\alpha$ (0\textsubscript{III}, 9.1.8, (i)), donc on
est ramené au cas où $X=\operatorname{Spec}(A)$ est affine.

On sait alors (1.8.3) que pour tout $\lambda\in L$ il existe un morphisme
$f_\lambda:X'_\lambda\to X$, où
$X'_\lambda=\operatorname{Spec}(A'_\lambda)$ est affine, tel que
$f_\lambda(X'_\lambda)=G_\lambda$. Il suffit donc de prouver le lemme
suivant~:

\begin{lemma}[1.9.3.2]
\label{IV.1.9.3.2-fr}
Soient $(A'_\lambda)_{\lambda\in L}$ une famille de $A$-algèbres,
$X=\operatorname{Spec}(A)$, $X'_\lambda=\operatorname{Spec}(A'_\lambda)$,
et soit $f_\lambda:X'_\lambda\to X$ le morphisme structural. Pour toute
partie finie $J$ de $L$, posons
$A'_J=\bigotimes_{\lambda\in J}A'_\lambda$ (produit tensoriel de
$A$-algèbres), $X'_J=\operatorname{Spec}(A'_J)$, et soit $f_J:X'_J\to X$ le
morphisme structural. On a alors
$f_J(X'_J)=\bigcap_{\lambda\in J}f_\lambda(X'_\lambda)$. Si
$A'=\varinjlim A'_J$, $J$ parcourant l'ensemble filtrant des parties finies
de $L$, $X'=\operatorname{Spec}(A')$, et si $f:X'\to X$ est le morphisme
structural, on a $f(X')=\bigcap_{\lambda\in L}f_\lambda(X'_\lambda)$.
\end{lemma}

\begin{proof}
La première assertion résulte de (I, 3.4.7)~; la deuxième résulte de (1.9.2,
(ii)), qui donne la relation $f(X')=\bigcap_Jf_J(X'_J)$.
\end{proof}

\begin{definition}[1.9.4]
\label{IV.1.9.4-fr}
Soit $X$ un espace topologique. On dit qu'une partie $E$ de $X$ est
\emph{pro-constructible} (resp. \emph{ind-constructible}) dans $X$ si, pour
tout $x\in X$, il existe un voisinage ouvert $U$ de $x$ dans $X$ tel que
$E\cap U$ soit intersection (resp. réunion) de parties localement
constructibles de $U$.
\end{definition}

Tenant compte de (1.2.7) et de (0\textsubscript{III}, 9.1.8 et 9.1.12), les
conditions équivalentes de la proposition (1.9.3) s'expriment donc en disant
que $E$ est \emph{pro-constructible} dans $X$, les ensembles localement
constructibles dans $X$ étant identiques aux ensembles constructibles dans
$X$ sous les hypothèses de (1.9.3).

\oldpage[IV-1]{244}

\begin{proposition}[1.9.5]
\label{IV.1.9.5-fr}
Soit $X$ un préschéma.
\begin{enumerate}
  \item[(i)] Pour qu'une partie $E$ de $X$ soit pro-constructible, il faut et
  il suffit que $X-E$ soit ind-constructible.
  \item[(ii)] Toute réunion finie et toute intersection de parties
  pro-constructibles de $X$ sont pro-constructibles. Toute intersection
  finie et toute réunion de parties ind-constructibles de $X$ sont
  ind-constructibles.
  \item[(iii)] Toute intersection (resp. réunion) de parties localement
  constructibles de $X$ est pro-constructible (resp. ind-constructible).
  Inversement, si $X$ est quasi-compact et quasi-séparé, toute partie
  pro-constructible (resp. ind-constructible) de $X$ est intersection (resp.
  réunion) de parties constructibles de $X$.
  \item[(iv)] Soit $(U_\alpha)$ un recouvrement ouvert de $X$. Pour qu'une
  partie $E$ de $X$ soit pro-constructible (resp. ind-constructible) dans
  $X$, il faut et il suffit que pour tout $\alpha$, $E\cap U_\alpha$ soit
  pro-constructible (rep. ind-constructible) dans $U_\alpha$.
  \item[(v)] Toute partie pro-constructible $E$ de $X$ est rétrocompacte
  dans $X$. Pour qu'une partie localement fermée $E$ de $X$ soit
  rétrocompacte dans $X$, il faut et il suffit qu'elle soit
  pro-constructible dans $X$.
  \item[(vi)] Soit $f:X'\to X$ un morphisme de préschémas~; pour toute partie
  pro-constructible (resp. ind-constructible) $E$ de $X$, $f^{-1}(E)$ est
  pro-constructible (resp. ind-constructible) dans $X'$.
  \item[(vii)] Soit $f:X'\to X$ un morphisme quasi-compact~; pour toute
  partie pro-constructible $E'$ de $X'$, $f(E')$ est pro-constructible dans
  $X$~; en particulier $f(X')$ est pro-constructible dans $X$.
  \item[(viii)] Soit $f:X'\to X$ un morphisme localement de présentation
  finie~; pour toute partie ind-constructible $E'$ de $X'$, $f(E')$ est
  ind-constructible dans $X$~; en particulier $f(X')$ est ind-constructible
  dans $X$.
  \item[(ix)] Supposons que $X$ soit quasi-compact et quasi-séparé~; alors,
  pour qu'une partie $E$ de $X$ soit pro-constructible, il faut et il suffit
  qu'il existe un schéma affine $X'$ et un morphisme (nécessairement
  quasi-compact) $f:X'\to X$ tels que $E=f(X')$.
\end{enumerate}
\end{proposition}

\begin{proof}
Les propriétés (i), (ii), (iv) et la première assertion de (iii) résultent
de la définition (1.9.4) et sont valables pour un espace topologique
quelconque, en utilisant la distributivité mutuelle de l'intersection et de
la réunion et le fait que si $T$ est localement constructible dans $X$,
$T\cap U$ est localement constructible dans $U$ pour tout ouvert $U$ de $X$.
Si $X$ est quasi-compact et quasi-séparé et $E$ est pro-constructible dans
$X$, il y a un recouvrement fini $(U_i)$ de $X$ formé d'ouverts affines tels
que $E\cap U_i$ soit intersection de parties constructibles $M_{i\lambda}$
de $U_i$ ($\lambda\in L_i$)~; les $M_{i\lambda}$ sont aussi constructibles
dans $X$ en vertu de (1.2.7) et (0\textsubscript{III}, 9.1.8, (ii)) et $E$ est
intersection des réunions finies $\bigcup_iM_{i,\lambda(i)}$ (où
$\lambda(i)\in L_i$ pour tout $i$), qui sont des parties constructibles dans
$X$~; cela démontre la seconde assertion de (iii). L'assertion (ix) résulte
de (iii) et de (1.9.3). Pour démontrer la première assertion de (v), on peut
se limiter au cas où $X$ est affine, et prouver alors que $E$ est
quasi-compact (0\textsubscript{III}, 9.1.1)~; mais comme $X$ est alors
quasi-séparé, il existe en vertu de (ix) un morphisme quasi-compact
$f:X'\to X$ tel que $E=f(X')$~; comme $X'$ est quasi-compact, il en est de
même de son image $E$ par une application continue.

Pour prouver (vi), on peut se borner, en vertu de (iv), au cas où $X$ est
affine~; la conclusion résulte alors de (iii) et de (1.8.2).

\oldpage[IV-1]{245}

De même, pour démontrer (vii), on peut se borner au cas où $X$ est affine~;
alors $X'$ est réunion finie d'ouverts affines $X'_i$ et $f(E')$ est réunion
des $f(E'\cap X'_i)$, si bien qu'on peut aussi supposer que $X'$ est affine,
en vertu de (ii)~; mais alors $E'=g(X'')$, où $g:X''\to X'$ est un morphisme
quasi-compact, en vertu de (ix), et l'on a $f(E')=f(g(X''))$, donc $f(E')$ est
pro-constructible puisque $g\circ f$ est quasi-compact.

On déduit de (vii) la démonstration de la seconde assertion de (v)~: en
effet, soit $E$ une partie localement fermée et rétrocompacte dans $X$, et
considérons un sous-préschéma de $X$ ayant $E$ pour espace sous-jacent (I,
5.2.1)~; l'injection canonique $j:E\to X$ est alors par hypothèse un
morphisme quasi-compact, et il résulte de (vii) que $E=j(E)$ est
pro-constructible dans $X$.

Enfin, pour prouver (viii), on peut encore supposer $X$ affine~; en outre, si
$(X'_\alpha)$ est un recouvrement ouvert affine de $X'$ tel que
$E'\cap X'_\alpha$ soit réunion de parties constructibles de $X'_\alpha$
((iii) et (iv)), $f(E')$ est réunion des $f(E'\cap X'_\alpha)$ et, en vertu
du théorème de Chevalley (1.8.4), chacune des $f(E'\cap X'_\alpha)$ est
réunion de parties constructibles de $X$, d'où la conclusion.
\end{proof}

\begin{env}[1.9.6]
\label{IV.1.9.6-fr}
\textbf{Exemples} (1.9.6). --- Toute partie \emph{finie} d'un préschéma $X$
est pro-constructible~: en effet, il suffit de considérer les parties
$\{x\}$ réduites à un seul point (1.9.5, (ii)) et $\{x\}$ est l'image de
$\operatorname{Spec}(k(x))$ par le morphisme canonique
$\operatorname{Spec}(k(x))\to X$, qui est quasi-compact~; d'où la conclusion
par (1.9.5, (vii)). Par contre, une partie $\{x\}$ réduite à un point n'est
pas nécessairement ind-constructible~; par exemple, soit $A$ un anneau
noethérien intègre ayant une infinité d'idéaux maximaux, et soit $x$ le point
générique de $X=\operatorname{Spec}(A)$~; si $\{x\}$ était ind-constructible
dans $X$, il serait constructible en vertu de (1.9.5, (iii)) et par suite
contiendrait un ouvert non vide de $X$ (0\textsubscript{III}, 9.2.2)~; mais
cela contredit l'hypothèse que $A$ possède une infinité d'idéaux maximaux, en
vertu du théorème d'Artin-Tate (0, 16.3.3).

Toute partie \emph{fermée} $Y$ d'un préschéma $X$ est pro-constructible, par
(1.9.5, (v)). Toute partie ouverte de $X$ est donc ind-constructible, mais
une partie ouverte $U$ de $X$ n'est pro-constructible que si elle est
\emph{rétrocompacte}, en vertu encore de (1.9.5, (v)).

Enfin, si $X$ est un préschéma \emph{noethérien}, donc quasi-séparé (1.2.8),
il résulte de (1.9.5, (iii)) et de (0\textsubscript{III}, 9.1.7) que les
parties ind-constructibles de $X$ sont les \emph{réunions de parties
localement fermées de $X$}.
\end{env}

\begin{theorem}[1.9.7]
\label{IV.1.9.7-fr}
Soient $X$ un préschéma quasi-compact, $E$ une partie ind-constructible de
$X$, $(F_\lambda)_{\lambda\in L}$ une famille de parties pro-constructibles
de $X$ telle que $\bigcap_{\lambda\in L}F_\lambda\subset E$~; alors il existe
une partie finie $J$ de $L$ telle que
$\bigcap_{\lambda\in J}F_\lambda\subset E$.
\end{theorem}

\begin{proof}
Notons que les ensembles $F=X-E$ et $F\cap F_\lambda$ sont
pro-constructibles, donc on est ramené au
\end{proof}

\begin{corollary}[1.9.8]
\label{IV.1.9.8-fr}
Soient $X$ un préschéma quasi-compact, $(F_\lambda)_{\lambda\in L}$ une
famille de parties pro-constructibles de $X$ dont l'intersection est vide~;
alors il existe une sous-famille finie $(F_\lambda)_{\lambda\in J}$ dont
l'intersection est vide.
\end{corollary}

\begin{proof}
Recouvrant $X$ par un nombre fini d'ouverts affines, et utilisant (1.9.5,
(iv)), on est ramené au cas où $X$ est affine~; en vertu de (1.9.5, (ix)), on
a $F_\lambda=f_\lambda(X'_\lambda)$,

\oldpage[IV-1]{246}

où $X'_\lambda=\operatorname{Spec}(A'_\lambda)$ est affine, et $f_\lambda$
est un morphisme $X'_\lambda\to X$~; alors, avec les notations de (1.9.3.2),
on a par hypothèse $f(X')=\varnothing$, donc $X'=\varnothing$~; mais cela
entraîne par (1.9.2, (i)) qu'il existe une partie finie $J$ de $L$ telle que
$X'_J=\varnothing$, donc
$f_J(X'_J)=\bigcap_{\lambda\in J}F_\lambda=\varnothing$.
\end{proof}

\begin{corollary}[1.9.9]
\label{IV.1.9.9-fr}
Soient $X$ un préschéma quasi-compact, $F$ une partie pro-constructible de
$X$, $(E_\lambda)_{\lambda\in L}$ une famille de parties ind-constructibles
de $X$ telle que $F\subset\bigcup_{\lambda\in L}E_\lambda$~; alors il existe
une partie finie $J$ de $L$ telle que
$F\subset\bigcup_{\lambda\in J}E_\lambda$. En particulier, de tout
recouvrement de $X$ par des parties ind-constructibles, on peut extraire un
recouvrement fini de $X$.
\end{corollary}

\begin{proof}
Cela résulte de (1.9.7) par passage aux complémentaires.
\end{proof}

\begin{proposition}[1.9.10]
\label{IV.1.9.10-fr}
Soit $X$ un préschéma. Pour qu'une partie $E$ de $X$ soit
ind-constructible, il est nécessaire que, pour tout $x\in X$, l'intersection
$E\cap\overline{\{x\}}$ soit un voisinage de $x$ dans
$\overline{\{x\}}$. Si $X$ est localement noethérien, cette condition est
suffisante.
\end{proposition}

\begin{proof}
Supposons $E$ ind-constructible, et soit $Y$ l'intersection de
$\overline{\{x\}}$ et d'un ouvert affine dans $X$, contenant $x$~; il y a
donc un sous-préschéma de $X$ ayant $Y$ pour espace sous-jacent (I, 5.2.1),
et si $j:Y\to X$ est l'injection canonique,
$E\cap Y=j^{-1}(E)$ est ind-constructible dans $Y$ (1.9.5, (vi)). Comme $Y$
est quasi-compact et séparé, $E\cap Y$ est donc réunion de parties
constructibles de $Y$ (1.9.5, (iii))~; par suite, il y a deux ouverts $U$,
$V$ dans $Y$ tels que
$x\in U\cap\complement V\subset E\cap Y$ (0\textsubscript{III}, 9.1.3).
Mais comme $x$ est point générique de l'espace irréductible $Y$, $V$ est
nécessairement vide et l'on a $U\subset E\cap Y$.

Supposons maintenant $X$ localement noethérien, et soit $E$ une partie de
$X$ vérifiant la condition de l'énoncé. En vertu de la définition (1.9.4),
on peut se borner au cas où $X$ est noethérien. Alors, pour tout $x\in E$,
$E\cap\overline{\{x\}}$ contient une partie non vide de la forme
$U\cap\overline{\{x\}}$, où $U$ est ouvert dans $X$~; cela montre que $E$
est réunion de parties localement fermées de $X$, donc est
ind-constructible (1.9.6).
\end{proof}

\begin{proposition}[1.9.11]
\label{IV.1.9.11-fr}
Soient $X$ un préschéma, $E$ une partie de $X$. Les propriétés suivantes sont
équivalentes~:
\begin{enumerate}
  \item[a)] $E$ est localement constructible.
  \item[b)] $E$ est ind-constructible et pro-constructible.
  \item[c)] $E$ et $X-E$ sont pro-constructibles.
  \item[d)] $E$ et $X-E$ sont ind-constructibles.
\end{enumerate}
\end{proposition}

\begin{proof}
Il suffit évidemment de prouver que $d)$ entraîne $a)$, et on peut se borner
au cas où $X$ est affine. Alors (1.9.5, (iii)), $E$ (resp. $X-E$) est réunion
d'une famille $(E_\lambda)$ (resp. $(E'_\mu)$) de parties constructibles de
$X$~; comme les $E_\lambda$ et les $E'_\mu$ forment un recouvrement de $X$,
il résulte de (1.9.9) qu'il y a des indices en nombre fini $\lambda_i$,
$\mu_j$ tels que les $E_{\lambda_i}$ et les $E'_{\mu_j}$ forment un
recouvrement de $X$~; cela implique que $E$ est réunion des
$E_{\lambda_i}$, donc est constructible.
\end{proof}

\begin{corollary}[1.9.12]
\label{IV.1.9.12-fr}
Soit $f:X\to Y$ un morphisme surjectif, qui est, soit quasi-compact, soit
localement de présentation finie. Pour qu'une partie $E$ de $Y$ soit
localement constructible (resp. pro-constructible, resp. ind-constructible),
il faut et il suffit que $f^{-1}(E)$ le soit dans $X$.
\end{corollary}

\oldpage[IV-1]{247}

\begin{proof}
On sait que la condition est nécessaire ((1.8.2) et (1.9.5, (vi)))~; pour
montrer qu'elle est suffisante, on est ramené au cas où $X$ est affine, en
vertu de (1.9.5, (iv))~; en outre, en vertu de (1.9.11), on peut se borner au
cas où $f^{-1}(E)$ est pro-constructible, \emph{ou} au cas où $f^{-1}(E)$ est
ind-constructible. Si $f$ est surjectif et quasi-compact, et $f^{-1}(E)$
pro-constructible, il résulte de (1.9.5, (vii)) que
$E=f(f^{-1}(E))$ est pro-constructible. Si $f$ est surjectif et localement de
présentation finie, et $f^{-1}(E)$ ind-constructible,
$E=f(f^{-1}(E))$ est ind-constructible par (1.9.5, (viii)), ce qui achève la
démonstration.
\end{proof}

\begin{env}[1.9.13]
\label{IV.1.9.13-fr}
Soient $X$ un préschéma, $\mathscr T$ sa topologie~; il résulte de (1.9.5,
(i) et (ii)) que les parties ind-constructibles (resp. pro-constructibles) de
$X$ sont les parties ouvertes (resp. fermées) pour une topologie sur $X$,
appelée \emph{topologie constructible} et que nous noterons
$\mathscr T^{\mathrm{cons}}$~; nous désignerons par $X^{\mathrm{cons}}$
l'ensemble $X$ muni de la topologie $\mathscr T^{\mathrm{cons}}$.
\end{env}

\begin{proposition}[1.9.14]
\label{IV.1.9.14-fr}
Soient $X$ un préschéma, $\mathscr T$ sa topologie,
$\mathscr T^{\mathrm{cons}}$ la topologie constructible sur $X$.
\begin{enumerate}
  \item[(i)] La topologie $\mathscr T^{\mathrm{cons}}$ est plus fine que
  $\mathscr T$.
  \item[(ii)] Les parties localement constructibles de $X$ sont identiques
  aux parties à la fois ouvertes et fermées de l'espace $X^{\mathrm{cons}}$.
  \item[(iii)] Pour tout morphisme $f:X\to Y$, l'application sous-jacente de
  $X^{\mathrm{cons}}$ dans $Y^{\mathrm{cons}}$ est continue~; on la note
  $f^{\mathrm{cons}}$.
  \item[(iv)] Si le morphisme $f:X\to Y$ est quasi-compact,
  $f^{\mathrm{cons}}$ est une application fermée~; en particulier, si $f$
  est quasi-compact et bijectif, $f^{\mathrm{cons}}$ est un homéomorphisme.
  \item[(v)] Si un morphisme $f:X\to Y$ est localement de présentation
  finie, $f^{\mathrm{cons}}$ est une application ouverte~; en particulier,
  si $f$ est bijectif et localement de présentation finie,
  $f^{\mathrm{cons}}$ est un homéomorphisme.
  \item[(vi)] Pour tout ouvert $U$ de $X$, la topologie induite par
  $\mathscr T^{\mathrm{cons}}$ sur $U$ est identique à la topologie de
  $U^{\mathrm{cons}}$.
\end{enumerate}
\end{proposition}

\begin{proof}
En effet, (i) résulte de ce que tout ouvert pour $\mathscr T$ est un ouvert
pour $\mathscr T^{\mathrm{cons}}$ (1.9.6), et (ii) est une traduction de
(1.9.11)~; (iii), (iv) et (v) traduisent respectivement (1.9.5), (vi), (vii)
et (viii). Enfin pour démontrer (vi), il suffit de remarquer que l'injection
canonique $j:U\to X$ est localement de présentation finie (1.4.3), donc
$j^{\mathrm{cons}}:U^{\mathrm{cons}}\to X^{\mathrm{cons}}$ est une
application continue ouverte.
\end{proof}

\begin{proposition}[1.9.15]
\label{IV.1.9.15-fr}
Soit $X$ un préschéma.
\begin{enumerate}
  \item[(i)] Pour que $X^{\mathrm{cons}}$ soit quasi-compact, il faut et il
  suffit que $X$ soit quasi-compact.
  \item[(ii)] Pour que $X^{\mathrm{cons}}$ soit séparé, il faut et il suffit
  que $X$ soit quasi-séparé, et alors $X^{\mathrm{cons}}$ est localement
  compact et totalement discontinu.
  \item[(iii)] Pour que $X^{\mathrm{cons}}$ soit compact, il faut et il suffit
  que $X$ soit quasi-compact et quasi-séparé.
  \item[(iv)] Tout point de $X$ admet, pour la topologie
  $\mathscr T^{\mathrm{cons}}$, un voisinage ouvert et compact.
  \item[(v)] Pour qu'un morphisme $f:X\to Y$ soit quasi-compact, il faut et il
  suffit que l'application continue
  $f^{\mathrm{cons}}:X^{\mathrm{cons}}\to Y^{\mathrm{cons}}$ soit propre.
\end{enumerate}
\end{proposition}

\begin{proof}
(i) Comme $\mathscr T^{\mathrm{cons}}$ est plus fine que $\mathscr T$, il est
clair que si $X^{\mathrm{cons}}$ est quasi-compact, il en est de même de $X$~;
la réciproque résulte de (1.9.9).

(ii) Supposons $X$ quasi-séparé, et montrons que $X^{\mathrm{cons}}$ est
totalement discontinu~: en effet, si $x$, $y$ sont deux points distincts de
$X$ et si par exemple $x\notin\overline{\{y\}}$, il existe un
\oldpage[IV-1]{248}
voisinage ouvert affine $U$ de $X$ ne contenant pas $y$~; comme $X$ est
quasi-séparé, $U$ est rétrocompact dans $X$ (1.2.7), donc $U$ et $X-U$ sont
localement constructibles dans $X$, et par suite ouverts dans
$X^{\mathrm{cons}}$ en vertu de (1.9.11), d'où notre assertion. Comme sur tout
ouvert affine $U$ de $X$, la topologie induite par $\mathscr T^{\mathrm{cons}}$
est celle de $U^{\mathrm{cons}}$ en vertu de (1.9.14, (vi)), il résulte de ce
qui précède que $X^{\mathrm{cons}}$ est localement compact, puisque $U$ est
ouvert dans $X^{\mathrm{cons}}$ et que $U^{\mathrm{cons}}$ est compact. Reste à
prouver que si $X^{\mathrm{cons}}$ est séparé, alors $X$ est quasi-séparé~;
considérons en effet un ouvert affine $U$ de $X$~; la topologie induite sur
$U$ par $\mathscr T^{\mathrm{cons}}$ est celle de $U^{\mathrm{cons}}$
(1.9.14, (vi)), donc $U$ est compact pour cette topologie induite, puisque
$U^{\mathrm{cons}}$ est compact par la première partie du raisonnement. Si
$V$ est un second ouvert affine dans $X$, $U\cap V$ est donc une partie
compacte de l'espace séparé $X^{\mathrm{cons}}$, étant intersection de deux
parties compactes de cet espace~; comme la topologie induite par
$\mathscr T^{\mathrm{cons}}$ sur $U\cap V$ est celle de
$(U\cap V)^{\mathrm{cons}}$ (1.9.14, (vi)) et que l'application identique
$(U\cap V)^{\mathrm{cons}}\to U\cap V$ est continue, on en conclut que
$U\cap V$ est quasi-compact pour la topologie induite par $\mathscr T$~; $X$
est par suite quasi-séparé en vertu de (1.2.7).

(iii) est un corollaire immédiat de (i) et (ii).

(iv) Pour tout $x\in X$, un voisinage ouvert affine $U$ de $x$ pour
$\mathscr T$ est aussi un voisinage de $x$ pour $\mathscr T^{\mathrm{cons}}$
et il est compact en vertu de (iii) et de (1.9.14, (vi)), la topologie induite
sur $U$ par $\mathscr T^{\mathrm{cons}}$ étant identique à celle de
$U^{\mathrm{cons}}$.

(v) Supposons $f$ quasi-compact~; alors on sait déjà (1.9.14, (iv)) que
$f^{\mathrm{cons}}$ est une application fermée. Soit d'autre part $y$ un
point de $Y$, et posons
\[
  Z=f^{-1}(y)=X\times_Y\operatorname{Spec}(k(y));
\]
$Z$ est quasi-compact (1.1.2, (iii)) et comme le morphisme canonique
$p:Z\to X$ est injectif, il résulte de (i) et de ce que l'application
$p^{\mathrm{cons}}$ est continue, que la topologie induite sur $f^{-1}(y)$
par celle de $X^{\mathrm{cons}}$ fait de $f^{-1}(y)$ un espace quasi-compact.
Cela prouve que $f^{\mathrm{cons}}$ est une application propre (Bourbaki,
\emph{Top. gén.}, chap. I\textsuperscript{er}, 3\textsuperscript{e} éd.,
\S~10, n\textsuperscript{o}~2, th.~1). Inversement, supposons que
l'application continue $f^{\mathrm{cons}}$ soit propre, et soit $V$ un ouvert
quasi-compact de $Y$~; si $h:V\to Y$ est l'injection canonique,
$h^{\mathrm{cons}}:V^{\mathrm{cons}}\to Y^{\mathrm{cons}}$ est continue et
injective et $V^{\mathrm{cons}}$ est quasi-compact par (i), donc la topologie
induite sur $V$ par celle de $Y^{\mathrm{cons}}$ fait de $V$ un espace
quasi-compact. L'hypothèse que $f^{\mathrm{cons}}$ est propre entraîne alors
que la topologie induite sur $f^{-1}(V)$ par celle de $X^{\mathrm{cons}}$ fait
de $f^{-1}(V)$ un espace quasi-compact (\emph{loc. cit.}, prop.~6), donc
$f^{-1}(V)$ est aussi un sous-espace quasi-compact de $X$, ce qui montre que
le morphisme $f$ est quasi-compact.
\end{proof}

\begin{env}[1.9.16]
\label{IV.1.9.16-fr}
Nous montrerons plus tard comment on peut, pour tout préschéma $X$, munir
l'espace $X^{\mathrm{cons}}$ d'un faisceau d'anneaux qui en fait un préschéma
dont les anneaux locaux sont tous des corps, identiques aux corps résiduels aux
points de $X$. De tels préschémas s'introduisent par exemple de façon naturelle
dans la traduction, en langage des schémas, des constructions de
Néron [32] dans sa théorie de la réduction des variétés abéliennes.
\end{env}

\oldpage[IV-1]{249}

\subsection{Application aux morphismes ouverts}
\label{subsection:IV.1.10-fr}

\begin{theorem}[1.10.1]
\label{IV.1.10.1-fr}
Soient $X$ un préschéma, $E$ une partie ind-constructible de $X$ (1.9.4), $x$
un point de $X$. Pour que $x$ soit intérieur à $E$, il faut et il suffit que
toute générisation ($0_{\mathrm I}$, 2.1.2) $x'$ de $x$ appartienne à $E$,
autrement dit (I, 2.4.2) que l'on ait
$\operatorname{Spec}(\mathcal O_{X,x})\subset E$.
\end{theorem}

\begin{proof}
La condition est évidemment nécessaire, tout voisinage de $x$ contenant les
générisations de $x$. Pour voir qu'elle est suffisante, on peut évidemment se
borner au cas où $X=\operatorname{Spec}(A)$ est affine, $x$ étant un idéal premier
$\mathfrak p$ de $A$. On sait alors (1.9.5, (ix)) qu'il existe un schéma affine
$X'=\operatorname{Spec}(A')$ et un morphisme $f:X'\to X$ tels que
$f(X')=X-E$. Posons $Y=\operatorname{Spec}(\mathcal O_{X,x})$ et
$Y'=X'\times_XY$~; l'hypothèse signifie (en
vertu de (I, 3.6.5)) que l'on a $Y'=\varnothing$, autrement dit
$A'\otimes_AA_{\mathfrak p}=0$. Comme
$A_{\mathfrak p}=\varinjlim A_t$, où $t$ parcourt $A-\mathfrak p$
($0_{\mathrm I}$, 1.4.5), on a $\varinjlim A_t'=0$, le foncteur
$\varinjlim$ commutant au produit tensoriel. Par suite (1.9.1), il existe un
$t\in A-\mathfrak p$ tel que $A_t'=0$, et $\mathrm D(t)$ est alors un
voisinage ouvert de $x$ dans $X$ contenu dans $E$.
\end{proof}

\begin{corollary}[1.10.2]
\label{IV.1.10.2-fr}
Soient $f:X\to Y$ un morphisme localement de présentation finie, et soit
$y\in Y$. Pour que $y$ soit intérieur à $f(X)$ il faut et il suffit que toute
générisation $y'$ de $y$ appartienne à $f(X)$.
\end{corollary}

\begin{proof}
On sait en effet (1.9.5, (viii)) que $f(X)$ est ind-constructible dans $Y$ et
il suffit d'appliquer (1.10.1).
\end{proof}

Nous dirons qu'une application continue $f:X\to Y$ est \emph{ouverte en un
point} $x\in X$ si l'image par $f$ de \emph{tout} voisinage de $x$ dans $X$
est un voisinage de $f(x)$ dans $Y$.

\begin{proposition}[1.10.3]
\EGAReaderTarget{ega4.c04.prop.1-10-3}{IV.1.10.3-fr}
Soient $f:X\to Y$ un morphisme, $x$ un point de $X$, $y=f(x)$. Considérons
les conditions suivantes~:
\begin{enumerate}
  \item[a)] $f$ est ouvert au point $x$.
  \item[b)] Pour toute générisation $y'$ de $y$, il existe une générisation
  $x'$ de $x$ telle que $f(x')=y'$.
  \item[b')] L'image par $f$ de $\operatorname{Spec}(\mathcal O_{X,x})$ est
  $\operatorname{Spec}(\mathcal O_{Y,y})$.
  \item[c)] Pour toute partie fermée irréductible $Y'$ de $Y$ contenant $y$,
  il existe une composante irréductible de $X'=f^{-1}(Y')$ contenant $x$ et
  dominant $Y'$.
\end{enumerate}
Alors on a les implications
$\mathrm{a)}\Rightarrow\mathrm{b)}\Leftrightarrow\mathrm{b')}
\Leftrightarrow\mathrm{c)}$. Si en outre $f$ est localement de présentation
finie, les quatre conditions sont équivalentes.
\end{proposition}

\begin{proof}
Par définition d'une générisation, l'image par $f$ de
$\operatorname{Spec}(\mathcal O_{X,x})$ est toujours contenue dans
$\operatorname{Spec}(\mathcal O_{Y,y})$, donc b) et b') sont équivalentes. Il
est immédiat
que c) entraîne b), car si $y'$ est une générisation de $y$,
$Y'=\overline{\{y'\}}$ est une partie irréductible de $Y$ contenant $y$~; si
$x'$ est le point générique d'une composante irréductible de $f^{-1}(Y')$ qui
contient $x$ et domine $Y'$, on a $f(x')=y'$ ($0_{\mathrm I}$, 2.1.5) et $x'$
est une générisation de $x$. Inversement, b) entraîne c)~: en effet, soit
$y'$ le point générique de $Y'$ et soit $x'$ une générisation de $x$ telle
que $f(x')=y'$~; soit $Z'$ une composante irréductible de $f^{-1}(Y')$
contenant $x'$ (donc $x$)~; comme $f(x')$ est déjà dense dans $Y'$, il en est
ainsi \emph{a fortiori} de $f(Z')$.

Pour voir que a) entraîne b'), on peut évidemment se borner au cas où
$X=\operatorname{Spec}(B)$ et $Y=\operatorname{Spec}(A)$ sont affines, $x$
étant un idéal premier
$\mathfrak q$ de $B$, de sorte que
$\mathcal O_{X,x}=B_{\mathfrak q}=\varinjlim B_t$,
\oldpage[IV-1]{250}
où $t$ parcourt $B-\mathfrak q$. On a, par suite (1.9.2),
\[
  f(\operatorname{Spec}(\mathcal O_{X,x}))
  =\bigcap_t f(\operatorname{Spec}(B_t))=\bigcap_t f(\mathrm D(t)),
\]
et par hypothèse, pour tout $t\in B-\mathfrak q$, $y$ est intérieur à
$f(\mathrm D(t))$, donc $f(\mathrm D(t))$ contient
$\operatorname{Spec}(\mathcal O_{Y,y})$~; d'où
$\operatorname{Spec}(\mathcal O_{Y,y})
\subset f(\operatorname{Spec}(\mathcal O_{X,x}))$, ce qui établit
notre assertion. Enfin, lorsque $f$ est localement de présentation finie, b)
implique a) en vertu de (1.10.2) appliqué à la restriction $U\to Y$ de $f$ à
un voisinage ouvert arbitraire $U$ de $x$.
\end{proof}

\begin{corollary}[1.10.4]
\label{IV.1.10.4-fr}
Soit $f:X\to Y$ un morphisme. Considérons les conditions suivantes~:
\begin{enumerate}
  \item[a)] $f$ est ouvert.
  \item[b)] Pour tout $x\in X$ et toute générisation $y'$ de $y=f(x)$, il
  existe une générisation $x'$ de $x$ telle que $f(x')=y'$.
  \item[b')] Pour tout $x\in X$, l'image par $f$ de
  $\operatorname{Spec}(\mathcal O_{X,x})$ est
  $\operatorname{Spec}(\mathcal O_{Y,f(x)})$.
  \item[c)] Pour toute partie fermée irréductible $Y'$ de $Y$, toute
  composante irréductible de $f^{-1}(Y')$ domine $Y'$.
\end{enumerate}
Alors on a les implications
$\mathrm{a)}\Rightarrow\mathrm{b)}\Leftrightarrow\mathrm{b')}
\Leftrightarrow\mathrm{c)}$. Si en outre $f$ est localement de présentation
finie, les quatre conditions sont équivalentes.
\end{corollary}

\begin{proof}
Dire que $f$ est ouvert signifie que $f$ est ouvert en tout point $x\in X$,
donc les implications
$\mathrm{a)}\Rightarrow\mathrm{b)}\Leftrightarrow\mathrm{b')}$ résultent des
implications analogues dans (1.10.3), ainsi que l'implication
$\mathrm{b)}\Rightarrow\mathrm{a)}$ lorsque $f$ est localement de présentation
finie. L'implication $\mathrm{c)}\Rightarrow\mathrm{b)}$ résulte aussi de
l'implication analogue dans (1.10.3)~; prouvons enfin que b) entraîne c). En
effet, soit $x'$ le point générique d'une composante
irréductible de $f^{-1}(Y')$, et montrons que $y'=f(x')$ est le point générique
de $Y'$. Soit $y''$ une générisation de $y'$~; il existe par hypothèse une
générisation $x''$ de $x'$ telle que $f(x'')=y''$, et comme
$x''\in f^{-1}(Y')$, on a nécessairement $x''=x'$, donc $y''=y'$, ce qui
achève la démonstration.
\end{proof}

\begin{center}
\emph{(A suivre.)}
\end{center}

\clearpage
\oldpage[IV-1]{251}
\phantomsection
\addcontentsline{toc}{section}{Bibliographie (suite)}
\section*{BIBLIOGRAPHIE (\emph{suite})}

\begingroup
\small
\setlength{\parindent}{0pt}
\setlength{\parskip}{3pt}
[30] M.~Nagata, \emph{Local rings}, \emph{Interscience Tracts}, 13 (1962),
Interscience, New York.

[31] M.~Nagata, \emph{A Jacobian criterion of simple points},
\emph{Ill. Journ. of Math.}, 1 (1957), p.~427--432.

[32] A.~Néron, \emph{Modèles minimaux des variétés abéliennes sur les corps
locaux et globaux}, \emph{Publ. Inst. Hautes Études Sci.}, n\textsuperscript{o}
21 (1964).

\clearpage
\oldpage[IV-1]{252}
\phantomsection
\addcontentsline{toc}{section}{Index des notations}
\section*{INDEX DES NOTATIONS}

$\dim_c(X)$, $\dim(X)$ ($X$ espace topologique) : 0, 14.1.2.\par
$\dim_x(X)$ ($X$ espace topologique, $x$ point de $X$) : 0, 14.1.2.\par
$\operatorname{codim}(Y,X)$ ($X$ espace topologique, $Y$ partie fermée de $X$) : 0, 14.2.1.\par
$\operatorname{codim}_x(Y,X)$ ($X$ espace topologique, $Y$ partie fermée de $X$, $x$ point de $X$) : 0, 14.2.4.\par
$N[T_1,\ldots,T_n]$ ($N$ $A$-module) : 0, 15.1.1.\par
$\dim(A)$ ($A$ anneau) : 0, 16.1.1.\par
$\operatorname{ht}(\mathfrak a)$ ($\mathfrak a$ idéal) : 0, 16.1.3.\par
$\dim_A(M)$, $\dim(M)$ ($M$ $A$-module) : 0, 16.1.7.\par
$P_{\mathfrak q}(M,n)$ ($M$ module de type fini sur un anneau semi-local noethérien $A$, $\mathfrak q$ idéal de définition de $A$) : 0, 16.2.1.\par
$\operatorname{prof}_A(M)$, $\operatorname{prof}(M)$ ($M$ $A$-module) : 0, 16.4.5.\par
$\operatorname{coprof}_A(M)$, $\operatorname{coprof}(M)$ ($M$ $A$-module) : 0, 16.4.9.\par
$\operatorname{dim.proj}_A(M)$, $\operatorname{dim.proj}(M)$,
$\operatorname{dim.inj}_A(M)$, $\operatorname{dim.inj}(M)$ ($M$ $A$-module) : 0, 17.2.1.\par
$\operatorname{dim.coh}(A)$ ($A$ anneau) : 0, 17.2.8.\par
$\operatorname{dim.proj}(\mathscr F)$, $\operatorname{dim.inj}(\mathscr F)$
($\mathscr F$ Module sur un espace annelé) : 0, 17.2.14.\par
$\operatorname{dim.coh}(X)$ ($X$ espace annelé) : 0, 17.2.14.\par
$E\times_B F$ ($B$, $E$, $F$ $A$-anneaux) : 0, 18.1.2.\par
$D_B(L)$ ($B$ anneau, $L$ $B$-bimodule) : 0, 18.2.3.\par
$Y\oplus_X Z$ ($X$, $Y$, $Z$ $A$-bimodules) : 0, 18.2.7.\par
$\operatorname{Exan}_A(B,L)$ ($B$ $A$-anneau, $L$ $B$-bimodule) : 0, 18.3.4.\par
$\operatorname{Exan}_{A/A'}(B,L)$ : 0, 18.3.7.\par
$\operatorname{Exal}_A(B,L)$, $\operatorname{Exal}_{A/A'}(B,L)$ ($A$ anneau commutatif, $B$ $A$-algèbre, $L$ $B$-bimodule) : 0, 18.4.1.\par
$\operatorname{Exalcom}_A(B,L)$, $\operatorname{Exalcom}_{A/A'}(B,L)$ ($A$ anneau commutatif, $B$ $A$-algèbre commutative, $L$ $B$-module) : 0, 18.4.2.\par
$\mathrm H_A^2(B,L)$, $\mathrm H_A^2(B,L)^s$ : 0, 18.4.3.\par
$\operatorname{Exantop}_A(B,L)$, $\operatorname{Exaltop}_A(B,L)$,
$\operatorname{Exalcotop}_A(B,L)$ : 0, 18.5.1.\par
$\operatorname{Exantop}_{A/A'}(B,L)$, $\operatorname{Exaltop}_{A/A'}(B,L)$,
$\operatorname{Exalcotop}_{A/A'}(B,L)$ : 0, 18.5.2.\par
$\operatorname{Hom.\ cont}_C(M,N)$ ($M$, $N$ $C$-bimodules topologiques) : 0, 18.5.3.\par
$\operatorname{D\acute{e}r}_A(B,L)$ ($B$ $A$-anneau, $L$ $B$-bimodule) : 0, 20.1.2.\par
$\operatorname{D\acute{e}r.\ cont}_A(B,L)$ : 0, 20.3.1.\par
$p_{B/A}$, $\mathfrak J_{B/A}$ ($B$ $A$-algèbre topologique commutative) : 0, 20.4.1.\par
$P^1_{B/A}$, $\varepsilon_{B/A}$ ($B$ $A$-algèbre topologique commutative) : 0, 20.4.2.\par
$\Omega^1_{B/A}$, $\widehat{\Omega}^1_{B/A}$ ($B$ $A$-algèbre topologique commutative) : 0, 20.4.3.\par
$\Omega^1_B$ ($B$ anneau topologique) : 0, 20.4.3.\par
$d_{B/A}$, $d(x)$, $dx$, $d_B$ : 0, 20.4.6.\par
$\Upsilon_{C/B/A}$, $\Upsilon_{C/B}$ ($A$, $B$, $C$ anneaux) : 0, 20.6.1.\par
$\chi_E$ ($E$ $B$-extension $A$-triviale de $C$ par $L$) : 0, 20.6.8.\par

\clearpage
\oldpage[IV-1]{253}
$\Upsilon^C_{B/A/\Lambda}$, $\Upsilon^C_{B/A}$ ($\Lambda$, $A$, $B$, $C$ anneaux) : 0, 20.6.14.\par
$\chi_{B/A}$, $\chi_B$ ($B$ $A$-algèbre) : 0, 20.6.24.\par
$\chi_E$ ($E$ $B$-algèbre topologique) : 0, 20.7.10.\par
$F_A$, $A^p$, $A^{(p)}$, $E^{(p)}$ ($A$ anneau de caractéristique $p>0$, $E$ $A$-module) : 0, 21.1.4.\par
$\pi_{B/A}$, $\Theta_{B/A}$, $\Xi_{B/A}$ ($A$, $B$ anneaux de caractéristique $p>0$) : 0, 21.3.2.\par
$\Delta(L/K,k/k_0)$, $d(L/K,k/k_0)$, $\Delta(L/K,k)$, $d(L/K,k)$ ($k_0$, $k$, $K$, $L$ corps) : 0, 21.6.1.\par
$\chi'_{B/A}$ ($A$, $B$ anneaux de caractéristique $p>0$) : 0, 22.4.6.\par
$\mathscr T^{\mathrm{cons}}$, $X^{\mathrm{cons}}$ ($X$ préschéma, $\mathscr T$ topologie de $X$) : IV, 1.9.12.\par
$f^{\mathrm{cons}}$ ($f$ morphisme de préschémas) : IV, 1.9.13.\par

\clearpage
\oldpage[IV-1]{254}
\phantomsection
\addcontentsline{toc}{section}{Index terminologique}
\section*{INDEX TERMINOLOGIQUE}

Algèbre augmentée des parties principales d'ordre 1 : 0, 20.4.2.\par
$A$-algèbre de Cohen : 0, 19.8.1.\par
Algèbre de présentation finie : IV, 1.4.1.\par
$A$-algèbre essentiellement de type fini : IV, 1.3.8.\par
$A$-anneau : 0, 18.1.1.\par
$A$-anneau augmenté sur $B$ : 0, 18.1.4.\par
Anneau biéquidimensionnel : 0, 16.1.4.\par
Anneau caténaire : 0, 16.1.4.\par
Anneau de caractéristique $p$ : 0, 21.1.1.\par
Anneau de Cohen : 0, 19.8.4.\par
Anneau de Cohen-Macaulay : 0, 16.5.1 et 16.5.13.\par
Anneau équicodimensionnel, équidimensionnel : 0, 16.1.4.\par
Anneau local géométriquement unibranche : 0, 23.1.7.\par
Anneau japonais : 0, 23.1.1.\par
Anneau local premier, anneau local complet premier : 0, 19.8.3.\par
Anneau régulier : 0, 17.1.1 et 17.3.6.\par
Anneau local unibranche : 0, 23.1.7.\par
Anneau universellement japonais : 0, 23.1.1.\par
Application ouverte en un point : IV, 1.10.2.\par
Augmentation : 0, 18.1.4.\par

Bimorphisme formel : 0, 19.1.2.\par

Chaîne dans un ensemble ordonné, chaîne de parties fermées irréductibles d'un espace topologique : 0, 14.1.1.\par
Chaîne saturée : 0, 14.3.1.\par
Codimension combinatoire, codimension d'une partie fermée d'un espace topologique : 0, 14.2.1.\par
Codimension en un point d'une partie fermée d'un espace topologique : 0, 14.2.4.\par
Condition des chaînes : 0, 14.3.2.\par
Coprofondeur d'un module sur un anneau local noethérien : 0, 16.4.9.\par
Corps $k_0$-admissible, corps admissible pour une extension : 0, 21.6.1.\par
Corps de multiplicité radicielle finie sur un sous-corps : 0, 19.6.6.\par

Défaut de $k_0$-admissibilité d'une extension : 0, 21.6.1.\par
Dérivation, $A$-dérivation : 0, 20.1.2.\par
Différentielles, 1-différentielles, différentielles absolues : 0, 20.4.3.\par
Différentielle extérieure : 0, 20.4.6.\par
Dimension combinatoire, dimension d'un espace topologique : 0, 14.1.2.\par
Dimension d'un espace topologique en un point : 0, 14.1.2.\par
Dimension de Krull, dimension d'un anneau : 0, 16.1.1.\par
Dimension d'un module : 0, 16.1.7.\par
Dimension cohomologique, dimension cohomologique globale d'un anneau : 0, 17.2.8.\par
Dimension cohomologique d'un espace annelé en un point, dimension cohomologique d'un espace annelé : 0, 17.2.14.\par

\clearpage
\oldpage[IV-1]{255}
Dimension injective, dimension projective d'un module : 0, 17.2.1.\par
Dimension injective en un point, dimension injective d'un faisceau : 0, 17.2.14.\par
Dimension projective en un point, dimension projective d'un faisceau : 0, 17.2.14.\par

Élément de $A$ non diviseur de 0 dans un $A$-module : 0, 15.1.1.\par
Élément $M$-régulier, élément $M$-quasi-régulier : 0, 15.1.4.\par
Ensemble ind-constructible, ensemble pro-constructible : IV, 1.9.4.\par
Épimorphisme formel : 0, 19.1.2.\par
$A$-équivalence de $A$-extensions : 0, 18.2.2.\par
Espace biéquidimensionnel : 0, 14.3.3.\par
Espace caténaire : 0, 14.3.2.\par
Espace équicodimensionnel : 0, 14.2.1.\par
Espace équidimensionnel : 0, 14.1.3.\par
$A$-extension d'un $A$-anneau $B$ par un $B$-bimodule : 0, 18.2.2.\par
Extension déduite d'une autre par un homomorphisme de bimodules : 0, 18.2.8.\par
Extension de Hochschild : 0, 18.4.3.\par
$B$-extension $A$-triviale : 0, 18.3.7.\par
Extension triviale type : 0, 18.2.3.\par
$A$-extensions $A$-équivalentes : 0, 18.2.2.\par

Famille $p$-libre, famille absolument $p$-libre : 0, 21.1.9.\par
Formellement étale (algèbre) : 0, 19.10.2.\par
Formellement inversible à droite (homomorphisme) : 0, 19.1.15.\par
Formellement inversible à gauche (homomorphisme) : 0, 19.1.5.\par
Formellement non ramifiée (algèbre) : 0, 19.10.2.\par
Formellement projectif (module) : 0, 19.2.1.\par
Formellement lisse (algèbre) : 0, 19.3.1.\par
Formellement lisse relativement à un anneau (algèbre) : 0, 19.9.1.\par

Géométriquement régulier sur un corps (anneau local) : 0, 19.6.5.\par
Groupe des classes de $A$-extensions : 0, 18.3.4.\par

Hauteur d'un idéal : 0, 16.1.3.\par
$A$-homomorphisme de $A$-anneaux : 0, 18.1.1.\par
Homomorphisme caractéristique d'une $B$-extension $A$-triviale : 0, 20.6.8.\par
Homomorphisme caractéristique d'une $A$-algèbre relativement à un anneau et à un idéal : 0, 20.6.24.\par
Homomorphisme caractéristique d'une algèbre de caractéristique $p>0$ : 0, 22.4.6.\par

Idéal d'augmentation : 0, 18.1.4.\par
Image réciproque d'un anneau augmenté : 0, 18.1.5.\par
Image réciproque d'une $A$-extension : 0, 18.2.5.\par

Longueur d'une chaîne : 0, 14.1.1.\par

Module biéquidimensionnel, caténaire, équicodimensionnel, équidimensionnel : 0, 16.1.7.\par
Module de Cohen-Macaulay : 0, 16.5.1 et 16.5.13.\par
Module d'imperfection : 0, 20.6.1.\par
Monomorphisme formel : 0, 19.1.2.\par
Morphisme de $A$-extensions : 0, 18.2.4.\par
Morphisme de présentation finie en un point, morphisme localement de présentation finie : IV, 1.4.2.\par
Morphisme de présentation finie : IV, 1.6.1.\par
Morphisme quasi-séparé : IV, 1.2.1.\par

$p$-anneau de Cohen : 0, 19.8.4.\par
$p$-base, $p$-base absolue : 0, 21.1.9.\par
Point maximal d'un préschéma : IV, 1.1.4.\par

\clearpage
\oldpage[IV-1]{256}
Polynôme de Hilbert-Samuel : 0, 16.2.1.\par
Préschéma de présentation finie sur un autre : IV, 1.6.1.\par
Préschéma quasi-séparé : IV, 1.2.1.\par
Produit fibré de $A$-anneaux : 0, 18.1.2.\par
Profondeur d'un module : 0, 16.4.5.\par

Somme amalgamée de bimodules : 0, 18.2.7.\par
Strictement formellement projectif (module) : 0, 19.2.3.\par
Suite $M$-régulière, suite $M$-quasi-régulière : 0, 15.1.7.\par
Suite $\mathscr F$-régulière, suite $\mathscr F$-quasi-régulière : 0, 15.2.2.\par
Système de paramètres pour un module sur un anneau semi-local noethérien : 0, 16.3.6.\par
Système de $p$-générateurs, système absolu de $p$-générateurs : 0, 21.1.9.\par
Système fondamental d'idéaux ouverts, de sous-modules ouverts : 0, 19.0.2.\par
Système régulier de paramètres d'un anneau local régulier : 0, 17.1.6.\par

Topologie constructible sur un préschéma : IV, 1.9.12.\par
Topologie sur un $A$-module déduite de la topologie de $A$ : 0, 19.0.2.\par
Trivial ($A$-anneau augmenté) : 0, 18.1.4.\par
$A$-triviale (extension) : 0, 18.2.3 et 18.3.7.\par

\clearpage
\oldpage[IV-1]{257}
\phantomsection
\addcontentsline{toc}{section}{Table des matières}
\section*{TABLE DES MATIÈRES}

\newcommand{\ivonecontentsline}[3]{\noindent\makebox[17mm][l]{#1}\parbox[t]{0.77\linewidth}{#2\dotfill}\makebox[12mm][r]{#3}\par}
\ivonecontentsline{Chapitre 0.}{\textbf{Préliminaires} (\emph{suite})}{5}
\ivonecontentsline{\S~14.}{Dimension combinatoire d'un espace topologique}{6}
\ivonecontentsline{14.1.}{Dimension combinatoire d'un espace topologique}{6}
\ivonecontentsline{14.2.}{Codimension d'une partie fermée}{8}
\ivonecontentsline{14.3.}{La condition des chaînes}{10}
\ivonecontentsline{\S~15.}{Suites $M$-régulières et suites $\mathscr F$-régulières}{12}
\ivonecontentsline{15.1.}{Suites $M$-régulières et suites $M$-quasi-régulières}{12}
\ivonecontentsline{15.2.}{Suites $\mathscr F$-régulières}{20}
\ivonecontentsline{\S~16.}{Dimension et profondeur dans les anneaux locaux noethériens}{22}
\ivonecontentsline{16.1.}{Dimension d'un anneau}{22}
\ivonecontentsline{16.2.}{Dimension d'un anneau semi-local noethérien}{25}
\ivonecontentsline{16.3.}{Systèmes de paramètres dans un anneau local noethérien}{28}
\ivonecontentsline{16.4.}{Profondeur et coprofondeur}{32}
\ivonecontentsline{16.5.}{Modules de Cohen-Macaulay}{36}
\ivonecontentsline{\S~17.}{Anneaux réguliers}{39}
\ivonecontentsline{17.1.}{Définition des anneaux réguliers}{39}
\ivonecontentsline{17.2.}{Rappels sur la dimension projective et la dimension injective des modules}{42}
\ivonecontentsline{17.3.}{Théorie cohomologique des anneaux réguliers}{46}
\ivonecontentsline{\S~18.}{Compléments sur les extensions d'algèbres}{51}
\ivonecontentsline{18.1.}{Images réciproques d'anneaux augmentés}{51}
\ivonecontentsline{18.2.}{Extensions d'un anneau par un bimodule}{54}
\ivonecontentsline{18.3.}{Le groupe des classes de $A$-extensions}{59}
\ivonecontentsline{18.4.}{Extensions d'algèbres}{64}
\ivonecontentsline{18.5.}{Cas des anneaux topologiques}{66}
\ivonecontentsline{\S~19.}{Algèbres formellement lisses et anneaux de Cohen}{69}
\ivonecontentsline{19.0.}{Introduction}{69}
\ivonecontentsline{19.1.}{Épimorphismes et monomorphismes formels}{71}
\ivonecontentsline{19.2.}{Modules formellement projectifs}{78}

\clearpage
\oldpage[IV-1]{258}
\ivonecontentsline{19.3.}{Algèbres formellement lisses}{79}
\ivonecontentsline{19.4.}{Premiers critères de lissité formelle}{86}
\ivonecontentsline{19.5.}{Lissité formelle et anneaux gradués associés}{90}
\ivonecontentsline{19.6.}{Cas des algèbres sur un corps}{100}
\ivonecontentsline{19.7.}{Cas des homomorphismes locaux : théorèmes d'existence et d'unicité}{104}
\ivonecontentsline{19.8.}{Algèbres de Cohen et $p$-anneaux de Cohen; application à la structure des anneaux locaux complets}{109}
\ivonecontentsline{19.9.}{Algèbres relativement formellement lisses}{114}
\ivonecontentsline{19.10.}{Algèbres formellement non ramifiées et algèbres formellement étales}{115}
\ivonecontentsline{\S~20.}{Dérivations et différentielles}{116}
\ivonecontentsline{20.1.}{Dérivations et extensions d'algèbres}{117}
\ivonecontentsline{20.2.}{Propriétés fonctorielles des dérivations}{119}
\ivonecontentsline{20.3.}{Dérivations continues dans les anneaux topologiques}{121}
\ivonecontentsline{20.4.}{Parties principales et différentielles}{123}
\ivonecontentsline{20.5.}{Propriétés fonctorielles fondamentales de $\Omega^1_{B/A}$}{128}
\ivonecontentsline{20.6.}{Modules d'imperfection et homomorphismes caractéristiques}{136}
\ivonecontentsline{20.7.}{Généralisations aux anneaux topologiques}{147}
\ivonecontentsline{\S~21.}{Différentielles dans les anneaux de caractéristique $p$}{153}
\ivonecontentsline{21.1.}{Systèmes de $p$-générateurs et $p$-bases}{154}
\ivonecontentsline{21.2.}{$p$-bases et lissité formelle}{157}
\ivonecontentsline{21.3.}{$p$-bases et modules d'imperfection}{160}
\ivonecontentsline{21.4.}{Cas des extensions de corps}{162}
\ivonecontentsline{21.5.}{Application : critères de séparabilité}{164}
\ivonecontentsline{21.6.}{Corps admissibles pour une extension}{167}
\ivonecontentsline{21.7.}{L'égalité de Cartier}{169}
\ivonecontentsline{21.8.}{Critères d'admissibilité}{171}
\ivonecontentsline{21.9.}{Modules de différentielles complétés dans les anneaux de séries formelles}{176}
\ivonecontentsline{\S~22.}{Critères différentiels de lissité formelle et de régularité}{182}
\ivonecontentsline{22.1.}{Relèvement de la lissité formelle}{183}
\ivonecontentsline{22.2.}{Caractérisation différentielle des algèbres locales formellement lisses sur un corps}{186}
\ivonecontentsline{22.3.}{Application aux relations entre certains anneaux locaux et leurs complétés}{191}
\ivonecontentsline{22.4.}{Résultats préliminaires sur les extensions finies d'anneaux locaux dont l'idéal maximal est de carré nul}{193}

\clearpage
\oldpage[IV-1]{259}
\ivonecontentsline{22.5.}{Algèbres géométriquement régulières et algèbres formellement lisses}{201}
\ivonecontentsline{22.6.}{Critère jacobien de Zariski}{205}
\ivonecontentsline{22.7.}{Le critère jacobien de Nagata}{209}
\ivonecontentsline{\S~23.}{Anneaux japonais}{213}
\ivonecontentsline{23.1.}{Anneaux japonais}{213}
\ivonecontentsline{23.2.}{Clôture intégrale d'un anneau local noethérien intègre}{217}
\ivonecontentsline{}{\textsc{Chapitre IV.} --- \textbf{Étude locale des schémas et des morphismes de schémas}}{222}
\ivonecontentsline{\S~1.}{Conditions de finitude relatives. Ensembles constructibles dans les préschémas}{224}
\ivonecontentsline{1.1.}{Morphismes quasi-compacts}{224}
\ivonecontentsline{1.2.}{Morphismes quasi-séparés}{226}
\ivonecontentsline{1.3.}{Morphismes localement de type fini}{228}
\ivonecontentsline{1.4.}{Morphismes localement de présentation finie}{230}
\ivonecontentsline{1.5.}{Morphismes de type fini}{233}
\ivonecontentsline{1.6.}{Morphismes de présentation finie}{234}
\ivonecontentsline{1.7.}{Amélioration de résultats antérieurs}{236}
\ivonecontentsline{1.8.}{Morphismes de présentation finie et ensembles constructibles}{238}
\ivonecontentsline{1.9.}{Ensembles pro-constructibles et ind-constructibles}{241}
\ivonecontentsline{1.10.}{Applications aux morphismes ouverts}{249}
\ivonecontentsline{}{Bibliographie (\emph{suite})}{251}
\ivonecontentsline{}{Index des notations}{252}
\ivonecontentsline{}{Index terminologique}{254}

\vspace{2em}
\begin{flushright}
\emph{Reçu le 15 mars 1963.}
\end{flushright}
\let\ivonecontentsline\relax
\endgroup

% FIN SOURCE ega4/ega4-1-fr.tex

\end{document}
